% ---
% title: "Biblia: Begleittexte"
% subtitle: Band 1
% author: Martin Luther und weitere Verfasser
% year: 1860
% language: de
% description: >-
%   Arbeitsausgabe der nichtbiblischen Begleittexte zur Biblia, Band 1
%   (Genesis–Hiob): Vorreden, Widmungen, Glossar, Kapitel-Einleitungen,
%   Summarien, Gebete und erklärende Glossen, ohne den fortlaufenden Bibeltext.
%   Bandweit erschlossener, teilweise scan-geprüfter OCR-Arbeitsentwurf.
% source: St. Louis 1860
% rights: "Gemeinfrei (Public Domain): erschienen 1860."
% draft: true
% ---
\documentclass{oll-book}
\begin{document}
\frontmatter
\begin{titelseite}
\quellseite{8}
\werktitel{BIBLIA,}\\[1em]
das ist\\[0.5em]
die ganze\\[0.5em]
Heilige Schrift\\[0.5em]
des\\[0.5em]
Alten und Neuen Testaments,\\[1em]
verdeutschet durch\\[0.5em]
Dr. Martin Luther,\\[1em]
mit dessen Vorreden und Randglossen, sowie mit den Summarien\\[0.5em]
M. Viti Dieterichs,\\[0.5em]
nebst den Vorreden und Schlußgebeten\\[0.5em]
Francisci Vierlings.\\[1em]
Erster Band.\\[0.5em]
Die fünf Bücher Mosis bis Hiob.\\[1em]
Unveränderter, unter Aufsicht der deutschen evang.-luth. Bibelgesellschaft zu St. Louis, Mo., besorgter Abdruck.\\[1em]
St. Louis, Mo.\\[0.5em]
Verlag von Aug. Wiebusch u. Sohn.\\[0.5em]
1860.
\end{titelseite}
\chapter{Hinweis zum Arbeitsstand}
Diese Arbeitsausgabe erschließt die nichtbiblischen Begleittexte des gesamten
ersten Bandes, von den Vorreden bis zu Hiob: Einleitungen, Kapitelübersichten,
Summarien, Gebete und erklärende Glossen. Der fortlaufende Bibeltext und reine
Querverweise sind ausgelassen; Bibelzitate innerhalb der Begleittexte bleiben
erhalten. Die Auswahl und der Wortlaut beruhen auf OCR und sind noch nicht
anhand sämtlicher Scans geprüft. Auch Auslassungen oder versehentlich
übernommener Bibeltext können trotz automatisierter Prüfungen noch vorkommen.
Diese Ausgabe ist kein scan-geprüfter Lesetext; siehe den Review-Bericht.
\tableofcontents
\mainmatter
\quellseite{10}

\chapter{Widmung an die sächsischen Fürsten}

Denen Durchläuchtigsten, Hochgebornen, Fürsten und Herren, Herrn Friedrichen, Regierenden Landes-Fürsten und Herrn, Herrn Albrechten, Herrn Bernharden, Herrn Heinrichen, Herrn Christianen, Herrn Ernsten, Herrn Johann Ernsten, Sämmtlichen Herren Gebrüdern, Herzogen zu Sachsen, Jülich, Cleve, und Berg, Landgrafen in Thüringen, Markgrafen zu Meißen, Gefürsteten Grafen zu Henneberg, Grafen zu der Mark und Ravensburg, Herren zu Ravenstein Meinen allerseits Gnädigsten Fürsten und Herren.

Durchläuchtigste, Hochgeborne Fürsten, Gnädigste Herren.

Es schreibet der Apostel S. Paulus an die Colosser am 3.: Lasset das Wort Christi reichlich unter euch wohnen in aller Weisheit, lehret und ermahnet euch selbst mit Psalmen und Lobgesängen und geistlichen lieblichen Liedern. Durch das Wort Christi verstehen Lombardus, der Cardinal Cajetanus, die Jesuiten, Cornelius a Lapide, und Eccius, die Lehre, und das Evangelium Christi. Daher auch die Centuriatores Magdeburgenses aus diesen Worten S. Pauli beweisen: die Historia der Evangelisten wäre neben den Büchern Altes Testaments zu der Apostel Zeiten öffentlich, wenn die Gläubigen zusammen kommen, vorgelesen worden. Centur. 1. lib. 2. cap. 6, fol. 494. Der Kirchen-Lehrer Chrysostomus und Theophylactus verstehen durch das Wort Christi, die heilige Schrift, wie bald soll angeführet werden. Diese Lehre nun von Christo, desselben Gnaden-Predigt, oder die heilige Schrift, soll unter denen zu Colossen wohnen. Das ist, diese Lehre soll nicht wie ein Gast mit Freuden angenommen, und hernach bald wieder weggelassen, oder, wie es bei dem ein wenig zu lange bleibenden Gaste meist zu geschehen pfleget, gar verachtet werden, sondern sie soll wohnen, das ist, die Lehre von Christo, die heilige Schrift, soll mit wahrem Glauben ergriffen, und beständig im Herzen bewahret werden, damit wir eine solche Freundschaft mit derselben machen, und sie uns so gemein und bekannt werde, als diejenigen, die bei uns in einem Hause wohnen, uns pflegen gar gemein und kundig zu werden. Und dieses soll geschehen reichlich, das ist, überflüssig und häufig.

Daß also die Art und Weise mit der Lehre Christi, oder der heiligen Schrift umzugehen, hiermit wird angedeutet. Der große majestätische Herr, der Herr der Heerschaaren, befiehlet dem Helden Josua ernstlich: Laß das Buch dieses Gesetzes nicht von deinem Munde kommen, sondern betrachte es Tag und Nacht, auf daß du haltest und thust aller Dinge nach dem, das drinnen geschrieben stehet; im Buch Josua am ersten Kap. v. 8., und vorher sagte er zu allen Israeliten: Diese Worte, die ich dir heute gebiete, sollst du zu Herzen nehmen, und sollst sie deinen Kindern schärfen, und davon reden, wenn du in deinem Hause sitzest, oder auf dem Wege gehest, wenn du dich niederlegest, oder aufstehest, 5 Buch Mosis am 6. Kap. v. 6. König David preiset denjenigen selig, welcher zum Gesetz, zur Lehre des Herrn Lust hat, und redet von seinem Gesetz Tag und Nacht, Psalm 1, v. 2. Unser Heiland Christus Jesus saget zu den Juden: Suchet in der Schrift,

\quellseite{11}

denn ihr meinet, ihr habt das ewige Leben drinnen, und sie ist, die von mir zeuget, Joh. 5, v. 39. Chrysostomus erkläret diese Worte reichlich also: nicht oben hin, sondern mit großem Fleiße.

Was nun Gott selber nennet: Das Gesetz nicht lassen vom Munde kommen; es Tag und Nacht betrachten; zu Herzen nehmen; es den Kindern schärfen; was König David heißt: zum Gesetz des Herrn Lust haben, davon reden Tag und Nacht. Was unser Heiland selber ausspricht mit: Suchen in der Schrift: das redet S. Paulus aus mit: Das Wort Christi reichlich unter sich wohnen lassen. Es wird aber nicht allein mit dem Worte reichlich gesehen auf die Art und Weise, wie man die heilige Schrift solle treiben; sondern auch auf den erbaulichen Nutzen, daß man sie also solle handeln, betrachten und damit umgehen, daß wir reich werden am Glauben und an der Liebe, auch in allen widrigen Begebenissen einen Trost nach dem andern zu unserm, und anderer Leute Besten können alsobald bei der Hand haben. Und dieses reichliche, überflüssige und zu unserer eigenen, auch anderer Leute Erbauung vorgenommenes Treiben der heiligen Schrift soll geschehen unter den Colossern: Unter euch, schreibet S. Paulus, ohne Unterschied des Standes, oder des Geschlechts, es mögen sein Lehrer, oder Zuhörer, Weiber oder Männer. Wie denn der Autor der Erklärungen, welche von etlichen Hieronymo, von andern Primasio, oder sonst jemandem zugereignet werden, diese Worte ausdrücklich von den Laien ausleget, wenn er schreibet: Hic ostenditur, verbum Christi non solum sufficienter, sed etiam abundanter habere debere etiam laicos et invicem semetipsos monere et docere. Das ist: Hier wird gewiesen, daß auch die Laien das Wort Christi nicht allein zur Bedürfniß und Genüge, sondern auch zum Ueberflusse sollen haben, und einer den andern erinnern und unterweisen. Das Chrysostomus unter den Worten: Unter euch, auch die γνωστικος, und βιωτικος, seculares, welche mit zeitlichen Sachen umgehen, in einem weltlichen Stande leben, welche Weib und Kinder haben, begreife, werden wir bald hören.

Dieses reichliche, überflüssige und fleißige Treiben der heiligen Schrift soll ferner geschehen: in aller Weisheit. Der bekannte Schullehrer Thomas, welchen die Römisch-Katholischen insgemein den englischen Lehrer nennen, vermeinet, S. Paulus weise hiermit die Gläubigen zu Colossen auf den rechten Ursprung, woraus alle Weisheit entspringe, und auch auf den Nutzen, welcher von dieser Weisheit zu hoffen, wenn er setzet: Cum dicit: Verbum Christi, monet ad Sapientiam; et primo docet Sapientiae originem, secundo sapientiae usum: ibi, docentes: Das ist: wenn er, Paulus, saget: Das Wort Christi, ermahnet er zur Weisheit; und lehret erstlich der Weisheit Ursprung, zum andern der Weisheit Nutzbarkeit, mit dem Worte: lehret. Wollte man dieses hinzusetzen, daß S. Paulus mit diesen Worten: in aller Weisheit, auch ein Absehen hätte auf die Weise, wie man mit Lesung, Betrachtung und Anführung der heiligen Schrift sollte umgehen, nämlich: behutsam, daß man dasjenige, was vorhergehe, was folge und dergleichen, wohl in Acht nehme; auch zugleich ein Ziel wollte setzen, wie weit man mit Betrachtung der heiligen Schrift gehen sollte, wird es der Aehnlichkeit des Glaubens nicht zuwider laufen. Denn es muß doch mit dem Nachsuchen in der heiligen Schrift Maß gehalten werden, daß wir dieselbe nicht nach unserm Verstande und Gutdünken einrichten, und das öfters verdammliche Quomodo? Wie? damit viele Gottes Allmacht, Weisheit und Wahrheit in Zweifel ziehen, von uns gehöret oder angemerket werden möge.

Endlich wird uns auch vorgestellet die End-Ursache, warum man Gottes Wort, die heilige Schrift, so fleißig treiben, andächtig betrachten und ihm so bekannt und läuffig machen solle. Lehret euch selbst, saget S. Paulus, das ist, unterweiset, unterrichtet euch mit einander. Und auf solche Weise werdet ihr in den Glaubens-Artikeln gegründet und bestätiget werden. Vermahnet euch selbst, schreibet S. Paulus weiter. Das ist: Ist jemand mit einem Fehl übereilet worden, ist er in eine Sünde gefallen, so helfet ihr, die ihr geistlich seid, und nicht aus dem Stande der Gnaden euch durch muthwillige Sünden bringen lassen, ihm wieder zurechte mit sanftmüthigem Geiste, nach S. Pauli Ermahnung in dem 6. Kap. der Epistel an die Galater. Oder wie S. Paulus im 10. Kap. an

\quellseite{12}

die Brüder v. 24. es erkläret: Lasset uns unter einander unser selbst wahrnehmen mit Reizen zur Liebe und guten Werken. Der Kirchenlehrer Chrysostomus, der um das vierhundertste Jahr nach Christi Geburt gelebet, leget vor angeführte Worte S. Pauli, wie sie erkläret worden, auch also aus, welche ich nur in unserer bekannten Muttersprache will anführen. Denn also redet er in der neunten Predigt über das dritte Capitel der Epistel an die Colosser seine Zuhörer zu Constantinopel öffentlich an: Höret zu, so viel euer Laien (βιωτικῶν, victum quaerentes) sein, und dem Weibe und Kindern vorstehet, wie auch euch (Paulus) beschielet, daß ihr vor allem die Schrift leset; und zwar nicht schlecht und oben hin, sondern mit großem Fleiße. Und bald hierauf: Höret alle ihr Laien (βιωτικῶν, victum quaerentes), ich ermahne euch, schaffet euch die Bibel, welche ist eine Arznei der Seelen. Wenn ihr nichts anderes wollet, so schaffet nur das Neue Testament, die Apostelgeschichte, die Evangelia, eure stetigen Lehrmeister. Wenn eine Traurigkeit einfället, so mache dich hierzu, als zu einer Apotheke der Arzneien. Daraus nimm dir einen Trost wider den Unfall. Es mag sich entweder ein Schaden zugetragen haben, oder ein Todesfall, oder ein Verlust der Unsrigen. Ja, mache dich nur nicht allein hierzu, sondern nimm alles an. Behalte es in deinem Verstande und Gedächtniß. Dieses ist die Ursache alles Uebels, daß man die Schrift nicht weiß. Soweit Chrysostomus.

Es wird aber daraus dieses unstreitig geschlossen: Sollen die Laien, diejenigen, welche Weiber und Kinder haben, vor allem die Schrift lesen, sie kaufen, und sich daraus wider alle Unfälle trösten, so muß ja folgen, daß man den Laien die Bibel in ihrer Muttersprache lassen, und sie zur fleißigen Lesung derselben anmahnen müsse und solle. Weil ihnen alles Uebels Ursache ist, die Schrift nicht wissen. Es bemühet sich zwar Joh. Paludanus in seinen also genannten Vindiciis Theologicis adversus verbi Dei Corruptelas part. alter. pag. 276. aus dem Stapletono hierauf zu antworten, wenn er schreibet: Der Apostel lehret deutlich, nicht daß die Schrift solle gelesen werden, sondern daß die Predigt des Worts, die ihnen geschehen, solle wiederholet und erneuert werden, indem sie sich unter einander ermahnen mit Psalmen etc. Denn dieses folget daselbst: Singet dem HErrn in euren Herzen. Dieses aber alles wird deutlich entweder auf die geistlichen Unterredungen unter einander mit den Brüdern, oder auf die besondern, oder Privatandachten gezogen. So weit Paludanus.

Aber vergeblich. Denn S. Paulus setzet zweierlei, welches die Gläubigen zu Colossen thun sollen, wie auch Thomas der Schullehrer angemerkt, und es droben schon angeführet, wenn Paulus spricht: Das Wort Christi, so ermahnet er zur Weisheit, und lehret erstlich den Ursprung der Weisheit, hernach auch den Nutz der Weisheit, in dem Worte „lehre.“ Aber Paludanus menget dieses beides unter einander, als ob es einerlei wäre, das Wort Christi reichlich unter uns wohnen lassen, und dasselbe lehren, oder daraus ermahnen. Die doch also unterschieden sind, als dasjenige, welches etwas wirket, und dasjenige, das da gewirket wird. Denn bei wem das Wort Christi reichlich wohnet, der kann auch andere daraus lehren und ermahnen. Es wird auch Paladano nicht nachgegeben, daß S. Paulus nicht rede von Lesung der Schrift, sondern nur von Wiederholung der Predigt des Worts. Denn wie kann doch das Wort Christi reichlich unter den Colossern wohnen, wenn sie die heilige Schrift nicht haben, und dieselbe nicht lesen? Die Wiederholung der Predigt des Worts gehöret auch mit zum reichlichen Wohnen des Worts Christi unter ihnen; aber es bestehet dieses reichliche Wohnen nicht allein in der Wiederholung des gepredigten Worts, sondern man muß auch die heilige Schrift fleißig und reichlich treiben. Welches nicht geschehen kann, wenn man sie, wie Chrysostomus befiehlt, nicht besitzet und sie nicht unsern stetigen Lehrmeister sein läßt. Welches Lehramt aber die heilige Schrift in einer fremden und unbekannten Sprache nicht kann verrichten, sondern in einer bekannten und auch den Laien läufigen Sprache. Dieses aber kann man Paladano willig nachlassen, daß S. Paulus unter dem reichlichen Wohnen des Worts Christi auch

\quellseite{13}

verstehe die geistlichen Unterredungen unter einander mit den Brüdern, die nach Chrysostomi Meinung sind κοσμικοί, βιωτικοί, seculares, laici, victui quaerendo intenti, Laien, die mit weltlichen Geschäften umgehen, und sich bemühen, wie sie sich und die Ihrigen mit gutem Gewissen hinbringen, oder auch Privatandachten, wenn sie nur aus der heiligen Schrift genommen. Denn solches zu unserm obigen Schlusse genug ist. Der Jesuit Justinianus will der Protestirenden Schluß unvermerket ablehnen, wenn er schreibet über unsern Text: Lasset das Wort Christi etc. Daß nicht jemand meine (welches die Novatores [die Protestirenden damit meinen] unbesonnen erdichten) es müsse die heilige Schrift, als leichte und einem Jeden offenstehend, von eines Jeden Hand angetastet werden, auch nicht einmal die Schuhflicker und Maurer ausgenommen, oder auch nicht, welches gar schön wäre, die Weiberchen: so sehet S. Paulus recht darauf: Lehret und vermahnet euch. Als wollte er sagen: einer unterrichte den andern. Welche aber unser Dolmetscher genennet hat: Vosmet ipsos; hat der Apostel griechisch genennet ἑαυτούς, das ist, euch unter einander, daß nämlich die Verständigeren die Unverständigeren unterweisen sollen. So weit Justinianus.

Aber dieses ist ein wunderlicher Schluß: Einer unterweise den andern. Derowegen soll die heilige Schrift nicht von eines Jeden Hand angetastet werden. Vielmehr gilt der Gegensatz: Einer unterweise den andern. Derowegen soll die heilige Schrift von einem Jeden angegriffen und gelesen werden. Der Schluß ist dieser: welche einander aus der heiligen Schrift sollen unterrichten, denen muß die heilige Schrift zu lesen nachgegeben werden.

Nun aber sollen die κοσμικοί und βιωτικοί, die seculares, welche ihre Nahrung suchen, welche Weiber und Kinder haben, (worunter auch die Gerber, Schuhflicker, Maurer und dergleichen Handwerker ohne einigen Zweifel gehören) einander aus heiliger Schrift unterrichten, derowegen muß diesen benannten Leuten und ihresgleichen nachgelassen werden die heilige Schrift zu lesen.

Der Vorspruch kann nicht geleugnet werden, als von demjenigen, der da wollte befehlen, man sollte mit dem Messer Brod schneiden, und zugleich verbieten, man sollte das Messer, als ein Instrument, dazu nicht brauchen. Man sollte aus einem Brunnen Wasser schöpfen, aber sich hierzu keines Eimers oder einigen anderen Instrumentes bedienen. Denn die heilige Schrift, die da von Gott gegeben, ist nütz zur Lehre, zur Strafe, oder Widerlegung, zur Besserung, zur Züchtigung, oder Unterweisung in der Gerechtigkeit, wie S. Paulus schreibet 2 Tim. 3, v. 16. Wer nun befiehlet, man solle aus der heiligen Schrift unterrichten, aber verbeut die heilige Schrift zu lesen, der befiehlet eine Verrichtung, und untersaget das ordentliche hierzu von GOtt selbst anbefohlene Mittel, die heilige Schrift, die uns doch kann unterweisen zur Seligkeit, durch den Glauben an Christo Jesu, und ist nütze zur Lehre und Unterrichtung. Wer aber also befiehlet und also verbeut, dem ist es mit dem Befehl kein Ernst, sondern spottet nur in dieser höchst wichtigen Sache, welche der Menschen Seligkeit angehet, der elenden Leute noch dazu.

Der Nachspruch ist enthalten in S. Pauli Worten: Lasset das Wort Christi reichlich unter euch wohnen, lehret und vermahnet euch selbst unter einander, und hat denselben der treffliche Lehrer Chrysostomus von denen κοσμικοῖς und βιωτικοῖς verstanden, wie dessen Worte kurz vorher angeführet. Derowegen bleibet auch der Schluß unumgestoßen: Dergleichen Leuten ist nachzulassen, die heilige Schrift in ihrer bekannten Muttersprache zu lesen.

Daß aber Justinianus ferner schreibet: Die Verständigen sollen die Unverständigen unterrichten, daraus folget nicht, daß die heilige Schrift nicht allen zu lesen solle verstattet werden. Denn auch unter den von Justiniano angeführten Schuhflickern und Maurern einer in der heiligen Schrift mehr erfahren als der andere, und also nach Justiniani eigenen Worten, der Verständigere den Unverständigern soll unterrichten. Thomas Stapletonus antwortet auf die Worte Chrysostomi: Ipsum moraliter, non dogmatice haec verba Pauli ad Scripturarum Lectionem traxisse. Aber dieser Unterschied inter moraliter et dogmatice, wie ihn Stapletonus machet, hat keinen Grund in der Theologie. Denn an den-

\quellseite{14}

Lehren hängen gleichsam die Anmahnungen, und folgen aus denselben. Denn wenn die Lehre von GOtt, daß er sei allmächtig, vorgetragen, so kann man daraus ermahnen, ein Jeder wolle in allen Nöthen sein Vertrauen auf diesen allmächtigen GOtt setzen. Wenn die Lehre aus der heiligen Schrift bewiesen, daß GOtt alle Sünde wolle allen bußfertigen Sündern vergeben, so kann man daraus eine Anmahnung ziehen: Alle bußfertigen Sünder wollten bei GOtt Vergebung aller ihrer Sünden suchen, und sich derer gewiß versichern. Und dieses gehet an in allen credendis, in allem, was uns wird vorgeleget aus heiliger Schrift zu glauben. Eben diese Bewandtnis hat es auch mit den faciendis, was GOtt wolle, daß wir thun sollen. Als: wenn der Befehl des wohlthätigen und gütigen GOttes aus dem 58. Kap. Esaia v. 7. kund gemacht: Brich dem Hungrigen dein Brod, und die, so im Elend sind, führe ins Haus. So du einen nacket siehest, so kleide ihn, und entzeuch dich nicht von deinem Fleisch; so folget ohne einigen Zwang die Ermahnung: Man wolle sich der Elenden, Hungrigen und Nackten erbarmen.

Ist derowegen kein Gegensatz zuzulassen, inter moraliter et dogmatice etwas vortragen. Denn weil S. Pauli Lehre ist: Lasset das Wort Christi unter euch reichlich wohnen, so nimmt Chrysostomus aus dieser Lehre die Anmahnung: Ihr Laien, die ihr Weiber und Kinder habt, kaufet euch die Schrift, denn sie ist eurer Seelen Arznei. Nehmet aus derselben wider alle Unfälle Tröstungen, leset die Schrift nicht oben hin, sondern mit großem Fleiße. Denn es ist alles Unglücks eine Ursache, die heilige Schrift nicht wissen. Ja diese Anmahnung ist in obiger Lehre Pauli enthalten und begriffen, daß also fast gar kein Unterschied zu machen in dem Sensu, und in der Meinung, zwischen Pauli Worten: Lasset das Wort Christi unter euch reichlich wohnen; und Chrysostomi Worten: Leset die Schrift nicht oben hin, sondern mit großem Fleiße.

Hat also Chrysostomus aus der Lehre Pauli diese Anmahnung rechtmäßig genommen, und bleibt demnach wahr, was droben geschlossen: Sollen die Laien die Bibel schaffen, und mit großem Fleiße lesen, so muß man ihnen dieselbe in ihrer Muttersprache lassen, und sie anmahnen, solche fleißig zu lesen, und in allen widrigen Begegnissen Tröstungen daraus zu nehmen. Kann man derowegen ohne großen Schrecken, ohne Besorgung göttlichen schweren Zorns, und ohne sonderbare Betrübniß nicht lesen, daß die römischen Bischöfe, welche sich doch vor Christi Statthalter ausgegeben, Pius der Vierte, Sixtus der Fünfte und Clemens der Achte, dieser Namen, sich geweigert, die Sünde zu erlassen denen Laien, welche die Bibel würden lesen. Ja, sie haben den Pfarrern, Beichtvätern und Ordinarien verboten, sie sollten keinen Laien die Absolution sprechen und die Sünde vergeben, wenn sie nicht zuvor die Bibel ausgehändigt. Wie unglücklich aber sind diejenigen Gemeinen, in welchen die Lesung der Bibel den Laien bei Strafe der vorbehaltenen Vergebung der Sünden wird verboten; so glücklich sind die Gemeinen, unter welchen das Wort Christi reichlich wohnet, beides, beim öffentlichen Gottesdienste, wie auch daheim in ihren Häusern, da sie es nicht allein frei und sicher dürfen lesen, sondern auch von ihren Seelsorgern und Lehrmeistern zum freundlichsten dazu angemahnet, und ihnen die großen Nutzbarkeiten, die ihnen daraus zuwachsen, vor Augen gestellet werden.

Dergleichen Glückseligkeit haben auch Durchlauchtigste Fürsten, Gnädigste Herren, E. E. E. E. E. E. Fürstl. Durchl. sämmtliche Fürstenthümer und Länder, daß, GOtt sei ewig dafür gedanket, in Städten und Dörfern das Wort Christi reichlich unter derselben Unterthanen wohnet, die Jugend zur Erlernung des Katechismi, der Psalmen und biblischen Sprüche, wie auch zur Lesung der heiligen Schrift, in den Schulen aller Orten fleißig angeführet, auch die Erwachsenen, daß das Gelernte nicht wieder ausfalle, zur Wiederholung werden angehalten.

Wie nun der auch weiland Durchlauchtigste Fürst und Herr, Herr Ernst, Herzog zu Sachsen, Jülich, Cleve und Berg etc. Euer Fürstl. Durchl. Hochgeehrtester Herr Vater, unsers mehr Christseeligen Andenkens, mein weiland Gnädigster Fürst und Herr, Seine herzliche Freude daran hatte, wenn Ihre Durchl. bei deren Unterthanen die wahre Gottesfurcht konnte befördern; also bezeugeten Hochgedachte Fürstl. Durchl. auch Ihr Gnädigstes Gefallen daran, da Anno 1673 im Juni zu Gotha auf dem Friedenstein, dem Fürstl. Residenz-Hause, unterthänigst ich mein Vorhaben mit Herausgebung einer Bibel eröffnete, indem Sie Ihre Gnädigste Einwilligung

\quellseite{15}

dazu in allen Gnaden gaben, und solche auch durch Dero daselbst wohlbestelltes Consistorium schriftlich auszufertigen Gnädigst befohlen. Welches auch geschehen.

Solche Bibel nun, welche durch hiesiges Waisenhauses Kosten zur Erziehung der Armen darin sich befindlichen Waisen aufgelegt worden, habe E. E. E. E. E. E. Fürstl. Durchlaucht unterthänigst überreichen, und zugleich meinen unterthänigsten Glückwunsch wegen ererbter Fürstenthümer und Landen Altenburgischen Theils, wie auch zur angetretenen Fürstlichen Regierung des Höchsten milden Segen in aller Unterthänigkeit öffentlich wiederholen wollen, meinen GOtt unaufhörlich anflehend, daß Er das Wort Christi reichlich in deren wohllöblichen Fürstenthümern und Landen ferner, und bis an den jüngsten Tag wolle lassen wohnen, und schnell laufen, Euer Fürstl. Durchl. bei glücklicher Regierung und allem Fürstlichen hohen Wohlwesen erhalten, derselben Freunde Freund, und Ihrer Feinde Feind sein und bleiben, auch niemals aus diesem Durchlauchtigsten Hause ermangeln lassen an einem Manne nach seinem Herzen, welcher vor Ihm aus und eingehe!

Welches alles daß es geschehen möge, wünschet von Herzen

\begin{unterschrift}

E. E. E. E. E. E. Fürstl. Durchl.\\Altenburg, den 30. März 1676.\\Unterthänigster Diener und Vorbitter bei GOtt\\Joh. Christfried Sagittarius, Dr.,\\Superint. Gener., Ober-Hofprediger und des Consistorii Assessor.

\end{unterschrift}

\quellseite{16}

\chapter{An den Gottseeligen Leser}

Bei Auflegung dieser Bibel, die mit Herrn \unsicher{Miti} Dieterichs Summarien, wie auch Herrn Francisci Bierlings Vorreden und Beschlüssen (dergleichen sonder Ruhm, noch nicht herauskommen, und also keiner, er sei auch wo oder wer er wolle, mit Recht sich kann beklagen, daß ihm eine einzige Bibel nachgedruckt worden, auch die Begierde, Gottes Ehre zu befördern, sowohl der lieben Stadt Altenburg, als andern evangelischen Städten, nicht unbillig wird zu gönnen sein), habe ich\footnote{1. Vielerley editiones in allerhand Formen mit einander conferiret, aber bei vorgefallenen Zweifeln mich nach derjenigen, welche zu Wittenberg im Jahr Christi 1556 in Folio durch Hans Lufften gedruckt, meistentheils gerichtet.\par 2. Habe ich bei Herrn Dieterichs Summarien dieselbe edition, welche im Jahr Christi 1585 zu Nürnberg in Folio allein, und ohne die biblischen Capitel, in den Druck gegeben, zum Grunde gelegt, und die andern hernach mit diesen Summarien herausgekommenen Bibeln gegen obige gehalten, und mit Verwunderung angemercket, daß in vielen Stellen des wohlverdienten Mannes Worte ausgelassen, und in einem andern, auch öfters widrigen Sinne, befunden worden, daß man mit allem Fug des theuren Lutheri Worte, welche in der Warnung vor der Bibel zu finden, auch allhier kann gebrauchen. Es ist mir oft widerfahren, daß ich der Nachdrucker Druck gelesen, also verfälschet gefunden, daß ich meine eigene Arbeit an vielen Orten nicht erkennet. Wer es nicht will glauben, der lasse die editiones gegen einander halten, da wird es mehr als ihm lieb ist, also befinden.\par 3. Sind die Vorreden und Beschlüsse Herrn Bierlings, welche Anno 1596 zu Breslau, in meiner liebsten Geburtsstadt, der ich alles Gute von Herzen wünsche, in Folio absonderlich herausgegeben, von mir durchgelesen, die Orthographie in Acht genommen, auch die Allegata aus der Bibel aufgeschlagen, und mit Versen an den meisten Orten versehen worden.}

Die Druckfehler, welche über allen angewandten Fleiß mit eingeschlichen, dem Sensui aber nichts benehmen, wolle der christliche und geneigte Leser entweder selbst corrigiren oder doch ins besten vermerken, auch diese Arbeit zu seiner ferneren Erbauung selig gebrauchen, und also das Wort Christi reichlich unter und bei ihm wohnen lassen. Ich begehre für alle aufgewandte Mühe, und darüber mir zugestoßenen Angelegenheiten nicht einen Pfennig, sondern Gott allein soll die Ehre, der christlichen Kirchen die Erbauung, und den hülflosen Waisen allhier der Nutz sein und bleiben. Hiermit sei ein Jeder der göttlichen Gnade zu zeitlicher und ewiger Wohlfahrt treulich befohlen. Altenburg, den 3. April Anno 1676.

\quellseite{17}

\chapter{Vorrede auf das Alte Testament}

Das Alte Testament halten etliche gering, als das dem jüdischen Volke allein gegeben, und nun fort aus sei, und nur von vergangenen Geschichten schreibe, meinen, sie haben genug am Neuen Testament, und geben für, eitel geistliche Sinn im alten Testament zu suchen. Wie auch Origenes, Hieronymus und viel hoher Leute mehr gehalten haben. Aber Christus spricht Joh. 5, 39.: Forschet in der Schrift, denn dieselbige gibt Zeugniß von mir. Und S. Paulus gebeut Timotheo 1 Tim. 4, 13.: Er solle anhalten mit Lesen der Schrift. Und rühmet Röm. 1, 2.: wie das Evangelium sei von Gott in der Schrift verheißen. Und 1 Cor. 15, 3. sagt er: Christus sei, nach Laut der Schrift, von Davids Geblüte kommen, gestorben und vom Tode auferstanden. So weiset uns auch S. Petrus mehr denn einmal hinter in die Schrift, 2 Petr. 1, 19.; 3, 15.

Damit sie uns je lehren, die Schrift des Alten Testaments nicht zu verachten, sondern mit allem Fleiß zu lesen, weil sie selbst das Neue Testament so mächtiglich gründen, und bewähren durchs Alte Testament, und sich darauf berufen. Wie S. Lucas, Apg. 17, 11. schreibet, daß die zu Thessalonich täglich forscheten die Schrift, ob sichs so hielte, wie Paulus lehrte. So wenig nun des Neuen Testaments Grund und Beweisung zu verachten ist, so theuer ist auch das Alte Testament zu achten. Und was ist das Neue Testament anders, denn eine öffentliche Predigt und Verkündigung von Christo, durch die Sprüche im Alten Testament gesetzet, und durch Christum erfüllet?

Daß aber diejenigen, so es nicht besser wissen, eine Anleitung und Unterricht haben, nützlich darinnen zu lesen, habe ich diese Vorrede nach meinem Vermögen, so viel mir Gott gegeben, gestellet. Bitte und warne treulich einen jeglichen frommen Christen, daß er sich nicht stoße an der einfältigen Rede und Geschichte, so ihm oft begegnen wird, sondern zweifle nicht daran, wie schlecht es immer sich ansehen läßt, es seien eitel Wort, Werk, Gericht und Geschicht der hohen göttlichen Majestät, Macht und Weisheit. Denn dies ist die Schrift, die alle Weisen und Klugen zu Narren macht, und allein den Kleinen und Albernen offen stehet, wie Christus sagt Matth. 11, 25.: Darum laß deinen Dünkel und Fühlen fahren, und halte von dieser Schrift, als von dem allerhöchsten, edelsten Heiligthum, als von der allerreichsten Fundgrube, die nimmer genug ausgegründet werden mag, auf daß du die göttliche Weisheit finden mögest, welche Gott hie so albern und schlecht fürgeleget, daß er allen Hochmuth dämpfe. Hier wirst du die Windeln und die Krippe finden, da Christus innen liegt, dahin auch der Engel die Hirten weiset, Luc. 2, 12. Schlechte und geringe Windeln sind es, aber theuer ist der Schatz, Christus, der drinnen liegt.

So wisse nun, daß dies Buch ein Gesetzbuch ist, das da lehret, was man thun und lassen soll, und daneben anzeiget Exempel und Geschichte, wie solche Gesetze gehalten oder übertreten sind. Gleich wie das Neue Testament ein Evangelium oder Gnaden-Buch ist, und lehret, wo mans nehmen soll, daß das Gesetz erfüllet werde. Aber gleich wie im Neuen Testament, neben der Gnaden-Lehre, auch viel andere Lehren gegeben werden, die da Gesetz und Gebot sind, das Fleisch zu regieren, sintemal in diesem Leben der Geist nicht vollkommen wird, noch eitel Gnade regieren kann. Also sind auch im Alten Testament, neben den Gesetzen, etliche Verheißung und Gnaden-Sprüche, damit die heiligen Väter und Propheten unter dem Gesetz im Glauben Christi, wie wir, erhalten sind. Doch wie des Neuen Testaments eigentlich Haupt-Lehre ist, Gnade und Friede durch Vergebung der Sünden in Christo verkündigen, also ist des Alten Testament eigentlich Haupt-Lehre, Gesetz lehren, und Sünde anzeigen, und Gutes fordern, Solches wisse im Alten Testament zu warten.

\quellseite{18}

Und daß wir zuerst auf Moses Bücher kommen, der lehret in seinem ersten Buch, wie alle Creaturen geschaffen sind, und, daß seines Schreibens meiste Ursache ist, wo die Sünde und der Tod herkommen sey, nämlich, durch Adams Fall, aus des Teufels Bosheit, 1 Mos. 3, 4. Aber bald darauf, ehe denn Moses Gesetz kömmt, lehret er, woher die Hülfe wiederkommen sollt, die Sünde und Tod zu vertreiben, nämlich, nicht durch Gesetz noch eigen Werk, weil noch kein Gesetz war, sondern durch des Weibes Samen, Christum, Adam, 1 Mos. 3, 15. und Abraham, 1 Mos. 12, 3. verheißen. Auf daß also der Glaube von Anfang der Schrift durch und durch gepreiset werde über alle Werke, Gesetz und Verdienst. Also hat das erste Buch Mose fast eitel Exempel des Glaubens und Unglaubens, und was Glaube und Unglaube für Früchte tragen, und ist fast ein evangelisch Buch.

Danach im andern Buch, da die Welt nun voll in der Blindheit versunken war, daß man schier nicht mehr wußte, was Sünde war, oder wo der Tod herkommen sey, bringet Gott Mosen herfür mit dem Gesetz, und nimmt ein besonders Volk an, die Welt an ihnen wieder zu erleuchten, und durchs Gesetz die Sünde wieder zu eröffnen. Und verfaßt also das Volk mit allerlei Gesetzen, und sondert sie von allen andern Völkern, lässet sie eine Hütten bauen, und richtet einen Gottesdienst an, bestellet Fürsten und Amtleute, und versorget also sein Volk, beide mit Gesetzen und Leuten, aufs allerfeinste, wie sie beide leiblich für der Welt, und geistlich für Gott sollen regieret werden.

Im dritten Buch wird insonderheit das Priesterthum verordnet mit seinen Gesetzen und Rechten, darnach die Priester thun, und das Volk lehren sollen. Da stehet man, wie ein priesterlich Amt nur um der Sünde willen wird eingesetzt, daß es dieselbige soll dem Volk kund machen, und für Gott versöhnen. Also, daß alle sein Werk ist mit Sünden und Sündern umgehen. Derhalben auch den Priestern kein zeitlich Gut gegeben, noch leiblich zu regieren befohlen oder zugelassen wird, sondern allein des Volkes zu pflegen in den Sünden, ihnen zugeeignet wird.

Im vierten Buch, da nun die Gesetze gegeben, Priester und Fürsten eingesetzt sind, die Hütten und Gottesdienst angerichtet, und alles bereit ist, was zum Volk Gottes gehört, hebt sich das Werk und Uebung an, und wird versucht, wie solche Ordnung gehen und sich schicken will. Darum schreibet dasselbige Buch von so viel Ungehorsam und Plagen des Volkes. Und werden etliche Gesetze verkläret und gemehret. Denn also findet sichs allezeit, daß Gesetze bald zu geben sind; aber wenn sie sollen angehen, und in den Schwang kommen, da begegnet nicht mehr denn eitel Hinderniß, und will niemand fort, wie das Gesetz fordert, daß dies Buch ein merklich Exempel ist, wie gar es nichts ist, mit Gesetzen die Leute fromm machen, sondern wie S. Paulus sagt: Daß Gesetz nur Sünde und Zorn anrichte, Röm. 4, 15.

Im fünften Buch, da nun das Volk um seinen Ungehorsam gestraft ist, und GOtt sie mit Gnaden ein wenig gelocket hatte, daß sie aus Wohlthat, da er ihnen die zwei Königreiche gab, beweget wurden, seine Gesetze mit Lust und Liebe zu halten, wiederholt Mose das ganze Gesetz mit allen Geschichten, so ihnen begegnet war (ohne was das Priesterthum betrifft), und erkläret also von neuem an alles, was beide zum leiblichen, und geistlichen Regiment eines Volks gehöret. Daß also Mose, wie ein vollkommener Gesetzlehrer, allenthalben seinem Amt genug thäte, und das Gesetz nicht allein gäbe, sondern auch dabei wäre, da mans thun sollt, und wo es fehlet, verkläret, und wieder anrichtet. Aber diese Verklärung im fünften Buch hält eigentlich nichts anders innen, denn den Glauben an GOtt, und Liebe zum Nächsten. Denn dahin gelangen alle Gesetze Gottes. Darum wehret Mose mit seinem Erklären alle dem, daß den Glauben an GOtt verderben mag, bis hinan in das zwanzigste Kapitel, und alle dem, daß die Liebe hindert, bis an des Buches Ende.

Hierbei ist nun zu merken aufs Erste, daß Mose das Volk so genau mit Gesetzen verfaßt, daß er keinen Raum lässet der Vernunft irgend ein Werk zu erwählen oder eigenen Gottesdienst erfinden. Denn er lehret nicht allein GOtt fürchten, trauen und lieben, sondern gibt auch so mancherlei Weise äußerlichen Gottesdienstes, mit Opfern, Geloben, Fasten, Kasteien \&c., daß niemand noth sey.

\quellseite{19}

etwas anders zu erwählen. Item, er lehret auch pflanzen, bauen, freien, streiten, Kinder, Gesinde und Haus regieren, kaufen und verkaufen, borgen und lösen, und alles was äußerlich und innerlich zu thun sei, sogar, daß etliche Satzungen gleich närrisch und vergeblich anzusehen sind.

Lieber, warum thut GOtt das? Einlich darum, er hat sich des Volks unterwunden, daß es sein eigen sein sollte, und er wollte ihr GOtt sein, darum wollt er sie also regieren, daß alle ihr Thun gewiß wäre, daß es für ihn recht wäre. Denn wo jemand etwas thut, da Gottes Wort nicht zuvor auf gegeben ist, das gilt für GOtt nicht, und ist verloren. Denn er verbeut auch im 5. Buch, Kap. 4, 2., Kap. 12, 32., daß sie nichts sollten zuthun zu seinen Gesetzen. Und im Kap. 12, 8. spricht er: Sie sollen nicht thun, was sie recht dünket. Auch der Psalter und alle Propheten drob schreien, daß das Volk gute Werke thäte, die sie selbst erwähleten, und von GOtt nicht geboten wären. Denn er will und kanns nicht leiden, daß die Seinen etwas fürnehmen zu thun, daß er nicht befohlen hat, es sei wie gut es immer sein kann. Denn Gehorsam ist aller Werk Adel und Güte, der an Gottes Wort hanget.

Weil denn nun dies Leben nicht kann ohn äußerlichen Gottesdienst und Weise sein, hat er ihnen fürgeleget solche mancherlei Weise, und mit seinem Gebot verfaßet, auf daß, ob sie ja müßten oder auch wollten GOtt irgend einen äußerlichen Dienst thun, daß sie dieser einen angreifen, und nicht einen eigenen erdächten, damit sie gewiß und sicher wären, daß solch ihr Werk in Gottes Wort und Gehorsam ginge. Also ist ihnen allenthalben gewehret, eigener Vernunft und freiem Willen zu folgen, Gutes zu thun, und wohl zu leben, und doch übrig genug Raum, Stätte, Zeit, Person, Werk und Weise bestimmt, und fürleget, daß sie nicht klagen dürfen noch fremder Gottesdienst Exempel nachfolgen müssen.

Aufs Andere ist zu merken, daß die Gesetze dreierlei Art sind. Etliche, die nur von zeitlichen Gütern sagen, wie bei uns die käiserlichen Gesetze thun. Diese sind von GOtt allgemein um der Bösen willen gesetzt, daß sie nichts ärgers thäten. Darum sind solche Gesetze nur Wehr-Gesetze, mehr denn Lehr-Gesetze. Als da Mose gehet, ein Weib mit einem Scheidebrief von sich zu lassen, 5 Mos. 24. 1. Item, daß ein Mann sein Weib mit einem Eiferopfer treiben, 4 Mos. 5, 14. und andere Weiber mehr nehmen mag, 5 Mos. 25, 5.; solches sind alles weltliche Gesetze. Etliche aber sind, die von äußerlichem Gottesdienst lehren, wie droben gesagt ist.

Aber diese beide gehen nun die Gesetze vom Glauben und von der Liebe, also, daß alle andern Gesetze müssen und sollen ihr Maß haben vom Glauben und von der Liebe, daß sie gehen sollen, wo ihre Werke also geraten, daß sie nicht wider den Glauben und die Liebe gehen, wo sie aber wider den Glauben und die Liebe geraten, sollen sie schlecht ab sein.

Daher lesen wir, daß David den Mörder Joab nicht tödtet, so er doch zweimal den Tod verdienet hatte, 2 Sam. 3, 27., 20, 10., und 2 Sam. 14, 11., gelobet er dem Weibe von Thekoa, ihr Sohn solle nicht sterben, ob er wohl seinen Bruder erwürget hätte. Item, Absalom tödtet er auch nicht, 2 Sam. 14, 21. Item, er selbst David aß von dem heiligen Brod der Priester, 1 Sam. 21, 6. Item, Thamar meinet, der König möchte sie geben Amnon ihrem Stiefbruder zur Ehe, 2 Sam. 13, 13. Aus diesen und dergleichen Geschichten siehet man wohl, daß die Könige, Priester und Obersten haben oft frisch ins Gesetz gegriffen, wo es der Glaube und die Liebe haben gefordert. Daß also der Glaube und die Liebe soll aller Gesetze Meisterin sein, und sie alle in ihrer Macht haben. Denn sintemal alle Gesetze auf den Glauben und Liebe treiben, soll keines nicht mehr gelten, noch ein Gesetz sein, wo es dem Glauben oder der Liebe will zuwider geraten.

Derhalben irren die Juden noch heutiges Tages fast sehr, daß sie so strenge und hart über etlichen Gesetzen Mose halten, und viel ehe Liebe und Friede lassen untergehen, ehe sie mit uns essen und trinken, oder dergleichen thäten, und sehen des Gesetzes Meinung nicht recht an. Denn dieser Verstand ist vonnöthen allen, die unter Gesetzen leben, nicht allein den Juden. Denn also saget auch Christus Matth. 12, 11., daß man den Sabbath brechen möchte, wo ein Ochs in eine Grube gefallen wäre, und ihm heraus helfen, welches doch nur eine zeitliche Noth und Schaden war. Wie vielmehr soll man frisch allerlei Gesetze brechen, wo es Liebes-Noth fordert, so anders dem Glauben und der Liebe nichts zuwider geschieht. Wie Christus sagt, daß David gethan hat, da er die heiligen Brode aß, Matth. 12, 4.

\quellseite{20}

Was ist aber, daß Mose die Gesetze so unordentlich untereinander wirft? Warum setzet er nicht die weltlichen auf einen Haufen, die geistlichen auch auf einen Haufen, und den Glauben und Liebe auf einen? Dazu wiederholet er zuweilen ein Gesetz so oft und treibet einerlei Wort so viel mal, daß gleich verdrossen ist zu lesen und zu hören? Antwort: Mose schreibt, wie sich’s treibet, daß sein Buch ein Bild und Exempel ist des Regiments und Lebens; denn also gehet es zu, wenn es im Schwang gehet, daß jetzt dies Werk, jetzt jenes gethan sein muß. Und kein Mensch sein Leben also fassen mag (so es anders göttlich sein soll), daß er diesen Tag eitel geistliche, den andern eitel weltliche Gesetze übe, sondern Gott regieret also alle Gesetze unter einander, wie die Sterne am Himmel, und die Blumen auf dem Felde stehen, daß der Mensch muß alle Stunde zu jeglichem bereit sein, und thun, welches ihm am ersten für die Hand kömmt: also ist Mose Buch auch untereinander gemenget.

Daß er aber so fast treibet, und oft einerlei wiederholet, da ist auch seines Amts Art angezeiget. Denn wer ein Gesetzvolk regieren soll, der muß immer anhalten, immer treiben, und sich mit dem Volk wie mit Eseln bläuen. Denn kein Gesetzwerk gehet mit Lust und Liebe abe, es ist alles erzwungen und abgenöthiget. Weil nun Mose ein Gesetzlehrer ist, muß er mit seinem Treiben anzeigen, wie Gesetzwerke gezwungene Werke sind, und das Volk müde machen, bis es durch solch Treiben erkenne seine Krankheit und Unlust zu Gottes Gesetz, und nach der Gnade trachte, wie folget.

Als Dritte ist das die rechte Meinung Mose, daß er durchs Gesetz die Sünde offenbare, und alle Vermessenheit menschlichen Vermögens zu Schanden mache. Denn daher nennet ihn S. Paulus, Gal. 2, 19. einen Amtmann der Sünde, und sein Amt ein Amt des Todes, 2 Cor. 3, 6. und Röm. 3, 20. und Kap. 7, 7. spricht er: Durchs Gesetz komme nicht mehr denn Erkenntniß der Sünde. Und Röm. 3, 20.: Durch Gesetzes Werke wird niemand fromm für Gott. Denn Mose kann durchs Gesetz nicht mehr thun, weder anzeigen, was man thun und lassen soll. Aber Kraft und Vermögen solches zu thun und zu lassen, gibt er nicht, und läßt uns also in der Sünde stecken. Wenn wir denn in der Sünde stecken, so dringet der Tod alsbald auf uns, als eine Rache und Strafe über die Sünde. Daher nennet S. Paulus die Sünde des Todes Stachel 1 Cor. 15, 56., daß der Tod durch die Sünde alle sein Recht und Macht an uns hat. Aber wo das Gesetz nicht wäre, so wäre auch keine Sünde, Röm. 4, 15. Darum ist’s alles Mose Amt Schuld, der reget und rüget die Sünde durchs Gesetz, so folget der Tod auf die Sünde mit Gewalt. Daß Mose Amt billig und recht ein Amt der Sünde und des Todes von S. Paulo genennet wird, 2 Cor. 3, 6.; denn er bringet nichts auf uns durch sein Gesetzgeben, denn Sünde und Tod.

Aber doch ist solches Sündenamt und Todesamt gut und fast von nöthen: denn wo Gottes Gesetz nicht ist, da ist alle menschliche Vernunft so blind, daß sie die Sünde nicht mag erkennen, Röm. 3, 20., Kap. 7, 7. 8., denn keine menschliche Vernunft weiß, daß Unglaube und an Gott verzweifeln Sünde sei, ja sie weiß nichts davon, daß man Gott glauben und trauen soll. Gehet also dahin in ihrer Blindheit verstocket, und fühlet solche Sünde nimmermehr. Thut dieweil sonst etwa gute Werke, und führet ein äußerlich ehrbares Leben. Da meinet sie denn, sie stehe wohl, und sei der Sachen genug geschehen. Wie wir sehen an den Heiden und Heuchlern, wenn sie auf ihr bestes leben. Item, so weiß sie auch nicht, daß böse Neigung des Fleisches und Haß wider die Feinde Sünde sei; sondern weil sie stehet und fühlet, daß alle Menschen so geschickt sind, achtet sie solches für natürlich und recht gut Ding, und meinet, es sei genug, wenn man nur äußerlich den Werken wehret. Also gehet sie dahin, und achtet ihre Krankheit für Stärke, ihre Sünde für recht, ihr Böses für gut, und kann nicht weiter.

Siehe, diese Blindheit und verstockte Vermessenheit zu vertreiben, ist Mose Amt noth. Nun kann er sie nicht vertreiben, er muß sie offenbaren und zu erkennen geben. Das thut er durchs Gesetz, da er lehret, man solle Gott fürchten, trauen, glauben und lieben, dazu keine böse Lust noch Haß zu einigem Menschen tragen oder haben. Wenn nun die Natur solches Recht hört, so muß sie erschrecken. Denn sie befindet gewiß weder Trauen noch Glauben, weder Furcht noch Liebe zu Gott, item, weder Liebe noch Reinigkeit gegen den Nächsten, sondern eitel Unglauben, Zweifeln, Verachtung und Haß zu Gott, und eitel bösen Willen und Lust zum

\quellseite{21}

Nächsten. Wenn sie aber solches findet, so ist der Tod alsobald vor Augen, der solchen Sünder fressen, und in die Hölle will verschlingen.

Siehe, daß heißt den Tod durch die Sünde auf uns bringen, und durch die Sünde uns tödten. Das heißt durch das Gesetz die Sünde regen, und für die Augen setzen, und alle unsere Vermessenheit in ein Verzagen und Zittern und Verzweifeln treiben, daß der Mensch nicht mehr kann thun, denn mit dem Propheten schreien: Ich bin von GOtt verworfen, Esaias 49, 14. oder wie man auf deutsch saget: Ich bin des Teufels, ich kann nimmermehr selig werden. Das heißt recht in die Hölle geführet. Das meinet S. Paulus mit kurzen Worten, 1 Cor. 15, 15., der Stachel des Todes ist die Sünde, aber das Gesetz ist der Sünden Kraft. Als sollte er sagen: Daß der Tod sticht und uns erwürget, machet die Sünde, die an uns erfunden wird, des Todes schuldig. Daß aber die Sünde an uns funden wird, und so mächtig uns dem Tod gibt, machet das Gesetz, welches uns die Sünde offenbaret und erkennen lehret, die wir zuvor nicht kannten und sicher waren.

Nun siehe, mit welcher Gewalt Mosis solches sein Amt treibet und ausrichtet. Denn daß er ja die Natur aufs allerhöchste schände, gibt er nicht allein solche Gesetze, die von natürlichen und wahrhaftigen Sünden sagen, als da sind die zehen Gebote, sondern machet auch Sünde, da von Natur sonst keine Sünde ist, und dringet und drücket auf sie mit Haufen Sünden. Denn Unglaube und böse Lust ist von Art Sünde, und des Todes werth. Aber daß man nicht soll gesäuert Brod essen auf Ostern, 2 Mos. 12, 15. und kein unrein Thier essen, 3 Mos. 11, 4. u. folg. sein Zeichen an dem Leib machen, 3 Mos. 19, 28. 21, 5. und alles, was das Levitische Priesterthum mit Sündern schaffet, das ist nicht von Art Sünde und böse, sondern wird allein darum Sünde, daß durchs Gesetz verboten ist, welches Gesetz wohl kann ab sein, aber die zehen Gebote mögen nicht also ab sein. Denn da ist Sünde, obschon die Gebote nicht wären, oder nicht erkennet wären. Gleichwie der Heiden Unglaube Sünde ist, ob sie es wohl nicht wissen noch achten, daß Sünde sei.

Also sehen wir, daß solche und so mancherlei Gesetze Mose nicht allein darum gegeben sind, daß niemand etwas eigenes dürfte erwählen Gutes zu thun, und wohl zu leben, wie droben gesagt ist, sondern vielmehr darum, daß der Sünden nur viel würden, und sich über die Maß häufeten, daß Gewissen zu beschweren, auf daß die verstockte Blindheit sich erkennen müßte, und ihr eigen Unvermögen und Nichtigkeit zum Guten müßte fühlen, und also durchs Gesetz genöthiget und gedrungen würde etwas weiters zu suchen, denn das Gesetz und eigen Vermögen, nämlich, Gottes Gnade in künftigen Christum verheißen. Denn es ist je alles Gesetz Gottes gut und recht, wenn er auch gleich hieße nur Mist tragen, oder Strohhulm aufheben. So muß aber der ja nicht fromm noch gutes Herzens sein, der solch gut Gesetz nicht hält, oder ungern hält. So vermag alle Natur nichts anders, denn ungerne halten, darum mußte sie hie am guten Gesetz Gottes ihre Bosheit erkennen und fühlen, und nach der Hülfe göttlicher Gnaden seufzen und trachten in Christo.

Darum, wo nun Christus kommt, da höret das Gesetz auf, sonderlich das Levitische, welches Sünde machet, da sonst von Art keine Sünde ist, wie gesagt ist. So hören auch die zehn Gebote auf, nicht also, daß man sie nicht halten noch erfüllen sollte, sondern Mosis Amt höret drinnen auf, daß es nicht mehr durch die zehn Gebote die Sünde stark machet, und die Sünde nicht mehr des Todes Stachel ist. Denn durch Christum ist die Sünde vergeben, GOtt versöhnet, und das Herz hat angefangen dem Gesetz hold zu sein, daß es Mosis Amt nicht mehr kann strafen, und zu Sünden machen, als hätte es die Gebote nicht gehalten, und wäre des Todes schuldig, wie es wohl vor der Gnade, und ehe denn Christus da war.

Das lehret S. Paulus, 2 Cor. 3, 14., da er spricht: daß die Klarheit im Angesicht Mose aufhöret, um der Klarheit willen im Angesicht Jesu Christi. Das ist das Amt Mose, das uns zu Sünden und Schanden machet, mit dem Glanz der Erkenntniß unserer Bosheit und Nichtigkeit, thut uns nicht mehr weh, schrecket uns auch nicht mehr mit dem Tode. Denn wir haben nun die Klarheit im Angesicht Christi, das ist, das Amt der Gnaden, dadurch wir Christum erkennen, mit welches Gerechtigkeit, Leben und Stärke wir das Gesetz erfüllen, Tod und Hölle

\quellseite{22}

überwinden. Wie auch die drei Apostel auf dem Berge Thabor Mosen und Eliam sahen, und doch nicht vor ihnen erschraken, um der lieblichen Klarheit willen im Angesichte Christi, Matth. 17, 4., Luc. 9, 33. Aber da Christus nicht gegenwärtig war, konnten die Kinder Israel die Klarheit und Glanz in Mose Angesicht nicht erleiden, darum mußte er eine Decke dafür thun, 2 Mos. 34, 30. 33.

Denn es sind dreierlei Schüler des Gesetzes: die ersten, die das Gesetz hören und verachten, führen ein ruchlos Leben, ohne Furcht, zu diesen kommt das Gesetz nicht. Und sind bedeutet durch die Kalbthiener in der Wüsten, um welcher willen Mose die Tafeln entzwei warf, und das Gesetz nicht zu ihnen brachte, 2 Mos. 32, 6. 19.

Die andern, die es angreifen mit eigener Kraft zu erfüllen, ohne Gnade, die sind bedeutet durch die, so Mose Antlitz nicht sehen konnten, da er zum andern mal die Tafeln brachte, 2 Mos. 34, 30., zu diesen kommt das Gesetz, aber sie leidens nicht, darum machen sie eine Decke darüber, und führen ein heuchelisch Leben mit äußerlichen Werken des Gesetzes, welches doch das Gesetz alles zu Sünden macht, wo die Decke abgethan wird. Denn das Gesetz erweiset, daß unser Vermögen nichts sei ohne Christus Gnade.

Die dritten sind, die Mosen klar ohne Decke sehen. Das sind die, die des Gesetzes Meinung verstehen, wie es unmöglich Ding fordere. Da gehet die Sünde in der Kraft, 1 Cor. 15, 56., da ist der Tod mächtig, da ist des Goliaths Spieß wie ein Weberbaum, und sein Stachel hat sechshundert Seckel Erz, daß alle Kinder Israel für ihm fliehen, ohne der einige David Christus unser Herr, erlöset uns von dem allen. Denn wo nicht Christus Klarheit neben solcher Klarheit Mose käme, könnte niemand solche Glänze des Gesetzes, der Sünde und des Todes Schrecken ertragen: diese fallen ab von allen Werken und Vermessenheit, und lernen am Gesetz nicht mehr, denn allein Sünde erkennen, und nach Christo zu seufzen. Welches auch das eigentliche Amt Mose und des Gesetzes Art ist, Gal. 3, 24.

Also hat Mose auch selbst angezeigt, daß sein Amt und Lehre soll währen bis auf Christum, und alsdenn aufhören, da er spricht, 5 Mos. 18, 15.: Einen Propheten wird dir der Herr dein GOTT erwecken aus deinen Brüdern, wie mich, den sollst du hören, etc. Dies ist der edelste Spruch, und freilich der Herr im ganzen Mose, welchen auch die Apostel hoch geführet und stark gebraucht haben, Joh. 1, 45., Kap. 6, 14.; Apost. 3, 22., Kap. 7, 37., das Evangelium zu bekräftigen, und das Gesetz abzuthun, und alle Propheten gar viel daraus gezogen. Denn weil Gott hier einen andern Propheten verheißet, den sie hören sollen, zwinget sichs, daß er etwas anders lehren würde, denn Mose, und Mose seine Macht ihm übergibt und weichet, daß man jenen hören solle; so kann jener Prophet nicht Gesetz lehren, denn das hat Mose aufs allerhöchste ausgerichtet, und wäre keine Noth um Gesetzes willen einen andern Propheten zu erwecken, darum ists gewiß von der Gnadenlehre und Christo gesagt.

Darum nennet auch S. Paulus Mose Gesetz das Alte Testament, 2 Cor. 3, 14., Christus auch, da er darum ein Testament eingesetzt, Matth. 26, 28.; und ist darum ein Testament, daß Gott darinnen verhieß und beschied dem Volk Israel das Land Canaan, wo sie es halten würden; und gab's ihnen, und ward bestätiget durch Schöps- und Bocks-Tod und Blut, 2 Mos. 24, 8. Aber weil solch Testament nicht auf Gottes Gnaden, sondern auf Menschen-Werken stund, mußte es alt werden und aufhören, und das verheißene Land wieder verloren werden, darum, daß durch Werke das Gesetz nicht kann erfüllet werden. Und mußte ein ander Testament kommen, das nicht alt würde, auch nicht auf unserm Thun, sondern auf Gottes Wort und Werken stände, auf daß es ewiglich währete, Hebr. 13, 20. Darum ists auch durch einer ewigen Person Tod und Blut bestätiget, und ein ewiges Land verheißen und gegeben. Das sei nun von Mose Büchern und Amt geredet.

Was sind aber nun die andern Bücher der Propheten und der Geschichten? Antwort: Nichts anders, denn was Mose ist, denn sie treiben allesamt Moses Amt, und wehren den falschen Propheten, daß sie das Volk nicht auf die Werke führen, sondern in dem rechten Amt Mose und Erkenntniß des Gesetzes bleiben lassen. Und halten fest darauf, daß sie durch des Gesetzes rechten Verstand die Leute in ihrer eigenen Untüchtigkeit behalten, und auf Christum treiben, wie Mose thut. Darum streichen sie auch weiter aus, was Mose von Christo gesagt hat, und zeigen an

\quellseite{23}

beiderlei Exempel, derer die Mose recht haben, und derer die ihn nicht recht haben, und aller beider Strafe und Lohn. Also, daß die Propheten nichts anders sind, denn Handhaber und Zeugen Mose und seines Amts, daß sie durchs Gesetz jedermann zu Christo bringen.

Aufs letzte sollte ich auch wohl die geistliche Deutung anzeigen, so durch das Levitische Gesetz und Priestertum Mose fürgelegt. Aber es ist sein zuviel schreiben, es will Raum und Zeit haben, und mit lebendiger Stimme ausgelegt sein. Denn freilich Mose ein Brunn ist aller Weisheit und Verstandes, daraus gequollen ist alles, was alle Propheten gewußt und gesagt haben, dazu auch das Neue Testament heraus fleußt, und darin gegründet ist, wie wir gehöret haben. Aber doch ein kleines kurzes Grifflein zu geben, denjenigen, so Gnade und Verstand haben, weiter darnach zu trachten, sei das mein Dienst.

Wenn du willst wohl und sicher deuten, so nimm Christum für dich, denn das ist der Mann, dem es alles ganz und gar gilt. So mache nun aus dem Hohenpriester Aaron niemand denn Christum alleine, wie die Epistel an die Hebräer thut, Kap. 7, 15 folg., welche fast allein genugsam ist, alle Figuren Mose zu deuten. Also ists auch gewiß, daß Christus selbst das Opfer ist, ja auch der Altar, der sich selbst mit seinem eigenen Blut geopfert hat. Wie auch dieselbe Epistel meldet, Kap. 7, 27., Kap. 13, 10. Wie nun der Levitische Hohepriester durch solch Opfer nur die gemachten Sünden wegnahm, die von Natur nicht Sünde waren: also hat unser Hohepriester Christus, durch sein selbst Opfer und Blut die rechte Sünde, die von Natur Sünde ist, weggenommen. Und ist einmal durch den Vorhang gegangen zu Gott, daß er uns versühne. Also, daß du alles, was vom Hohenpriester geschrieben ist, auf Christum persönlich, und sonst auf niemand deutest.

Aber des Hohenpriesters Söhne, die mit dem täglichen Opfer umgehen, sollten auf uns Christen deuten, die wir für unsern Vater Christo im Himmel sitzend, hie auf Erden mit dem Leibe wohnen, und nicht hindurch sind bei ihm, ohne mit dem Glauben geistlich. Derselben Amt, wie sie schlachten und opfern, bedeutet nichts anders, denn das Evangelium predigen, durch welches der alte Mensch getödtet, und Gott geopfert durchs Feuer der Liebe, im heiligen Geist verbrannt und verzehret wird, welches gar wohl reucht für Gott, das ist, es macht ein gut, rein, sicher Gewissen für Gott. Diese Deutung trifft S. Paulus, Röm. 12, 1., da er lehret, wie wir unsere Leiber sollen opfern Gott zum lebendigen, heiligen, angenehmen Opfer. Welches wir thun (wie gesagt) durch stetige Uebung des Evangelii, beide mit Predigen und Glauben. Das sei diesmal genug zur kurzen Anleitung, Christum und das Evangelium zu suchen im Alten Testament.

\quellseite{24}

\chapter{Vitus Dieterich, an den christlichen Leser}

Ich bitte alle frommen Christen, die diese Summaria lesen, sie wollen meine Arbeit und Dienst ihnen gefallen lassen, und zum besten deuten. Denn ichs je christlich gemeinet, und ferner nicht gedacht habe, denn wie durch kurzen und einfältigen Unterricht auch die Ungelehrten möchten zur heiligen Schrift recht geleitet werden. Derhalben ich die scharfe und subtile Auslegung mit Fleiß gemieden, und allenthalben mich an dem habe genügen lassen, daß verständlich, einfältig, und der eigentlichen Deutung der Schrift gemäß gewesen. So hat es sich keineswegs schicken wollen, daß ich es hätte weitläufiger gemacht, weil ich anfänglich diese Arbeit für mich genommen habe, daß solche Summaria des Alten Testaments in meiner Kirchen, vor den Kapiteln gelesen würden, und jetzund viel andere Kirchen solchen Brauch auch angenommen haben, und diese Summaria zum kurzen Unterricht lesen. Denn es muß zwischen solcher Lection und der Predigt ein Unterschied sein. In der Predigt soll man genau allem nachgehen, und nichts dahinten lassen. In solcher Lection aber ists genug, daß man nur auf die vornehmsten Punkte sehe, und die Lehre herausnehme, so am verständlichsten und zur Besserung am meisten dienstlich. Wie ich denn hoffe, es sei hiemit treulich geschehen.

Daß ich aber des Ehrwürdigen Herrn, und meines lieben Präceptors, Doct. Luthers Summaria über den Psalter, habe darunter gesetzt, ist der Ursache geschehen, daß ich wohl wußte, ich würde es nicht besser machen, und sind dieselben bereits unter den Leuten gewesen. Wo dergleichen Summaria über die ganze Bibel vorhanden gewesen, wollte ich dieser Arbeit gar \unsicher{fein} müßig gestanden, und mich an dem, daß andere gemacht haben, gern begnügen lassen. Weil aber dergleichen nicht vor der Hand war, achte ich, es werde niemand solches mein Fürnehmen übel deuten noch strafen können, sintemal ich anders nichts, denn Erweiterung des Worts und Besserung der Kirchen damit gesucht habe. So ist es auch niemand gewehret noch verboten, der es kann, daß ers besser mache. Ohne allein, daß ich wohl weiß und sehe, daß Gott die muthwillige Verachtung und schreckliche Lästerung seines Worts an den Papisten dermaßen strafet, daß sie schier, wie die Juden, mit sehenden Augen starr blind sind, und Gottes Wort ihnen zum lautern Strick gerathen ist, wie man an ihrem Gemäsch und Lästern täglich sehet. Derhalb nicht zu besorgen, daß sie in diesem oder andern etwas verbessern. Besorgen muß man, wie sie von Tag zu Tag je länger je ärger und verstockter werden, also werde der Satan sie treiben, je länger je mehr Muthwillen zu üben, mit Verderbung ihrer selbst, und anderer, die Gottes Wort geringer achten, denn des Pabstes Geifer und Menschensatzung.

Der ewige barmherzige Gott, der solches Licht seines Worts uns so reichlich läßt scheinen, wolle mit Gnaden uns dabei erhalten, und dem Satan, der beide, durch den Türken und den Pabst, seine willigen Werkzeuge, damit umgehet, wie ers dämpfen oder fälschen möge, mit Ernst wehren, um seines lieben Sohnes und unsers Herrn Christi Jesu willen, durch seinen heiligen Geist; dem sei Lob und Ehre in Ewigkeit, Amen.

\quellseite{25}

\section{Widmung an Albrecht, Markgraf zu Brandenburg}

Dem Durchläuchtigen und hochgebornen Fürsten und Herrn, Herrn Albrechten, Markgrafen zu Brandenburg, in Preußen, zu Stettin, Pommern, der Caßuben und Wenden Herzogen, Burggrafen zu Nürnberg und Fürsten zu Rügen etc. etc.

Meinem gnädigsten Herrn:

Gnade, Friede und Barmherzigkeit, von Gott unserm Vater und unserm Herrn Jesu Christo.

Es ist eine sehr nütze und feine Ordnung jetzund hin und wieder in den Kirchen, so das Evangelium Christi angenommen, und von dem Gräuel des Papstthums abgetreten sind, daß man neben dem Predigen, alle Tage etliche Kapitel aus dem Alten und Neuen Testament von Wort zu Wort dem Volk vorlieset. Denn es will den Christen hoch vonnöthen sein, daß sie in steter Uebung des Worts bleiben, und nicht allein in der Historien der heiligen Schrift läufig und bekannt, sondern auch in den Predigten der Propheten erfahren sein.

Also siehet man, daß es im Judenthum eine sonderliche Ordnung gehabt hat, daß die Bücher Mose alle Sabbather für der Gemeine sind gelesen worden. Item, das Exempel Christi im Tempel, Luc. 4., zeigt auch an, daß man in den Synagogen der Propheten Bücher gelesen, und die Juden davon disputiret haben. Und der Apostel Sanct Paulus vermahnet seine Colosser, Cap. 3., daß sie das Wort Christi reichlich unter ihnen sollen wohnen lassen, in aller Weisheit, daß ein jeglicher sich selbst dadurch lehren und vermahnen möge.

Warum aber solches alles unter den Juden erstlich also geordnet, und hernach in der christlichen Kirchen auch gehalten, und heutiges Tages wieder in Brauch geführet ist, kann jedermann leichtlich urtheilen, der nur ein wenig geistlicher Erfahrung hat. Denn weil der Anfechtungen sehr viel sind, so ohne Unterlaß vom Teufel, der Welt, und unserm eigenen Fleisch erreget werden, ists unmöglich, daß man sich dagegen aufhalten, wehren und trösten möge, denn allein durch das Wort Gottes.

Darum siehet man auch, daß in der Schrift allenthalben Gottes Wort so herrlich gepreiset wird. Und sonderlich der 119. Psalm, neben andern schönen Namen, die er dem Wort Gottes gibt, nennet ers auch ein Licht, nach dem man gehen, und im Finstern sich für dem Anstoßen bewahren soll. Und hat hernach der heilige Apostel S. Petrus, 2 Petri 1., solchen Namen auch gebraucht. Denn weil wir hie auf Erden, unsers verderbten und sündlichen Fleisches, und anderer Aergerniß halben (gleichwie im Finstern) wandeln, und es sehr leicht geschehen kann, daß wir in Sünde oder Irrthum, entweder selbst fallen oder geführet werden, ist kein anderer Befehl, Rath noch Mittel, wider solche Fahr, denn daß die Christen dies Licht stets in Händen tragen, und nach demselbigen nicht allein ihre Wege und Wandel, so viel das ganze Leben betrifft, richten, sondern auch den Glauben rein und lauter, wider alle irrige Meinung erhalten.

Nun ist aber im Alten Testament viel Ding, in welche der gemeine Mann sich nicht richten kann, es sei denn, daß er zuvor eine Anleitung habe. Sonderlich aber siehet mans in den Propheten, daß unterweilen auch gelehrte Leute zu schaffen haben, daß sie die Predigt recht in einander richten, und die rechte Meinung treffen mögen. Denn der heilige Geist, ob er wohl durch die Propheten alles was von Christo und der Seligkeit zu wissen vonnöthen ist, genugsam hat angezeigt, so hat ers doch verborgener heimlicher Weise gethan.

Denn dies ist fürnehmlich unsers Herrn Christi Amt, daß er durch sein Evangelium vom Willen Gottes, Vergebung der Sünden, und unsrer Seelen Seligkeit uns lauter und klar, ohne alles Verdecken, unterrichten soll. Wie der heilige Johannes der Täufer von ihm sagt: Niemand hat Gott je gesehen. Der eingeborne Sohn, der in des Vaters Schooß ist, der hats uns verkündiget.

\quellseite{26}

Weil nun solches in den Propheten und dem Alten Testament dunkel ist angezeigt, und auch sonst viel Historien vorfallen, aus welchen sich nicht jedermann genugsam kann richten, ja da auch Aergerniß bei zu besorgen, wo man durch rechten Unterricht nicht zuvorkommt, habe ich meiner Kirchen hie zu gut, solche Arbeit auf mich genommen, und in das Alte Testament, auf alle Kapitel kurze Summaria gestellet, auf daß der gemeine Mann solche gewöhnliche Lection mit mehr Frucht hören, und allzeit eine kurze Lehre aus den Historien oder Predigten der Propheten fassen könnte, welche zur Besserung des Lebens und Unterricht des Glaubens nütz und dienstlich wäre.

Nun wirds nicht fehlen, solche Arbeit, so ich in meiner Kirchen hie zum besten habe vorgenommen, werde etlichen aus der Ursache nicht am besten gefallen, daß solche Summaria nicht alles fassen, das in den Historien und in der Propheten Predigten, in einem jeglichen Kapitel begriffen ist. Weil ich aber vornehmlich dahin habe müssen sehen, wie es solchem Kirchengebrauch am gelegensten wollte sein, hat es sich in keinen Weg wollen schicken, weil die Kapitel sonst lang genug sind, daß man allerwege eine lange Predigt und vollkommene Auslegung wollte daran hängen.

Habe mir also an dem genügen lassen, daß allerwege aufs kürzeste das angezeigt würde, welches zu merken am nöthigsten und nützlichsten ist. Denn auch neben dem, daß man solche Lectionen nicht soll zu lang machen, ists auch für den gemeinen Mann nicht, daß man alles weitläufig handeln, und die Kapitel in viel Stücke theilen wollte. Es ist genug, wo mit kurzen Worten allerwege das Fürnehmste angezeigt wird. Denn die vollkommene Erklärung gehört in die Predigt, nicht in solche Lection.

Es könnte auch wohl geschehen, daß man von etlichen Sprüchen und Historien schärfer und subtiler disputiren möchte, denn ich in diesen Summarien habe wollen thun, ich habe aber mich durch und durch beflißen, daß ich allenthalben dem jungen und einfältigen Volk zu gut, die gemeine und richtige Meinung anzeigte, die ihnen sein verständlich und nützlich wäre. Denn subtil disputiren, zu dem, daß es für den gemeinen Mann nicht taugt, kanns auch wohl in etliche Wege schädlich sein.

Um solche und andere dergleichen Ursachen mehr habe ich oft die Kapitel von einander getheilet, und zwei Summaria darauf gestellet, nicht allein darum, daß die Kapitel an ihnen selbst zu lang sind, wo man eine Ordnung mit dem Lesen hält, wie hie zu Nürnberg, sondern auch, daß oft die Händel darnach sind, daß es wohl bedarf, daß der Zuhörer ein jeglich Stück insonderheit bedenke, was es auf ihm habe, und wie ers ihm zur Besserung brauchen soll.

Item, ich habe auch viel Kapitel ausgelassen, und keine Summaria darüber gestellet, um der Ursachen willen, daß sie meines Erachtens, ohne sonderliche Frucht in der Kirchen heutiges Tages gelesen werden. Als da Mose den Tabernakel oder Hütten, mit seinem Zubehör, Salomon den Tempel und des Königs Haus bauet. Item, da Mose lehret, wie man mit dem Opfern umgehen, und die hohen Feste halten soll. Item, da man die Stämme des Volks Israel zählet rc. Denn ob man gleich solches in eine Allegorien ziehen, und eine heimliche Deutung daraus flechten will, wie es sich oft an etlichen Orten fein reimen läßt, so ists doch ohne Frucht, daß man öffentlich in der Kirchen den bloßen Text lieset.

Sonderlich aber halte ich keineswegs mit den Wiederläufern, die ihnen träumen lassen, die höchste Weisheit sei im Bau des Tabernakels und Tempels heimlich angezeigt. Eben wie die Juden auch narren, und halten das erste und zehnte Kapitel des Propheten Hesekiels für den heimlichsten Schatz der ganzen heiligen Schrift.

Wir Christen sollen das für das größte, und (wie es S. Paulus nennet) für das Geheimniß achten, so von Anbeginn der Welt ist verborgen gewesen, daß wir durch den Sohn Gottes vom Tod und der Sünden am Kreuz erlöset, und aus Gnaden, um Christi Willen, selig sind worden. Nach solchem Geheimniß ist das größte, daß Gott seinen Willen in den zehn Geboten offenbaret hat. Bei dieser Weisheit und Geheimniß sollen wir bleiben, und solches für das Höchste achten, wiederum alles anders, was in der heiligen Schrift und sonst ist, dem nirgend gleich rechnen. Weil nun das erste und zehnte Kapitel Hesekiels, ob es gleich eine nützliche Lehre in sich hat, dennoch so dunkel ist, daß man sich nicht wohl durchaus recht darein kann schicken, habe ichs in seinem Werth bleiben, und unter andern lassen liegen, welche ich für unnütz-

\quellseite{27}

thig achte, daß mans in der Kirchen lesen soll. Wer aber die Deutung solches Gesichtes begehrt zu wissen, kann es in der Vorrede Dr. Luthers über den Ezechiel lesen.

Nun bin ich Anfangs nicht Willens gewesen, solche Summaria in Druck zu geben, habe es auch eine zeitlang in meiner Kirchen allein brauchen lassen, aber weil sie doch endlich in fremde Hände kommen, und ich von vielen darum bin angesucht worden, daß ichs in Druck wollte geben, sintemal sie nicht allein dem gemeinen Mann, sondern auch vielleicht den schlechten einfältigen Kirchendienern, hin und wieder auf dem Lande, möchten zu Nutz kommen, habe ichs in die Länge nicht können weigern, ob, wie Hoffnung ist, der christlichen Kirchen ein wenig damit möchte gedienet werden.

Ich habe aber, Gnädigster Fürst und Herr, solche Arbeit fürnehmlich E. F. G. darum wollen dediciren, daß ich weiß, daß E. F. G. ein sonderlich Herz und Liebe zum Wort hat, und viel darauf gewagt, daß die Kirchen in E. F. G. Landen göttlich und recht würden angericht't. Denn gewiß ist es, wo E. F. G. nicht so treulich dazu geholfen und gerathen hätte, es wäre ganz Preußen noch heutiges Tages unter dem Pabstthum, und müßte des Evangelii eben sowohl, als die anstoßenden Herrschaften beraubt sein.

Daß aber nun GOttes Wort allenthalben in E. F. G. Landen recht geprediget, die Sacramente nach der Einsetzung Christi, den Leuten zum Trost und Seligkeit, recht gehandelt werden und GOttlob die Frucht des Worts nicht außen bleibt, sondern alle Tage viel Christen werden, und in Christo selig entschlafen, solches fürstlichen, ja bischöflichen und besten Werks ist E. F. G. durch Gottes Hülfe eine sonderliche Ursache.

Darum auch alle Christen E. F. G. billig lieb und werth haben, und für einen christlichen frommen Fürsten mit aller Wahrheit rühmen und loben, dagegen die andern, so um ihres Bauchs willen dem Pabstthum anhangen, und das heilige Evangelium, wo sie können, unbillig verfolgen, wohl werth sind, daß sie alle Welt hasse, und von ihnen sage, wie sie es mit ihren gottlosen Werken beweisen.

Sonderlich aber, gnädigster Fürst und Herr, hat mich zu solchem auch dieß verursacht, daß E. F. G. neulich mich ganz gnädigst hat lassen ersuchen, ob ich mich bei E. F. G. zu Kirchendiensten begeben wollte. Weil ich aber vorhin mit Diensten hie in meinem Vaterlande verpflichtet, und das Predigtamt sonderlich dieß auf ihm hat, daß niemand seinen Beruf ehe soll noch kann wechseln, denn man wolle sein nicht, habe ich mit solchem Dediciren E. F. G. meinen schuldigen Dienst und geneigten Willen wollen anzeigen, ob damit E. F. G. Kirchen, weil ichs gegenwärtig mit dem Leibe nicht kann, etwas möchte gedienet werden. Bitte derhalben E. F. G. wolle solches im besten annehmen, und mich E. F. G. gnädig lassen befohlen sein. Gott bewahre E. F. G. zu seinen Ehren, zu E. F. G. Seligkeit, und des ganzen Landes Wohlfahrt, Amen. Datum Nürnberg, am Tage Simonis und Judä, 1540.

Euer Fürstl. Gnaden

\begin{unterschrift}

unterthäniger und williger\\Vitus Dieterich,\\zu Nürnberg Prediger in der \unsicher{Sebalder} Pfarrkirchen.

\end{unterschrift}

\quellseite{28}

\begin{titelseite}

Vorreden und Beschlüsse\\[0.35em]

über die Kapitel der Bücher des Alten und Neuen Testaments,\\[0.35em]

wie dieselbigen sammt den Summarien Herrn Veit Dieterichs seligen, in der Kirche Gottes zu Breslau in den gewöhnlichen Kirchenlectionen, die Werktage über, zum Frühgebet, verlesen werden.\\[0.35em]

Gesetzt, und jezund neulich in Druck gegeben durch,\\[0.35em]

Franciscum Bierling,\\[0.35em]

Diener des göttlichen Worts daselbst, bei S. Maria Magdalena.\\[0.35em]

CUM PRIVILEGIO SENATUS.\\[0.35em]

ANNO 1676.

\end{titelseite}

\chapter{Widmung an den Rath der Stadt Breslau}

Den Edelen, Gestrengen, Ehrenvesten, Hochweisen und Wohlbenamten Herrn

J. Rathmannen

der Kaiserl. und Königl. Stadt Breslau, haltend und verwaltend die Hauptmannschaft Breslauischen Fürstenthums, Neumärkischen und Namslischen Weichbilder,

Meinen Großgünstigen Herren, Patronen, Förderern:

Gnade, Friede und Segen, von Gott dem Vater, durch seinen liebsten Sohn, unsern Herrn Jesum Christum, neben kräftigem Trost des heil. Geistes, sammt meinem andächtigen Gebet und unterthänigen Diensten, mit Fleiß zuvor.

Edle, gestrenge, ehrenveste, hochweise und wohlbenamte großgünstige Herren, es sind nunmehr acht und zwanzig Jahr verlaufen (dafür dem Vater aller Gnade und Barmherzigkeit, GOtt unserm Herrn, von dem alle gute Gabe und alle vollkommene Gabe von oben herab kommt, in wahrer Demuth, von ganzem Herzen, Lob, Ehre, Preis und Dank gesagt sei, durch JEsum Christum, seinen Sohn und unsern Mittler und Fürsprecher), daß ein ehrbarer, hochweiser und wohlbenamter Rath, bei den dreien Pfarrkirchen dieser löblichen Stadt anordnen hat lassen, daß neben den gemeinen öffentlichen Predigten auf die andern ledigen Tage der Wochen soll gehalten werden die Lection oder das Vorlesen der heiligen Biblia, der Bücher des Alten und Neuen Testaments, von den hochgelehrten Propheten und auserwählten Aposteln geschrieben, die von der christlichen Kirche Canonici libri, das ist, bewährte Bücher genennet werden, darum, daß die allein zu Zeugen, die Lehre der christlichen Religion zu bestätigen, sind zugelassen worden. Und dieses mit den Summarien des ehrwürdigen und hochgelehrten Herrn Magisters Veit Dieterichs (in GOtt seligen), welcher eine zeitlang ein fleißiger Schüler des werthen Mannes GOttes, Herrn Dr. Martin Luthers, und hernach Prediger zu Nürnberg in St. Sebalds-Kirche gewesen. Darinnen aufs kürzeste angezeigt wird, was am nöthigsten und nützlichsten ist dem jungen Volk und gemeinen Manne, aus allen Capiteln gemeldeter biblischer Schriften, zu wissen und zu lernen, darnach sie ihr Leben richten und solche seine Lehre zu ihrer Seelen Seligkeit brauchen können.

Diese nöthige und gute Ordnung hat ein ehrbarer, hochweiser und wohlbenamter Rath, als eine christliche Obrigkeit, die ihrer treuen Unterthanen zeitliche und ewige Wohlfahrt gern

\quellseite{29}

befördern wollte, billig machen lassen, und ist freilich auch alles Lobes werth, und als eine besondere große Gnade und Wohlthat GOttes zu erkennen, und mit Dankbarkeit allzeit zu preisen.

Denn es soll ja die theure, werthe, liebe heilige Biblia, und also die göttlichen Schriften der heiligen Propheten und Aposteln, dem Volke Gottes, sonderlich in den gewöhnlichen ordentlichen Kirchenversammlungen, täglich und fleißig vorgelesen werden.

Erstlich darum, daß GOtt, unser Herr, solches ernstlich befohlen hat, als der am besten verstehet, was uns Menschen und seinen Kindern auf Erden gut, nöthig und nütze ist. Gottes, des Vaters, Befehl ist, der also im 31. Kapitel des fünften Buches Mose aufgeschrieben stehet: Und Mose schrieb dies Gesetz und gabs den Priestern, den Kindern Levi, die die Lade des Bundes des HERRN trugen, und allen Ältesten Israel. Und gebot ihnen, und sprach: Je über sieben Jahr, zur Zeit des Erlaßjahrs, am Fest der Laubhütten, wenn das ganze Israel kommet zu erscheinen vor dem HErrn, deinem Gott, an dem Ort, den er erwählen wird, sollst du dies Gesetz vor dem ganzen Israel ausrufen lassen vor ihren Ohren, nämlich vor der Versammlung des Volks, beide der Männer und Weiber, Kinder und deines Fremdlings, der in deinem Thor ist; auf daß sie hören und lernen, damit sie den HErrn, ihren Gott, fürchten, und halten, daß sie thun alle Worte dieses Gesetzes, und daß ihre Kinder, die es nicht wissen, auch hören und lernen, damit sie den HErrn, euren Gott, fürchten, alle eure Lebtage. Diese Worte lehren, daß es Gottes endlicher und ernstlicher Wille sei, daß man die göttliche Lehre in den Büchern, welche man die heilige Schrift und Biblia nennet, verfasset, dem christlichen Volk fleißig vorlesen, bei demselben treiben, lehren, wiederholen und erklären soll, welches nicht allein in sieben Jahren einmal, sondern für und für alle Jahre und Tage im Schwange gehen soll; insonderheit aber auf alle Sabbathe, Sonntage und Feiertage. Der Sohn Gottes, unser HErr JEsus Christus, hat auch solches befohlen und geboten; denn also sagt er, Luc. 16.: Sie haben Mosen und die Propheten, laß sie dieselbigen hören. Und abermals Johannes im 5. Kapitel: Suchet in der Schrift (vernehmt durch stetiges Lesen, andächtiges Anhören und fleißige Betrachtung), denn diese ists, die von mir zeuget. Diese Mandate und Befehle unseres HErrn Christi soll ihm ein jeder Christ also einbilden, als rufe er ihm diese alle Tage zu, und soll derwegen keinen Tag dahin gehen lassen (sofern er nicht durch wichtige und nöthige Geschäfte verhindert wird), daß er nicht Gottes Wort sollte anhören. GOtt, der heilige Geist, befiehlet und vermahnet das auch durch Paulus, den Apostel und theuren Rüstzeug Gottes, Col. 3. und sagt: Lasset das Wort Christi unter euch reichlich wohnen, in aller Weisheit. Diese Worte sind ein solch göttliches Gebot, daß Gottes Wort bei uns soll in vollem Schwange gehen, und uns gar wohl bekannt sein, auch immer für und für bleiben. Denn wohnen ist an einem Ort eine gute lange Zeit sein und bleiben, da man nicht bald wieder abzieht. Es soll uns die heilige Biblia gar wohl bekannt sein, nicht anders, als diejenigen Leute, so bei uns und in einem Hause, oder ja hart neben uns wohnen, mit welchen wir täglich viel umgehen, Gespräche halten und zu thun haben. Es soll uns Gottes Wort also gemein sein, wie das tägliche Brod, welches wir alle Tage müssen haben und genießen, und dessen wir keinen Tag zur Nothdurft können entbehren. Denn wie das tägliche Brod eine leibliche Speise ist, die den Leib erhält und ernähret: also ist Gottes Wort eine geistliche Speise der Seelen, die uns muß erhalten zum ewigen Leben. Daraus folget, daß wir abermals einen starken Befehl haben, daß wir derwegen täglich Gottes Wort lesen oder anhören sollen.

Die andere Ursache, warum alle Tage Gottes Wort in den Kirchenversammlungen soll vorgelesen werden, ist die Nothdurft, daß es alle Menschen, jung und alt, in allen Ständen wohl bedürfen. Denn die Menschen sind von Natur zur Sünde nur sehr geneigt. So setzt der Teufel täglich an die Menschen mit sonderlichen Tücken und Stücken, daß er sie, zu ihrem Verderben, in Sünde und Schande locke, führe und stürze. Die Aergernisse der Welt-Kinder sind auch groß und überall gemein. Damit nun jedermann wider solche Gefahren und Unfälle verwahret werde, ists ja nöthig, daß täglich zur Warnung Gottes Wort vorgelesen werde, damit sich niemand sein sündlich Fleisch und Blut, den bösen neidischen Teufel und die ärgerliche Welt in Sünden an-

\quellseite{30}

führen und zu seinem zeitlichen und ewigen Schaden verführen lasse. Denn das göttliche Wort muß vorleuchten in allem unserm Vornehmen, Rath und That, Thun und Lassen, sollen wir anders unser ganzes innerliches und äußerliches Leben und Wandel recht anstellen; auch unsere bösen Affecten und Zuneigungen dämpfen und im Zaum halten, und unserer Vocation und Beruf treulich vorstehen wider alle Aergernisse und Hindernisse. Denn die heilige Schrift ist die einige Norm und Form, Regel und Richtschnur, nicht allein des rechten seligmachenden Glaubens, sondern auch aller Gottseligkeit, als die da lehret, wie man Gott in rechtem Glauben und gutem Gewissen, wahrhaftig und gefällig dienen soll. Solches bezeuget der königliche Prophet David in Frage und Antwort, da er spricht im 119. Psalm: Wie wird ein Jüngling des Herrn Weg unsträflich gehen? Wenn er sich hält nach seinen Worten und Geboten. Und sagt ferner: Ich behalte dein Wort in meinem Herzen, auf daß ich nicht wider dich sündige. Dein Wort ist meines Fußes Leuchte, und ein Licht auf meinem Wege. So ist auch Gottes Wort eine starke Verwahrung und Antidotum wider die Sünde, wider die Anreizung des Teufels und wider die Aergernisse der gottlosen Epicurer und Sünder, die sie den Christen an den Weg, darüber sie gehen müssen, legen und eingraben. So sind auch unter den wahren Christen allzeit betrübte Herzen, die das Joch des heiligen Kreuzes tragen, jetzt am Leibe, jetzt an ihrer Nahrung, in diesen ganz kümmerlichen und geschwinden Zeiten, bald an ihren Ehegatten und Kindern, und die mit Schwermuth und Anfechtung beladen sind, daß sie in ihrem Nothstande unter Gottes gewaltige Hand sich demüthigen, und solch ihr Kreuz sich nicht unterdrücken lassen, sondern täglichen nöthigen Trost haben, mit Geduld ihre Seele fassen und in wahrem Glauben der verheißenen gewissen Hülfe Gottes erwarten. Derwegen erfordert's abermals die Noth solcher und anderer Kreuzträger und Kreuzträgerinnen, daß ihnen zu Trost Gottes Wort täglich vorgelesen werde. Denn die heilige Schrift ist auch das rechte Trostbuch, daraus allein angefochtene, bekümmerte und geängstigte Herzen zu jeder Zeit rechten, gewissen und beständigen Trost schöpfen und erlangen mögen, durch Gnade des heiligen Geistes, deß Buch es ist, als des rechten himmlischen Kanzlers. Davon zeugt abermals David, der israelitische König und Prophet Gottes im 119. Psalm und sagt: Wenn dein Gesetz nicht wäre mein Trost gewesen, so wäre ich vergangen in meinem Elende. Und der heilige Prophet Jeremias im 15. Kapitel danket dem Herrn für solche Kraft und Nutzbarkeit seines göttlichen Wortes, und spricht ihm also zu: Indeß enthalte uns dein Wort, wenn wirs kriegen; und dasselbe dein Wort ist unsers Herzens Freude und Trost. Von des göttlichen Wortes kräftiger und heilsamer Wirkung zeuget auch der heilige Apostel Paulus im 15. Kapitel seiner Epistel an die Römer, mit davon folgendem Lobsspruch: Was zuvor geschrieben, das ist uns zur Lehre geschrieben, auf daß wir durch Geduld und Trost der Schrift Hoffnung haben. Ferner ist auch dies eine große Nothdurft aller Menschen, weil wir leben und wissen nicht, wie lange, und weil wir sterben, und wissen nicht, wann; daß ein jeder, wie mit dem Namen, also auch mit der That, ein wahrer Christ sei und bleibe, da er anders nicht soll in Sünden sterben und ewig verderben. Also ist auch die Frage, wer denn ein rechter Christ sei? Antwort: Ein rechter Christ ist der Mensch, welcher in täglicher wahrer christlicher Buße lebet; der alle Tage Gott, den himmlischen Vater, im Namen Jesu Christi, seines lieben Sohnes, bittet um Vergebung aller seiner Sünden; der auch ein gläubig Herz in seinem Leibe augenblicklich herumträgt, daß er mit Gott versöhnet sei, und habe Gnade erlanget um des Herrn Christi Gehorsams und Leidens willen; der fleißig betet für sich, die Seinen und für die gemeine Noth der Christenheit; der unter dem von Gott auferlegten Kreuz geduldig ist; der in seinem Beruf vor wissentlichen, muthwilligen und vorsätzlichen Sünden sich hütet, und dagegen züchtig gegen ihn selbst, gerecht gegen seinen Nächsten und gottseelig gegen Gott lebet, so viel ihm möglich, durch Hülfe des heiligen Geistes, und der, da er aus menschlicher Schwachheit oder Verführung der Welt und des Teufels einen Fehltritt thut, mit bösen Gedanken, Worten und Werken sündiget, bald solches in wahrem Glauben Gott abbittet, und der im Herrn Christo endlich selig entschläft und aus dieser Welt scheidet. Diese Bewegungen, Werke und Tugenden des wahren Christenthums kommen aus Gottes Wort, und werden dadurch gestärket und vermehret von Tag zu Tag in denen, die Gottes Wort mit Ernst, Andacht und Fleiß hören, und im Gehorsam desselbigen sich üben. So erfordert's demnach wiederum die höchste Nothdurft aller Menschen, daß sie aus Gottes Wort davon gelehret und dazu vermahnet werden.

\quellseite{31}

Daß aber die heilige Schrift und Gottes Wort also von wahrem Christenthum lehre und unterrichte, bekräftigen folgende Zeugnisse, insonderheit unsers, der Heiden, Apostels S. Pauli:

Röm. 3. sagt er: Durch das Gesetz kommt Erkenntniß der Sünden.

Röm. 10.: So kommt der Glaube aus der Predigt, das Predigen aber durch das Wort Gottes.

2 Cor. 5.: Gott hat uns mit ihm selber versöhnet durch Jesum Christ, und das Amt gegeben, das die Versöhnung prediget. Denn GOtt war in Christo, und versöhnte die Welt mit ihm selber, und rechnete ihnen ihre Sünde nicht zu, und hat unter uns aufgerichtet das Wort von der Versöhnung.

Röm. 10.: Wer den Namen des HErrn wird anrufen, soll selig werden. Wie sollen sie aber anrufen, an den sie nicht glauben? Wie sollen sie aber glauben, von dem sie nichts gehört haben? Wie sollen sie aber hören ohne Prediger?

Röm. 15.: (Damit dieser nöthige Spruch abermal wiederholet werde). Auf daß wir durch Geduld und Trost der Schrift Hoffnung haben.

2 Tim. 3.: Alle Schrift von Gott eingegeben, ist nütz zur Lehre, zur Strafe, zur Besserung, zur Züchtigung in der Gerechtigkeit, daß ein Mensch Gottes sei vollkommen, zu allem guten Werk geschickt.

Röm. 1.: Das Evangelium von Christo ist eine Kraft Gottes, die da selig macht alle die daran glauben. Und 1 Cor. 1.: Dieweil die Welt durch ihre Weisheit Gott nicht erkannte, gefiel es Gott wohl, durch thörichte Predigt selig zu machen die, so daran glauben.

Die dritte Ursach, daß man auch heute noch Gottes Wort täglich vorlesen, und sich darin üben solle, sollen sein die löblichen Ordnungen und besondern Exempel der Heiligen Gottes, und theuren werthen Männer in der heiligen christlichen Kirchen, als:

Im 12. Kapitel des 1. Buchs Mose wird gesagt von dem Patriarchen und Altvater Abraham, daß er eine Hütten aufgerichtet hat, und bauete dem HErrn einen Altar, und predigte vom Namen des HErrn: Das ist, er hat die Seinen alle Tage gelehret von dem verheißenen Messias, und sie erinnert seiner Wohlthaten, wie solches zu dieser Zeit geschieht durch das tägliche Vorlesen, und die Predigten göttlichen Worts.

Das 13. Kapitel der Geschichte der Apostel meldet deutlich, daß im Judenthum eine sonderliche Ordnung gewesen, daß das Gesetz oder die Bücher Mosis und der Propheten alle Sabbathter in der Synagogen und Judenschulen sind vorgelesen worden.

Dieses bestätiget auch das Exempel unsers HErrn Jesu Christi, Luc. 4., denn da er gen Nazareth kömmt, da er empfangen und erzogen war, und am Sabbath in die Schule ging, und daselbst nach seiner Gewohnheit aufstehet und lesen will, da wird ihm das Buch des Propheten Jesaia gereicht, daraus er hat eine tröstliche Predigt gethan, von der gegenwärtigen Gnadenzeit, davon der gemeldete Prophet im 61. Kapitel geweissaget hat.

In Geschichten der heil. Apostel im 15. Kapitel wird abermal von solcher jüdischen Ordnung Zeugniß gegeben. Denn also spricht S. Jacobus: Moses hat von langen Zeiten her in allen Städten, die ihn predigen, und wird alle Sabbathter in den Schulen gelesen. Und zuvor im 13. Kapitel wird auch gedacht der Lectio des Gesetzes und der Propheten, welche in den Judenschulen am Sabbath üblich, und in stetem Gebrauch gehalten worden sind. Aus diesem allen ists offenbar, daß, wie auch vor Christi Geburt, also dergleichen zu seiner Apostel Zeiten und Lehrtagen, bei dem Volk Gottes allwegen der Gebrauch gewesen, daß man das Gesetz Mosis, sowohl der heiligen Propheten Schriften, beides im Tempel zu Jerusalem und in allen Synagogen und Schulen, im ganzen Lande, nicht allein alle Feste und Sabbathte, sondern auch täglich dem Volk vorgelesen und ausgelegt hat, wie auch heutiges Tages bei den Judenschulen, wie sie jetzt geschaffen, zu sehen.

Diese nothwendige Uebung der heil. göttlichen Lehre haben auch der Apostel unsers HErrn Christi, treue Successoren und Nachkommen, die alten christlichen Kirchenlehrer, behalten, und haben die Schriften der heil. Propheten und Apostel in ihren Gotteshäusern und öffentlichen Ver-

\quellseite{32}

sammlungen, beides auf die Fest- und Feiertage, und sonst wo man zusammenkommen ist, der christlichen Gemeine und Versammlung vorgetragen, vorgelesen und ausgelegt. Solches bezeugen:

Justinus, der Lehrer und Märtyrer Christi, welcher gelebet hat um das 150. Jahr nach Christi Geburt, und sagt in der andern Verantwortung für die Christen: Am Sonntage kommen alle, die in Städten und auf dem Lande wohnen, zusammen an einem Ort, und werden der Propheten und Apostel Schriften gelesen, so viel daß eine Stunde leiden kann. Darnach, wenn der, so gelesen hat, aufhöret, erinnert und vermahnet der Vorsteher das Volk, daß sie das Gute, so vorgelesen, wollten nachthun.

Clemens Alexandrinus, der berühmt gewesen ist um das 200. Jahr, sagt im 7. Buch Stromatum, daß der Christen Opfer gewesen sein die Gebete und Lobgesänge, und die Lectionen der Schrift, so für Essens geschehen.

Mit diesen stimmet Tertullianus, der gelebt hat um das 230. Jahr, im 39. Kapitel seines Apologetici adversus Gentes. Wir kommen, spricht er, zusammen zur Predigt und Erklärung der Schrift, wenn uns die Gelegenheit der gegenwärtigen Zeit zwinget, etwas zu vermahnen oder zu erinnern. Wahrlich, wir speisen unsern Glauben mit dem heiligen Worte, wir richten die Hoffnung auf, wir stärken das Vertrauen, und mehren nichts desto weniger die Zucht mit stetigem Einbläuen der Gebote.

Mehr Zeugnisse der ersten Kirchen anzuziehen, ist nicht vonnöthen. Und weil Justinus in Asia, als nämlich in dem Städtlein Nicopoli, Palästina, geboren, desgleichen zu Rom und anderswo in Europa, gelebet hat; Tertullianus in Afrika; Clemens zu Alexandria in Egypten, ist daraus offenbar, daß zu der Zeit an allen Orten, so viel den Gebrauch und die Weise die heilige Bibel zu lesen anlanget, eine wunderliche und sonderliche Einigung gewesen, und Gleichförmigkeit gehalten worden sei.

Von diesem alten und löblichen Brauch haben unter andern Dienern der Kirchen die Lectores auch ihren Namen gehabt, von ihrem Amt, daß sie die heilige Schrift in der christlichen Versammlung haben gelesen und vorgetragen, welchen das Amt zu lesen vom Bischofe mit Uebergebung des heiligen Buchs der Biblia, in Gegenwärtigkeit der ganzen Gemeine befohlen und vertraut ward.

Hieher gehören die Exempel vortrefflicher und heiliger Leute, die fleißig in Gottes Wort sich geübet, und die heilige Biblia täglich gelesen haben.

David, der tapfere Held in Frieden und Streit, der herrliche König in Krieg und Sieg, da er jetzt vor dem Tyrannen Saul flohe, und in unwegsamen Orten in der Irre ging, jetzt beim Könige Gath im traurigen Elende war, oder als er sich für seinem Sohne Absalom, der einen bürgerlichen Aufruhr wider ihn erwecket, fürchtete, hernach, als er draußen im Felde stritt oder daheim das Regiment versahe, war er dennoch auf das Gesetz und Gottes Wort, welches er auf Gottes Befehl angeschrieben und allzeit bei sich hatte, beflissen, ja auch in den alltraurigsten Empörungen und Unruhen hat er selbst sehr viele Psalmen geschrieben und gesungen.

Dergleichen haben gethan der mühselige Salomon, der geschäftige Hiskia, der überladene Josias, die Könige in Juda.

Der edle Kämmerling oder Schatzmeister der Mohren-Königin Candaces, zu welchem, wie glaublich ist, auch das Evangelium kaum von Christo kommen war, brachte die Zeit zu mit der Lection des Propheten Jesaia, nicht allein daheim, unangesehen, daß er beschweret genug und unmüßig war, wegen seines verdrießlichen Amts, belangend die Rechenschaft des Einnahmens und Ausgebens, sondern auch da er über Feld reisete, und auf dem Wagen geführet ward, in der Geschichte der Apostel im 8. Kapitel.

Vom Kaiser Theodosio Secundo melden die Historien, daß er mit eigener Hand die Bücher des Neuen Testaments abgeschrieben, und täglich zur Morgenzeit daheim ein Stück in demselben gelesen hat: Und daß ich dies auch melde, so hat er auch Psalmen und Gebetlein mit seiner Gemahl, Eudoxia genannt, und mit seinen Schwestern zu singen und zu recitiren oder zu sprechen gepflegt.

Vom Alfred, dem Könige aus England, welcher die berühmte Hochschule zu Oxford oder Ochsenfurt gestiftet hat, zeuget die Chronica, daß er die ganze Zeit des Tages in drei

\quellseite{33}

Theile getheilet habe, und täglich acht Stunden mit Lesen, Beten und geistlichen Betrachtungen zugebracht, und denn so viel mit der Regierung des Reichs. Die andern Stunden mit dem, was auf Pflegung des Leibes gehöret.

Eben auf diese Weise hat Carl der Große, der erste römische deutsche Kaiser, sein Gemüth geweidet mit dem Studiren, zuvor aber mit der Lection der heiligen Schrift. Auch hat er unter Essens allezeit etwas aus der heiligen Bibel, dazu auch aus den Schriften des Kirchenlehrers Augustini ihm lassen lesen.

Alphonsus, der hochweise und gelehrte, auch der Lehrer und Gelehrte liebhabende König in Sicilien und zu Neapolis, hat die heilige Biblia mit den Explicationibus oder Erklärungen Nicolai de Lyra, dreizehn mal ausgelesen.

Also haben fleißig die heilige Biblia selber gelesen, und haben ihnen und den Ihren lesen lassen, Herr Christianus III., der christliche löbliche König in Dänemark. Also Frau Margaretha, die heilige, fromme Fürstin zu Anhalt, geborne Schlesische Herzogin von Münsterberg, und andere viel mehr fürstliche und andere christliche Personen.

Aus diesen und dergleichen Exempeln folget, weil alle wahre Christen die Schriften der heiligen Propheten und Apostel herzlich gern täglich mit Andacht gelesen, und lesen haben lassen, daß also öffentlich im gemeinen die heilige Biblia auch billig alle Tage soll vorgelesen werden, und daß jedermann die Lection gern und mit allem Fleiß auch anhören, und derer abwarten soll.

Und dies sind allein drei fürnehme Ursachen, welche alle Menschen, so sich der christlichen Religion rühmen, mit Ernst, gleich als mit einer Posaune und Feldgeschrei, zur Anhörung der Lection der heiligen Biblia, aufmuntern, reizen und treiben sollen, also, daß sich auch niemand daran hindern lassen soll: weil keines Standes, keiner Verrichtung irgend eines Amts, keiner Arbeit, keiner vorfallenden Sache, von den gemeinen und eigenen Händeln, Entschuldigung billig genugsam sein kann, daß sie irgend einen Christen los mache vom Anhören oder ja Lesen der heiligen biblischen Bücher.

Weil nun aber erzähleter und anderen Ursachen wegen, auch ein ehrbarer, hochweiser, wohlbenamter Rath, als eine christliche und gottselige Obrigkeit, angeordnet hat die Lection der heiligen Biblia, ward rathsam, nöthig und nützlich befunden, damit zu fleißiger andächtiger Anhörung, und nöthigem heilsamen Gebrauch göttlichen Worts die lieben Zuhörer angebracht und verursacht würden, daß allezeit ein jedes Kapitel der heiligen Biblia mit einem sonderlichen Prolog oder mit einer Vorrede angefangen, und daß ein jedes Summarium des wohl gemeldeten seligen Herrn M. Dieterichs mit einem Epilogo, Voto, Wunsch oder Gebetlein sollte beschlossen werden. Dieses ins Werk zu richten, hat der ehrwürdige und hochgelehrte Herr Lucas Pollio, das theure Werkzeug Gottes und wohlverdiente Pfarrherr bei unserer Kirchen zu S. Maria Magdalena, damals mir vertraut und auferleget, daß ich solche Prologos oder Vorreden, desgleichen die Epiloge und Beschlüsse ordentlich stellen und fassen sollte. Solches ich auch im Namen Gottes auf mich genommen, und habe diese Arbeit mit Rath und Anweisung wohl gemeldeten meines seligen Pastors und lieben Mitbruders in Christo, auch von vielen Jahren treuen Freundes, und nun wegen seiner geliebten Schwester, meiner ehelichen Hausfrauen, Herrn Schwagers, also zusammengebracht und geschrieben.

Gleichwie aber auch der theure Mann Gottes, Herr Philippus Melanchthon (dessen Auditor und Schüler ich gewesen bin etliche Jahr, und habe seine heilige Cygneas Cantiones angehört, welches ich unter die vornehmsten und größten Wohlthaten rechne, die mir mein lieber GOTT bisher mein Lebenlang aus Gnaden und Barmherzigkeit erzeiget hat), alle seine Scholas, Lehre und Vermahnungen also und dahin richtet, auf daß er die wahre Gottesfurcht in uns, seine Discipul und Zuhörer, pflanzte, und uns zu täglichem, herzlichem, gläubigem Gebete gewöhnte, auch nöthigen Trost weisete, weil er allezeit sagte: wir würden betrübte gefährliche Zeiten und Läufte erleben; (wie denn seine Weissagung nicht ohne sonder große Schmerzen, Kreuz und Noth vieler frommer Herzen bisher erfüllet worden und noch erfüllet wird), also habe ich die gemeldeten Prologe gestellt, daß der meiste Theil Vermahnungen sind zu täglicher Buße, das ist, zur Furcht Gottes, zur Bekehrung, zum Glauben, Gebet und Uebungen der Tugenden im Decalogo oder zehen Geboten Gottes befohlen; viel Prologe und Vorreden habe ich auch accordirt und vereinigt mit den Summarien, daß sie in specie, gewisse nöthige

\quellseite{34}

Lehrpuncte des Texts und der Kapitel mit erklärten und zu besonderen Haupt-Tugenden reizen; damit abermal Gottesfurcht, Gebet und Trost in den Herzen der Zuhörer, der sichern und furchtsamen, der betrübten und fröhlichen erwecket, gepflanzet und erbauet würde. Damit auch die Lehrpuncte, wie die in den Summarien des Herrn Veit Dieterichs seligen berühret werden, mit gläubigem Herzen, und mit freudigem angenehmen Willen, die ausständige Zeit des Tages, und das ganze Leben über, im Gehorsam Gottes zu leben, haften und Nutzen bringen möchten, derwegen, wie alle ordentlichen Predigten mit einem christlichen Wunsch beschlossen worden, daß das angehörte Wort Gottes bleibe und bestehe in einem guten und feinen Herzen, und hundertfältige Früchte in Geduld und allen Gott wohlgefälligen Tugenden bringe, also habe ich auf die gemeldete Summaria ein Votum und Beschluß gestellet, zur Application, wie es dieselben erfordern, und der gütige Leser solches wohl erkennen wird können. Und habe mich beflißen, so viel möglich, sonderlich um der Einfältigen willen, daß ich gewöhnliche, deutliche und verständliche Worte gebraucht.

Erwähnte Arbeit ist also, durch Gottes väterlichen reichen Segen gerathen und angewandt worden, daß diese Lectionen nicht allein fleißig und andächtig unsere christliche Kirchenversammlung angehört, sondern daß auch viele von unserer gottseligen Bürgerschaft, nachdem der Geist Gottes in ihnen gewirket, einer jetzt des Stückes, bald ein anderer eines andern, Abschrift begehrt haben; daraus denn auch weiter folget, daß von etlichen Jahren her ist angehalten worden, daß solches alles in Druck verfertiget würde, damit solche nöthige und heilsame Lehren und Erinnerungen auch zu Hause den Kindern und dem Gesinde zu Nothdurft und Nutzen möchten vorgelesen werden.

Dieses (das den begehrten Abdruck anlanget), edle, gestrenge, ehrenveste, wohlweise, wohlbenamte großgünstige Herren, an eure Gnaden und Herrschaften billig, aus schuldiger Pflicht und Unterthänigkeit, gelangen hat sollen, und da solches in Gehorsam von mir geschehen, haben eure Gnaden und Herrschaften, auf darüber vor gehaltenen Rath, günstig bald bewilliget, daß solches Buch in dieser Stadt-Druckerei solle publiciret und verfertigt werden.

Daß nun solches also gerathen, danke ich von Herzen dem treuen lieben Gott, der es also gesegnet, daß es eure Gnaden und Herrschaften und vielen frommen Herzen gefallen; und lasse auch nun im Namen meines Herrn Christi diese Arbeit, so gut, als er sie gegeben, in offenem Druck ausgehen, auf daß also diese löbliche und vortreffliche Stadt und christliche Gemeine einen Dank von mir habe für die vielfältigen und großen Wohlthaten, die ich allhier empfangen habe: anfänglich in meiner Jugend, in der ich von der Neiße (die meine liebe Geburts-Stadt ist), hieher ad studia pietatis et bonarum artium kommen bin. Da ich denn bei ehrbaren und vornehmen Leuten ehrliche und gar bequeme Paedagogias gehabt, und denn darauf von einem ehrbaren, hochweisen wohlbenamten Rath, unangesehen, daß ich ein Fremder war, zu ihrer Gnaden und Herrschaften Alumno angenommen, und zu Wittenberg Verlag gehabt durch Beförderung des gestrengen, edlen, ehrenvesten und wohlbenamten Herrn Anton Band, Hauptmann etc., in Gott seligen, sowohl auch durch Hülfe des edlen, ehrenvesten und wohlbenamten Herrn Hans Morenberg, des Rathes, und Scholarchae oder Schul-Herrns, auch seligen, welches ich mit dankbarem Herzen und Gemüthe wohlgemeldeten Herren und ihren Erben zu Ehren billig rühme, mit herzlichem Wunsch, daß der barmherzige und gütige Gott diese Wohlthaten an wohlgemeldeten Herren, meinen großen Patronen und Förderern, ehrbaren Geschlechtern und Nachkommen, allen und jeden, allzeit, bis an den jüngsten Tag, reichlich, an Leib und Seele, zu zeitlicher und ewiger Wohlfahrt, vergelten wolle. Darnach, auf daß ich auch für die viele Ehre, Freundschaft und große Wohlthaten, so diese 31 Jahre, weil ich in dem ehrwürdigen Kirchenamt diene, mir, meinen lieben, den ersten und andern Eheweibern und Kindern, von vielen frommen Herzen erzeiget sind worden, mich etlicher Maßen dankbar erfinden lasse. Denn ich ja auch die wiederum ehren soll, welche mir Ehre und Gutes bewiesen haben. Wie solches die Vernunft lehret, daß es billig sei, und die heilige Schrift, daß es geschehen solle, ernstlich befiehlet, mit Ankündigung greulicher Strafen gegen die, so undankbar erfunden werden.

Diese Vorreden und Beschlüsse aber über den ersten Theil der heiligen Biblia, nämlich über das Alte Testament, haben eure Gestrengen und Herrschaften, mit schuldiger Reverenz und Ehrerbietung ich sollen und wollen dediciren und zuschreiben, damit ich für meine Person, Gott zu

\quellseite{35}

Ehren und Dank, öffentlich wahrhaftig Zeugniß mit gebe, daß eure Gnaden und Herrschaften auch rechte Pfleger und Vormünder, und treue Säugammen und Wärterinnen seien der Kirchen des Herrn Christi, so unter Euer Gnaden und Herrschaften löblichen Gemeine und lieben Unterthanen, ihm der Sohn Gottes sammlet; und daß Eure Gnaden und Herrschaften als christliche Herren und Regenten auch wahre Theophili und Gottes Liebhaber seien. Wie denn der heilige Evangelist Lucas sein Evangelium und Scripta auch seinem Theophilo dediciret und zugeschrieben hat.

Gelanget demnach an Eure Gnaden und Herrschaften sämmtlich und sonderlich, mit unterthänigem und demüthigem Fleiße, meine Bitte, Eure Gnaden und Herrschaften wollen ihnen diese meine wohlgemeinte Dedication, auch Declaration und Erzeigung meiner schuldigen Dankbarkeit günstig gefallen lassen, die mit gutem Willen annehmen; auch mein, meines lieben Eheweibes, ihres Sohnes und meiner Kinder, und sonderlich nach meinem tödtlichen Abgang, Vater und Förderer sein und bleiben, und also uns in ihre Gunst und Patronium, zu väterlicher Beförderung, lassen befohlen sein. Das wollen wir mit unserm andächtigen Pater noster und Gebet, für Euer Gnaden und Herrschaften und gemeiner Stadt Wohlfahrt, Gedeihen, zeitlichen und ewigen Segen treulich zu flehen, nach höchstem und besten Vermögen zu jeder Zeit beflissen sein.

Helffe der treue, barmherzige und gütige Gott, daß alle Menschen und Herzen, welche die Lection der heiligen Biblia, der Summarien darüber, und dieses Buches anhören werden, daß sie je länger je mehr erleuchtet werden, und daß sie ihres christlichen Berufes gewisser, in gottseligem Wandel fleißiger, zum Creuz williger, in Anfechtung getrösteter, zum Gebet ernster und brünstiger und zu allen guten Werken geschickter werden, und also zu Gottes Ehre, und allen, die solche Lection anhören werden, zur Seligkeit gedienet werde in Christo Jesu, unserm Herrn, bis zu dem gewünschten jüngsten Tage, Amen.

Auch bitte ich herzlich, der allmächtige Gott und treue Vater wolle ihm E. G. und Herrschaften, diese löbliche Stadt, Kirchen und Schulen und die ganze ehrbare Gemeine, in seine Gnade und Schutz lassen befohlen sein; wolle Eure Gnaden und Herrschaften lange Zeit frisch und gesund, bei glückseligem Regiment, gesegneter Haushaltung und aller Wohlfahrt Leibes und der Seelen, um seines geliebten Sohnes, unsers Herrn und Heilands Jesu Christi willen, für und für erhalten, schützen und bewahren, und nach vollendter Ritterschaft dieses elenden Lebens, die himmlische Ehren-Krone ewiger Freude und Seligkeit verleihen, Amen. Datum im Jahr Jesu Christi, unsers Herrn und Seligmachers Geburt 1596, den 18. Februar, an welchem Tage das auserwählte Rüstzeug Gottes, der Herr Dr. Martin Luther, seliger Gedächtniß, als der treue Gott sein vorgehabtes Werk mit großer Kraft und Stärke durch ihn ausgerichtet, zu Eisleben, in seiner Geburtsstadt, vor 50 Jahren in Christo, dem Herrn, fröhlich und selig entschlafen, seines Alters im 63. Jahr.

\begin{unterschrift}

E. G. und Herrschaften\\Unterthäniger Diaconus der Kirchen zu S. Maria Magdalena Franciscus Bierling.

\end{unterschrift}

\chapter{Erklärung alterthümlicher und dunkler Ausdrücke}

\quellseite{36}

Aesern (nach Luther's Schreibart: euern), Spr. 17, 9.; wiederholen, wiedererregen.

Astersabath, Luc. 6, 1.; der erste Sabbath nach dem zweiten Tage des Osterfestes.

Aufmutzen, Sir. 13, 27.; hervorziehen, ans Licht bringen, groß machen.

Aufsetzen, 2 Kön. 18, 29.; 19, 10.; 2 Chron. 32, 15.; überreden, betrügen.

Aufsetzen, Mich. 12, 9.; aufstacheln.

Aufzug, 1 Macc. 8, 26.; Ausflucht.

Barte, Ps. 74, 6.; ein breiteres Beil, eine Art.

Crethi und Plethi, 2 Sam. 8, 18.; Läufer und Scharfrichter, die Leibwache Davids.

Dreiling, Jes. 40, 12.; ein drei Finger breites Maß.

Ehrlich, 4 Mos. 16, 2.; Luc. 14, 8.; vornehm, angesehen.

Endlich, Spr. 21, 5.; 22, 29.; Luc. 1, 39.; was zu Ende kommt und zu Ende bringt, eilig, emsig, tüchtig.

Enthalten, Ps. 51, 14.; aufrecht halten. Sich enthalten, Apok. 1, 13.; sich aufhalten.

Erwegen (erwägen), sich des Lebens, Weish. 17, 15.; 2 Cor. 1, 8.; das Leben verloren gehen.

Fast, sehr, überaus.

Ferge, Hes. 27, 27.; Fährmann, Schiffer.

Feser (Fäser), Jes. 5, 7.; die Haare oder Fasern an den keimenden Reben oder andern Pflanzen, welche in die Erde gesenkt, zuerst Wurzel fassen.

Fladdernholz, Hes. 27, 5.; Mastholz.

Geilen, Luc. 11, 8.; auf eine unverschämte Art, mit geifernder Stimme, um etwas betteln.

Geren, Hes. 16, 8.; Hagg. 2, 13.; die Schleppe des Kleides, der Schooß; auch der Flügel, der Zipfel des Mantels.

Glum, Jes. 32, 2.; schlammig, unrein; im Gegensatz von lauter, B. 14.

Gnätze, 3 Mos. 14, 56.; Krätze.

Greten, Hes. 16, 25.; aufspreizen, die Füße auseinander sperren.

Hellig, Jer. 2, 25.; athemloß, keuchend, lechzend.

Hengel, 1 Kön. 22, 34. Nach Luther's Erklärung: „Da das Schwert anhängt, von der Achsel über bis auf die Hüfte.“ Andere erklären: die Gelenke und Fugen am Panzer.

Hümpler, Spr. 26, 10.; ein Stümper, schlechter Arbeiter.

Igel (Egel), Spr. 30, 15.; der Blutigel.

Keulig, 1 Kön. 7, 41.; kugelig, rund.

Köden, Jes. 28, 7.; herausreißen.

Kogel, Hes. 23, 15.; langherabhängende Kopf binden.

Kolk, 3 Mos. 11, 36.; eine Grube, worin sich Wasser sammelt.

Koller, Jes. 3, 23.; Halskragen, Brusttuch; auch Leibgürtel, Apok. 19, 12.

Kollern, 1 Sam. 21, 13.; sich unsinnig und toll gebärden, sich auf dem Boden herumwälzen.

Koller, 2 Kön. 8, 15.; Decke, Bettdecke.

Krebs, Weish. 5, 19.; Eph. 6, 14.; 1 Thess. 5, 8.; Brustharnisch, Panzer (von der krebsförmigen Gestalt so genannt).

Kreuel, 2 Mos. 27, 3.; 1 Sam. 2, 13.; Fleischgabel.

Legel, 1 Sam. 16, 20.; Jer. 13, 12.; 48, 12.; nach Adelung: ein rundes, hölzernes Gefäß, in Gestalt einer Tonne.

Leibzucht, 1 Macc. 10, 54.; Leibrente, Unterhalt auf Lebenszeit.

Librarei, 2 Macc. 2, 13.; Bücherei, Büchersammlung.

Läden (läden, nach Luther's Schreibweise: leden), Hiob 21, 11.; Ps. 29, 6. u. ö.; springen, hüpfen; Apok. 9, 5.; mit den Füßen ausschlagen.

Löden, Hof. 7, 14.; heulen, plärren.

Lotterbube, Apok. 17, 18.; Schwätzer, unnützer Mensch.

Maul, 2 Kön. 13, 29.; 18, 9.; 1 Kön. 1, 33.; Maulthier.

Melechet des Himmels, Jer. 7, 18. u. ö.; Himmelskönigin, wie Luther in den ersten Ausgaben selbst übersetzt hatte.

Meye, Bar. 6, 8., im schlimmen Sinne; Richt. 5, 30., im guten, eine ehrliche Frauenperson.

Müller, Pred. 12, 3. 4.; Mühlmagd; die Stimme der Müllerin: der schwache Ton der Mühle.

Nehmen, 4 Mos. 11, 5.; eine Art Kürbis. (Nach dem Hebr.: Wassermelonen.)

Nadeln, 1 Kön. 7, 31.; Nadeln, Hesek. 1, 11.; Büchsen, erhabene Zierrathen von Metall, bes. auf Schilden.

Racha, Matth. 5, 22.; loser, nichtswürdiger Mensch.

Rachus, Jer. 15, 13.; 17, 3.; Hes. 23, 46.; Gegenstand feindlicher Plünderung.

Rechtfertigen, Apok. 12, 19.; vor Gericht laden, anklagen und bestrafen.

Saff, 2 Mos. 12, 8.; 4 Mos. 9, 11.; eine Brühe von bittern Kräutern.

Sangen, Ruth 2, 14.; 1 Sam. 17, 17.; gesengte, geröstete Aehren.

Schauer (Scheuer), Hiob 27, 18.; Obdach gegen böse Witterung.

Schemen, Ps. 39, 7.; Spr. 27, 19.; Schattenbild.

Scheuer, Sir. 50, 10.; ein schimmerndes Gefäß; Becher.

Schier, Ps. 119, 87.; Apok. 17, 15.; 1 Tim. 3, 14.; bald, schnell.

Schröter, Jer. 48, 12.; Leute, welche einen Stoff ausschroten, d. i. zertheilen und stückweise veräußern.

Seettig, 1 Tim. 6, 4.; (von Seuche) krank, siech.

Sichermahl, 1 Sam. 20, 20.; sicheres festes Ziel (im Gegensatz des Schießens ins Ungewisse).

Sinflut, 1 Mos. 7, 6.; Sir. 39, 27. u. ö.; große, weitverbreitete Flut.

Sponde, Amos 3, 12.; Bettgestelle.

Spunden, 1 Kön. 6, 9. u. ö.; decken, täfeln.

Stöcken, Weish. 2, 19.; in den Stock legen, zwingen (Ps. 105, 18.); peinigen.

Sträter, Hof. 6, 9. nach Luther's Erklärung: Strauchdiebe, Strauchräuber.

Tartische, 1 Kön. 10, 17.; ein langer Schild.

Tendchen, 5 Mos. 14, 5.; eine Art von Gemsen oder wilden Ziegen.

Thelding, a) Jagdübung, der Termin zur gerichtlichen Handhabung; b) die Verhandlungssache; dann unnützes Hin- und Herreden, Geschwätz. Hiob 35, 16.; Hes. 22, 28. Daher Narrenheiding Eph. 5, 4., loses, unnütziges Gerede.

Thramen, 1 Kön. 6, 6.; Balken.

Thürstig, thürstlich, 1 Mos. 34, 25.; Hiob 12, 6. u. ö.; kühn, beherzt, dreist.

Treusen, Mich. 2, 6.; schelten, strafen.

Unrath, Matth. 26, 8.; nach Luther's Erklärung: Es ist Schade drum, sie sind unräthlich damit umgegangen.

Verschleifen, 2 Sam. 14, 14.; schleifen, sich verlaufen.

Verschleißen, Sir. 14, 18.; zerreißen, verschleißen, verderben.

Versprechen, a) Richt. 9, 23.; verwünschen, verfluchen; b) Marc. 7, 2.; tadeln, sich über etwas aufhalten.

Wind, Spr. 30, 31.; Windhund, Windspiel.

Zauen, sich, 2 Sam. 5, 24.; sich sputen, eilen.

Zerlecken, zerlecken, Pred. 12, 6.; leck werden, Ritze bekommen.
% Scan 37: no retained apparatus prose; source blocks omitted.
% Scan 38: no retained apparatus prose; source blocks omitted.
% Scan 39: no retained apparatus prose; source blocks omitted.
\part{Das erste Buch Mose}

\chapter{Das 1. Kapitel}

\section{Vorrede}

\quellseite{40}

Aller guten und nützlichen Werke Anfang, Mittel und Ende stehet allein in Gottes Händen. Derowegen weil wir heute die Lection oder das Vorlesen der heiligen Biblia anfangen, so laßt uns herzlich beten, der allmächtige und gütige Gott wolle unsere Augen öffnen, und unsere Herzen erleuchten, daß wir die Wunder sehen in seinem Gesetz. Wie der königliche Prophet David im 119. Psalm redet; das ist: daß wir also sein heilig Wort täglich anhören, lernen, behalten und gebrauchen, auf daß wir seinen väterlichen Willen und seine Gnade und Barmherzigkeit daraus erkennen, und daß Glaube, Liebe, Trost, Geduld, Zucht, Gehorsam und Unterthänigkeit in uns allen dadurch gepflanzet werde, wachse und zunehme. Wer nun solches begehret, der spreche von Herzen, Amen. Das hilf uns allen, o lieber Gott und Vater, Amen. Darauf höret an in Gottes Namen das erste Kapitel des ersten Buchs Mose.

\section{Schöpfung der Welt}

a Geist) Wind ist zuweilen noch nicht gewesen, dann muß es den heiligen Geist deuten.

b Gut) das ist, nützlich, fein, löblich.
\quellseite{41}

c Unterthan) was ihr bauet und arbeitet auf dem Lande, das soll nur eigen sein, und die Erde soll euch hierin dienen, tragen und geben.

d Speise für die Menschen und Thiere etc.

\section{Summarium}

Ist eine Historia, wie Gott durch sein ewiges, allmächtiges Wort Himmel und Erden, sammt allen Kreaturen erschaffen, und am sechsten Tage den Menschen, der solches alles genießen und gebrauchen sollte, gemacht, und zum Herrn darüber gemacht habe. V. 27.

Auf daß wir lernen Gottes Güte in solchem Geschöpf erkennen, und unser Herz in allerlei Nöthen damit trösten, daß wir so einen allmächtigen und gnädigen Gott haben, der uns so reichlich versorget, und so väterlich mit allem, was wir bedürfen, bedacht hat. Solchen gnädigen Willen Gottes haben auch die Heiden aus der Schöpfung und Kreatur erkennet, und derselben Gott für gut und gnädig gehalten. Aber Gottes Wort hält den Christen einen höhern und größern Trost für, dabei sie müssen schließen, daß Gott gnädig sei; sintemal er nicht allein alles den Menschen zu Nutz erschaffen, sondern auch seinen eingebornen Sohn zum Opfer für die Sünde, und in den Tod hat gegeben. Darum sollte niemand an solches gnädigen Gottes Hülfe, Gunst und Rettung zweifeln, es sei gleich die Noth so groß sie immer wolle.

Merke hie den Artikel unsers Glaubens. Gott ist, ehe einige Kreatur ist, das ist, er ist ewig, und hat bei sich das Wort, das ist auch ewig und Gott gleich; denn es ist, ehe alle Kreaturen sind. Und weil durch solches Wort alle Kreaturen erschaffen werden, muß folgen, daß solches Wort allmächtig ist. Zum dritten nennt Moses hie den Geist des Herrn, der alles Geschöpf lebendig macht und erhält. Daß also Moses klar anzeiget, GOTT Vater, GOTT Sohn, oder das Wort, und GOTT der Heilige Geist sei der einige, ewige, allmächtige GOTT, der Anfang und das Ende, wie unser Christen-Glaube ausweiset.

\section{Gebet}

Wir danken Gott herzlich für die Erschaffung aller Kreaturen; der wolle ferner die Werke seiner Hände segnen und erhalten, und uns in seiner Gnade zum ewigen Leben bewahren, Amen.
\chapter{Das 2. Kapitel}

\section{Vorrede}

\quellseite{42}

Gott hat uns Menschen erschaffen, daß wir ihm sollen gleich sein in Gerechtigkeit, Heiligkeit, Barmherzigkeit, Wahrheit und in andern Tugenden, die in ihm wesentlich leuchten, und von welchen die heiligen zehn Gebote lehren. Aber, ach weh uns allen, wie ungleich sind wir Gott! Das lässet uns erkennen, und Gott abbitten, im wahren Glauben und Vertrauen auf das Versöhn-Opfer Jesu Christi, unsers HErrn. Wir wollen auch, so lange wir leben, unsere fürnehmste Sorge sein lassen, daß wir durch Hülfe des heiligen Geistes Gott, unserm Schöpfer, von Tage zu Tage gleichförmiger werden, und also gegen uns selbst züchtig, gegen unsern Nächsten gerecht, und gegen Gott gottselig leben, Tit. 3. Amen. Das soll auch erinnern das 2. Kapitel.

\section[Vom Sabbath, des Menschen wesentlichen Stücken, Paradiesgarten, Gottes]{Vom Sabbath, des Menschen wesentlichen Stücken, Paradiesgarten, Gottes Gebot, Ehestand}





a Pison) ist das große Wasser in Indien, das man Ganges heißt.

b Hevila) ist Indien-Land.

c Gihon) ist das Wasser in Aegypten, das man Nil heißt.


d Hiddekel) ist das Wasser in Assyrien, das man Tigris heißt.

e Euphrath) aber ist das nächste Wasser in Syrien, das man Euphrat heißt.




*Um ihn sey) das ist, kein Thier nahm sich des Menschen an, um ihn zu sein, daß ihm hülfe sich mehren und nähren etc.
\section{Summarium}

\quellseite{43}

Ist eine Historia, wie Gott durch sein ewiges, allmächtiges Wort Himmel und Erden, sammt allen Kreaturen erschaffen, und am sechsten Tage den Menschen, der solches alles genießen und gebrauchen sollte, gemacht, und zum Herrn darüber gemacht habe. V. 27.

Auf daß wir lernen Gottes Güte in solchem Geschöpf erkennen, und unser Herz in allerlei Nöthen damit trösten, daß wir so einen allmächtigen und gnädigen Gott haben, der uns so reichlich versorget, und so väterlich mit allem, was wir bedürfen, bedacht hat. Solchen gnädigen Willen Gottes haben auch die Heiden aus der Schöpfung und Kreatur erkannt, und derhalben Gott für gut und gnädig gehalten. Aber Gottes Wort hält den Christen einen höhern und größern Trost für, dabei sie müssen schließen, daß Gott gnädig sei; sintemal er nicht allein alles den Menschen zu Nutz erschaffen, sondern auch seinen eingebornen Sohn zum Opfer für die Sünde, und in den Tod hat gegeben. Darum sollte niemand an solches gnädigen Gottes Hülfe, Gunst und Rettung zweifeln, es sei gleich die Noth so groß sie immer wolle.

Merke hie den Artikel unsers Glaubens. Gott ist, ehe einige Kreatur ist, das ist, er ist ewig, und hat bei sich das Wort, das ist auch ewig und Gott gleich; denn es ist, ehe all Kreaturen sind. Und weil durch solches Wort alle Kreaturen erschaffen werden, muß folgen, daß solches Wort allmächtig ist. Zum dritten nennt Moses hie den Geist des Herrn, der alles Geschöpf lebendig macht und erhält. Das also Moses klar anzeiget, GOTT Vater, GOTT Sohn, der das Wort, und GOTT der heilige Geist sei der einige, ewige, allmächtige GOTT, der Anfang und das Ende, wie unser Christen-Glaube ausweiset.

\section{Beten}

Wir danken Gott herzlich für die Erschaffung aller Kreaturen; der wolle ferner die Werke seiner Hände segnen und erhalten, und uns in seiner Gnade zum ewigen Leben bewahren, Amen.

d Speise für die Menschen und Thiere etc.

c Unterthan) was ihr bauet und arbeitet auf dem Lande, das soll euer eigen seyn, und die Erde soll euch hierin dienen, tragen und geben.
\quellseite{44}
allen Kreaturen; da kann die Vernunft anders nicht richten, denn es sei keine Hülfe mehr im Himmel und auf Erden. Aber hier aus diesem Exempel sollen wir lernen, daß Gott auch helfen wird, wenn er gleich von allen Kreaturen verlassen sind. Doch thut er solche Hülfe um seines Sohnes willen, den er alhier dem menschlichen Geschlecht zu senden zugesaget etc.

d Derselbe hieß ist das erste Evangelium und Verheißung von Christo, geschehen auf Erden, daß er sollte Sünde, Tod und Hölle überwinden, und uns von der Schlangen Gewalt selig machen. Daran glaubet Adam mit allen seinen Nachkommen; davon er zum Christen, und selig worden ist von seinem Fall.

e (Stechen) plagen, kreuzigen und martern. Denn so gehets auch: Christus zertritt dem Teufel seinen Kopf, das ist, sein Reich, des Todes, Sünd' und Hölle. So sticht ihn der Teufel in die Fersen, das ist, er tödtet und martert ihn und die Seinen leiblich.

f Heva) Kai heißt Leben; daher kömmt Heva oder Hava, Leben oder Lebendige.

\section{Summarium}

Ist eine Historia, wie der Mensch durch den Teufel von Gottes Wort abgeführet, und in die Sünde und Strafe Gottes durch den Ungehorsam gefallen sei, uns zur Warnung, daß wir lernen fest an Gottes Wort halten, und uns vor Sünden hüten, welche vom Teufel ihren Anfang hat, und den Tod, sammt allem andern Jammer, über die Menschen gebracht hat. Dagegen soll man auch den Trost merken, welchen hie Gott Adam und Eva gibt, daß des Weibes Same, Christus, der Schlange den Kopf zertreten, das ist, dem Teufel seine Macht nehmen, und durch sein Leiden uns aus dem Tode und Sünde zur Gerechtigkeit und ewigen Leben helfen soll. Solchen Trost sollen wir fest behalten und uns das Fersenstechen des Teufels, das ist, allerlei Widerwärtigkeit dieses Lebens nicht ärgern lassen; denn im Paradies hat es angefangen, und muß für und für also gehen. Der Teufel gönnet uns die Seligkeit nicht; wiederum will Gott uns in Sünden nicht verderben lassen, und schaffet uns Hülfe durch seinen Sohn, der für uns am Kreuze stirbt und bezahlet. Solchen Trost und Hülfe sollen wir mehr achten, denn des Teufels Fluch und Feindschaft.

\section{Gebet}

Lob, Ehr, Preis und Dank sei Jesu Christo, unserm Heilande, der des Teufels Werk zerstöret hat; der erhalte uns im wahren Glauben und guten Gewissen zum ewigen Leben, Amen.
\chapter{Das 4. Capitel}

\section{Vorrede}

\quellseite{45}

Wie Kain ein Mörder worden; also ist kein Mensch auf Erden, er trägt bei sich ein Stück von einem Mörder. Denn gleich wie der Teufel ein Lügner und Mörder ist von Anfang, Joh. 8., das ist, von seinem Falle her: also hat er unter andern auch diese zween giftige Pfeile in aller Menschen Herzen, Fleisch und Blut geschaffen, und hat ihnen das Schandmal angehängt, daß sie nicht allein von Natur zu Lügen, sondern auch zu Zorn, Haß und Todtschlag geneigt sind. Derwegen ein jedes gute Achtung auf sich gebe, und lasse die Sünde ja nicht herrschen in seinem sterblichen Leibe, ihr Gehorsam zu leisten in ihren Lüsten. Röm. 6. Dies lehret auch das 4. Capitel.

\section{Kains Brudermord. Seine Nachkommen}

a Ei, Gott sei gelobt, da habe ich erlanget den HErrn, den Mann, den Samen, der dem Satan, der höllischen Schlange, den Kopf zertreten soll; der wird thun.

b Hebräisch lautet Thür so viel, als man das offen stehet oder aufgethan wird. Marc. 7, 34.: Sephahtah, thue dich auf etc. Und ist die Meynung: die Sünde lieget und ruhet, wie ein Docht lieget und ruhet. Aber sie lieget in der Thür, das ist, sie wird offen stehen, oder offenbar werden, ob der Sünder wohl eine Zeit lang ihrer dahin gehet, als sei die Sünde todt und schlaf.

c Redet mit Habel das ist, Scham halber mußte er sich äußerlich stellen und reden mit seinem Bruder, weil er gestraft ward, ob er wohl im Herzen ihn zu tödten gedachte. Also ist Kain aller Heuchler und falscher Heiligen Vater.
\quellseite{46}
*c. 5, 6.

d fing man an) nicht daß zuvor nicht auch Gottes Name wäre gepredigt; sondern nachdem durch Kains Bosheit der Gottesdienst gefallen war, ward er diesmal wieder aufgerichtet und irgend ein Altärlein gebauet, dahin sie sich versammleten, Gottes Wort zu hören und zu beten.

\section{Summarium}

Das ist eine Historia, wie der leibliche Segen mit Adam und Eva gegangen, und darneben die Sünde sich auch gewaltig an Kain erzeuget habe, der seinen Bruder Abel erschlug.

Solches Exempel gehet noch immerdar in der Welt. Denn der Abel, das ist die rechte christliche Kirche, die auf Gottes Wort und seine Barmherzigkeit in Christo sich gründet, muß sich von Kain, das ist, von der falschen Kirche und den Werktheiligen leiden. Aber es soll doch nicht ungerochen bleiben. Denn Gott will das unschuldige Blut suchen, und dräuet gar ernstlich, daß alle, so sich hier vergreifen, mit bösem Gewissen und allerlei Unglück ihr Leben lang sollen geplaget werden.

Hier ist sonderlich zu merken die Predigt, die (der HErr) Adam dem Kain thut, und warnet, er soll der Sünde den Raum nicht zu lang lassen, sondern an sich halten, sie würde ihn sonst verführen. Bist du fromm, spricht er, so bist du angenehm, das ist, halt dich in Gottesfurcht, und hüte dich für der Sünde, so wird Gott ob dir halten. Bist du aber nicht fromm, das ist, wenn du der Sünde folgest, und willst auf Gottes Furcht und Wort nicht sehen, so ruhet die Sünde, das ist, du hast Lust und Freude daran; aber es währet nicht lange, sie schläft in der Thür, an einem sehr unruhigen Ort, da muß sie aufwachen. Alsdann folget unermeßliche Angst, Schrecken, Jammer und der Tod selbst; wie man hier an Kain und an Judas stehet. Darum warnet (der HErr) Adam, laß du ihr nicht ihren Willen, das ist, wehre dich, widerstehe der Sünde. Und wie es Paulus heißt, würge sie. Da soll uns zu helfen des Weibes Saame, der will bei uns halten, und uns wider die Sünde helfen. Wir für uns würden zu schwach sein.

\section[Notum]{\unsicher{Notum}}

Helffe der treue Gott, daß ein jedes unter uns seine Natur und böse Zuneigung erkenne, die Sünde meide und meide, und dagegen christlich und ehrbar lebe, zu Lob und Preis seinem heiligen Namen, Amen.

\chapter{Das 5. Kapitel}

\section{Vorrede}

Wenn ein Mensch noch so lange lebet, so muß er doch endlich sterben. Dies lehren auch die gewöhnlichen Sprüchworte: Der Tod ist eine unvermeidliche Noth; Sterben ist unser Erbe; und: Wider Todes Gift und Macht, hilft keines Krautes Kraft und Saft. Derwegen laßt uns wachen und beten, damit, wenn wir auch durch den zeitlichen Tod aus dieser Welt scheiden, wir im Glauben und mit gutem Gewissen einschlafen, und am jüngsten Tage zum ewigen freudenreichen Leben mit allen Auserwählten von Todten auferstehen. Dazu höret an das 5. Kapitel.

\section[Geschlechtsregister der Patriarchen von Adam an bis auf]{Geschlechtsregister der Patriarchen von Adam an bis auf Noah}
\chapter{Das 4. Kapitel}

\section{Vorrede}

\quellseite{47}

Wie Kain ein Mörder worden: also ist kein Mensch auf Erden, er trägt bei sich ein Stück von einem Mörder. Denn gleich wie der Teufel ein Lügner und Mörder ist von Anfang, Joh. 8., das ist, von seinem Falle her: also hat er unter andern auch diese zween giftige Pfeile in aller Menschen Herzen, Fleisch und Blut geschaffen, und hat ihnen das Schandmaal angehängt, daß sie nicht allein von Natur zu Lügen, sondern auch zu Zorn, Haß und Todtschlag geneigt sind. Derwegen ein jedes gute Achtung auf sich gebe, und lasse die Sünde ja nicht herrschen in seinem sterblichen Leibe, ihr Gehorsam zu leisten in ihren Lüsten. Röm. 6. Dies lehret auch das 4. Kapitel.

\section{Kains Brudermord. Seine Nachkommen}

a Ei, Gott sei gedankt, da habe ich erlanget den HErrn, den Mann, den Samen, der dem Satan, der höllischen Schlange, den Kopf zertreten soll; der wirds thun.

b Hebräisch lautet Thür so viel, als was offen stehet oder aufgethan wird, Marc. 7. 34.: Hephatha, thue dich auf etc. Und ist die Meinung: die Sünde lieget und ruhet, wie ein Ochs lieget und ruhet. Aber sie lieget in der Thür, das ist, sie wird offen stehen, oder offenbar werden, ob der Sünder wohl eine Zeit lang sicher dahin gehet, als sei die Sünde todt und schlafe.

* Redet mit seinem Bruder) das ist, Scham halber mußte er sich äußerlich stellen und reden mit seinem Bruder, weil er gestraft ward, ob er wohl im Herzen ihn zu tödten gedachte. Also ist Kain aller Heuchler und falscher Heiligen Vater.
\quellseite{48}

d Fing man an) nicht daß zuvor nicht auch Gottes Name wäre gepredigt; sondern nachdem durch Kains Bosheit der Gottesdienst gefallen war, ward er dazumal wieder aufgerichtet und irgend ein Altärlein gebauet, dahin sie sich versammleten, Gottes Wort zu hören und zu beten.

\section{Summarium}

Das ist eine Historia, wie der leibliche Segen mit Adam und Eva gegangen, und darneben die Sünde sich auch gewaltig an Kain erzeuget habe, der seinen Bruder Abel erschlug.

Solches Exempel gehet noch immerfort in der Welt. Denn der Abel, das ist die rechte christliche Kirche, die auf Gottes Wort und seine Barmherzigkeit in Christo sich gründet, muß sich von Kain, das ist, von der falschen Kirche und den \unsicher{Werth[e}iligen] leiden. Aber es soll doch nicht ungerochen bleiben. Denn Gott will das unschuldige Blut suchen, und dräuet gar ernstlich, daß alle, so sich hier vergreifen, mit bösem Gewissen und allerlei Unglück ihr Leben lang sollen geplaget werden.

Hier ist sonderlich zu merken die Predigt, die (der HErr) Adam dem Kain thut, und warnet, er soll die Sünde den Zaum nicht zu lang lassen, sondern an sich halten, sie würde ihn sonst verführen. Bist du fromm, spricht er, so bist du angenehm, das ist, halt dich in Gottesfurcht, und hüte dich für der Sünde, so wird Gott ob dir halten. Bist du aber nicht fromm, das ist, wenn du der Sünde folgest, und willst auf Gottes Furcht und Wort nicht sehen, so ruhet die Sünde, das ist, du hast Lust und Freude daran; aber es währet nicht lange, sie schläft in der Thür, an einem sehr unruhigen Ort, da muß sie aufwachen. Aldann folget unermäßliche Angst, Schrecken, Jammer und der Tod selbst; wie man hier an Kain und an Judas siehet. Darum warnet (der HErr) Adam, laß du ihr nicht ihren Willen, das ist, wehre dich, widerstehe der Sünde. Und wie es Paulus heißt, würge sie. Da soll uns zu helfen des Weibes Same, der will bei uns halten, und uns wider die Sünde helfen. Wir für uns würden zu schwach sein.

\section{Beten}

Helfe der treue Gott, daß ein jedes unter uns seine Natur und böse Zuneigung erkenne, die Sünde meide und meide, und dagegen christlich und ehrbar lebe, zu Lob und Preis seinem heiligen Namen, Amen.

\chapter{Das 5. Capitel}

\section{Vorrede}

Wenn ein Mensch noch so lange lebet, so muß er doch endlich sterben. Dies lehren auch die gewöhnlichen Sprichworte: Der Tod ist eine unvermeidliche Noth; Sterben ist unser Erbe; und: Wider Todes Gift und Macht, hilft keines Krautes Kraft und Saft. Derwegen laßt uns wachen und beten, damit, wenn wir auch durch den zeitlichen Tod aus dieser Welt scheiden, wir im Glauben und mit gutem Gewissen einschlafen, und am jüngsten Tage zum ewigen freudenreichen Leben mit allen Auserwählten von Todten auferstehen. Dazu höret an das 5. Capitel.

\section[Geschlechtsregister der Patriarchen von Adam an bis auf]{Geschlechtsregister der Patriarchen von Adam an bis auf Noah}
\quellseite{49}
wegen laßt uns in diesen gefährlichen Zeiten fleißig beten, und ein jedes lebe in seinem Beruf und Stande, in wahrer Buße christlich und unsträflich. Wird ja Gott in folgender Zeit die Gottlosen mit den wohlverdienten Strafen heimsuchen, so wird er auch die Gläubigen und Frommen in solchem allgemeinen Uebergange in Gnaden schützen, trösten und zum ewigen Leben bewahren. Davon höret auch an das 7. Kapitel.

\section{Die Sündflut bricht ein}








\section{Summarium}

Ist eine Historia, wie Gott um der Sünde willen durch die Sündflut alles vertilget, was sich auf Erden regt, B. 12.-17. An dem frommen Noah aber beweiset er seine Gnade und erzeiget sich so gnädig und väterlich gegen ihn, daß er ihn mit den Seinen nicht allein erhält im Kasten, sondern auch für sie sorget, daß sie zu essen haben, und die Creaturen, so dem Menschen zu seinem Brauch dienstlich waren, erhalten werden.

Wie nun dies eine Sündflut heißt, darum, daß die Sünde also gräulich gestraft ist: also heißt die heilige Taufe eine rechte Sündflut, durch
\quellseite{50}
welche die Sünde nicht wie dort gestraft, sondern erlauft und vergeben wird durch den Glauben an Christum, welcher die geistliche Arca ist, daran wir erhalten werden.

\section{Psalm}

Der starke barmherzige GOtt tröste und errette alle, die in Noth sind, lindere uns auch die wohlverdienten Strafen, und schütze und erhalte uns gnädiglich in diesen gefährlichen Zeiten, Amen.

\chapter{Das 8. Kapitel}

\section{Vorrede}

Ach wie oft hätte uns Gott auch strafen müssen, wenn er nach unserm Verdienst mit uns hätte handeln sollen! Wer erkennet aber diese Gottes Gnade, Güte und Barmherzigkeit? Dergemahnet Gott, daß er sich so gnädig und väterlich gegen uns erzeigt, und daß er die wohlverdienten Strafen aufschiebt. Lasset uns auch in wahrer Buße leben. Denn auf muthwillige Sünden folget doch endlich große Strafe. Denn Gott kann bald also zornig werden, als gnädig er ist, Sirach 5. So hat kein Zorn über die Gottlosen ein Aufhören; also erinnern wir uns heute abermals der Sündflut. Davon auch lehret das 8. Kapitel.

\section{Die Sintflut nimmt ein Ende}

a Oelblatt) das Blatt bedeutet das Evangelium, das der heilige Geist in der Christenheit hat predigen lassen; denn Oel bedeutet Barmherzigkeit, und Friede, davon das Evangelium lehret.
\quellseite{51}
wegen laßt uns in diesen gefährlichen Zeiten fleißig beten, und ein jedes lebe in seinem Beruf und Stande, in wahrer Buße christlich und unsträflich. Wird ja Gott in folgender Zeit die Gottlosen mit den wohlverdienten Strafen heimsuchen, so wird er auch die Gläubigen und Frommen in solchem allgemeinen Uebergange in Gnaden schützen, trösten und zum ewigen Leben bewahren. Davon höret auch an das 7. Kapitel.

\section{Die Sintflut bricht ein}

\section{Summarium}

Ist eine Historia, wie Gott um der Sünde willen durch die Sintflut alles vertilget, was sich auf Erden regt, B. 12.–17. In dem frommen Noah aber beweiset er seine Gnade und erzeiget sich so gnädig und väterlich gegen ihn, daß er ihn mit den Seinen nicht allein erhält im Kasten, sondern auch für sie sorget, daß sie zu essen haben, und die Kreaturen, so dem Menschen zu seinem Brauch dienstlich waren, erhalten werden.

Wie nun dieß eine Sintflut heißt, darum, daß die Sünde also gräulich gestraft ist: also heißt die heilige Taufe eine rechte Sündflut, durch
\quellseite{52}
welche die Sünde nicht wie dort gestraft, sondern erlöst und vergeben wird durch den Glauben an Christum, welcher die geistliche Arca ist, daran wir erhalten werden.

\section{Psalm}

Der starke barmherzige GOtt tröste und errette alle, die in Noth sind, lindere uns auch die wohlverdienten Strafen, und schütze und erhalte uns gnädiglich in diesen gefährlichen Zeiten, Amen.

\chapter{Das 8. Capitel}

\section{Vorrede}

Ach wie oft hätte uns Gott auch strafen müssen, wenn er nach unserm Verdienst mit uns hätte handeln sollen! Wer erkennet aber diese Gottes Gnade, Güte und Barmherzigkeit? Derwegen danket Gott, daß er sich so gnädig und väterlich gegen uns erzeigt, und daß er die wohlverdienten Strafen aufschiebt. Lasset uns auch in wahrer Buße leben. Denn auf muthwillige Sünden folget doch endlich große Strafe. Denn Gott kann bald also zornig werden, als gnädig er ist, Sirach 5. So hat sein Zorn über die Gottlosen kein Aufhören; also erinnern wir uns heute abermals der Sündflut. Davon auch lehret das 8. Capitel.

\section{Die Sündflut nimmt ein Ende}



a Oelblatt) das Blatt bedeutet das Evangelium, das der heilige Geist in der Christenheit hat predigen lassen; kein Oel bedeutet Barmherzigkeit, und Friede, davon das Evangelium lehret.
\section{Das 1. Buch Mose. 8. 9}

\quellseite{53}

\section{Summarium}

Ist eine schöne Historia, wie die Sintflut wieder abgenommen und Noah sammt den andern, die mit ihm waren, und allen Thieren frisch und gesund aus dem Kasten gingen, und Gott den Segen auf ein neues beide über Menschen und Vieh ausgesprochen habe; durch welchen Segen alles, bis an der Welt Ende, erhalten würde. Sonderlich aber ist dies Exempel wohl zu merken, daß Noah dem HErrn opfert und für seine Wohlthat danket. Und Gott solch Dankopfer ihm läßt so wohl gefallen, daß er verheißet, er wolle fortan seinen Zorn nimmer so heftig über alles Fleisch gehen lassen, sintemal der Mensch von Natur böse und durch die Sünde also verderbet ist, daß er seiner Natur halben nicht anders kann, denn Unrecht thun; solches beweget den gnädigen Gott nicht zum Zorn noch Ungnade, sondern daß er sich unser erbarmet und Sünde vergibt.

\section{Votum}

GOtt der Vater aller Barmherzigkeit vergebe uns alle unsere Sünde, und erhalte uns in aller Gefahr und Noth, mit Leib und Seele zum ewigen Leben, um JEsu Christi seines lieben Sohnes willen, Amen.

\chapter{Das 9. Kapitel}

\section{Vorrede}

Die heiligen Engel im Himmel und alle Creaturen und Geschöpfe am Firmament, in der Luft und auf Erden, bleiben in ihrer Ordnung, zu der sie sind erschaffen, und sind Gott gehorsam: allein der Mensch, der doch die edelste Creatur ist, ist GOtt ungehorsam; wie solches GOtt selber bezeuget im 14. Psalm, da er also saget: Sie sind alle abgewichen und allesamt untüchtig, da ist keiner, der Gutes thue, auch nicht einer. Derwegen läßt uns lernen, was da sei des HErrn Wille, auf daß wir also leben nach seinem Wort durch Hülfe des heiligen Geistes, damit wir auch in Gottes Gnade und Segen leben, zeitlich und ewig, um unsers HErrn Jesu Christi Verdienst und Fürbitte willen. Dazu höret ferner an das 9. Capitel.

\section{GOtt bestätiget seinen Bund durch den Regenbogen}

a Durch Menschen) hier ist das weltliche Schwert eingesetzt, daß man die Mörder tödten soll.
\section{Summarium}

\quellseite{54}

Dies ist eine Historia, wie Gott insonderheit wiederholet den Segen über den Menschen, und setzet ihn zum Herrn über alle andere Kreatur, daß er derselben mächtig sein und zu seinem Nutzen brauchen soll.

Allein verbietet er das Blut, zum Zeichen, daß man Menschenblut nicht vergießen soll; denn hier stehet das Recht, welches Gott selbst spricht: Wer Menschenblut vergießt, deß Blut soll wieder durch Menschen, das ist, durch die Obrigkeit vergossen werden.

Darnach bekräftiget Gott durch das Zeichen des Regenbogens seine Verheißung, daß er fortan dergleichen nicht brauchen wolle.

\section{Der andere Theil des 9. Capitels}

\section{Summarium}

Dies ist ein Exempel, aus welchem Kinder und Unterthanen lernen sollen, wie sie ihre Eltern und Obrigkeit in allen Ehren sollen haben; auf daß der Fluch nicht auf sie erbe, welchen Noah über seinen Sohn Ham hie spricht, darum, daß er ihn unehret. Denn es ist nicht allein Sünde, daß er seine Lust an dem hat, daß er den Noah also entblößet siehet, der böse Bube will auch die andern Brüder dahin bringen, daß sie des alten Vaters lachen, ihn verspotten und nichts auf ihn halten noch geben sollen. Denn Ham thut zweierlei Sünde: den Vater verachtet er und hält sich für weiser und klüger. Da ist unmöglich, daß in einem solchen Herzen könnte einiger rechter Gehorsam sein, wie das Gott von allen Kindern fordert, und greuliche Strafe wider die ergehen läßt, so ihn nicht leisten.

Dagegen aber, wer seine Eltern und Oberherrn in Ehren hat, soll sich des Segens trösten, welchen Noah über seine anderen zween Söhne, Sem und Japheth, spricht. Denn Gott will es mit geistlichem und leiblichem Segen bezahlen, wie das vierte Gebot auch verheißt, daß Gott solchen Gehorsam und Ehrerbietung mit langem glückseligen Leben bezahlen wolle.

\section{Votum}

Helfe Gott uns allen, daß wir gegen unsere Obrigkeit und Eltern uns also verhalten, und in Gottesfurcht also unser Leben beschließen, auf daß wir für Christi, unseres HErrn, Richterstuhl wohl bestehen, Amen.

\chapter{Das 10. Capitel}

\section{Vorrede}

Dieses Kapitel wird nicht gelesen, denn es ist nur ein Geschlechts-Register und Namen-Buch der Nachkommen des Noah, welches dazu dienet, daß man daraus sehe die Fortpflanzung des frommen und gottesfürchtigen Geschlechts Noahs und der Verheißung seines Samens. Wiewohl, was von Nimrod gesagt wird, zu dem Ende zu merken ist, weil er, nachdem er das Reich und Gewalt überkommen, zu einem Tyrannen worden und die rechtgläubige Kirche sehr geplaget hat. Denn daß er hat gehabt den Namen eines gewaltigen Jägers für dem HErrn, ist nichts anderes, denn daß er die rechtgläubige Kirche unterdrückt habe, und nur darauf sei bedacht gewesen, wie er seine Macht und Gewalt vermehren und erweitern möge.
\section{Geschlecht und Nachkommen der Söhne Noah}

\quellseite{55}

a) Peleg auf deutsch eine Zertheilung.

\chapter{Das 11. Kapitel}

\section{Vorrede}

Vom heiligen alten Tobia stehet geschrieben Tob. 14.: Je älter er worden, daß er je mehr in Gottesfurcht habe zugenommen. Das Widerspiel stehet man heut am meisten Theil der Menschen. Desselbige aber ist dieses: Je älter die Leute werden, je ärger sie auch werden. Derwegen weil wir Gottes Wort haben, auch das täglich hören, so wollen wir auch in Kraft des heiligen Geistes nach Gottes Wort leben, wollen das Böse hassen und lassen, und das Gute lieben und üben, auf daß wir Gottes Zorn und Strafe auf uns nicht laden, sondern in seiner Gnade und unter seinem Schutze leben und sterben, Amen. Darauf höret nun an das 11. Kapitel.

\section{Thurm zu Babel. Verwirrung der Sprachen. Geschlecht Sems}
\quellseite{56}



a Babel) auf deutsch eine Vermischung oder Verwirrung.







\section{Summarium}

Dies ist eine Historia, wie die Welt, sobald sie ein wenig gewachsen, wieder sei böse, frech und hoffärtig worden, darum sie Gott den Ländern nach mit der Sprache getheilet und ihr Fürnehmen gehindert hat. Solche Theilung der Sprachen ist eine greuliche Sündenstrafe gewesen. Denn da ist so bald mit der Sprache das Predigtamt gefallen und ist so bald Abgötterei eingerissen. Denn ein jeder ihm einen eigenen Gottesdienst gemacht hat. In des einigen Sems Linie und Geschlecht ist die rechte Sprache, das Predigtamt und Gottesdienst blieben. Zum selben hat sich Japhet gehalten. Aber Ham, als der große Haufe, ist mit Macht in die Abgötterei gerathen. Gleich aber, wie Gott hier die Sünde strafet durch viele Zungen: also beweiset er am Pfingstage seine Gnade durch viele Zungen, daß man in allen Sprachen Gottes Gnade hört und das Predigtamt in allen Zungen gehet.

Darnach erzählet Mose das Geschlecht Sem, auf daß wir wissen die Ankunft des Erzvaters
\quellseite{57}
Abraham, da nachmals die Verheißung auf uns beruhet. Da stehet man, wie die Welt so bald abgenommen und die heiligen Leute nicht mehr so lange gelebet haben, wie Noah und seine Voreltern.

\section{Votum}

Behüte und regiere uns der treue GOtt alle, damit wir an Tugend und Ehrbarkeit nicht abnehmen, sondern zunehmen und zeitlich und ewig leben in Heiligkeit und Seligkeit, um Jesum Christi, unsers HErrn und Heilandes willen, Amen.

\chapter{Das 12. Kapitel}

\section{Vorrede}

Gleich wie die Sonne mitten im Winter und in der größten Kälte am weitesten von uns scheinet zu sein, da sie doch in ihrem Zirkel näher bei uns ist, als wenn die Hitze im Sommer am größten ist: also ist unser HErr und Helfer, der treue liebe Gott, uns auch am nächsten, wenn wir meinen, er sei gar weit von uns. Derowegen sagt man im gemeinen Sprüchwort: Je größer Noth, je näher Gott. Denn wenn Rath und Hilfe bei den Menschen aus ist, da lässet sich der große Nothhelfer und Rathgeber erst recht sehen mit Trost, Hilfe und Errettung. Davon lehret der König David im 34. Psalm also: Der HErr, spricht er, ist nahe bei denen, die zerbrochenes Herzens sind, er hilft denen, die zerschlagenen Gemüth haben. Des erfreuet und tröstet euch, ihr betrübte und nothleidende Herzen. Und höret darauf an das 12. Kapitel.

\section{Abram wird berufen und ziehet nach Egypten}
\quellseite{58}

\section{Summarium}

Hier gehet an die Legende des heiligen Patriarchen Abrahams, welchen Gott fordert aus dem Lande Haran ins Land Canaan, und verheißet ihm, er wolle ihn zum besondern Volk machen, und das Land Canaan eingeben; und zu seiner Zeit ihm einen solchen Samen erwecken, der ein gebenedeyter Same sei, das ist, durch welchen alle Welt aus dem Fluch zum Segen und Gottes Gnade soll gebracht werden.

Hie soll man erstlich lernen, daß Gott gnädig sei. Darnach soll man Acht haben auf das Exempel Abrahams, wie er Gott trauet und seinem Befehle folget und Gottes Wort prediget, und dennoch das Kreuz nicht außen bleibet, sintmal er Theurung halben in Egypten fliehen und da seines Weibes halben in Gefahr kommen muß. Aber weil er Gott glaubet und vertrauet, so hilft er ihm aus aller Noth; wie durch seine ganze Legende zu sehen ist; und strafet, die ihn beleidigen, wie er hie Pharao und sein Haus plaget, darum, daß er Sara ihm zum Weibe wollte nehmen.

\section{Votum}

Der gerechte starke Gott strafe noch heut alle muthwillige und unbußfertige Verfolger seiner Kirche, und tröste und errette alle gläubige Kreuzträger, zu Ehren seinem heiligen Namen, Amen.

\chapter{Das 13. Kapitel}

\section{Vorrede}

Wer Gott treulich von ganzem Herzen und rechtschaffen dienet, dem dienet Gott wieder auch nutzbarlich. Gott dienen aber heißt, Gottes Wort folgen und gehorsam sein, beides im Glauben gegen ihn und in der Liebe gegen den Nächsten. Derwegen weil wir Gottes auch keinen Augenblick ohne unsern Schaden entbehren können, und weil wir Gottes zu unsern Augen alle Augenblicke bedürfen: so laßt uns durch Hülfe des heiligen Geistes Gottes Willen thun, und ihm dienen ohne Furcht unser Leben lang in Heiligkeit und Gerechtigkeit, die ihm gefällig ist, Luc. 1. Und darauf höret an das 13. Kapitel.

\subsection{Abram scheidet sich von Lot, seinem Vetter}
\quellseite{59}

\section{Summarium}

Hält uns erstlich ein schön Exempel für mit Abraham und Lot, wie Gott seinen Segen reichlich mittheilet denen, so ihn fürchten, nach dem Spruch Christi: Trachtet am ersten nach dem Reich Gottes und nach seiner Gerechtigkeit, so wird euch das andere alles zufallen; denn euer Vater weiß, was ihr nothdürftig seid, Matth. 6, 33.

Darnach wie Uneinigkeit zu fliehen und sonderlich die Blutfreunde, alles, so Widerwillen erregen kann, meiden und eher den Vortheil aus den Händen geben sollen, denn daß man Ursache zu Zank und Widerwillen geben sollte, wie der fromme Abraham mit Lot thut. Solcher Fleiß ist sonderlich vonnöthen; denn wer Zank suchet und nach Einigkeit nicht mit besonderm Fleiß trachtet, der wird sein allenthalben genug finden. Aber selig sind die Friedfertigen.

\section{Votum}

Gott der Allmächtige steure in aller Welt dem Zanke und Hader und allen andern Sünden, und regiere uns, daß wir leben nach seinem Willen, ihm zu Ehren, uns zu Nutz und zu unsers Nächsten Wohlfahrt, Amen.

\chapter{Das 14. Kapitel}

\section{Vorrede}

Gott züchtiget mit väterlicher Ruthe Land und Leute, daß sie sollen Buße thun, auf daß er sie nicht mit dem Garaus und ewiger Pein strafen dürfe. Das bezeuget St. Paulus 1 Cor. am 11. und saget: Wenn wir gerichtet, das ist, gestrafet werden, so werden wir von dem Herrn gezüchtiget, auf daß wir nicht sammt der Welt verdammt werden. Derowegen weil Gott mit allerlei Krankheiten, mit seltsamen Gewittern, mit wunderlichen Fällen und auf andere mancherlei Weise uns zur Buße rufet, so wollen wir unsere Sünde erkennen, die Gott abbitten in wahrem Glauben an den Herrn Christum, und wollen gottsfürchtig leben, daß uns Gott mit endlichem Verderben nicht strafen dürfe.

Darauf höret an das 14. Capitel.

\subsection{Abram errettet Lot; wird von Melchisedech gesegnet}
\quellseite{60}

\section{Summarium}

\unsicher{Sie} gehet an die Legende des heiligen Patriarchen Abrahams, welchen Gott fordert aus dem Lande Haran ins Land Canaan, und verheißet ihm, er wolle ihn zum besondern Volk machen, und das Land Canaan eingeben; und zu seiner Zeit ihm einen solchen Samen erwecken, der ein gebenedeyter Same sei; das ist, durch welchen alle Welt aus dem Fluch zum Segen und Gottes Gnade soll gebracht werden. Sie soll man erstlich lernen, daß Gott gnädig sei. Darnach soll man Acht haben auf das Exempel Abrahams, wie er Gott trauet und seinem Befehle folget und Gottes Wort prediget, und dennoch das Kreuz nicht außen bleibt, finemal er Theurung halben in Egypten fliehen und da seines Weibes halben in Gefahr kommen muß. Aber weil er Gott glaubet und vertrauet, so hilft er ihm aus aller Noth; wie durch seine ganze Legende zu sehen ist; und strafet, die ihn beleidigen, wie er hie Pharao und sein Haus plaget, darum, daß er Sara ihm zum Weibe wollte nehmen.

\section[Notum]{\unsicher{Notum}}

Der gerechte starke Gott strafe noch heut alle muthwillige und unbußfertige Verfolger seiner Kirche, und tröste und errette alle gläubige Kreuzträger, zu Ehren seinem heiligen Namen, Amen.

\chapter{Das 13. Kapitel}

\section{Vorrede}

Wer Gott treulich von ganzem Herzen und rechtschaffen dienet, dem dienet Gott wieder auch nutzbarlich. Gott dienen aber heißt, Gottes Wort folgen und gehorsam sein, beides im Glauben gegen ihn und in der Liebe gegen den Nächsten. Derowegen weil wir Gottes auch keinen Augenblick ohne unsern Schaden entbehren können, und weil wir Gottes zu unserm Nutzen alle Augenblicke bedürfen: so laßt uns durch Hülfe des heiligen Geistes Gottes Willen thun, und ihm dienen ohne Furcht unser Leben lang in Heiligkeit und Gerechtigkeit, die ihm gefällig ist, Luc. 1. Und darauf hört an das 13. Kapitel.

\subsection{Abram scheidet sich von Lot, seinem Vetter}
\quellseite{61}

\section{Summarium}

Hält uns erstlich ein schön Exempel für mit Abraham und Lot, wie Gott seinen Segen reichlich mittheilet denen, so ihn fürchten, nach dem Spruch Christi: Trachtet am ersten nach dem Reich Gottes und nach seiner Gerechtigkeit, so wird euch das andere alles zufallen; denn euer Vater weiß, was ihr nothdürftig seid, Matth. 6, 33.

Darnach wie Uneinigkeit zu fliehen und sonderlich die Blutfreunde, alles, so Widerwillen erregen kann, meiden und eher den Vortheil aus den Händen geben sollen, denn daß man Ursache zu Zank und Widerwillen geben sollte, wie der fromme Abraham mit Lot thut. Solcher Fleiß ist sonderlich vonnöthen; denn wer Zank suchet und nach Einigkeit nicht mit besonderm Fleiß trachtet, der wird sein allenthalben genug finden. Aber selig sind die Friedfertigen.

\section{Votum}

Gott der Allmächtige steure in aller Welt dem Zanke und Hader und allen andern Sünden, und regiere uns, daß wir leben nach seinem Willen, ihm zu Ehren, uns zu Nutz und zu unsers Nächsten Wohlfahrt, Amen.

\chapter{Das 14. Kapitel}

\section{Vorrede}

Gott züchtiget mit väterlicher Ruthe Land und Leute, daß sie sollen Buße thun, auf daß er sie nicht mit dem Garaus und ewiger Pein strafen dürfe. Das bezeuget St. Paulus 1 Cor. am 11. und saget: Wenn wir gerichtet, das ist, gestrafet werden, so werden wir von dem Herrn gezüchtiget, auf daß wir nicht sammt der Welt verdammt werden. Derowegen weil Gott mit allerlei Krankheiten, mit seltsamen Gewittern, mit wunderlichen Fällen und auf andere mancherlei Weise uns zur Buße rufet, so wollen wir unsere Sünde erkennen, die Gott abbitten in wahrem Glauben an den Herrn Christum, und wollen gottfürchtig leben, daß uns Gott mit endlichem Verderben nicht strafen dürfe. Darauf hört an das 14. Kapitel.

\subsection{Abram rettet Lot; wird von Melchisedek gesegnet}
\quellseite{62}

\textit{Zu Vers 18.}a Trug Brod) nicht daß ers opferte, sondern daß er die Gäste speisete und ehrte, dadurch Christus bedeutet ist, der die Welt mit dem Evangelio speiset.

\section{Summarium}

Dies ist eine Historia, wie Gott durch leibliche Strafe endlich die Menschen zur Buße fordert, daß sie sich bessern und frömmer sollen werden. Denn ehe denn er Sodoma und Gomorra ließ mit höllischem Feuer verbrennen, vermahnet er sie zur Buße und Besserung durch Kriege.

Solche Strafe muß der fromme Lot auch mit tragen. Aber endlich wird er errettet durch Abraham, welcher hie zum Exempel stehet aller Obrigkeit, daß sie die Ihren schützen und doch nicht eigennützig sein sollen.

Daß Abraham dem König Melchisedek den Zehnten gibt, welcher ein Priester Gottes war, ist zum ersten ein Exempel, daß weltliche Obrigkeit schuldig ist zu helfen, daß die Kirchen und der rechte Gottesdienst erhalten werde. Denn sie muß Schutz und Unterhaltung haben. Und selig sind die, so ihr Geld, Gut und Gewalt zu solchem köstlichen guten Werke brauchen, daß das Predigtamt dadurch gefördert und erhalten werde.

Danach hat es seine Deutung auf das Priesterthum Christi, wider das Gesetz, wie die Epistel an die Ebräer zeuget, Kap. 7, 10. und David Ps. 110, 4., der Christum einen Priester nennet, nicht nach der Ordnung oder Weise Levi, sondern Melchisedek, der nicht mit Opfern und Sünden zu thun hat, wie der Priester im Gesetz, sondern segnet und nimmt den Fluch von uns. Der Name Melchisedek heißt auf deutsch so viel, als ein König der Gerechtigkeit. Solchen Namen führet unser lieber HErr Christus, der durch Vergebung der Sünden uns gerecht machet. Es ist aber gewiß dieser Melchisedek Sem gewesen, des Noah Sohn, welcher dazumal der älteste Patriarch und der oberste Regent und Lehrer in der Kirche gewesen ist, und derhalben den Namen hat, er sei des Allerhöchsten Priester. Dieses Sems Nachkommen ist Abraham, und hat Abraham von ihm geerbet die Verheißung von Christo.

\section{Votum}

Der treue barmherzige GOtt verleihe, daß wir alle in wahrer Buße leben und sterben; und segne und behüte ferner unsere Stadt, und das ganze gemeine Vaterland, um Jesu Christi unsers HErrn willen, Amen.

\chapter{Das 15. Kapitel}

\section{Vorrede}

Die großen Heiligen haben auch großes Kreuz und müssen eines nach dem andern tragen. Aber sie sind davon, auch mit großer Freude, endlich erlediget worden. Von beiden sagt der Prophet David im 34. Psalm also: Der Gerechte muß viel leiden, aber der HErr hilft ihm aus dem allem. Derowegen wer ein Kreuzträger ist, der lasse ihm Gottes guten väterlichen Willen wohlgefallen, sei geduldig und bete. Denn Gott

2*
\quellseite{63}
zürnet nicht ewiglich und seine rechte Hand kann alles ändern, Ps. 77.; und kann auch das größte Unglück bald in Freude verwandeln. Dieses lehret auch das 15. Kapitel.

\subsection{Abrams Glaube und Gerechtigkeit wird gepriesen}

\textit{Zu Vers 9.}a Bringopfer.

\textit{Zu Vers 11.}b Gevögel) das Gevögel und der rauchende Ofen und der feurige Brand bedeuten die Egypter, die Abrahams Kinder verfolgen sollten; aber Abraham scheucht sie davon, das ist, Gott erlöset sie um der Verheißung willen, Abraham versprochen. Daß aber er nach der Sonne Untergang erschrickt, bedeutet, daß Gott seinen Samen eine Zeit lang verlassen wolle, daß sie verfolget würden, wie der HErr selbst hie betet. Also gehet es auch allem Gläubigen, daß sie verlassen und doch erlöset werden.

\section{Summarium}

Zeiget erstlich an, wie der fromme Abraham ohne große Anfechtung nicht sei gewesen, daß Gott die Verheißung, daß er ihn zum großen Volk wollte machen, so lange verzog. Darum, da Gott zu ihm spricht: Fürchte dich nicht, ich bin dein Schild etc.; spricht er darauf: Was willst du mir geben, ich gehe dahin ohne Kinder; als wollte er sagen: Soll ich zum großen Volk werden, so mußt du mir, lieber HErr, ja einen Sohn geben. Als ihm Gott aber verheißet, er soll einen Sohn haben, überwindet er solche Anfechtung, daß er so alt und seine Sara unfruchtbar war, und glaubet Gott. Solcher Glaube, spricht Moses, ist ihm zur Gerechtigkeit gerechnet. Darnach am Ende ist eine Weissagung, wie die Verheißung aufgezogen und Abrahams Nachkommen zuvor in Egypten müssen geplagt werden, auf daß sie sich an solchem Unfall nicht ärgerten, sondern fest an der Zusagung Gottes hielten, daß Land Canaan würde noch ihr eigen werden, und daß Christus, der gebenedeite Same, Gottes Verheißung nach, gewißlich kommen würde.

\section{Psalm}

Wir danken GOtt, der seiner Christenheit unter dem Kreuz beigestanden ist allezeit;
\quellseite{64}
der wolle uns auch bewohnen, und in aller Noth trösten und erhalten, hochgelobt in Ewigkeit, Amen.

\chapter{Das 16. Kapitel}

\section{Vorrede}

Wenn ein böses Ding ein allgemeiner Gebrauch wird, so pflegen solche auch die, so fromme Leute sonst sind, nicht für so gar sträflich zu halten: wie heute viele große Laster, als Trunkenheit, Hoffart, Wucher und dergleichen gering geachtet werden, weil die beim größten Haufen im Schwange gehen. Aber Sünde ist und bleibt Sünde, und wenn sie noch so gemein wäre. Derowegen wer also handeln und wandeln will, daß er gegen Gott bestehe, der soll nicht nach der Welt bösen Exempeln, sondern nach Gottes ernstem Willen und unwandelbarem Wort leben. Denn das allein soll unser Fuß Leuchte sein und ein Licht auf unserm Wege, saget der königliche Prophet David Ps. 119. Dazu höret auch an das 16. Kapitel.

\subsection{Hagar gebieret dem Ismael}

\textit{Zu Vers 2.}a Bauen) das ist, Kinder kriegen.

\textit{Zu Vers 11.}b Ismael) heißt, Gott erhöret.

\section{Summarium}

Zum ersten ist ein Exempel, wie die Vernunft Rath suchet, ihr selbst zu helfen, wo Gott die Hülffe zu lange verzeucht und nicht so bald hilft, als man hofft. Denn Sarai, weil sie je länger je älter wird, und doch unfruchtbar bleibt, beginnt sie an Gottes Verheißung zu wanken, denket, es werde mit ihr nichts draus, weil sie unfruchtbar und nun alt sei, und hält an bei Abram, daß er die Hagar, ihre Magd, zum Weibe nehme, auf daß Gottes Verheißung nicht verhindert werde. Es hat aber Gott solches lassen also geschehen zur sondern Deutung; wie Paulus zu den Galatern Kap. 4, 22. davon disputirt, von zweien Kirchen, die eine so mit Werken umgehet und auf eigen Verdienst bauet; die andere, so an Gottes Zusagung und Barmherzigkeit sich hält.

Darnach ist hie eine feine Hauslehre, wie Hagar, da sie schwanger war, ihrer Frau Sarai Strafe nicht leiden wollte und lief davon. Aber der Engel heißt sie wieder in den Dienst gehen und ihrer Frau unterthan sein, und sich für ihr demüthigen; zum Exempel allen Ehehalten oder Dienstboten, daß sie mit Demuth allen Gehorsam ihrer Herrschaft leisten sollen. Denn sol-
\quellseite{65}
Dieß will Gott haben und gefällt ihm wohl, will auch Glück dazu geben; so dagegen, wo das Gesinde untreu und ungehorsam ist, Gott es mit allerlei Unglück strafen will.

\section{Datum}

Verleihe der gütige treue Gott, daß wir uns nach seinem Wort gegen ihn und gegen unsern Nächsten gehorsamlich verhalten, und daß wir an Leib und Seele seinen Segen haben und behalten, Amen.

\chapter{Das 17. Kapitel}

\section{Vorrede}

Der Glaube und der Gehorsam gegen Gott sind die zwei Stücke, welche die heilige Schrift fürnehmlich lehret, und auch von allen Menschen, die zeitliche Wohlfahrt und ewige Seligkeit erlangen sollen, erfordert. Derowegen läßt uns alle Tage fleißig zur Kirche gehen und beten, daß die Anhörung göttlichen Wortes uns auch zu täglicher Stärkung unsers Glaubens und Besserung unsers Lebens fruchtbarlich diene und gelange. Also wird sich Gott gegen uns väterlich erzeigen, da wir uns gehorsamlich gegen sein Wort verhalten. Darauf höret auch an das 17. Kapitel.

\subsection{Isaaks Verheißung wird mit der Beschneidung bestätigt}

\textit{Zu Vers 5.}a Abram heißt hoher Vater, Abraham aber, der Haufen Vater; wiewohl dieselben Haufen nur mit einem Buchstaben angezeigt werden in seinem Namen, nicht ohne Ursache.

\textit{Zu Vers 7.}b Nachkommen) Ledorotham, daß ist, so lange ihr Ding währen wird. Denn Moses hiermit deutet, daß ihr Ding solle endlich aufhören und ein anderes kommen.

\section{Summarium}

Ist eine Historia, wie Gott Abraham tröstet und die Verheißung mit einem neuen Bundeszeichen der Beschneidung bekräftiget, daß er aller der Gott sein, sie segnen und behüten wolle, so solches Bundeszeichen an ihrem Leibe tragen.

Die hebt sich erstlich der Unterschied an, daß Gott den Abraham und seine Nachkommen als ein sonder Volk von allen andern Völkern absondert und ihn und seine Nachkommen, die Juden, für der Welt und an ihrem Leibe mit einem sondern Gemerk oder Zeichen bezeichnet; auf daß man wüßte, daß Christus allein aus diesem beschnittenen Volk und sonst von keinem Volk sollte geboren werden.

Zum andern hat solche Beschneidung auch ihre Bedeutung gehabt: denn Gott hat damit wollen anzeigen, wie unsere Natur sündhaftig sei und dürfe, daß sie geistlich beschnitten würde; wie die Propheten auch predigen, daß man die Vorhaut des Herzens beschneiden, das ist, in Gottesfurcht wandeln und die Sünde mit allen bösen Lüsten töten soll.

Zum dritten ist diese Beschneidung gewesen,
\quellseite{66}
wie es Paulus Röm. 4, 11. nennet, ein Siegel der Gerechtigkeit; denn dabei hat Abraham und alle seine Nachkommen sich können erinnern, daß sie Gottes Volk sind; sintemal solches Bundeszeichen, als ein Sacrament, die Verheißung hätte. Gott wollte solcher beschnittenen Leute sich sonderlich annehmen, sie schützen und segnen als die Seinen.

\section{Summarium}

Ist endlich ein Exempel, daß wir lernen sollen, daß Gottes Verheißungen wahr sind und wir daran nicht sollen zweifeln; denn bei Gott ist nichts unmöglich; wie es Abraham unmöglich däucht, daß er von seiner Sarah noch sollte einen Sohn zeugen.

Darnach ist's ein Exempel eines christlichen Gehorsams in Abraham, daß er nicht lange das von disputiret, warum er und sein Gesinde sich beschneiden sollen; sondern weil es Gott befohlen hätte, thut ers und folget mit einfältigem Glauben, B. 23. Denn es für der Welt ein sehr närrisch Ansehen hat, daß der alte Mann sich und sein ganzes Haus beschneiden soll. So kann Vernunft keinen sondern Nutzen noch Frucht solches Werkes erdenken. Aber um solches bekümmert sich Abraham gar nichts und gehet sein straß in vollem Gehorsam und willigem Herzen dahin, läßt ihm an dem genügen, daß es Gott also geheißen und befohlen hat. Also sollen wir in allen Dingen auch thun und nur dahin sehen, ob es Gottes Befehl sei, und weiter nichts fragen noch uns bekümmern, warum ers befohlen oder wozu es nütze sei. Solchen einfältigen Gehorsam will Gott von uns allen haben.

\section{Votum}

Der barmherzige gütige Gott verleihe, daß wir seinem Willen in Liebe und Leid treulich folgen, und seiner gnadenreichen Verheißung zu zeitlicher und ewiger Wohlfahrt fruchtbarlich genießen, Amen.

\chapter{Das 18. Kapitel}

\section{Vorrede}

Gottes Wort hat diese Kraft, wo man es mit Ernst handelt, höret und betrachtet, daß es nimmer ohne Frucht abgehet, sondern allzeit neuen Verstand und Andacht erweckt und reine Herzen und Gedanken machet. Auch ist dies gewiß: wie sich der Mensch gegen das Wort Gottes verhält in seinem Leben, also verhält sich Gottes Wort wiederum gegen den Menschen in seinem Absterben. Derowegen wollen wir Gottes Wort lieben und lernen, und bitten, daß Gott auch dadurch in uns täglich wirken wolle Verstand, Glauben und Gehorsam, und daß wir daraus Freude und Trost haben unterm Kreuz, und in unserer Hinfahrt aus diesem zeitlichen kümmerlichen Leben, Amen. Darauf höret ferner an das 18. Kapitel.

\subsection[Dem Abraham wird nochmals Isaac verheißen, und die]{Dem Abraham wird nochmals Isaac verheißen, und die Verheißung Sodoms geoffenbaret}
\quellseite{67}

\textit{Zu Vers 2.}a Nieder] vor einem fällt er nieder und redet auch als mit einem, und doch mit dreien. Da ist die Dreifaltigkeit von Gott angezeigt.

\textit{Zu Vers 14.}b So ich lebe] Gott redet als ein Mensch, 1 Mos. 3, 9.: Adam, wo bist du? Item 1 Mos. 11, 7.: Ich will hinab fahren und sehen. 1 Mos. 18, 21.: Ich will sehen, ob sie sei. Denn dies Wort (so ich lebe), setze es, wo du willst, so lautet es doch nicht, als wenn er von der Frucht oder Eltern wollte sagen: Du sollst einen Sohn haben, so er lebet, oder, so ihr lebet. Meinst du, er wisse nicht, ob er Sohn oder die Eltern leben werden? Ja, wo sie nicht ewig leben würden, was wäre die Verheißung?

\section{Summarium}

Ist eine holdselige Historia, wie Gott zu Abraham kommet und sein Gast wird, um \unsicher{meinen} darum, daß er ihn im Glauben stärke, daß er nicht zweifele, er werde einen Sohn von seiner Sarah haben, denn dem HErrn sei nichts unmöglich, V. 10. Auf daß wir auch lernen auf Gottes Wort sehen, und uns gewiß auf seine Zusagung verlassen. Sonderlich aber ist hie zu merken, wie der Text sagt, Abraham habe drei Männer gesehen, und redet doch je als mit einem allein, und je als mit dreien, V. 2.

Item, der HErr sprach zu Abraham etc., V. 13.

Solches ist eine gewisse Anzeigung, daß nur ein einiger Gott sei, aber drei unterschiedliche Personen, als GOtt Vater, GOtt Sohn, GOtt heiliger Geist. Wie denn die Offenbarung im Neuen Testament solches klärlich mit bringt. Solche Zeugnisse aber des Alten Testaments dienen dazu, daß wir lernen, die Patriarchen haben auch solches gegläubet und gewußt wie wir, und sind eben in solchem Glauben selig worden.

Darnach ist ein fein Exempel Abrahams halben, der so gastfrei und gegen fremde und unbekannte Gäste so freundlich und diensthaftig ist, V. 4, nach der Regel St. Pauli Röm. 12, 13.: Herberget gerne.
\quellseite{68}

\textit{Zu Vers 25.}c In Hebräisch: sollst du sonst alle Welt richten und nicht recht thun.

\section{Summarium}

Ist erstlich eine köstliche Historia, in welcher Gott seine Barmherzigkeit reichlich läßt sehen, wie er gerne alle Strafe abwenden und schonen wolle, so nur die Menschen sich begehreten zu bessern; ja daß er um wenig frommer Leute willen eines ganzen Landes gerne verschonen wollte.

Widerum aber ist es schrecklich, daß die Welt so arg und boshaftig soll sein, daß unter so viel Leuten in einem ganzen Lande nicht zehn fromme und gottesfürchtige Menschen sollen erfunden werden. Darum muß die Strafe über sie gehen und Abrahams Fürbitte umsonst sein, weil keine Besserung des Lebens folget. Denn beten und sich bessern muß beides bei einander sein.

Zum andern siehet man hie, wie die Fürbitte für andere gegen Gott so ein trefflich Ding ist. Aber wer solcher Fürbitte genießen will, muß sich dargegen recht schicken und die Fürbitte durch Unbußfertigkeit nicht hindern.

\section{Gebet}

Der allmächtige barmherzige Gott komme auch zu uns, und verleihe, daß wir glauben, beten und leben, daß sein Name allzeit durch uns geheiliget, gelobet und gepreiset werde, hie zeitlich und dort ewiglich, Amen.

\chapter{Das 19. Kapitel}

\section{Vorrede}

Betet herzlich, daß Gott der heilige Geist in uns kräftig wirken wolle, auf daß die erschreckliche Historia vom Untergang Sodoma und Gomorra heut so tief in unsere Herzen sich senke, daß wir der nimmermehr vergessen, und auf daß wir diese Geschichte zu unserer täglichen Besserung fruchtbarlich gebrauchen. Denn Gott dem HErrn niemand so lieb ist, da er sündigt, und in Sünden muthwillig verharrret, er straft ihn, daß ers an Leib und Seele, zeitlich und ewig fühlet. Aber, o besser nicht geboren, denn allhie Gottes Zorn tragen, und dort ewig sein verloren. Darauf höret an das 19. Kapitel.

\subsection{Vertilgung Sodoma. Lots Errettung und Sündenfall}
\quellseite{69}

\textit{Zu Vers 22.}a Zoar) heißt klein.
\quellseite{70}

\section{Summarium}

In dieser Historie stehet man, wie Gott an seine Heiligen denkt, und sie errettet von der Strafe, so über die Gottlosen kömmt. Denn die Frommen sollen ihrer Frömmigkeit genießen, wie Lot.

Wiederum stehet man, weil die Sünde so überhand genommen hat, und die Sodomiter keine Warnung wollten zu Herzen nehmen, daß Gott hat müssen solchen Ernst brauchen, auf daß wir lernen uns fürchten und für Sünden hüten; sintemal sie Gott nicht kann ungestraft hingehen lassen. Merke aber wie sicher die Sünder sind, wenn ihr Verderben am allernächsten ist: die Sünde des Lots, da sie gewarnet werden, hielten es für ein Gelächter und gedachten nicht zu weichen, B. 14.

Daß die Frau Lots sich umsiehet und zur Salzsäule wird, B. 26. deutet Christus im Evangelio Luc. 17, 21. dahin, daß man sich das Zeitliche nicht lasse hindern, sondern auf Gottes Wort sehe und demselbigen folge; sintemal wir der Strafe sonst nicht können entrinnen. Solche Deutung nimmt der Herr gewöhnlich davon, daß es mit Lots Weib nicht ein schlechtes fürwitziges hinter sich sehen gewesen ist, sondern daß sie ihr die Welt auch hat lassen lieb sein, und Gottes ernstliches Dräuen nicht für wahr gehalten hat. Solchen Unglauben hat Gott mit Ernst strafen, und das Gedächtnis andern zur Abscheu, öffentlich wollen bleiben lassen.

\section{Summarium}

Ist eine Historia, welche nicht darum geschrieben ist, noch gelesen wird, daß wir solchem Exempel, Lots und seiner Töchter, folgen, sondern unsere eigene Bosheit recht sollen erkennen lernen, und desto mehr Gott für Augen haben, und unsere bösen Lüste dämpfen; sintemal auch die Heiligen sich nicht allezeit bewahren können, sondern bisweilen allzu grob fallen.

Der Text zeiget an, daß die Töchter Lot solches aus Einfalt fürgenommen, und Lot unwissend gethan habe, B. 31. 33. Aber damit sind sie beide nicht entschuldiget. Und sonderlich stehet man hie, was für gräuliche Sünden aus der Trunkenheit erwachsen: denn wo Lot nüchtern gewesen, würde er nimmermehr solche Sünde begangen haben. Ob er nun wohl als ein trauriger, bekümmerter Mann, der so vieler Seelen Verderben gesehen, so viel guter Freunde und unter denselben sein liebes Weib alle auf einen Tag, und schier in einem Augenblicke verloren hatte, mehr vor Leid und Kummer, denn von Wein ist voll worden; denn bekümmerten Leuten schadet der Trunk viel leichter: so ist die Trunkenheit an ihr selbst ein schädliches Laster, und verursacht viel Jammer und Noth; darum alle Christen sich dafür hüten sollen.

\section{Gebet}

Helffe der treue gütige GOtt, daß alle Welt an Sodoma und Gomorra sich spiegele, Buße thue, und Gottes Zorn und Strafe in diesem und künftigen Leben entfliehe, Amen.

\chapter{Das 20. Capitel}

\section{Vorrede}

Keine Sünde ist so groß, gräulich und schädlich, ein jeder unter uns würde solche schon lange begangen und gethan haben, da uns Gott nicht aus väterlicher Güte und Barmherzigkeit
\quellseite{71}
behütet hätte. Dafür laßt uns ihm herzlich danken und bitten, er wolle uns ferner regieren und führen, daß wir durch unsere Sünde ihn nicht erzürnen, des Herrn Christi nicht spotten, den heiligen Geist nicht betrüben, die lieben Engel nicht verjagen, uns und den Unsern, auch unserm lieben Vaterlande, verderbliche Strafen nicht zuziehen; und laßt uns derhalben täglich mit dem Könige David seufzen, Psalm 119. Thue wohl deinem Knecht, daß ich lebe, und dein Wort halte; also: thue wohl deiner Magd, daß ich lebe und dein Wort halte. Zu diesem höret an das 20. Kapitel.

\subsection[Sarah wird dem Abraham von Abimelech genommen, und]{Sarah wird dem Abraham von Abimelech genommen, und mit Gewinst wieder gegeben}

\textit{Zu Vers 10.}a Geschehen weil du ein Prophet bist, magst du was gesehen haben, daß ichs verdienet habe mit meinen Sünden.

\textit{Zu Vers 13.}b Wandern Gott hieße mich in die Irre ziehen, als wären viel, und doch ein Gott.

\textit{Zu Vers 16.}c Ihre Strafe hieß die Heiligen werden säuberlich und mit Gewinst gestraft, als hie Sarah wird gestraft, daß sie Abraham hatte Bruder genennet, und erhält große Wohlthat.

\section{Summarium}

Ist eine Historia, wie Abraham seiner Sarah halben noch einmal in Gefahr kommt, daß sie ihm vom König Abimelech genommen wird, V. 2. Aber hie beweiset sich Gott wiederum, wie er die Seinen behüten wolle, und sonderlich wie er dem Ehebruch feind, und denselben mit Unfruchtbarkeit und dem Tod wolle strafen, V. 18. Daneben ist sonderlich wohl zu merken, wie es eine Gabe Gottes und von Gott zu bitten sei, daß er für Sünden uns bewahren wolle; wie Gott zu Abimelech sagt: Ich habe dich behütet, daß du nicht wider mich sündigest,
\quellseite{72}
und habe das nicht zugeben, daß du Sarai, Abrahams Weib, berühretest, V. 6. Denn unsere Natur ist viel zu schwach, daß sie sich aller Anfechtung und Sünde erwehren sollte, wo Gott nicht mit sonderer Gnade hält und hütet, und wir nicht täglich um solchen Schutz und Hülfe bitten; wie Christus im Vaterunser auch befiehlt und lehret, daß wir beten sollen: Nicht einführe uns in Versuchung, das ist, lieber Herr Gott, wo wir in Versuchung und Anfechtung geraten, sei gnädig, und führe uns wieder heraus, daß wir nicht darinnen liegen bleiben und verderben. Solcher Bitte bedürfen wir überaus wohl.

\section{Votum}

Gott der heilige Geist helfe, daß wir in täglichem Gebet fleißig und andächtig sein; und stärke uns an Leib und Seele, daß wir recht gläubig und christlich leben, bis an unser selig Ende, Amen.

\chapter{Das 21. Kapitel}

\section{Vorrede}

Freilich ist die heilige Biblia wie ein großer schöner Garten, darin viel und allerlei Arten Bäume stehen, davon man auch allerlei gut Obst und Früchte haben kann, die zur Speise, Arznei und Labsal dienen. Denn die heilige Biblia gibt fleißigen Zuhörern und andächtigen Herzen reichen Unterricht, Vermahnung, Warnung, Verheißung und Trost, die nützlich sein in Leib und Leid, Noth und Tod. Derowegen laßt uns die heilige Biblia also anhören und lernen, auf daß wir täglich zu unserer Besserung, und zu unserm Trost, etwas mit zu Hause nehmen, allezeit behalten, fruchtbarlich im Leben und Sterben gebrauchen. Folgendes höret an das 21. Kapitel.

\subsection{Isaaks Geburt. Austreibung Ismaels. Abrahams Bund mit Abimelech}

\textit{Zu Vers 17.}a Hagar) mercke hier auf Hagar, wie die des Gesetzes und glaubensloser Werck Figur ist, Gal. 4, 22, und dennoch sie Gott zeitlich belohnet und groß machet auf Erden.
\quellseite{73}

\section{Summarium}

Ist eine Historia, wie Sarah in ihrem Alter einen Sohn träget. Auf daß wir sehen, wie bei Gott nichts unmöglich und sein Wort allein wahr sei, und wir uns desto gewisser darauf verlassen.

Hie hebet sich nun die Spaltung der Kirchen: Ismael bedeutet, wie St. Paulus Gal. 4, 21. sagt, das Volck, so durch die Werke des Gesetzes will selig werden. Aber da wird nichts aus: der Ismael muß aus dem Haus mit seiner Mutter und sich am zeitlichen Segen begnügen lassen, Isaac aber bleibet im Hause und ererbet den Segen, B. 10. Das sind die, so nicht auf ihre Werke, sondern auf Gottes Güte und Barmherzigkeit ihr Vertrauen setzen.

Ueber das ist hie auch ein fein Exempel, wie Gott Ismaels Gebet erhöret, der nun ein Knabe bei vierzehn Jahren alt war und sein Leben rettet, B. 17.; auf daß die jungen Leute auch lernen, mit Ernst beten und fromm sein. Denn Gott will gerne sie schützen und in aller Noth helfen.

\textit{Zu Vers 31.}b Versaba) heißt auf deutsch Schwörbrunn oder Eidsbrunn, möchte auch wohl Siebenbrunn heißen.

\section{Summarium}

Am Ende dieses Kapitels ist ein Exempel, daß man in bürgerlichen Sachen der Einigkeit zu gut, ohne Sünde möge Eid schwören, wie Abraham und Abimelech thun, B. 23. Darnach stehet, Abraham habe geprediget vom Namen des HErrn, des einigen Gottes, B. 34., das ist, von der Verheißung von Christo: wie Gott durch Christum in unserm Fleisch sich offenbaren und ein ewiges Reich werde anrichten. Darum sagt Christus im Evangelio, Abraham habe seinen Tag gesehen und sich gefreuet.

\section[Notum]{\unsicher{Notum}}

GOtt der HErr erleuchte uns, daß wir uns auch in wahrem Glauben unsers HErrn JEsu Christi freuen; in unsern Nöthen fleißig beten, und auf sein Wort uns beständig verlassen, bis zum ewigen Leben, Amen.

\chapter{Das 22. Kapitel}

\section{Vorrede}

Eines jeden Christen Menschen Leben soll eine tägliche Buße sein. Dieses lehret der heilige Apostel Paulus, da er sagt, 1 Tim. 1.: Uebe eine gute Ritterschaft, habe den Glauben und gut Gewissen. Derowegen wollen wir heute abermal unsere Sünde erkennen und sonderlich diese, daß wir alle schwach im Glauben sind, und daß wir noch nie Gott dem HErrn den Gehorsam geleistet haben, den wir ihm schuldig sind. Und wollen ihn herzlich bitten, daß er uns, aus Gnaden um Christi willen, alle unsere Sünde vergebe, unsern Glauben uns vermehre und stärke und uns erfülle mit Früchten der Gerech-
\quellseite{74}
tigkeit, auf daß wir sind lauter und unanstößig bis auf den Tag Christi unsers Herrn, Amen. Daneben höret an das 22. Kapitel.

\subsection{Aufopferung Isaaks. Verheißung von Christo}

\textit{Zu Vers 2.}a Moria) Moria heißt Gottesfurcht, Reverentia Dei, Cultus Dei. Denn die Altväter Adam, Noah, Sem, auf demselbigen Berge, Gott gehret, gefürchtet, gedienet haben. Wir Deutschen heißen es vielleicht den heiligen Berg, oder da man Gott dienet mit Loben, Beten und Danken.

\textit{Zu Vers 8.}b Ersehen) Gott siehet und weiß wohl, wo das Schaf sei, läßt ihn dafür sorgen, er siehet es besser, denn wir.

\textit{Zu Vers 14.}c Siehet) Hebrei dicunt, Dominus videbitur. Sed nos Hieronymum secuti Rabinos Grammaticos cum suis punctis et Cametz hoc loco negligimus, et sine punctis dicimus: Der HErr siehet, das ist, Gott sorget für alles und machet. Etiamsi sensus ille: Dominus videbitur, sit pius valde, quod Deus appareat, ubi verbum ejus docetur, quod Rabini Grammatici non intelligunt, Heb. 6, 14.

\textit{Zu Vers 18.}d Gehorschet) hie wird Abraham nicht gerecht durch seinen Glauben für sich, sondern erweiset solche Herrlichkeit seines Samens; denn er zuvor gerecht ist, wie oben Cap. 15, 6.
\quellseite{75}

\section{Summarium}

Dies ist ein schön Exempel eines sonderlichen großen Glaubens und völligen Gehorsams gegen Gott, daß Abraham Gottes Gebot folget und seines einigen Sohnes nicht verschonet; er will, wie Gott ihm befohlen hatte, ihn aufopfern, V. 6., so doch die Verheißung alles zeitlichen und geistlichen Segens allein auf dieses Sohnes Leben stund.

Auf daß wir auch lernen, uns ganz Gott ergeben und dennoch gewisse Hülfe erwarten; wie Abraham thut: läßt ihm seinen Sohn nicht zu lieb sein und denkt, Gott könne seine Zusage vollstrecken, wenn Isaac schon aufgeopfert werde. Solch gewiß Vertrauen auf Gottes Zusage ist Gottes höchster Dienst und Wohlgefallen.

Merke aber, daß der Text sagt: Gott versucht Abraham, V. 1., das ist, es war keine Meinung nicht: daß Abraham sollte seinen Sohn tödten. Denn Todtschlag ist verboten; er wollte nur dabei Abrahams Glauben, Gehorsam und Furcht erfahren.

\section{Votum}

Gott regiere uns alle, wie er den heiligen Abraham geführet hat; auf daß wir auch unser Licht leuchten lassen, daß die Menschen unsere gute Werke sehen, und Gott allezeit gelobet und gepreiset werde, Amen.

\chapter{Das 23. Kapitel}

\section{Vorrede}

Der Mensch lebet, und, wenn es ihm Gott noch so lange gönnet, so muß er doch endlich sterben. Wohl aber allen, die im Herrn entschlafen. Derowegen weil wir alle leben und wissen nicht, wie lange, und weil wir alle sterben und wissen nicht, wann; so wollen wir alle täglich Buße thun und also glauben, beten, handeln und wandeln, ob wir wohl zeitlich sterben müssen, daß wir doch mit Gott ewig leben, um Jesu Christi, unseres Heilandes, fröhlichen Geburt, bittern Leidens, schmählichen Todes, siegreichen Auferstehung und herzquickenden Himmelfahrt willen. Darauf höret nun ferner an das 23. Kapitel.

\subsection{Sarahs Tod und Begräbniß}

\textit{Zu Vers 2.}a Hebron) Hebron ist Kirjath Arba (spricht Moses), das ist, die Vierstadt; denn die hohen Hauptstädte waren vor Zeiten alle Arba, das ist, in vier Theile getheilet, wie Rom, Jerusalem und Babylon auch, 1 Mos. 10, 10.

\textit{Zu Vers 15.}b Sekel) Sekel ist ein Gewicht, an der Münze ein Drittheil gülden; denn vor Zeiten man das Geld so wog, wie man jetzt mit Gold thut.
\quellseite{76}

\section{Summarium}

Das ist eine Historia, wie Abraham seine liebe verstorbene Sarah beweint und beklaget und ihr ein eigen Begräbniß kaufet. Zum Zeugniß, daß er fest glaube, ob er wohl ein Fremdling im Lande sei, kein Theil darin habe, auch nicht eines Fußes breit, daß gleichwohl sein Same nach ihm dasselbe soll zu eigen haben.

Ist auch ein fein Exempel hie, wie im Kaufen und Verkaufen fromme Leute ehrbarlich und aufrichtig gehandelt haben. Wie solches noch wohl stehet, aber leider an solcher Ehrbarkeit täglich viel abgehet.

\section{Votum}

Helfe der treue GOtt, daß wir also leben und sterben, auf daß wir alle auferstehen zum ewigen Leben, das uns erworben hat JEsus Christus unser Heiland, hochgelobt mit dem Vater und heiligen Geist in Ewigkeit, Amen.

\chapter{Das 24. Kapitel}

\section{Vorrede}

Nach Gottes Wort ist der größte Schatz auf Erden, in einem gewissen und rechten Beruf stehen, und nichts anfangen oder fürnehmen ohne Gottes Befehl. Und wer derowegen seinen christlichen Beruf also treibet, daß er Gott um Segen bittet, ist fleißig und treu, und lebt in Gottesfurcht, dem gelinget es, daß er Gott dafür genug zu danken hat. Dieses lehret unser HErr JEsus Christus, da er sagt, Matth. am 7.: Bittet, so wird euch gegeben werden; und abermal, da er spricht: Trachtet am ersten nach dem Reich Gottes und nach seiner Gerechtigkeit, so wird euch das andere alles zufallen, Matth. am 6. Dessen ist ein Exempel das 24. Kapitel.

\subsection{Isaak heyrathet die Rebekka}
\quellseite{77}

\section{Summarium}

Ist eine Historia, in welcher wir sehen, wie die Eltern sich fleißigen sollen, daß ihre Kinder ehrlich und bei frommen Leuten heirathen. Und ist sehr tröstlich, daß Abraham und sein frommer Knecht Gott zum Heirathsmann oder Ehestifter nehmen und bitten, daß er den Isaac wohl versorgen wolle, B. 3. Und daß sie die Ehe so groß und ehrlich für Gott achten, daß Gott seine Engel dazu verordnen werde, daß es alles glücklich zugehe, B. 7. Dermaßen wird die Ehe bei den Papisten nicht geachtet, die es für einen solchen Stand achten, da Gott nichts mit zu thun habe, ja da man Gott nicht recht könne dienen.
\quellseite{78}

\textit{Zu Vers 47.}a Spangen) diese güldene Spange ist gewesen ein halber Zirkel auf der Stirn bis zu beiden Ohren; darum heißt ers jetzt Ohrenringe, jetzt Stirnpangen. Und siehet, als habens beide Mann und Weibsbilde getragen zum Schmuck. Wie unten Kap. 35, 4. Sprüchw. 11, 22. Circulus aureus in naribus suis. Das sagen wir Deutschen, die Sau gefrönet. Inde diadema Regum et lamina summi Sacerdotis in fronte etc.

\textit{Zu Vers 49.}b Saget mirs) er handelt zuvor mit Mutter und Bruder um die Braut. Daraus man siehet, daß heimliche Verlöbniß ohne Vorwissen der Eltern nicht recht ist.

\textit{Zu Vers 53.}c Würze) köstliche Früchte.

\textit{Zu Vers 57.}d Dirne) die Braut soll ungezwungen zur Ehe gegeben sein, von den Eltern, dazu auch gefragt werden um ihren Willen.
\quellseite{79}

\section{Summarium}

Ist ein fein Exempel, wie der Rebekka Freundschaft mehr auf Gottes Schickung stehet, denn auf Geld und Gut. Darum sind sie willig und lassen die Tochter in ein fremdes Land dem Isaak gerne folgen.

Wie sich nun der Bräutigam und der Werber dagegen halten, ist auch fein zu merken: Isaak gehet aufs Feld und betet, V. 63. Das soll ein Bräutigam thun, soll anders Glück und Heil, wie wohl es bedarf, dabei sein. So danket der Werber auch Gott, daß er seine Reise also gefördert und Glück dazu gegeben hat, V. 56. In Summa, du siehest, wie diese fromme, heilige Leute alles mit Gottesfurcht in rechtem Glauben und mit ernstem Gebet anfangen; darum ist auch Gottes Segen so reichlich bei ihnen. Wir sollten aber aus solcher Historien uns bessern und dergleichen und auch zu all unserm Thun gottseilig und christlich schicken. Aber jedermann thut das Widerspiel: Gott fürchten wir nicht, um seinen Segen bitten wir nicht, schelten, fluchen ist das erste; darum stehets und gehets so übel in aller Welt.

\section{Gebet}

Verleihe der treue GOtt, daß wir auch fleißig und herzlich beten und gottsfürchtig leben, und also aus Gnaden seinen Segen an Leib und Seele haben und behalten, bis zum freudenreichen ewigen seligen Leben, Amen.

\chapter{Das 25. Kapitel}

\section{Vorrede}

Mitten im Leben sind wir mitten im Tode. Denn wie wir von wegen der Sünden alle Augenblick mancherlei und tödtlichen Krankheiten

unterworfen sind: also ist kein Unfall so groß und schwer, er kann uns bald zu Haus und Hof kommen und uns unversehens dahin reißen. Darum bedürfen wir allezeit der Gnade Gottes, daß er uns behüte und erhalte. Das aber solches uns widerfahre, so sollen wir auch alle in christlicher Buße leben, auf daß uns der Tod in unsern Sünden nicht übereile, daß wir in Sünden zum ewigen Verderben nicht umkommen und sterben. Ferner höret an das 25. Kapitel.

\subsection{Abrahams andere Ehe. Ismaels und Isaaks Söhne}
\quellseite{80}

a Er fiel, mancherlei Deutung kann hie sein. Meine Deutung ist diese, daß Ismael ein herrlicher Mann gewesen sei, daß zu seinem Ende kommen und alle seine Brüder und Freunde, und ist für denselben ehrlich und löblich gestorben.

\section{Summarium}

Dies ist eine Historia, wie Abraham gestorben sei, v. 7.; und daß sein Sohn Isaak eben die Anfechtung hat müssen tragen, wie sein Vater, daß er zwanzig Jahr in der Ehe gewesen, und dennoch keinen Erben hat gehabt, darum, daß Rebekka unfruchtbar war, v. 21.

Weil nun aber die Verheißung auf Isaaks Erben stund, läßt er sich solches Aergerniß nicht abschrecken, sondern hält sich an Gottes Gebet, und erlanget einen, Zwillinge, Jakob und Esau, v. 26., welche gleich wie Ismael und Isaak bedeuten die zwo Kirchen: eine, so auf Gottes Barmherzigkeit traut; die andere, so auf eigen Werk und Frömmigkeit bauet.

Daß Esau die Herrlichkeit der Erstgeburt um eine Schüssel voll Linsen verkauft, v. 31., welche fürnehmlich hatte die Verheißung von Christo, ist ein sonderlich Exempel des Unglaubens, daß er solche Verheißung für nichts geachtet hat; Jakob aber, als ein gottseliger Mensch, achtet dieselbige hoch und theuer, darum erlanget er sie auch, und wird solches Glaubens halben Jakob vom Propheten Hosea Kap. 12, 3. herrlich gepriesen. Denn daß er in der Geburt den Esau bei den Fersen hält, v. 26., solches deutet Hoseas auf den Glauben Jakobs, daß er gern die Erstgeburt und die Verheißung vom Herrn Christo gehabt und seinen Bruder daran wollte gehindert haben.

\section{Votum}

Der barmherzige treue GOtt stärke unsern Glauben, daß wir sein heilig Wort lieben und üben, christlich leben, alles Creutz geduldig tragen und endlich selig sterben, hochgelobt in Ewigkeit, Amen.
\quellseite{81}

\chapter{Das 26. Capitel}

\section{Vorrede}

Der Welt Art und Gebrauch ist, daß sie die Christen verfolgt, beleidigt und beschädigt: Gottes Art und Weise dagegen ist, daß er die Christen vertheidigt, schützet und segnet. Derowegen wer den Welt-Kindern verhalten muß, der befehle es Gott, räche sich nicht selbst, bete und sei in seinem Beruf fleißig. So wahr als Gott im Himmel lebet, er wird dem Gerechten helfen und seine Stärke sein in der Noth; er wird die Gläubigen versorgen und wird sie nicht ewiglich in Unruhe lassen, Ps. 55. Dieses lehret auch das 26. Kapitel.

\section{Isaaks Wanderschaft, Wohnung und Widerwärtigkeit}

a Esek heißt Unrecht, wenn man jemand Gewalt und Unrecht thut.

b Sitna heißt Widerstand; daher der Teufel Satan heißt, ein Widerwärtiger.

c Rehoboth heißt Raum und Breite, das nicht enge ist.
\quellseite{82}

d Saba heißt ein Eid und Schwur, oder die Fülle, der aber heißt ein Brunn.

\section{Summarium}

Dies ist eine Historia, in welcher wir sehen, wie Isaak Theurung halber an einen fremden Ort ziehen, V. 1., und daselbst nicht allein seiner Rebekka halben in Gefahr leben, sondern auch allerlei Ueberlaß von den Philistern leiden muß, V. 7. Wir sollen daraus lernen, wie unser Gott seine Christen in der Welt unter bösen Buben in mancherlei Weise plagen läßt, aber doch ihnen auch väterlich hilft und sie schützet. Denn je mehr die Welt die Frommen will dämpfen, je mehr er sie segnet und ihnen hilft.

\section[Gebetum]{\unsicher{Gebetum}}

Gott steure und wehre den Verfolgern seiner Christenheit, und segne und behüte alle, welche die bedrängten Christen fördern, nähren und schützen, und thue solches um Jesu Christi seines lieben Sohns willen, Amen.

\chapter{Das 27. Kapitel}

\section{Vorrede}

Gott ist ein Brunnquell alles Guten denen, die ihn fürchten; dagegen ist er ein verzehrend Feuer allen, die ihn verachten. Derowegen wer von Gott will gesegnet sein, der diene ihm in wahrem Glauben, und dem Nächsten in wahrer Liebe, daß ist, er vertrauhe auf Gottes Verheißung und beschädige nicht seinen Nebenchristen; sondern rede und thue, das ihm in seinem Anliegen nöthig ist. Also wird er von Gott aus Gnaden bekommen und haben, was ihm und den Seinen gut und nützlich ist; und dagegen wird ihn Gott aus Barmherzigkeit behüten und schützen, wider alles, was ihm und den Seinen böse und schädlich ist. Das unterrichtet auch das 27. Kapitel.

\section[Jacob wird von seinem Vater gesegnet, von Esau]{Jacob wird von seinem Vater gesegnet, von Esau angefeindet}
\quellseite{83}

\section{Summarium}

Dies ist eine Historia, wie Esau zwei heidnische Weiber nimmt, ohne Zweifel wider seiner Eltern Rath und Willen, V. 1., und der Segen auf den frommen Jakob kommt, welchen Segen zuvor Esau aus dem Unglauben verachtet und um ein Essen Linsen ihm verkauft hatte, und zur Warnung; denn wer Gott und seinen Segen verachten kann, demselben soll der Fluch zu Theil werden.

Daß aber Jakob durch seine Mutter mit einer List zum Segen gefördert wird, V. 6., geschieht aus dem Befehl des Herrn, da er sagt: Der Größere wird dem Kleinern dienen.
\quellseite{84}
Maße sehr betrübt; und sprach zu seinem Vater: Segne *mich auch, mein Vater.

a) Untertreten) Ebr heißt eine Fußsohle. Daher kommt Jakob oder Jacob, ein Untertreter, oder der mit Füßen tritt, und bedeutet alle Gläubigen, die durch das Evangelium die Welt, das Fleisch und den Teufel mit Sünde und Tod unter sich treten durch Christum.

\section{Summarium}

Dies ist eine Historia, wie der Unglaube und die Verachtung der Verheißung Gottes am Esau leiblich gestraft wird. Denn auf dem erstgebornen Sohn stund das Regiment und die Verheißung von Christo. Weil nun Esau solches beides nicht größer achtet, denn ein Essen Linsen und Bauchfülle, nimmt Gott solchen Segen von ihm und gibt ihn den Jakob.

Dies ist nun die Ursache, daß Esau den Jakob verfolget. Wie in der Welt je und je geschehen ist, daß die rechte christliche Kirche, so sich Gottes Güte tröstet, von den Werk- und Bauchheiligen muß verfolget werden.

\section{Votum}

Der allmächtige treue GOtt steure allen Feinden seiner Christenheit, und erhalte seine Gläubige unter seinem Schutz, in Leibes und der Seelen Gefahr, bis zum ewigen Leben, Amen.

\chapter{Das 28. Kapitel}

\section{Vorrede}

Wenn die Kinder ihren treuen Eltern und die Eltern ihrem frommen Gott gehorsam sein, so gelinget es beiden. Denn die Kinder, so ihre Eltern ehren, die ehret Gott wieder an Leib und Seele, in diesem und künftigen Leben; also, die Eltern, die in Gottesfurcht leben, sich treulich nähren und Gottes Segen recht gebrauchen, die genießen es auch zeitlich und ewig, beides aus Gnaden, um Jesu Christi willen. Derowegen Kinder und Eltern fleißig auf Gott und seine Gebote sehen und denen treulich und beständig folgen sollen. Davon höret auch an das 28. Kapitel.

\section{Jakob ziehet auf seiner Reise die Himmelsleiter}
\quellseite{85}

a Deinen Samen) hie wird dem dritten Patriarchen Christus verheißen, der Heiland aller Welt, und das künftige Evangelium von Christo in allen Landen zu predigen durch die Engel auf der Leiter vorgebildet.

b Heilig) heilig heißt hie metuendus, terribilis, Nota, da man GOtt fürchten und ehren sollte, als er daselbst will gefürchtet und geehrt sein. Daher auch derselbe Berg, Morija, timor, reverentia, cultus Dei. heißt, sup. Kap. 22. 2. Denn Gottesfurcht ist der höchste Gottesdienst. Und ist hie angezeigt, wo Gottes Wort ist (wie Jakob hie hört) da ist Gottes Haus, da stehet der Himmel offen mit allen Gnaden.

c Mein Gott sein) nicht daß er vorhin nicht sein Gott gewesen sei, sondern er gelobet einen Gottesdienst aufzurichten, da man predigen und beten sollte; da will er den Zehnten zu geben den Predigern, wie Abraham dem Melchisedec den Zehnten gab.

\section{Summarium}

Dies ist ein fein Exempel, wie Jakob seinen Eltern folget und nach ihrem Willen sich begehret zu verheirathen, V. 1. Denn in solchem Falle die Kinder sonderlich auf der Eltern Willen sehen und sich befleißigen sollen, daß sie wider ihrer Eltern Willen und Vorwissen nichts fürnehmen. Wie man siehet, daß selten Glück dabei ist, wo die Kinder wider der Eltern Willen heirathen und die Eltern hoch betrüben. Denn Gott will Eheleuten Glück geben, wo sie Gott fürchten, Pf. 128, 1. 4. Wo man aber den Eltern nicht folget und wider ihren Willen thut, da kann keine Gottesfurcht sein, wird auch kein Glück da sein.

Darnach ist eine andere Historia, wie Gott dem Jakob erscheinet in dem Traum und bestätiget mit einer neuen Verheißung den Segen von Christo, und dem Lande Canaan, welches Esau verachtet hatte, V. 10. Darum verpflicht sich Jakob wieder gegen Gott, er wolle alle seine Hoffnung und Trost allein auf solchen gnädigen Gott setzen, ihm dienen und an dem Ort ein Gottes-
\quellseite{86}
haus bauen, und mit dem Zehenten aller seiner Güter den rechten Gottesdienst am selben Ort fördern und erhalten, V. 20. Und zum Exempel, daß wir auch lernen Gottes Segen also erkennen, denselbigen zu Erhaltung seines Worts und des rechten Gottesdienstes brauchen.

\section{Datum}

Helfe der treue barmherzige GOtt, daß alle Kinder und Eltern ihm allzeit gehorsam seyn, und segne beide an Leib und Seele, zu zeitlicher und ewiger Wohlfahrt, Amen.

\chapter{Das 29. Kapitel}

\section{Vorrede}

Wie man im Alten Testament alle Morgen frisch Holz zum Altar des HErrn bringen und damit das Feuer, daß es nimmermehr verlöschte, erhalten mußte: also sollen wir auch alle Morgen und Abend neu Holz vom Baum des Lebens, welcher ist unser HErr JEsus Christus, das ist, von seinem heiligen Wort, in unser Herz legen, damit das Feuer des Glaubens, der Liebe gegen Gott und unsern Nächsten, und des Gebets, welches der heilige Geist in uns anzündet, nicht ausgelöscht und vertilget werde. Wie denn ein Feuer auslöschen muß, wenn man kein Holz anleget. Denn wir sollen unser Licht brennen lassen, vermahnet Christus unser HErr, Luc. 12. Wir sollen uns auch unter einander brünstig lieb haben, lehret Petrus 1 Petr. 1. Und wir sollen mit Ernst Gott anrufen und mit dem Geist und Sinn beten, unterrichten David und Paulus, Ps. 145. 1 Cor. 14. Ferner höret an das 29. Kapitel.

\section{Jakob erwirbt durch seinen vierzehnjährigen Dienst zwei Weiber}
\quellseite{87}

\section{Summarium}

Dies ist eine Historia, wie Jakob auch in Ehestand kommen sei, V. 20. Auf daß man an der heiligen Patriarchen Exempel sehe, welche doch Gottes Wort geprediget, und mit Opfern als rechte Priester umgegangen und Gott gedienet haben, daß der Ehestand eine heilige Weihwohnung sei; welchen doch der Papst als unheilig seinen Geistlichen verboten hat. Zum andern ist ein fein Exempel einer sonderlichen Geduld und Einfalt des Jakobs: welcher, ob ihm wohl die andere Schwester ward zugelegt, solches Unrecht duldet, und noch sieben Jahr um die Rahel dienet, so ihm verlobet war, V. 28. Zum dritten stehet man hie, wie Kinder Gottes Gaben sind, welche Gott auch darum der Lea gibt, auf daß er sie ihrem Jakob lieb und werth mache, weil sie ungeschaffener war, denn die Rahel, V. 31. Denn Gott will, daß zwischen Eheleuten alle Liebe und Freundschaft gehen soll; gleich wie der böse Feind alles dahin wendet, daß er zwischen Eheleuten Uneinigkeit anrichten und die Herzen trennen möge.

\section{Votum}

Der gnädige gütige Gott segne alle christliche Eheleute an Kindern und Nahrung, verleihe auch ihnen Geduld unter ihrem Hauskreuz, und helfe, daß sie und alle ihre lieben Kinder in seiner Furcht leben und sterben, Amen.

\chapter{Das 30. Kapitel}

\section{Vorrede}

Das Alte Testament ist ein Gesetzbuch, das da lehret, was man thun und lassen soll, und daneben anzeiget Exempel und Geschichte, wie solche Gesetze gehalten oder übertreten sind. Gleich wie das Neue Testament ein Evangelium oder Gnadenbuch ist, und lehret, wo man es nehmen soll, auf daß das Gesetz erfüllet werde. Wie aber Moses nun rufet dem ganzen Israel und sprach zu ihnen: Höre Israel die Gebote und Rechte, die ich euch heute für euren Ohren rede, und lernet sie, und behaltet sie, daß ihr darnach thut, 5 Mos. 5.: also sagt unser Herr Christus zu allen seinen Christen: Sie haben Moses und die Propheten, laß sie dieselben hören, Luc. 16. Derowegen lasset euch die heilige Biblia, und also Gottes Wort, treulich befohlen sein und betet, daß ihr darnach züchtig, gerecht und gottselig allezeit lebet, Amen. Und darauf höret ferner an den ersten Theil des 30. Kapitels.

\section{Jakobs Kinder und Reichthum}
\quellseite{88}

a Naphtali heißt verwechselt, umgewandt, umgekehrt, wenn man das Widerspiel thut. Ps. 18, 27.: Mit den Verkehrten verkehrst du dich.

b Dudaim) fragst du selbst, was Dudaim sind: es sollen Lilien, es sollen Beeren sein, und niemand weiß, was es sein sollen. Es heißen etliche Judenkirschen, die in der Weizen-Ernte reif sind.

\section{Summarium}

Dies ist eine Historia, wie Gott Jakob gesegnet und ihm Kinder gegeben hat; daß aber die zwo Schwestern so peinlich nach Kindern thun, und eine die andere neidet, geschicht alles darum, daß sie es für einen sonderlichen Segen und Gnade von Gott achten, wenn Gott Eheleuten Kinder gibt. Sonderlich aber, weil dazumal der verheißene Same Christus zu warten war, durch welchen alle Völker sollten gesegnet werden. Solche Hoffnung hat ohne Zweifel die frommen gottsfürchtigen Matronen getrieben, daß sie groß Verlangen nach der Leibesfrucht gehabt, ob doch dermaleinst Gott die gnädige Stunde wollte geben, daß dem Teufel seine Gewalt genommen und allen Menschen von Sünden und Tod geholfen würde.

\section{Votum}

Der Sohn Gottes, unser Herr Jesus Christus, lasse ihm alle gottseligen Eheleute befohlen sein, und behüte sie für dem Teufel und für allem Herzleid, Noth und Schaden, Amen.

\section{Der andere Theil des 30. Kapitels}

\section{Vorrede}

Wer Glück und Gedeihen zu seiner Nahrung haben will, der trachte am ersten nach dem Reich Gottes und nach seiner Gerechtigkeit, der bete fleißig um Gottes Segen dazu, der habe auch das Vertrauen zu Gott, daß seine Arbeit nicht vergeblich sein soll, und der treibe daneben seinen
\quellseite{89}
Beruf ohne des Nächsten Betrug und Schaden. Denn fleißig arbeiten und sich treulich nähren, aber ohne Gottesfurcht, ohne Gebet, ohne wahren Glauben an Gott, und mit anderer Leute Nachtheil, das bringet der Nahrung endlichen Verderb und Untergang. Das ist gewiß und Amen. Ferner höret an den andern Theil des 30. Kapitels.

Meinen Fuß) Das ist, ich habe müssen laufen und rennen durch dünn und dick, daß du so reich würdest. Mein Fuß hat müssen thun.

Unde pedes Evangelizantium pacem et cursus verbi seu ministerii das ist, daher die Füße derer, die das Evangelium und den Frieden predigen und der Lauf des Wortes- oder Predigtamtes.

d Du mußt dich hie nicht irren lassen, daß Moses das kleine Vieh jetzt Ziegen, jetzt Lämmer, jetzt Böcke heißt, wie dieser Sprache Art ist. Denn er will so viel sagen, daß Jakob habe alles weiße Vieh behalten und alles bunte und schwarze Laban gethan. Was nun bunt von dem einfärbigen Vieh käme, das sollte sein Lohn sein. Es ward Laban froh, und hatte die Natur für sich, daß von einfärbigen nicht viel bunte natürlich kommen; aber Jakob half der Natur mit Kunst, daß die einfärbigen viel bunte trugen.

\section{Summarium}

Ist eine Historia, wie Jakob als ein Hausvater auch für sein Weib und Kinder sorget, und ein Gedinge mit Laban macht seines Lohnes halben, doch also, daß er alles stellet, nicht auf seine Arbeit, sondern auf Gottes Segen, und zum Exempel, daß wir also unsere Nahrung suchen sollen, daß weder Geiz noch Uebervortheilen bei uns sey, sondern neben unsrer treuen Arbeit alles auf Gottes Segen gestellet werde. Solche göttliche Fürsorge des heiligen Patriarchen Jakobs sollten die unachtsamen vertrumbischen Hausväter merken, welche nichts für sich sparen, alles anwenden und nach Weib und Kind, welche doch Gott ihnen zu nähren befohlen, auch Geschicklichkeit und andern Vortheil dazu verliehen, nichts fragen. Eine gräuliche Sünde ist solches und wird mit mancherlei Unrath gestraft, wie vor Augen.
\quellseite{90}
Zum andern ist auch ein fein Exempel hie, wie um eines frommen Knechts willen Gott alles Glück dem Herrn und ganzem Hause widerfahren läßt.

\section{Votum}

Regiere und führe der treue GOtt die Herrschaft und das Gesinde, daß ein jedes seines Berufs also abwarte, auf daß es GOtt zu Lob, ihrem Nächsten zu gut, und ihnen selbst zur Wohlfahrt diene und gelange, Amen.

\chapter{Das 31. Kapitel}

\section{Vorrede}

Wie wir uns gegen Gott halten, nach seinen Geboten: also hält sich Gott gegen uns mit seinen Verheißungen. Denn weil in allen Ständen, hohen und niedrigen wider seine Gebote sehr gesündigt wird, so entzieht er uns auch an vielen Orten seine Zusagungen, und also an Leib, Seele und Nahrung, Segen und Wohlfahrt. Derowegen wer da will, daß ihn Gott benedeie und behüte, wie er versprochen hat, der lebe und wandele auch in seinem Beruf, Amt, Handel, Ober und Handthierung, wie ihm Gott befohlen hat. Daher sagt der Prediger Salomon im 8. Kapitel: Wer das Gebot hält, der wird nichts Böses erfahren. Daneben höret an den ersten Theil des 31. Kapitels.

Jakob ziehet mit Weib und Kind von seinem Schwiegervater Laban.

* Stahl das Herz, Herz stehlen ist, betrüglich geredt, so viel, als etwas thun hinter eines andern Wissen.

\section{Summarium}

Ist ein schön Exempel, darin wir sehen, wie Gott alle Treue den Frommen reichlich belohnen, wiederum alle Untreue damit will strafen, daß er seinen Segen dazu geben will. Denn ob wohl Laban den frommen Jakob nicht mit Treuen mei-
\quellseite{91}
net, weil aber Jakob mit allen Treuen ihm dienet, läßt ihn Gott solcher Treu genießen, Laban muß aber seiner Untreue entgelten und soll ihm kein Vortheil helfen. Auf daß wir auch lernen treulich und ungefährlich handeln, so anders Gottes Segen soll bei uns sein. Es ist in aller Welt, an allen Orten, sehr viel armes Volk, das Nahrung halben weder hinter sich noch für sich kann. Wer solches Unfalls Ursache gewiß will wissen, mag es hie lernen: untreulich dienen sie, sind nachlässig, unachtsam, ungehorsam. Solche Sünde muß darnach an ihnen gestraft werden, wenn sie zu Haus ziehen und auch in den Stand gerathen, da sie Gesinde halten und mit denselben sich nähren müssen.

\section{Votum}

GOtt vergebe allen ihre Sünde, die Vergebung begehren und bitten; und verleihe, daß ein jedes willig und fleißig wandele in seinem Beruf, wie es GOtt in seinem Wort, treulich und ernstlich, lehret und erfordert, Amen.

\section{Der andere Theil des 31. Kapitels}

\section{Vorrede}

Treue in allerlei christlichen Händeln und Gewerben belohnet Gott mit Segen als ein barmherziger, gütiger Vater: Untreue strafet Gott mit Ungefallen als ein gerechter, wahrhaftiger Richter. Derowegen ein jedes in seinem Beruf und Nahrung treulich und ungefährlich handeln soll, so will ihn Gott auch segnen und für Gefahr behüten; oder will ihn unter dem Kreuz trösten und erhalten. Dies lehret der königliche Prophet David im 128. Psalm, und spricht: Wohl dem, der den Herrn fürchtet und auf seinen Wegen gehet. Du wirst dich nähren deiner Hände Arbeit, wohl dir, du hast es gut. Und St. Petrus in Geschichten der Apostel im 10. Kapitel sagt also davon: In allerlei Volk, wer Gott fürchtet und recht thut, der ist ihm angenehm. Zu diesem höret auch an den andern Theil des 31. Kapitels.
\quellseite{92}

\section{Summarium}

An ein schön Exempel, in welchem wir sehen, wo man anders fromm und treu ist, daß Gott die Seinen nicht lassen, sondern für Unglück bewahren will. Denn ehe Jakob weiß, daß ihm sein Schwäher Laban nacheilet, der Meinung, daß er ihm will Schaden thun, nimmt sich Gott des frommen Jakobs an, strafet und bedrohet Laban, daß er mit Jakob nicht anders denn freundlich handeln solle. Darum fleißige man sich nur mit allem Ernst, daß man fromm sey und mit gutem Gewissen mit den Leuten handle; so wird Gott sich dermaßen auch gegen uns erzeigen und uns schützen und retten, wie den gottseligen frommen Jakob.

\section{Gebet}

Wir danken Gott für allen Segen, den er von Anfang seinen Gläubigen gegeben hat; der segne und behüte uns alle ferner, doch nach seinem gnädigen väterlichen Willen und Wohlgefallen, Amen.

\chapter{Das 32. Capitel}

\section{Vorrede}

Freilich ists wahr und recht gesagt: Je größer Noth, je näher Gott. Denn solches lehret Gottes Wort, dieses bezeugen auch die Exempel der heiligen Gottes. Dergleichen erfahren wir solches alle Tage an den Gläubigen. Derowegen sei niemand unter seinem Kreuze ungeduldig, wie lange es auch währet und wie schwer solches ist; sondern ein jeder Kreuz-Bruder und jede Kreuz-Schwester bete und gebrauche gebührliche und christliche Mittel. Der

b Furcht) Jakob nennet die Gott Isaaks Furcht darum, daß Isaak gottfürchtig war und Gottes Diener.

c Gilead) heißt ein Zeugenhaufe, und bedeutet die Schrift, da viele Zeugnisse von Gott häufig innen sind.
\quellseite{93}
Gott, der andern bisher wunderbarlich geholfen, wird euch auch in seiner Noth \unsicher{verlassen}, denn er ist Gott über alle Noth. Und darauf höret an das 32. Kapitel.

\section{Jakobs Furcht, Gebet und Kampf}

\section{Summarium}

Das ist ein schön Exempel, wie Jakob sich für seinem Bruder Esau demüthiget, und weil er sich für ihn besorgen muß, seine brüderliche Mittel suchet, den Zorn damit abzulehnen, V. 3. Am allerersten hebt er die Sache mit Gott an, und danket ihm um seinen reichen Segen: Da ich, spricht er V. 10., über diesen Jordan ginge, hatte ich mehr nicht, denn diesen Stab, aber jetzt bin ich zwei Heere worden, das ist, ich habe einen großen Haufen Rind und Gesinde, und darnach einen großen Haufen Viehes, das alles dein Segen und Gabe ist und ich billig dir dafür danke. Darnach betet er, und erinnert Gott seiner Zusagung, daß er ihn gnädig für dem Zorn Esau behüten wolle, V. 11. Im letzten braucht er auch diese Mittel, daß er mit Geschenken des Bruders Zorn gedenkt zu versöhnen, V. 14. Also siehet man, wo die heiligen Leute ihrer Vernunft gebraucht haben, daß sie alles aufs einfältigste und zum Besten gethan haben.
\quellseite{94}

a Rang) im Hebräischen kommt ringen vom Staub her, als wenn zween mit einander ringen, daß der Staub sich erhebt, und dicke um sie wird, und lautet so viel: Es staubet ein Mann mit ihm, das ist, ein heftiger Kampf war es; das soll niemand verstehen, denn die Erfahrung.

b Israel) Israel kommt von Sara und heißt kämpfen oder überwältigen; daher auch Sar ein Fürst oder Herr, und Sara eine Fürstin oder Frau heißt, und Israel ein Fürst oder Kämpfer Gottes, das ist, der mit Gott ringet und angewinnet. Welches geschieht durch den Glauben, der so fest an Gottes Wort hält, bis er Gottes Zorn überwindet, und Gott zu eigen erlanget zum gnädigen Vater.

c Pniel) Pniel oder Pnuel heißt Gottes Angesicht oder Erkenntniß. Denn durch den Glauben im Streit des Kreuzes lernt man Gott recht erkennen und erfahren; so hats denn keine Noth mehr, so gehet die Sonne auf.

\section{Summarium}

Ist eine Historia, wie ein Engel mit Jakob ringt, und Jakob obsieget. Es zeuget aber Hoseas Kap. 12, 4. klar, daß dieser Engel nicht ein gemeiner Engel, sondern unser Herr Christus, der Sohn Gottes, selbst gewesen sei, und mit Jakob gerungen habe. Was aber solch Ringen sei gewesen, ist ungewiß, wiewohl aus dem Text abzunehmen ist, daß Jakob nicht allein mit der Leibes-Stärke, sondern viel mehr mit dem Glauben und Herzen gekämpft habe, und fest an der Verheißung Gottes gehalten, daß ihn Gott wieder heimbringen und seinem Samen nach ihm segnen wolle. Solcher Glaube hat dem Engel obgesieget, das ist, wie Jakob geglaubet, und der Zusagung Gottes sich getröstet hat, also ist es ihm auch widerfahren; Gott hats ihm nicht können abschlagen, er hat ihm müssen gewähren. Denn sein Wort und Zusagung ist ja, und kann nicht fehlen, wo man mit Herzen sich darauf verläßt.

\section{Votum}

Der treue barmherzige GOtt stärke unsern Glauben, und stehe uns bei, auf daß wir auch nach seiner Verheißung an Leib und Seele gesegnet, und zum ewigen Leben erhalten werden, Amen.

\chapter{Das 33. Kapitel}

\section{Vorrede}

Gegen Gott sollen wir uns als Kinder verhalten; und gegen unsern Nächsten sollen wir uns als Brüder und Schwestern erzeigen. Kinder aber ehren ihre Eltern; und einander dienen treue Brüder und Schwestern. Derowegen wollen wir Gott unsern Vater allezeit mit Mund, Herzen und guten Werken loben und preisen; und gegen unsern Nächsten, als unsern Brüdern und Schwestern, wollen wir uns allezeit freundlich und dienstlich, mit Rath und That, beweisen. Beides lehret der heilige Apostel Paulus. Denn 1 Cor. 10. sagt er: Was ihr thut, so thut es alles zu Gottes Ehre. Und an die Galater am 5. Kapitel spricht er: Durch die Liebe diene einer dem andern. Davon höret auch an das 33. Kapitel.

\section{Versöhnung Jacobs mit Esau}
\quellseite{95}

a Mählig) merkt daß Rechtgläubige und Werckheilige nicht können mit einander wandeln. Denn die Gläubigen fahren säuberlich mit stillem Geist, aber die Werckheiligen fahren stark mit Vermessenheit ihrer Werke in Gottes Gesetzen.

b Oder Schafe.

\section{Summarium}

Das ist ein fein Exempel, wie freundlich sich Jakob stellet, und demüthiget für seinem ältern Bruder Esau, V. 3. Und zur Lehre, daß wir nicht allein gegen leiblichen Brüdern, sondern auch ein jeder Christ mit dem andern aufs freundlichste leben sollen.

Daß aber Jakob einen Acker kauft im Lande Canaan, und da eine Hütte bauet und dem HErrn einen Altar zurichtet, V. 19., ist erstlich eine Anzeigung, wie er an der Hoffnung fest gehalten; und obgleich die Verheißung mit dem Lande Canaan sich lang verzogen hat, dennoch nicht hat gezweifelt, es werde seinem Samen zum Erbe werden. Darum, da er mehr nicht kann, kaufet er einen sondern Ort zum Gottesdienst, zum Zeichen, daß noch Gott sein Volk darein bringen und seinen Dienst darin anrichten werde. Darnach ists ein Exempel, wie die alten Väter Gottes Wort treulich geprediget und rechten Gottesdienst erhalten, und ohne Zweifel von den Heiden darum auch viel haben leiden müssen. Denn der Altar hat zu den zweien Stücken müssen dienen: daß erstlich daselbst Gottes Wort verkündiget und die Leute vom künftigen gebenedeyten Samen sind unterwiesen worden, und darnach die äußerlichen Gottesdienste mit dem Opfer am selben Ort sind ausgerichtet worden.

\section{Votum}

Verleihe der gütige und barmherzige GOtt, daß heute christliche Eltern auch sein Wort fördern helfen, und daß sie mit ihren Kindern allezeit in aller Gottseligkeit leben, und, wenn es GOtt gefällt, auch selig sterben, Amen.

\chapter{Das 34. Kapitel}

\section{Vorrede}

Also gehet es den Kindern Gottes auf Erden: wenn ein Kreuz bei ihnen aufhöret, so ist es ein Anfang eines andern und größern, und treibet immer eins das andere. Und widerfähret ihnen gleich, wie die Heiden schreiben von ihrem Hercule, welcher, wenn er dem gräulichen und schrecklichen Thiere, Hydra genannt, kaum ein Haupt hatte abgehauen, so wuchsen flugs drei andere Häupter an derselbigen Stelle: also ists ein unruhig und mühselig Ding um eines gottesfürchtigen Menschen Leben. Denn er ein stetiger Kreuzträger sein muß, und wenn er kaum eine Anfechtung überwunden hat, so ist bald da eine andere. Dessen ist ein sonderlich Exempel der heilige Patriarch Jakob, von welches großem Kreuze einem meldet das 34. Kapitel.
\quellseite{96}

\section{Die Schwächung der Dina verursachet ein Blutbad}
\quellseite{97}

\section{Summarium}

Hier ist eine Historia, wie der fromme Jakob an seinen Kindern viel Jammer und Herzeleid erlebt hat: daß ihm seine Tochter geschändet, V. 2., und die zween Söhne einen jämmerlichen Mord haben angerichtet, V. 25. Aber dies ist der Kreuze eines, das bisweilen Gott auch auf fromme Leute pflegt zu legen, da es doch weder an Zucht, noch guten Exempeln mangelt, wie man in der Historie Davids wohl siehet.

Neben dem siehet man, wie Gott der Hurerei so feind ist, weil er eine ganze Stadt um solcher Sünden willen erwürgen läßt. Derohalben denn weltliche Obrigkeit mit allem Ernst solchem Laster wehren und in dem Fall Gottes Zorn und Strafe, der es nicht will ungestraft lassen, zuvor soll kommen. Aber leider, es gehet solche Sünde je länger je gräulicher, und schier ohne alle Scheu und Schande; darum muß es allenthalben desto übler und schwerlicher zugehen.

\section{Votum}

GOtt steure und wehre allen Lastern; und regiere die liebe Jugend, daß sie zu seinem Ehren, und ihren Eltern zur Freude und Trost, allzeit christlich und züchtig leben, Amen.

\chapter{Das 35. Kapitel}

\section{Vorrede}

Es ist heuer nicht neu, daß einem Christen oft alle Wochen, ja alle Tage ein neu Kreuz zu Hause und Hof kömmt. Denn wie man im Sprüchwort sagt: Je lieber Kind, je größer Ruthe; also: Je frommer Christ, je mehr, größer, länger und schwerer Kreuz er tragen muß. Davon redet der heilige Paulus also in Geschichten der Apostel im 14. Kapitel: Durch viel Trübsal muß man in das Reich Gottes gehen. Und an die Ebräer am 12. Kapitel wird gesagt: Welchen der HErr lieb hat, den züchtiget er; er stäupet aber einen jeglichen Sohn, den er aufnimmt. Wohl aber denen, die ihr Kreuz in Geduld tragen. Denn wer sein Kreuz trägt in gläubiger Geduld, der lebt und bleibt in Gottes väterlicher Huld. Davon höret auch an das 35. Kapitel.

\section{Rahel stirbt an der Geburt}

a Ohren-Spangen) Lunulas heißt man güldene Haar-Bänder. Nun sind's Perlen-Borten geworden, Sprüchw. 11, 23.
\quellseite{98}

b Trankopfer) das war Wein; wie das in den folgenden Büchern genugsam gesehen wird.

c Benoni) heißt meines Schmerzens Sohn.

d Benjamin) heißt der Rechten Sohn.

\section{Summarium}

Ist eine Historia, wie Jakob ein recht bischöflich Amt führet, daß er die Abgötterei unter den Seinen ausrottet, und nach Gottes Befehl den rechten Gottesdienst anrichtet, B. 2. Zum Exempel aller Obrigkeit, daß sie wisse, es sei ihr fürnehmstes Amt und höchstes Werk, über dem rechten Gottesdienst halten, und Abgötterei mit Gewalt ausrotten.

Zum andern siehet man, wie der fromme Vater Jakob wiederum zwei neue und schwere Kreuze erfahren muß: daß ihm seine liebe Rahel in Kindesnöthen stirbt, B. 18., und sein ältester Sohn Ruben sein Kebsweib beschläft, B. 22. Denn Gott will seine Heiligen auf mancherlei Weise hier auf Erden üben und versuchen. Solche Exempel mögen wir uns zum Trost merken, wenn dergleichen uns begegnet, daß wir gedenken, wir sind nicht allein, es sei dergleichen auch den großen trefflichen Heiligen widerfahren.

\section{Votum}

Der treue barmherzige GOtt lasse ihm alle christliche Kreuzträger zu Trost und Hülfe befohlen sein, und behüte uns alle für allem Unfall Leibes und der Seele, Amen.

\chapter{Das 36. Kapitel}

\section{Vorrede}

Dies Kapitel meldet nur, wie Esau gewachsen und zum großen Volke worden, und auf Erden mit Gewalt und Herrlichkeit sehr habe zugenommen. Ist derohalben unnöthig, daß man die bloßen Namen in der Kirche lese; es geht aber dem Sprüchwort nach: Unkraut wächst mit Haufen, und verdirbt nicht; da es dagegen mit dem frommen Jakob, der doch den Segen und die Verheißung hatte, sehr geringlich sich anließe und schwächlich fort ginge; Ursache der Satan, der ist den Frommen Feind und hindert, wo er kann.

\section{Geschlechtsregister Esaus}
% Scan 99: no retained apparatus prose; Bible text and other excluded blocks are inventoried in selection.json.
\quellseite{100}
hieß Mehetabel, eine Tochter Matreds, die Mesahabs Tochter war.

\chapter{Das 37. Kapitel}

\section{Vorrede}

Das 36. Kapitel des ersten Buches Mose lehret, wie Esau gewachsen und zum großen Volke worden, und auf Erden mit Gewalt und Herrlichkeit habe zugenommen (denn Unkraut wächst mit Haufen und verdirbet nicht); da es dagegen mit dem frommen Jacob, der doch den Segen und die Verheißung hatte, sehr geringlich sich anließe und schwächlich fortginge. Denn der Satan ist den Christen feind und hindert, wo er kann. Das lässet aber Gott darum geschehen, auf daß der Glaube geübet und geprüft werde, daß wir lernen allein am Wort hangen, und was sichtbar und greiflich ist, nicht ansehen, sondern uns daran halten, daß man nicht siehet, und daß wir den Trost behalten, davon St. Paulus sagt 2 Cor. 4.: Wir leiden Verfolgung, aber wir werden nicht verlassen. Wir werden unterdrückt, aber wir kommen nicht um. Darauf höret ferner an das 37. Capitel.

\section{Joseph wird aus Neid von seinen Brüdern verkauft}
\quellseite{101}

a Vater) das war Jaakob, Jakobs Vater und Josephs Großvater, welcher nach dieser Geschichte noch zwölf Jahre gelebet hat.

\section{Summarium}

Hie hebt sich an die Historia vom Joseph, in welcher wir sehen, wie dem frommen Jakob seine Söhne ein neues Herzeleid zurichten, daß sie Joseph verkaufen und sprechen, es habe ihn ein wildes Thier zerrissen, V. 20.

Dies soll uns ein sonderlich Exempel sein, an welchem wir lernen und sehen sollen, wie Gottes Rath durch Menschen Anschläge nicht möge gehindert werden. Joseph wird darum von seinen Brüdern verkauft, auf daß sein Traum nicht erfüllet werde, und sie ihn nicht müssen als einen Herrn anbeten. Aber eben mit solchem Rath fördern sie Joseph zu solcher Herrlichkeit, zu welcher er nimmermehr hätte mögen kommen, wo er bei dem Vater wäre blieben. Derohalben sollen wir nicht so bald verzagen, wenn das Kreuz kommt und es übel zugehet. Man harre nur; es kann Gott treffliche und große Dinge dadurch ausrichten, so man in der Furcht Gottes bleibet (und durch die Sünde nicht Ursache zur Strafe gibt), wie Joseph thut.

\section[Notum]{\unsicher{Notum}}

Der treue Gott verleihe solches heute über allen christlichen Kreuzträgern; und segne und behüte unsere liebe einheimische und ausländische Jugend, an Leib und Seele, um JEsu Christi, seines geliebten Sohnes willen, Amen.

\chapter{Das 38. Kapitel}

\section{Vorrede}

Der heilige Apostel Paulus warnet und spricht 1 Cor. 10.: Wer sich lässet dünken, er stehe, mag wohl zusehen, daß er nicht falle. Hat diese Warnung jemals vonnöthen gethan, so thut sie ja zu dieser gegenwärtigen Zeit vonnöthen, da die Aergernisse so gar groß und die Teufel so er
\quellseite{102}
schrecklich grimmig sind. Und: haben große hochbegabte Leute gestolpert, und schwere Sündfälle gethan, so mögen sich wohl vorsehen andere, die weniger und geringere Gaben haben, daß sie nicht auch in grobe und schwere Sünde gerathen. Derowegen, ach laßt uns wachen und beten, daß wir nicht in Anfechtung und Sünde wider das Gewissen fallen. Darauf höret weiter an das 38. Kapitel.

\section{Juda Blutsschande mit der Thamar}

a Hirten Mag auch heißen Freund, darnach die Puncta im Hebräischen sich setzen lassen. Denn Juda hat ja müssen Weide haben, vielleicht auch eines Freundes nicht gerathen mögen. Puncta können so wohl fehlen als treffen.

b Hirten oder Freund, welches du willst.
\quellseite{103}

\section{60	Das 1. Buch Mose 38. 39}

\section{Summarium}

Es ist eine Historia, welche am allermeisten von Mose geschrieben ist um des HErrn Christi willen, welches Geschlecht von Juda und Thamar herkommt. Nun war Thamar eine Heidin, und ist doch unsers HErrn Christi Großmutter, anzuzeigen, daß auch die Heiden Gottes Volk sollen werden. Ja Gott ist nicht allein der Juden, sondern auch der Heiden Gott. Nun werden aber solche Historien von den heiligen Vätern nicht darum geschrieben, daß man ihren Sünden folgen soll, sondern daß wir an ihnen erkennen sollen unsere verderbte Natur, und derhalben nicht sicher sein, sondern beten, und in Gottesfurcht wandeln.

Darnach, daß man lerne wider das böse Gewissen sich mit Gottes Barmherzigkeit trösten; sintemal Gott solche Sünder nicht hat weggeworfen, sondern in seiner Linie behalten. Aber hie gehört auch zu, daß man sich bessere und von Sünden ablasse, wie Juda thut. Denn seine zween Söhne Er und Onan, da sie sich nicht wollen bessern, tödtet sie der HErr ihrer Sünden halber, V. 7. 10.

\section{Votum}

Der barmherzige GOtt regiere uns, damit wir nicht Aergerniß anrichten, und helfe, daß wir uns seiner väterlichen Barmherzigkeit in allen Nöthen freuen und trösten, und alle selig werden, Amen.

\chapter{Das 39. Capitel}

\section{Vorrede}

Wie zweierlei Christen sind, was den Glauben anlangt, nämlich, starke und schwache: also sind auch zweierlei Christen, was das Leben und die Werke betrifft. Etliche sind stark, etliche sind schwach. Starke sind, welche sich im Leben also verhalten, daß sie mit äußerlichen Sünden und Gebrechen, die man augenscheinlich sehen kann, sich nicht beflecken. Schwache sind, die bisweilen, nicht aus bösem Vorsatz, sondern aus angeborner menschlicher Blödigkeit, und von ihrem eigenen Fleisch und Blut übereilet, fallen in Gebrechen, die sie denn auch alsobald wiederum erkennen, herzlich Reue und Leid darüber haben, Gott im Namen Jesu Christi um Verzeihung bitten, auch, mit Hülfe und Beistand des heiligen Geistes, forthin ihren Wandel vorsichtiger anzustellen, damit sie niemand muthwilliger Weise ärgerlich sein, sich befleißigen. So wolle Gott die Starken in ihrer Frömmigkeit bestätigen und erhalten, und die Schwachen stärken, auf daß sie erfüllet werden mit Früchten der Gerechtigkeit, die durch Jesum Christum geschehen, zur Ehre und Lobe Gottes, Amen. Darauf höret nun an das 39. Kapitel.

\section{Josephs Dienst, Keuschheit und Gefängniß}
\quellseite{104}

a In dem du) Sonst mußte er auch die Frau versorgen, eben so wohl als das ganze Haus, ohne daß er nicht bei ihr schlafen sollt.

\section{Summarium}

Es ist eine sehr schöne Historia, in welcher wir am ersten lernen, wie Gott fromme Leute segnen und Glück und Heil um ihretwillen auch andern geben will, V. 3. 5. Darnach wie Joseph sich keusch und rein hält, da ihm doch Ursach, Uebels zu thun, auf das nächste bracht ward, V. 8. Solch Exempel, daß er Gott für Augen hat und für Sünde sich hüte, ist der Anfang alles Segens und Gnade, so Gott dem Joseph hat widerfahren lassen. Denn wie im ersten Gebot stehet, will Gott die, so ihn lieben, das ist, so sein Wort für Augen haben, in tausend Glied segnen. Aber man lerne hie auch die Weise, so Gott mit seinen Heiligen hält, welche er will segnen, und etwas Sonderes aus ihnen machen. Denn er steckt sie zuvor in Leiden und Unglück bis über die Ohren, und ob sie gleich unschuldig sind, stellet er sich, als kenne er sie nicht, und frage nicht darnach, wie man hie an Joseph siehet, der fromm und unschuldig ist, und dennoch muß er ins Gefängniß, daß böse unzüchtige Weib bleibet in Ehren und Würden. Aber man halte nur fest über Gottes Wort und Willen und leide geduldig. Denn gewiß ist es, ob man schon über einer rechten Sache gestöcket und geplöcket wird, so soll doch die Wahrheit endlich an Tag kommen, und Gott will endlich dem gerechten Theil helfen, wo man anders seine Hülfe durch Geduld erharren will. Darum gibt Gott auch Joseph im Gefängniß Glück und Heil.

\section{Votum}

Gott der Vater aller Barmherzigkeit lasse ihm uns alle befohlen sein; und helfe, daß die Jugend mit uns, in Lieb und Leid christlich lebe, und zeitlichen und ewigen Segen, auch mit uns habe und behalte, Amen.

\chapter{Das 40. Kapitel}

\section{Vorrede}

Der König und Prophet David sagt im 34. Psalm: Der Gerechte, das ist, der gläubige Mensch, muß viel leiden. Und St. Paulus spricht in den Geschichten der Apostel im 14. Kapitel: Durch viel Trübsal müssen wir ins Reich Gottes eingehen. Wenn man nun glauben könnte, daß Gott den Menschen, die er am liebsten hat, das größte, schwerste und langwierigste Kreuz zuschickt, so würde viel Klage und Ungeduld unter dem Kreuz dahinten bleiben. Derowegen wer heute ein Kreuzträger ist, der gedenke, wie jener erfahrene, versuchte und christliche Kriegsmann recht und wohl gesagt
\quellseite{105}
hat: Es muß niemand mehr Gefahr und Unglück erwarten und aufsehen, als der, welcher Gott und dem Kaiser treulich dienet. Er freue sich auch, daß Gott unter dem Kreuz trösten und Geduld verleihen wolle, und daß auf das Kreuz endlich Hülfe und ewige Freude folgen werde. Dies erinnert auch das 40. Kapitel.

\section{Joseph leget den Gefangenen ihre Träume aus}

\section{Summarium}

Sie ist ein fein Exempel, welches dazu dienet, daß man sich in Gottes wunderbare Werke recht lerne schicken, V. 3. Denn Joseph muß in den Kerker und ist doch unschuldig. Solches wird ihm ohne Zweifel haben wehe gethan; aber weil er solch Unrecht geduldig leidet, und sein Vertrauen auf Gott setzet, er werde letztlich helfen und retten, muß solcher Kerker und Gefängniß ihm dazu dienen, daß man Gottes Gabe an ihm lerne erkennen, um welcher willen er hernach zu großen Ehren kömmt. Solcher Gabe braucht er also, daß er doch Gott allein die Ehre gibt.

Aber er muß hie zum andernmal die Untreue der Welt erkennen lernen, daß der Schenk sein vergißt und er noch zwei Jahr muß im Gefängniß bleiben, V. 23. Auf daß wir lernen fest auf Gottes Hülfe trauen, ob er sie gleich lange verzeucht.
\quellseite{106}

\section{Votum}

Der GOtt alles Trostes tröste alle die das Kreuz tragen, und behüte uns für Gefängniß, und für zeitlicher und ewiger Noth, um JEsu Christi, seines lieben Sohns und unsers HErrn willen, Amen.

\chapter{Das 41. Kapitel}

\section{Vorrede}

Wer droben im Himmel mit den Heiligen und Auserwählten in Freuden leben will, derselbe Mensch soll allhie auf Erden, nach dem Exempel der Heiligen und Auserwählten, in Gottesfurcht leben. Denn, die dem Teufel zu Willen leben in dieser Welt zeitlich, die werden mit ihrem Unwillen, beim Teufel auch in der Hölle, leben und leiden ewiglich. Derowegen fürchtet den HErrn und dienet ihm treulich und rechtschaffen. Und so euch der Satanas euer Widersacher mit Versuchung und Anfechtung zusetzt, so thut ja nicht seinen Willen, sondern widerstehet ihm feste im Glauben, und leget nimmer ab den Harnisch Gottes, das ist, lebet in täglicher Buße, daß ihr bestehen könnt gegen die listigen Anläufe des Teufels und das Feld behalten möget. Dies betrachten wir heute auch bei dem 41. Kapitel.

Joseph wird durch Auslegung der Träume Pharaos zu fürstlicher Hoheit erhoben.

a Bei mir nicht) will sagen: ich bins nicht, der die Träume gedeutet hat oder könne; Gott ists, der es durch mich gethan hat, kann dirs auch thun.
\quellseite{107}

\section{Summarium}

Es ist eine Historia, in welcher wir an Joseph lernen: ob Gott schon die Frommen eine Zeit lang läßt leiden, daß er sie doch endlich errettet, und ihnen das Unglück zum besten gedeihen muß. Denn wo Joseph nicht wäre im Kerker gewesen, wäre mans nicht von ihm innen worden, daß er solche Gabe hätte, die Träume zu deuten, wäre auch nicht zu solchen Ehren kommen. Darum sollen wir lernen, daß wir uns am Kreuz nicht ärgern, sondern auf Gottes Hülfe sehen, und derselben gewiß warten. Neben diesem ist auch eine nützliche Hauslehre hie, V. 47., daß die Hausväter, wenn Gott Segen zur Nahrung gibt, solchen Segen nicht verschwenden sollen, wie doch gemeiniglich geschieht. \unsicher{Steermann} genennt, weil man viel hat, man soll viel umbringen und aufheben lassen. Aber Gottes Segen sollen wir aufheben und auf künftige Noth für uns und auf andere sparen. Denn weil der meiste Theil Gottes Segen mißbrauchet und undankbar ist, folget gemeiniglich auf eine große Wohlfahrt große Theurung, da soll ein vernünftiger Haushalter, wie Joseph denn hie that, sich aufrüsten.

b) Abrech) was Aelters heiß, lassen wir die Zänker \unsicher{suchen} an den jüngsten Tag, wollen \unsicher{vielmehr} verstehen wie es verdeutscht ist.
\quellseite{108}

\section{Summarium}

Ist ersichtlich ein fein Exempel, daß wir, wie Joseph, unsere Gaben nicht sollen brauchen zu eigener Hoffahrt und Nutz, sondern zu des Nächsten Wohlfahrt. Denn darum hat Gott solchen Traum Pharao sehen und die Deutung Joseph sehen lassen, weil eine solche Hungersnoth fürhanden war, daß die Leute sich darein schicken und der guten Jahre desto besser zum Vorrath brauchen sollen.

Darnach ist ein fein Exempel einer feinen fürstlichen Tugend, daß Pharao den hochbegnadeten Mann Joseph theuer achtet, ihm einen Vater des Landes nennt und den nächsten nach dem Könige ordnet, B. 43. Denn in Regimenten soll man sonderlich darauf Achtung geben, daß tapfere, geschickte Leute in Ehren gehalten werden. Wo aber leichtfertiges, liederliches Gesinde bei den großen Herrn lieb und werth ist, da hauset man selten wohl.

\section{Votum}

Verleihe GOtt, daß Obrigkeit und Unterthanen nach seinem Willen leben, und daß ein jedes endlich von allem Uebel erlöset, mit ihm in ewiger Sicherheit und Freude lebe, Amen.

\chapter{Das 42. Kapitel}

\section{Vorrede}

- Auf Gottesfurcht folget ein fröhlich Gewissen und Wohlfahrt; auf die Sünde folget dagegen ein böses Gewissen und Unglück oder Strafe. Denn von der Furcht Gottes saget also Sirach im 1. Capitel: Die Furcht des HErrn macht das Herz fröhlich und gibt Freude und Wonne ewiglich. Wer den HErrn fürchtet, der wird endlich den Segen behalten. Von der Sünde aber spricht David im 32. Psalm: Der Gottlose hat viel Plage; und Esaias im 57. Kapitel: Die Gottlosen haben nicht Friede, spricht mein Gott; derowegen laßt uns in der Furcht des HErrn leben und die Sünde fliehen, daß wir nicht mit der äußerlichen Züchtigung Seelen-Angst und des Gewissens Plagen und Nag[e]n ewig müssen tragen. Ferner höret an das 42. Kapitel.

\section{Reise der Söhne Jakobs in Egypten, ohne Benjamin}
\quellseite{109}

\section{Summarium}

Es ist eine feine Historia, wie Gott seine Räthe heimlich und wunderlich führet, bis er es dahin bringet, darauf sein Wort verrichtet hat. Denn hie sehet man, wie wunderbarer es sich spielet, ehe denn der Traum erfüllet wird, welchen Joseph seinen Brüdern gesagt, und sie ihn gern gehindert und ihn darum verkauft hatten. Aber eben damit sie es gedenken zu hindern, fördern sie es. Denn Gottes Gedanken müssen für sich gehen; der Menschen Gedanken aber müssen hinter sich gehen.
\quellseite{110}

\section{Summarium}

Es ist ein fein Exempel an den Brüdern Joseph, wie sie ihr Gewissen naget und sie jetzt an ihre Sünde gedenken, da sie es längst sollten vergessen haben. Aber die Sünde läßt sich nicht vergessen, sondern muß das Herz unruhig machen, es geschehe über lang oder über kurz. Darum mag man sich dafür hüten. Es war bisher unter den Brüdern Joseph alles still gewesen; aber jetzt mußte Simeon gefangen liegen und mit den andern das Gewissen und Sorge tragen, Gott wollte das Blut Josephs von ihnen fordern. Das Geld in Säcken, V. 35, macht ihnen auch Beisorge; ihre Untugend mußte gestraft werden. Endlich beklagt sich auch Jakob, ihr Vater, daß sie ihn um seine Söhne bringen, V. 36. Also wird ihnen mancherlei Weise ihre Sünde ins Gewissen eingetrieben, alles uns zur Warnung, daß wir’s bedenken und nicht sündigen, noch in Sünden sicher sein.

\section{Votum}

Gott der heilige Geist regiere uns, daß wir Gottes Wort folgen, und für Sünden und hüten, auf daß wir nicht Sündenstrafen leiden dürfen, sondern in Gottes Gnade leben und sterben, Amen.

\chapter{Das 43. Kapitel}

\section{Vorrede}

Gottes Wort wird täglich gelehret; aber viel hören es also, daß sie sich dennoch mit Sünden in ihrem Beruf, Handel, Erwerb und Haushaltung nähren, und daß sie unter dem Kreuz wenig zu ihrem Trost an Gottes Verheißung gedenken. Derowegen sehet zu, wie ihr Gottes Wort höret, und seid nicht allein Hörer, sondern auch Thäter desselben, und tröstet euch auch damit. Denn bei denen, die den Bund halten und darnach thun, bleibet Gottes Gnade ewiglich, sagt David im 103. Psalm. Und St. Paulus Röm. 2. spricht: Für Gott sind nicht gerecht, die das Gesetz hören, sondern die das Gesetz thun. Und Luc. 12. sagt auch Christus: Der Knecht, der seines Herrn Willen weiß und hat nicht nach seinem Willen gethan, der wird viel Streiche leiden müssen. Und Ebr. 4. wird gesagt: Das Wort der Predigt half jene nichts, da nicht glaubten die, so es höreten. Dazu höret an das 43. Kapitel.

\section{Reise der Söhne Jacobs in Egypten mit Benjamin}
\quellseite{111}

\section{Summarium}

Es ist zum ersten ein Exempel, wie auch heilige Leute, die die gemeinen Strafen der Welt, als Hunger, Krieg und andere dergleichen, müssen mit tragen, und es GOtt dennoch also mit ihnen schaffet, daß sie sonderlich Hülfe und Schutz finden, B. 1.

Darnach ist’s auch ein fein Exempel, wie Jakob treulich und ehrbarlich handelt, daß er nicht auf seinem Nutz, sondern auf die Billigkeit siehet, und das Geld so in Säcken gefunden ward, heißt wiederbringen, B. 12. So man solches nach der jetzigen Welt Art wollte rechnen, würde man wenig Leute finden, die es ihm nachthäten; sogar ist jetzt alles im Geiz und Eigennutz ersoffen.
\quellseite{112}

\section{Summarium}

Dies ist ein Exempel, wie Gottes Zusagungen müssen erfüllet werden und unverhindert bleiben, es thue die Welt dazu, was sie wolle. Denn hie wird Josephs Traum recht erfüllet, von den Garben und Sternen, weil seine Brüder für ihm sich neigen und niederfallen, und hilft sie nichts, was sie, solchem vorzukommen, fürgenommen hatten.

Darnach ist auch wohl zu merken die Freundlichkeit, so Joseph gegen seine Brüder brauchet, B. 27. Denn auch in heiligen Leuten die Natur ihre angeborne Eigenschaft behält, daß Brüder sollen gegen einander freundlich sein. Darum ist es eine unnatürliche, teuflische Untugend, wo Brüder untereinander uneins oder unfreundlich sind.

\section[Notum]{\unsicher{Notum}}

Helfe der treue GOtt, daß wir alle einander als Brüder und Schwestern herzlich lieben, damit er aus Gnaden uns wieder liebe, und die Werke seiner Liebe erzeige, um JEsu Christi, seines geliebten Sohns, unsers HErrn willen, Amen.

\chapter{Das 44. Kapitel}

\section{Vorrede}

Gott züchtiget und strafet uns, nicht, daß er uns verderbe, sondern, daß er uns dadurch zur Buße berufe, auf daß wir also gebracht werden zum Erkenntniß der begangenen Sünden, und daß dadurch den künftigen Sünden gewehret, und daß wir daneben zur Gottseligkeit gebracht und geführt werden. Davon sagt das Buch der Weisheit im 2. Kapitel also: Da strafest säuberlich, so da fallen, und erinnerst sie mit Zucht, wenn sie sündigen, auf daß sie von der Bosheit los werden und an dich, HErr, glauben. So pflegt man auch sonst zu sagen: Das Kreuz entweder den Menschen probirt, oder demütiget, oder reiniget, oder Gottes Ehre verkläret. Derowegen ein jeder sein Creuz in Geduld tragen und zur Buße und Uebung des Glaubens gebrauchen soll. Davon lehret auch das 44. Kapitel.

\section{Josephs Brüder werden hart geängstiget}
\quellseite{113}

\section{Summarium}

Sie ist eine sehr feine Historia, wie Joseph seinen Brüdern, ehe er sich zu erkennen gibt, einen neuen Schrecken einjaget mit dem gestohlenen Becher, B. 2. Er thut aber solches der Meinung, auf daß er ihnen die Sünde wohl eintreibe und das Gewissen weidlich rühre, daß sie aus Neid ihn erstlich hätten tödten wollen, und darnach unter die Fremden verkauft, und ihren frommen Vater so schmerzlich betrübet hätten. Solche Sünde hatte ihr Gewissen nicht dermaßen daheim gesucht, wie jetzt geschieht.

\section{Votum}

Der treue Gott erleuchte uns durch den heiligen Geist, daß wir unsere Sünde erkennen, die Gott abbitten, fliehen und meiden, damit wir in wahrer Buße christlich leben und selig sterben, Amen.
\quellseite{114}

\chapter{Das 45. Kapitel}

\section{Vorrede}

Das fünf und vierzigste Kapitel des ersten Buchs Mose ist ein recht tröstlich Kapitel und möchte wohl ein Brunnquell reichen Trostes genannt werden. So helfe der treue, liebe Gott, daß dadurch erquicket werden aller derer Herzen, die auch ihrer Sünden halben Leid tragen, und derer, die auch Kreuzträger und Kreuzträgerin sind. Denn wahrhaftig, Gott tröstet unter dem Kreuz, und hilft endlich von allem Kreuz. Und Gott versucht also seine Gläubigen und schicket ihnen ja zu Anfechtung und allerlei Noth und Widrigkeit, aber doch also, daß er sie wiederum erquicket und aufrichtet mit großem, starkem Trost, und daß er sie erfreuet mit nötiger und großer Hülfe. Davon höret auch an das 45. Kapitel.

\section{Joseph gibt sich seinen Brüdern zu erkennen}

a Lasset euch euren Hausrath nicht hindern, was ihr nicht verkaufen könnt in solcher theuren Zeit, das laßt hinter euch.
\quellseite{115}

\section{72 Das 1. Buch Mose 45. 46}

† Das hebräische Wort heißt eigentlich anders thun und anders werden, Thren. 2, 18. Fleto et non desinat pupilla oculi tui, neque quiescat pupilla oculi tui. Weine, und laß deine Augen nicht anders thun, Kap. 3, 49. Ich kann nicht anders thun. Lex Tapug. Hab. 1, 4. Es gehet anders denn recht, Recht gehet anders, gilt nichts. Also auch Jakob gedenkt viel anders, als sie es ihm erzählen. Sic Jacob longe aliud sentit, quam illi narrant.

\section{Summarium}

Es ist eine sehr feine Historia, wie Joseph sich seinen Brüdern zu erkennen gibt, B. 3, und thut wie ein rechtschaffener Prediger thun soll. Nachdem seine Brüder bekümmert und ihrer Sünden halben erschrocken waren, daß er sie wieder tröstet und ihre Sünden, so sie an ihm begangen hatten, zum Besten deutet. Er schilt sie nicht, gedenkt auch ihres Neides und Hasses nicht. Denn solches alles hatte zuvor ihr eigen Gewissen schon gethan, darum war es Zeit Tröstens. Sondern zeiget ihnen aufs allerfreundlichste an, wie Gott durch ihr Verkaufen ihn dahin gefordert habe, daß er sie und ganz Egypten von der Hungersnoth erretten sollte, B. 6. Aber solches hat Joseph selbst hernach müssen lernen; da er verkauft und hernach lange Zeit gefangen war, wußte er’s noch nicht. Auf daß wir auch lernen des Guten hoffen, und nicht kleinmüthig seien, wenn es uns übel gehet. Denn wir wissen nicht, was Gott damit meinet, und daß es uns zum Besten gedeyen soll, wie Paulus sagt Röm. 8, 28.: Den Gottseligen soll es alles zum Besten dienen.

\section{Votum}

Gelobet und gepreiset sei der gütige, treue GOtt, daß er betrübt und tröstet, daß er Kreuz aufleget, auch Kreuz tragen hilft und wegnimmt; der wolle uns auch in aller Noth allzeit trösten und helfen, Amen.

\chapter{Das 46. Kapitel}

\section{Vorrede}

Von Gottes Verheißungen fällt nicht ein einzig Wort auf die Erde, sondern sie werden alle endlich und vollkommenlich erfüllet. Denn also sagt St. Paulus 2 Cor. 1.: Alle Gottes-Verheißungen sind ja in Christo, und sind Amen in Christo, Gott zu Lobe. Derwegen laßt uns nach Gottes Wort und Geboten gehorsamlich leben und wandeln, wie es eines jeden Beruf und Stand erheischt; so wird sich Gott nach seinen Zusagungen und Verheißungen gegen uns väterlich verhalten und erzeigen, wie es unsre Nothdurft und Wohlfahrt in diesem und künftigem Leben erfordert. Denn die Gottseligkeit ist zu allen Dingen gut und hat die Verheißung dieses und des zukünftigen Lebens, sagt St. Paulus 1 Tim. 4. Darauf höret auch an das 46. Kapitel.

\section{Jakobs Reise nach Egypten zu seinem Sohne Joseph}
\quellseite{116}

\section{Summarium}

Hier ist eine Historia, wie Jakob sei in Egypten gezogen, wie Gott oben im 15. Capitel dem Abraham zuvor saget, sein Same soll vierhundert Jahre in Egypten wohnen, und denn allererst die Verheißung vom Lande Canaan erlangen. Nun ist’s tröstlich zu hören, daß nur siebenzig Menschen in Egypten ziehen, v. 27, und in vierhundert Jahren zu einem solchen großen Volke werden, daß sechsmal hundert tausend Mann, ohne Weib und Kinder, aus Egypten in das Land Canaan ziehen. Es ist auch ein fein Exempel für die Kinder, daß Joseph, ob er wohl ein Fürst war, seinem alten Vater nachreiset und sich auf das freundlichste gegen ihn stellet, v. 29.

\section{Votum}

Der treue Gott erfülle auch an unserm Leib und Seele seine Verheißungen, zu zeitlicher und ewiger Wohlfahrt; und helfe, daß alle Kinder ihre Eltern ehren, wie er es in seinem Wort befohlen hat, Amen.
\quellseite{117}

\chapter{Das 47. Kapitel}

\section{Vorrede}

Himmel und Erde, auch aller Kaiser, Könige und Fürsten Schlösser vermögen nicht, daß ihm Gott daraus eine Wohnung mache; aber in dem Menschen, der sein Wort hält, will Gott wohnen, sammt seinem Sohne und dem heiligen Geiste. O, ein herrlicher Trost! Die Gottes Wort glauben und darnach christlich leben, sollen der ganzen heiligen Dreifaltigkeit Tempel und Ruhhaus sein; da doch wir Menschen von Natur arme Sündenwürmlein, Kinder des Zorns, leibeigene Knechte Satanä sein und aller Teufel Fußschemel ewig sein sollten. Wer wollte nun nicht Gottes Wort gerne hören, fleißig lernen, mit Freuden glauben und christlich gebrauchen? Denn sein wir Behausungen Gottes, so wird auch Gott seine Wohnung wider allen Sturm und Anlauf schützen, wird uns in seiner Noth verlassen und wird uns im ewigen Leben wieder beherbergen, wie wir ihn beherberget haben in unsern Herzen. Und darauf höret an das 47. Kapitel.

\section[Jakob wohnet im Lande Gosen. Die Theurung ist]{Jakob wohnet im Lande Gosen. Die Theurung ist groß in Egypten}

a Wohnen) zur Herberge, Gast sein, Fremdling sein. Non cives aut domestici hujus mundi.

b Kinder) Quia nos senes mali propter pueros omnibus bonis fruimur. Wir alten Narren essen mit den Kindern, nicht sie mit uns. Ipsi Domini, nos procuratores.

\section{Summarium}

Das ist ein fein Exempel einer schönen fürstlichen Tugend am Könige Pharao, daß er um Josephs willen seinem ganzen Geschlecht hilft, V. 6. Damit beweiset Gott wiederum, obschon die gemeinen Weltstraßen über Böse und Gute gehen, daß dennoch die Frommen einen Vortheil haben sollen.

Darnach ist wohl zu merken, wie der alte Jakob von seinem und unser aller Leben auf Erden für dem Könige Pharao redet, V. 9. Daß es nur eine kurze Wallfahrt sei, da wir keine bleibende Stätte haben, und dennoch allerlei Unglück versuchen und tragen müssen. Darum sollen wir immerdar ein Sehnen und Verlangen haben nach dem rechten Vaterlande, daß wir desselben in Geduld und gutem Gewissen erwarten.
\quellseite{118}
* Neiget sich (eig) er lag im Bette krumm, richtet sich hoch auf, neiget sich zu Häupten, betet und danket Gott; dieweil that Joseph den Eid.

\section{Summarium}

Es ist eine Historia, in welcher uns eine fürstliche, herrliche Tugend wird angezeigt: obwohl ganz Egypten der Theurung halben an Geld und Vieh entblößet, die Aecker verkauft und der Fünfte von dem Gewächs dem Könige bringet wird; daß dennoch die Priester nicht beschweret, sondern denselbigen von Hofe ihr Aufenthalt gereicht wird, V. 26. Denn dies soll das fürnehmste Werck der frommen Obrigkeit sein, daß sie ob denen halten, so zu Kirchenämtern gebraucht werden. Was aber Jakob begehrt, man soll ihn begraben im Lande Canaan, V. 29, geschieht zum Anzeigen, daß er fest an der Verheißung Gottes gehalten: daß Land Canaan soll seines Samens eigen sein, ob es sich gleich lang verziehen würde.

\section{Votum}

Gott der Vater aller Barmherzigkeit verleihe, daß wir das kurze mühselige Leben recht gebrauchen und beschließen, damit wir hernach in Ewigkeit mit ihm leben in himmlischer Freude und Wonne, um JEsu Christi Verdienst willen, Amen.

\chapter{Das 48. Kapitel}

\section{Vorrede}

Die Kinder sollen ihre Eltern im Herzen aus Demuth hoch und groß, als Gottes Statthalter achten und halten, weil ihnen Gott die Ehre gibt und sie gleich neben sich an die Stelle setzt. Sie sollen ihnen von Gottes wegen schuldigen Gehorsam leisten, und sollen thun, was ihnen lieb, und lassen, was ihnen leid ist, sofern es nicht wider Gott ist; sie sollen sich mit Geduld von ihnen strafen und züchtigen lassen; sie sollen sie im Fall der Noth unterhalten und versorgen; und da sie Alters oder Krankheit halben wunderlich sind, sollen sie ihre Fehle und Mängel zu gut halten, dieselben mit Sems und Japhets Mantel bedecken. In Summa: die Kinder sollen sich nach ihren Eltern richten, und sie allwege für klüger, weiser, verständiger, denn sich halten, und sollen alle
\quellseite{119}
ihre Befehle dahin deuten, daß sie es gut mit ihnen meinen. Also soll man eben Vater und Mutter mit Worten und Geduld, auf daß ihr Segen über die Kinder komme, Sirach 3. Davon unterrichtet auch das 48. Kapitel.

\section{Jakobs Testament wegen Ephraim und Manasse}

a Stück heißt im Hebräischen Sichem, und dieselbe Stadt meinet er hie.

\section{Summarium}

Dies ist ein fein Exempel, wie Joseph seinem kranken Vater nachreiset und ihn besucht, B. 1. Und begehret seines Segens über seine Kinder, daß sie nach Jakobs, Isaaks und Abrahams Namen genennet werden, das ist, daß sie der Verheißung theilhaftig werden, welche von Gott
\quellseite{120}
den heiligen Vätern, Jakob, Isaak und Abraham, geschehen war. Daß aber Jakob den erstgeborenen Sohn an des jüngsten Statt stellt, V. 14., geschieht wie oben gemeldet, auf daß man lerne, Gott wolle mit uns nicht nach Verdienst, noch nach dem äußerlichen Ansehen, sondern nach lauter Gnade handeln.

\section{Votum}

Der treue liebe GOtt erleuchte und regiere alle Kinder, daß sie ihre Eltern ehren, damit sie GOtt wieder ehre; und segne uns alle, als unser lieber Vater, beide zu Leibes und Seelen Nothdurft und Wohlfahrt, Amen.

\chapter{Das 49. Kapitel}

\section{Vorrede}

Wie geringe Hoffnung sein kann, wenn der Same verdirbt und Schaden nimmt, daß eine gute Ernte darauf erfolgen könne: also kann auch nichts Gutes dem Vaterland und gemeinem Nutz daraus erwachsen, wenn die zarte Jugend, (welche die Alten des gemeinen Vaterlandes Hoffnung genannt) in der Blüthe verderbet und übel auferzogen wird. Derwegen hat Jakob der Patriarch seine Kinder nicht allein treulich die Zeit seines Lebens, sondern auch auf seinem Todbette unterrichtet, wie sie sich gottselig halten sollen. Dergleichen haben gethan David, und der alte Tobias an ihrem Todbette, 1 Buch der Könige 2., Tob. 4. Auch die Mutter der sieben Söhne, die Antiochus der Tyrann jämmerlich, um Gottes Worts willen, hat hinrichten lassen, 2 Maccab. 7. Wohl nun allen Eltern, die ihre Kinder aufziehen in der Zucht und Vermahnung zu dem Herrn. Dies erinnert auch das 49. Kapitel.

\section{Jakobs Weissagung und Abschied}

a Ruben sollte der Erstgeburt Würde haben, nämlich das Priesterthum und Königreich. Nun aber wird beides von ihm genommen, und Levi das Priesterthum, und Juda das Königreich gegeben. \unsicher{Dieß bedeutet die Synagoga, die das Bette Jakob, das ist, die Schrift befleckt mit falscher Lehre, darüber sie verloren hat Priesterthum und Königreich.}

b Die fängt an der Segen von Christo, der von Juda geboren sollte werden. Und heißt ihn Silo, das ist, der glückselig sein und frisch durchbringen sollte mit Geist und Glauben, daß zuvor durch Werke faul und unselig Ding war. Darum nennen wir Silo einen Held. Denn der vorige Theil dieses Segens betrifft den König David; und ist sonst in allen Segen nichts mehr von Christo, sondern alles andere ist von zeitlichem Heil, das den Kindern Israel gegeben ist, als daß Sebulon sollte am Meer wohnen, bis gen Sidon, und Isaschar mitten im Lande vom Meer wohnen und doch zinsbar gewesen ist den Königen von Assyrien.

c Den Segen Dan hat Simson erfüllet.
\quellseite{121}

\section{78. Das 1. Buch Mose 49}

d Gad hat seinen Segen ausgerichtet, da sie vor Israel herzogen.

e Asser hat gut Getreideland innen gehabt.

f Naphtali Segen ist erfüllet durch Debora und Barak.

\section{Summarium}

Hier ist eine Weissagung auf alle Stämme der Kinder Israel, wie es ihnen bis auf die Zukunft Christi gehen soll. Und ist hie sonderlich zu merken, daß Gott den Ehebruch an Ruben, B. 4., und den Todschlag an Simeon und Levi B. 6., so hart strafet, daß alle ihre Nachkommen nach ihnen solches entgelten, und die Strafe nun tragen müssen. Solches Exempel sollen alle Väter und Mütter fleißig merken, und ihrer Kinder halb sich in Gottes Furcht halten, auf daß nicht der Fluch auf sie erbe, welchen die Eltern mit ihren Sünden verdienet haben.

Die fürnehmste Weissagung aber ist von Juda, welcher Stamm zum Regiment durch David kommen, und dabei blieben ist bis auf die Maccabäer. Darum gehet die Prophezeiung also: Wenn es mit dem Königreich Juda werde ein Ende sein, und die Propheten aufhören, daß man alsdann des HErrn Christi gewarten soll. Da sollen die Heiden auch zum Reich Gottes kommen; welches Reich soll so herrlich sein, daß man Esel an die Weinstöcke binden, und die Kleider in Wein waschen soll, B. 11. Solches ist auf das einfältigste zu verstehen von der Lehre, welche im Neuen Testament geprediget wird. Da gehet Gottes Gnade in vollem Schwang, und auf das allerreichlichste, Vergebung der Sünde, und ewiges Leben wird allen Menschen angeboten, wie die Gleichnisse im Evangelio vom Hochzeit-Mahl ausweisen. Es wird aber auch an andern Orten in den Propheten das Evangelium dem Wein verglichen, dadurch die Herzen fröhlich werden, und der alte Adam in uns getödtet wird.

g Der Segen Joseph geht auf das Königreich Israel und ist ganz vom leiblichen Regiment gesagt, daß die Töchter (das ist die Städte im Lande) wohl regieret worden zeitlich, und viele Propheten und große Leute zu Ephraim hatten. Und wiewohl sie oft angefochten worden, gewannen sie doch. Und dieses Königreich war im Geschlecht Ephraim. Also bleibt der geistliche Segen und Reich auf Juda, und das leibliche Reich auf Ephraim.

h Benjamins Segen hat erfüllet der König Saul und die Bürger zu Gaba.

Kap. 50. v. 1. Da fiel Joseph auf seines Vaters Angesicht, und weinete *über ihm, und küssete ihn.

\section{Summarium}

In des Josephs zweier Söhne Segen, welche nach der Verheißung Juda von Christo das Größte haben. Denn das Königreich Israel ist mächtig und groß gewesen, und hat Herrn und
\quellseite{122}
köstliche Steine, das ist Könige und herrliche Propheten gehabt. Aber da sie denselben nicht folgen wollten, mußte ihr Land von Heiden verwüstet, und sie gefangen weggeführt werden.

Und zum Exempel, daß wir Gottes Wort und die Edelsteine seiner Kirchen, die frommen treuen Lehrer sollen lernen hoch halten, und ihnen folgen, sonst wird es nicht ungestraft bleiben.

\section{Votum}

Wir danken Gott, daß wir sein Wort und treue Kirchenlehrer haben; der wolle sie und uns segnen und behüten, und helfen, daß uns sein Wort ein Lehr-, Nähr- und Wehr-Wort sei allezeit, Amen.

\chapter{Das 50. Capitel}

\section{Vorrede}

Wie uns der gnädige und gütige Gott seine Verheißungen treulich hält, also sollen wir uns gegen ihn nach seinen Geboten gehorsamlich verhalten. Nun vergibt uns Gott aus Barmherzigkeit um Jesu Christi seines lieben Sohns willen, alle unsere manchfältige und schwere Sünden. O eine große unaussprechliche Gnade und Wohlthat! Derowegen sollen wir auch unsern Nächsten vergeben, was er uns zu Leid oder Schaden, mit Worten oder Werken gethan hat, und sollen ihm auch dienen und rathen nach unserm höchsten Vermögen; denn also lehret uns der Herr Christus, Luc. 6. und sagt: Seid barmherzig, wie auch euer Vater barmherzig ist. Und Matth. 5. spricht er: Liebet eure Feinde, thut wohl denen die euch hassen, bittet für die so euch beleidigen und verfolgen, auf daß ihr Kinder seid eures Vaters im Himmel. Davon höret auch an das 50. und letzte Kapitel.

\section{Jakob wird begraben, Joseph stirbt}

\quellseite{123}

\section{Summarium}

Ist erstlich ein Exempel, wie Joseph des Landes Brauch hält, mit seines Vaters Tod, V. 2. Daß er ihn aber ins Land Canaan führet, und befiehlt, wenn seine Nachkommen aus Egypten ziehen, daß sie seine Gebeine auch sollen mitführen, V. 25., geschieht alles anzuzeigen, daß er fest an Gottes Verheißung halte vom Land Canaan, auf daß seine Nachkommen an solcher Verheißung auch nicht zweifeln.

Darnach ist ein schön Exempel hie, daß Joseph seinen Brüdern nicht allein von Herzen vergibt, und sie nicht entgelten läßt, sondern deutet solche Mißhandlung zum besten, V. 19., auf daß auch wir lernen Böses mit Gutem bezahlen, und von Herzen einer dem andern vergeben.

\section{Votum}

GOTT regiere unsere Herzen, daß wir nicht allein unserm Nächsten herzlich vergeben, sondern auch das Böse mit Gutem belohnen lernen. Und solches um unseres Heilandes willen, Amen.

\part{Das zweite Buch Mose}

\chapter{Das 1. Capitel}

\section{Vorrede}

Gleich wie im ersten Buch Mose wiederholet und erkläret wird die Verheißung von dem gebenedeyten Weibesamen und Schlangentreter Christo unserm HErrn, und werden darin angehänget herrliche Verheißungen, auch schöne Exempel, wie die Patriarchen in wahrem Glauben an die Verheißung sich gehalten haben; also lehret das andere Buch Mose, wie Gott die heiligen zehn Gebote wiederholt hat auf dem Berg Sinai, und erzählet schreckliche Exempel, wie Gott die gestraft, welche seine Gebote muthwillig übertreten haben. So helfe nun der hochgelobte Sohn Gottes, daß wir das Buch Mose also anhören, auf daß wir uns alle daraus bessern, damit unser Leben Gott wohlgefalle, und unserm Nächsten zu Nutz diene, und daß wir zeitlich und ewig gesegnet werden, Amen. Darauf höret nun an das 1. Capitel.

\section{Der Kinder Israel Dienstbarkeit und Drangsal in Egypten}
\quellseite{124}



\section{Summarium}

Ist eine Historia, in welcher man fürnehmlich siehet, wie der Teufel Gottes Segen begehret zu hindern, und des Josephs Treu und Wohlthat vergessen, und seine Nachkommen vom Tyrannen Pharao greulich beschweret werden.

Auf daß man lerne, daß die rechte Kirche oder Gottes Volk in der Welt allweg sich haben leiden müssen. Aber daneben hilft Gott, daß, je mehr sie gedrückt werden, je mehr sie sich mehren und ausbreiten.

Am Ende aber ist ein fein Exempel von den Hebammen, welche dem tyrannischen Befehl Pharao nicht wollten gehorsam sein, und darum von Gott gesegnet werden, V. 17. Denn in unbilligen Sachen ist man der Obrigkeit keinen Gehorsam schuldig, sondern man muß auf Gottes Gehorsam fürnehmlich sehen, und der Menschen Gehorsam dahinten lassen; nach dem Spruch Petri, Apost. 5, 29.: Man soll Gott mehr gehorchen denn den Leuten.

\section{Votum}

Der allmächtige und barmherzige GOtt erhalte seine Kirche, und thue wieder Gutes zeitlich und ewig allen, die seiner Gemeine und bedrängten Christen Gutes erzeigen, hochgelobt in Ewigkeit, Amen.

\chapter{Das 2. Capitel}

\section{Vorrede}

Wie kein edlerer Schatz auf dieser Welt zu finden ist, weder Gottes Wort, als darin allein die seligmachende Weisheit, Erkenntniß und Furcht Gottes, lebendiger Trost und ewiges Heil verborgen ist: also mögen und sollen sich die Leute wohl für die seligsten achten und halten, denen Gott solchen theuren Schatz vertrauet und gönnet, und dagegen für die unseligsten, die denselben nicht haben oder nicht haben wollen. So danket Gott alle herzlich für die besondere, ja für die höchste Gnade, daß ihr auch habt die heilige Biblia, und daß euch die Schriften der heiligen Propheten und Apostel zu Lehr, Trost und Warnung fürgelesen werden; auch wendet nöthigen Fleiß an, daß zu eurem großen Nutz solches geschehe, in Kraft des heiligen Geistes. Und höret jetzt an das 2. Capitel.

\section{Moses Geburt, Auferziehung, Flucht und Heirath}
\quellseite{125}

a Masa heißt ziehen; daher heißt Mose gezogen, nämlich aus dem Wasser.

b Gersom heißt ein Fremdling oder Ausländer.

†Non est in Hebræo.

‡(Elieser) heißt Gott mein Helfer.

\section{Summarium}

Die Historia von Mose hebt sich hie an, in welcher wir sehen, wie Gott der Welt und des Teufels Anschläge wunderbarlich brechen und zu nicht machen kann; weil eben des Königes Tochter Mosen erhält, V. 6., welcher des jüdischen Volks Herzog sein, und Egypten ganz wüst machen sollte.

Darum deutet das 11. Capitel B. 23. an die Ebräer, diese Historia also, und spricht:

Durch den Glauben ward Moses, da er geboren war, drei Monat verborgen von seinen Eltern, darum, daß sie sahen, daß er ein schön Kind war, und fürchteten sich nicht für des Königs Gebot. Durch den Glauben wollte auch Moses, da er groß ward, nicht mehr ein Sohn heißen der Tochter Pharao, und erwählete viel lieber mit dem Volk Gottes Ungemach zu leiden, denn die zeitliche Ergötzung der Sünde zu haben. Und achtete die Schmach Christi für größern Reichthum denn die Schätze Egyptens, denn er sahe an die Belohnung.

\section{Datum}

Gott unser Vater, der Mose erhalten und regieret hat, der wolle auch uns mit Gnaden bedenken, und wolle heut auch der Welt und des Teufels Anschläge hindern und brechen, Amen.
\quellseite{126}

\chapter{Das 3. Kapitel}

\section{Vorrede}

Die heilige Biblia ist das beste Lehrbuch und Trostbuch. Ein Lehrbuch ist die heilige Biblia: denn durch Gottes Wort und dessen Uebung und Betrachtung zündet Gott an und erwecket rechte Erkenntniß Gottes und wahren Glauben, unterrichtet auch, wie man Gott in rechtem Glauben und gutem Gewissen wahrhaftig und gefällig dienen, desgleichen wie ein Christ sein Leben und Stand recht anstellen und führen soll. Ein Trostbuch ist auch die heilige Biblia: denn daraus angefochtene, bekümmerte, gefängstigte Herzen zu jederzeit rechten, gewissen und beständigen Trost schöpfen und erlangen mögen, durch Hülfe des heiligen Geistes, daß Buch die Biblia ist, als des rechten himmlischen Kanzlers. So höret an das 3. Kapitel.

\section[Mose wird berufen, die Kinder Israel aus Egypten]{Mose wird berufen, die Kinder Israel aus Egypten zu führen}

a Ich werde seyn wenn ihr dahin kommet, so will ich bei euch seyn, und mich so erzeigen, daß ihr erkennen sollt, daß ichs sei.
\quellseite{127}

\section{Summarium}

Ist ein tröstlich Exempel, wie Gott dem Tyrannen nicht ewig zusehen, sondern seinem Volk, das unbillig beschweret ward, helfen und ihr Gebet erhören will, um seiner Verheißung willen. Darum offenbaret er sich Mose im brennenden Busch, B. 2., und fordert ihn zu solchem Werk und sagt, sein Name heiße Gott des Abraham, Gott des Isaak, Gott des Jakob, B. 6. Das ist, ein Gott, der sich durch seine Verheißung mit diesen Vätern verbunden habe, und ihnen zugesagt, daß er ihre Nachkommen nicht allein aus Egypten erlösen, und ins gelobte Land bringen wolle, sondern auch aus ihren Nachkommen den Samen kommen lassen, durch welchen alle Völker auf Erden sollen gesegnet, das ist, von Sünde und Fluch ledig, und ewig selig werden.

\section{Votum}

Der starke allmächtige GOtt erhalte auch bei uns und in der ganzen Welt seine Kirche; und verleihe, daß wir deren gläubige und gehorsame Gliedmaßen sein und allezeit bleiben, Amen.

\chapter{Das 4. Kapitel}

\section{Vorrede}

Von Natur haben wir allesamt erkaltete Herzen, die zu Gottes Wort keine Lust haben. Fabeln und Mährlein höret man viel lieber, da sperret man Maul und Nasen auf und ist niemand keine Zeit zu lang, man höret solches mit Lust und Liebe. In der Kirche aber wartet man nicht gerne lange. O man höret sich Gottes Wort und der Predigen bald überdrüssig. Dies betrachtet, und weil Gottes Wort eine Werkstatt ist des heiligen Geistes und der Wagen, darauf der heilige Geist zu uns kommt und bei uns eingehet, zu unserer Besserung und Trost, so sollt ihr nimmermehr Gottes Wortes überdrüssig werden. Und derowegen höret auch jetzt an das 4. Kapitel.

\section[Mose wird in seinem Beruf mit der Gabe]{Mose wird in seinem Beruf mit der Gabe Wunder zu thun gestärket}
\quellseite{128}

\section{Summarium}

Dies ist eine historia, wie Moses sich gewehret und dem Beruf nicht folgen wollte, V. 1. Darum, daß er wohl sahe, wie mancherley Gefahr es auf ihm hatte, so er wider des Königs Willen sollte die Juden aus dem Lande führen. Aber Gott tröstet ihn, er wolle bei ihm sein und ihn lehren, was er sagen soll, V. 3. Gibt ihm auch Macht, drei Wunderwerke zu thun, mit der Schlange, ausfätzigen Hand, und dem Blut; auf daß er ein gewiß Zeugniß hätte, was er thäte, daß ers aus Gottes Befehl thäte. Applicatio. Solches Exempel ist tröstlich allen denen, so Gottes Beruf vor ihnen haben, daß sie unverzagt sein und Gott sollen sorgen lassen, wie sie es ausrichten, was er ihnen befohlen hat. Gleich wie man siehet, wo der Beruf nicht ist, es sei Macht und Vermögen so groß es wolle, so kann man doch nichts Fürnehmliches ausrichten.

*Blutbräutigam) das ist, sie ward zornig, und sprach: Es kostet Blut, daß du mein Mann bist und muß mein Kind beschneiden; welches sie ungern that, als das eine Schande war unter den Heiden. Bedeutet aber des Gesetzes Volk, welches gerne wollte Gott haben, aber es will das Kreuz nicht leiden, noch den alten Adam beschneiden lassen, bis es thun muß.

\section{Summarium}

Ist erstlich ein tröstlich Exempel, wie Gott den Tyrannen wegnimmt, der Mose nach dem Leben stund, V. 19. Darnach, wie Moses, da er seinem Beruf nachkommen wollte, in die Gefahr kommt, daß ihn Gott tödten wollte, V. 23., welches darum geschah, eben wie oben mit Jacob, daß er über Gottes Verheißung etwas leiden, und mit dem Glauben darüber fechten müßte, auf daß er hernach desto fester stünde, und stärker würde. Denn unversuchte Leute taugen zum Predigtamt eben so wenig, als zu andern Händeln.
\quellseite{129}

\section{Votum}

GOtt der Vater des Lichts, von dem alle gute Gaben von oben herab kommen, segne unsern Beruf, unsere Aemter und alle christliche Händel und Orber, und wende ab alle Hindernisse und Gefahr, Amen.

\chapter{Das 5. Kapitel}

\section{Vorrede}

Gott strafet die Gottlosen und hilft den Frommen. Er gibt aber den Gottlosen eine zeitlang Glück und Wohlfahrt, ehe er ihnen nach ihrem Verdienst im Zorn ablohnet, und die Frommen läßt er das Kreuz auch eine zeitlang wohl drücken und fühlen, ehe denn er ihnen nach seiner Verheißung in Gnaden aushilft. Derowegen sollen wir uns in Gottes Weise schicken lernen, sollen Buße thun, geduldig sein, beten und Gottes Hülfe hoffen und erharren. So wird es denn zuletzt wohl gehen denen, die fromm bleiben und sich recht halten: aber die Gottlosen zuletzt ausgerottet und wie das Gras abgehauen, und wie das grüne Kraut verwelken werden, Psalm 37. Ferner höret an das 5. Kapitel.

\section[Mose wird von Pharao verachtet, und das Volk]{Mose wird von Pharao verachtet, und das Volk noch mehr geängstiget}

a Das ist, deine armen Leute müssen unrecht haben und Sünder sein, man suchet Schuld zu deinem Volk.
\quellseite{130}

\section{Summarium}

Sind zwei sonderliche Exempel: das erste, wie die Tyrannen und Gottes Feinde sich stellen, wenn ihr Verderben am nächsten ist, nämlich, daß sie von Gott nichts wollen wissen, V. 2., und mit den Leuten je länger je unbarmherziger umgehen, auf daß sie ihr eigen Urtheil forttreiben und zum Ende bringen, V. 6.

Das andere Exempel ist, daß wir an den Kindern Israel sollen lernen, wie sich Gott dazu pflege zu schicken, wenn er will helfen, nämlich, daß ers hinter sich anfühet, und wenn man der Hülfe gewartet, allererst am tiefsten ins Unglück fällt. Darum gilt es hie hoffen und harren, bis die Hülfe hernach kommt; wie es den Kindern Israel ging, ob es sich gleich im Anfang sehr übel dazu schickte.

\section{Votum}

Der treue GOtt, der ein GOtt ist alles Trostes und der Geduld, wirke und erhalte Geduld allen, die jetzt das liebe Kreuz tragen; hindere auch und stürze alle boshaftigen und unwilligen Feinde seines heiligen Namens, Amen.

\chapter{Das 6. Kapitel}

\section{Vorrede}

Wir müssen bekennen, daß wir alles Kreuz mit unserm sündhaftigen Wesen und Leben wohl verdienet haben. Wir sollen bekennen, daß es eine Strafe ist von Gott, viel geringer, denn unsre Sünden sind, Judith 8. Da bekennen sollen wir, daß wir auch die größten zeitliche, und darnach auch die ewige Strafe verdienet hätten, wenn Gott mit uns ins Gerichte gehen und mit uns handeln wollte nach seiner Gerechtigkeit und unserm Verdienst: wie wir denn solches alle bekennen müssen, wir haben so heilig gelebt, als wir immer vermocht. Und dieweil Gott unser dort mit der ewigen wohlverdienten Strafe um seines lieben Sohnes willen verschonen will, sollen wir uns nicht weigern allhier ein wenig Kreuz zu leiden, sondern sollen mit dem Propheten Micha sagen, Mich. 7.: Ich will den Zorn des HErrn tragen, denn ich habe ihm gesündiget. Davon höret an das 6. Kapitel.

\section[Mose bekommt von Gott einen neuen Befehl, beschreibet]{Mose bekommt von Gott einen neuen Befehl, beschreibet sein Geburtsregister}

a Nicht offenbaret) die Patriarchen haben Gott wohl erkannt, aber eine solche öffentliche gemeine Predigt war zu der Zeit von Gott noch nicht aufgegangen, wie durch Mose und Christum geschehen ist.

b Meine Hand gehabet) das ist, geschworen, wie oben. 1 Mos. 22, 16.
\quellseite{131}


\section{Summarium}

Ist eine Historia, wie Gott Mosen und das Volk tröstet, auf daß sie fest an Gottes Zusagung halten, und durch die Beschwerung der Frohndienste sich nicht sollten dahin bewegen lassen, daß sie darum ihrem Gott und Mosi seinem Diener, ungehorsam würden, V. 6.

Aber wie der Text die sagt, V. 9., sie hörten ihn nicht für der Angst und Seufzen, und für harter Arbeit.

Also gehet es noch, daß man sich des Unglücks halben allerweg mehr bekümmert, denn man sich Gottes Gnade und Gottes Worts tröstet; so doch Gottes Wort billig sollte von uns am größten geachtet werden, als der höchste Schatz und Trost in allerlei Jammer und Anfechtung.

\section{Votum}

Der barmherzige Gott lasse ihm alle befohlen sein, die noch das Creuz tragen, und in Noth und Gefahr stecken, und helfe ihnen an Leib und Seele, daß sie ihn ewig mit uns dafür loben und preisen, Amen.

\chapter{Das 7. Capitel}

\section{Vorrede}

Wir sind in dieser Welt in des Teufels Reich, nicht anders denn wenn ein Pilgrim in eine Herberge käme, da er wüßte, daß sie alle im Hause Räuber wären. Wenn er dahin kommen würde, würde er sich dennoch rüsten und aufs beste als er könnte vorsehen, und nicht viel schlafen: also sind wir jetzund auf Erden, da der böse Geist ein Fürst ist, und hat der Menschen Herzen in seiner Gewalt und thut durch sie was er will. Ferner, wie der Teufel einen Hirsch umbringet, wenn er ihm auf den Kopf springet und ihm das Gehirn ausfrißt, oder greift ihn bei der Kehle, und beißet sie ihm entzwei: also ist auch der Teufel, wenn er einen Menschen in seine Gewalt bekommt, so kann man seiner nicht leichtlich los werden, er führet den Menschen aus einer Sünde in die andere, und thut ihm Schaden an Leib
\quellseite{132}
und Seele. Derwegen wollen wir wachen und beten, daß der Teufel unser ja nicht mächtig werde, und uns nicht führe in Sünde, Schand und Schaden. Dieß betrachten wir auch bei dem 7. Kapitel.

\section{Verwandlung des Wassers in Blut}

a Gott kann Pharao weder mit Plagen, noch mit Wunder, noch mit Dräuen bessern, er muß ihn lassen ersaufen im Meer.
\quellseite{133}

\section{Summarium}

Dies ist eine Historia, wie Pharao der Tyrann sich an kein Zeichen kehret, sondern in seinem Unglauben und Ungehorsam bleibet, V. 4. Solches aber muß dazu dienen, auf daß Gott ihn desto härter angreife, und Gottes Ehre und Kraft desto herrlicher geoffenbaret werde in ganz Egypten. Eben also gehet es mit den gottlosen Tyrannen noch: je mehr sie sich wider Gottes Wort und Willen setzen, je größer und tröstlicher ist’s, daß sie am Ende darüber zu Boden gehen; Gottes Wort aber bleibet ewiglich.

\section{Votum}

Der HErr Zebaoth, der starke GOtt, steure allen Tyrannen, und tröste und erhalte dagegen auch alle, die seinen heiligen Namen anrufen, bekennen, loben und preisen, Amen.

\chapter{Das 8. Kapitel}

\section{Vorrede}

Gottes Strafen sollen unsere Bußprediger sein; und wenn Gott züchtiget, so soll man nicht ärger, sondern frömmer werden. Denn wie St. Paulus sagt, Röm. 6.: Unser alter Mensch wird gekreuziget, daß der sündhaftige Leib aufhöre, daß wir hinfort der Sünde nicht dienen. Eben dies meinet auch St. Petrus, da er schreibt, 1 Epist. Kap. 4.: Wer am Fleisch leidet, der höret auf von Sünden, daß er hinfort, was noch hinterstelliger Zeit im Fleisch ist, nicht der Menschen Lüsten, sondern dem Willen Gottes lebe. Derowegen wer von Gott gestraft wird, der thue Buße und lasse sich den Teufel und seine Werkzeuge nicht verführen, sondern folge Gottes Wort in Gehorsam, so wird ihm folgen Gottes Segen zur Wohlfahrt. Darauf höret an davon das 8. Kapitel.

\section{Egypten wird mit Fröschen, Läusen und Ungeziefer geplaget}

\section{Summarium}

Ist ein Exempel, wie die Tyrannen verzagt werden, wenn Gottes Gericht über sie kommt, V. 8. Aber so bald sie Luft haben, und die Strafe nachläßt, wieder hinfallen und bleiben wie sie gewesen sind, V. 15.

Daß aber der Egypter Zauberer etliche Zeichen nachthun, und endlich an den Läusen zu Schanden werden, V. 18., ist ein Zeugniß, daß solche Teufels-Künstler mehr nicht können, denn Gott aus seinem Zorn ihnen selbst, oder andern zur Aergerniß und Strafe verhänget, aber sie
\quellseite{134}
müssen endlich an dem Geringsten zu Spott und Schande werden.

a Ungeziefer) das die Griechen heißen Cynonomia, sind alle bösen Würmer, so da Schaden thun im Felde, Raupen, Fliegen, Zweyfalter, Ameisen, Käfer, Brenner und dergleichen Geschmeiß, das Bäume und Gewächse verderbet.

\section{Summarium}

Ist ein Exempel, wie die Tyrannen in der Noth sich füßlich demüthigen, aber doch immer ihrem eigenen Willen und nicht dem Wort Gottes folgen; wie man an unsern Papisten auch sieht: da hilft keine Strafe, kein Ueberweisen, kein Vermahnen, sondern eben wie Pharao, bleiben sie verstockt und ungebessert, ob sie gleich selbst bekennen müssen, daß ihr Wesen mit Gottes Wort nicht übereinstimmt. Auf solche Pharaonische Sünde gehört auch ein Pharaonisch plötzliches Verderben.

\section{Votum}

Der treue GOtt behüte uns für Pharaonischer Verstockung, und regiere uns durch den heiligen Geist, daß wir unsern Willen brechen, und seinem heiligen Willen nach seinem Wort folgen, bis zum seligen Todesstündlein, Amen.

\chapter{Das 9. Kapitel}

\section{Vorrede}

Des alten Adams oder des alten Menschen Kreuzigung steht in erster Buße, daß man Reue habe über die Sünde und sich fleißig für derselben hüte, und also das Fleisch im Zaum halte. Das kömmt manchem wohl so sauer an als der Tod, sonderlich wenn man von böser alter Gewohnheit soll ablassen, seinem eigenen
\quellseite{135}
Sünde abbrechen und steuren, böser Gesellschaft sich verziehen, der Welt Ungunst auf sich laden, ein Gewinnlein fahren lassen, so fort an. Derowegen weil wir, als die wir Christum angehören, unser Fleisch kreuzigen sollen, samt den Lüsten und Begierden: so laßt uns Gott um Hülfe darzu anrufen; denn die Natur kann ihr selber nicht widerstreben noch abbrechen. Darauf höret auch an das 9. Kapitel.

\section{Pestilenz. Schwarze Blattern. Hagel}

\section{Summarium}

Ist ein Exempel, wie Gott mit Strafen nicht aufhören wolle, wo man von Sünden nicht ablassen und seinem Wort nicht folgen will. Wiederum, wo man folget, da will er schützen und behüten, ob es gleich sonst allenthalben übel zugeht; wie er seine Juden behütet, da sonst das ganze Egypten um der Sünde und Ungehorsams willen mußte leiden.
\quellseite{136}

\section{Summarium}

Ist ein sonderlich Exempel wie oben, daß die verstockte Herzen sich fälschlich demüthigen, als wollten sie fromm werden, begehren Gnade und Fürbitte. Aber es ist ihnen nicht länger Ernst, bis das Unglück fürüber ist. Ob wohl solches frommen Herzen wehe thut, sollen sie doch diesen Trost daraus nehmen, daß die Gottlosen mit solcher falscher Demuth ihnen selbst ihre Sache nur ärger machen, und zu größerer Strafe und Zorn Gott verursachen.

\section{Votum}

Gott unser lieber Vater helfe uns allen, daß wir rechtschaffene Buße thun, und in der Frömmigkeit beständig fortfahren, und schütze und erhalte in diesen gefährlichen Zeiten seine Christenheit und alle Gläubige, Amen.

\chapter{Das 10. Capitel}

\section{Vorrede}

Wie Gott mit Strafen nicht aufhöret nach den Dräuungen des Gesetzes, wo man von Sünden nicht ablässt, und seinem Wort zu einem christlichen Leben nicht folget: also wo man folget, und einen ehrbaren heiligen Wandel führet, da beschützet und behütet Gott nach den Verheißungen des Evangelii, ob es gleich sonst allenthalben übel zugehet. Derwegen wer nicht will gestraft sein, der thue Buße und lebe beständig in Gottesfurcht; und wen Gott erhalten soll unter künftigem Unglück, der beharre auf dem Wege der Gerechtigkeit. Denn, wie an die Ebräer am 2. gesagt wird: Eine jegliche Uebertretung und Ungehorsam hat empfangen seinen rechten Lohn. Also spricht David Psalm 37.: Laß vom Bösen und thue Gutes, und bleibe immerdar. Dies lehret auch das 10. Kapitel.

\section{Heuschrecken und Finsterniß}
\quellseite{137}
* Die Heuschrecken heißen hie nicht Hagab auf Hebräisch, wie an etlichen Oertern, sondern Arbe. Es sind aber vierfüßige fliegende Thiere, und rein zu essen, wie Hagab, Lev. 11, 12. aber uns unbekannt, ohne daß sie Heuschrecken gleich sind.
\quellseite{138}

\section{Summarium}

Dies ist eine Historia, wie über die andern Plagen Gott Egyptenland mit den Heuschrecken verderbet, B. 13., und darnach mit Finsterniß geschlagen habe, B. 21., auf daß man lerne, wo man Gottes Wort nicht folgen will, daß da kein Aufhören der Strafe zu hoffen sei. Darum will auch Gott, daß solcher Exempel bei Kindes Kindern soll gedacht werden, auf daß man beides lerne: sich für solchem Gott fürchten, der die so ihm ungehorsam sind, dermaßen strafen will; und auf solchen Gott trauen, der seinem Volk so wunderlich helfen kann.

Es ist auch hie ein Exempel, wie Pharao je länger je verstockter wird, und die Kinder Israel nicht allein nicht ziehen, sondern auch Moses will tödten lassen, B. 18.

\section{Votum}

Gott, unser aller lieber Vater, behüte uns für Pharaonischer Blindheit, und verleihe, daß wir deinem Wort gehorsamlich folgen, und daß wir die gegenwärtigen Strafen zu wahrer Buße gebrauchen, Amen:

\chapter{Das 11. Kapitel}

\section{Vorrede}

Was Gott den Unbußfertigen dräuet, das läßt er endlich wahrhaftig über sie ergehen; und da schonet er keines, wer er auch ist vor der Welt. Denn es ist kein Ansehen der Person für Gott, und bei ihm trägt niemand vor sein großes Geschlecht, ehrliche Freundschaft, ansehnliches Alter, hohe Weisheit, vielfältige Erfahrung, große Kunst, herrliche Aemter, mächtige Gewalt und dergleichen, was den Menschen in der Welt bei den Leuten ein Ansehen machen kann. Deswegen so wollen wir unsere Sünde erkennen, die Gott abbitten, davon abstehen, und ein Christen-Leben führen, daß wir nicht zeitlichen und ewigen Jammer über uns führen. Denn um Jesu Christi willen vergibt Gott den Bußfertigen alle Sünden und lindert ihnen aus Gnaden die wohlverdienten Strafen. Darauf höret auch an das 11. Kapitel.

\section{Auszug aus Egypten befohlen}

a (Großer Mann) das saget er darum, daß es Wunder ist, wie er nicht getödtet ist von den Egypten; sie haben sich müssen eines Aergern und Aufruhrs fürchten.

\section{Summarium}

Dies ist eine Historia, wie Moses Pharao dräuet mit der letzten Strafe, B. 4. Aber Pharao achtet solches Dräuen nicht, bis ers erfährt, B. 9. Wie denn die Gottlosen allwege nicht eher glauben, bis der Glaube ihnen in die Hände kommt, und sie es erfahren. Daß aber der HErr heißet silbernes Geschirr von den Egypten entlehnen, und dasselbe ihnen entwenden, B. 2., ist darum geschehen, auf daß den Kindern Israel ihre Arbeit belohnet würde, welche ihnen vom Könige mit unbilliger Beschwerung ward aufgelegt. Denn Gott will, daß die Regimente billiger Weise geführt, und alle Tyrannei verhütet, und sonderlich einem jeden seine Mühe und Arbeit belohnet werde. Wo solches weltliche Obrigkeit nicht will thun, muß es Gott selbst schaffen, aber durch ein hartes Mittel. Denn die Egyptier hätten die Juden leichter können bezahlen, denn daß sie ihnen alle ihre Schätze und Kleinode anhängen, und darnach Leib und Leben auch darüber verlieren, wie im rothen Meer geschehen ist.

Digitized by Google
\quellseite{139}

\section{Votum}

Helfe GOtt, daß wir mit den Gottlosen sein Wort und Träumungen ja nicht verachten, sondern denen Glauben geben, und seinem Wort zu täglicher Besserung mit allen Frommen folgen, Amen.

\chapter{Das 12. Capitel}

\section{Vorrede}

Gott ist den Unbußfertigen ein eifriger Gott, den Gläubigen aber ist er ein gütiger Gott. Dies lehret er selber im 2 Buch Mose im 20. Kapitel, da er sagt: Ich der HErr dein Gott, bin ein eifriger Gott, denen, die mich hassen; und thue Barmherzigkeit an viel Tausenden, die mich lieb haben und meine Gebote halten. Derowegen wer unter den mannigfaltigen Strafen, damit der gerechte Gott jetzt die Welt wegen der überhäuften großen Sünden, Schande und Lastern heimsucht, will Schutz, Trost und Hülfe haben, der bitte heute abermal, Gott wolle ihm gnädig sein und alle seine Sünde vergeben; er halte sich auch in wahrem Glauben an Jesum Christum, unsern HErrn, und lebe nach seinem Willen; also wird er sein und bleiben ein Verborgener des HErrn, und wird mit Freuden singen mit dem König und Propheten David: Der HErr decket mich in seiner Hütten zur bösen Zeit; er verbirget mich heimlich in seinem Gezelt: denn der HErr verbirget die Seinen heimlich bei ihm, und verdecket sie in seiner Hütten, Ps. 27. und 31. Davon höret an das 12. Kapitel.

\section[Stiftung des Osterlamms. Erwürgerung der Erstgeburt. Des Auszugs]{Stiftung des Osterlamms. Erwürgerung der Erstgeburt. Des Auszugs Anfang}

a Zusammen kommet) das ist, sie sollen predigen, das Benedikte und Gratias beten für das Passah und Erlösung aus Egypten. Darum gebraucht er das Wörtlein Kara, welches heißt predigen, lesen in der Versammlung.
\quellseite{140}

\section{Summarium}

Ist eine Historia, wie und aus was Ursache Gott den Juden das Osterfest habe eingesetzt.

Zum ersten, daß es soll sein ein ewig Gedächtniß der herrlichen Erlösung aus Egypten. Denn Gott will, daß man seiner Wohlthat nicht vergessen, sondern ihm dafür danken und derselbigen in allen Nöthen zum Trost brauchen soll.

Zum andern, daß solch Osterfest auch soll ein Fürbild sein der geistlichen Erlösung von Sünde und Tod, durch das Blut Christi, des unschuldigen Lämmleins, wie St. Paulus solches deutet 1 Cor. 5, 7.: Daß unser Osterlamm sei Christus, für unsere Sünde geschlachtet. Solches Osterlamm sollen wir täglich essen, das ist, durch den Glauben uns seines Todes und Verdienstes täglich trösten; und den Sauerteig ausfegen, das ist, der Sünden wehren, und ihr nicht folgen.
\quellseite{141}

b) Passah) heißt ein Gang, darum, daß der HErr durch Egyptenland des Nachts ging und schlug alle Erstgeburt todt. Bedeutet aber Christus Sterben und Auferstehen, damit er von dieser Welt gegangen ist, und in demselben Sünde, Tod und Teufel geschlagen und uns aus dem rechten Egypten geführet hat zum Vater; das ist unser Passah oder Ostern.

\section{Summarium}

Ist ein tröstlich Exempel, in welchem wir sehen, wie Gott die Seinen endlich errettet und der Tyrann Pharao mit seinem harten Sperren nichts anderes hat ausgerichtet, denn daß er sich und das ganze Land Egypten in alles Unglück gebracht, und jetzt nach solchem Schaden ohne Dank das thun muß, daß er zuvor, ehe einiger Schaden geschehen war, zu Dank nicht hat wollen thun. Also muß Gott die Tyrannen zum Gehorsam bringen, wenn sie sein Wort so gar nicht achten und sich für ihm nicht demüthigen wollen. Wer auf die Historien und tägliche Erfahrung siehet, wird solcher Exempel viel finden, daß es den großen Herren an ihren Anschlägen fehlet und Gottes Wille mit Macht ungewehret fortgehet.

\section{Gebet}

Der treue GOtt erbarme sich heut auch seiner verfolgten Christenheit, und steure und wehre allen bösen Anschlägen und Tyrannen, um JEsu Christi seines Sohns, unseres HErrn willen, Amen.

\chapter{Das 13. Kapitel}

\section{Vorrede}

Wohl denen, die an Gottes Wohlthaten gedenken und ihm dafür dankbar sein. Denn je öfter und fleißiger ein Mensch Gott danket für die empfangene Wohlthaten, je mehr ihm Gott ferner, so lange er lebet, Gutes erzeiget. Dieses lehret David, der in der Person Gottes, Ps. 50. also saget: Wer Dank opfert, der preiset mich, und das ist der Weg, daß ich ihm zeige das Heil Gottes, das ist, mit Danksagung bauet man unserm HErrn Gott den Weg und lockt ihn, daß er mit Wohlthun desto öfter und reichlicher wiederkomme. Derwegen lobsinget eurem Gott, und danket ihm täglich für seine Gnade und Barmherzigkeit, wie er solche euch und den Euern, an Leib, Seele und Nahrung, erzeiget hat. Darauf höret nun an das 13. Kapitel.

\section{Heiligung der Erstgeburt. Des Auszugs Fortgang}

a) Abib) ist der Monat, den wir April heißen; denn die Hebräer heben ihr Neujahr an nach der Natur, wenn die Ding wieder neu grünet und wächst und sich fühlet. Darum heißet er auch Mensis novorum, daß denn Alles neu wird.
\quellseite{142}

b Ungesäuert Brod) so hart wird der Sauerteig verboten, daß man ja das lautere Evangelium und Gottes Gnade, nicht unser Werk und Gesetz soll predigen nach der Auferstehung Christi. Wie St. Paulus 1 Cor. 5, 6. auch zeiget. Und ist solch Essen nichts Anderes, denn glauben an Christum.

c Schilfmeer) die Griechen heißen es das rothe Meer, von dem rothen Sand und Boden; aber die Hebräer heißen es das Schilfmeer, von dem Schilf.

d Gerüstet) hebräisch gefünfft. Was das sei, lassen wir Andere suchen; ob es sei, daß sie bei fünf Haufen gezogen, oder bei fünf neben einander gegangen sind, oder was solch fünf sei. Denn von der fünften Rippe, so das Harnisch erreicht am Leibe (wie die Juden hie \unsicher{fabeln}), verstehen wir nichts.

\section{Summarium}

Ist ein neuer Befehl Gottes, wie man neben dem Opferfest alle Erstgeburt dem HErrn heiligen und mit einem Opfer lösen soll, V. 1. Zum Gedächtniß, daß Gott sein Volk aus Egypten habe ledig gemacht, da er alle ihre Erstgeburt durch den Engel in einer Nacht hat würgen lassen.

Solches ist uns zur Lehre fürgeschrieben, daß wir ja die Werke Gottes nicht aus der Acht lassen, sondern derselbigen gedenken, Gott dafür danken und uns damit in allerlei Anfechtung trösten sollen.

Das Feuer und Wolkensäule, in welcher Gott für seinem Volk herzog, V. 21., sind ein Fürbild des Wortes Gottes und des heiligen Geistes, dadurch wir hie auf Erden, gleich wie in einer Wüste geführt, getröstet und erhalten werden.

\section{Gebet}

GOtt der heilige Geist regiere uns durch sein Wort, daß wir alle Sünden fliehen, und alle Tugenden üben, und die Wallfahrt dieses Lebens in wahrem Glauben und gutem Gewissen seliglich vollenden, Amen.
\quellseite{143}

\chapter{Das 14. Kapitel}

\section{Vorrede}

Gott tröstet und erhält die Frommen und schrecket und stürzet die Gottlosen. Derowegen wer in Sünden lebet, der thue Buße, daß ihn Gott nicht strafen dürfe. Wer aber in Gottesfurcht lebet, und ist daneben ein betrübter Kreuzträger, oder stecket in Noth und Gefahr, der sei geduldig und bete; Gott wird ihn wahrhaftig erretten und erhalten. Diesen Trost lehret David im 93. Psalm und saget: Die Wasserwogen im Meer sind groß und brausen gräulich; der HErr ist noch größer in der Höhe; und St. Johannes, da er spricht 1 Epist. 4. Cap.: Der in euch ist, der ist größer, denn der in der Welt ist. Dies erinnert auch das 14. Kapitel.

\section[Des Auszugs Vollendung. Der Egypter Untergang im rothen]{Des Auszugs Vollendung. Der Egypter Untergang im rothen Meere}


a Heil) Hülfe.

b Schreiest du) merckt hie ein trefflich Exempel, wie der Glaube kämpfet, zappelt und schreyet in Nöthen und Fährlichkeit, und wie er sich an Gottes Wort bloß hält und von Gott Trost empfähet und überwindet.
\quellseite{144}

e (Erleuchtet) das ist, es war ein Wetterleuchten in den dicken Wolken.


\section{Summarium}

Ist eine treffliche Historia, in welcher wir sehen, wie Gott die Seinen, ehe er sie von den Egyptern gar erledigt, zuvor in einen neuen Schrecken und Gefahr kommen läßt, V. 5. Wie man sonst auch siehet, daß Gott nicht eher Errettung schaffet, es sei denn zuvor mit aller menschlichen Hülfe und Mittel verloren.

Da gehet des Glaubens Kampf an, daß man in solcher Noth allein auf Gottes Barmherzigkeit, Allmächtigkeit und Verheißung sich verlasse. Wie denn Moses seinen Juden auch hie prediget: Fürchtet euch nicht, und sehet zu, was der HErr heute für ein Wunderwerk an euch thun werde, V. 13.

Also kehret es sich denn um: das Volk Gottes ist am ersten erschrocken, und sieget doch endlich durch die Hülfe Gottes. Die Egypter sind am ersten sicher und freudig, endlich liegen sie im rothen Meer, da es am tiefsten ist. Dies soll das Ende sein aller Feinde Gottes und Tyrannen.

\section{Votum}

Wir danken Gott herzlich für beides, daß er an den Israeliten und Egyptern gethan hat; der strafe und stürze ferner alle Unbußfertige, und erhalte die Frommen in Gnaden, und segne und erhalte sie ewiglich, Amen.

\chapter{Das 15. Kapitel}

\section{Vorrede}

Nach dem Glauben ist die Danksagung der fürnehmste und vortrefflichste Gottesdienst. Derowegen wie der Vogel am Gesang und der Topf am Klang erkannt wird: also ist der Mensch ein rechter Christ, der gern betet und Gott für seine Wohlthaten herzlich danket. So wollen wir nun auf Gottes Barmherzigkeit und Gaben Achtung geben, und als wahre Christen wollen wir ihm auch täglich dafür danken. Und wollen also diesen Gottesdienst alhie anfangen, weil wir ihn fürnehmlich in jenem Leben mit den lieben Engeln Gottes treiben, und mit ihnen fröhlich und ewig singen werden: Heilig, heilig, heilig ist der HErr Zebaoth, Es. 6. Das lehret auch das 15. Kapitel.

\section{Lobgesang Mosis. Bitteres Wasser süß gemacht}
\quellseite{145}

a Mara) heißt bitter, und bedeutet Leiden und Anfechtung, welche durch das Kreuz Christi im Glauben auch süß werden, Matth. 11, 29. 30. Mein Joch ist süß.

\section{Summarium}

Ist ein Danklied für die herrliche Erlösung am rothen Meer, welche Moses rühmet für ein sonderliches Trostbild, daß sich Gott nicht allein mit großer Kraft wider seine Feinde leget, sondern auch mit rechter Gnade gegen seinem Volk hat sehen lassen. Auf daß auch wir auf solchen Gott unser Vertrauen setzen, und uns seiner Hülfe in allerlei Noth trösten lernen.

Aber wie bald sie solcher Wohlthat vergessen haben, zeigt an das Volk am dritten Tage hernach, welches, da es ihnen an Wasser mangelt, wider Mosen murrete, V. 24. Aber Gott hilft auch in derselbigen Noth, V. 25.; und verheißet durch Mosen ferner weiter zu helfen, V. 26., so fern man sich seines Worts und Befehls halten will.
\quellseite{146}
Amen.

Lob, Ehr und Preis sei GOtt für alle seine Wohlthaten; der regiere und befestige uns, daß wir ihm ferner täglich mit Mund, Herz und gottseligem Wandel danken bis zu dem ewigen Leben, Amen.

\chapter{Das 16. Kapitel}

\section{Vorrede}

Zwei Dinge unterscheiden uns fürnehmlich von den Heiden und Unchristen: als nämlich Gottes Wort und ein heiliges Leben. Beides bestätiget der königliche Prophet David, da er sagt im 93. Psalm: Dein Wort ist eine rechte Lehre, Heiligkeit ist die Zierde deines Hauses ewiglich. Dies Wort Gottes haben wir, und das wird gepredigt und vorgelesen, daß wir darnach recht gläuben, christlich leben und geduldig leiden sollen in Lieb und Leid, Noth und Tod. Derowegen wollen wir Gottes Wort also anhören, lernen und gebrauchen, daß wir gute Werke draus machen; wir wollen auch in allerlei löblichen Tugenden uns üben, nach Erheischung eines jeden Beruf und Amt, durch Hülfe und Regierung des heiligen Geistes, daß wir uns als wahre Christen gegen Gott und Menschen verhalten. Darauf höret ferner an das 16. Kapitel.

\section{Die Wachteln und das Manna wird gegeben}

a Man) heißt auf hebräisch eine Gabe. Bedeutet, daß uns das Evangelium ohne unser Verdienst und Zuthun, aus lauter Gnaden vom Himmel gegeben wird, wie das Man auch gegeben ward.
\quellseite{147}

\section{Summarium}

Sie gehet an der Juden Historia, darinn man siehet, wie ein ungläubig und störrig Volk sie gewesen sind. In Egypten schrieen sie sehr über die Arbeit und unbillige Beschwerung. Sie wünschen sie, daß sie wieder drinnen wären, um eines Bissen Fleisches und Stück Brodes willen, V. 2. Und denken nicht einmal daran, daß sie einen solchen Gott haben, der sich schon bewiesen habe, daß er in allerlei Nöthen könne und wolle helfen. Ja sie vergessen all seiner Wunder und sagen, Mose und Aaron haben sie auf nichts Gutes aus Egypten geführet, V. 3.

Aus dieser Historie sollen wir unsere Art und Natur erkennen lernen, wie wir so gar verderbet sind; aber Gott wiederum so barmherzig ist, daß er gerne vergibt, und Fleisch und Brod den ungläubigen murrenden Juden die Fülle gibt.


b Zeugniß das ist an dem Ort, da man opfert und betet, und der Predigtstuhl war, ehe die Hütte war gemacht.


\section{Summarium}

Es ist ein fein Exempel, wie der Unglaube nichts ausrichtet mit dem, da er die Sache verhoffet zu fördern. Gott wollte, wie wir auch im Vater Unser beten, man sollte am täglichen Brod sich genügen lassen. Aber etliche wollten fürsammlen, denen verdarb das Man, V. 18. Etliche wollten am Sabbath sammeln, die fanden nichts, V. 27. Auf daß wir beides lernen, am ersten nicht geizig sein, sondern am täglichen Segen uns genügen lassen, und für das Morgende nicht sorgen.

Zum andern, daß wir uns die Nahrung am rechten Gottesdienst nicht lassen hindern. Wie Christus auch befiehlt: Suchet am ersten das Reich Gottes, so soll euch das andere alles zufallen, Matth. 6, 33.
\quellseite{148}

\section{Votum}

Helfe der treue GOtt, daß wir fleißig thun was wir sollen, daß er auch als ein lieber Vater treulich thue, was wir zu seinen Ehren und unserer Nothdurft begehren und wollen, hochgelobet in alle Ewigkeit, Amen.

\chapter{Das 17. Kapitel}

\section{Vorrede}

Der weise Mann Sirach vermahnet ein jedes unter uns zur Danksagung und spricht im 32. Kapitel: Danke für alles dem, der dich geschaffen und mit seinen Gütern gesättiget hat. Und St. Paulus vermahnet jedermann zum Gebet und sagt Ephes. 6.: Betet stets in allem Anliegen, mit Beten und Flehen im Geist; aber, ach wir müssens alle bekennen, undankbar sind wir, wenn uns Gott geholfen hat, und nachlässig sind wir im Gebet, wenn wir Gottes Hülfe bedürfen. Derowegen laßt uns solche Unart erkennen und davon abstehen: lasset uns Gott für seine Wohlthaten danken; denn er verdient es um uns alle Augenblick. Lasset uns auch fleißig beten; denn zu dieser Zeit Gottes Ehre und unsere höchste Nothdurft solches erfordert. Davon unterrichtet auch das 17. Kapitel.

\section[Die Kinder Israel werden aus einem Felsen getränket,]{Die Kinder Israel werden aus einem Felsen getränket, überwinden die Amalekiter}




\section{Summarium}

Es ist ein Exempel, wie die Welt neben allem, was sie von Gott weiß und gesehen hat, alsobald sie einen Mangel spüret, unwillig wird, Gottes sich nicht tröstet und schlecht nichts leiden will; ja alle Schuld auf die Prediger und ihre Obrigkeit schiebet. Daß Moses sich hie bezeugen muß, sie möchten ihn mit Steinen zu Tode werfen, V. 4. Aber Gott beweiset hie wiederum, wie er mehr nach seiner Barmherzigkeit, denn nach dem wir werth sein, mit uns handeln wolle.
\quellseite{149}
Zum andern ist ein fein Exempel hie, mit Moses Gebet, daß man des Sieges von Gott gewarten und durchs Gebet Gott suchen soll, V. 11. Daß man aber Mose seinen Arm unterstützt, V. 12., bedeutet, wenn das Gebet soll stark und rechtschaffen gehen, daß man Buße thun und mit dem Leben sich bessern soll, daß beides beisammen stehe, ein ernstliches Gebet und christliches Leben.

\section{Votum}

Der barmherzige gütige GOtt helfe, daß wir in täglicher Buße leben und fleißig beten, daß GOtt uns auch beistehe, und allen unsern Feinden steure und wehre, auch uns unsere Nothdurft aus Gnaden beschere. Amen.

\chapter{Das 18. Kapitel}

\section{Vorrede}

Wo in einer Stadt oder Lande die Kirchen mit treuen Predigern und das weltliche Regiment mit treuer Obrigkeit begnadet ist, da wohnet Gott und da ist Glück, Heil und Seligkeit. Denn christliche Prediger sind nützliche Amtleute im geistlichen Reiche Gottes. Und fromme Obrigkeiten sind heilsame Amtleute im weltlichen Reiche Gottes. Diese Kleinode hat auch unsere Stadt und Vaterland. Derowegen laßt uns Gott dafür danken und Gottes Wort, daß unsere Seelsorger lehren, und der Obrigkeit, die uns mit göttlichen Gesetzen regieret, also gehorsam sein, daß gemeldete Gnadenschätze bis auf den jüngsten Tag uns erhalten werden, Amen. Dies unterrichtet auch das 18. Kapitel.

\section{Jethro gibt Mose einen guten Rath}
\quellseite{150}

\section{Summarium}

Ist eine Historia, wie Moses das Regiment mit mehr Leuten bestellet und er fürnehmlich des Predigens und Gottesdienstes wartet, V. 24. Denn es ist nicht eines Mannes Arbeit, einem solchen Haufen fürzustehen.

Es ist aber hie sonderlich fein abgemalet, was für Leute zum Regiment zu brauchen seien, V. 21. Nämlich redliche Leute, die Ehre und Tugend lieb haben, und fest darob halten. Darnach gottesfürchtige Leute, die wahrhaftig und nicht geizig sind. Denn wo Gottesfurcht nicht ist, da wird wenig Glück sein. Und wer nicht wahrhaftig noch treu ist, da ist man übel mit versorget in Händeln. Was aber geizige und eigennützige Leute sind, die taugen gar nichts zum Regiment. Denn der Eigennutz blendet sie, daß sie auf den gemeinen Nutz nicht können sehen.

\section{Votum}

Der barmherzige Gott erhalte und segne unsere liebe Obrigkeit, daß sie auf Erden also regiere, damit sie mit dem Himmel-Könige Christo, um seines Verdienstes willen, regiere in Ewigkeit, Amen.

\chapter{Das 19. Kapitel}

\section{Vorrede}

Daß Gott die heiligen Zehen Gebote auf dem Berge Sinai wiederholet hat, ist so eine große Wohlfahrt, daß sie keines Menschen Zunge genugsam ausreden kann.

Derowegen danken wir heute Gott für solche Wohlfahrt; der wolle uns allen den heiligen Geist geben, daß wir durch seine Regierung nach den heiligen Zehen Geboten leben, die darin verbotene Sünde fliehen und die darin befohlenen guten Werke thun und üben. Denn, wie die Kriegsleute einen Bestallungs-Brief haben, darauf sie schwören: also sind unsere Bestallungs-Briefe die heiligen Zehen Gebote, welche in den Schriften der Propheten und Apostel erkläret werden; und daß wir nach diesem Bestallungs-Brief leben wollen, haben wir alle in unserer Taufe geschworen, und derowegen wer dawider handelt, der wird an Gott der heiligen Dreifaltigkeit meineidig. Dieß betrachtet auch bei dem 19. Kapitel.

\section{Vorbereitung, das Gesetz Gottes anzuhören}
\quellseite{151}

\section{Summarium}

Ist eine Historia, wie GOtt sich geoffenbaret habe, daß er das Gesetz auf dem Berge Sinai gegeben hat. Und reimet sich solche Offenbarung fein mit der Predigt des Gesetzes. Denn gleich wie das Donnern und Blitzen, wo es so gräulich ist, erschreckt und furchtsam macht, v. 16.: also wird auch ein menschlich Herz, wo es des Gesetzes Predigt recht trifft, erschrocken seiner Sünden und Gottes Zorn halben. Darum heißt es St. Paulus eine Predigt, dadurch die Sünde offenbaret, die Herzen erschreckt und getödtet werden.

Um dieser Ursache willen hat Gott seinen Sohn gesandt und unter das Gesetz gethan, daß er uns von solcher Last durch sein Genugthun erlösete, und durchs Evangelium wieder tröstete und lebendig machte. Wie denn das Neue Testament seine eigene und neue Pfingsten hat, darin Gott vom Himmel sich anders hat hören lassen, denn hie; nicht mit Donner und Blitz, sondern mit einem feinen sanften Wind und mit schwebenden Feuerzungen; dafür erschraken die Herzen nicht, wie am Berge Sinai, sondern wurden fröhlich und muthig, und preiseteten und lobeten Gott, unangesehen, was die Welt dazu that.

\section{Votum}

GOtt der Vater erleuchte und regiere unsere Herzen durch den heiligen Geist, daß wir nach seinen heiligen Zehn Geboten heilig leben, und um Jesu Christi willen an Leib und Seel gesegnet und selig werden, Amen.

\chapter{Das 20. Kapitel}

\section{Vorrede}

Unter den größten Dingen, die im menschlichen Geschlecht geschehen sind, ist wahrhaftig eines, die Offenbarung Gottes, da er mit himmlischer Stimme sein Gesetz oder die heiligen Zehn Gebote wiederholet hat auf dem Berg Sinai. Für solche Offenbarung sollen wir allezeit, und also heute abermal Gott danken; sollen auch Gottes Wort für wahrhaftig halten und Gottes Geboten folgen, mit Thun und Lassen, in allem, was darinnen geboten und verboten ist. Denn Gottes Gesetz ist ein Kerker, geordnet zu schützen die auswendige Zucht und Ehrbarkeit; und die Zehn Gebote sind wie zehn Ketten und Fesseln dadurch unsere
\quellseite{152}
Glieder gebunden werden, daß sie sich nicht muthwillig an Gott und dem Nächsten vergreifen. Darnach, Gottes Gebote sollen auch unser Kompaß, Richtschnur und Winkelmaß sein, darnach wir ein christlich und unsträflich Leben führen sollen. Darauf höret davon an das 20. Kapitel.

\section{Mose empfängt die heiligen zehn Gebote GOttes}

a Das ist, was du zu thun hast.

\section{Summarium}

Ist eine schöne Predigt, welche Gott selbst vom Himmel herunter gethan hat, in welcher er seinem Volk anzeiget, wie sie sich halten sollen, wenn sie nach seinem Willen leben wollen. B. 1. Aber hie müssen wir unsere verderbte Natur kennen lernen. Denn weil wir solchen Willen Gottes wissen, und doch denselben weder wollen noch können thun: daher könnens, daß wir gegen
\quellseite{153}
Gott nimmermehr kein gut Gewissen können haben, wir müssen uns für Sünder bekennen und Gottes Zorn, Ungnade und Strafe gewarten. Und ist hie kein ander Mittel, dadurch uns zu helfen ist, denn daß wir Vergebung durch Christum glauben. Sonderlich aber ist hie zu merken, daß Gott so fleißig verbeut, man soll keinen Gottesdienst fürnehmen ohne und außer seinem Wort und Befehl.

\section{Votum}

Lob, Ehr, Preis und Dank sei GOtt für seine heiligen Zehen Gebote; der lehre uns durch seinen heiligen Geist dieselben nicht allein mit dem Munde zu bekennen, sondern daß wir auch nach denselben allzeit gehorsamlich mit der That leben, Amen.

\chapter{Das 21. Kapitel}

\section{Vorrede}

An Unterricht, wie man gegen Gott und den Menschen leben solle, mangelts nicht; daran mangelts, daß wenig solchem Unterricht gemäß leben. Derowegen laßt uns Gott danken, daß er uns lehret, wie wir ihm mit Thun und Lassen schuldigen Gehorsam leisten sollen; und schaue ein jedes zu, daß es also lebe, auf daß es nicht als ein muthwilliger Uebertreter Gottes Gebots sterbe und verdammt werde. Denn sind doch die unvernünftigen Kreaturen Gott ihrem Schöpfer gehorsam: als auf Gottes Befehl stehen die Wasser und Sonne stille, der Himmel gibt Himmelsbrod und Wachteln. Der harte dürre Steinfels gibt Wasser. Der Wallfisch speiet Jonam aus. Das Feuer muß die drei Männer im feurigen Ofen nicht verbrennen. Die Löwen müssen Daniel nicht fressen. Sollten wir vernünftigen Kreaturen nicht viel mehr Gott gehorsam sein? Müßten wir doch härter denn die Felsen, wilder denn die Löwen, ja gräßlicher, denn der Teufel selbst sein, der auf Christi Befehl weichet und ihm gehorsam leistet. Darauf höret nun an das 21. Capitel.

\section{Ordnung vom Gesinde, Item, über Mörder und Todtschläger}
\quellseite{154}

\section{Summarium}

Es ist eine Ordnung der Gerichtshändel, so in Todtschlägen oder in andern Verletzungen bei den Juden haben müssen gehalten werden. Solche Ordnungen gehen uns weiter nichts an, denn sie den natürlichen Rechten gemäß sind, was außer demselben zum Regiment verordnet ist, das mag man nach jedes Landes und Ortes Gelegenheit setzen oder fallen lassen. Sonderlich aber ist es fein, daß Gott verbeut, es soll keine Freiheit noch Geleit die helfen, so fürstlich und mit List einen Todtschlag begehen, V. 14. Item, daß die Kinder, so Vater und Mutter entweder schlagen oder ihnen fluchen, sollen getödtet werden, V. 15. Solche Satzungen stimmen mit den natürlichen Rechten; derhalben sie sollen gehalten und solche unartige Leute nicht geduldet werden.

\section{Votum}

Gott steure den unmorrischen Teufeln und blutgierigen Menschen, und verleihe uns, daß wir nach allen seinen Geboten thun und lassen, was er uns geboten und verboten hat, Amen.

\chapter{Das 22. Kapitel}

\section{Vorrede}

Wir sollen nicht allein Hörer, sondern auch Thäter sein des Wortes Gottes. Denn wir unserm lieben Gott in der heiligen Taufe Gehorsam gelobt und geschworen haben. So ist er auch unser zu Leib und Seel, Ehr und Gut mächtig, und wir ohne ihn nicht einen Finger regen, oder einen Odem holen, viel weniger größere Thaten üben können. So daß auch reichlich und wohl um uns Gott mit viel und großen Wohlthaten verdienet, daß wir ihm zur Dankbarkeit Gehorsam leisten. Darwegen helfe Gott, der heilige Geist, uns allen, daß wir also Gottes Wort anhören, daß wir nach seinen Geboten christlich leben auf Erden, damit wir nach seinen gnädigen Verheißungen leben in Seligkeit im Himmel. Darauf hört ferner an das 22. Kapitel.

\section{Strafe des Diebstahls und anderer Sünden}
\quellseite{155}

a Götter) heißen die Richter, darum, daß sie an Gottes Statt nach Gottes Gesetz und Wort, nicht nach Eigendünkel, richten und regieren mußten; wie Christus zeuget.

b Dringen) dringen und wuchern sind zweierlei: dringen ist, wenn du deinem Nächsten zwingest, zu bezahlen mit deinem Vortheil und seinem Nachtheil. Wuchern weiß man wohl, was es sei.

c Fülle) heißet er alle harte Früchte, als da sind Korn, Gerste, Aepfel, Birnen, da man Speise von macht.

d Thränen) heißet er alle weiche Früchte, da man Saft und Tran von macht, als da sind Weintrauben, Oele.
\quellseite{156}

\section{Summarium}

Es ist eine Ordnung der Gerichtshändel, so sich zutragen im siebenten Gebot, mit Diebstahl. Darnach sind viel seiner bürgerlicher Satzungen, daß man die Fremdlinge nicht übersehen, V. 22., Wittwen und Waisen handhaben, armen Leuten ohne Wucher leihen und ihr nöthiges Pfand oder Hausrath, daß sie zur Nahrung nicht gerathen können, über Nacht nicht behalten soll, V. 25. Denn mit solchem Leihen würde ihnen übel gedienet. Item, daß man der Obrigkeit nicht fluchen, noch sie lästern, sondern sie fürchten und ehren soll.

Item, daß man den Gottesdienst und die Personen, so dazu verordnet, mit dem, daß jedermann Gott bescheret hat, gern fördern und erhalten, die Gotteslästerer aber mit dem Tode strafen soll, V. 28.

\section{Votum}

Gott versiegele sein Wort in unsern Herzen, und helfe, daß wir uns gegen ihn kindlich, und gegen unsern Nächsten brüderlich allzeit verhalten, in Kraft des heiligen Geistes, Amen.

\chapter{Das 23. Kapitel}

\section{Vorrede}

Die zehen Gebote Gottes sollen aller Menschen täglicher Sündenspiegel und Laftafel sein, also, daß ein jeder seine Sünde daraus erkennen lerne, und die Gott abbitte, in wahrer Reue und Glauben. Und weil die Gebote Gottes lehren, was man dagegen Gutes thun soll, daß vom Bösen bewegen ein jedes lasse, und in guten Werke sich übe. Daher singen wir: Das Gesetz darum ist gegeben, daß es sei ein Spiegel zart, darinnen wir sehen unsre sündige Art, die in unserm Fleisch ist verborgen. Und: Die Gebote all uns gegeben sind, daß ein jedes soll lernen wohl, wie man für Gott leben soll. So verleihe Gott, daß wir alle seiner Gebote also täglich recht und christlich gebrauchen, Amen. Davon predigt auch das 23. Kapitel.

\section{Von Festen und Feiertagen}

a Versäume gern] Id est, relinquo mea interim et damnum accipio. Das ist: ich versäume indessen das Meine und leide Schaden.

b (Gedenken) ihr sollt von keinen Heiligen predigen, als von Göttern, die euch helfen, noch dafür danken, sondern Gott allein. Denn gedenken heißt so viel als predigen, rühmen, danken, Gottesdienst pflegen, ut: hoc tacite in mei commemorationem. 3. B.: daß thut zu meinem Gedächtnis.

c Dreimal) das ist das Osterfest im April, Pfingsten im Brachmonat, und das Laubhütten-Fest im Weinmonat. Davon lies im 23. Kapitel des dritten Buches Mose.
\quellseite{157}

d Des Jahres Ausgang heißt er den Weinmonat, daß alsdann aus ist mit Frucht wachsen und sammeln.

e Das Blut rc.) das ist, du sollst das Opferlamm nicht schlachten noch opfern, ehe denn alles gesäuerte Brod aus deinem Hause kommt.

\section{Summarium}

Ist zum Theil eine Auslegung des achten Gebotes, daß man wahrhaftig sei, recht richte und mit Geld und Gut sich nicht überwiegen lasse. Darnach ordnet er die drei fürnehmsten Feste im Jahr, B. 14., da alle Juden zum Tabernakel und Tempel kommen mußten; und sonderlich gebet Gott, man soll vor ihm nicht leer erscheinen, B. 16., das ist, man soll zur Kirche und zur Erhaltung des rechten Gottesdienstes gern geben, und nach Vermögen dazu helfen.

\section{Summarium}

Ist eine schöne Vermahnung, damit er die Predigt von den zehen Geboten beschließt. Auf daß, weil Gott durch sein Wort seinen Willen hat offenbaret, wir demselbigen sollen folgen. Thun wir es, so wolle er Brod und Wasser, das ist, Essen und Trinken segnen, alle Krankheit abwenden, gesundes Alter verleihen und allenthalben segnen, B. 25, 26. Thun wir es aber nicht, so soll die Strafe auch nicht außen bleiben.

Da stehet man nun, was es für Ursache sei, daß so viel arme, kranke, unglückselige Leute auf Erden sind: nämlich, daß der meist Theil solches durch eigene Sünde und gottlos Leben verdienet. Darum, wer solchem Unglück entlaufen, und wie der 34. Psalm B. 13. sagt, gute Tage haben will, der fürchte den HErrn, thue Gutes und fliehe das Böse, so will Gott wiederum thun, was er soll und wir begehren.

\section{Votum}

Der treue, barmherzige GOtt helfe, daß wir ja seinem Willen nicht widerstreben, sondern nach demselben allzeit in Gehorsam christlich leben und selig sterben, Amen.
\quellseite{158}

\chapter{Das 24. Capitel}

\section{Vorrede}

Gleich wie bei dem Berge Sinai stund sechsmal hundert tausend Mann und hörten die Gesetzpredigt Gottes, zitterten und bebten, gelobten auch an, sie wollten Gott gern gehorsam sein in allem; aber es ward darnach bald vergessen, und geriethen in allerlei Sünden, daß nur zween kamen in das gelobte Land: also gehet es stets. Nicht alle, die getauft sind, nicht alle, die Gottes Wort hören, nicht alle, die da gebrauchen das Sacrament des wahren Leibes und Blutes des HErrn Christi, auch nicht alle, die sich des heiligen Geistes rühmen, erlangen das ewige Leben. Denn viele ihr Fleisch und Blut sich auch verführen lassen, leben und sterben in Sünden. Derowegen sollen wir zusehen, wie wir fürsichtiglich wandeln, und sollen verständig sein, was da sei des HErrn Wille, und sollen zugesagten Gehorsam Gott, nach seinem Wort, treulich leisten. Davon höret auch an das 24. Capitel.

\section[Moses steiget, nach Bestätigung des Bundes, wieder auf]{Moses steiget, nach Bestätigung des Bundes, wieder auf den Berg Sinai}

a Seine Hand) er schreckte sie nicht mit Donner und Blitz, wie zuvor das Volk erschrecket ward, Kap. 20, 18.

\section{Summarium}

Ist eine Historia, wie das Volk sich verpflicht-, Gottes Gesetze und Willen zu halten, V. 4. Aber wie solches erfüllet sei, werden wir in der Historia weiter hören; daß auch St. Petrus zeuget in den Geschichten der Apostel Kap. 15, 10.: Das Gesetz sei ein Joch, welches weder die Apostel noch die Väter gehalten haben. Darum spricht die Epistel an die Ebräer: Der Bund, so durch der Farren Blut gemacht, sei nicht kräftig gewesen. Und hat eben solcher Ursache halben Gott einen andern Bund aufgerichtet, der nicht auf äußerlichem Thun und Lassen, sondern auf Gottes Gnade beruhen soll, durch das Blut unsers Seligmachers Jesu Christi, wenn wir an ihn glauben und seines Verdienstes uns trösten.
\quellseite{159}

\section{Votum}

Gott der heilige Geist erleuchte und regiere uns, daß wir auch thun, was wir Gott zugesagt haben, und daß wir alle Sünde hassen und Tugenden üben, auf daß wir zeitlicher und ewiger Strafe entfliehen, Amen.

\chapter{Das 25. Kapitel}

\section[Inhalt der folgenden sieben Kapitel: Von dem Heiligtum]{Inhalt der folgenden sieben Kapitel: Von dem Heiligtum und was mehr dazu gehöret hat}

\section{Vorrede}

Diese folgenden sieben Kapitel lehren anders nichts, denn wie Gott Mose auf dem Berge habe Befehl gegeben, von sonderlichem Gottesdienst aufzurichten, der unter den Juden bis auf Christus Reich sollte gehalten werden.

Denn hie ordnet er erstlich, wie man die Hütte oder Kirche bauen, V. 8., die Lade des Bundes, V. 10., samt dem Gnadenstuhl, V. 17., und Cherubim zurichten, V. 18., den güldenen Tisch mit den Schaubroden, V. 23., den güldenen Leuchter, V. 31., den Vorhang samt anderem Geräthe, das man dazu bedurfte, machen soll.

Darnach ordnet er die Priester, wie sie gekleidet sein, was und wenn sie opfern sollen, auch wie man sie weihen soll. Solches alles, ob es jetzt schon nicht mehr im Brauch ist, ist doch derhalben sein zu wissen, daß man unsern Papisten desto baß begegnen kann, welche auch ihre Kirchen, Pfaffen, Kelch, Meßgewand und andere Kirchengefäße weihen; aber doch mit einer weit andern Meinung, denn Moses.

Moses hatte Befehl von Gott, er sollte es so machen und anders nicht; darum weiset ihm auch Gott auf dem Berg ein Muster, da er alles vor ihm stehen sahe, wie es sein sollte.

Unsere Papisten aber haben außer und ohne Gottes Befehl Pfaffenweihe, Kirchenweihe und andere Gerätheweihe aus ihrem eigenen Kopf erfunden und fürgenommen, können auch keine andere Ursache anzeigen des Papstthums halben, daß der Papst soll das oberste Haupt der Kirchen sein; item der Pfaffen und Kirchenweihe etc., denn daß im Alten Testament also gewesen sei. Wo es nun im Neuen Testament nicht auch sollte sein, sagen sie, das Alte Testament wäre herrlicher, denn das Neue. Aber wir wissen eine andere Herrlichkeit, nämlich, daß der Sohn Gottes oberster Priester ist, und mit seinem Opfer uns versöhnet hat, der weihe seine Kirche mit dem heiligen Geist und gibt tüchtige Kirchendiener. Derohalben soll solches päpstliches Wesen durchaus von niemand anders gedacht werden, denn wie andere Menschenkünste und Ordnung, so ohne alles Gewissen mögen, ja sollen geändert werden; sonderlich, wenn man sie für nöthig achten will, wie der Papst thut. Dies ist das erste, das man bei diesen Kapiteln soll merken.

Das andere ist, daß Gott allen solchen Unkost allgemein auf das Volk schlägt, daß sie zusammen schließen und ein jeder das Seine geben soll, daß es also möge gebauet und angerichtet werden, V. 20. Uns zur Lehre, daß wir auch gerne dazu geben und helfen sollen, was zur Förderung des rechten Gottesdienstes und Kirchen-Ämtern vonnöthen ist.

Das dritte ist die heimliche Deutung mit solchem Bau und Ordnung. Die Hütte ist die christliche Kirche, in derselben liegt das Zeugniss, das ist, die heilige Schrift und das Wort Gottes.

Der Gnadenstuhl darin, das ist unser lieber Herr Christus Jesus, durch welches Verdienst wir mit Gott versöhnet sind. Derselbe ist auch der Hohepriester, welcher sich geopfert hat für unsere Sünde, und ist von seinem himmlischen Vater gesalbet oder gewehet mit seinem heiligen Geist, welchen er ausgegusst durchs Evangelium, daß wir auch gesalbte Priester werden und Gott dienen durch den rechten Glauben, mit der Predigt von seiner Güte und Gnade, mit reinem Gebet, mit Opfern unser selbst und mit täglichem Töden des alten Adams.

Solches alles sind die sonderlichen herrlichen Gaben des Neuen Testaments, welches Gott den Juden durch solchen Brauch und Ordnung, gleich als durch ein Fürbild hat anzeigen wollen; wie die Epistel an die Ebräer auch lehret. Dieses sei zu einer einfältigen Anleitung von diesen Kapiteln genug gesagt.

\section{Freiwillige Opfer zur Stiftshütte}

a Fürbild) das ist der Hauptstrich einer, darin Mose zeuget, daß sein Regiment soll aufhören und nicht das rechte ewige Wesen, sondern ein Fürbild und Beispiel sein des Reiches Christi. Und ohne allen Zweifel hat hieran Christus mit den zwei Jüngern zu Emmaus geredet, Luc. 24, 27.
\quellseite{160}

b Dir zeugen) das ist, dabei, als bei einem gewissen Zeichen und Zeugniß, will ich dich wissen lassen, daß ich da bin gegenwärtig, daß ich daselbst reden werde.

c Centner) dreißig Pfund Goldes.

d Und siehe zu) er wiederholts hie zum andern Mal, daß er droben in diesem Capitel sagt: Es soll alles viel Fürbild und Figur sein mit Mose, auf daß es, zweimal gesagt, ja gewißlich geglaubt werde.

\chapter{Das 26. Kapitel}

\section{Bau der Stiftshütte}
% Omitted: numbered Bible verse, excluded; printed folio or catchword marker; running head; not text apparatus.
\quellseite{162}

\chapter{Das 27. Kapitel}

\section{Vom Brandopferaltar, Vorhof und heiligen Oel}

a Hörner) aufgerichtete kleine Säulen, oben mit Knöpfen oder Blumen.

b Moed) das hebräische Wort Moed haben wir nicht anders wissen noch wollen deuten. Es soll aber so viel heißen als ein gewisser Ort oder Stätte, wie eine Pfarrkirche oder Stift, dahin das Volk Israel kommen und Gottes Wort hören sollten, damit sie nicht ihrer eigenen Andacht nach hin und wieder liefen, auf Bergen, in Gründen und andern Oertern, Gott zu opfern.

\chapter{Das 28. Kapitel}

\section{Priesterliche Kleiderzierde Aarons und seiner Söhne}
\quellseite{163}

a Licht) was das gewesen sey, weiß man jetzt nicht mehr. Hebräisch heißt es Urim und Thummim, 4 Mos. 27, 21.
\quellseite{164}

b Füllen) dies Füllen ist eine hebräische Sprache, der man muß gewöhnen. Und war das, wie im folgenden Capitel stehet, daß in der Weihe den Priestern die Hände mit Opfer gefüllet wurden für dem HErrn.

\chapter{Das 29. Kapitel}

\section{Einweihung und Opfer der Priester}
\quellseite{165}

a (Ein Feuer des HErrn) das Wörtlein „Feuer des HErrn“ brauchet Mose sehr viel, und wir sollten's gewöhnen, daß es so viel heiße, als ein Opfer des HErrn. Als wenn du von ferne ein Feuer sähest, und fragest, was ist das? So spräche man, es ist Gottes Feuer, das ist, man opfert dasselb Gott. Per synecdochen vel aliam figuram.

b Entsündigen) das ist, absolviren und lossprechen, wie Ps. 51, 9. Asperges me Isopo, das ist: Entsündige und absolviere mich mit Isop.

c Hin) ist fast eine halbe Metze.
\quellseite{166}

\chapter{Das 30. Capitel}

\section[Vom Räuchaltar, der aufgesetzten Steuer, Handfaß, Salböl und]{Vom Räuchaltar, der aufgesetzten Steuer, Handfaß, Salböl und Räuchwerk}

a Dieser Altar ist außen für dem Fürhang im Heiligen, und nicht hinter dem Fürhang im Allerheiligsten gestanden.

b Lampen zurichtet) Benefacit, depurando Evangelium ab opinionibus, daß es sein brennen kann. Das ist, wenn man das Evangelium von Menschenfälschungen reinigte.

c (Gera) gilt bei uns drei Löwenpfennige, oder ein Dreigröschlein. Ein halber Seckel ist ein Ort eines Gülden.
\quellseite{167}

\chapter{Das 31. Kapitel}

\section{Bestellung der Werkmeister. Sabbathfeier. Gesetztafeln}
\quellseite{168}

\chapter{Das 32. Kapitel}

\section{Vorrede}

Gott belohnet das Gute und strafet das Böse; und auf große Sünden folgen große Strafen, und solches ist die Belohnung der schändlichen Sicherheit. Das lehret Gott, da er Ezech. 8. also sagt: Ich will wider sie mit Grimm handeln, mein Auge soll ihrer nicht verschonen und will nicht gnädig sein. Und abermal Malachi am 3. Kapitel also: Ich will ein schneller Zeuge sein wider die, so mich nicht fürchten, denn ich bin der HErr, der nicht leuget, spricht der HErr Zebaoth. So ist dies auch wahr: Wie ein Bock härter stößt, wenn er weit aushole[t]: also strafet Gott desto schwerer und heftiger, wenn er die Strafe aufschiebt. Derowegen, weil ohne das Jammer und Noth genug ist in allen Winkeln der Welt, in allen Ständen, so laßt uns Sicherheit fliehen, beten und gottesfürchtig leben, auf daß wir mit unsern Sünden Gott nicht mehr erzürnen und zu Strafen weiter Ursache geben. Darauf höret auch an das 32. Kapitel.

\section{Abgötterei mit dem gegossenen Kalbe}

a Verzog das redet Moses mit einem Wort also, als habe er das Volk verlassen mit Schanden, und sey gestorben, daß sie nicht wissen, wie sie nun thun sollen. Gleich als wenn uns Gott eine kleine Zeit verläßt, denken wir, er lasse uns in Schanden stecken, müssen anderwärts Hülfe suchen; da wird denn solch Kalb unser Gott.

b Entwarf das ist, er machte es ihnen für, was sie für ein Bild machen sollten. Das bedeutet, daß Menschenlehre dem Volke vorbilden, was sie für Werke thun sollen, damit sie Gott dienen. Denn hie siehst du, daß sie in diesem Kalbe vermahnet haben dem rechten Gott zu dienen, weil Aaron rufen läßt: Es sei des HErrn Fest, und bauet ihm einen Altar.

\section{Summarium}

Dies ist eine erschreckliche Historia, wie das Volk Israel, welches so große Wunder in Egypten und in der Wüste gesehen und erfahren, und sich gegen Gott, seine Gebote zu halten, so theuer verpflichtet hatte, so bald alles vergessen und in greuliche Abgötterei gefallen ist und ein gegossen Kalb angebetet hat, V. 4.
\quellseite{169}
Uns zum Exempel, wie es Paulus selbst deutet, 1 Cor. 10, 6., daß wir nicht sicher noch vergessen sein, sondern in Gottes Furcht stehen sollen und lernen, daß es, so bald Gott die Hand abzieht, mit uns aus ist. Darum ist vonnöthen, daß wir stets bitten, daß Gott durch seinen Geist und Wort uns erhalten, und nicht in Anfechtung führen wolle.

Darnach ist ein fein Exempel an Mose, wie er für das Volk bittet und das Gebet allein auf Gottes Güte und Verheißung stellet, und Gott solches annimmt, V. 11.

c Das Wort Phrea lautet und heißt schier dem Deutschen gleich, frei; will so viel sagen: Aaron hatte das Volk frei, los und ledig gemacht von Gottes Wort und Gehorsam, wie die Menschen-Lehre thut, machet das Volk wild, frei, los und bereit zu aller Abgötterei, und thut doch der Meinung, als wollten sie den Leuten aufhelfen, und besser machen. Und ist doch (wie es im Hebräischen heißt) ein Schemez à Schama, ein Gehör oder Geschwätz, da nichts hinter ist.

d Nemo evadet poenam peccati sui. Das ist, Niemand wird der Strafe seiner Sünde entgehen.

\section{Summarium}

Ist eine Historia, wie Moses um des Aergernisses willen heißt die Kinder Levi erwürgen und erstechen, was sie antreffen, V. 27. Auf daß er dem Zorn und der Strafe Gottes fürkomme, und den größten Haufen errette. Denn solches gehöret der Obrigkeit zu, daß sie öffentliche Aergerniß strafe; auf daß Gott nicht zu größerem Zorn beweget werde und den ganzen
\quellseite{170}
hausen darum angreife, so man öffentliche Aergernisse ungestraft läßt. Das aber Moses die steinernen Tafeln zerbricht, V. 19. ist eine Deutung, daß menschliche Natur das Gesetz Gottes nicht vermag zu halten, es komme denn Gott, und helfe nicht allein mit seinem heiligen Geiste, sondern auch mit Vergebung der Sünden, um Christi willen.

\section{Votum}

GOtt sei uns allen gnädig und barmherzig, und helfe, daß wir für Abgötterei und allen andern Sünden uns hüten, und nach seinen heiligen Geboten gehorsamlich bis zu unserm seligen Todesstündlein leben, Amen.

\chapter{Das 33. Kapitel}

\section{Vorrede}

Kein recht Gebet ist verloren. Denn Gott erhöret alle unsere Seufzer und schreibet sie auf; und wenn wir nur anfangen zu schreien: Vater, so gibts in dem Himmel einen mächtigen starken Schall: Sieh bin ich mein Sohn, meine Tochter, was willst du? Derhalben ist ein christlich Gebet gar sehr kräftig, mächtig und geschäftig, und es dringet durch alle Himmel, und richtet aus drüber und nicht drunter, denn wir je gebeten haben. Daher das gläubige Gebet wohl mag genennet werden der edelste Schatz der Christenheit, die beste Wehre und der höchste Trost, und die bewährteste Arznei in Nöthen, und der schönste Gottesdienst. So laßt uns alle Tage oft und fleißig im Geist und in der Wahrheit beten, Gott anrufen und danken. Dies erinnert auch das 33. Kapitel.

\section[Mose bittet für das Volk, und begehret des]{Mose bittet für das Volk, und begehret des Herrn Herrlichkeit zu sehen}

a Den Rücken Mose sehen alle Werckheilige, die das Gesetz nicht verstehen, noch unter Augen kommen.

b Angesicht, heißt die Wolke und feurige Säule, darin Gott gegenwärtig war.
\quellseite{171}

c Mein Angesicht) also heißt Gottes Angesicht nicht die Wolke noch Säule, sondern er selbst, wie er spricht: Der Mensch wird nicht leben, so mich siehet. Aber es ist alles von Christo gesagt, welcher sollte erscheinen in der Menschheit hernach, wenn Moses Reich ein Ende hätte. Indes soll Moses im Fels stehen und den vorbeygehenden Christum sehen und predigen, bis er komme; daß also Israel Gottes Wort immer habe, bis auf Christum.

\section{Summarium}

Ist eine feine Historia, wie Moses nirgend hinziehen will, es ziehe denn Gott mit, B. 15. 16. Solches soll eine christliche Obrigkeit allzeit thun, daß sie durch's Gebet Gott zu einem Vorgeher, Geleitsmann und Gehülfen nehme; sonst wird wenig Glück dabei sein, wo Gott mit uns nicht dabei ist.

Daß aber Moses begehren, Gottes Herrlichkeit zu sehen, und Gott sich solches weigert, B. 18. 19., ist uns zur Lehre fürgeschrieben, daß wir uns an der Offenbarung, so Gott in Christo und durch sein Wort uns hat widerfahren lassen, sollen genügen lassen. Denn außer diesem Anblick der Gnaden, welchen wir in Christo und im Wort haben, kann man Gottes Herrlichkeit nicht sehen und leben. Unsere Natur kanns nicht erdulden noch tragen, bis in jenem Leben, da das Fleisch nicht mehr wird sterblich sein, da werden wir Gott von Angesicht zu Angesicht sehen.

\section{Votum}

Der treue, barmherzige Gott sei und stehe auch bei uns in diesen gefährlichen Zeiten, blicke uns an mit Gnaden unter unserm Kreuz, und helfe, daß wir seine Herrlichkeit sehen in alle Ewigkeit, Amen.

\chapter{Das 34. Kapitel}

\section{Vorrede}

Ein jeder Mensch soll täglich gedenken an Gottes Gerechtigkeit und Barmherzigkeit. An Gottes Gerechtigkeit sollen wir gedenken, auf daß wir uns für Sünden hüten; an seine Barmherzigkeit aber, auf daß wir wider die Sünde Trost haben. Denn wie Gott, als ein eifriger Gott, die Sünde an allen Unbußfertigen strafet: also vergibt er die Sünde, als ein barmherziger Gott, allen, die ihn darum bitten. Von diesen beiden lehret also die heilige Schrift: Jeremias sagt im 5. Kapitel: Eure Mißthat und Sünden wenden das Gute von euch. Und David im 103. Psalm spricht: Wie sich ein Vater über Kinder erbarmet, so erbarmet sich der HErr über die, so ihn fürchten. Darauf höret weiter an das 34. Kapitel.

\section{Neue Gesetztafeln. Bund Gottes. Moses glänzendes Angesicht}

a Predigt) Mose berichtet hie das Geheimniß, daß Christus der HErr sey, der für ihn (das ist, für den Volk des Geistes) werde übergeben und predigen, daß alle Welt sündig, allein aus Gnaden selig werden müsse, wie er lehret Kap. 33, 19. sagt: Wem ich gnädig bin, dem bin ich gnädig, das ist, es soll mirs keiner mit Werken abverdienen, es soll Gnade sein, und kein ander. Wie denn des Evangelii Predigt gleich also lehret.
\quellseite{172}

b Halsstarrig Volk) Opera nostra non faciunt, si misericordia non comitetur. Das ist: unsere Werke werden nicht thun, wenn uns nicht Barmherzigkeit begleitet. Wir sind wohl so böse.

\section{Summarium}

Ist eine Predigt, in welcher die zehn Gebote wiederholet, sammt etlichen andern Geboten von den fürnehmsten Festen. Und ist hie sonderlich zu merken, wie das erste Gebot wird heraus gestrichen, daß Gott barmherzig, gnädig, geduldig, und von großer Gnade und Treue sei, die Sünde vergeben wolle, weil wir doch für ihn nicht können unschuldig sein. Darum soll man sich solches Gottes allein trösten und auf ihn trauen. Wiederum ist er ein eifriger Gott, der Abgöt%
terei,\quellseite{173} Mißglauben und gottlos Wesen nicht leiden, sondern bis ins dritte und vierte Glied strafen will, V. 7. Darum soll man solchen Gott allein fürchten und sein Wort und Willen halten.

\section{Summarium}

Ist eine Historia, wie Moses, da er vom Berge gegangen und das Gesetz gebracht hat, unter seinen Augen geleuchtet habe, daß ihn niemand hat ansehen können. Wie denn solcher Ursache halben die Maler ihn mit Hörnern malen. Es sind aber nicht Hörner, sondern ein Glanz gewesen, der von seinem Angesicht gegangen und den Menschen ist unleidlich gewesen, derhalben er eine Decke hat müssen fürhängen.

Solches ist eine seine Deutung, wie es St. Paulus 2 Cor. 3. angzeigt, daß man das Gesetz nicht recht verstehen noch thun kann, wo der Geist Gottes nicht recht im Herzen ist, welcher allein durch Christum uns geschenkt wird.

Denn gleich wie die Juden nur den Fürhang für Moses Angesicht haben gesehen und den Glanz nicht leiden noch sehen können: also sind alle die, so sich unterstehen das Gesetz zu thun aus eigenem Vermögen. Die Werke mag man wohl nachthun dem äußerlichen Wesen nach; aber das Herz bleibt unrein, einen Weg wie den andern. Wo aber Gottes Geist durchs Evangelium ist, da werden die Herzen geändert und das Gesetz recht erfüllet. Denn ob schon dem Werken etwas fehlet, wird es doch durch Vergebung der Sünde in Christo erstattet.

\section{Votum}

GOtt gebe und erhalte in uns allen seinen heiligen Geist, daß wir durch ihn die heiligen zehn Gebote recht verstehen lernen, und daß wir nach denselben allzeit in Heiligkeit und Gerechtigkeit leben, Amen.

\chapter{Das 35. Kapitel}

\section{Vorrede}

Wie kein Mensch nimmer aus der Kirche gehen soll, er habe denn zuvor gebetet, und etwas Nützliches aus Gottes Wort gelernt: also welchen Gott gesegnet hat, wenn er aus der Kirche geht, der soll sein Almosen im Kasten Gottes auch einlegen, den Hospitalen und hausarmen Leuten zur Beisteuer. Wohl denen, die solches thun aus wahrem Glauben. Denn Gott wird an ihrem Leib, Seele und Nahrung zeitlich bleiben und im künftigen Leben aus Gnaden mit ewiger Glorie, Freude und Herrlichkeit vergelten. Davon lehret der heilige Prophet Jesaias im 58. Kapitel, da er spricht: Entzeuch dich nicht von deinem Fleisch (das ist, verlasse deinen Nächsten nicht mit Almosen und anderer nöthiger Hülfe): alsdann wird dein Licht hervor brechen, wie die Morgenröthe, und deine Besserung wird schnell wachsen, und deine Gerechtigkeit wird vor dir hergehen, und die Herrlichkeit des HErrn wird dich zu sich nehmen. Dies lehret auch das 35. Kapitel.

\section{Sabbathsruhe. Freiwillige Steuer. Beruf der Werkmeister}

a Was Gott droben vom 25. Kapitel an bis aufs 32. Mose befohlen hat von der Wohnung und alle ihrem Geräthe zu machen, das zeiget Mose dem Volk an in diesem Kapitel.
\quellseite{174}

\section{Summarium}

Dies ist ein schön Exempel, in welchem wir sehen, wie die Juden so reichlich und gern Gold und Silber, und was sie haben, dazu geben, daß die Hütte oder Kirche gebaut und der rechte Gottesdienst gefördert und angerichtet würde, V. 21. Und zur Lehre, daß wir mit dem, so uns Gott gegeben hat, auch gern sollen helfen, auf daß das Predigtamt, die Schulen, sammt anderem, so zur Unterhaltung des Kirchendienstes vonnöthen, möge erhalten werden. Denn wer solches thut, der bauet recht an Gottes Tempel und leget sein Geld also an, daß Gottes Ehre und der Menschen Seligkeit dadurch gefördert werde.

Vom 10. Vers bis auf den 20. werden nur
\quellseite{175}
diese Worte gelesen: Und wer unter euch verständig ist, der komme, und mache, was der HErr geboten hat.

\section{Votum}

Helffe der treue barmherzige Gott, daß all unser Thun und Lassen zu seinem Ehren, und unsers Nächsten Nutz und Wohlfahrt gelange, und daß ein jedes mit Gebet und Hülfe, zur Erhaltung der wahren Gottesdienste, treulich diene, Amen.

\chapter{Das 36. Kapitel}

\section{Vorrede}

Diese vier Kapitel ist nicht vorzunehmen, daß mans von Wort zu Wort lese. Denn es ist nur eine Erzählung, wie man die Kirche sammt aller Zugehörung gebauet habe. Da können wir keine andere Lehre aus nehmen, denn diese: daß alles, was zur Kirchen und Gottesdienst im Alten Testament gebräucht hat, allein nach Gottes Wort, wie es Gott befohlen hat, ist gemacht worden. Denn eigen erwählter Gottesdienst außer und ohne Gottes Wort kann und will Gott nicht leiden. Darum gehet alle Klage in den Propheten darüber, daß das Volk auf Gottes Befehl nicht thun, sondern ein jeder ihm selbst ordnen und fürnehmen will, wie er Gott dienen möge.

Gleich wie unter dem Papstthum auch geschehen ist, da sind die rechten Gottesdienste, Glaube, Liebe, Beten, im Kreuz und Leiden geduldig sein, dahinten blieben, deren hat niemand gedacht. Aber Mönche und Nonnen werden, sonderlich kleiden, fasten, zu heiligen Wallfahrten gehen und dergleichen, davon Gott nicht ein Wort befohlen hat, dieß hat man für sonderliche Gottesdienste gehalten und fürgenommen; so doch ein jeder Hausvater in seinem Haus das für keinen Dienst achtet, noch ihm gefallen läßt, was sein Knecht ohne Befehl aus eigenem Willen fürnimmt; sondern das heißt dem Hausherrn gedienet, so man thut, was er befohlen hat.

Darum beschreibet hie Moses nach der Länge wiederum, auf daß seine Juden und wir lernen, wie er alles, was zu machen gewesen ist, nach Gottes Wort und Befehl hat machen lassen; auf daß wir in allen Dingen auch auf Gottes Wort sehen, und demselben Wort nach alles handeln und anrichten. Das ist die fürnehmste Lehre dieser vier Kapitel.

Darnach ist zu merken, wie Moses hie anzeiget, daß es eine Gabe Gottes sei, wenn einer ein Handwerk wohl kann, V. 1. Darum soll man es erkennen, und Gott dafür danken, und solcher Geschicklichkeit nicht mißbrauchen, wie doch gemeiniglich geschiehet. Denn wo sonderbare Gaben sind, da denket niemand, daß er zu Gottes Reich und des Nächsten Noth mit helfen soll; sein selbst Ehre, Geld, Gut und anderes suchet schier jedermann damit. Es versündigen aber die Leute sich also dadurch, daß Gott weder Glück noch Segen dazu geben will. Wie man stehet an vielen Künstlern, daß sie gemeiniglich arme Bettler sind, nicht der Kunst halben, sondern daß sie ihr nicht recht brauchen.

Zum dritten wird auch hie wiederholet, wie das Volk so gar reichlich habe zugetragen, mit Gold, Silber und anderem, daß Moses verbieten muß und heißen aufhören, V. 6. Uns zum Exempel, daß wir auch gern und reichlich zum Gottesdienst geben, und zu Unterhaltung der Kirchen willig helfen sollen.

\section{Die Stücke des Heiligthums werden verfertiget}
\quellseite{176}

\chapter{Das 37. Kapitel}

\section{Von etlichem Geräthe des Heiligthums}
\quellseite{177}

\chapter{Das 38. Kapitel}

\section{Ansehnliche Summe des angewandten Goldes, Silbers und Erzes}

a Weibern) diese Weiber waren die andächtigsten Wittwen und Weiber, die mit Fasten und Beten für der Hütte Gott dieneten, wie 1 Kön. 2, 22. gezeuget. Und Paulus 1 Tim. 5, 5. beschreibt. Wie auch St. Lucas die heilige Prophetin Hanna rühmet, Luc. 2, 37.
\quellseite{178}

\chapter{Das 39. Kapitel}

\section{Priesterliche Kleiderzierde}
% Omitted: numbered Bible verse, excluded; printed folio or catchword marker; running head; not text apparatus.
\quellseite{180}

\chapter{Das 40. Kapitel}

\section{Stiftshütte aufgerichtet und eingeweihet}
\quellseite{181}

\section{Summarium}

Es ist eine Historia, wie Moses alles mit der Kirchen, und was dazu gehöret, angerichtet habe. Und ist solches, wie oben auch gemeldet, sein zu wissen, um des Anspiels willen, so der Papst und Bischöfe ohne Gottes Befehl mit Kirchen, Kleidern und Priesterweihen, auch mit Glockentaufen und andern angerichtet haben. Denn hie wiederholets der Text achtmal, da er allwege spricht: Moses habe solches gemacht, wie ihm der HErr geboten hat, B. 16. 19. 21. 23. 26. 27. 29. 32.

Weil nun solches alles eine mosaische Ordnung, und im Neuen Testament aufgehoben und abgethan ist, kann jedermann leichtlich abnehmen, daß solche \unsicher{Papstweihe}, ohne Gottes Wort fürgenommen, gleich so wenig zu halten oder zu achten sei, als andere Menschen-Fündlein.

\part{Das dritte Buch Mose}

\chapter{Das 1. Kapitel}

\section{Vorrede}

Das 1., 2., 3., 4., 5., 6., 7., 8. u. 9. Kapitel werden in der Kirche nicht gelesen.

Dies Buch hat den Namen in der lateinischen Kirche, und heißt Leviticus, darum, daß es fürnehmlich das Priestertum fället, und ordnet mit Gesetzen und Rechten, nach welchen die Priester mit dem Opfer sich halten und das Volk lehren sollen.

Da stehet man fürnehmlich, wie das Priestertum nur um der Sünden willen eingesetzt ist, daß dieselbigen dem Volk erstlich angezeigt und darnach wieder versühnet sollen werden. Also, daß eines Priesters Amt anders nichts ist, denn mit Sünden und Sündern für Gott umgehen. Wie denn unser lieber HErr Christus darum auch den Namen hat, Ps. 110, 4., und ein Priester heißt, daß er sich der Sünder angenommen und für die Sünde sich selbst geopfert hat am Stamme des Kreuzes.

Daß man aber einen rechten Verstand habe von dem Priestertum und dem Opfer des Alten Testaments, muß man zum Eingang wissen, daß das Priesteramt und ganze Gesetz, so fern es das äußerliche Regiment und Kirchenordnung belanget, dazu gedienet hat, daß es nicht allein einen äußerlichen Unterschied zwischen Heiden und Juden machte, sondern daß es auch eine Uebung wäre des äußerlichen Gehorsams gegen Gott, welcher Gehorsam keine Vergebung und zeitliche Belohnung hätte. Aber zur Gerechtigkeit, so für Gott gilt, und zu dem ewigen Leben hat solches Priesteramt und Opfer nicht gedienet. Nach dem Spruch Joh. 1, 17.: Das Gesetz ist durch Mosen gegeben, die Gnade aber und Wahrheit ist durch Christum worden. Darum hat es auch ein Ende nehmen und das Gesetz sammt den Priestern und dem Opfer aufhören müssen, da Christus kommen ist; weil es nur ein Schatten und Fürbild der künftigen Güter in Christo gewesen ist, wie zum Ebräer am 10. B. 3. stehet.

Denn der rechte Priester, unser HErr Christus selbst, der opfert für der Welt Sünde nicht Kälber noch Ochsen, sondern sein eigen Leib und Blut. Durch dies einige Opfer werden wir von Sünden erlöset, wenn wirs mit Glauben fassen und unser Vertrauen darauf setzen. Durch dies einige Opfer werden wir mit Gott wahrhaftig versühnet, daß wir zur ewigen Kindschaft wieder kommen, welche wir durch die Sünde verloren haben.

Daß nun im Gesetz auch Sündopfer gewesen sind, ist nicht dahin zu rechnen, als hätten sie die Sünden für Gott dadurch abgelegt und versühnet. Denn dazu gehöret, wie gesagt, nicht Kühe- und Kälber-Blut, sondern das Blut unsers HErrn Jesu Christi. Und haben die heiligen Propheten und Patriarchen alle, von Adam bis auf Johannes, in solchem Fall allein auf den verheißenen Weibes-Samen und sein Opfer gewiesen, daß man desselbigen sich annehmen und wider die Sünde und Verdammniß trösten soll. Daß aber außer diesem Opfer die Juden, als ein sonder Volk, ihre eigene Opfer gehabt, die man auch Sündopfer nennet, solche Opfer haben dazu gedienet, daß die Juden am ersten damit für der Gemeine bekennet haben ihre Uebertretung am Gesetze, und sich darnach durch solch Sündopfer gegen der Gemeine wieder versühnet. Denn wo sie solche Ordnung nicht gehalten hät%
ten,\quellseite{182} wären sie Heiden und nicht mehr Juden geachtet worden. Neben demselbigen haben die Gläubigen solche Sündopfer angesehen als ein Fürbild des rechten Opfers, dadurch sie und die ganze Welt, der Verheißung nach, von Tod und Sünden durch Christum sollten erlöset werden. Durch solchen Glauben haben alle Heiligen im Alten Testament, vor und unter dem Gesetz, Vergebung der Sünden erlangt und sind selig worden. Das sei genug vom Unterricht und Eingang dieses Buches, von dem Opfer und Priester-Amt, davon die ersten neun Kapitel melden.

\section{Gesetz der Brandopfer}

\chapter{Das 2. Kapitel}

\section{Gesetz der Speisopfer}
\quellseite{183}

\chapter{Das 3. Kapitel}

\section{Gesetz von Dankopfern}

a Speise) das vom Feuer aufgefressen wird.

\chapter{Das 4. Kapitel}

\section{Gesetz von Sündopfern}
\quellseite{184}

a Ärgerte) das ist, mit Lehren oder Leben ungefähr zur Sünde und Schuld Ursache gäbe.
\quellseite{185}

\chapter{Das 5. Kapitel}

\section{Gesetz vom Schuldopfer}
\quellseite{186}

a Geweiheten) Sacrilegium, raptum scilicet ex hostiis per impudentiam. Das ist, Vergreifung am Geweiheten, das nämlich entwendet ist von den Opfern aus Versehen.

\section{Kap. 6. v. 1}

\chapter{Das 6. Kapitel}

\section{Vom Brand-, Speis- und Trankopfer}

a Nec extinguantur, nec alienus sit.
\quellseite{187}

b Aber alle das) hie siehst du, daß Moses klärlich zweierlei Sünd-Opfer oder zweierlei Brauch des Sünd-Opfers setzet: eines, da man das Blut nicht in das Heilige bringet zu versühnen, solches möchten sie essen. Das andere, da man das Blut in das Heilige bringet zu versühnen, solches müßte man nicht essen, sondern außer dem Lager alles verbrennen. Davon in der Epistel an die Ebräer 13, B. 11. Quorum animalium sanguis. Et supra, cap. 4. v. 5.

\chapter{Das 7. Kapitel}

\section{Vom Schuld- und Dankopfer}

a Lob-Opfer, Dank-Opfer) diese zwei Opfer sind in ein Opfer gerechnet. Dank-Opfer heißt, wenn sie Schafe, Ochsen etc. geschlachtet haben. Lob-Opfer, wenn sie Fladen und Kuchen (wie ein Speis-Opfer) dazu gethan hatten. Und nennet also eins das andere, daß es heißt Fleisch des Lob-Opfers, das ist, neben dem Lob-Opfer; item, Lob-Opfer des Dank-Opfers.
\quellseite{188}

\chapter{Das 8. Kapitel}

\section{Einweihung der Priester}
% Omitted: numbered Bible verse, excluded; printed folio or catchword marker; running head; not text apparatus.
\quellseite{190}

\chapter{Das 9. Capitel}

\section{Das erste Opfer Aarons wird vom Feuer verzehret}
\quellseite{191}

\chapter{Das 10. Kapitel}

\section{Nadab und Abihu vom Feuer getödtet}

\section{Summarium}

Ist ein schrecklich Exempel, und zur Warnung fürgeschrieben, auf daß jedermann sich hüte und nicht nach seinem Gutdünken, noch eigener Andacht, sondern allein nach dem Wort und Befehl Gottes handle, in den Sachen, die Gottes Ehre und den rechten Gottesdienst betreffen. Denn da die Söhne Aarons solches nicht thäten, mußten sie darum sterben, V. 2.

a Ich werde geheiliget) Qui mihi ministrant, sciant me esse sanctum, non prophanum. Das ist: die mir dienen, sollen wissen, ich sei heilig, nicht unheilig.

b Sein Blut) weil sein Blut nicht ins Heilige gebracht ist, sollte es nicht verbrannt, sondern gegessen worden sein. Welches Blut aber hinein gebracht wär, mußte man nicht essen, sondern alles verbrennen, v. 4. 5. c. 6, 30. und in c. 16. 26. 27.
\quellseite{192}
Darnach ists eine feine Warnung, da Mose Aaron und seinen Söhnen gebeut, daß sie nicht sollen Wein noch stark Getränke trinken, wenn sie in die Hütten des Zeugnißes gehen, auf daß sie nicht sterben, V. 8. Also vermahnet auch Christus seine Christen, daß sie sollen nüchtern sein. Luc. 19, 34. Und sonderlich St. Paulus 1 Tim. 3, 3. lehret: Ein Bischof soll nicht ein Weinsäufer sein. Denn solche Unart soll an keinem Christen und sonderlich an keinem Kirchendiener sein.

\section{Votum}

Der gnädige und treue GOtt regiere uns alle, daß wir in Gehorsam seiner Gebote wandeln, und daß wir uns für allen Sünden hüten, die zu Unehren seinem heiligen Namen, und uns zu Schand und Schaden gelangen mögen, Amen.

\chapter{Das 11. Kapitel}

\section{Vorrede}

Wie wir Gott für seine vielfältigen großen Wohlthaten selten herzlich danken: also gebrauchen wir solche selten recht. Gott hat uns durch JESUM Christum seinen lieben Sohn und unsern treuen Heiland befreiet von der Beschwerung der levitischen Gesetze Mosi und andern; wer danket Gott dafür? Also mißbraucht fast jedermann diese Freiheit, daß der Sünden alle Winkel der Welt voll sind. Derowegen laßt uns Gott danken, daß wir Freie sind, und laßt uns leben, nicht als hätten wir die Freiheit zum Deckel der Bosheit, sondern als die Knechte Gottes, 1 Pet. 2. Denn wir sind nicht zur christlichen Freiheit gekommen, daß uns frei stehe zu thun und zu lassen, was wir wollen, und also die Freiheit mit dem Evangelio zum Schand-Deckel der Sünden zu gebrauchen; sondern daß wir Gott dienen sollen, wie er uns solches in seinem Gesetz und Wort befohlen hat. Davon sagt St. Paulus Röm. 6. also: Nun ihr frei worden seid von der Sünde, seid ihr Knechte worden der Gerechtigkeit. Dies unterrichtet auch das 11. Kapitel.

\section{Vom Unterschied reiner und unreiner Thiere}

a Diese vier Thiere sind in unsern Landen nicht, wiewohl gemeiniglich Arbe und Hagab für Heuschrecken gehalten werden, die auch vierfüßige Vögel sind. Aber es ist gewisser, diese hebräische Namen zu gebrauchen, wie wir mit Halleluja und andern fremden Sprachen Namen thun.
\quellseite{193}

\section{Summarium}

Dies Kapitel sammt etlichen andern, so hernach folgen, sind von etlichen sonderlichen Gebräuchen, welche die Juden aus Gottes Ordnung und Befehl halten mußten, als mit dem Essen, daß ihnen etliche Speise verboten war; item mit der Unreinigkeit am Leibe, daß sie nicht durften unter die Leute gehen etc. Solches sind eitel Beschwerung gewesen, damit Gott dies Volk nicht allein hat wollen unterscheiden von den Heiden, sondern auch äußerlich beschweren, daß sie zum ersten genugsam Ursache hätten, den Gehorsam gegen Gott zu erzeigen; darnach daß sie nicht Raum noch Weile hätten, ihnen eigene Gottesdienste zu erwählen.

Wir Christen aber sollten Gott von Herzen danken, daß wir durch Christum von solchen äußerlichen Beschwerungen erlöset sind; und nicht allein solches zu halten nicht mehr schuldig, sondern auch von der höchsten Last des Gesetzes gefreiet sind, nämlich vom Fluch und Tode, so um der Sünden willen auf uns liegt.

\section{Votum}

GOtt sey Lob, Ehr und Preis für alle seine geistliche und leibliche Wohlthaten;
\quellseite{194}
der verleihe, daß mir durch diese seine Wohlthaten im Glauben gestärket werden, und unser Leben derowegen wohl bessern, Amen.

\chapter{Das 12. Kapitel}

\section{Vorrede}

Der Kinder Geburt ist ein unaussprechlich groß Mirakel und Wunderwerk. Aber wie wenig und selten danken Gott viele Menschen, daß er sie im Mutterleibe erschaffen und erhalten hat, und daß er sie auch mit vollkommenen Gliedmaßen, frisch und gesund in diese Welt hat lassen geboren werden. Derowegen wollen wir solche und andere Gottes Wohlthaten erkennen, und ihm alle Tage dafür danken, und darum ferner anrufen. Denn die Gesundheit des Leibes und der Glieder desselben nicht sein ein natürlich Gut, oder das ungefähr entstehet, wie viel solche Menschen es dafür halten; sondern sind eine besondere edle Gabe Gottes. Und schreiben die Anatomici und gelehrten Aerzte, daß der menschliche Leib so viel Glieder habe, als Tage im Jahre sind. Sollen wir denn nun nicht Gott danken, und bitten täglich für die Gesundheit unsers Leibes und unserer Glieder? Dies erinnert auch das 12. Kapitel.

\section{Ordnung der Kindbetterinnen}

\section{Summarium}

Das ist eine mosaische Ordnung von den Sechswöchnerinnen oder Kindbetterinnen, welcher wir im Neuen Testament eben so wenig verbunden sind, als den andern von Essen und unreiner Speise. - Darum ists ein Mißbrauch und unwesen, daß man das Einsegnen im Papstthum, nach den Sechswochen oder Kindbett hat für nöthig erachtet.

Daß man aber die Sechswöchner halten, geschieht nicht um Moses Ordnung willen, sondern daß es die Natur und Noth also erfordert, daß eine Kindbetterin nach der Geburt sich innen halten und eine Zeit lang ruhig sein, und mit gebührlichem Essen und Trinken ihr auswarten soll; nicht allein ihret halben, daß sie wieder zu Kräften komme, sondern auch des Kindes halben, daß das nicht verwahrloset werde, weil es noch schwach ist, und seine Nahrung von der Mutter haben muß.

\section{Votum}

Wir danken Gott herzlich für alle Wohlthaten, die er uns durch unsere liebe Eltern erzeiget hat; der segne und behüte alle schwangere Hausmütter und Sechswöchnerinnen, mit ihren Kinderlein, und verleihe, daß wir mit ihnen allen selig werden, Amen.

\chapter{Das 13. Kapitel}

\section{Kennzeichen des Aussatzes an den Menschen und Kleidern}
\quellseite{195}

a Die ganze Haut) dieser Aussatz heißt rein, denn es ist ein gesunder Leib, der sich also selbst reiniget, als mit Pocken, Masern und Krätze geschiehet, da durch den ganzen Leib das Böse heraus schlägt; wie wir Deutschen sagen, es sei gesund etc.
% Omitted: numbered Bible verse, excluded; printed folio or catchword marker; running head; not text apparatus.
\quellseite{197}

\chapter{Das 14. Kapitel}

\section{Reinigung des Aussatzes}

a Log, Log ist ein klein Maßlein, auf Hebräisch also genannt, aber noch ungewiß, wie groß es sei.
% Omitted: Bible narrative, running matter, or non-apparatus block; numbered Bible verse, excluded; printed folio or catchword marker; running head; not text apparatus.
\quellseite{199}

\chapter{Das 15. Kapitel}

\section[Manns- und Weibspersonen, mit unreinem Fluß behaftet, wie]{Manns- und Weibspersonen, mit unreinem Fluß behaftet, wie sie zu reinigen}
\quellseite{200}

\chapter{Das 16. Kapitel}

\section{Jährliches Versöhnopfer}
\quellseite{201}

a Unrein) das ist, ganz Israel ist unrein und in Sünden, und Gott rechnet noch mitten unter viel Sündern, und heiliget sie; das ist, keine Heiligkeit genießen sie, und heißen heilig um seinetwillen, sonst ist alles unrein und Sünde mit ihnen, das ist Christus in seinem Volk.

\chapter{Das 17. Kapitel}

\section{Ort der Opfer bestimmt. Blut zu essen verboten}

a Schuldig) da siehst du, daß er \unsicher{nicht will einen selbst erwählten Gottesdienst} haben außer seinem geordneten Gottesdienst. Darum nennet er einen solchen Opferer einen Mörder, wie er 1 Mos. 9, 6. auch thut.
\quellseite{202}

b Priester nicht sie selbst, aus eigener Wahl und Andacht.

\chapter{Das 18. Kapitel}

\section{Vorrede}

Wie wir uns für andern Sünden hüten sollen: also die, so sich in Ehestand begeben wollen, sollen \unsicher{Ehetrei} fliehen und meiden, das ist, sie sollen nicht zu nahe in die Blutsfreundschaft oder Schwägerschaft freien. Denn Unzucht treiben ist große Sünde, die Ehe brechen ist größer; in Blutschanden leben ist die größte. Aber da sind viel evangelische, eigenwillige Leute, die keiner guten Ordnung, ja auch Gottes Gebot nicht gehorsamen, und ziehen also Blutschulden und Gottes Strafen auf Land und Leute. Wie denn dieser Sünden Mose gedenkt als der fürnehmsten Ursachen einer der Sündflut, 1 Mos. 6. Solches zu verhüten, soll jedermann bekannt sein das 18. Kapitel des dritten Buches Mose.

\section{Verbotene Grade der Blutsfreundschaft im Heirathen}
\quellseite{203}

a Moloch) war ein Abgott, dem sie ihre eigenen Kinder zu Dienst verbrannten; wie Manasse that, der König Juda, und meineten, Gott damit zu dienen, wie Abraham, da er Isaak, seinen Sohn, opfert. Aber weil ihnen Gott das nicht befohlen hatte, wie dem Abraham, war es unrecht. Darum spricht hie Gott, daß sein Name dadurch entheiliget werde; denn es geschah unter Gottes Namen und war doch teuflisch. Wie auch jetzt die Kloster-Gelübde und andere Menschen-Aufsätze viele Leute verderben, unter göttlichem Namen, als sei es Gottesdienst.

\section{Summarium}

Das ist eine Ordnung, die da belanget das sechste Gebot, in welcher angezeigt wird, welche Personen zum ehelichen Beiwohnen erlaubt oder verboten sind.

Solche Ordnung ist nicht zu achten wie andere Gesetze, so von Mose geboten, nun aber im Neuen Testament aufgehoben sind; sondern sie zeucht sich auf die Natur. Denn der Text meldet klar, wie solches unordentliches Beiwohnen auch unter den Heiden eine Sünde sei, V. 24. Und daß um solcher Sünden willen, als Ehebruch und andere Unzucht, Land und Leute ausgeräumet und verwüstet müssen werden, V. 25. Darum soll jedermann sich für Unzucht hüten, und sonderlich der Blut-Freundschaft schonen, weil es Gott so theuer verboten und so greulich auch unter den Heiden gestraft hat.

\section{Votum}

Helfe der treue GOtt, daß wir seinen heiligen Geboten folgen, und als Gehorsame an Leib und Seel von ihm gesegnet werden, um JEsu Christi seines lieben Sohns Gehorsam und Verdiensts willen, Amen.

\chapter{Das 19. Kapitel}

\section{Vorrede}

Wenn zusammen sollte gerechnet werden, wie oft die meisten Menschen in der ganzen Welt an die heiligen zehn Gebote gedenken, oder wie
\quellseite{204}
oft sie dieselbigen bei ihrem Gebet sprechen: ach, die Summa würde klein und geringe sein. Aber wie Gott durch Mosen oft die heiligen zehen Gebote hat wiederholen und erklären lassen: also will Gott, daß alle Menschen die heiligen zehen Gebote täglich sprechen und betrachten sollen, daß sie dadurch zur Buße gereizt werden, daß ist, daß sie ihre Sünde erkennen, die Gott abbitten, des Herrn Christi im wahren Glauben sich trösten, und darauf ein Gott wohlgefälliges und ihrem Nächsten nützliches Leben führen. Darauf höret davon an das 19. Kapitel.

\section{Auslegung der zehn Gebote, sammt andern Gesetzen}

a) Gefallen könnte) das ist: Macht's, wie es gebeut, nicht wie euch dünket oder gefällt nach eigener Andacht. Denn er will schlechthin keine selbst erwählte noch eigen erdachte Weise haben. Darum sehet zu, daß ihr also opfert, daß ihm gefalle, und nicht weiter erzürnet werde durch eure eigene Weise. Nam praeter juge sacrificium non erat necessarium vovere et offerre ei praeterea pro peccatis; illa pacifica sint libera. Das ist: denn außer dem täglichen Opfer war nicht nöthig, Gelübde und Opfer zu thun überdieß noch für die Sünden; jene Dankopfer sollen freiwillig sein.
\quellseite{205}

b Vorhautschneiden ist hie so viel, als drei Jahr barren, wie er selbst deutet und spricht: Drei Jahr sollt ihr sie achten für unbeschnitten etc.

\section{Summarium}

Dies ist eine Predigt, in welcher Gott etliche Stücke aus den zehen Geboten wiederholet und ausstreicht, nämlich, wie man den Nächsten lieben, ihm dienen und allerlei Freundschaft beweisen soll. Item, wie man sich für Aberglauben, Wahrsagen und dergleichen Sünden hüten soll, V. 31.

An solche Gebote hänget er alle Wege diesen Beschluß: Ich bin der HErr euer Gott; damit er anzeiget, daß wir Gott Gehorsam schuldig sind, in allem, was er haben will, es sei so gering es wolle. Auf daß wir allein auf sein Wort und Befehl sehen, und uns nichts an solchem Gehorsam hindern lassen.

\section{Votum}

GOtt stärke unsere schwachen und furchtsamen Herzen, und verleihe, daß wir in Noth und Tod seinen Geboten gehorsam sein, und in wahrem Glauben, Liebe, Hoffnung und Geduld leben und sterben. Amen.

\chapter{Das 20. Kapitel}

\section{Vorrede}

Gott, als unser Schöpfer und lieber Vater, meinet ja treulich mit uns Menschen, auf daß er uns segnen könne und nicht strafen dürfe. Denn beides thut er in seinem Wort: er lehret, für was Sünden wir uns in allen Ständen hüten sollen, und zeiget an, daß er die Sünden nicht wolle ungestraft lassen, sondern die an allen Unbußfertigen mit zeitlicher und ewiger Plage und Marter belohnen. Derowegen, wer Gottes Treue erkennet, der lebe täglich nach seinem (im Wort) offenbarten Willen; und wer ihm selbst Gutes gönnet, der hüte sich für allen Sünden, so hütet er sich auch also vor den Strafen der Sünden. Sonst heißt es, wie Gott saget in diesem dritten Buch Mose am 26. Kapitel: Werdet ihr mir nicht gehorchen und nicht thun diese Gebote alle, und werdet meine Satzungen verachten, und eure Seele meine Rechte verwerfen, daß ihr nicht thut alle meine Gebote, so will ich euch auch solches thun: ich will euch heimsuchen mit Schrecken, Schwulst und Fieber, daß euch die Angesichter verfallen und der Leib verschmachten soll. Darauf höret ferner an das 20. Kapitel.

\section{Strafen unterschiedlicher Sünden}

a Das ist, seiner Kinder.
% Omitted: numbered Bible verse, excluded; printed folio or catchword marker; running head; not text apparatus.
\quellseite{207}

\section{Summarium}

Die ist eine Predigt, in welcher Gott etliche Gesetze wiederholet, wie man den Abgöttern nicht opfern, Wahrsagen und Zeichendeuten nicht brauchen soll, Gott wolle sonst selbst, wo es schon die Leute nicht thäten, solche Sünde strafen, V. 4. 5.

Sonderlich aber ist zu merken, daß er befiehlt, mit dem Tod zu strafen, wo jemand den Eltern fluchet, V. 9., seine Ehe bricht, V. 10., oder unnatürliche Unzucht treibet, V. 15. Item, daß er die, so zu nahe in die Blutsfreundschaft sich verheirathen, mit Unfruchtbarkeit strafen und ohne Kinder wolle sterben lassen, V. 20. 21. Solche Gebote ziehen sich fast alle auf die Natur, darum binden sie uns gleich so wohl als die Juden.

\section{Votum}

Der treue barmherzige GOtt, der da befiehlet, was wir thun sollen, der verleihe Gnade und Hülfe durch den heiligen Geist, daß wir auch seinem Befehl folgen und nachkommen mögen, und in seinem Gehorsam leben und sterben, Amen.

\chapter{Das 21. Kapitel}

\section{Wie sich ein Priester zu verhalten}

a Krone) auf Hebräisch Nezer, hiervon die Nazarei heißen, daß sie sich enthielten, und sonderten von weltlicher Speise, 4 Mos. 6, 3. Und Simson von Mutterleib ein Nazareus hießet, Richt. 13, 5. Und unser HErr Christus auch daher der rechte Nazareus heißet von Nazareth, Matth. 2, 23.; als der von aller Sünde rein und heilig und kein Unheiliges noch Ungeweihtes an ihm ist. Wiewohl die Juden auf Neid und Feinheit ihn nicht Nazar, sondern Nozri, das ist, den Verächter und Verstörten, wie einen Schänder nennen, so sie doch selbst die rechten Nazirim sind, in alle Welt zerstreuet und verstöret. Und mich dünkt, daß sich St. Paulus Röm. 1, 1. einen Nazareum nennet, da er sich einen Abgesonderten rühmet zum Evangelio, vom Gesetz c. 1 Tim. 3, 2. Tit. 1, 7.
\quellseite{208}

b (ungewöhnlich und seltsam) das sind allerlei Ungestalten. Als so die Nasen zu groß, zu klein, zu krumm, zu breit, item, das Maul krumm, schartig, blöcket \&c., das einen für Andern verstellet.

\chapter{Das 22. Kapitel}

\section{Von des Opfers Beschaffenheit}

a (Enthalten) das ist, wie die Nazaräer sich heiliglich abhalten und meinen, Nazaräer etc., nämlich, daß sie nicht essen sollen vom Opfer, sie seien denn ganz rein und heilig.
\quellseite{209}

b Fehl) über dies Stück klagt Maleachi 1, 8. sehr hart; denn was nicht taugt noch gut ist, das gibt man Gott und seinen Dienern, und wills also ertragen. Er aber doch solches hie als unangenehm verdammt, denn er hat lieb einen fröhlichen Geber.

c Wandelbar) als da nur ein Ohr oder keins, das ein Auge oder keins, daß ein oder mehrere Füße mangelt, oder sonst ungestaltet und ungeschaffen ist; es heißt: Wer geben will, der gebe was Gutes, oder lasse es anstehen.

\chapter{Das 23. Kapitel}

\section{Vorrede}

Den Weltkindern sind alle Tage Sünden-Tage; dagegen sind den Christen alle Tage Sabbath-Tage, das ist, solche Tage, an denen sie, wo sie nur können, Gottes Wort hören, an denen sie auch Gottes Wohltaten betrachten, ihm dafür danken, für sich, die Ihren, und für die Noth der ganzen Christenheit beten, von Sünden ruhen und die Tage mit heiligen Gott wohlgefälligen Werken, wie es ihr Beruf und Stand erfordert, zubringen. Derowegen, weil wir auch Christen sein, so wollen wir dergleichen alle Tage, als Sabbaths-Tage, heilig halten und christlich begehen. Dazu uns denn helfe unser lieber HErr JEsus Christus, hochgelobet in Ewigkeit, Amen. Darauf höret ferner an das 23. Capitel.

\section{Ordnung der vornehmsten Feste}

a Dienstbarkeit) als da sind die Werke, so man an den Werktagen thut, Nahrung zu suchen, da man Getreide und Vieh zu brauchet. Aber Haus-Geschäfte und Gottesdienst ist nicht verboten, als kochen, kehren, kleiden etc.
\quellseite{210}

b Blasens zum Gedächtnis] solch Blasen mit einem Horn geschah, daß man damit Gottes und seiner Wunder gedachte; wie er sie erlöset hatte, davon predigte und dankte. Wie bei uns durchs Evangelium Christi und seiner Erlösung gedacht und geprediget wird.
\quellseite{211}

\section{Summarium}

Hie ist eine Ordnung von den fürnehmsten Festen, so Gott durchs ganze Jahr den Juden hatte befohlen zu halten. Welche Feste alle mit einander dahin gestiftet sind, auf daß man Gottes Güte und Wunderwerke nicht aus dem Gedächtnis lasse, sondern Gott dafür soll dankbar sein, und sich seiner Hülfe auch im künftigen Unglück trösten.

Den Sabbath nennt er das fürnehmste Fest, V. 3. Denn allewege der siebente Tag dahin gestiftet ist, daß man Gottes Wort hören und Gottes Werk betrachten und Gott für einen Schöpfer Himmels und der Erden erkennen und loben soll, von welchem alle Menschen Leib und Leben haben und alles gewarten, bitten und hoffen sollen, was zur Erhaltung Leibes und Lebens und endlich zur Seligkeit vonnöthen ist. Solches heißt den Sabbath nicht allein feiern, sondern auch heiligen. Es ist aber solches Fest oder Sabbathtag nicht allererst von Mose geordnet, sondern bald nach Vollendung der Kreaturen, ehe denn Adam und Eva sind gefallen, von Gott selbst aufgesetzt und zu solchem Gottesdienst geheiliget, das ist, dazu verordnet und darauf geordnet, daß man ihn mit heiligen Dingen soll zubringen. Solche Satzung hat Moses verneuert. Aber da sind der Propheten Zeugniß und Predigt immer daneben gegangen, im Neuen Testament: wenn der verheißene Same kommen werde, werde solcher mosaische Sabbath fallen und alle Tage Sabbath sein, für und für, wie Christus auch sagt: Solches thut zu meinem Gedächtnis, so oft ihrs thut. Und Paulus: Verkündiget den Tod des HErrn, bis er kömmt, das ist, alle Tage für und für rühmet Gottes Gnade und Wohlthat in Christo. 1 Cor. 11, 26.

1. Darnach waren andere sonderliche Feste. Das erste war das Osterfest, V. 6. Daran man bedenken sollte, wie Gott sein Volk wunderbarlich aus Egypten geführet und den Tyrannen Pharao mit seinem Volk gestürzet hat, und Gott für solche Erlösung dankbar sein, und sich in anderem Unglück seiner Hülfe auch trösten. Neben solchem Dank war dies Fest auch eine Erinnerung der künftigen Erlösung, durch das rechte Osterlamm, den Sohn Gottes, unsern HErrn Jesum Christum, der nicht allein die Juden, sondern alle Welt von dem schädlichen Pharao und Egypten, das ist, von Sünde, Tod und Hölle erlösen sollte. Wie St. Paulus solches selbst deutet 1 Cor. 5. V. 7.

2. Das andere war das Fest der Erstlinge, V. 10., daß allewege, wenn der Schnitt anging, man eine Garbe von einem jeden Gewächs im Tempel opfern mußte, sammt einem Brandopfer, Speisopfer und Trankopfer; und darnach fünfzig Tage zählen, bis die Feld-Arbeit gar verrichtet ward. Da mußte aus einem jeden Hause Brod von neuem Getreide in dem Tempel geopfert werden. Solches geschah, auf daß die Leute lernten Gott für seinen Segen danken und erkennen, daß es Gottes Gabe sei, wenn der Feld-Bau wohl geräth, und Gott allerlei, was wir bedürfen, reichlich läßt wachsen. Man soll aber solchen Segen nicht allein dahin brauchen, daß wir alles vollauf und genug haben, sondern fürnehmlich damit helfen den äußerlichen Gottesdienst erhalten und den Kirchendienern ihren Theil davon geben, auf daß sie sich desto füglich erhalten, ihres Amts warten und zur Erkenntniß Gottes und Seligkeit Vielen helfen können. Solcher Ursache halben hat Gott dies Fest geordnet.

3. Darnach war das Posaunenfest, V. 24., welches man hielt am ersten Tage des siebenten Monats nach Ostern, welches bei uns ungefährlich ist der Weinmonat, wenn alle Frucht des Feldes, Korn, Wein, Oel und was da wächst,
\quellseite{212}
eingeführet sind. Solches Fest scheinet auch, daß es am fürnehmsten dahin hat sollen dienen, auf daß man Gott für seinen Segen dankbar wäre.

4. Am zehenten Tage nach solchem Fest war das Fest der Versühnung, B. 27., da mußten die Juden fasten und feiren. Denn es gehört sich, wenn Gott seinen Segen gibt, daß wir erkennen, wie wirs nicht werth sind und wohl alles Unglück unserer Sünden halben verdienet hätten, wo Gott sonst nicht so gnädig wäre. Darum mußte man am selbigen Feste die Sünde bekennen und Gott um Versühnung und Gnade bitten. Aus diesem Fest haben die Papisten die Weise genommen, daß sie zur Genugthuung für die Sünde den Leuten sonderliche Festtage haben aufgelegt. Aber es ist solches Fasten schlechter Gottesdienst, ja auch wohl ein Greuel für Gott; sintemal Gott davon uns nichts geboten, sondern Vergebung der Sünden im Namen Jesu Christi allein zu predigen befohlen hat, und uns gar nichts angehet, was Gott in solchem Fall den Juden befohlen hat. Was aber das für eine Versühnung gewesen, ist im Anfang dieses Buches gemeldet.

5. Das letzte Fest war am fünften Tage nach dem Versühnfest und hieß das Fest der Laubhütten, B. 34. Da mußten die Juden unter keinem Dach, noch in ihren Häusern weder essen, schlafen, noch anders thun, sondern mußten unter dem offenen Himmel bleiben sieben ganze Tage, zum Gedächtniß, daß sie Gott in der Wüste ohne solchen Segen, Weins, Korns und anderes Gewächs wunderbarlich erhalten und ernähret hat.

Auf daß wir auch lernen, obgleich solche Feste mit andern Moses-Ordnungen gefallen sind, daß wir dennoch Gott um alle seine Wohlthaten dankbar sein und seinen Segen gern zu Unterhaltung der Kirchendiener und des rechten Gottesdienstes brauchen sollen. Sonderlich aber haben wir im Neuen Testament, wie oben auch gemeldet, alle Tage unsern Sabbath, unsere Ostern- und Versühn-Feste, weil Christus für uns gestorben ist und sein Verdienst durch sein Wort uns geschenket und zu eigen gibt. Für solche Wohlthat, neben dem zeitlichen Segen, sollen wir Gott alle Tage (weil wir hie gleich wie die Juden in der Wüste, ehe sie in das gelobte Land und die rechte Heimath kommen, in unserer Pilgerschaft sind) danken und ihn loben, bis wir endlich gar selig werden und in die ewige Hütte kommen, Amen.

\chapter{Das 24. Kapitel}

\section{Vorrede}

Wenn sonst keine Sünde auf Erden wäre, denn die große Gotteslästerung und der gräuliche Mißbrauch göttlichen Namens, so wäre es Sünde genug, Gottes Zorn und alle Strafen zu erwecken und auf die Welt zu ziehen; und wäre freilich nicht Wunder, daß Gott nicht allein mit Pestilen, Krieg und Theurung die Welt strafe, sondern auch wohl mit Schwefel und höllischem Feuer auf sie schienen und schlacken ließ. Derowegen wer an dieser Sünde schuldig ist, der thue heute Buße, und wer dieser Sünde wehren kann, der thue es bald, um Gottes willen. Denn versucht sei die Zunge, die lästerlich redet wider den, der sie geschaffen hat, sagt Gordius der Märtyrer. Dagegen mag man auch sagen: Selig sei die Hand, welche die Gotteslästerung strafet, und die es wehren soll und kann. Davon höret auch an das 24. Kapitel.

\section{Vom Leuchter. Schaubrode. Strafe der Gotteslästerer und Todtschläger}

a Denkbrode) das sind die Schaubrode, welche die Kuchen heißen, darum, daß sie breit waren wie Kuchen. Und sind Denkbrode, darum, daß sie damit Gottes gedenken, und von ihm predigen sollen. Gleich wie Christus uns befehlet, daß wir sein gedenken, das ist, seinen Tod verkündigen und predigen sollen, 1 Cor. 11.

b Feuer) das ist Opfer.
\quellseite{213}

\section{Summarium}

Am ersten ist da ein Exempel, wie das Volk von dem Thron das Licht im Tempel erhalten und Oel dazu habe geben müssen, B. 2. Denn Gott will, daß jedermann etwas dazu geben soll, daß rechter Gottesdienst erhalten werde.

Sonderlich aber reimet sich die geistliche Deutung fein. Das Licht im Tempel ist das Predigtamt und die Kirchendiener, wie Christus seine Apostel Matth. 5, 14. auch nennt, das Licht der Welt. Solche sollen ihren Unterhalt von der Gemeine haben, wie die Juden zum Licht das Oel mußten geben.

Darnach ist ein ernst Exempel, wie Gott gebeut den zu steinigen, der bei Gottes Namen geflucht hatte, V. 13., auf daß wir Gottes Namen, Gott zu Lob und Ehre, lernen recht brauchen, mit Predigen, Beten und Bekennen, und nicht also zu Unehren mißbrauchen, wie leider jetzt in der Welt geschiehet und wir derhalben erstlich alle insgemein, und darnach ein jeder in seinem Hause mancherlei Unglück müssen tragen. Denn es ja eine gräuliche, schreckliche Sünde ist, daß wir aus schändlicher böser Gewohnheit unsers lieben Herrn Christi, seiner Taufe, Sacrament, Leiden und Marter, welches alles um unsrer Seligkeit willen geschehen, dahin mißbrauchen, daß wir andern Böses dadurch wünschen. Und ist gewiß, solche Sünde verdammt Leib und Seele an vielen, daß sie hie all ihr Glück und Segen durch solches Fluchen und Lästern selbst verschlagen und ihnen das höllische Feuer verdienen. Darum lerne ein jeder Christ dafür sich hüten, sonderlich aber soll weltliche Obrigkeit mit harter Strafe solchem Aergerniß wehren und Gott seine Ehre retten, so viel möglich.

\section{Votum}

GOtt regiere die Obrigkeit, daß sie die Gotteslästerung ernstlich strafe, und regiere uns, daß wir seinen heiligen Namen nicht schänden, sondern denselben im Herzen und mit Worten und Thaten loben und preisen, Amen.

\chapter{Das 25. Kapitel}

\section{Vorrede}

Daß jetziger Zeit allerlei Händel und Gewerbe stecken bleiben, dessen ist nicht eine kleine Ursache, daß man allerlei Betrug in Käufen und Verkaufen gebraucht; und, da Gott segnet, daß der meiste Theil so wenig der Armuth davon mittheilet. Derowegen, ob schon die Zeit schwer ist, dennoch nähre sich ein jeder mit gutem Gewissen, und lasse ihm arme Leute befohlen sein. Also wird ihm an Gottes Segen zu täglichem Brod und an Trost, unter gemeinem Kummer, nicht mangeln. Von diesem sagt Salomon in seinen Sprüchen im 19. Kapitel also: Wer sich des Armen erbarmet, der leihet dem HErrn, der wird ihm wieder Gutes vergelten. Von jenem aber spricht also David im 128. Psalm: Du wirst dich nähren deiner Hände Arbeit; wohl dir, du hast es gut. Dies unterrichtet auch das 25. Kapitel.

\section{Feier- und Jubeljahr}
\quellseite{214}

a Speise) das ist, gemein sein und nicht einsammeln noch ausschütten.
\quellseite{215}

b Ueberfaß) Wucher heißt er, so mit Geld geschiehet; Ueberfaß, wenn der arme Mann muß kaufen oder annehmen die tägliche Waare, so theuer der Geizhalß will, weil ers haben muß zur Noth.

\section{Summarium}

Hier ist eine Ordnung etlicher äußerlicher Händel, welche alle dahin gerichtet sind, daß der Armen nicht vergessen und sie mit Kaufen und Verkaufen nicht übervortheilet, noch beschweret werden. Denn daß bei den Juden allewegen das siebente Jahr das Feld ganz und gar feiern mußte, geschah, wie der Text meldet, v. 6., auf daß alles, was von sich selbst wüchse, jedermann gemein wäre und die armen Leute sammt dem Vieh ihre Nahrung desto besser haben möchten.

Digitized by Google
\quellseite{216}
Also ist die Ordnung mit dem fünfzigsten Jahr auch gewesen: obschon die Leute verdarben und ihre Aecker und Güter Armuth halber verkaufen mußten, daß sie doch im Jubel-Jahr wieder zu Gütern kamen, V. 10. Solches ist eine überaus gute und gleichmäßige Ordnung gewesen, daß die Güter nicht haben können ganz verruckt werden.

Sonderlich aber ists fein, daß Gott will, daß in Kaufen und Verkaufen allein die Billigkeit und kein beschwerlicher Vortheil soll gesucht werden, V. 16. Und gehet, man soll die Armen mit Wucher und Aufsatz nicht beschweren und sonst auch allewege einen Bruder gnädiger halten, denn einen Fremden, V. 25. 35. 39. 47.

\section{Votum}

Helfe der treue GOtt, daß diese Lehren in uns Nutz schaffen, und daß ein jedes also lebe, auf daß es an Sünden arm, und an Tugend reich sei und bleibe, und daß wir alle um Christi willen das ewige Jubeljahr im Himmel halten mögen, Amen.

\chapter{Das 26. Capitel}

\section{Vorrede}

Gott will, daß kein Mensch verloren werde, sondern daß sich ein jeder bekehre, seiner Barmherzigkeit und Güte genieße und selig werde. Derowegen locket er zur Buße durch ernste Predigten und Strafen. Solche Treue und Gnade Gottes lasset uns erkennen und wohl gebrauchen; daß ein jedes von seinen Sünden ablasse und dagegen thue, was es Gott zu Ehren, seinem Nächsten zu Nutz und ihm selbst zur Wohlfahrt, nach Gottes Wort, thun soll und kann. Denn je treulicher Gott warnet, nimmt man die Warnung nicht an und führet fort in Sünde, Leben und Unbußfertigkeit, je härter und grausamer er darauf strafet, wie solches die heilige Schrift und die Erfahrung bezeuget. Davon höret auch an das 26. Capitel.

\section{Gedrohter Fluch, und verheißener Segen}
\quellseite{217}

\section{Summarium}

Ist eine herrliche und ernste Predigt, in welcher Gott sein Volk vermahnet, V. 3. Weil sie nun wissen, was sie thun und lassen, wie sie sich mit Opfern und dem Gottesdienst halten sollen, daß sie bei solchem Befehl bleiben und nichts dran ändern, noch unterlassen sollen. Und verheißet ihnen, wo sie solches thun, er wolle mit reichem Segen auf dem Felde, mit Frieden im Lande und mit allem Glück sie überschütten. Wiederum, wo sie es nicht thun, da soll Krankheit, Hunger, Krieg, schädliche Thiere und alles Unglück sie treffen ohne Aufhören, V. 14.

Da sehet ihr nun, wo solcher Jammer in der Welt alle herkömmt; darum mag sich jedermann wohl bessern und frömmer werden, so wir anders solcher Strafen wollen überhoben sein.

a Ueber den andern) wie es geschieht in der Flucht, Furcht und Schrecken.

b Gefallen) das ist: gleich wie sie Lust an ihren Sünden und Ekel an meinen Rechten hatten: also werden sie wiederum Lust und Gefallen haben an der Strafe, und sagen: ach wie recht ist uns geschehen; Dank hab unsere verfluchte Sünde, daß haben wir nun davon. O recht lieber Gott, o recht. Und das sind Gedanken und Worte einer Erkenntniß und Buße, die sich selbst aus Herzensgrund hassen und anklagen lehrt, \unsicher{wozu sich, was hab ich gethan?} daß gefällt denn Gott, daß er wieder gnädig wird. Darum haben wir das Wort Missethat verbüßt die Strafe der Missethat, solchen Verstand zu geben, sonst lautet's, als sollten sie Gefallen an der Missethat haben. Eben so ist's auch zu verstehen: dem Lande gefällt seine Feier, das ist, es spricht: Gott habe Recht in der Strafe, daß es wüste liegen muß um des Volks willen, nach dem es sich sehnet. Solcher Weise redet auch Jes. 40, 2. Dimissa est iniquitas, id est, placita et accepta poena.
\quellseite{218}

\section{Summarium}

Es ist eine heftige und schreckliche Predigt, darin Gott dräuet den unbußfertigen Sündern, daß Gott mit der Strafe nicht aufhören, sondern je mehr man mit Sünden fortfähret, je mehr er mit der Strafe anhalten wolle. Auf daß man lerne, daß Gottes Zorn mit nichts mag abgelegt werden, es sei denn, daß man sich bessere und von Sünden ablasse, und Gnade begehre.

Am Ende ist ein sehr tröstlich Stück, daß Gott solche harte Ruthen über die Sünder nicht darum gehen lasse, daß er sie gar verderben, sondern allein, daß er sie zur Buße treiben und frömmer machen und wieder zu Gnaden annehmen wolle, V. 42. Denn weil Gott anders nicht, denn gnädig ist, will er seinen Gnaden-Bund nicht gar um unserer Sünden willen aufheben. Darum sollen wir in Sünden nicht verzagen, sondern davon ablassen, und uns mit Gottes Barmherzigkeit trösten.

\section{Gebet}

Erleuchte und regiere uns der gütige treue GOtt, daß wir für allen Sünden uns hüten, und dagegen gottselig und ehrbar allezeit leben, auf daß wir der Sünden wegen nicht zeitlich und ewig an Leib und Seele gestraft werden, Amen.

\chapter{Das 27. Kapitel}

\section{Vorrede}

Wir sind alle für Gott Missetäter. Ach eine große schwere Sünde! Denn wir alle den Eid oft und gräulich gebrochen, den wir in der Taufe Gott gethan und da wir zugesagt haben, daß wir wider unser Fleisch und Blut, wider den Teufel und wider die Welt streiten, Sünde meiden und dagegen in vollem Gehorsam der heiligen zehn Gebote Gott dienen wollen. Solches läßt uns erkennen, Buße thun und Gott leisten, durch unser ganzes Leben, was wir ihm schuldig sind. Denn halten wir, was wir Gott verheißen haben, so wird Gott wieder halten, was er zugesagt hat. Darauf höret auch an das 27. und letzte Capitel.

\section{Von Gelübden und Zehnten}
\quellseite{219}

\section{Summarium}

Dies ist eine Ordnung, wie es mit den Gelübden und den Zehnten soll gehalten werden. Weil nun unsere Päpisten viel von ihren Gelübden rühmen, welche ihre Pfaffen, Mönche und Nonnen thun, soll man diesen gewissen und wahren Unterricht davon wissen.

Die Gelübde im Alten Testament sind fast alle dahin gerichtet, daß den Leuten, so zum Kirchendienst verordnet waren, damit Gott gedienet würde, wie man hie aus Mose klar siehet, daß sie die Häuser, Acker, Vieh, ihre eigenen Leiber eine Zeit lang, verlobet haben und bisweilen mit Geld wieder abgelöset, welches Geld alles dem Tempel und Gottesdienst ist zu Nutz kommen, B. 7.

Solche Gelübde haben dreierlei sonderliche Anhänge gehabt. Erstlich ist es den Juden nicht frei erlaubt gewesen zu geloben, was und wohin sie wollten, sondern es ist hie in Mose das Ziel gesetzt, und mit Gottes Befehl ausgedrückt, was und wohin man geloben soll.

Zum andern sind solche Gelübde nicht ewig gewesen, sondern haben nur eine Zeit lang gewähret, und nach der bestimmten Zeit wieder aufgehöret, ohne daß Gott etlichen Personen insonderheit hat auferlegt, ihr Leben lang sich zu halten, daß anderen nur eine Zeit lang geboten, wie man siehet mit Simson und andern.

Zum dritten waren solche Gelübde dermaßen, daß es allwegen in der Oberherrn Hand stund, dieselbigen zu brechen. Kinder und Gesinde konnten nichts geloben, wider Wissen und Be%
willigen\quellseite{220} des Vaters oder Haußherrn. Ein Weib durfte nichts geloben, ohne Begunst und Verwilligung des Mannes, wie man klar in Mose siehet.

Wenn wir nun unserer Papisten Gelübde gegen diese halten, so ist es ganz und gar mit ihnen umgekehret. Denn erstlich sollen nach des Papstes Recht Vater und Mutter nicht Macht haben, daß sie ihrer Kinder und Ehehalter Gelübde brechen mögen, sondern in solchem Fall billiget es der Papst, daß der Kinder Wille fortgehe und die Eltern ihnen müssen weichen.

Zum andern sind sie nicht auf eine gewisse Zeit bestimmt, sondern bleiben ewig.

Zum dritten, und welches das größte ist, geschehen solche Papst-Gelübde ohne Befehl Gottes, aus eigenem Fürnehmen, ja auch wohl wider Gottes Befehl und Ordnung, wie wir hernach hören werden.

Nun, solche Gelübde des Alten Testamentes, welche nach Gottes Wort und Willen geschehen, sind dermaßen gewesen, daß dennoch die Propheten, wie man im Psalter und sonst siehet, immerdar dies Gelübde weit darüber gerühmet haben, daß man Gott preisen, ihn für gnädig und barmherzig halten und um seiner gnädigen Hülfe willen von Herzen danken soll. Wie es denn der heil. Geist ein Gelübde heißt, Ps. 22. V. 26.: Dich will ich preisen in der großen Gemeinde; ich will mein Gelübde bezahlen, für denen, die dich fürchten. Und im 61. Psalm V. 9.: Ich will deinem Namen lobsingen ewiglich, daß ich mein Gelübde bezahle täglich.

Dies Gelübde des Dankens und Lobens setzen die Propheten weit über die andern Gelübde des Gesetzes. So dagegen unsere Papisten heutiges Tages, auf daß sie ihre Gelübde erhalten mögen, dies recht hohe Gelübde mit Füßen treten, Gottes Wort verfolgen, seinen Namen lästern und schänden, wie man für Augen siehet.

Neben dem, daß ihre Gelübde an ihnen selbst unnütz und wider Gottes Ordnung sind. Denn daß sie die Ehe verschwören, geschieht ja wider Gottes Ordnung, welcher Mann und Weib geschaffen und seinen Segen ihnen verheißen hat. Solch Gelübde aber ist darum desto ärger, weil es so viel Aergerniß, Schande und Laster mit bringet, daß es ein groß Wunder ist, daß Gott nicht längst mit allen Klöstern und Stiften ist umgegangen, wie mit Sodom und Gomorra. Wie können sie sich aber rühmen, daß sie Gehorsam schwören, weil sie aus aller ordentlichen und von Gott bestätigten Gewalt sich ziehen, und niemand, denn ihrem Abgott zu Rom, wollen gehorsam sein.

Fürs dritte ist sehr schimpflich, daß sie rühmen das Gelübde der Armuth, weil jtzund die höchste Klage ist, man ziehe ihre Güter ein, nehme ihnen ihr Einkommen. Wenn sie denn Armuth geloben und dieselbige halten, wie reimet es sich denn, daß sie klagen: man nehme ihnen ihre Güter.

Also ist es am Tage, daß unserer Papisten Gelübde wider Gottes Ordnung und Befehl, ärgerliche und erdichtete Gelübde sind. Darum soll niemand dieselbigen sich binden lassen, sondern eine freie offene Hand haben, dieselbigen zu brechen, ja wissen, wo mans nicht brechen, sondern halten will, daß man wider Gott, sein Wort und Ehre thut. Sintemal solche Gelübde ohne allen Gottes-Befehl geschehen, und wir wissen, daß man Gott mit allen Menschen-Satzungen nicht kann noch soll dienen.

Im Alten Testament waren die Gelübde recht. Ursach, sie geschahen aus Gottes Befehl. Derselbige Befehl aber gehet uns im Neuen Testament nichts an, da sollen wir bleiben bei dem Gelübde, das wir in der Taufe gethan haben: wie wir dem Teufel und allen seinen Werken entgegen sein, und uns an Gottes Güte und Gnade, so er in seinem Sohn Christo Jesu uns bewiesen, halten wollen. Bei dem Gelübde sollen wirs lassen bleiben, und haben unser Leben lang genug damit zu thun, dürfen uns nach andern nicht umsehen. Denn es bringt mit sich öffentliches Bekenntniß des Glaubens, Kreuz, rechte Anrufung im Kreuz, Hoffnung, Glaube und dergleichen anderes.

Darum, wenn die Papisten mit der Schrift einher prallen, es stehe in Mose, im Psalter, in Propheten, was man Gott gelobe, soll man halten; so antworte du und sprich: Ja es ist wahr, wenn mans verstehet vom Gelübde in der Taufe, denn das ist ein ewiges Gelübde. Aber man besehe Mosen und die Propheten fleißig, so wird man finden, daß sie von ihren Gelübden, so im Gesetze gebräuchlich, reden; dieselbigen Gelübde gehen uns nichts an. Darum binden sie uns eben so wenig, als andere Sprüche oder Gebote, die da verbieten, daß man Schweinen-Fleisch, Hasen, Aale und dergleichen nicht essen soll.

Eben also ists eine große Thorheit, daß die Papisten mit den Sprüchen, so auf jene Gelübde gerichtet, ihre Gelübde gedenken zu erhalten, von welchen Gott nie keinen Befehl ihnen gegeben hat, ja welche wider Gottes Wort und Ordnung und derhalben ärgerlich und verboten sind.

Von bürgerlichem Gelübden und wie fern dieselben zu halten, handelt man hie nicht, sondern von den Gelübden allein, die da sollen ein Gottesdienst sein. Da haben wir im Neuen Testament keines, denn das Gelübde von der Taufe. Zu andern guten Werken dürfen wir uns nicht geloben, wir sinds vorhin zu halten schuldig, sintemal Gott sie uns allen hat auferlegt und zu halten geboten. Das sei genug zum Bericht von den Gelübden.

\quellseite{221}

\part{Das vierte Buch Mose}

\chapter{Das 1. Kapitel}

\section{Vorrede}

Dies vierte Buch hat den Namen in der lateinischen Sprache, daß es heißt das Buch der Zahl. Denn nachdem die Hütten oder Kirche zugerichtet, und aller Gottesdienst angerichtet und verordnet war, befahl Gott, man sollte alle Stämme zählen, V. 3., und einem jeden seinen Ort im Lager bestimmen, auf daß die Wohnung Gottes mitten unter ihnen wäre, und sonderlich, daß die Leviten müßten, wenn man reiste, was jedermann zu verwahren hätte; solches ist im 1., 2., 3. u. 4. Kapitel verzeichnet.

\section{Zahl der streitbaren Männer in Israel}
\quellseite{222}

\chapter{Das 2. Capitel}

\section{Ordnung der Lager im Volk Israel}
\quellseite{223}

\chapter{Das 3. Capitel}

\section[Zählung der Leviten und ihr Amt. Der Erstgebornen]{Zählung der Leviten und ihr Amt. Der Erstgebornen Auslösung}
% No apparatus blocks retained on this page. See selection.json for block-by-block exclusion reasons.
\quellseite{225}

\chapter{Das 4. Kapitel}

\section{Besondere Amtsverwaltung und Zahl der Leviten}

a Gelbe Decke) hieraus scheint, daß die innersten Teppiche sind nicht bunt, sondern eine jegliche hat ihre eigene Farbe allein gehabt, welche gelb ist, die ist ganz gelb gewesen, und welche weiß, ganz weiß etc.
% No apparatus blocks retained on this page. See selection.json for block-by-block exclusion reasons.
\quellseite{227}

\chapter{Das 5. Kapitel}

\section{Reinigung des Lagers. Versöhn- und Eiferopfer}
\quellseite{228}

\section{Summarium}

Dies Kapitel ordnet, wenn einer den andern beleidiget hat, wie man den Schaden widerlegen und den fünften Theil über die Hauptsumma soll geben, auf daß der, so beleidiget ist, versühnet werde.

Darnach ist eine Ordnung, wie sich ein Mann halten soll gegen seinem Weibe, wenn sie Unzucht halben verdächt ist, V. 11. Und zeigt Gott hie sonderlich an, wie er dem Ehebruch feind sey, und ihn mit Unfruchtbarkeit des Leibes auch andern hartem Unglück strafen wolle.

\section{Votum}

Helfe der treue barmherzige GOtt, daß wir alle seine Ordnung allezeit fleißig halten, niemand Schaden thun, und daß wir Junge und Alte, in Zucht, auch Keuschheit, in und außer dem Ehestande leben, Amen.

\chapter{Das 6. Kapitel}

\section{Vorrede}

Undank ist der Menschen Untergang. Nun müssen wir alle bekennen, daß die Wohlthaten, so uns Gott aus Gnaden und Barmherzigkeit erzeiget hat und noch täglich erzeiget, so wenig können gezählet werden, als der Menschen Haar auf dem Haupt, als die Sterne am Himmel und als die Sandkörnlein am Ufer der Wasser. Wer wollte denn nicht Gott danken, loben und preisen täglich für solche seine unzählige, große und nützliche Wohlthaten? Danket demnach dem HErrn eurem Gott von ganzem Herzen, und ehret seinen Namen ewiglich, Ps. 86. Dieses vermahnet auch Sirach im 40. Kapitel, da er also saget: Singet löblich, und lobet den HErrn in allen seinen Werken, preiset seinen Namen herrlich, danket ihm und lobet ihn mit Singen und Klingen, und sprechet also im Danken: Alle Werke des HErrn sind sehr gut, und was er gebeut, das geschiehet zu rechter Zeit. Darauf höret ferner an den ersten Theil des 6. Kapitels.

\section[Gesetz der Nazaräer oder Verlobten. Formular das Volk]{Gesetz der Nazaräer oder Verlobten. Formular das Volk zu segnen}

a Auf Hebräisch heißt dies Neser, und wer sie hält, heißt Nasir, welchem nach auch unser HErr Jesus Christus Nazarenus heißt, und er der rechte Nasir ist.
\quellseite{229}

\section{Votum}

Wir danken GOtt herzlich für alle seine vielfältigen und unaussprechlichen Wohlthaten; der wolle uns ferner an Leib und Seele, zeitlich und ewiglich, segnen und behüten, um JEsu Christi, seines lieben Sohns und unsers HErrn willen, Amen.

\section{Der andere Theil des 6. Capitels}

\section{Vorrede}

Der Ursprung und Quellbrunn aller Sünden ist, wenn man Gottes Wort hinterrück wirft; wie dagegen aller guten Werke Ursprung ist, wenn man Gott Gehorsam leistet, und sein Wort Norm und Form sein läßt, darnach wir unser Leben anstellen: also, die erste Weihe zum Verderben ist, wenn man die Zucht hasset, ungestraft sein will und Gottes Wort verwirft. Wenn man aber die Sünde erkennet und sich den Geist Gottes willig fragen läßt, so ist es der erste Antritt zur Seligkeit. Derwegen danken wir heute abermal Gott billig und herzlich, daß er uns sein heilig Wort hören läßt, und bitten, er wolle uns durch den heil. Geist erleuchten und helfen, daß wir dasselbe zu täglicher Besserung und Erquickung fruchtbarlich gebrauchen mögen, Amen. Darauf höret an den andern Theil des 6. Capitels.
\quellseite{230}

\section{Summarium}

Ist eine Ordnung, wenn das Volk beisammen ist, mit was Worten der Priester sie segnen soll, V. 22. Denn wie man siehet, Gott will, daß man bei seinem Dienst thue, daß man es alles nach seinem Wort und Befehl thun und aus eigenem Gutdünken nichts fürnehmen soll.

Nun fasset aber dieser Segen drei Stücke: Daß er am ersten sagt, V. 24.: Der HErr segne dich und behüte dich, daß gehet auf den leiblichen Segen, daß Gott Leib und Leben, Nahrung und alles segnen und darnach für aller Gefahr und Schaden, so der Teufel und die arge Welt zufügen kann, behüten wolle. Denn es liegt nicht an unserer Arbeit noch Fleiß, wiewohl es beides sein soll, sondern an Gottes Segen und Behüten.

Das andere: Der HErr lasse sein Angesicht leuchten über dir, und sei dir gnädig, V. 25. Das gehet auf den geistlichen Segen, daß Gott gnädig sein und sein Wort rein und lauter geben wolle, auf daß wir ihn recht erkennen und nach seinem Wort uns halten mögen.

Das dritte: Der HErr erhebe sein Angesicht über dich, und gebe dir Frieden, V. 26. Dies gehet auf den geistlichen Schutz; denn wo Gottes Wort und Geist ist, da feiert der Teufel nicht, sicht die Leute immer an, ob er sie davon möchte abführen. Da ist denn vonnöthen, daß Gott sein Angesicht über uns hebe, das ist, daß er auf uns schaue und in solchen Anfechtungen erhalte, Friede und Sieg gebe wider den Teufel und alles Unglück.

Dabei sollen wir auch lernen, wo solcher leiblicher und geistlicher Segen und Schutz ist, daß man Gott von Herzen dafür danken soll und erkennen, daß wirs nicht von uns selbst, sondern von Gott haben. Darnach, daß man auch täglich bitten soll, daß Gott solche Gaben uns erhalten und mehren wolle.

\section{Votum}

Wir danken GOtt herzlich für alle seine Wohlthaten; der wolle uns ferner geben an Leib und Seele, was uns gut und nöthig ist, und wolle uns auch behüten für allem, was uns an Leib und Seele böse und schädlich ist, Amen.

\chapter{Das 7. Kapitel}

\section{Vorrede}

Eine jegliche Obrigkeit, vom höchsten Amt bis zum geringsten auf Erden, soll sich theilen als in drei Orte, und soll sehen, erst über sich, darnach in sich, endlich neben sich. Erstlich soll die Obrigkeit sehen über sich, auf den, der da sagt in Sprüchen Salomonis am 8. Kapitel: Mein ist Rath und That; und ist der Erb-Herr über alle Heiden, Ps. 82. Mit desselben Hülfe und Segen sollen sie ihr Amt anfangen, um den Schlüssel zur Regierung bitten, und Gottes Ehre, Wort und Kirche helfen, fördern und schützen. Zum andern soll die Obrigkeit sehen und gehen in sich selbst, das ist, in ihr Gewissen, und bedenken, wie sie regieret, was sie ordnet und wie sie ihr Amt ausführet, daß sie das Gewissen zum Zeugen habe, recht und fleißig geführten Regiments. Zum dritten soll auch die Obrigkeit neben sich sehen, das ist, eine jede auf ihre befohlenen Unterthanen, daß sie gegen den Frommen Friedrich und gegen den Bösen Ernst sei. Dies erinnert auch an einem Stück ein Theil des 7. Kapitels.

\section[Opfer und Geschenke der Fürsten zur Einweihung der]{Opfer und Geschenke der Fürsten zur Einweihung der Stiftshütte}
\quellseite{231}

a) Löffel) oder ein Köpflein oder sonst inwendig rund wie die Salbkrüglein.
% No apparatus blocks retained on this page. See selection.json for block-by-block exclusion reasons.
\quellseite{233}

\section{Summarium}

Ist ein fein Exempel, wie die zwölf Fürsten den Gottesdienst anheben, und die Ersten sind, so mit ihrem Opfer und reichlichem Geschenk zum Gottesdienst helfen. Denn gleich wie die Obrigkeit sonst in allen andern Sachen ein gut Exempel geben, und aller Tugend und Ehrbarkeit sich fleißigen soll: also ist's sonderlich der Obrigkeit Amt, sich Gottes Wort und des rechten Gottesdienst mit Ernst annehmen, und den Unterthanen ein gut Exempel fürtragen. Bei solchen Regenten ist sich alles Glücks und Heils zu versehen.

Widerum wo gottlose Obrigkeit ist, die entweder Gottes Wort nicht achtet, wie gemeiniglich geschiehet, oder auch wohl verfolget, da wird alles Unglück sich zeitlich genug finden, wie die Exempel vor Augen sind.

\section{Votum}

Wir danken Gott herzlich für unsere liebe Obrigkeit, unter deren Schutz mir Gottes Wort und den rechten Gottesdienst haben, der wolle sie gnädiglich segnen, schützen und erhalten, zu Lob und Preis seines Namens, und zur Wohlfahrt der Unterthanen, Amen.

\chapter{Das 8. Kapitel}

\section{Einweihung der Leviten}

a Vorwärts) das ist, die Seiten oder Schnauzen an den Lampen soll er gegen dem Vorhang richten, daß es für dem Leuchter und Vorhang Licht sei.

b Sünd-Wasser) entsündigen ist so viel als abspülen oder lossprechen; daher das Wasser, damit sie abgespület wurden, heißt Sünd-Wasser.
\quellseite{234}

\section{Summarium}

Das achte Kapitel hält in sich die Ordnung mit der Leviten Weihe oder Heiligung. Wie aber solchem Exempel nach der Papst mit seinen Pfaffen gehandelt hat, ist für Augen. Denn seine Pfaffen müssen alle geweyhet und gesalbet sein. So doch solches weder Christus befohlen, noch die Apostel geordnet haben. Es ist aber kein gerathen. Des HErrn Christi und der Apostel Pfaffen hatten eine Weihe, die heißt die Weihe des heil. Geistes, dieselbe hat der Papst und seine Pfaffen nicht. Darum sollen sie ohne Gottes Befehl ihnen eine eigene erdenken, die doch sehr schlecht ist, und ein geringes Ansehen auch für der Welt hat. Denn was soll es zum Amt helfen, daß man sie mit Oele schmieret. In Summa, der Papst als ein Affe, hat ohne Befehl das hiernach gethan, welches die Juden aus Befehl Gottes gethan haben, daß sie keinen zum Gottesdienst haben genommen unter fünf und zwanzig Jahren, und daß sie im fünfzigsten Jahr sind frei gewesen; wiewohl der Papst solchs mit der Freiheit auch geendet hat.

\chapter{Das 9. Kapitel}

\section{Vom Paßah, und Zeichen zum Lagern oder Fortziehen}
\quellseite{235}

\section{Summarium}

Im neunten Kapitel ist nichts Sonderliches zu merken, denn wie Moses preiset, daß die Kinder Israel mit ihrem Ziehen und Stillliegen sich nach Gottes Wort und Befehl gehalten haben. Solches sollten wir lernen, daß es die größte Kunst und Tugend ist, davon der Papst mit seiner Kirchen nichts weiß, sintemal sie alle ihre Gottesdienste aus eigenem Gutdünken ohne Gottes Befehl und Wort haben angericht.

\chapter{Das 10. Kapitel}

\section{Gebrauch der Trommeten. Der Israeliten Fortreisen}
\quellseite{236}

a Richteten) indes die Kahathiter das Heiligthum hernach trugen, waren hierbei vorhin die Gersoniter und Merariter, und richteten die Wohnung auf, daß die Kahathiten mit der Lade die Wohnung bereit fanden.

\section{Summarium}

Wie ist sonderlich zu merken das Gebet Mos., V. 35, welches allen Regenten zum Exempel ist vorgeschrieben, daß sie neben allem andern, das sie ihres Amts halben thun sollen, auch beten und bitten sollen, auf daß Gott seine Feinde stürzen und die Seinen gnädiglich beschützen und segnen wolle.

Denn obwoll die Obrigkeit mit weltlicher Gewalt die Bösen strafen und die Frommen beschützen kann, so dürfen sie doch fürnehmlich der Hülfe Gottes dazu, daß es wohl von Statten gehe, und der Teufel solches nicht hindere, wie er wohl im Sinne hat, und zu thun sich unterstehet, und ihm bei dem meisten Theil geräth, daß sie ihres Amts vergessen, die Sünde nicht strafen, wie es Gott befohlen hat, und die Tugend und die Ehrbarkeit nicht schützen noch handhaben.

\section{Votum}

Verleihe GOtt, der HErr aller Herren, daß die Obrigkeit im Gebet und allen andern Stücken ihres Amts sich recht und christlich verhalte, damit alle fromme Unterthanen ihm dafür Lob und Dank sagen mögen allezeit, immer und ewiglich, Amen.
\quellseite{237}

\chapter{Das 11. Kapitel}

\section{Vorrede}

Der heilige Chrysostomus lehret von Kraft des gläubigen Gebets also: Solche große Herrlichkeit und Wirkungen hat das Gebet in sich, daß es aus Menschen Gottes Tempel macht. Und wie Gold, Silber, Edelsteine und Marmor der Könige Häuser schön und herrlich machen: also schmückt und zieret das Gebet einen Christen, und macht aus ihm einen schönen, herrlichen, gezierten Tempel und Palast Christi, darin er wohnet; und welchen der Himmel nicht umgreift, der gibt sich in die Seele, in welcher fleißig und emsig Gebet und Seufzen ist. Ferner sagt auch Chrysostomus: Viel Heilige oder Gläubige Gottes haben durch die Macht oder Kraft ihres Gebets erhalten nicht alleine Land und Leute, sondern auch die ganze Welt. Derwegen laßt uns auch fleißig und in wahrem Glauben alle Tage beten. Solches unsers herzlichen Gebets werden wir gewiß genießen, zu unsrer geistlichen und leiblichen Wohlfahrt, um Jesu Christi unsers Herrn Verdienst und Vorbitt willen. Darauf höret nun an das 11. Kapitel.

\section{Das lüsterne Volk wird gestraft}

\section{Summarium}

Dies ist ein schrecklich Exempel, wie die Juden des Himmelbrods in der Wüsten überdrüssig worden sind, und wünschen um des Fleisches und andere Essens willen, daß sie wieder in Egypten wären, haben vergessen aller Noth und Jammers, den sie darin erlitten haben, und Gott sie davon erlöset, V. 1. Solche Unbarkeit straft Gott mit Feuer vom Himmel. Darum sollen wir lernen, daß wir um zeitlichen Mangels willen nicht auch der herrlichen und geistlichen Güter vergessen, wie jetzo sehr gemein ist. Denn daß wir der Sünden los und zur rechten Erkenntniß Gottes und wahren Gottesdienst kommen sind, solches halten viel Leute für ein Geringes und Schlechtes; und klagen dieweil über den zeitlichen Mangel, daß
\quellseite{238}
mehr Mißwachsung der Früchte, Krankheit, Krieg und anders ist, denn unter dem Papstthum, so doch solches nicht des Evangelii, sondern unsere Sünde, Unbußfertigkeit und greuliche Undankbarkeit schuldig ist.

Es ist auch ein fein Exempel, wie Gott Mosen und des Volks Gebet erhöret, und mit der Strafe aufhöret.

\section{Summarium}

Ist ein fein Exempel, wie Gott sich bewilliget dem Volk, ihrem Begehr nach, Fleisch zu geben. Aber Moses achtet, es sei unmöglich. Denn er rechnet nur den Haufen, und denket nicht, was Gott vermöge zu thun. Darum sollen wir an solchem Exempel lernen, wie leicht es Gott ist, uns zu geben und zu thun, was uns gleich unmöglich dünket. Denn Gottes Hand ist nicht verkürzt, das ist, Gott ist allmächtig; ob bei uns gleich etwas unmöglich ist, sollen wir auf seine Gewalt und Allmächtigkeit sehen, und uns derselben in allerlei Noth und Mangel trösten.

Gott strafet hie zum andernmal die Juden, um ihrer Undankbarkeit und Murrens willen wider Mosen und den HErrn.

\section{Votum}

GOtt, der HErr, behüte uns für Undankbarkeit, Unglauben, Ungeduld, und allen andern Sünden, und verleihe uns allen Gnade, daß wir täglich fleißig und fruchtbarlich beten, und christlich leben, Amen-

13*
\quellseite{239}

\chapter{Das 12. Kapitel}

\section{Vorrede}

Gott ist tausendmal geneigter die Sünde zu vergeben, denn wir darum zu bitten, und siehet nicht darauf, wie groß die Sünden, und wie viel derselben sind. Ja er wendet auch ab, oder lindert gnädiglich die mit Sünden verdieneten Strafen. Hieher gehört, daß Gott selber saget durch den Propheten Esaiam im 65. Capitel: Es soll geschehen, ehe sie rufen will ich antworten, wenn sie noch reden, will ich hören. Und im 1. Kapitel spricht Esaias: Wenn eure Sünde gleich blutroth ist, soll sie doch schneeweiß werden, und wenn sie gleich ist wie Rosinfarbe, soll sie doch wie Wolle werden. Deß sollen sich trösten alle bußfertige Sünder und betrübte Kreuzträger; sie sollen auch Buße thun, und also Gottes väterlicher, gewisser Hülfe in beständiger Geduld erwarten. Davon höret an das 12. Kapitel.

\section{Die murrende Mirjam wird aussätzig}

a (Geplagter) elender, der viel leiden mußte, Ps. 132, 1. Gedenke Davids und alles seines Leidens, Ps. 18, 36. Ps. Sal. 18, 12. Ante gloriam passio. Das ist, vor der Herrlichkeit Leiden.

\section{Summarium}

Ist eine Historia, wie Aaron und Mirjam Mose Schwester sich ihres Standes und Wesens überheben, und Mosen verachten, V. 2. Solche Hoffart gefiel Gott übel, darum straft er sie beide mit harten Worten und Mirjam mit dem Aussatze, auf daß wir lernen und demüthigen, und niemand verachten. Sonderlich aber, daß wir die in Ehren halten, so Gottes Amt führen und mit dem Wort umgehen. Darnach ist ein fein Exempel, wie Aaron seine Sünde erkennet, und Mosen bittet, daß er gegen Gott solche Sünde abbitten und seiner Schwester wieder vom Aussatz helfen wolle, V. 12. Denn das gehöret darauf, wer unrecht gethan hat, daß ers ihm lasse leid sein, nimmer thue, und um Gnade bitte, so will Gott die Strafe auch gnädig abwenden, wie man hie siehet.

\section{Votum}

Der treue, barmherzige GOTT wolle uns auch gnädig sein, und nicht strafen, wie wirs alle verdienet haben, und wolle uns regieren, daß wir ihn lieben, und alle die ehren, denen wir Ehre schuldig sein, Amen.

\chapter{Das 13. Kapitel}

\section{Vorrede}

Gott ist allmächtig, der kann alles thun, und mehr, denn wir in unsern Nöthen glauben, und
\quellseite{240}
ihm trauen können. So ist er auch wahrhaftig, was er zusagt, das hält er fest und erfüllets endlich. Denn von Gottes Allmacht sagt David im 135 Psalm: Alles was der HErr will, das thut er, im Himmel, auf Erden, im Meer und in allen Tiefen. Und St. Paulus Ephes. 3.: Gott kann überschwänglich thun, über alles das wir bitten oder verstehen. Von Gottes Wahrheit aber spricht auch David Psalm 89.: Gnade und Wahrheit sind für des HErrn Angesichte. Und St. Paulus Röm. 3.: Gott ist wahrhaftig und alle Menschen falsch. Derwegen wollen wir in diesen gefährlichen Zeiten Gott vertrauen, und ihm uns und die Unsern befehlen, er wird uns wohl schützen und erhalten, und sollt es allen Teufeln in der Hölle leid sein. Dies unterrichtet auch das 13. Kapitel.

\section{Von Kundschaftern des Landes Canaan}

a Escol heißt eine Traube, daraus wird der Traubenbach.
\quellseite{241}

\section{Summarium}

Ist eine Historia, wie Moses zwölf Männer ausschickt, das Land Canaan zu besichtigen, V. 4., und dieselbigen das Volk verzagt machen, daß sie daran zweifelten und besorgten, sie würden den Leuten zu schwach sein, und das Land nicht mögen einnehmen, V. 32. Solches ist ein fein Exempel, in welchem man siehet, wie die Vernunft, so bald sie von der Zusagung Gottes abfällt und sich nicht daran hält, daß Gott allmächtig und gnädig sei, an ihr selbst verzagen und verzweifeln muß.

Darum sollen wir uns lernen trösten der Hülfe und Verheißung Gottes, es stehe unserthalben wie es wolle, daß doch unser Vertrauen und Trost allein auf Gottes Wort, Güte und Allmächtigkeit stehe; so soll es nicht Noth haben.

\section{Votum}

Der allmächtige und treue GOtt stärke unsern schwachen Glauben, und helfe, daß wir ihm in Noth und Tod beständig trauen, und zeitlich und ewig an Leib und Seele erhalten werden, Amen.

\chapter{Das 14. Kapitel}

\section{Vorrede}

Es ist uns nichts auf Erden und in dieser Welt, (was das äußerliche Leben betrifft) höher und größer vonnöthen, denn die wahre brüderliche Liebe und Treue, als deren wir so wenig gerathen mögen, als der natürlichen Sonne. Denn, gleich wie wir ohne die Sonne nicht mögen leben, also müßten wir auch mit einander zu Grunde gehen, wo die Liebe und Treue aufhörete, und dagegen alle Untreue und Ungerechtigkeit überhand nähme. Diese Liebe sollen wir an unsern Nächsten üben, fürnehmlich aber mit dem Gebet für ihn in seiner Gefahr und Noth.

Von diesem der h. Chrysostomus also saget: daß einer für sich bittet, dazu treibet ihn seine eigene Noth, daß er aber für seinen Nächsten bittet, dazu zwinget ihn die Liebe. Es ist aber für Gott ein lieblicher Gebet, dazu uns die brüderliche Liebe, denn dazu uns die Noth treibet. Dies soll auch erinnern das 14. Kapitel.

\section{Das murrende Volk wird geschlagen}
\quellseite{242}

\section{Summarium}

Dies ist eine feine Historia, in welcher man siehet erstlich, wie bald der Pöbel zu bewegen ist, daß er von Gott und seinem Wort abfällt. Sintemal die Juden aller Wunderwerke vergessen, so Gott bei ihnen gethan hatte, als könnte oder wollte ihnen Gott nimmer helfen, und begehren wieder in Egypten, V. 4.

Zum andern ist hie eine schöne Vermahnung wider solches Verzagen und Unglauben, daß man den Gedanken nicht Raum geben, sondern mit Ernst widerstehen soll, die uns zum Unglauben treiben, wie die Juden hie, die da zaghaftig werden, daß sie gedenken, die Feinde sind ihnen zu stark und mächtig, V. 11.; sondern daran sollen wir gedenken in aller Noth, daß Gott bei uns sein, und mit uns wider unsere Feinde streiten will. Solches sollte uns mehr trösten, denn alle Welt mit alle ihrem Vermögen uns schrecken kann. Aber wer also auf Gott und seine Güte sich und andere tröstet, der muß gewarten, daß man ihn steinige, wie die Juden hie Mose zu thun gedachten, V. 13. Aber Gott wachet über die Seinen und errettet sie.
\quellseite{243}

\section{Summarium}

Ist ein fein Gebet, daß Gott schonen, und das Volk um der Sünden willen nicht gar vertilgen wolle. Dies Gebet soll uns eine Form sein, wie wir unsere Fürbitte gegen Gott sollen stellen, auf daß man in aller Fürbitte und allem Gebet erstlich Gottes Ehr und Namen bedenke. Zum andern, daß man auf Gottes Verheißung sehe, und sich seiner Barmherzigkeit tröste, weil er als ein gnädiger Gott will gerühmet sein, daß er solche Gnade gegen uns beweisen wolle. Solch Gebet kann nicht fehlen. Darum saget Gott Mose zu, er wolle vergeben. Aber gleichwohl dräuet er, und zur Warnung: es soll ihr keiner ins Land Canaan kommen, und die Kinder, so solcher Sünde unschuldig waren, sollen der Väter Schuld mit tragen vierzig Jahre; auf daß wir ja sehen, wie Gott der Sünde feind sei, und lernen uns dafür hüten, auf daß nicht unsere Kinder derselben nach uns entgelten müssen.

\section{Summarium}

Ist erstlich ein Exempel, wie Gott die Kundschafter, so zu solcher Sünde dem Volk hatten Ursach gegeben, mit dem Tode strafet. Caleb aber und Josua blieben lebendig, warum daß sie solcher Sünde gewehret, und auf Gottes Hülfe das Volk vertröstet hatten. Denn Gott will beides thun, böse Buben strafen, und die Frommen segnen und erretten.

Zum andern ist noch ein Exempel hie, in welchem wir sehen, wie die menschliche Natur so gar Gott und seinem Willen entgegen ist. Da die Juden Gottes Hülfe am meisten sich sollten getröstet haben, und ins Land Canaan ziehen, verzagten sie, und wollten mit nichten hinein. Jetzt, so es ihnen Gott verbeut und saget, sie sollen vierzig Jahre in der Wüste bleiben, wollen sie auch nicht, sondern wollen wider Gottes Willen ins Land und kriegen, V. 40. Aber weil sie es wider Gottes Befehl und Willen fürnehmen, werden sie darüber geschlagen, V. 45. Auf daß wir lernen, Gott und seinem Worte folgen, sonst soll weder Glück noch Heil bei uns sein.

\section{Votum}

GOtt, unser lieber Vater, um JEsu Christi, unsers Mittlers, willen, durch den heiligen Geist, helfe, daß wir uns an andern Leuten spiegeln, Buße thun, und Gottes Wort allzeit gehorsam sein mit Thun und Lassen, Amen.

\chapter{Das 15. Kapitel}

\section{Vorrede}

Es ist nichts leichter, als die Sünde, und nichts schwerer als die Strafen, so auf die Sünde den Unbußfertigen folgen. Derwegen wer gesündiget hat, der thue Buße, und hüte sich ferner für willentlichen muthwilligen Sünden. Denn auf viele und große muthwillige Sünden folgen viele und große verderbliche Strafen. Denn um die Sünden es so ein arg Ding ist, daß dieselbigen von der Erde gen Himmel schreien, und für Gottes Angesicht gleich treten, und die Rache und Strafe erfordern, daß er endlich seinen Zorn darauf muß ergehen lassen, wo man sich nicht durch rechte Buße zu ihm bekehret. Darauf höret auch an das 15. Kapitel.
\quellseite{244}

\section{Entheiligung des Sabbaths wird gestraft}
\quellseite{245}

a Frevel) das ist die Sünde, so nicht will Sünde sein, sondern recht haben, wie die Ketzer- und Rotten-Sünde, welche Gottes Wort und Gesetze ändert, und nach ihrem Willen deutet.

\section{Summarium}

Ist am ersten ein ernstes Gebot, daß Gott die Sünde, so aus Unwissenheit geschehen, will versühnen lassen; aber muthwillige Sünden, so fürstlich geschehen, will er nicht ungestraft lassen und sagt, V. 22: Wer solches thue, der habe den HErrn geschmähet, sein Wort verachtet, und sein Gebot zerstöret. Darum sollen solche frevele und muthwillige Sünder schlecht ausgerottet werden. Solche Predigt soll man wohl merken, und sich desto fleißiger vor Sünden hüten.

\section{Votum}

Gott, der Brunnquell aller Gnaden, erbarme dich unser, und regiere uns, daß wir nach seinen Geboten in allen Ständen christlich leben, handeln und wandeln, damit wir zeitlichen und ewigen Strafen entfliehen, Amen.

\chapter{Das 16. Kapitel}

\section{Vorrede}

Gott fürchten, der Obrigkeit und Seelsorgern gehorsam sein, und den Nächsten lieben, sollen alle, die zeitlich gesegnet und ewig erhalten sein wollen, um Christus willen aus Gnaden; dagegen, wer zeitlich und ewig Unglück, an Leib und Seele, Gut und Nahrung begehret, der sündige nur und verharre in Sünden; Herzeleid genug wird er zur Strafe haben in diesem und künftigen Leben. Weil aber die Sünde über unser Fleisch und Blut ist Meister worden, wütet und tobet derselben für und für in uns, daß ihr leicht mag eine Gelegenheit gegeben werden, so bricht sie heraus in allerlei greuliche Laster. So laßt uns täglich und fleißig beten, Gott wolle uns regieren, daß die Sünde über uns nicht herrsche, und uns in zeitlichen und ewigen Jammer nicht verführe. Darauf höret ferner an das 16. Kapitel.

\section{Gott übet Rache an der aufrührischen Rotte}
\quellseite{246}

\section{Summarium}

Ist eine Historia, wie etliche Leviten ihnen einen Anhang machten wider Aaron, und wollten auch Priester sein, und eben so viel Gewalt haben als Mose und Aaron, V. 1. Aber wie es geraten ist, wird hernach folgen; auf daß jedermann lerne, bei dem zu bleiben, dazu ihn Gott berufen hat, und desselbigen mit Fleiß auszuwarte, und nicht anders noch Höheres begehre. Denn solches ist außer dem Beruf gegangen; darum weil es wider Gottes Willen ist, kann kein Glück dabei sein, wie man siehet, daß alle Aufrührer ein bös Ende nehmen.
\quellseite{247}

\section{Summarium}

Ist eine Historia, wie Gott um der Aufrührer willen das ganze Volk wollte vertilgen, aber Moses bittets ihm ab, daß allein die Aufrührer gestraft werden, V. 22. Und hie stehet man, wie eine greuliche Sünde es ist, außer dem Beruf gehen und sich selbst ohne Beruf eindringen zum Kirchen-Amt, weil Gott Weib und Kind mit nimmt, und die Erde alles verschlingen läßt, V. 32. Solch Urtheil Gottes sollen wir fleißig merken, und uns darnach bessern lernen. Darum befiehlet Gott, man soll die Rauchpfannen aufheben zum Gedächtniß, daß jedermann bei seinem Beruf bleibe, und ohne Beruf nichts fürnehme.

\section{Summarium}

Ist ein schrecklich Exempel, in welchem man siehet, wie gar nichts andern Gottlosen hilft, weder Strafe noch andere Vermahnung. Denn Gott selbst hatte sich erzürget mit der Strafe wider die Aufrührer, daß die Erde sie sammt ihren Weibern und Kindern verschlungen hatte. Solches ist allein Gottes Werk gewesen, wie alle Vernunft muß zeugen. Aber das blind verstockte
\quellseite{248}
Volk empöret sich auf ein neues wider Mosen und Aaron, und gibt ihnen die Schuld, als hätten sie jene umbracht; wollens derhalben strafen.

Aber Gott bedeckt Mosen und Aaron mit der Wolke, daß ihnen niemand nichts thun kann, V. 42., auf daß man sehe, wie Gott die Seinen schützen und behüten wolle. Die Bösen aber strafet er, daß ihrer vierzehn tausend und sieben hundert umkommen. Aber gleichwohl ist’s tröstlich, daß Gott auch solchen Zorn ihm läßt abbitten, und der Plage durch Gebet gewehret wird, V. 48. Solche Exempel sind sonderlich wohl zu merken, auch der Ursach halben, daß man sehe, wie Gott gnädig ist, und das Gebet erhöret, und den Zorn will fallen lassen.

\section{Votum}

Der gnädige und barmherzige GOtt erhöre auch unser Gebet, regiere uns, daß wir unserer Obrigkeit und Seelsorgern Gehorsam leisten, und lindere uns die wohlverdienten Strafen, um JEsu Christi, unsers HErrn willen, Amen.

\chapter{Das 17. Capitel}

\section{Vorrede}

Fürwitzige Leute sind unglückselige Leute. Derwegen niemand andern in ihre Aemter greifen, sondern ein jeder soll seines Berufs warten, in der Furcht Gottes, und soll Gott bitten, daß er seinen Segen und glückselige Verrichtung dazu geben wolle. Also behält ein jeder ein gut Gewissen, und also wird die Arbeit seines Berufs gebenedeiet, daß sie Gott zu Ehren, ihm selber zu Nutz und dem Nächsten zur Wohlfahrt gelanget. Derwegen ein jeglicher Christ täglich in seinem Beruf betrachten soll folgende des h. Apostels Pauli Regeln und Unterricht, 1 Cor. 7.: Ein jeglicher, lieben Brüder, worinnen er berufen ist, darinnen bleibe er bei Gott. Und 1 Thess. 4. saget er: Wir ermahnen euch lieben Brüder, ringet darnach, daß ihr stille seyd und das Eure schaffet. Darauf höret auch an das 17. Kapitel.

\section{Aarons Priesterthum wird durch den grünenden Stab bestätiget}

\section{Summarium}

Ist eine feine Historia, wie Gott mit einem neuen Wunderwerk an der Ruthe Aarons bezeuget, daß das Priesteramt allein bei Aaron soll bleiben, V. 8., auf daß niemand sich ferner desselbigen anmaßen, und um solches Frevels willen umkommen sollte. Darum heißt er auch die Ruthe Aarons aufheben und bewahren, V. 10. Solches alles dienet uns dahin, daß wir uns fremder Aemter und Berufs entschlagen, und sonderlich mit Kirchendiensten ein jeder seines Berufs sich halten, und denselbigen mit Fleiß auswarten soll.

\section{Votum}

Helfe der treue GOtt, daß wir Hörer und Thäter seines Worts sein und bleiben, und segne die Arbeit unsrer Hände, zu Lob seinem heiligen Namen, und unserm Nächsten zur Freude und Wohlfahrt, Amen.
\quellseite{249}

\chapter{Das 18. Kapitel}

\section{Vorrede}

Die uns das Himmelbrod fürtragen und mittheilen, die sollen wir mit leiblichem Brod und Nahrung versorgen. Denn dieses hat Gott ernstlich befohlen, wie der h. Apostel Paulus bezeuget, 1 Cor. 9. und saget: Wisset ihr nicht, daß, die da opfern, essen vom Opfer? Und die des Altars pflegen, genießen des Altars? Also hat auch der Herr befohlen, daß, die das Evangelium verkündigen, sollen sich vom Evangelio nähren. So soll auch das darum geschehen, damit also Gottes Wort erhalten werde. Denn wenn man das heilige Evangelium verachtet und die Diener desselben Noth leiden läßt, so nimmt Gott sein Wort endlich hinweg, das ist gewiß. Wehe aber denen Leuten, zeitlich und ewiglich, die Gottes Wort nicht haben. Daneben höret an das 18. Kapitel.

\section{Amt und Unterhalt der Priester und Leviten}
\quellseite{250}

a) Bund) im Hebräischen heißt es ein Salzbund: daß, wie das Salz erhält das Fleisch unverweslich, also soll auch dieser Bund unverrücklich sein, so redet die Schrift auch 2 Chron. 13, 5. Gott hat das Reich David gegeben, und den Seinen mit einem Salzbund.

b) Fülle) das ist Most, Wein, Oel und dergleichen.

c) Besten) hebräisch heißt es das Fette; daher auch die Reichen die Fettern hießen, das ist die Besten im Volk, Ps. 17, 10. und 22, 12.

\section{Summarium}

Ist eine Ordnung der Personen, so des Kirchendiensts warten sollen, V. 5., daß das Priesteramt bei Aaron und seinen Kindern bleiben, die andern Leviten aber, ein jeder seinem Beruf nach, sonst in andern zum Kirchendienst helfen sollen, V. 6., und daß solchen Kirchenpersonen das ganze Volk Versorgung thun, und ihre Nahrung ihnen schaffen soll. Denn außerhalb des Zehnten, alles was von Vieh und Früchten des Feldes geopfert ward, an solchem hatten die Priester und Leviten auch ihren Theil.

Solches ist uns alles zum Exempel fürgeschrieben, daß jedermann mit den Seinen dazu helfen soll, auf daß die Kirchendienste recht bestellt und erhalten werden mögen; wie Christus Luc. 10. selbst spricht, da er redet mit seinen Aposteln, die er aussandte das Evangelium zu predigen: Ein Arbeiter ist seines Lohns werth. Und St. Paulus, Gal. 6.: Man soll denen das Zeitliche nicht versagen, welche das Geistliche den Leuten mittheilen.

\section{Votum}

GOtt erhalte uns sein heiliges Wort, und regiere uns, daß wir nach demselben unsern Beruf und Aemter christlich verrichten, und uns gegen seine treuen Diener allzeit auch dankbarlich erzeigen, Amen.

\chapter{Das 19. Capitel}

\section{Vorrede}

Der edelste und beste Schatz aller Menschen hie auf dieser Welt ist Vergebung der Sünden. Der edelste und beste Schatz aber in jener Welt wird sein das ewige freudenreiche Leben. Dieser Schatz werden wir theilhaftig durch den Glauben an Jesum Christum unsern Heiland, der uns solche durch sein Blut und Tod erworben hat. Dieses lehret der h. Apostel Paulus, als da er sagt, Col. 1. und Eph. 1.: An Christo haben wir (vernehmt, die Gläubigen) die Erlösung durch sein Blut, nämlich die Vergebung der Sünde. Und denn daß er spricht Röm. 6.: Die Gabe Gottes ist das ewige Leben in Christo Jesu unserm HErrn, vernehmt der uns solches verdient hat, und derer allein die Gläubigen theilhaftig werden. Demnach laßt uns Buße thun und an Christum glauben, so haben wir
\quellseite{251}
Vergebung der Sünden, und bekommen das ewige Leben, daß ist je gewißlich wahr. Dies soll auch erinnern das 19. Kapitel.

\section{Von der röthlichen Kuh und dem Sprengwasser}

\section{Summarium}

Dies Kapitel redet von einer sonderern Ceremonie, welche die Juden mit Sprengwasser von dem Sündopfer gehabt haben, sich damit zu reinigen oder zu entsündigen, V. 4. Was nun aber für Reinigkeit durch solches Wasser die
\quellseite{252}
Juden gehabt, meldet die Epistel an die Ebräer im 9. und nennet's eine leibliche Reinigkeit. Aber durch das Blut Christi, spricht er, der sich selbst ohne Wandel durch den heiligen Geist geopfert hat, werden unsere Gewissen gereiniget von den todten Werken, zu dienen dem lebendigen Gott. Darum ist es nur eine äußerliche Ceremonie gewesen, welche Gott befohlen hat. Aber die Deutung ist sehr fein gewesen: denn die rothe Kuh so zu solchem Sündopfer gehört, V. 2., hat den Sohn Gottes unsern HErrn Christum bedeutet, der ein Opfer für der ganzen Welt Sünde worden ist. Daß man aber über die Asche der verbrannten Kuh ein frisches lebendiges Wasser gegossen, und darnach die Menschen sich damit besprengt haben, V. 9., solches ist ohne Zweifel die Deutung, daß durch die heilige Taufe das Leiden unsers lieben HErrn Jesu Christi, sammt Vergebung der Sünden und ewigem Leben, ausgetheilt solle werden, wie Christus spricht, Marc. 16, 16.: Wer glaubt und getauft wird, der wird selig. Solche Taufe ist das rechte Weihwasser, das wir im neuen Testament haben, und durch jenes bei den Juden ist gedeutet worden.

Aber der Papst hat ihm ein anderes Weihwasser in seiner Kirche erdacht, damit, überredet er die Leute, soll man täglich Sünde abwaschen. Aber wo ist Gottes Wort und Befehl? In Summa, es ist nichts Gutes an allem was der Papst je und je in der Christenheit ohne Gottes Wort hat eingeführet. Man weiß aber, was für Abgötterei, Aberglaube, Zauberei und andere Gotteslästerung aus dem geweihten Wasser gefolget. Darum es billig nicht ein geweihetes, sondern ein verfluchtes und heillos Wasser ist, dadurch das Blut Christi gelästert, und die Leute schändlich betrogen sind, als fürchte sich der Teufel dafür. Wer nun wissen will wo es her komme, der hörets hie in Mose: Aber was gehet solches den Papst und uns an, was die Juden für Weise gehabt haben.

\section{Votum}

Gott, der HErr, steure allen abgöttischen Greueln, und helfe, daß wir uns des Blutes unsers HErrn Jesu Christi, und unserer Taufe mit Freuden trösten, in aller Noth und wider den Tod, im Leben und Sterben, Amen.

\chapter{Das 20. Kapitel}

\section{Vorrede}

Der Unglaube ist die größte Sünde, ja der Unglaube ist ein Quellbrunn und Ursprung aller Sünden. Denn wie die Haupttugend ist der Glaube: also ist wiederum die Hauptsünde der Unglaube. Und alle Sünden erlangen Vergebung, allein der Unglaube nicht, und verdammt keine Sünde, wo nicht der Unglaube an Christum dazu kömmt. Derwegen prüfe ein jedes unter uns sein Herz; und da es des Unglaubens wegen nicht auch wohl gestraft werden, so höre es Gottes Wort, und brauche die heiligen Sacramente, wie solche Christus geordnet; bete auch täglich: Ich glaube lieber HErr, hilf meinem Unglauben, Marc. 9., und: HErr, stärke uns den Glauben, Luc. 17. Dies erinnert auch das 20. Kapitel.

\section{Wasser aus dem Felsen. Mirjam und Aaron sterben}

a. Quia debebit semel percutere. Das ist, weil er nur einmal schlagen sollte.
\quellseite{253}

b Dubitatio est peccatum. Sed significat mysterium, quod populus legis non potest per suum Doctorem salvus fieri, et quod sub Christi tempus Moses incredulus, idest, populus, corruit. Das ist, Zweifel ist Sünde. Aber es zeigt ein Geheimniß an, daß das Volk des Gesetzes nicht kann durch seinen Lehrer selig werden, und daß zur Zeit Christi der ungläubige Moses, das ist, das Volk, dahinfallen wird.

\section{Summarium}

Ist eine Historia, wie das Volk sich wieder versündiget mit Murren wider Mosen und Aaron, darum daß sie kein Wasser hätten, und Mose und Aaron um solches Murrens willen besorgen. Gott möchte sein Wort nicht vollstrecken, und sein Wasser aus dem Fels kommen lassen, V. 3. Solchen Unglauben zeigt Gott für eine große Sünde an, und drohet Mose und Aaron, sie sollen solcher Sünde halben auch nicht in das gelobte Land kommen, V. 12. Denn Gott will das Vertrauen bei uns haben, daß wir uns seiner Güte und Gnade sollen trösten, wenn wir gleich nicht werth sind noch verdient haben. Er will nicht, daß wir wie Moses und Aaron denken sollen, er werde denn allererst gnädig sein, wenn wir fromm und solcher Hülfe würdig sind. Daß aber Moses ins gelobte Land nicht kömmt, hat seine Deutung; denn Gesetz und Evangelium, Werk und Gnade, sind wider einander; so ist das gelobte Land ein Bild der ewigen Seligkeit, da gehöret nicht Moses, sondern Christus zu, nicht Verdienst, sondern Gnade und Barmherzigkeit.

\section{Votum}

Der barmherzige GOtt mehre unsern Glauben, und verleihe, daß wir den durch gute Werke leuchten lassen, und gebe uns das Ende des Glaubens, das ewige Leben, Amen.
\quellseite{254}

\chapter{Das 21. Kapitel}

\section{Vorrede}

Das Bild der ehernen Schlange, welche Moses nach Gottes Befehl in der Wüste erhöhet hat, soll allen Christenherzen wohl bekannt sein. Derwegen laßt uns dessen Deutung und Erklärung anhören. Dieses Bild aber erkläret der HErr Christus selbst, daß er dadurch für gebildet und gemeinet sei, wie zu lesen Johannes am 3.: Denn also saget er: (Gleich wie Moses in der Wüste eine Schlange erhöhet hat, also muß des Menschen Sohn auch erhöhet werden, auf daß alle, die an ihn glauben, nicht verloren werden, sondern das ewige Leben haben. Allhier ist zu merken, worinnen denn die erhöhete Schlange eine Figur und Deutung ist Jesu Christi, nämlich darinnen: Erstens, gleich wie die eherne Schlange wohl allerdings eine Form und Gestalt hat einer rechten, wahrhaftigen Schlange, aber doch ist sie ohne alles tödtliche Gift gewesen: also hat Jesus Christus wohl aller Dinge menschliche Natur und Affecten an sich, doch ist er ohne Sünde und Gift der Sünden, wie St. Paulus sagt 2 Cor. 5.: Der von keiner Sünde wußte, Jesus Christus, den hat Gott für und zur Sünde gemacht, auf daß wir in ihm würden die Gerechtigkeit, die für Gott gilt. Und Röm. 8.: Gott sandte seinen Sohn in der Gestalt des sündlichen Fleisches, und verdammte die Sünde durch Sünde, auf daß die Gerechtigkeit, im Gesetz erfordert, in uns erfüllet würde. Darnach in dem Stücklein ist auch die erhöhete Schlange ein Vorbild Jesu Christi: daß gleich wie der, der von den rechten feurigen Schlangen gebissen war, und doch diese erhöhete eherne Schlange mit fröhlichem Vertrauen hat können ansehen, der ist von dem tödtlichen Gift entledigt und genesen: also wer diesen am Kreuz und durch seine Auferstehung erhöheten Jesum Christum im Glauben wird ansehen, dem soll der alten Schlangen, das ist des Teufels und Todes giftiger Biß, sogar nicht schaden, daß auch alle, die an ihn glauben, nicht sollen verloren werden, sondern haben das ewige Leben. Doch ist auch heilsam dies zum Unterschied zu wissen: Moses Schlange hat die Kraft zu heilen und Gift zu vertreiben, nicht ab intra, sed ab extra, das ist, nicht in und von sich selbst, sondern von dem lebendigen Wort Gottes, der solche Kraft dadurch hinein gesprochen hatte; Jesus Christus aber hat die lebendige Kraft in und von sich selbst, und ist uns fürgestellet zu einem Gnadenstuhl durch den Glauben in seinem Blut, Röm. 3.

\section[Der Biß feuriger Schlangen, nach Anschauung der ehernen]{Der Biß feuriger Schlangen, nach Anschauung der ehernen Schlange geheilet}

a Harma) heißt ein Bann.

b Feurige) darum heißen sie feurige, daß die Leute von ihnen gebissen, durch ihr Gift feuerroth wurden, und für Hitze starben, wie an der Pestilenz oder Carfunkel.

c Arnon) ist derselbe hohe Fels, der in Besupha, das ist in Wolken, Winden und Wetter stehet. Unten fließen Bäche, und er an denselben seinen Bächen hin währet bis gen Ar.
\quellseite{255}

d Stäbe mögen auch die heißen die Fürsten selber; wie im Zeichen der Könige in Egypten ein Stab heißt, darum, daß sie das Volk regieren, wie ein Hirte das Vieh mit seinem Stabe.

e Töchtern) das ist, die Dörfer und Flecken um die Stadt gelegen.

\section{Summarium}

Ist eine feine Historia, wie das Volk wider Mosen murrete, und verdrossen und unwillig ward, daß sie so lange in der Wüste bleiben mußten, V. 5. Solchen Unwillen strafet Gott mit giftigen Schlangen; aber da Moses für sie bittet, befiehlt Gott, daß man eine eherne Schlange machen soll, und verheißet, wer die ansehe, dem soll das Gift nicht schaden, V. 9.

Solche Historia deutet unser lieber HErr Christus Joh. 3, 13. auf sich und sagt, die eherne Schlange hab ihn bedeutet, denn gleich wie dieselbe einer rechten Schlange gleich gewesen, aber doch nicht allein sein Gift gehabt, sondern auch wider das Gift geholfen hat. Also ist Christus ein Mensch wie andere Menschen, aber er hat keine Sünde, denn er ist Gott, und vom heiligen
\quellseite{256}
Geist empfangen, und von einer Jungfrau geboren. Er hilft aber wider das Gift der Sünden allen denen, so ihn mit Glauben am Holz des Kreuzes ansehen, daß er für ihre Sünden bezahlet und genug gethan habe.

Sie hebt sich auch die Historia an, wie die Kinder von Israel wider die Heiden jenseit des Jordans gestritten, und ihre Länder eingenommen haben.

\section{Votum}

Lob, Ehr, Preis und Dank sei dem HErrn JEsu Christo für alle seine Wohlthaten, die er uns durch seinen Gehorsam und Kreuz erworben hat; der wolle uns an Leib und Seele heilen, und alle selig machen, Amen.

\chapter{Das 22. Kapitel}

\section{Vorrede}

Der Geiz wird von den Alten genennet eine Stiefmutter der Gerechtigkeit, und eine Hauptsach aller Laster. St. Paulus nennet den Geiz eine Abgötterei und eine Wurzel, das ist, einen Ursprung und Quellbrunn alles Uebels, 1 Tim. 6. Die Alten haben denselben Geiz abermals genennet eine unersättliche Trunkenheit, und ein solch Laster, das bei und in den Menschen nimmer alt werde; denn gemeiniglich, je älter einer wird, je karger und geiziger er wird, und je näher die Leute zur Herberge (das ist, zum Grabe) haben, je ernstlicher sie nach der Zehrung pflegen zu trachten. Derwegen sich niemand den Geiz einnehmen soll lassen, sondern ein jedes soll den meiden und meiden, in Betrachtung, daß Gott für uns sorgen, und keinen Mangel uns will leiden lassen. Denn Gottes Ordnung ist für und für, wo er das Leben gibt, daß er auch schaffet, daß man es erhalten könne. Darauf höret auch an das 22. Kapitel.

\section{Bileam soll den Israeliten fluchen. Seine Eselin redet}
\quellseite{257}

\section{Summarium}

Ist eine Historia, wie der Moabiter König den Propheten Bileam holen ließ, daß er den Juden fluchen sollte, V. 5. Sie ist ersichtlich zu merken, daß es scheinet, als sei Bileam ein rechter Prophet gewesen, der nicht mit Zauberei umgangen, sondern aus dem Geist Gottes geweissaget habe; wie aus seiner Weissagung unten wohl gespüret wird, und hie klar stehet, daß Gott mit ihm geredet habe.

Zum andern, daß er Gott zweimal fraget, V. 11. 19., und Gott ihm erstlich heißt, er soll nicht ziehen, V. 12., und zum andern Mal heißt, er soll ziehen, V. 20. Das ist also zu verstehen, daß Bileam am ersten Befehl und an der ersten Antwort sich sollte haben genügen lassen, und weiter nichts gefraget haben. Denn Gott ist in seinem Wort beständig. Darum was er einmal sagt, dabei soll man bleiben. Daß aber Bileam den ersten Befehl verachtet, und zum andernmal fraget, thut er unrecht, und versucht Gott um seines Geitzes willen. Darum versucht ihn Gott wieder, und heißt ihn ziehen. Denn wer sich Gottes Wort nicht will warnen lassen, der mag ihm selbst zum Aergerniß und Schaden fortfahren, und sehen, wie es ihm endlich gerathe. Gott soll man nicht versuchen. Es heißt aber Gott versuchen, wo man bei einem Wort nicht bleibet, und nicht einfältiglich glaubt, sondern weiter fragt, und denkt, Gottes Wort sei wie eines Menschen Wort, das sich ändern und beugen läßt.
\quellseite{258}

\section{Summarium}

Ist eine wunderbärliche Historia. Nachdem der Prophet Bileam sich des ersten Befehls Gottes, um seines Geizes und Menschengunst willen, nicht hat gehalten, sondern Gott versucht, und zum andernmal gefragt, strafet ihn Gott; wie es St. Petrus in seiner andern Epistel bezeuget durch die stumme Eselin, darauf er reitet, daß er unrecht habe gethan, daß er um Geldes und unrechten Lohnes willen sich wider den ersten Befehl Gottes hat lassen aufbringen, gleich wie ihn der Engel auch strafet. Aber wer einmal Gottes Wort nicht hat wollen folgen, der ist wohl werth, daß er ihm selbst zum Schaden und Aergerniß fortfahre, und darüber zu Boden gehe. Wie dieser Ursach halber diese Historia beschrieben ist, uns zur Warnung, ob wir gleich hohe Gaben haben, daß wir doch beim Wort auf das einfältigste bleiben, und uns davon nicht sollen abweisen lassen.

\section{Votum}

Gott regiere uns, daß wir ihn ja nicht versuchen, sondern daß wir seinem Wort gehorsamlich folgen, für Geiz und andern Sünden uns hüten, damit wir von ihm gesegnet werden, zeitlich und ewiglich, an Leib und Seele, Amen.

\chapter{Das 23. Kapitel}

\section{Vorrede}

Wem Gott wohl will, dem können der Teufel und böse Menschen nicht schaden. Denn also sagt Christus, Matth. 16.: Die Pforten der Hölle sollen sie (vernehmt meine Gläubigen) nicht überwältigen. Und David im 118. Psalm spricht: Der HErr ist mit mir, darum fürchte ich mich nicht, was können mir Menschen thun? Derwegen erwarte ein jedes seines Berufs und befehle sich täglich in Gottes Schutz und Segen, und lasse den Teufel Hulfstricke legen, und die bösen Menschen suchen und Unglück stiften, wie sie wollen und mögen; Gott behütet doch die Seinen, und wo die bösen Flüche ausgehen, da geben sie wieder ein, mit Unglück und Schaden. Dies vermahnet auch das 23. Kapitel.

\section{Bileams Fluch in einen Segen verwandelt}

a Kirchen Baal.

b (Herr) hie leugt Bileam, daß er will zum Herrn gehen, welcher ihm bereits gesagt hatte, daß er nicht sollte gehen; sondern er gehet zu seinen Zauberern unter des Herrn Namen.

c (Spruch) heißt die Oracula, das ist, solche Worte, die er nicht von ihm selbst redet; sondern die ihm Gott in den Mund gab. Als wenn ein Gottloser den Text des Wortes Gottes spricht, daß wider ihn selbst und die Seinen ist.
\quellseite{259}

d Menschen lügen, und müssen auch zuweilen fehlen, da sie nicht halten können, denn sie sind selbst ihres Lebens ungewis.

e Mühe und Arbeit heißt die Schrift Abgötterei, oder falschen Gottesdienst, und was ohne Glauben geschieht, Ps. 10, 7. Unter seiner Zunge ist Mühe und Arbeit.

f Trommeten des Königs) das ist, die lieblichen Trommeten Gottes ihres Königs, der sie zu machen befohlen hatte; darum sie unüberwindlich waren im Streit. Mag auch wohl heißen das Wort Gottes, so in diesem Volke lauter und öffentlich gelehret ward.

\section{Summarium}

Dies ist eine feine Historia, in welcher man sieht, wie es allen Menschen und dem Teufel selbst unmöglich ist, wo Gott seinen Segen gibt, daß man denselben könne wenden oder ändern. Wie denn Bileam nicht anders kann, denn daß er die Kinder Israel segne, V. 8. 20. Darum wer in Gottes Furcht und seinem Wort und Willen lebt, der darf sich für nichts besorgen. Denn ob die Leute schon wollten Schaden thun, so ist doch Gott auf jener Seite und behütet.

Sonderlich aber ist hie zu merken, daß Bileam die Juden fürnehmlich darum lobet und segnet, daß weder Mühe noch Arbeit, das ist, keine Abgötterei noch falscher Gottesdienst, sondern Gott und seine Trommeten, das ist, sein Wort unter ihnen ist, V. 21. Wo das ist, da darf man sich für Zauberei und Fluchen, ja für dem Teufel selbst nicht fürchten; wie es die Erfahrung noch heutiges Tages mit bringt. Wo man aber Gott nicht fürchtet, an sein Wort sich nicht hält und in Sünden lebt, da ists nicht Wunder, ob schon aus Gottes Verhängnis der Teufel, oder böse Leute, Schaden anrichten; denn sie können ohne das anders nichts, denn Schaden thun.

\section{Votum}

Der allmächtige Gott steure und wehre dem Teufel und bösen Menschen; erhalte uns in seiner Furcht, und behüte uns für allem Unglück, an Leib und Seele, bis zum ewigen, seligen Leben, Amen.

\chapter{Das 24. Kapitel}

\section{Vorrede}

Unglauben und Geiz fallen noch heute viele Menschen, die da meinen, sie werden selig werden, und werden doch im Abgrund der Hölle in
\quellseite{260}
Einigkeit müssen verlieb nehmen, und Qual leiden. Dies lehret der heilige Täufer Johannes, da er spricht, Joh. 3.: Wer dem Sohn nicht glaubet, der wird das ewige Leben nicht sehen, sondern der Zorn Gottes bleibet über ihm. Und St. Paulus Eph. 5. sagt also: Kein Geiziger, welcher ist ein Götzendiener, hat Erbe an dem Reiche Christi und Gottes. Darnach schaue ein jeder zu, wie er sich in Kaufen und Verkaufen, und in allerlei Handel und Wandel halte, und wie und was er glaube, und niemand betrüge sich selbst. Denn was Gott an andern gestraft hat, das wird er ihm an mir und euch nicht gefallen lassen. Davon lehret auch das 24. Kapitel.

\section{Bileams Weissagung von dem Stern aus Jakob}

a Hieraus merket man, daß Bileam droben allzeit zur Zauberei gegangen, unter Gottes Namen; aber der HErr ist ihm immerdar begegnet und hat die Zauberei gehindert, daß er hat müssen das rechte Wort Gottes fassen und aussprechen, anstatt der Zauberei.

b Amalek war der erste unter den Heiden, so die Kinder Israel angefochten, 2 Mos. 17, 8., aber durch Saul vertilget, 1 Sam. 15, 7.
\quellseite{261}

c Chittim) sind die aus Europa, als der große Alexander und Römer, welche auch zuletzt untergehen. Und zeuget hie die Weissagung, daß alle Königreiche auf Erden eins nach dem andern untergehen müssen, neben dem Volk Israel, welches ewig bleibt um Christi willen.

\section{Summarium}

Dies ist eine Weissagung des Propheten Bileams, wie seine Macht den Kindern Israel soll fürstehen, finstermal Gott bei ihnen ist, und will sie gesegnet haben, V. 4.

Daß er aber meldet von dem Stern aus Jakob und dem Scepter Israel, V. 17., das geht auf den Herrn Christum, welches Reich alle andre Reiche dämpfen, und alles unter sich bringen soll, nicht mit dem Schwerdt oder weltlicher Gewalt, sondern allein durchs Wort.

Aber ehe Christus kommt, sagt er, sollen die Assyrer gewaltig werden, und alles unter sich bringen, bis der große Alexander aus Chittim die Assyrer auch dämpfe, und endlich alle Gewalt an die Römer komme, und Christus geboren werde, V. 24. Solches ist eine treffliche Weissagung, und ohne Zweifel aus dem h. Geist. Darum Bileam ein rechter Prophet gewesen und Gott erkennet hat, aber durch Geiz und Unglauben ist er gefallen, wie andere große Leute mehr. Denn Weissagen ist eine solche Gabe des h. Geistes, die wohl Leuten widerfahren kann, die dennoch nicht selig werden.

\section{Votum}

Der treue, barmherzige GOtt regiere uns alle, daß wir für Sünden uns fleißig hüten, und im wahren Glauben gegen ihn, und in christlicher Liebe gegen unsern Nächsten beständig leben und selig sterben, Amen.

\chapter{Das 25. Kapitel}

\section{Vorrede}

Wie alle Menschen Abgötterei und Unzucht fliehen sollen, da sie nicht Leib und Seele zu büßen wollen: also soll die Obrigkeit Abgötterei abschaffen, und die Unzucht ernstlich strafen, damit Gott Land und Leute derwegen nicht strafen dürfe. O beiderseits eine nöthige Warnung. Demwegen laßt uns nicht leben, daß uns dieß Leben gereue, und uns und andern zum Verderben gelange; sondern laßt uns leben, daß wir um Christi willen alhier und in Ewigkeit in Gottes Gnade und selig leben, und daß wir unsern Nächsten durch Kraft des heiligen Geistes viel nützen und dienen. Darauf höret ferner an das 25. Kapitel.

\section{Abgötterei und Hurerei wird ernstlich gestraft}
\quellseite{262}

\section{Summarium}

Es ist ein schrecklich Exempel, wie das Volk Israel der Heiden Abgötter anbetet, und mit den heidnischen Weibern Hurerei treibet, V. 2. Solche Sünde strafet Gott mit einer so grausamen Plage, daß ihrer vier und zwanzig tausend darüber todt blieben, V. 4.; auf daß wir lernen, wie Gott beiden Sünden, der Abgötterei und Hurerei feind sei, und wolle sie greulich strafen.

Darnach ist ein nützlich Exempel mit Pinehas, des Aarons Enkel, der den Hurer sammt der Hure ersticht, V. 7. Denn da sehen wir, wo die Obrigkeit solche Sünde mit Ernst strafet, daß Gott die gemeine verdiente Strafe abwenden will, welche Strafe sonst müßte folgen, wo dergleichen Sünden von der Obrigkeit nicht gestraft würden. Denn wo die Obrigkeit mit Strafen faul und nachlässig ist, da wird Gott verursacht, die Sünde zu strafen, und gebet gemeiniglich darnach solche Strafe über den ganzen Haufen, wie man in Exempeln siehet.

\section{Votum}

Gott regiere uns alle durch seinen heiligen Geist, daß wir uns für Abgötterei und aller Unzucht hüten, und helfe, daß die Obrigkeit alle Laster ernstlich strafe, ihm zu Ehren, und den Menschen zur Besserung, Amen.

\chapter{Das 26. Kapitel}

\section{Das jüdische Volk wird von neuem gezählet}

a) Zeichen) ein Schreck-Zeichen, daran sie gedächten und sich für gleicher Sünde hüteten.
% No apparatus blocks retained on this page. See selection.json for block-by-block exclusion reasons.
\quellseite{264}

\chapter{Das 27. Kapitel}

\section{Vorrede}

Wir Menschen sind zur Ordnung von Gott erschaffen; aber keine Kreatur oder Geschöpf auf Erden lebet unordentlicher, als die Menschen. Darüber klagt der heilige Prophet Jeremias im 8. Kapitel, und sagt: Ein Storch unter dem Himmel weiß seine Zeit, eine Turteltaube, Kranich, Schwalbe, merken ihre Zeit, wenn sie wiederkommen sollen; aber mein Volk will das Recht des HErrn nicht wissen. Weil aber Gott gewisse Stände geordnet hat, und jedem Stande in seinem Wort Regeln und Weise vorgeschrieben, wie ein jeder Mensch in seinem Stande christlich und unsträflich leben soll; deswegen laßt uns Gottes Ordnung treulich folgen, damit er uns aus Gnaden um Christi willen, nach diesem elenden Leben, in das ewige und freudenreiche Leben ordne und versetze. Darauf höret auch an das 27. Capitel.

\section[Gesetz von Erbtöchtern. Josua an Mosis Statt zum]{Gesetz von Erbtöchtern. Josua an Mosis Statt zum Fürsten des Volks geordnet}
\quellseite{265}

a Deine Herrlichkeit) das wird vielleicht eine sonderbare Weise gewesen sein, daß Mose dem Josua den Stab oder die Hände aufs Haupt gelegt hat. Gleich wie man die Könige salben, oder, wie man die Lehen zu empfahen pflegt; welches alles muß eine Weise und Gepränge haben.

b Das Licht) das ist, das Licht auf der Brust des Hohenpriesters, 2 Mos. 28, 30. Daher sagen etliche, wenn Gott habe auf des Priesters Frage geantwortet, das hat sollen Ja sein, so habe das Licht Glänzen von sich gegeben. Es haben aber hernach die Könige also gebraucht, wenn sie Gott um Rath fragen; als 1 Sam. 28, 6. c. 30, 7.

\section{Summarium}

Am ersten ist eine seine nöthige Ordnung, wie es mit den Testamenten oder mit dem Bescheiden der Güter soll gehalten werden, wenn sich ein Todsfall begibt. Da ordnet es Gott also, daß das Gut allweg auf den nächsten Bluterben fallen soll, V. 6. 11.

Darnach ist ein sehr fein Exempel, als Gott Mosen abfordert von diesem Leben, und sagt, er müsse sterben, V. 13., da ist sein höchstes Anliegen und Begehren, daß Gott seinem Volk einen andern Fürsteher ordnen wolle, damit sie nicht ohne Regiment, wie die Schafe ohne einen Hirten seien, V. 16. Solchem Exempel soll alle Obrigkeit nachfolgen, und für ihre Unterthanen auch also sorgfältig sein und bitten.

\section{Votum}

Helfe der liebe GOtt, daß beide, Obrigkeit und Unterthanen, seinem Wort sich gemäß halten, in christlicher Regierung und gehorsamer Unterthänigkeit, und beschere uns allen zu seiner Zeit ein seliges Sterbestündlein, Amen.

\chapter{Das 28. Capitel}

\section{Gesetz von mehrerlei Opfern wiederholet}
\quellseite{266}

a Wein) ist hie sicera, hoc est, inebriativum, id est, purum, non dilutum, sed merum et efficax. Quia verbum purum inebriat vere animas. Das ist, berauschendes Getränk, d. i. rein, nicht verdünnt und geschwächt, sondern lauter und kräftig. Weil das reine Wort die Seelen wahrhaft trunken macht.

\chapter{Das 29. Capitel}

\section{Noch andere Opfer}
% No apparatus blocks retained on this page. See selection.json for block-by-block exclusion reasons.
\quellseite{268}

a Ezere(th) als auf dem Kirchhof versammelt und beschlossen, ut liceat audire maledictiones etc., daß man die Versuchungen hören könne. 5 Mos. 28, 15.

\chapter{Das 30. Kapitel}

\section{Von Gelübden, wie sie verbinden, oder nicht verbinden}

a Ihre Seele) das ist, wenn sie sich verbünde zu fasten, oder sonst etwas zu thun, mit ihrem Leibe (Gott zu Dienst; daß Seele hie heißt so viel als der lebendige Leib, wie die Schrift allenthalben braucht.

\chapter{Das 31. Kapitel}

\section{Vorrede}

Ach der großen Sicherheit und Sünden dieser letzten Welt. In allen Predigten wird gesagt vom Zorn Gottes wider die Sünde; von dem zeuget Gott selbst auch, erschrick vom Himmel durch Wunderzeichen, und durch ungewöhnliche Wetter, dann auf Erden durch allerlei wunderliche Unfälle. Wer fürchtet sich aber für Gottes Zorn? Wer bessert deswegen sein Leben? Sondern es gehet wie der heilige Prophet Jeremias im 5. Kapitel klaget, und sagt: HErr du schlägest sie, aber sie fühlens nicht, du plagest
\quellseite{269}
ke, aber sie bessern sich nicht, sie haben ein härter Angesicht denn ein Fels, und wollen sich nicht bekehren. Derwegen laßt uns Buße thun, und also leben, daß Gott nicht vollends seinen Zorn wie Wasser auf uns ausgieße, und durch seine Strafen, in zeitlichen und ewigen Verderb stürze. O HErr Christe, das hilf uns allen, Amen. Darauf höret nun an das 31. Kapitel.

\section[Israel erhält wider die Midianiter den Sieg, und]{Israel erhält wider die Midianiter den Sieg, und große Beute}

\section{Summarium}

Hie ist eine Historia, wie das Volk Israel die Midianiter geschlagen, und ihr ganzes Land geplündert haben, darum daß sie das Volk Israel zur Sünde gereizt und verführt hatten. So sieht man den Zorn Gottes wider die Sünde, daß er alles heißt erstechen und erwürgen, ausgenommen die Mägdlein, so noch nicht mannbar waren, der waren an der Zahl zwei und dreißig tausend, v. 7. 35., und zum Exempel und zur Besserung, auf daß wir lernen, daß Gott endlich pflegt die Sünde zu strafen, wo man nicht ablassen noch sich bessern will.
\quellseite{270}

\section{Summarium}

Dies ist ein fein Exempel, wie Gott befiehlt, daß man von der Beute, so an den Midianitern gewonnen war, den Leviten auch ihren Theil geben soll zum Zeugniß, daß der Sieg vom HErrn kommet, und die Leute denselben brauchen sollen zur Ehre Gottes und Förderung des Gottesdienstes. Wie am Ende die Hauptleute auch thun, und eine große Summe Goldes schenken zur Hütte, v. 48., zum Bekenntniß und Dank, daß ihnen Gott den Sieg gegeben und sie für Unglück behütet hätte. Das sind keine gottesfürchtige Kriegsleute gewesen, darum

15*
\quellseite{271}
hat ihnen Gott auch Glück gegeben. An solchen mangelt es uns wider den Türken. Denn unser Kriegs-Volk fragt gemeiniglich nach der Kirche und predigt nicht viel, spielen, Unzucht treiben, plündern und rauben, gottloslästern und fluchen ist ihr Brauch. Darum ist weder Glück noch Heil bei ihnen, Gott erbarme es.

\section{Votum}

Gott erbarme es, und streite selber für seine arme Christenheit wider den Türken, und alle andere Feinde, und erhalte uns den Land- und Stadtfrieden gnädiglich, bis an den lieben jüngsten Tag, Amen.

\chapter{Das 32. Capitel}

\section{Vorrede}

Gott hilft und segnet darum viel Leute, daß die so viel haben, denen rathen und dienen sollen, die nichts oder wenig haben. Aber dies kehren die Weltkinder um, und suchen immer weiter ihren Nutzen und Vortheil, wo auch der Nothdürftige bleibe. Derwegen trachte also ein jeder auf sein Bestes, daß er seines Nächsten nicht vergesse, sondern dem wo und wie er kann, helfe. Denn Nothdürftige verlassen, bringt Gottes Fluch und Strafe; dagegen wohl dem, der sich des Dürftigen annimmt, denn wird der Herr erretten zur bösen Zeit, sagt David im 41. Psalm. Ferner höret an das 32. Capitel.

\section{Der Besitzung des Landes Canaan wird ein Anfang gemacht}
\quellseite{272}

\section{Summarium}

Ist eine historia, wie Ruben, Gad und der halbe Stamm Manasse, das Land jenseit des Jordans einnehmen und besitzen. Doch also, daß sie den andern sollten helfen kriegen, und nicht Ursach geben, daß sie verzagt würden, V. 19. 25. Denn wo sie allein ihren Vortheil gesucht und sich der andern weiter nicht hätten wollen annehmen, sagt Moses, daß es eine große Sünde wäre, welche Gott zu seiner Zeit strafen würde. Und hie gehet an, daß Gott seine Verheißung vollstreckt, und sein Volk in der Heiden Lande einsetzet.

\section{Votum}

Dem treuen barmherzigen GOtt danken wir für alle seine Wohlthaten; der wolle seiner Verheißung nach an Leib und Seele uns segnen, und helfen, daß wir unsern Nächsten allzeit willig und nach höchstem Vermögen dienen, Amen.

\chapter{Das 33. Capitel}

\section{Register der Reisen und Lagerstätten des Volks Israel}
\quellseite{273}
\quellseite{274}

\chapter{Das 34. Capitel}

\section[Grenzen des gelobten Landes. Wie und durch welche]{Grenzen des gelobten Landes. Wie und durch welche es auszutheilen}
\quellseite{275}

\chapter{Das 35. Kapitel}

\section{Vorrede}

Das Leben ist des Menschen edelstes und herrlichstes Kleinod in dieser Welt, wie denn der Satan selber bekennt, Hiob 2. und saget: Alles was der Mensch hat, läßt er für sein Leben. Nun ist aber die Welt nichts anders, denn ein Stall voller Mörder, die andere tödten mit Herzen, Geberden, Worten und Werken. Derwegen verbeut Gott ernstlich allen Todtschlag, und befriedet uns als der Liebhaber aller Menschen wider aller Leute Frevel und Gewalt, und will, daß keiner dem andern an Leib und Leben schaden solle, und da es jemand thäte, dräuet er, daß er wieder solle gestraft werden an Leib und Leben zeitlich, und da er ohne Buße stirbt, auch ewiglich. Dies lehret auch das 35. Kapitel.

\section{Von den Städten der Leviten, Freistädten und Todtschlag}

Messen, das ist Geometrica proportione gerechnet, nämlich, die Vorstadt soll rings umher, an der Stadt, tausend Ellen weit sein, und eine jegliche Seite der Stadt zwei tausend Ellen lang. Das heißt auf deutsch: Die Vorstadt soll halb so weit sein, als eine Seite der Stadt lang ist. Sie sei viereckt und viereckt, oder wie sie kann, so soll man sie messen und in vier Seiten theilen; und darnach sie groß und klein ist, wird die Vorstadt auch groß oder klein.
\quellseite{276}

\section{Summarium}

Das ist eine Ordnung, daß man soll etliche Freistädt(e) machen für sie, so einen Todtschlag begangen haben, V. 6., nicht erlauben, daß man die Todtschläger befreien sollte, sondern daß man ordentlicher und redlicher Weise sich erkunden möchte, ob der Tod fürseßlich oder ohngefähr geschehen wäre, und also unschuldiges Blutvergießen verhütet würde.

Und ist hie sonderlich zu merken, daß Gott sagt, das Land werde durch Todtschlag verunreinigt, und solche Sünde möge nicht versöhnet werden, denn der Thäter werde von der Obrigkeit auch getödtet, V. 34.; denn man werde mit Geld noch andern, solche Sünde soll verrichten, noch versöhnen lassen. Und wo gleich ein Todtschlag nicht fürseßlich, sondern ohngefähr sich begeben hätte, müßte dennoch derselbe Thäter flüchtig sein, und sich in der Freistädte einer enthalten, V. 24., und dürfte nicht ehe zu seinem Erbe kommen, bis der Hohepriester stürbe, V. 24. Also hat Gott mit sonderlichem Ernst über dem fünften Gebot gehalten, und soll weltliche Obrigkeit solchen Ernst wohl bedenken, und derhalben ihrem Amt nach keinen Todtschlag ungestraft noch abbitten lassen; denn unschuldiges Blut schreiet zu Gott, und wo es an dem Rechtschuldigen nicht gerochen wird, muß der ganze Haufen und ein ganzes Land solche Sünde tragen und darum verhalten.

\section{Votum}

Wir danken Gott, dem leutseligen Menschenfreunde, daß er unserm Leib und Leben beschützet; der steure allen mörderischen Menschen, und regiere uns, daß wir niemand an seinem Leib und Leben weder Schaden noch Leid thun, Amen.

\chapter{Das 36. Kapitel}

\section{Vorrede}

Gott hat im jüdischen Volke die Erbe und Güter unterschieden, und geboten, daß die liegenden Gründe von einem Stamm an den andern
\quellseite{277}
 nicht erben sollten; denn Gott wollte, daß jedermann wissen sollte, aus welchem Stamm er wolle, daß der verheißene Messias sollte geboren werden. Daraus folget, daß Gott einem jeden das Seine gönnet, daß er mit gutem Gewissen erworben oder ererbet hat, und daß er nicht wolle, daß alles gemein sein sollte, wie davon die Wiedertäufer lehren. Und derwegen auch ein jeder Christ mit gutem Gewissen sein Eigenes haben und besitzen kann, doch daß er sein Herz auf den ungewissen Reichthum nicht stelle, sondern auf den lebendigen Gott traue, und daß er das, so ihm Gott bescheret, auch Gott zu Ehren, ihm selbst und andern zu Nutz anwende und gebrauche. Darauf höret an das 36. und letzte Kapitel.

\section[Durch ungleiches Heirathen sollen die Erbtheile nicht verrückt]{Durch ungleiches Heirathen sollen die Erbtheile nicht verrückt werden}

\section{Summarium}

Dies ist eine feine bürgerliche Ordnung, daß die Töchter, so eigen Erbe hatten, nicht durften in fremde Stämme heirathen, V. 6., auf daß die Theilung des Landes, wie sie Gott geordnet hatte, unverrückt bliebe.

\part{Das fünfte Buch Mose}

\chapter{Das 1. Kapitel}

\section{Vorrede}

Das fünfte Buch Mose hat nach der griechischen Sprache den Namen, und heißt Deuteronomium, zu deutsch, das wiederholte Gesetz. Darum daß Moses in diesem Buch wiederholet, erstlich die fürnehmsten Historien, und darnach fast das ganze Gesetz, welches er hie besser erkläret, und weitläufiger ausleget. Sonderlich aber, was die erste Tafel betrifft, wie man Gott von Herzen fürchten und ihm vertrauen soll, und sich seines Willens und Worts allein halten. In dem ersten Theil dieses Kapitels wird wiederholet die Historia aus dem andern Buch Mose am 18. Kapitel, in welcher sonderlich zu
\quellseite{278}
merken ist, wie Mose die vermahnet, welchen das Regiment befohlen ist, daß sie keine Person ansehen, den Armen wie den Reichen hören, und sich für Niemand scheuen sollen. Denn richten sei Gottes Amt; darum soll man es treulich und rechtschaffen führen, und auf Gott mehr denn auf die Menschen sehen. Solche Vermahnung sollen alle Regenten ihnen lassen sonderlich befohlen sein. Darauf höret nun an das 1. Capitel.

\section{Gottes Gutthaten. Israels Undankbarkeit}
\quellseite{279}

\section{Summarium}

Ist eine Historia aus dem 13. Kapitel des vierten Buches Mose, wie das Volk GOtt nicht vertrauete, daß er bey ihnen sein und für sie streiten würde. Darum strafet sie Gott, daß sie vierzig Jahre in der Wüste bleiben, und daselbst alle sterben mußten, und keiner in das gelobte Land kommen, auf daß wir lernen, wie man Gott nicht höher unehren, noch ehe zu Zorn bewegen kann, denn so man an ihm und seiner Hülfe und Macht zweifelt. Darum tröstet Moses so treulich, und will, daß man doch hinter sich denke, wie Gott in Egypten und am rothen Meer geholfen, und sein Volk, wie ein Vater seine Söhne, getragen und geleitet habe. Aber solches Predigen war gar vergeblich, gleich wie es auch vergebens war, da er erboth, sie sollten wider die Amoriter nicht streiten, und sie darum geschlagen wurden.

\section{Votum}

Helffe der treue, barmherzige GOtt, daß die Obrigkeit ihr Amt recht und glücklich führe; und stärcke unsern Glauben, daß wir ihm allzeit beständig vertrauen, in Noth und Tod, Amen.
\quellseite{280}

\chapter{Das 2. Kapitel}

\section{Vorrede}

Unser HErr Jesus Christus spricht Johannis am 5. Kapitel: Suchet in der Schrift, denn dieselbige zeuget von mir. Und St. Paulus gebeut Timotheo 1 Tim. 4, 13., er solle anhalten mit Lesen der Schrift, und rühmet Röm. 1., wie das Evangelium sei von Gott in der Schrift verheißen. So weiset uns auch St. Petrus mehr denn einmal zurück in die Schrift; damit Christus und beide Apostel uns je lehren die Schrift des Alten Testaments nicht zu verachten, sondern mit allem Fleiß zu lesen und anzuhören, weil sie selbst das Neue Testament so mächtiglich gründet und bewähren durchs Alte Testament, und sich darauf berufen. So wenig nun des Neuen Testaments Grund und Beweisung zu verachten ist, so theuer ist auch das Alte Testament zu achten. Und was ist das Neue Testament anders, denn eine öffentliche Predigt und Verkündigung von Christo, durch die Sprüche im Alten Testament gesetzt, und durch Christum erfüllet? Darauf höret ferner an das 2. Kapitel.

\section{Vom Sieg wider die Amoriter}
\quellseite{281}

\section{Summarium}

Dies ist eine Historia, in welcher Moses wiederholet, wie des HErrn Wort erfüllet sei, in dem, daß alle die, so an Gottes Hülfe verzagt hatten, in der Wüste die vierzig Jahr sind abgestorben, V. 15. Auch meldet Moses hie etliche Völker, so aus ihrem Lande vertrieben, und andere dagegen eingesetzt sind; auf daß wir lernen, wo dergleichen Strafen ergehen, daß es Gottes Werk sei, der endlich Land und Leute pfleget zu vertilgen, wenn seine Gottesfurcht noch Besserung will folgen.

\section{Summarium}

Ist eine Historia, wie Gott seine Verheißung erfüllet, und der Amoriter König geschlagen, und das Volk Israel sein Land einnehmen läßt. Und ist hie sonderlich zu merken, wie Moses den Amoritern Friede anbeut, daß sie unbeschädigt hätten mögen bleiben, wo sie den Juden den Durchzug erlaubt hätten. Aber die Welt, wenn das Stündlein kommt, muß zu ihrem Verderben selbst Ursach geben. Darum geschichts oft, wie mit den Amoritern hie, daß mancher könnte Friede und Ruhe haben, in Würden und Aufnehmen bleiben, aber durch unnöthiges Gezänk und andern Muthwillen machen sie ihnen und andern Leuten Unruhe, und bringen sich selbst in alles Unglück.

Am Ende ist ein sehr feiner Spruch, da Moses spricht: Es war keine Stadt, die sich für uns schützen konnte, der HErr unser Gott gab uns alles für uns, V. 36. Denn wo Gott seine Hand mit anlegt, da muß es fortgehen. Wiederum, wo er die Hand abzieht, da hilft weder Macht, Geld noch Gestreng. Darum soll man für allen Dingen, sonderlich aber, wo man Krieg führet, darnach trachten, daß man Gott zum Freunde habe, so darf man sich für aller Welt Feindschaft und Macht nicht besorgen.

\section{Votum}

GOtt verleihe, daß wir alle in seiner Furcht leben, steure allen unsern Feinden, und sei und bleibe unser GOtt in Noth und Tod, um JEsu Christi, seines lieben Sohnes, willen, Amen.
\quellseite{282}

\chapter{Das 3. Kapitel}

\section{Vorrede}

Gott ist heut so stark, gewaltig und mächtig, als er von Ewigkeit her gewesen. Daß er aber den Feinden seiner Christenheit so viel gönnet zu seiner Christen Schaden und Verderben, das geschieht nicht derwegen, daß er, Gott der HErr, nicht schützen und helfen könne, sondern unsere Sünden sind deß Ursach, als die Gott strafen soll und muß, und daß man Gott nicht recht bittet um seinen Schutz und Hülfe. Dies lehret Jeremias im 2. Kapitel und spricht: Es ist deiner Bosheit Schuld, daß du so gestraft wirst, und deines Ungehorsams, daß du gestraft wirst. Und St. Jakob in seiner Epistel im 4. Kapitel sagt: Ihr bittet und krieget nicht, darum daß ihr übel bittet. Derwegen daß Gott den Türken und andern unseren Feinden steure, so wollen wir uns bekehren, auch in täglicher Buße leben, und Gott in wahrem Glauben anrufen, daß er uns wider unsere Feinde behüten und erhalten wolle. Dies erinnert auch das 3. Kapitel.

\section{Sieg wider Og, den König zu Basan}
\quellseite{283}

a Haus(e) Kirchen oder Tempel.

\section{Summarium}

Ist eine Historia, in welcher Moses wiederholet, wie Gott den König Og zu Basan in der Kinder Israel Hände gegeben hat. Da stehet man wieder einmal, daß der Sieg Gottes Gabe ist. Am Ende ist ein sonderlich Exempel, wie Moses der treffliche und heilige Mann, Gott darum bittet, er wolle ihn mit ins gelobte Land lassen, v. 23. Aber Gott wird zornig auf ihn, und will ihn nicht erhören, v. 26.; auf daß wir am ersten lernen, wie eine schwere Sünde der Unglaube sei, mit welchem sich Moses bei dem Haderwasser versündiget, und Gott also erzürnet; darnach daß wir lernen unsern Willen in Gottes Willen setzen und geduldig sein, wenn wir um etwas bitten, und dasselbige uns nicht widerfähret; sintemal es Gott gegen Mosen und andere Heiligen auch gethan, und sie in der Anfechtung hat stecken lassen.

\section{Votum}

GOtt stärke unsern Glauben, behüte uns für Ungeduld unter dem Kreuze, und steure und wehre allen Feinden seiner armen Christenheit, Amen.

\chapter{Das 4. Capitel}

\section{Vorrede}

Der königliche Prophet David lehret im 1. Psalm, daß dieß wahre Christen seien, die da Lust haben zum Gesetz des HErrn, und reden von seinem Gesetz Tag und Nacht. Das Gesetz aber nennet David Gottes Wort, das uns Gott (und damit weise und verständig zu machen) offenbaret, und darinnen er uns sein Wesen und gnadenreichen Willen zu erkennen gegeben hat; daraus wir auch lernen, wie wir christlich leben, selig sterben und fröhlich am jüngsten Tage zum ewigen Leben auferstehen sollen. Derwegen soll ein jeder Christ zu diesem Gesetze und Wort Gottes Lust, Freude, Gefallen allezeit und sein Leben lang haben, damit für und für umgehen, und sich darinnen erlustigen, auch solches zu Besserung seines Lebens und Stärkung seines Glaubens fleißig gebrauchen. Ach das helfe uns allen die heilige hochgelobte Dreifaltigkeit, Amen. Darauf höret ferner an das 4. Capitel.

\section{Vermahnung zum Gehorsam des Gesetzes}

a Bewahren) denn Menschen-Lehre hindert Gottes Gebot und führet von der Wahrheit, Tit. 1, 14.
\quellseite{284}

\section{Summarium}

Bisher hat Moses etliche Geschichten wiederholet. Jetzt fähet er an, wie ein treuer Prediger, weil er doch sterben und sein Volk hinter ihm lassen muß, prediget und vermahnet, und warnet sein Volk, daß sie erstlich bei Gottes Wort fest bleiben, darnach demselben ernstlich folgen, und endlich der gnädigen Hülfe und wunderbarlichen Werke Gottes nicht vergessen, sondern dieselbigen ihren Nachkommen predigen und kund thun sollen, V. 2.

Sonderlich aber meldet er, wie Gott am Berge Sinai sich offenbaret und selbst seinen Bund und Gesetz dem Volk fürgetragen habe; demselbigen nach sollen sie sich halten, und keine anderen Götter anbeten, auch kein Bildniß einiger Creatur machen, dieselbigen anbeten, und ihnen zu dienen, wie solches im ersten Gebot ausdrücklich befohlen ist, V. 10.

Wir sollen aber merken, daß Moses die Bildnisse ferner nicht verbeut, denn daß man sie nicht soll anbeten. Sonst kann man ohne Sünde Bildnisse haben zur Erinnerung, wie die Juden die eherne Schlange, Aarons Ruthen, die Cherubin und anderes hatten, und wir in unsern Kirchen, das Bildniß des Leidens, der Auferstehung Christi und anderer Historien haben, nicht zum Anbeten, sondern zur Erinnerung. Solches ist recht gut und nützlich, und in der
\quellseite{285}
heiligen Schrift nirgend verbieten, wie doch die Bildstürmer ohne Grund fürgeben, und ohne allen Unterschied allerlei Bildniß aus der Kirche werfen, so doch wahrhaftige Bilder gleich als eine Schrift sind, welche die Ungelehrten lesen, und sich davon bessern können.

b Das ist, um euretwillen.

c Faciebus suis, scilicet, nube et columna. Das ist, nämlich mit der Wolke und Feuersäule.
\quellseite{286}

\section{Summarium}

Ist eine ernste Warnung, weil Gott durch sein Wort ihnen offenbaret, was sie thun und nicht thun sollen, daß sie sich demselbigen nach mit Fleiß sollen halten, oder Gott werde sie um ihrer Sünde willen aus dem Lande vertreiben, und den Heiden dienen lassen. Doch tröstet er, wo man sich bekehre und Buße thue, daß Gott gnädig und sich ihrer wieder annehmen werde. Jerobalben ist solches wie eine Prophezeyung; denn mit den Juden ist es hernach also gegangen: da Israel Abgötterei trieb, kam der König von Assyrien, und führete sie hinweg, und verwüstete das Land. Das Königreich Juda besserte sich ob solcher Strafe nicht, da that der König von Babel ihnen dergleichen. Solches sind ehe Warnungen, daß wir Gottes Wort für Augen haben und dawider nicht thun sollen.

Sonderlich aber rühmet Moses dieses hie für die größte Wohlthat, die sie erkennen und dankbar dagegen sein sollen, daß Gott sie zum eigenen Volk ihm erwählet und ihnen sein Wort gegeben, und sie so wunderbarlich geführet und erhalten habe, auf daß sie sollen Gott gehorsam sein; so werde es ihnen und ihren Kindern wohl gehen, und Glück und Heil haben, V. 29. Solches alles sollen wir zur Lehre annehmen, und uns auch daraus lernen bessern.

\section{Votum}

Der treue, barmherzige GOtt regiere uns durch sein Wort und den heiligen Geist, daß wir alle in seiner Furcht leben, erhalte uns auch das Evangelium Christi, unsers HErrn, und gebe uns ferner seinen Segen, zu zeitlicher und ewiger Wohlfahrt, Amen.

\chapter{Das 5. Capitel}

\section{Vorrede}

Die zehn Gebote Gottes, so zu reden, sind anders nichts, denn ein göttlich Land- und Stadtrecht, darnach sich alle Christen richten sollen. Denn nachdem ihm Gott der HErr außer auf Erden auch eine Stadt und Bürgerschaft, das ist, eine Christenheit sammeln und erhalten will, so will er dieselbige auch mit gewissen Statuten und Satzungen versehen, darnach sich die Bürger in der Stadt Gottes, das ist die Christenheit, richten sollen. Dies Land- und Stadtrecht, welches sind die heiligen zehn Gebote, hat Gott selbst durch seinen göttlichen Finger in zwo steinerne Tafeln geschrieben. Die erste Tafel begreift die ersten drei Gebote und lehret, wie wir uns gegen dem Landfürsten selbst, das ist, gegen Gott dem HErrn halten sollen, mit dem Herzen nach dem ersten Gebot, mit dem Munde oder Worten nach dem andern Gebot, mit dem Werke gegen sein Wort, nach dem dritten Gebot. Die andere Tafel begreift die letzten sieben Gebote, und lehret, was wir unserm Nächsten zu Ehren und Gehorsam, am Leibe, Weibe, Gut, Ehr, Haus und Hausgesinde, schuldig sind. Dieses sollen wir betrachten bei dem 5. Kapitel.

\section{Wiederholung der zehn Gebote Gottes}
\quellseite{287}

\section{Summarium}

Hie fähet Moses an, und wiederholet die Predigt der zehn Gebote, in welcher Gott seinem Volk anzeigt, was man thun und lassen soll gegen jedermann, wenn man recht thun und selig leben will. Da gehöret nun zu, daß man erstlich solche Gebote recht lerne und wisse, darnach daß man Gott um seine Gnade und heiligen Geist bitte, daß er ein solches Herz geben wolle, auf daß wir solchen Willen also mögen vollbringen, und der Sünde widerstehen; denn wo Gottes Geist nicht ist, da wird solcher Wille Gottes nimmermehr recht erfüllet. Wiederum aber haben wir die Verheißung, daß Gott seinen
\quellseite{288}
Geist denen will geben, so ihn darum bitten, Luc. 11.

\section{Votum}

Diese Verheißung erfülle an uns allen der treue, barmherzige GOtt, und helfe, daß wir nach allen seinen Geboten in rechtem Gehorsam allezeit wandeln bis zum ewigen, himmlischen Leben, Amen.

\chapter{Das 6. Kapitel}

\section{Vorrede}

Wir Christen sind gefreiet vom Gesetz Gottes, was anlanget die Verfluchung und Verdammniß; aber mit nichten sind wir frei, was den Gehorsam belangt. Denn das ist der ewige unveränderbare Wille Gottes, daß die vernünftigen Creaturen, Engel und Menschen, dem Gesetz Gehorsam leisten sollen. Und solchen Gehorsam will ihm Gott gefallen lassen, um seines Sohns unsers Herrn Jesu Christi willen, und will alle die, so sich üben in Gehorsam nach dem Gesetz, segnen an Leib und Seele, zeitlich und ewiglich. So betet, o ihr Christenmenschen täglich und herzlich, Gott wolle euch durch den heiligen Geist regieren und leiten, daß ihr in treuem Gehorsam der heiligen zehn Gebote, gegen Gott im Himmel, und gegen den Menschen auf Erden, leben und sterben möget. Wie ihr denn solches auch in der Taufe eurem Gott und Herrn habt versprochen und zugesaget. Dies erinnert auch das 6. Kapitel.

\section{Erklärung des ersten Gebots, von der Liebe GOttes}

*1 Mos. 18, 19. Schärfen immer treiben und üben, daß sie nicht verrosten noch verdunkeln, sondern stets im Gedächtniß und Wort als neu und helle bleiben. Denn je mehr man Gottes Wort handelt, je heller und neuer es wird, und heißt billig je länger je lieber. Wo mans aber nicht treibet, so wird es bald vergessen und unkräftig.
\quellseite{289}

\section{Summarium}

Dies ist eine Vermahnung, da Moses insgemein die zehn Gebote mit beschließet, daß man Gott fürchten, und diese seine Gebote also hören und behalten soll, daß man sie thue und ins Werk bringe, V. 2. Sonderlich aber vermahnet er, wenn sie ins gelobte Land kommen, und da allerlei Segen und Reichthum finden werden, daß sie des HErrn Werk und Gebot nicht vergessen, sondern dieselbigen ihren Kindern schärfen, das ist, stets davon reden, und im frischen Gedächtniß behalten sollen, V. 6., auf daß sie Gottes Huld behalten, und sich nicht versündigen, wie es doch gemeiniglich gehet; denn wenn es mit uns am besten gehet, und wir alles genug haben, geschieht es gemeiniglich, wo wir nicht sonderlich uns in Gottes Furcht halten, daß wir sicher, stolz und muthwillig werden, und Gottes vergessen. Darum will es vonnöthen sein, daß man in Gottes Furcht wandele, und stets sein Wort übe und treibe, wie Moses hie befiehlet.

\section{Votum}

Helffe der treue, barmherzige GOtt, daß wir nach seinen heiligen Geboten gehorsamlich folgen, auf daß uns an Leib und Seele sein väterlicher Segen folge, in diesem gegenwärtigen elenden und künftigen seligen Leben, um Jesum Christi seines lieben Sohnes willen, Amen.

\chapter{Das 7. Kapitel}

\section{Vorrede}

Gott hat Gefallen an allen denen, die ihn fürchten, und solche haben keinen Mangel, und es wird ihnen wohl gehen in der letzten Noth. Aber dagegen wird Gott ein schneller Zeuge sein wider die, so ihn nicht fürchten, und denen wird es endlich übel gehen. Beides lehret Sirach im 1. Kapitel und saget: Wer den HErrn fürchtet, dem wirds wohl gehen, und wenn er Trost bedarf, wird er gesegnet. Wer aber ohne Furcht fährt, der gefällt Gott nicht, und seine Freiheit wird ihn stürzen. Demwegen wollen wir durch Gottes väterliche Hülfe also leben, damit wir nicht Gott und seine Gebote verachten, sondern wollen den HErrn fürchten, und ihm allezeit dienen, treulich und rechtschaffen, Josua 24. Darauf höret auch an das 7. Kapitel.

\section{Freundschaft mit den Heiden und deren Götzendienst verboten}
\quellseite{290}

\section{Summarium}

Sie ist ein sonderlicher nöthiger Befehl, daß das Volk Israel ganz und gar keine Gemeinschaft mit den Heiden habe, v. 2., sondern sie rein ausrotten und verwüsten soll, mit allen ihren Kirchen, Götzen und Bildern, v. 5., um der Ursache willen, daß sie an ihnen auch nicht lernen und gewöhnen Abgötterei zu treiben. Denn es ist gar leichtlich versehen, daß der Teufel uns Menschen in seine Stricke bringet, wenn wir bei und mit den Gottlosen stets umgehen.

Solche Lehre und Vermahnung ist hoch vonnöthen heutiges Tages, auf daß man des Papsts Geschmeiß mit unnöthigen Kirchen und Klöstern, mit unnützen Büchern und lügenhaftigen ärgerlichen Bildern, und was dergleichen mehr ist, aus dem Weg räume, und nicht also \unsicher{bege…} um des künftigen Aergernisses willen, daß folgen kann, wenn man solches Rattengeniße duldet und nicht ausfeget.

\section{Summarium}

Ist eine sehr schöne und tröstliche Predigt, wie man es so reichlich werde genießen, wenn man Gott fürchtet und seinen Geboten und Worten folget, nämlich, daß Gott Segen des Leibes, beides an Menschen und Vieh geben, alles auf dem Felde reichlich wachsen lassen, allen Krankheiten wehren, Glück und Sieg wider die Feinde geben wolle.

Die sollen wir mit Fleiß lernen und merken, woher doch komme allerlei solcher Unrath, Hunger, Krankheit, Kriege und dergleichen, nämlich davon, daß man Gott nicht fürchtet, und seinem
\quellseite{291}
Wort und Willen nicht folget. Darum sollen wir gern fromm sein, und für dem Bösen uns hüten, weil es Gott so reichlich belohnet, und wo mans nicht thut, so mancherlei Weise strafen will.

\section{Votum}

Unser HErr JEsus Christus, welcher der Weg und die Wahrheit ist, regiere uns, daß wir aller Gottlosen Gemeinschaft fliehen, und helfe, daß wir seinen Geboten folgen, auf daß er auch heut Menschen und Vieh segne und behüte zu unserer Wohlfahrt, Amen.

\chapter{Das 8. Kapitel}

\section{Vorrede}

Wenn es den Menschen übel gehet, so gedenken sie an Gott und beten; aber wenn es ihnen wohl gehet, so vergessen sie gemeiniglich Gottes, und sind undankbar. Solches sollen wir bedenken, und Gott täglich für seine Wohlthaten danken, mit Mund, Herz und christlichem Leben, damit uns der gerechte Gott mit neuem Kreuz und Strafen nicht dazu treiben dürfe. Denn es soll also heißen, wie Gott spricht im 50. Psalm: Rufe mich an in der Noth, so will ich dich erretten, so sollst du mich preisen. Das ist so viel gesagt: Ich stecke dich Mensch darum in Noth, auf daß ich dir Ursache gebe zum Gebet, und ich Ursache habe dir zu helfen, und dir damit auch ferner Ursache gebe, mich zu loben und zu preisen. Dieses soll auch erinnern das 8. Kapitel.

\section[Das Volk wird ermahnet, der Gutthaten GOttes nicht]{Das Volk wird ermahnet, der Gutthaten GOttes nicht zu vergessen}
\quellseite{292}

\section{Summarium}

Ist eine schöne und sehr nöthige Vermahnung, daß das Volk Gottes Gebote halten, und ja seiner Werke und Wunderthaten dann nicht vergessen wolle, wenn sie ins Land kommen, und es ihnen wohlgehet, V. 2. Denn dies ist eine sehr gemeine und gefährliche Anfechtung, daß man unsers HErrn Gottes vergisset und hoffärtig wird, wenn es wohlgehet, darum schickt Gott den Seinen hie auf Erden so mancherlei Trübsal und Anfechtung zu, auf daß er sie in der Furcht und Demuth erhalte, wie Moses hie sagt, V. 16.: Gott habe sein Volk in der Wüste auf mancherlei Weise versucht, auf daß sie seiner nicht vergessen sollten, noch gedenken, sie hätten alles von ihnen selbst.

\section{Votum}

GOtt erbarme sich unser, und gedenke in der Strafe an seine Barmherzigkeit. Er wolle auch verleihen, daß wir an seine Wohlthaten allzeit gedenken, und die loben, preisen und recht gebrauchen, Amen.

\chapter{Das 9. Kapitel}

\section{Vorrede}

Sollte uns Gott nicht mehr zu unserer Leibes- und der Seelen-Nothdurft geben, als daß wir verdienen, ach weh uns, da würde uns nichts mangeln denn alles, und solches überall. Derwegen, weil uns Gott unser Leben lang so große und unzählige Wohlthaten erzeiget hat, und weil er uns noch täglich alle Augenblicke an Leib und Seele Gutes thut, da wir doch eitel Strafe verdienen, so laßt uns Gottes Güte und Barmherzigkeit erkennen, und ihm dafür danken. Dazu vermanet der heilige Prophet Jeremias im 33. Kapitel und saget: Danket dem HErrn unserm Gott Zebaoth, daß er so gnädig ist, und thut immerdar Gutes. Dies lehret auch das 9. Kapitel.

\section{Ursachen der Einnehmung des Landes Canaan}
\quellseite{293}

\section{Summarium}

Dies ist die Vermahnung, wenn sie ins Land kommen und die Heiden vertrieben haben, daß sie Gott die Ehre geben sollen, daß er solches ihnen um seiner Zusagung willen und nicht darum, daß sie es verdienet hätten oder so fromm wären, habe widerfahren lassen, V. 4. Und hie wiederholet er, wie mancherley ihre Väter sich in der Wüste an Gott versündiget haben, auf daß wir lernen, wie wir nicht allein aus Gnaden selig werden, sondern auch alles, was uns Gott in diesem zeitlichen Leben gibt, daß es sei uns aus Gnaden geschenkt, und sich niemand einiges Verdienstes für Gott rühmen dürfe.

\section{Votum}

Lob, Ehr, Preiß und Dank sei GOtt durch JEsum Christum, unsern lieben HErrn, für alle seine großen und vielfältigen Wohlthaten; der wolle uns, die Unseren und die ganze Christenheit, ferner geistlich und leiblich segnen und behüten, als ein treuer, lieber Vater, Amen.

\chapter{Das 10. Kapitel}

\section{Vorrede}

Gott vertrauen und seinem Wort folgen, das sind die Dienste, die Gott von uns fordert, zur Dankbarkeit für seine Wohlthaten, die er
\quellseite{294}
uns reichlich erzeiget. Demwegen weil mir alle bekennen müssen, daß uns Gott ja große und unzählige Wohlthaten bisher mitgetheilet hat, so wollen wir seinem Wort so viel als möglich gehorsam sein, und wollen ihm treulich dienen im Glauben, und unserm Nächsten in wahrer Liebe, daß wir also allzeit dankbar erfunden werden, und daß wir auch mit dem königlichen Propheten David sagen können, aus dem 79. Psalm: Wir dein Volk und Schafe deiner Weide, danken dir ewiglich, und verkündigen deinen Ruhm für und für. So gehöret auch hieher, daß der heilige Apostel Johannes sagt, 1 Epist. Kap. 4.: Das Gebot haben wir von Gott, daß, wer Gott liebet (vernehmt, und die Liebe durch Dankbarkeit beweisen will), daß er auch seinen Bruder liebe. Dazu vermahnet auch das 10. Kapitel.

\section[Die Erneuerung des Bundes soll Israel zur Liebe]{Die Erneuerung des Bundes soll Israel zur Liebe GOttes reizen}
\quellseite{295}

\section{Summarium}

Ist eine Predigt, in welcher Moses etliche mehr Historien wiederholet, und endlich vermahnet, weil Gott sich so gnädig erzeiget, und ihnen nun das Land auch zu eigen geben wolle, daß sie ja wiederum Gott wollen dienen, der anders nichts fodert, denn daß man ihn fürchten und vertrauen, und seinem Worte folgen wolle, V. 12. Insbesondere befiehlt er ihnen die Fremdlinge, daß sie ihnen dienen, und alles Guts thun sollen, weil sie in Egypten auch sind Fremdlinge gewesen, V. 18.

\section{Votum}

Verleihe GOtt, unser lieber Vater, daß sein Wille unser Wille sei, auf daß wir ihm dankbar sein für alle seine Wohlthaten, und daß wir in Gehorsam leben nach seinen Geboten, auch zu Trost unsern nothleidenden Mitchristen, Amen.

\chapter{Das 11. Capitel}

\section{Vorrede}

Was hilfts, wenn die Lehre ist göttlich, und die Menschen leben darauf teuflisch? Also was hilfts, wenn man in der Kirche mit reiner Lehre bauet, und die Leute reißens wieder ein daheim mit gottlosem Leben? Derwegen laßt uns zuschauen, wie wir Gottes Wort hören, und laßt uns dies gebrauchen, auf daß wir sammt den Unsern nicht über unsere Vergessenheit klagen, und über unsern Ungehorsam ewig wehklagen dürfen. Denn wer Gottes Wort verwirft mit Ungehorsam, den wird Gott auch verwerfen mit Zorn, aus seiner Gnade in Noth und Herzeleid; dagegen selig sind zeitlich und ewig, alle die da hören die Worte der Weissagung, und behalten, was darinnen geschrieben ist, Offenb. Joh. 1., und darnach christlich und ehrbar leben. Darauf höret an das 11. Kapitel.

\section[Ursachen, warum GOtt zu lieben und seine Gebote]{Ursachen, warum GOtt zu lieben und seine Gebote zu halten}

a In pedibus eorum: Ich habe gelaufen und gerannt, daß ichs zusammen brächte. Non sterti, nec otiosus fui: Ist mir sauer worden, erjagt, erraunt, erlauffen.
\quellseite{296}

\section{Summarium}

Ist eine sehr schöne Vermahnung, in welcher Moses fast die fürnehmsten Historien wiederholet, in welcher wir sehen, wie GOtt seinen Zorn wider die Gottlosen, und seine Gnade seinem Volk erzeiget hat, V. 3., auf daß sie solche Exempel nicht vergessen, sondern daran lernen sollen GOtt fürchten, und ihm vertrauen, und seinem Wort folgen. Thun sie das (spricht er) V. 13., so werde Gott seinen Segen ihnen reichlich widerfahren lassen; thun sie es aber nicht, so werde sie Gott mit Hunger und Kummer und allerlei Unglück strafen. Solche Vermahnung sollten wir zu Herzen nehmen und uns bessern, so würde Gott nicht also mit unserer Zeit, Krieg, Pestilenz und anderm Unglück unsern Ungehorsam heimsuchen, wie jetzt geschieht und noch mehr geschehen wird, wo wir also in allen Sünden und Ungehorsam fortfahren.

\section{Votum}

Der treue, barmherzige GOtt regiere uns, daß wir von Sünden ablassen, damit wir zu strafen nicht Ursache geben; und helfe, daß wir also leben, auf daß es ihm ein Gefallen sei, und daß er uns aus Gnaden segne, zeitlich und ewiglich, Amen.

\chapter{Das 12. Kapitel}

\section{Vorrede}

Wir danken unserm HErrn Christo herzlich, daß er uns Heiden in seinen Schaffstall und christliche Kirche auch berufen hat, und daß er unser Licht, Trost und Heil sein will. Denn also sagt von Christi Wohlthaten der Prophet Jesaias im 40. Kapitel: Er wird seine Heerde weiden, wie ein Hirte, er wird die Lämmer in seine Arme sammeln und in seinem Busen tragen, und die Schafmütter führen. Und der HErr Christus lehret von ihm selbst auch also, Joh. 10.: Ich gebe meinen Schafen (spricht er) das ewige Leben, und sie werden nimmermehr umkommen, und niemand wird sie mir aus meiner Hand reißen. Lasset uns aber zuschauen, daß wir als seine Schafe seine Stimme und Wort hören, dem glauben, und mit einem
\quellseite{297}
unsträflichen Leben und Wandel folgen, bis zum seligen Todesstündlein. Wie abermals der HErr Christus lehret, Joh. 10. und sagt: Meine Schafe hören meine Stimme, und ich kenne sie, und sie folgen mir. Das helfe uns allen der hochgelobte Sohn Gottes, Christus Jesus, Amen. Darauf höret weiter an das 12. Kapitel.

\section{Ort und Weise des wahren Gottesdienstes}
\quellseite{298}

\section{Summarium}

Hie ist ein ernstlich Gebot, daß die Juden, wenn sie ins Land kommen, ganz und gar nichts sollen lassen bleiben von allem, das die Heiden zu ihrer Abgötterei gebraucht haben, V. 2., allein darum, daß sie nicht Ursach haben gleiche Abgötterei zu treiben. Solches sollen wir auch bedenken, und nichts dulden, was noch vom Papstthum fürhanden ist, und uns zum Aergerniß kann gereichen.

Darnach befehlet Moses insonderheit, wer Gott wolle mit Opfern und andern dienen, daß ers ja nirgend thun soll, denn an dem Ort, welchen ihm Gott selbst dazu erwählen würde, V. 5. Denn Gott kann man nicht dienen, man thue denn den Befehl seines Worts, wie ers ordnet und fordert.

\section{Votum}

GOtt der heilige Geist erleuchte und regiere uns, daß GOttes Wort unserer Füße Leuchte sei, auf daß wir dem in Thun und Lassen gehorsamlich folgen, bis wir durch das finstere Thal des Todes eingehen in den Freudensaal des ewigen Lebens, Amen.

\chapter{Das 13. Kapitel}

\section{Vorrede}

Nachdem Mose mit Fleiß, wie ein treuer Prediger die Seinen vermahnet hat, wie sie Gott fürchten, und in seinem Wort und Willen wandeln sollen, befiehlt er hie, wie man mit den Predigern handeln und umgehen soll, so sie etwas anders lehren und predigen: nämlich, daß man sich an kein Zeichen noch Wunderwerk soll kehren, sondern allein an GOttes Wort halten, V. 3. Was wider das Wort ist, es seien Wunderwerke wie es wollen, soll mans für Teufels-Gespenst und lauter Verführung halten, und solche verführerische Prediger ohne alle Erbarmung tödten, V. 5.

Wo man auf solchen Befehl im Papstthum gesehen, und demselben gefolget hätte, würde der Teufel mit falscher Lehre und falschen Zeichen nicht so viel ausgerichtet, noch die Abgötterei so tief eingerissen haben. Aber weil man die Lehre und das Wort hat lassen fallen, hat der Teufel gute Sache gehabt, und durch lügenhaftige Zeichen den Irrthum wohl gefördert.

\section{Strafe der falschen Propheten, der Verführer und Verführten}

a. Andern) dieser Prophet lehret wider die alte und bestätigte Lehre, und will (wie er sagt, andere Götter) die ersten nicht lassen bleiben, welchem man nicht soll glauben, wenn er gleich Zeichen thät. Aber im Kap. 18, 22. wird der Prophet verdammt, der nicht wider die erste, sondern neue Lehre fürgibt, der soll Zeichen thun, oder nicht gehöret werden.
\quellseite{299}

b Erbarmen, das heißt, Gott über alles lieben.

\chapter{Das 14. Kapitel}

\section{Vorrede}

Gleichwie das vierte Gebot der Ehrenkranz ist, welchen Gott der HErr den Eltern aufgesetzt, gleichwie Gottes Wort und Ordnung des weltlichen Regiments das güldene Band ist, welches Gott der HErr der Obrigkeit hat angehängt, also ist das auch das schöne Kleinod, die güldene Kette und der Freudenbrief, welchen Christus Jesus allen treuen Lehrern und Predigern verehret hat: Wer euch höret, der höret mich, wer euch verachtet, der verachtet mich, Luc. 10. Es sollen aber die Zuhörer ihre christliche Seelsorger, beide ihres Amts und ihrer Person halben, also hören und ehren, daß sie auch dieselben ernähren helfen sollen. Hieher gehöret St. Pauli Vermahnung, die aus dem 6. Kapitel an die Galater also lautet: Der unterrichtet wird mit dem Worte, der theile mit allerlei Gutes dem, der ihn unterrichtet. Irret euch nicht, Gott läßt sich nicht spotten. Das erinnert auch zum Theil das 14. Kapitel.

\section{Trauer-, Speis- und Zehntordnung}
\quellseite{300}

\section{Summarium}

In diesem Kapitel wiederholet Moses etliche Gesetze, so den Juden insonderheit aufgelegt waren. Sonderlich aber treibt er das Gebot mit dem Zehnten, V. 22. und andern, so zu Unterhaltung der Priester und Leviten von GOtt geordnet war, daß man es treulich ausrichten und gern geben soll, wie er droben auch sagt: Hüte dich und verlaß den Leviten nicht, so lange du auf Erden lebest; auf daß wir lernen, wie es Gottes Ordnung und Wille sei, daß jedermann mit dem Seinen gern dazu helfen soll, auf daß die Kirchendiener mögen erhalten werden.

\section{Votum}

Der treue GOtt erhalte uns bei seinem heiligen Worte, und behüte unsere Kirchen und Schulen für falscher Lehre und falschen Gottesdiensten, um der Ehre seines heiligen Namens, und wegen unserer Seelen Seligkeit, Amen.
\quellseite{301}

\chapter{Das 15. Kapitel}

\section{Vorrede}

Freilich, was man den Armen gibt, deß wird niemand ärmer; denn solches wird Gott auf Wucher geliehen, sagt Salomon in seinen Sprüchen im 19. Kapitel. Wohl nun denen an Leib, Seel und Nahrung, ja auch wohl ihren Kindern, ganzem Geschlecht und allen Nachkommen, an geistlichen und leiblichen Gütern, die sich der Armen erbarmen. Denn wahrhaftig des armen Hand ist unsers Herrn Christi Schatzkasten, und ihre eigene Sparbüchse. Derwegen, wer der Armut dienen und helfen kann, der lasse es an ihm nicht mangeln; denn Gott wird ihm wieder dienen und helfen, wo er seiner am meisten bedürfen wird, und wird ihn erquicken auf seinem Sterbebette, Psalm 41. Davon höret auch an das 15. Kapitel.

\section{Vom Erlaßjahr, verkauften Knechten, und Erstgeburt}
\quellseite{302}

\section{Summarium}

Dieß ist eine Ordnung, daß die Juden im siebenten Jahr allwege alle Schuldregister sollen abthun, und nichts mehr an ihre Brüder fordern, V. 2. Weil aber zu besorgen war, sie würden desto weniger mitleiden, wenn das siebente Jahr nahe wäre, vermähnet Mose hie auf das allerschärfste, daß man darum die Armen nicht wolle lassen Noth leiden, V. 8.; denn damit versündiget man sich. Wo man aber gern und willig helfe, da wolle Gott mit reichem Segen kommen, und Glück geben in allem, was man fürnimmt.

Da sollen wir nun auch lernen den Armen die Hand aufzuthun, und gerne helfen, weil es Gott so wohlgefällt, und er so reichen Segen dagegen verheißt.

\section{Votum}

GOtt regiere uns alle, daß wir uns der Armen Noth treulich annehmen, und die es gethan und noch thun werden, die segne GOtt an Leib, Seele und Nahrung, und tröste sie in der letzten Noth, Amen.

\chapter{Das 16. Kapitel}

\section{Vorrede}

Der Mensch ist erschaffen zur Freude in Gott, ist auch nach dem Fall, durch welchen wir in Traurigkeit, Angst und Noth gerathen, wieder erlöst durch den Sohn Gottes, daß er allweg soll fröhlich sein in dem HErrn. So wird der heilige Geist auch gegeben, daß er der Tröster sei der Gläubigen, und in ihren Herzen Freude wirke in betrübtem Zustande. Derjenige welcher Mensch in Noth stecket, oder heimliches Leiden und ein betrübt Herz hat, der ängste sich selber nicht, sondern bete: Gott der Vater wolle ihm sehen lassen sein treues Vaterherz; der HErr Christus wolle ihm zeigen seine heilige tröstliche Wunden; der heilige Geist wolle sein Herz salben mit Freuden-Oel; und sei gewiß, daß ihn Gott segnen, behüten, erretten und erhalten wolle, mit Leib und Seele zur ewigen himmlischen Wonne, Freude und Seligkeit, Amen. Dies erinnert auch das 16. Kapitel.

\section{Von den drei jährlichen Hauptfesten}
\quellseite{303}

\section{Summarium}

Ist eine Ordnung mit den fürnehmsten Festen, Ostern, V. 1., Pfingsten, V. 8. und Laubhütten, V. 13., auf welche alle Mannes-Namen mußten an den Ort kommen, da der Tabernakel oder Tempel war, zum Gedächtnis der Wunderwerke und Wohltat Gottes. Und ist sonderlich zu merken, daß er befiehlet, man soll fröhlich sein für dem HErrn, sintemal er sie so gnädig regieret und geführet, und darnach mit dem Segen so reichlich begnadet hat, V. 11. Solches gnädigen Gottes soll man sich freuen, und allen Trost an ihm haben, und andere Unfälle sich nicht lassen kümmern noch traurig machen.

\section{Votum}

Der treue GOtt tröste und erhalte in wahrem Glauben alle betrübte Herzen in der ganzen Welt, und verleihe, daß wir uns Gottes und seiner Güte und Barmherzigkeit freuen, allweg in Noth und Tod, Amen.

\chapter{Das 17. Capitel}

\section{Vorrede}

Wir danken dem gütigen barmherzigen Gott, der zur Erhaltung der menschlichen Gesellschaft, auch gemeines Nutzes und des Landes und Stadtfriedens, die Obrigkeit geordnet, und bisher zum Zeugniß seiner göttlichen Providenz und väterlichen Fürsorge für uns, erhalten hat. Weil aber durch den Propheten Daniel geweissaget ist, daß am Ende der Welt die Regimente auf schwachen Beinen stehen werden, so bitten wir auch von Herzen, Gott wolle nicht mit uns handeln nach unsern Sünden, und uns nicht vergelten nach unserer Missethat, sondern wolle die wohlverdiente Strafe lindern, und wolle christliche Regimente und Regenten durch seine Gnade geben und erhalten, damit ihm eine Kirche gesammlet, und gemeiner Frieden erhalten werde, auf daß die Unterthanen ein geruhiglich und stilles Leben führen mögen, in aller Gottseligkeit und Ehrbarkeit, Amen. Dies lehret auch das 17. Capitel.

\section[Bestellung, Wahl und Amt der Richter, sonderlich eines]{Bestellung, Wahl und Amt der Richter, sonderlich eines Königs}
\quellseite{304}

\section{Summarium}

Ist eine Ordnung, wie die, so im weltlichen Regiment Richter sind, sich recht halten, und nicht Geschenke nehmen sollen, V. 19., und sonderlich ob dem fest halten sollen, daß keine Abgötterey geduldet werde, V. 21. Was im selbigen Fall mit zweien oder dreien Zeugen kann überwiesen werden, daß man nicht schonen, sondern tödten soll, V. 5.; denn das ist der fürnehmsten Werke eins, daß der Obrigkeit von Gott ist auferlegt, daß sie ob dem Wort Gottes und dem rechten Gottesdienst mit Ernst halten, und allem Aergerniß, so viel möglich ist, mit dem Schwert wehren soll.

\section{Summarium}

Ist ein Befehl, wenn sie zum Könige wählen, und wie sich derselbige halten soll. Da ist sonderlich zu merken, daß er verbeut, der König soll nicht viel Rosse noch Weiber haben, und nicht Gold und Silber sammeln, V. 16.; denn diese drei Stücke sind, damit sich die Obrigkeit gemeiniglich versündiget. Denn viel Rosse, das ist, die große Gewalt, verursachen Frevel und Muthwillen, und machen stolze Herren. Viel Weiber oder Unzucht, machen ruchlose Leute, die beide des Regiments und Gottes nicht achten. Gold und Silber sammeln, geschieht nimmermehr ohne der armen Unterthanen Schaden, so von geizigen Oberherrn übersetzt und ausgesogen werden. Darum verbeut Moses, als ein vernünftiger Regent, solches, und sagt, man solls nicht thun. Das aber soll der König thun, daß er das Wort Gottes stetig lese, und daraus lerne Gott fürchten, V. 18. Welche Obrigkeit solches thut, wird sich für den obgemeldeten dreien Stücken wohlwissen zu hüten.

\section{Votum}

GOtt erhalte unsere liebe Obrigkeit in langem Leben und beständiger Gesundheit,
\quellseite{305}
und helfe, daß sie ferner also regiere, auf daß es Ihm zu Lobe, und den Unterthanen zu Nutz und Wohlfahrt gelange, Amen.

\chapter{Das 18. Kapitel}

\section{Vorrede}

Die Zauberei ist die schöne Tugend, darüber auch der Teufel in der Hölle lachet; und ist wohl zu weisen, mit was Namen dies Laster nur könne genugsam genennet werden, daß ein Mensch, der durch Christum aus des Teufels Reich und Rachen gerissen, auch in der Taufe dem Teufel abgesagt, und Gott dem Herrn zu dienen geschworen hat, nun allererst alle solcher Sünde und Blindes vergessen, sich bloß von Gott abreißen und sich muthwillig und fürsetzlich, dem Teufel zu Dienst, übergeben. Diese überaus greuliche große Sünde läßt uns fliehen, und niemand sage, wie solche Leute den Teufel pflegen fromm zu machen: Eintracht doch alles und eitel gute Worte; denn nimmt ein einfältig, doch deutlich Gleichniß: ein gülden Becher ist ja gut und schön; also ein sammetner Rock auch; aber mit dem Becher Mistpfützen schöpfen, und mit dem sammetnen Rock sich in Koth legen, das ist so viel desto ärgere und größere Schande. Also hier, die Worte sind gut, desto ärger ists, daß sie zur Zauberei, wider Gott, sollen mißbraucht werden. Darauf höret nun an das 18. Kapitel.

\section{Vom Priesterrecht, Abgötterei, und Christo, dem rechten Propheten}

a Hie wird klärlich eine andere Predigt verheißen, denn Moses predigt, welche nicht kann das Gleiche sein, das ganzsam durch Mose gegeben, darum muß es das Evangelium sein. Und dieser Prophet niemand denn Jesu Christus selbst ist, der solche neue Predigt auf Erden gelehrt hat.
\quellseite{306}

b Vermessenheit hie redet Mose von den Propheten, so neue Lehre über die alte und vorige Lehre fürsetzen. Solchen soll man ohne Zeichen nicht glauben. Denn Gott allzeit sein neu Wort mit neuen Zeichen bestätiget. Aber droben Kap. 13, 2. redet er von den Propheten, so wider die alte bestätigte Lehre predigen; diesen soll man nicht glauben, wenn sie gleich Zeichen thun, ut supra.

\section{Summarium}

Hier ist endlich ein Befehl von den Leviten, weil sie des Gottesdienstes sollen warten, daß jedermann sie ihm lassen befohlen sein; darnach warnet er für Abgötterei und sonderlich für Zauberei und Wahrsagern, und sagt klar, wie das vermessene Sünden sind, um welcher willen Land und Leute müssen gestraft und verwüstet werden, V. 10. Am Ende ist eine Verheißung, wie Gott noch einen andern Propheten, denn Mosen, mit einer andern Lehre, denn des Gesetzes, schicken werde, V. 25. Dies ist unser lieber HErr JEsus Christus mit seinem Evangelio, den soll man hören und sein Wort annehmen. Wie es der heilige St. Petrus selbst deutet in den Geschichten der Apostel am 3. Kapitel V. 23.: Wer es nicht thut und andern Lügen-Propheten will gehorchen, deß Seele soll ausgerottet werden.

\section{Votum}

Unser lieber HErr JEsus Christus helfe, daß wir sein Wort hören und auch thun mögen, und regiere uns, daß wir uns für Abgötterei, Zauberei, falscher Lehre und allen andern Sünden allzeit hüten. Amen.

\chapter{Das 19. Kapitel}

\section{Vorrede}

Gott will ernstlich, daß, welcher Menschen-Blut vergießt, daß der wieder von ordentlicher Obrigkeit soll getödtet werden. Daher sagt unser HErr Christus Matth. 26.: Wer das Schwert nimmt, der soll durchs Schwert umkommen. Darum hat Gott auch allezeit mit sonderlichem Ernst über dem fünften Gebot gehalten. Solchen Ernst soll weltliche Obrigkeit wohl bedenken, und derhalben ihrem Amt nach keinen Todtschlag ungerochen noch abbitten lassen. Denn unschuldiges Blut schreiet zu Gott, wie Abels Blut; und wo es an dem Rechtsschuldigen nicht gerochen wird, so muß ein ganz Land solche Sünde tragen, und darum dem Zorn und Strafen Gottes herhalten. Dieß soll auch erinnern das 19. Kapitel.

\section{Ordnung der Freistädte. Falscher Zeugen Strafe}
\quellseite{307}


\section{Summarium}

Moses erholet hie das Gesetz von Freistädten, dahin sie fliehen sollten, so wider ihr Wissen und Willen ohngefähr einen Todtschlag hätten begangen, V. 3.; auf daß sonderlich die weltliche Obrigkeit lerne, wie so viel daran gelegen ist, daß niemand unschuldig umgebracht werde. Und wiederum, daß der Schuldige in seinem Wege soll verschonet werden. Wie der Text hie ernstlich sagt: Deine Augen sollen des Todtschlägers nicht verschonen, sondern du sollst das unschuldige Blut aus Israel thun, daß es dir wohlgehe, V. 13. Wenn nun durch die Obrigkeit Mord und Frevel nicht gestraft werden, daselbst kanns nicht wohl gehen, sondern es müssen sich allerlei gemeine Plagen finden. Denn wie Mose von den falschen Zeugen hernach sagt: Es müssen solche Sünden gestraft werden, auf daß die andern, so es hören, sich fürchten, und nicht mehr solche böse Stücke zu thun fürnehmen.

\section{Votum}

Helfe der treue GOtt, daß wir der Obrigkeit unterthan seyn, und daß die Obrigkeit christlich und wohl regiere, zur Rache über die Uebelthäter, und zu Lobe den Frommen, Amen.

\chapter{Das 20. Kapitel}

\section{Vorrede}

Wie der Land- und Stadtfriede ist ein halber Himmel und ein irdisch Paradeis; also ist der Krieg eine halbe Hölle, oder ja ein Vorbote und Vorläufer der Höllen. Denn wie der Friede eine Gabe Gottes ist, daraus alle andern zeitlichen Güter fließen und quellen: also, wenn Gott mit Krieg straft, so gibt immer ein Unglück dem andern die Hände. Derwegen danken wir Gott heute herzlich, daß wir guten Frieden bisher haben, und laßt uns zuschauen, daß wir also leben und also beten, daß wir den güldenen Frieden auch ferner behalten. Denn die Furcht des HErrn gibt reichen Frieden und Heil, sagt Sirach im 1. Kapitel. Ferner höret davon an das 20. Kapitel.

\section{Kriegsregeln}

a Ein Haus einweihen, heißt auf hebräisch, wenn man drein zieht, drinnen zu wohnen.
\quellseite{308}
Ist's doch Holz auf dem Felde, und nicht Mensch, daß es vor dir ein Bollwerk sein möge. Es ist Holz auf dem Felde, und nicht in der Stadt, es thut dir nichts, und ist dir nütze. Hic sensus congruit allegoria, non esse pugnandum contra eos, qui non sunt contra nos, sed pro nobis. Der Sinn des Gleichnisses ist: man soll nicht gegen die streiten, die nicht gegen uns, sondern für uns sind.

\section{Summarium}

Ist erstlich eine feine Ordnung, wie man dem Kriegsvolk im Streit soll zusprechen, daß sie sich nicht entsetzen sollen für den Feinden, sondern ihre Hoffnung darauf setzen, daß Gott, welcher ihnen aus Egypten und allerlei Noth treulich geholfen habe, bei ihnen sein, für sie mit ihren Feinden streiten und ihnen helfen wolle. Solcher Trost ist gewiß und fehlet nicht, V. 1. Aber es gehöret auch dazu, daß man eine gute Sache habe, denn zum Unrecht will Gott niemand helfen. Welche Obrigkeit aber eine gute Sache hat, die mag Schwert, Büchsen und Harnisch brauchen; aber das Vertrauen soll allein auf Gottes Hülfe stehen, so wird sich der Sieg endlich auch wohl finden.

Darnach sehet Moses etliche andere seine und billige Ordnung, so in Kriegsläuften zu gebrauchen sind, V. 10. Deren zum Theil man sich noch könnte brauchen. Aber es ist leider bei unserm Kriegsvolk seine Zucht noch Gottesfurcht. Moses verbeut, sie sollen keinen fruchtbaren Baum in der Feinde Land beschädigen. Denn, spricht er, es ist ein Holz und nicht ein Mensch, V. 19. So verwüsten unsere Krieger muthwillig auch der Fremden Güter, Weinberge, und verderben, was sie können. Daß man wohl gedenken mag, es sind nicht Menschen, sondern lauter Teufel. Was sollte denn für Glück bei ihnen sein.

\section{Votum}

Der treue, gütige GOtt wende ab die wohlverdienten Strafen, erleuchte die Kriegsleute zu seiner Furcht und Gerechtigkeit, und behüte uns ja für Krieg und Blutvergießen, um JEsu Christi, seines lieben Sohnes und unsers HErrn, willen, Amen.

\chapter{Das 21. Kapitel}

\section{Vorrede}

Es sind nicht allein Mörder die, so mit wehrhafter Faust, oder sonst mit Gewalt und Gift,
\quellseite{309}
einen Menschen todtschlagen und ermorden; sondern die Kinder, welche ihre Eltern übel halten, oder mit ihrem Ungehorsam und Schandleben ängsten, daß sie für Sorge und Grämniß sterben müssen, das sind auch Vater- und Mutter-Mörder. Wehe aber beiden. Denn kein Mörder bleibt ungestraft, und es ist auch nie keinem ungehorsamen Kinde glücklich und wohl gegangen, und wenn sie sich nicht bekehren, so ist ihr letzter Lohn höllisch Feuer und ewige Verdammniß. Denn welche solches thun, werden das Reich Gottes nicht erben, sagt St. Paulus zum Galatern im 5. Kapitel. Dies lehret auch das 21. Kapitel.

\section[Vom unbekannten Todtschlage; gefangenen Weiberpersonen; Recht des Erstgebornen;]{Vom unbekannten Todtschlage; gefangenen Weiberpersonen; Recht des Erstgebornen; ungehorsamen Sohne; Gelübden}
\quellseite{310}

\section{Summarium}

Ist eine Ordnung, wie es erstlich soll gehalten werden, wenn einer im Felde erschlagen wird, da man nicht weiß, wer es gethan hat, daß die Priester und Obrigkeit zusammen kommen, und Gott bitten sollen, daß er solche Sünde nicht am ganzen Volk strafen wolle, weil man den Thäter nicht haben und strafen kann, V. 1. Da sehen wir abermal, wie eine große Sünde es ist, todtschlagen und den Todtschläger nicht strafen. Darnach sind etliche andere Ordnungen, unter welchen sonderlich zu merken ist, wie Gott die ungehorsamen Kinder, so Vater und Mutter übel halten, bestraft, daß man sie tödten soll; auf daß andere sich bessern und lernen fromm und gehorsam sein.

\section{Votum}

Der starke GOtt steure allen blutgierigen Menschen, behüte uns für derselben Zorn und Mord, und regiere die aufwachsende liebe Jugend, daß sie ihren Eltern gehorsam sei und folge, und an Leib und Seele gesegnet werde, Amen.

\chapter{Das 22. Kapitel}

\section{Vorrede}

Unser Seligmacher JESUS Christus sagt Matth. am 5.: Selig sind die Barmherzigen, denn sie werden Barmherzigkeit erlangen. Und: Selig sind die reines Herzens sind, denn sie werden Gott schauen. Daraus folget, daß Gott denen, so unbarmherzig gegen ihrem Nächsten sein, auch keine Barmherzigkeit erzeigen wolle; und die Unzucht lieben und üben, daß sie zeitlich an Leib, Seele, Gut und Ehre sollen gestraft werden, und, da sie in solchen Sünden ohne wahre Buße sterben, daß sie auch in Ewigkeit von Gottes Angesicht, zu ewiger Pein und Marter, verstoßen sein sollen. Dergleichen läßt uns Barmherzigkeit an unserm Nächsten beweisen, und alle Unreinigkeit unsers Leibes fliehen und meiden. Davon höret auch an das 22. Kapitel.

\section{Mancherlei Gesetze verzeichnet}
\quellseite{311}

\section{Summarium}

Ist eine Ordnung, wie man verloren Vieh und anderes bewahren, und dem wieder zustellen soll, der es verloren, V. 1. Sonderlich aber ist hie zu verwundern, daß Gott befiehlt und verheißet: Wer ein Vogelnest ausnimmt, und die Mutter fliegen läßt, daß demselben Menschen wohl gehen und er langes Leben haben soll, V. 6. Denn Gott will, daß wir barmherzig sollen sein, auch gegen das unvernünftige Vieh, und hat ein Ungefallen dran, wo man tyrannisch handelt, wenn es gleich nur um ein armes Vöglein zu thun ist. Wir müssen aber bei diesem Befehl annehmen, so es Gott wohlgefällt, wo man mit den unvernünftigen Thierlein freundlich umgeht, wie viel ein größer Wohlgefallen er an dem haben, und mit reichlichem Segen ers belohnen wolle, wo ein Mensch dem andern freundlich und behülflich ist.

Es ordnet auch Moses hie, wie man Ehebrecher und Jungfrauschänder steinigen und tödten soll, V. 22. Auf daß man lerne, wie große Sünde dies sei, obs wohl weltliche Obrigkeit nicht so heftig strafet, wiewohl vonnöthen wäre. Und ist gar kein Zweifel dran, daß deshalb viel gemeiner Strafen täglich über uns gehen, daß solche Laster so gar keine Strafe haben.

\section{Votum}

Gott breche aller Unbarmherzigen Muth und Sinn, und steure dem Schandgeist und aller Unzucht, und regiere uns und die liebe Jugend, daß wir uns für allerlei Sünde und Schandlastern hüten, und dagegen züchtig und christlich leben, Amen.

\chapter{Das 23. Kapitel}

\section{Vorrede}

Es ist kein größerer Menschenfeind auf Erden nach dem Teufel, denn ein Geizhals und Wucherer; und sollen derwegen fromme Herren solche Wucherer und Geizhälse ansehen als leibhaftige Teufel und sich für ihnen segnen, wo man sie höret oder siehet, und sollen gedenken, daß Türken, Tartarn und andere Heiden eitel Engel sind gegen einem Wucherer. Desgleichen sollte man die Jugend lehren und gewöhnen, daß sie erschrecken und pfui dich sagen für dem Namen Wucherer, als für dem ärgsten Teufel. Weil dem also, so helfe Gott, daß die, welche solcher Sünden sich schuldig wissen, rechte wahre Buße thun, und regiere alle, die an diesem Hauptlaster unschuldig sein, daß sie nimmermehr in diese Sünde, zu ihrem und der Ihren Unheil, gerathen, Amen. Dies erinnert auch das 23. Kapitel.
\quellseite{312}

\section{Wer in die Gemeine des Herrn gehöre, oder nicht}

a Bösen) daß du selbst nicht unsträflich seiest, und also den Sieg zur Strafe verlierst und geschlagen werdest, wie zur Zeit Eli und Saul geschah.

\section{Summarium}

Ist eine Ordnung von mancherlei Sachen, unter welchen sonderlich diese zu merken sind, wie Moses befehlet, wenn man wider die Feinde ziehen soll, daß jedermann sich für dem Bösen hüten soll, v. 9. Denn weil man Gottes Gnade und Hülfe zum Sieg bedarf, soll man sich nicht versündigen und Gott nicht Ursache zur Strafe geben; wie man stehet, wie es mit unserm Kriegsvolk leider gehet, da alle Unzucht und Gotteslästerung
\quellseite{313}
 auf das greulichste geübt wird; darum ist auch so wenig Glück dabei.

Darnach befiehlt er, daß kein Jude am andern soll wuchern, V. 19. Viel mehr sollte solches kein Christ gegen den andern thun. Einmal wir durch Christum zu Brüdern unter einander gemacht und einer dem andern um christlicher Liebe willen, und nicht um Eigennutz willen, gern dienen und helfen sollte, wie denn alle Christen thun. Die es aber nicht thun, sind keine Christen, ob sie gleich den Namen haben; müssen derhalben zum Teufel und in das höllische Feuer.

\section{Votum}

Der starke und eifrige GOtt steure allen, die muthwillig sündigen, und helfe, daß ein jeder gegen seinen Nächsten also handele, daß sich GOtt wieder seiner aus Gnaden erbarme, und ihm helfe an Leib und Seele, bis zum ewigen Leben, Amen.

\chapter{Das 24. Capitel}

\section{Vorrede}

Welcher Mensch sich mit Schaden seines Nächsten nähret, der erarbeitet wahrhaftig nicht, und je glückseliger es ihm am Anfange und eine Zeitlang dazu gehen mag, je mehr Verderben und größerer Schaden endlich folget. Dies lehret mit seiner Vermahnung und Warnung der heilige Apostel Paulus, da er 1 Thess. 4. also sagt: Das ist der Wille Gottes, daß niemand zu weit greife, nicht übervortheile seinen Bruder im Handel; denn der HErr ist ein Rächer über das alles. Derwegen ein jeder seinem Nächsten in seinem Beruf, wo und wie er kann, diene und helfe; er bete auch, und sei gottsfürchtig daneben, so wird Gottes Segen bei seiner Nahrung sein und bleiben. Ferner höret an das 24. Kapitel.

\section{Gesetze, vornehmlich den häuslichen Stand betreffend}
\quellseite{314}

\section{Summarium}

In eine Ordnung von dem Scheidebrief, V. 1., davon Christus klärlich sagt, Moses habe ihn den Juden erlaubet, nicht daß er recht sei, sondern daß sie so ein störrig, unverträglich Volk sind. Darnach ordnet er, wie mans mit Leihen und Pfand nehmen dermaßen halten soll, daß den Armen damit geholfen und sie nicht mehr beschweret werden, V. 10., und sonderlich gebeuget er, man soll die Pfande über Nacht im Hause nicht behalten, der die armen Leutlein ohne ihren Schaden nicht entbehren können, V. 13. Solche und dergleichen andere Lehren, so hie angezeigt sind, sollten wir bedenken und uns darnach halten, so würde der Untreue des Geizes, Schindens und Kratzens weniger werden. Weil es aber nicht geschiehet, sehen wir, daß ein Unglück nach dem andern sich findet.

\section{Votum}

Helffe der treue GOtt, daß sich ein jeder in seinem Beruf treulich und christlich nähre, und seinem Nächsten gern und nutzbarlich diene, damit ihm GOtt wieder diene zu aller Nothdurft, um Christi JEsu willen, Amen.

\chapter{Das 25. Capitel}

\section{Vorrede}

Der heilige Apostel Paulus vergleicht die Geizigen den Abgöttern. Denn gleich wie ein abgöttischer Mensch seinen Schöpfer verläßt, und seine Hoffnung und Zuversicht setzet und gründet auf die Creaturen: also thut auch ein Geizhals, welcher Gold und Silber für seinen Gott hält. Das ist aber sehr schrecklich zu gedenken, will geschweigen zu reden, daß ein Mensch eine stumme tote Kreatur seinem Schöpfer soll vorziehen. Zu dem, so ist das Gottes Ordnung für und für, wo er das Leben gibt, daß er auch schafft, daß man es erhalten könne. Dieweil er denn solches dem Viehe thut, wie viel mehr will er solches den Menschen und sonderlich seinen Christen thun, welchen er nicht allein Leib und Leben, sondern auch seinen eingebornen Sohn gegeben hat, daß sie nicht allein zeitlich, sondern in Ewigkeit leben sollen. Dergleichen für dem Geiz sich jedermann hüten soll. Darauf höret auch davon an das 25. Kapitel.

\section{Etliche Stücke jüdischer Polizeiverordnung}
\quellseite{315}

\section{Summarium}

Ist erstlich eine Ordnung, wie man die mit Schlägen strafen soll, so bußfällig und dennoch den Tod nicht verdienet hätten. Denn Sünde soll nicht ohne Strafe hingehen. Darnach ordnet er anders mehr, und ist sonderlich zu merken, wie er bräuet, wer unrecht Gewicht, Maß, Ellen brauchet, daß ihn Gott strafen werde. Wer aber recht Gewicht und gute völlige Maße gebe, den werde Gott segnen, V. 15. Solche Sprüche sollten die Kaufleute in die Kräme an Waage und Gewicht schreiben und malen lassen, und sich darnach halten. Aber der leidige Geiz stopfet ihnen die Ohren, und hält ihnen die Augen zu, daß sie es verachten. Darum verursachen sie Gott, daß es unter ihnen wüste umgehet, jetzt da einen, jetzt dort einen am Leib, Gut und andern solcher Sünden halben strafen muß. Es wird aber endlich, wo Besserung nicht folget, auch das höllische Feuer über sie kommen, wie Paulus sagt: Die Geizigen werden das Reich Gottes nicht erben.

\section{Votum}

Verleihe der treue GOtt, daß die Obrigkeit die Sünden strafe, und daß für Betrug, Geiz, und anderen Sünden wir uns hüten, und daß wir dem Nächsten brüderlich rathen und dienen, um unsers lieben Bruders Christi Jesu willen, Amen.

\chapter{Das 26. Capitel}

\section{Vorrede}

Die heilige Schrift lehret, daß die Kinder Gottes sein Erstlinge seiner Kreaturen, Jacobi 1, 18. (Gleich wie man aber im alten Testament die Erstlinge allerlei Früchte dem HErrn geheiliget und ergeben, oder Gott zur Dankbarkeit geopfert hat: also sollen wir uns auch im neuen Testamente Gott heiligen und opfern als Erstlinge seiner Kreaturen und seine neugborne Kinder, das ist, wir sollen unsern Leib und alle innerliche und äußerliche Gliedmaßen also gebrauchen, daß Gott täglich von uns geehret und gehorsamlich dienet werde, und daß wir nach seinem Willen leben, ein jeder nach seinem Beruf, Amt und Stande. Davon lehret St. Paulus Röm. 12, 1. also: Euer Leiber gebet zum Opfer, das da lebendig, heilig und Gott wohlgefällig sei, welches sei euer vernünftiger Gottesdienst. Dies erinnert auch das 26. Kapitel.

\section{Von den ersten Früchten und Zehnten}
\quellseite{316}

a Leide) Gottes Opfer soll fröhlich, heilig und rein sein; darum soll nichts in Traurigkeit davon gegessen, nichts in Unreinigkeit davon genommen, nichts den Götzen oder Todten davon gegeben sein.

\section{Summarium}

Dieß ist die Ordnung von den Erstlingen der Früchte, von Zehnten, V. 2., in welcher wir sehen, wie Gott immerdar befiehlt, daß man seinen Segen mit Dank annehmen, und denselbigen dazu brauchen soll, daß zum ersten die Leviten, das ist die, so die Kirchendämter verwalten, versorget, und darnach Wittwen, Waisen und Fremdlingen damit geholfen werde. Wer nun sein Gut also brauchet, der mag sich trösten, wie die Juden hie bitten, daß GOtt solches sehen, und sie und das Land mildiglich segnen werde, V. 12., da sonst bei den Geizigen und die sich der Dürftigen nicht annehmen, weder Glück noch Heil ist.

\section{Votum}

Der treue GOtt steure allem Geiz und Unbarmherzigkeit; und die des Armuths und seiner Kirchen sich treulich annehmen, die segne GOtt an Leib und Seele, zeitlich und ewig, um JEsu Christi seines lieben Sohns willen, Amen.

\chapter{Das 27. Capitel}

\section{Vorrede}

Was Gott den Unbußfertigen dräuet in seinem Wort, das läßt er wahrhaftig ergehen und geschehen, wie lange auch die Strafe außen bleibt. Derwegen, ach, laßt uns des heiligen Wortes Gottes Flüche und Dräuungen nicht mit tauben Ohren und steinernen Herzen, sondern mit Andacht fleißig anhören; und erschrecket alle für
\quellseite{317}
Gottes Zorn, hütet euch für allerlei Sünden, und lebet, wie es euch Gott ernstlich in seinem Gesetz und Wort befohlen hat. So wird nicht Gottes Fluch, sondern sein Segen über euch kommen, und also wird Gott euer lieber Vater, ihr seine Kinder, der HErr Christus euer treuer Bruder, ihr auch seine Brüder und Schwestern, der heilige Geist euer beständiger Lehrer und Tröster, ihr seine Tempel und Wohnungen sein und bleiben. Denn also sagt der HErr Christus auch selber Joh. 14, 23.: Wer mich liebet, der wird mein Wort halten, und mein Vater wird ihn lieben, und wir werden zu ihm kommen und Wohnung bei ihm machen. Darauf höret nun das 27. Kapitel.

\section{Denksteine der Kinder Israel. Segen und Fluch}

\section{Summarium}

Ist eine Ordnung, wie die Leviten öffentlich verfluchen sollen, alle, die wider das Gesetzes
\quellseite{318}


HErrn thun, und die ganze Gemeine soll dabei stehen und Amen dazu sagen, V. 15. In solchem Fluch werden unter andern sonderlich genennet, ungehorsame Kinder, so ihren Eltern fluchen, V. 16.; auch weltliche Obrigkeit, die den Fremdlingen, Wittwen und Waisen nicht zum Rechten verhelfen, sondern Unrecht thun, V. 19. Und ist hie sonderlich zu merken, daß solcher Fluch allein auf die Sünden gehe, so ohne Mittel wider Gott und das Gewissen, oder wie mans sonst nennet, wider das natürliche Recht sind.

\section{Votum}

Der Vater aller Barmherzigkeit, GOtt, unser lieber HErr, sei uns allen gnädig, helfe, daß wir in seinen Geboten wandeln, und erhalte uns in seinem Segen, an Leib und Seele, bis zum ewigen freudenreichen Leben, das uns erworben hat unser HErr JEsus Christus, Amen.

\chapter{Das 28. Capitel}

\section{Vorrede}

Das Wort Gottes ist der unvergängliche Same, wie St. Petrus spricht, dadurch wir neu geboren werden. Sintemal der Glaube, der uns zu Kindern Gottes macht, kömmt aus der Predigt des Evangelii, dadurch wir in unsern Herzen, Verstand und Willen erneuert werden, daß, wie wir vorhin unserer ersten Geburt nach, ohne alle Erkenntniß und ohne alle Furcht Gottes lebten und unser Herz voll böser Lust und Begierde waren: also gewinnen wir, der andern Geburt nach, ein herzliches Mißfallen über die Sünde, und ein Gefallen, daß wir Gott fürchten und lieben, ihm vertrauen, daß wir uns befleißen eines gottseligen Lebens und Wandels, und daß wir Lust und Freude am Gesetz des HErrn haben. So helfe Gott, daß zu unserer christlichen Besserung sonderlich viel diene das 28. Kapitel des fünften Buches Mos.

\section{Verheißener Segen, gedrohter Fluch}

a Korb) das ist alles, was du gegenwärtiglich brauchest, und was du beiseit legest zu behalten.

\section{Summarium}

Es ist die Beschluß-Predigt. Denn nachdem Moses wiederholet hat die fürnehmsten Stücke, was sie thun und lassen sollen, wenn sie Gott gefallen und recht thun wollen, beschleußt er

18*
\quellseite{319}
nun und sagt erstlich, wie Gott sie segnen und solchen Gehorsam gegen seinem Wort reichlich bezahlen wolle, mit allerlei Gaben und Gnaden, V. 1. Auf daß wir auch lernen, wer da begehrt, daß er Glück und Heil haben möge, daß er sich nach Gottes Wort und Willen halte. Wie Christus sagt Matth. 6, 33.: Trachtet am ersten nach dem Reiche Gottes und seiner Gerechtigkeit, so wird euch das andere alles zufallen.
\quellseite{320}

b Alii infantes recens natos. D. i. Andere: neugeborne Kinder.

\section{Summarium}

Ist das andere Stück der Bußpredigt Mose, in welcher er warnet, daß man sich vor Sünden hüten wolle, sintemal Gott dieselbigen mit allerlei Plagen strafen will. Da sollt du lernen, was die rechte Ursache sei, daß Krieg, Pestilenz und andere Krankheit, Hunger und Kummer in der Welt ist, nämlich, daß solches unsere eigene Sünden verursachen. Darum sollten wir von Sünden ablassen, so würden solche Strafen auch ausbleiben.
\quellseite{321}

\section{Summarium}

Dies ist eine helle Anzeigung, warum solche grausame Plagen kommen sollen, nämlich, daß man der Stimme des HErrn nicht gehorchet, und nicht thut, was Gott in seinem Gesetz befohlen hat. Und ist sonderlich hie zu merken, daß Gott an solchen leiblichen Strafen nicht gesättiget ist, sondern auch die höchsten und geistlichen Strafen dräuet, daß man andern Göttern, Holz und Stein, mit unruhigem, erschrockenen Gewissen dienen soll, wo man dem rechten Gott nicht dienen und seinem Wort nicht hat mit gutem Gewissen folgen wollen, v. 64.

\section{Votum}

Verleihe der treue, barmherzige GOtt, daß wir dieser Lehren nimmer vergessen, sondern daß wir nach diesen allezeit leben, in Gerechtigkeit und Heiligkeit, und also thun was recht ist, daß uns GOtt Gutes thue, wie er uns verheißen hat, Amen.

\chapter{Das 29. Capitel}

\section{Vorrede}

Wer Gottes gutes heiliges Wort höret, und doch böse und sündliche Werke thut, den bedrohet Gott mit Bösem. Denn wie sich ein Mensch gegen Gottes Wort hält, also hält sich Gott wiederum gegen denselbigen Menschen. Derwegen wer nicht will, daß ihn Gott verachten soll, der verachte nicht Gottes Wort. Und wer da will, daß ihm Glück und Segen folgen soll, der folge auch Gottes Wort, daß er nach demselbigen gehorsamlich lebe. Denn die Gottes Wort hören und thun, oder wer den Willen thut Gottes, des Vaters im Himmel, derselbige ist des HErrn Christi, Matth. 12, 50., das ist, er will ihnen Liebe und Treue erzeigen zu aller Wohlfahrt, mehr denn das Geschwister einander, ja mehr, als die Eltern ihren Kindern Gunst und Dienst erzeigen können. Darauf höret ferner an das 29. Kapitel.

\section{Erneuerung des Bundes}
\quellseite{322}

a Eurem Thun) ohne Gottes Wort ist all unser Thun Narrheit.

b Es gehet mir wohl) das ist der ruchlosen Leute Wort und Gedanken: ei die Hölle ist nicht so heiß; es hat nicht Noth, der Teufel ist nicht so greulich, als man ihn malet. Welches alle Werckheiligen frech und durstiglich thun, ja noch Lohn im Himmel gewarten.

c Weil ich wandele) das ist, weil ich so thue, und meine Abgötterei treibe, so ist eitel Glück da; wer Gott dienet, dem gehets nimmermehr wohl, Apost. 8, 23.
\quellseite{323}

d Das Geheimniß will so sagen: Uns Juden hat Gott für allen Völkern auf Erden seinen Willen geoffenbaret, und was er im Sinn hat. Darum sollen wir auch desto fleißiger sein.

\section{Summarium}

Das ist eine Verneuerung des Bundes, weil sie alle Wunderwerke in Egypten gesehen, V. 2., und hernach erfahren haben, wie sie Gott in der Wüste erhalten, V. 5., ihnen wider die Heiden Sieg gegeben, und sie in das verheißene Land eingesetzt hat, daß sie sich Gottes Willens und seines Wortes halten sollen, und sonst keinem andern Gott dienen, sonst soll alles Unglück über sie kommen, V. 18. Solches soll uns eine Verwarnung sein, daß wir uns für Sünden hüten und nach Gottes Wort thun sollen, sonst wird dergleichen Fluch über uns auch kommen.

\section{Votum}

GOtt, unser treuer, lieber Vater, lasse seinen Segen über uns und all den Unsern bleiben, und helfe, daß wir nach seinem Wort thun und leben, wie er uns befohlen hat, und wir ihm solches schuldig sind als seine Kinder, Amen.

\chapter{Das 30. Kapitel}

\section{Vorrede}

Gottes Barmherzigkeit hat noch kein Ende, sondern sie ist alle Morgen neu, sagt der heilige Prophet Jeremias in Klagliedern im 3. Kap. V. 22. Und die heilige Jungfrau Maria singet schön und tröstlich in ihrem Magnificat, Luc. 1, 50.: Gottes Barmherzigkeit währet immer für und für, bei denen, die ihn fürchten. Diese Barmherzigkeit Gottes sollen wir also gebrauchen, daß ein jedes seine Sünde erkennen, Gott dieselbigen alle herzlich abbitten, und sich trösten soll, obschon seiner Sünden viel und die größten sind, daß Gottes Barmherzigkeit tausendmal mehr und größer ist, der ihm alle seine Sünde vergeben will; doch, daß niemand auf Gottes Barmherzigkeit muthwillig sündige. Denn wer auf Barmherzigkeit sündiget, dem wird mit Unbarmherzigkeit endlich abgelohnt. Darauf höret an das 30. Kapitel.

\section[Die Bußfertigen erlangen Gnade, die Widerspenstigen Rache und]{Die Bußfertigen erlangen Gnade, die Widerspenstigen Rache und Strafe}

a Zu gut denn die Gottlosen haben auch wohl Ehre und Gut, oft mehr, denn die Heiligen, aber zu ihrem und Andern Verderben.
\quellseite{324}

\section{Summarium}

Ist eine tröstliche Verheißung, wie Gott wiederum zu Gnaden sein Volk wolle annehmen, wenn sie nur in sich selbst schlagen und lassen ihnen ihre Missethat leid sein, und begehren sich zu bessern, V. 2. Solche Verheißung gehet fürnehmlich auf das Neue Testament. Obwohl das Volk ihrer Sünden halben mancherlei soll gestraft werden, daß dennoch Gott Christum senden, und durch den heil. Geist ihre Herzen beschneiden und sie recht fromm machen werde. Darum vermahnet er am Ende, weil er beides ihnen habe fürgelegt, das Gute, so denen nachfolgen soll, die Gott lieben, und sein Wort halten, und den Fluch, darauf die Sünden gehet, daß sie sich zum Guten halten und fromm sein sollen.

\section{Votum}

Der treue, gütige GOtt sei uns allen gnädig, und vergebe uns alle unsere Sünde, und strafe uns ja nicht nach Verdienst, sondern segne und behüte uns nach seiner großen Barmherzigkeit, um Jesu Christi, unsers Heilandes und Mittlers, willen, Amen.

\chapter{Das 31. Capitel}

\section{Vorrede}

Gott schicket uns Kreuz und Unglück zu, daß er den zukünftigen Sünden wehre und uns des Fleisches Kitzel vertreibe. Denn wenn wir ohne Kreuz und Widerwärtigkeit sein, und es uns nach unserm Wunsch und Willen gehet, so wird unser Fleisch fürwitzig, muthwillig und sicher, und gewinnet Lust und Liebe zu sündhaften Dingen; aber wenn wir unter dem Kreuze liegen, so vergehet uns die Geilheit. Daß also gewiß ist, wir sind frömmer, wenn es uns übel gehet. Und dies beweiset das Exempel des Königs Davids. Wie heilig war doch dieser Herr, da ihn Saul und darnach Absalon verfolgte; aber da er Ruhe und Friede hatte, fiel er und sündigte schändlich. Derwegen wir die Vater-Ruthe unsers lieben Gottes küssen und in Geduld das Kreuz tragen sollen. Davon höret auch an das 31. Capitel.

\section{Mose sagt sein Amt auf; ordnet Josua an seine Stelle}
\quellseite{325}

\section{Summarium}

Sie ist ein fein Exempel, wie Moses erstlich das Volk, darnach seinen Diener Josua, der fortan ihr Hauptmann sein sollte, tröstet, daß sie beide auf GOtt ihr Vertrauen setzen, und sich für den Feinden nicht entsetzen sollen, V. 3.

Denn wer GOtt auf seiner Seite hat, der wird auch wohl mitten im Tode sicher sein; dagegen wo GOtt nicht dabei ist, und die Hand abzeucht, da kann weder Geld noch Gut, weder Gewalt noch großes Vermögen helfen oder retten. Solchen Trost soll alle Obrigkeit fürnehmlich haben und durch fleißiges Gebet täglich suchen.
\quellseite{326}

\section{Summarium}

Ist ein fein Exempel, wie Gott gern aller Ursache vorkommen wollte, daß sein Volk nicht sündigte. Aber Art läßt von Art nicht, sondern es gehet, wie der HErr hie sagt: Wenn wir wohl stecken und genug haben, so vergessen wir unsers HErrn Gottes und seines Worts. Darum muß er mit allerlei Unglück hinter uns her sein, auf daß wir nicht sicher werden und in Sünde fallen.

Eben um dieser Ursache willen befiehlt er droben, daß man dies Buch für dem ganzen Volk lese, und das Lied, so hernach folget, singen soll, V. 24. Denn wo uns Gottes Wort nicht erhält, ists nicht möglich, daß wir anders können, denn sicher und gottlos sein, wenn es uns wohl gehet und keinen Anstoß haben.

\section{Votum}

Der treue Gott helfe, daß wir nach seinem Wort in Gehorsam leben, ihm vertrauen, und unter allem Kreuz von ihm Trost, Geduld und Hülfe haben und behalten, bis zur fröhlichen Erlösung, Amen.

\chapter{Das 32. Kapitel}

\section{Vorrede}

Weil die Kinder sind die edlen Pflänzlein, welche Gott in seinem Lustgarten des Ehestandes sollen erzogen werden; so sollen die Eltern nicht gedenken, es sei genug und alles gethan, wenn die Kinder getauft sind. Es ist wahr, sie sind als eine wilde Art abgebrochen und durch die heilige Taufe eingepfropfet und einverleibet der Unschuld Jesu Christi. Aber da sollen die Eltern thun, als fleißige Gärtner und Haushalter des HErrn. Wie ein milder Stamm noch immer ausschlägt in Wasserreiser, und der Gärtner dem Bösen kaum wehren, die guten Reiser aber schwerlich fortbringen kann: also gehet es auch mit den Kindern. Die alte Adams-Art will immer herfür, wie die Schobreiser. Soll nun das Gute, das in sie gepflanzet, nicht zum Abnehmen kommen, so müssen die Eltern ihnen wehren, und mit christlicher Zucht die Uebermaße abschneiden, und solches, weil der Baum noch zart und jung ist. Ach dieses helfe allen lieben Eltern der treue Gott, Amen. Darauf höret auch an das 32. Kapitel.

\section{Moses Lied und Lobgesang}

a Allein) laßt unserm Gott allein Gott sein, und alle Ehre haben und keinem andern.

b Fels) die hebräische Sprache heißt Gott einen Fels, das ist, ein Trotz, Trost, Hort und Sicherung, allein, die sich auf ihn verlassen und ihm vertrauen.

c Gerecht) bei den Gottlosen muß Gott immer unrecht haben und sich meistern lassen, Matth. 11, 19.: Die Weisheit muß sich lassen rechtfertigen von ihren Kindern, die wissen alles besser; wie es Gott macht, so taugts nicht.
\quellseite{327}

d Nicht kannten) hatten sie keinen Befehl noch Gottes Wort gehabt, sondern solchen neuen Gottesdienst aus eigener Andacht erwählet haben.

\section{Summarium}

Sie ist ein Lied, in welchem Moses das Volk lehret, wie Gott ihnen so gnädig gewesen und so große Hülfe und Wohlthat ihnen bewiesen hat, daß ers zum eigenen Volk angenommen, aus Egypten geführet, ihnen das Gesetz gegeben, und endlich ins gelobte Land gebracht hat, V. 7. Darum vermahnet er nach dem ersten Gebot, daß sie solchen Gott ehren, auf seine Hülfe allein sich verlassen und sonst keinem fremden Gott dienen sollen. Aber er weissaget, daß sie bei solcher Lehre und Pflicht nicht bleiben werden; sondern wenn sie satt, dick und fett sein werden, werden sie solchen Gott fahren lassen und Abgötterei treiben, V. 15. Und zum Exempel, daß wir nimmer fleißiger beten und weniger sicher sein sollen, denn wenn wir Friede haben und es uns wohl geht.

e Beräuben, schrecken) das ist: außen sollen sie Wittwen und Waisen werden durch Schwerdt und Gefängniß der Kinder, Männer und Weiber beraubt. Was aber innen bleibt, soll durchs Schrecken, das ist, durch Hunger, Pestilenz, Aufruhr jämmerlich umkommen.

f Rath) sie achten Gottes Wort nicht, wissen alles besser.

g Versiegelt) sie glauben nicht, bis sie es erfahren. Denn es ist für ihren Augen verborgen.
\quellseite{328}

b Custoditum et neglectum, id est, servatam rem et contemptam, quia tempore pacis contemptum est, quod tempore male libenter colligeres, si adesset pretiosum et vile. Das ist, das Verwahrte und Verwahrlosete, d. i. das Beobachtete und Verachtete; weil zur Zeit des Friedens und Wohlergehens verachtet ist, was man zur Zeit des Unglücks und der Noth gern zusammenhielte, wenn es da wäre, Kostbares und Geringes.

i Ueber dem Blut) das sind drei Strafen des Schwerts. Die erste, daß ihrer viel erschlagen werden. Die andere, daß sie gefangen geführet werden. Die dritte, daß ihr Haupt soll bloß werden, das ist, ihr Königreich und Priesterthum soll von ihnen genommen werden, welches durchs Haar auf dem Haupt bedeutet ist.

\section{Summarium}

Ist der andere Theil dieses Liedes, in welchem Moses weissaget, wie sie um Abgötterei und anderer Sünde willen sollten mit allerlei Jammer und Unglück geplaget und endlich ganz und gar verstöret und vertilget werden, und Gott ihm ein ander Volk, die Heiden, werde auserwählen, wie denn geschehen ist. Alles darum, daß sie Gott und seines Worts nicht achteten. Uns zum Exempel. Denn so Gott seines eigenen Volks nicht hat verschonet, dürfen wir nicht denken, wenn wir in gleiche Sünde fallen, daß wir nicht auch gleichen Zorn und Strafe tragen sollen.

\section{Summarium}

Es ist ein fein Exempel, wie Mose sonderlich den Eltern befiehlet, daß sie ihre Kinder zu Gottes Wort halten, ihnen dasselbe für sagen und mit Ernst zu halten befehlen sollen. Denn es sei ein Wort des Lebens, sintemal die es verachten und nicht mit Fleiß halten, darum sterben und verderben müssen. Solche Lehre ist trefflich hoch vonnöthen. Denn die Furcht Gottes, so allein aus dem Wort Gottes kömmt, ist der beste Zuchtmeister; wo dieselbe nicht ist, werden die Kinder ungerathen, und leben ihren Eltern zu Schanden und Schaden. Solchem Jammer ist nicht besser zu wehren, denn man hebe in der Zeit an, die Kinder zur Kirche und Gottes Wort zu gewöhnen, wie Moses hie heißt. Darum ist ein Exempel, wie Moses seiner Sünden halber gestraft wird, daß er nicht in das gelobte
\quellseite{329}
Land kommen soll, V. 48. Wie der Psalm solches auch anzeigt, Gott habe die Sünde seiner Heiligen gestraft; auf daß wir Gott und seinen Zorn lernen fürchten, und uns mit Ernst für Sünden hüten.

\section{Votum}

Regiere uns alle der barmherzige GOtt, daß wir für Sünden uns hüten, und in wahrer Buße leben und sterben. Derselbe treue GOtt helfe auch, daß alle Eltern ihre Kinder wohl ziehen, ihm zu Lob, ihnen selbst zu Ehren, und zur Wohlfahrt den Kindern, Amen.

\chapter{Das 33. Kapitel}

\section{Vorrede}

Ein christlicher Wunsch ist ein christliches und kräftiges Gebet; denn solche herzliche Wünsche sind gegründet auf die Zusagung Christi: Wahrlich, wahrlich, ich sage euch, was ihr den Vater bitten werdet in meinem Namen, das wird er euch geben, Joh. 16, 23. Und wenn man solchen Wunsch und Gruß im Glauben zu allem Dank, mit Freuden aufnimmt, so können sie uns nicht fehlen; denn eines Christen Mund redet nicht vergebene tote und unkräftige Worte, sondern Gott wirket dadurch. Aber wie ihrer viel jetzt einander grüßen, und was sie einander wünschen, das erfähret man oft in betrübten Exempeln. Derwegen wie wir uns selbst alles Gute gönnen, und von Gott bitten, also sollen wir unsern Nächsten auch lieben, und sollen ihm alle Wohlfahrt gönnen, und daß sie ihm widerfahre, wünschen und Gott darum bitten. Dies soll auch erinnern das 33. Kapitel.

\section{Moses Baltesegen und Weissagung}

a Amt) er war nicht König, hatte auch nichts davon, und hielt doch das Volk zusammen, also, daß es ein Haupt hatte, die einen König, und nicht zerstreuet in der Irre ginge.

b Recht) das ist, wie 2 Mos. 28, 30. stehet, daß Heiligthum auf dem Brustlappen. Will also sagen: dein priesterlich Amt sei glücklich für Gott und den Menschen, mit Beten und Lehren.

c Liebliche) das ist, der Tempel und Jerusalem und Königreich war in Benjamin.

d Edle Früchte) das ist vom Königreich Israel gesagt, welches hoch gesegnet ward, mit allem, das Himmel, Sonne, Mond, Erde, Berg, Thal, Wasser und alles zeitliche Gut trug und gab, dazu auch Propheten und heilige Agenten hatte.
\quellseite{330}
• Auszugs) dies hat Debora und Barak gethan, Richt. 5, 12.

* Den Segen Gad hat der König Jehu aufgerichtet, 2 Kön. 10, 27., da er Baal vertilget und das Volk weiter zurechtbrachte, und schlug zween Könige todt, dazu auch Jesebel. Ihr Lehrer ist Elia, der Prophet, der in Himmel genommen und verborgen; denn er war ein Bürger aus Gilead, im Stamme Gad.

§ Gottes Wohnung ist im Volk Israel von Anfang, unter daß sie bleiben ewiglich; als die Hüttlein unter den Armen oder Flügeln der Cherubim.

\section{Summarium}

Ist eine letzte Predigt, in welcher Mosis sein Volk, die zwölf Stämme einen nach dem andern, für seinem Ende segnet, Glück und Heil wünscht, wie im 1 Buch Mos. 49. der Patriarch Jakob auch gethan hat. Aber der Segen hie gehet fürnehmlich auf das zeitliche Priestertum und Königreich, und den zeitlichen Segen, so bei den Juden bis auf Christum, nach Gelegenheit der Zeit, bleiben sollte; wie solches die Historia nach einander fein anzeiget.

\section{Votum}

Der Sohn GOttes, unser Hoherpriester, wolle uns auch segnen, geistlich und leiblich, und geben, was wir sammt den Unsern zeitlich und ewig bedürfen, und wolle solches thun um der Ehre seines heiligen Namens willen, Amen.

\chapter{Das 34. Kapitel}

\section{Vorrede}

Wie andere vor uns gelebt haben und gestorben sind, also die wir jetzund leben, werden auch sterben müssen. Wir leben aber und wissen nicht wie lange, also sterben wir, und wissen nicht wann. Derwegen so sollen wir uns auf ein selig Todtsündlein bereiten, und alle Tage in christlicher Buße leben, auch GOtt bitten, daß er in unserer letzten Noth uns auch beistehe, tröste und stärke wider alle Anfechtungen, und daß er uns erhalten wolle in wahrem Glauben, mit Leib und Seele, zum ewigen Leben; daß uns also der zeitliche Tod eine Thür und Pforte sei, dadurch wir mögen eingehen, aus des HErrn Jesu Christi Gnadenreich dieser Welt, in sein Reich der Glorie und Herrlichkeit, zu ewiger Freude und Seligkeit, Amen. Darauf höret ferner an das 34. und letzte Kapitel des fünften Buchs Mose.

\section{Moses Tod, und was darauf erfolget}

\quellseite{331}
\section{Summarium}

Ist eine Historia, wie Moses in der Moabiter Land gestorben, v. 5., und vom HErrn begraben sei, v. 6., und sein Grab hat niemand können erfahren bis auf diesen heutigen Tag, v. 8. Da sollen wir erstlich lernen, wider der Papisten Irrthum: wo Gott wollte, daß man die verstorbenen Heiligen anbeten, ihren Gräbern, Gebeinen oder Asche Ehre thun sollte, würde er Moses, als seines treuesten und trefflichsten Dieners Grab und Gebeine nicht selbst verborgen haben.

Zum andern ist hie zu merken das große Wunderwerk, daß Moses hundert und zwanzig Jahr alt worden, und doch weder am Gesicht, noch sonst am Leibe etwas habe abgenommen. Zu einem Zeugniß, auf daß wir lernen, wer Gott und sein Wort für Augen hat, daß Gott denselbigen mit Leib und Leben erhalten, für Krankheit behüten, und sonst in andere Wege mehr seinen Segen reichlich widerfahren lassen.

\section{Votum}

Also hat euer Liebe das fünfte und letzte Buch Mose gehört, in welchem sonderlich zu merken ist, wie treulich er uns alle vermahnet, daß wir fromm sein, und GOttes Gebot sollen mit Fleiß halten, so werde GOtt auch mit seinem Segen bei uns sein. Wer es aber nicht thut, sondern GOtt und sein Wort verachtet, und sündlich lebet, der soll darum hie und dort gestraft werden. GOtt verleihe seine Gnade, daß wir durch solche Lehre mit Ernst uns bessern und frömmer werden, Amen.

\part{Das Buch Josua}

\chapter{Das 1. Kapitel}

\section{Vorrede}

Euer Liebe haben gehört in den fünf Büchern Mose, daß Gott erstlich im Anfang Himmel und Erde dem Menschen zu Nutz erschaffen hat, und daß der Mensch hernach durch den Teufel im Paradies verführet, in die Sünde und Ungehorsam gefallen sei. Solche sündliche Natur hat sich für und für hernach in der Welt gereget, je länger je mehr, bis Gott mit der Sündflut alle Menschen um der Sünde willen ersäufet hat, aufgenommen den frommen Noah, sammt allem was in der Arch oder Kasten bei ihm war. Von denselbigen acht Menschen ist hernach die neue Welt kommen; dieselbige hat eben der vorigen Welt Exempel nach gelebet, ist gefallen von Gott, seinem Wort und Willen, auf Abgötterei und in allerlei Sünde; denn Fleisch und Blut kann nicht anders, wo der heilige Geist nicht sonderlich hält und behütet. Da erwählet ihm Gott Abraham den Propheten und verhieß ihm, er wollte sein Gott sein, und seines Samens nach ihm; und wollte ihm nicht allein das Land Canaan zu eigen
\quellseite{332}
gehen, sondern auch den gebenedeiten Samen Christum den HErrn, aus seinem Fleisch lassen geboren werden, durch welchen alle Völker auf Erden sollten gesegnet, das ist, von Gottes Zorn und dem Fluch, so sie durch die Sünde verdienet hätten, ledig werden und zu Gnaden kommen, wo sie anders an ihn glaubten. Auf solche Zusage hänget Gott dem Abraham und seinem Samen ein sonder Zeichen, die Beschneidung an, bei welchem sie erkennen sollten, daß sie Gottes Volk wären, und etwas Sonderes für allen andern Völkern auf Erden.

Als nun des frommen Abrahams Nachkommen Hungers halben aus Canaan entweichen mußten in Egypten, und da eine Zeit lang bauseten, aber hernach von den Egyptiern unbillig und hart bedrängt und beschweret worden, da nahm sich Gott allererst ihrer, als seines Volks an, plaget die Egypter, wie das andere Buch Mose zeuget, den König mit aller seiner Macht ersäufet er im rothen Meer, und erlediget also sein Volk.

Von dem an unterscheidet er sie erst recht von allen andern Völkern, nicht allein durch die Beschneidung, sondern auch durchs Gesetz und Eröffnung seines göttlichen Willens, was sie thun und lassen, wie sie ihm auch äußerlich mit Opfern und Anderm dienen sollten. Daß also dies Volk sein eigen Königreich und Priesterthum, oder Kirchenordnung hätte, von Gott selbst gestiftet und geordnet; daher sind in den fünf Büchern Mose manche schöne Predigt, auch seine Exempel, wie man Gott und seinem Wort folgen soll, und wie es Gott entweder belohnet, oder gestraft habe, wenn man nach oder wider sein Wort gethan hat.

Auf diese fünf Bücher folget nun die ganze Historia, wie die Juden, Abrahams Nachkommen, durch Gottes Hülfe das Land eingenommen, und sich darnach im Lande unter den Richtern und Königen gehalten haben, bis sie endlich um ihrer Sünde willen vom Könige aus Assyrien und vom Könige zu Babel, gar aus dem Lande sind vertrieben und weggeführt worden.

Solche Historien, ob sie schon nun längst vergangen sind und zum alten Testament gehören, sind sie doch in mancherlei Weg nütze, daß man sie unter den Christen lesen und lernen soll; denn da findet man nicht allein viel schöne Lieder und Predigten, die sich bisweilen zutragen vom rechten Gottesdienst und allerlei guten Werken, sondern auch die feinsten Exempel, die wohl und fleißig zu merken sind, aus welchem wir zweierlei lernen sollen. Erstlich Gott fürchten, als der die Sünde nicht hat ungestraft lassen hingehen, und darnach auch Gott trauen, als der die Seinen nicht hat stecken lassen, sondern allwege zu gelegener Zeit seine Hülfe mittheilet. Solches ist an ihm selbst nützlich und gut, sintemal ein Christ sein ganzes Leben über stets in der Uebung des Glaubens und Furcht Gottes leben soll. Zum andern ists auch sonderlich tröstlich, daß man in solchen Historien siehet, wie die feinsten, frömmsten und heiligsten Leute bisweilen gestrauchelt, und Unrecht thun, und darnach mancherlei Weise durch Kreuz und Leiden sind geläutert, und doch endlich aus Gottes Gnade solche Sünden ihnen vergeben sind worden.

Zum letzten soll uns je nichts lustigers und lieblichers sein, denn daß wir in solchen Historien lernen mögen, nicht allein uns selbst kennen, und bessern, sondern auch gleich eine Erquickung haben, wenn wir der jetzigen bösen Welt satt und überdrüssig sind, daß wir in die vorige Welt unter die frommen heiligen Leute sehen, und da einen Trost schöpfen, und an ihrem Exempel lernen geduldig sein, und der Welt Bosheit überwinden mit Geduld und Hoffnung eines bessern Lebens, nach diesem sündhaftigen elenden Leben, Amen.

\section{Josua wird in seinem Beruf gestärket}

a Wer nach Gottes Wort sich richtet, der handelt weislich und glücklich; wer nach seinem Kopf führet, der handelt unweislich und vergeblich.
\quellseite{333}
\section{Summarium}

Dies ist eine Historia, wie der HErr den Krieger Josua beruft, daß er soll nach Mose-Abgang des Volks Israel Hauptmann sein; tröstet und lehret ihn, wie er sich halten soll; denn es gehöret ein großer Muth dazu, wer ein Regent sein und den Sachen allenthalben recht thun soll. Darum befehlet Gott mit Fleiß, daß Josua sich fest an Gottes Wort und Gebot halten, und demselbigen folgen soll. Laß das Gesetz des HErrn (spricht er V. 8.) nicht von deinem Munde kommen, alsdenn wird dir gelingen in allem das du thust, und wirst weislich handeln können. Solches soll insonderheit die Obrigkeit zu Herzen nehmen, so sie anders glücklich und recht regieren will; denn wo man Gottes Wort nicht folgen, sondern dasselbige verfolgen will, da wird weder Glück noch Heil sein können.

\section{Votum}

Der treue, barmherzige GOtt erhalte uns sein heiliges Wort, und regiere durch den heiligen Geist Obrigkeit und Unterthanen, daß sie seinem Wort gehorsamlich folgen, ihm zu großem Lob, Ehr und Preis, und den Menschen zu gemeiner heilsamen Wohlfahrt, Amen.

\chapter{Das 2. Kapitel}

\section{Vorrede}

Gott will ein gnädiger Vater sein allen Armen und Gläubigen, und soll ihm kein Mensch zu geringe noch zu verachtet sein. Wiederum will Gott ein zorniger Richter sein über die Gottlosen und Ungläubigen, und soll ihm keiner zu groß, zu mächtig, zu klug und zu reich sein; beides bezeugen die geistlichen und weltlichen Historien alter Geschichte und die tägliche Erfahrung. Derwegen wollen wir uns demütigen und in wahrem Glauben und Gehorsam an Gott halten, so genesen wir in Gefahr, in Noth und Tod. Wer aber Gott und seinem Wort nicht trauet noch folget, der wird auch den Lohn haben, aber mit Schand und Schaden. Dies soll auch erinnern das 2. Kapitel.

\section{Die Kundschafter werden bei der Rahab errettet}
\quellseite{334}
\section{Summarium}

Sie ist ein schön Exempel eines feinen Glaubens, welchen die Epistel an die Ebräer auch rühmet: daß Rahab für gewiß achtet, daß die Kinder Israel Gottes Volk sind und sollen wider die Heiden siegen, und ihr Land einnehmen, weil Gott in Egypten und am rothen Meer so große Wunderzeichen mit ihnen gethan habe. Darum verbirgt sie die Kundschafter, V. 6., und wird durch sie auch erhalten, da die ganze Stadt dagegen keinen Glauben an Gottes Wort

19 *
\quellseite{335}
hatte, und derhalben alles darin mußte erwürget werden.

\section{Votum}

Gott, der Vater unsers Herrn Jesu Christi, stärke unsern schwachen Glauben, und helfe, daß wir durch den Glauben überwinden den Teufel, die Welt, auch unser eigen Fleisch und Blut, und daß wir alle an Leib und Seel erhalten werden, zeitlich und ewig, Amen.

\chapter{Das 3. Kapitel}

\section{Vorrede}

Gleich wie die levitische Bundeslade zur Zeit Josua mitten im Jordan gestanden, den Strom desselben aufgehalten, und den Israeliten, nach ihrer schweren und gefährlichen vierzigjährigen Pilgrimschaft einen sichern und freien Paß und Gasse ins gelobte Land verschaffet: also da unser Herr Christus hat sollen getauft werden, und in sein Predigtamt treten, kömmt er als der lebendige Gnadenthron, und stellet sich selbst mitten in den Jordan, und will unser Weg und Steg zum ewigen Vaterland sein, und schützen und aufhalten alle Schüsse und Güsse der Flut göttliches Zorns, und des Wassers der Trübsalen abzuwehren, daß uns die Ströme des rothen Drachens und Satans nicht verschw emmen, noch die Meereswellen der Widerwärtigkeit allzu hart über unsere Seele gehen, und uns gar tilgen sollen, so wir nur auf ihn sehen, und uns an ihn in wahrem Glauben halten. Das sollen wir uns auch erinnern bei dem 3. Kapitel.

\section{Israel geht trockenes Fußes durch den Jordan}

a Diese Geschichte bedeutet, daß uns Christus durch die Apostel vorgetragen im Evangelio, leitet ins Himmelreich durch den trockenen Jordan, der doch diesmal am vollesten war. Das ist, daß Gesetz, das uns mit Werken ängstet und treibet, höret auf durch Evangelium, daß unsere Gewissen frei, fröhlich und sicher für ihm sind und allein Christo im Glauben folgen.
\quellseite{336}
\section{Summarium}

Ist eine Historia des herrlichen Wunderwerks, daß Gott den Jordan still stehen und sein Volk mit trockenem Fuß ließ hindurch gehen V. 15., auf daß sie an solchem Wunderwerk lernten, Gott wollte mit ihnen sein und sie derhalben desto getroster hinan gingen, und das Land einnähmen.

Die geistliche Deutung ist diese, daß Christus gleich wie die Archa uns wird fürgetragen im Evangelio. Demselben gehen wir nach im Glauben durch den Jordan ins gelobte Land, das ist, durch dies armselige Leben ins ewige Leben, da sonst, wo Christus nicht wäre, der Jordan, das ist, die Welt sammt dem Teufel, Sünd und allem Unglück uns würde ersäufen.

\section{Votum}

Unserm HErrn JEsu Christo sei Lob, Preis und Dank für alle seine Wohlthaten; der wolle auch heute fürgehen allen, die durch den Jordan des Trübsals wandern müssen, und wolle uns allesamt in diesem Elend und im zeitlichen Tode zum ewigen Leben erhalten, Amen.

\chapter{Das 4. Kapitel}

\section{Vorrede}

Die christliche Kinderzucht ist leider jetziger Zeit gar ungemeine. Weh aber über alle, weh allen Eltern, die ihre Kinder nicht in der Zucht und Vermahnung zu dem HErrn auferziehen. Denn kein Löwe ist so grausam, kein Bär so grimmig, kein Wolf so wüthend, keine Schlange so giftig, keine Gluckhenne so hungrig, sie versorgen ihre Jungen nach ihrer Art, mit aller Nothdurft. Sollte denn der Mensch, der doch die alleredelste, herrlichste und die vernünftigste Creatur nach den Engeln sein soll, sein Fleisch und Blut mit nöthiger Zucht verwahrlosen und verlassen? Derwegen die Eltern, und wem mehr die Kinderzucht befohlen ist, sollen also die Jungen erziehen, daß sie Gott zu Ehre, ihrem Nächsten zu Nutz, und ihnen selber zur Wohlfahrt leben, auf daß sie auch selber an den Kindern Ehre und Freude erleben, und nach diesem betrübten Leben mit ihnen ewige, freudenreiche Gemeinschaft, sammt der hochgelobten Dreifaltigkeit, allen auserwählten Engeln und Christgläubigen haben mögen; und sollen wissen, so manche Mühe sie mit der Kinderzucht haben, so manchen Gottesdienst begehen sie damit, wo ferne sie aus dem Glauben solche Werke thun. Dies lehret auch das 4. Kapitel.

\section{Denckzeichen des wunderbaren Durchgangs durch den Jordan}

a Also das ist, die Priester, so die Lade trugen, stunden, wie sie geordnet waren, und thaten nichts anders.
\quellseite{337}
\section{Summarium}

Ist ein Befehl, wie man zwölf Steine zum Gedächtniß des großen Wunderwerks soll aufrichten, v. 3.; denn Gott will sonderlich, daß man seine Wohlthat nicht vergessen, sondern dieselbigen dem jungen Volk fleißig einpflanzen, und es dabei erhalten soll, auf daß man in andern fürfallenden Nöthen auf seine Hülfe bauen und ihm trauen lerne. Darum sollen die Eltern hie sich ihres Amts erinnern, und mit sonderlichem Fleiß die Kinder zu Gottes Wort und Furcht halten.

\section{Votum}

Helfe der treue GOtt, daß die Eltern in der Kinderzucht fleißig und treu, und daß die Kinder gehorsam seyn und sich gerne ziehen lassen, auf daß die Eltern und Kinder GOtt zeitlich und ewig mit einander loben und preisen, Amen.

\chapter{Das 5. Kapitel}

\section{Vorrede}

Gott verläßt niemand, man verlasse ihn nur nicht; denn je größer Noth, je näher GOtt ist, wenn man ihn in wahrem Glauben und Vertrauen auf seine Verheißungen anruft um Trost und Hülfe in der Noth. Derwegen wollen wir auch mit Freuden fleißig thun, was einem jeden in seinem Beruf gebühret, und wollen beten, daß Gott helfen, schützen und erretten wolle. So werden wir auch Gottes Segen haben, und von ihm behütet und erhalten werden. Denn also tröstet Gottes Wort im fünften Buch Mose im 4. Kapitel V. 31.: Der HErr dein Gott ist ein barmherziger GOtt, er wird dich nicht verlassen, noch verderben, wird auch nicht vergessen des Bundes, den er deinen Vätern geschworen hat. Davon höret auch an das 5. Kapitel.

\section{Beschneidung und Passah im gelobten Lande verrichtet}

a (Steinerne Messer) scharf, als am Stein gewetzt, Ps. 89, 44. Avertisti petram gladii ejus. Die Schärfe seines Schwerts ist stumpf, daß es nicht schneidet.
\quellseite{338}
b) Sangen) am Feuer gedörrte oder versengte Aehren, tostatas spicas.

\section{Summarium}

Dies ist erstlich eine Anzeige, was Gott mit dem Wunderwerk verrichtet habe, daß der Jordan ist ausgetrocknet, nämlich, daß sein Volk dadurch getröstet, und den Feinden das Herz genommen würde. Eben um solcher Ursache willen erscheinet ein Engel Josua insonderheit, V. 13., auf daß er ja allenthalben sehe und merkte, daß Gott wolle bei ihm seyn, und er deßhalben desto freudiger und getroster seinem Beruf soll nachkommen.

\section{Votum}

Gott lasse ihm alle christliche Obrigkeit befohlen sein, verleihe ihnen Hülfe und den Unterthanen Schutz, und segne auch die Werke unserer Hände zu seinem Lob, und unsers Nächsten Wohlfahrt, um Jesu Christi, unsers Herrn, willen, Amen.

\chapter{Das 6. Kapitel}

\section{Vorrede}

Gott widerstrebet den Gottlosen und hilft den Gläubigen. Denn obschon die Kinder Gottes auch ihre Mühseligkeit und Kreuz auf Erden haben, und oft mehr und größeres, denn die Sünden-Knechte; dennoch folgt ihnen ein gewünschter Ausgang, die Gottlosen aber gehen endlich gar zu Grunde und Boden. Davon sagt David im 75. Psalm V. 9. also: Der HErr hat einen Becher in der Hand und mit starkem Wein voll eingeschenkt, und schenket aus demselben; aber die Gottlosen müssen alle trinken und die Hefen aussaufen. Derwegen wer Gott mit seinem Segen auf seiner Seite haben will, der thue Buße, bete und lebe nach seinem Wort und Willen unsträflich. Dies erinnert auch das 6. Kapitel.

\section{Eroberung und Zerstörung der Stadt Jericho}
% Scan page 339: No non-biblical apparatus, gloss, or structural title/synopsis retained in OCR blocks.
\quellseite{340}
\section{Summarium}

Ist eine Historia, daß die Stadt Jericho jenseit des Jordans am ersten genommen, und mit allem verbrannt sei worden, V. 2. Hie soll man erstlich lernen, wie Gott solch greulich Urtheil über diese Stadt um der Sünden willen hat gehen lassen, daß Alt und Jung, Mann und Weib erstochen ward, V. 21., auf daß wir uns für Sünden hüten, und dergleichen Urtheil nicht auch verdienen. Zum andern, daß die Mauern von ihnen selbst umfallen, und das Volk Israel doch nicht mehr thut, denn daß es um die Stadt gehet mit der Arche, V. 30., solches ist ein Exempel durch welches Gott sie hat wollen lehren, wie er für sie streiten und ihnen beistehen wolle. Wie denn aller Sieg von Gott kommt, und alles an Gottes Segen liegt.

\section{Votum}

Der allmächtige, barmherzige GOtt streite heute auch für seine Gemeine, und segne das Kirchen-, weltliche und Hausregiment, helfe doch, daß ein jeder in seinem Beruf christlich und ehrbar lebe, Amen.

\chapter{Das 7. Capitel}

\section{Vorrede}

Gottes Wort lehret es, die Historien und aller Welt Geschichte bezeugen, und die tägliche Erfahrung bestätiget, daß Gott die Sünder endlich ernstlich und heftig strafet. Aber, ach der Welt Blindheit und Unsinnigkeit, die sich dies nicht überreden und bessern will lassen. Derwegen wollen wir unser und der Unsern verschonen, wollen die Sünde hassen und lassen, und Tugend lieben und üben, daß solches sei Gott ein Wohlgefallen, den heiligen Engeln eine Freude, unserm Nächsten Nutz, uns zur Wohlfahrt diene, und daß den Teufeln eine Pein und Marter sei ewiglich, Amen. Darauf höret ferner an das 7. Capitel.

\section{Achan wegen Diebstahls bestraft}
\quellseite{341}
\section{Summarium}

Hie ist ein Exempel, wie Gott bisweilen an einem ganzen Haufen strafet, was etliche insonderheit haben unrecht gethan, V. 35., auf daß die weltliche Obrigkeit lerne ihres Amts fleißig warten, und das Uebel mit höchstem Ernst strafen, denn thut mans nicht und läßt die Strafe unterwegen, so muß hernach der ganze Haufe darum herhalten, wie man siehet.

a Spangen, wie eine Zunge gestaltet.

\section{Summarium}

Ist ein sonderlich Exempel, in welchem wir sehen, wie hart und mit großem Ernst Gott die Sünde strafet, und endlich an den Tag will bringen, man berge und decke sie so lang man immer wolle. Denn weil Gott verboten hatte, daß niemand vom Raub etwas sollte nehmen, und Achan dennoch Gold und Silber genommen und vergraben hatte, mußte es an Tag kommen, und er mit Weib und Kindern, mit Vieh und aller seiner Habe verbrennet werden, V. 26. Solche heftige Strafe soll uns allen eine Warnung sein, daß wir uns für Sünden hüten und wider Gottes Wort nichts thun noch fürnehmen sollen, auf daß uns dergleichen nicht auch widerfahre.

\section{Votum}

Regiere der treue GOtt die Obrigkeit, daß sie die Sünden ernstlich strafe, und uns regiere GOtt, daß wir uns für Sünden fleißig hüten, auf daß wir allem Unglück entfliehen, und bestehen mögen am jüngsten Tage für des Menschen Sohn, Amen.

\chapter{Das 8. Kapitel}

\section{Vorrede}

Ein jeder Beruf und Handthierung hat seine Mittel, dadurch sie sollen getrieben werden. Diese soll ein jeder fleißig üben und gebrauchen, und den Segen dazu allein von Gott bitten und hoffen, nach dem gemeinen Sprüchwort: Bete und arbeite. Denn wer sein Vertrauen auf die Mittel allein setzt, als auf seinen Fleiß und Arbeit, der hat wohl viel Mühe, aber seinen Nutzen hat er davon, wie Gottes Wort und die Erfahrung bezeuget. Dagegen der Segen des HErrn
\quellseite{342}
macht reich ohne Mühe, sagt Salomon Sprüche 10, 22. Und St. Paulus spricht von denen, die mit Gott und Ehren ihres Berufs abwarten: Eure Arbeit ist nicht vergeblich in dem HErrn, 1 Cor. 15, 58. Davon unterrichtet auch das 8. Kapitel.

\section{Belagerung und Eroberung der Stadt Ai}
\quellseite{343}
\section{Summarium}

Ist eine Historia, wie Ai und Bethel von Josua gewonnen sind, V. 2. Daß aber der HErr befiehlt, daß Josua mit List sie aus der Stadt locken, und darnach von hinten zufallen sollen, ist eine Anzeigung, daß, ob wohl der Sieg von Gott kömmt, und alles am Segen Gottes liegt, dennoch wir allerlei Mittel, so von Gott verordnet sind, dahin brauchen sollen, und gleichwohl unsere Hoffnung auf Gott und nicht auf die Mittel setzen. Also soll ein Kriegsmann seinen Vortheil nicht begeben, Büchsen und Harnisch brauchen, aber auf solche Mittel sich nicht verlassen, sondern seine Zuversicht allein zu Gott haben.

\section{Votum}

Gott regiere uns alle, daß wir ihm allein die Ehre geben, in allem unsern Vornehmen und Thun, und segne und behüte einen jeden in seinem christlichen Beruf und Stande, daß es zu seinem Lob und gemeinem Nutz zum Besten diene und gelange, Amen.

\chapter{Das 9. Kapitel}

\section{Vorrede}

Die Alten nennen Eidschwur eine Arznei; denn wie man die Arznei nicht oft und ohne sonderer Noth nicht gebrauchen soll, also sollen wir ohne große erhebliche Noth des Eides nicht gebrauchen. Wenn er aber zur Stärkung der Wahrheit, Gott zu Ehren, uns und dem Nächsten zur Wohlfahrt, auch zur Erhaltung Friedens und Einigkeit und Rettung der Unschuld dienet, und ist nicht wider Gottes Wort, so kann und soll man die Eide gebrauchen und thun, und sonst soll man sich für unnützen und unnöthigen Schwüren hüten. Denn schwören ohne Noth (spricht man) ist der Seelen Tod. Derwegen warnet uns Sirach ernstlich dafür und sagt im 23. Kapitel V. 9.: Gewöhne deinen Mund nicht zum Schwören, denn gleich als ein Knecht, der oft gepräget wird, nicht ohne Striemen ist, also kann einer der oft schwöret, nicht ohne Sünden sein. Dies lehret auch das 9. Kapitel.

\section{Die Gibeoniter kommen mit List in den Bund}
% Scan page 344: No non-biblical apparatus, gloss, or structural title/synopsis retained in OCR blocks.
\quellseite{345}
\section{Summarium}

Hie ist eine sonderliche Historia, in welcher wir sehen, wie man den geschworenen Eid soll halten, obs gleich mit unserm Nachtheil geschehen muß, und ob wir übervortheilt sind. Das also keine Ausflucht ist, was man geschworen und zugesagt hat, das soll man halten, ohne so es wider Gott und sein Wort ist.

Darneben siehet man auch hie, wie die, so durch List mit ihrem Vortheil den Eid suchen, ungestraft nicht bleiben. Denn die Gibeoniter sammt allen ihren Nachkommen müssen diese Strafe tragen, daß sie Holzhauer und Wasserträger unter den Kindern Israel müssen bleiben, V. 27. Hie mag man bedenken, was für eine Strafe werde erfolgen, wo man einen falschen Eid schwört, oder den geschworenen Eid um Geldes oder Guts willen bricht. Denn weil der Name Gottes zum Eid gebraucht, und der Eid dadurch für kräftig gehalten wird, will Gott seinen Namen von niemand unehren, noch die Leute dadurch betrügen lassen.

\section{Votum}

Der barmherzige GOtt regiere uns, daß wir im Thun und Lassen nach seinem Wort uns halten, und also nach dem Eide, den wir ihm in der Taufe gethan, christlich leben, damit es zu seinen Ehren und zu unserer Wohlfahrt gelange, Amen.

\chapter{Das 10. Kapitel}

\section{Vorrede}

Die Gott kennen, die ihn fürchten, zeitlich und ewiglich wohl zu thun verheißt, also dräuet er wiederum zeitlich und ewig zu strafen, die so ihn verachten. Von Gottesfürchtigen sagt David also Psalm 33, 18.: Siehe des HErrn Auge stehet auf die, so ihn fürchten, die auf seine Güte hoffen, daß er ihre Seele errette vom Tode, und ernähre sie in der Theurung. Und abermals spricht er Psalm 147, 11.: Der HErr hat Gefallen an denen, die ihn fürchten, die auf seine Güte hoffen. Von Gottesverächtern sagt Gott Mal. 3, 5. also: Ich will zu euch kommen und euch strafen, und will ein schneller Zeuge sein wider die, so mich nicht fürchten. Und im Buch der Weisheit Cap. 1, 8. stehet diese Dräuung: Der kann nicht verborgen bleiben, der das Unrecht redet, und das Recht, so ihn strafen soll, wird sein nicht fehlen; denn des Gottlosen Anschläge müssen vor Gericht, und seine Reden müssen für den HErrn kommen, daß seine Untugend gestrafet werde. Davon höret auch an das 10. Kapitel.

\section{Wunderbarer Sieg Josuas wider die Amoriter}
\quellseite{346}
\section{Summarium}

Ist eine sonderliche Historia und ein herrlich großes Wunderwerk, daß Gott der Stimme Josua, des einigen Mannes gehorchet und läßt Sonne und Mond einen ganzen Tag am Himmel stehen, daß sie nicht fortgehen, V. 13.; daß also das Volk desto mehr Zeit hätte, sich an seinen Feinden zu rächen. Solches Wunderwerk dienet dahin, daß man lerne, wie Gott über den Seinen, die ihn fürchten, halten wolle, und thun was sie begehren, wie der 145. Psalm V. 19. saget. Neben dem, daß auch solch Wunderwerk dem Volk einen Muth, und den Feinden einen Schrecken macht, weil sie sahen, daß Gott mit Josua war.
\quellseite{347}
\section{Summarium}

Ist eine Historia, wie Josua etliche Könige und Städte geschlagen und vertilget habe. Da sehet man, wo Gott mit seinem Segen ist, daß alles schleunig von der Hand gehet. Weil aber Gott seinen Segen allein den Frommen hat zugesagt, sollen wir uns seines Worts und Willens mit Ernst halten, so wird es uns auch glücklich gehen. Wie denn hie allenthalben von Josua gerühmet wird, er hab alles verbannet und ausgerottet, wie der HErr ihm geboten hat. Darum ist auch Glück bei ihm.

\section{Votum}

GOtt, unser lieber, treuer Vater, regiere und leite uns, daß wir seinen heiligen Geboten folgen, und segne und behüte uns an Leib und Seele, in allen Ständen, zu zeitlicher und ewiger Wohlfahrt, um JEsu Christi willen, Amen.

\chapter{Das 11. Capitel}

\section{Vorrede}

Gott pfleget den Verzug der Strafe mit Schwerheit der Strafe einzubringen, das ist, je länger Gott verzucht die Gottlosen zu strafen, je greulicher er sie darnach, wenn sie ihr Maß gefüllet, und er seine Zeit ersehen, zu strafen pfleget. Dieses sollte uns genug Ursache sein, christlich zu leben und für Sünden sich zu hüten. Aber die Welt will gestraft sein, darum ringet sie nach Schlägen und mißbraucht Gottes Langmuth und Geduld erschrecklich zu greulichen Sünden. Derwegen wird sie gestraft, und wird mehr und größere Strafe leiden müssen. Gott erbarme sich unser gnädiglich, Amen. Darauf höret auch davon an das 11. Kapitel.

\section{Ausrottung der Canaaniter}
\quellseite{348}
\section{Summarium}

Das ist ein fein Exempel, in welchem man siehet, wie Gott den Gottlosen eine Zeit lang zusieht, aber wenn sie sich nicht wollen bessern, und das Stündlein kommt, müssen sie ihnen selbst Ursache geben zu ihrem Verderben. Wie Josua hie sagt: Es geschah aber vom HErrn, daß ihr Herz verstocket wäre, daß sie sich unterstünden mit Israel zu streiten, V. 26., das ist, ihre Sünden hattens verdient und waren reif zur Strafe, darum mußten sie selbst Ursache dazu geben. Also stehet man, daß kein größerer Vortheil im Kriege wider die Feinde ist, denn so sie durch ihre Sünde Gott wider sich haben; wie heutiges Tages des Türken höchste Macht und großer Vortheil ist, daß die Obrigkeit zum Theil Gott und sein Wort verfolget, und sonst auch kräftlich ist, und wir insgemein so ruchlos sind in Geiz, Unzucht, Büberei, Haß, Neid und anderen Sünden, und sehr wenig sich an Gottes Wort bessern, da muß denn wohl Unglück folgen.

\section{Votum}

Der treue GOtt erleuchte unsere Augen und Herzen, daß wir uns an anderer Leute Unglück spiegeln, von Sünden abstehen, gerecht und heilig leben, auf daß wir nicht auch zeitlich und ewig gestraft werden, Amen.
\quellseite{349}
\chapter{Das 12. Kapitel}

\section{Besiegte Könige}

\section{Summarium}

Josua wiederholet hie in einer Summa alle Könige und Städte, so Mose und Josua geschlagen und eingenommen haben, nämlich, ein und dreißig Könige, V. 24. An solcher Summa mag man lernen, wie die Sünde und das gottlose Wesen einen Riß in der Welt und großen Regimenten machet, wie es Krieg und alles Unglück erreget. Darum sollen wir in Gottesfurcht wandeln, und seines Willens und Wortes uns fleißigen, so will er uns für solchem Unglück behüten.

\chapter{Das 13. Kapitel}

\section{Austheilung des Landes wird vorgenommen}
% Scan page 350: No non-biblical apparatus, gloss, or structural title/synopsis retained in OCR blocks.
\quellseite{351}
\chapter{Das 14. Capitel}

\section{Von Calebs Erbtheil}

a Kirjath Arba) das ist die Arbastadt, von Arba einem Riesen also genannt, wie wir Deutschen sagen, Karlstadt, Arnstadt etc.

\section{Summarium}

Ist ein fein Exempel wie Caleb, Josua Gesell, der Kundschafter einer, darum, daß er dem Gott Israel treulich gefolget hatte, von Gott dermaßen ist behütet worden, daß er sich rühmet, er sei fünf und achtzig Jahr alt, und sei dennoch allerdings so stark, als da er vierzig alt sei gewesen, V. 11. Wer aber ein gottlos Leben führet, und nicht wandelt in der Furcht GOttes, da träuet die Schrift, daß dieselbigen ihre Tage nicht auf die Hälfte bringen werden; wie wir zu unserer Zeit Exempel genug, und schier ein jeder in seinem eigenen Hause und an ihm selbst sehen.

\section{Votum}

GOtt, der Vater aller Barmherzigkeit, behüte uns für gottlosem Thun und Wandel, und helfe, daß unser Leben ihm ein wohlgefällig Leben sei, und daß wir Segen von ihm haben an Leib und Seele, in diesem und künftigem Leben, Amen.

\chapter{Das 15. Capitel}

\section{Grenze und Städte des Stammes Juda}
% Scan page 352: No non-biblical apparatus, gloss, or structural title/synopsis retained in OCR blocks.
\quellseite{353}
a Denn David hat sie hernach mit der Zeit vertrieben, 2 Sam. 5, 7.

\chapter{Das 16. Kapitel}

\section{Grenzstätte der Kinder Ephraims}

\chapter{Das 17. Kapitel}

\section{Vom Erbtheil der andern Hälfte des Stammes Manasse}
\quellseite{354}
\chapter{Das 18. Capitel}

\section{Grenzstätte des Stammes Benjamin}
\quellseite{355}
\chapter{Das 19. Kapitel}

\section{Der übrigen sechs Stämme und Josua's Erbtheil}
% Scan page 356: No non-biblical apparatus, gloss, or structural title/synopsis retained in OCR blocks.
\quellseite{357}
\chapter{Das 20. Kapitel}

\section{Verordnung der Freistädte}

\section{Summarium}

Ist eine Ordnung mit den Freistädten, V. 2., auf daß wer unwissend oder wider seinen Willen hätte einen Todtschlag gethan, nicht übereilet würde, wie zuvor im Mose weiter davon gesagt ist.

\chapter{Das 21. Kapitel}

\section{Versorgung der Leviten}
% Scan page 358: No non-biblical apparatus, gloss, or structural title/synopsis retained in OCR blocks.
\quellseite{359}
\chapter{Das 22. Kapitel}

\section{Vorrede}

Die heilige Biblia muß neben und bei dem Schwerte sein, soll ein Regiment glücklich und wohl von statten gehen; denn die Obrigkeit soll nicht allein Hüter der andern, sondern auch der ersten Tafel sein, und soll nicht allein ihrer lieben getreuen Unterthanen Leib und Gut beschirmen, sondern auch daran sein, daß die Religion und der rechte Gottesdienst wohl bestellt sei. Denn wie eine Laterne ein Licht bewahret und erhält, daß es nicht liederlich kann ausgeblasen oder ausgelöscht werden, und, wenn auch das Licht in der Laterne nicht brennet, so ist sie zu lauter nichts nütze: also sollen auch christliche Regenten das helle Gnadenlicht Christum und sein göttliches Wort lassen lehren und predigen, daß der lebendige persönliche Gnadenthron in der ganzen Welt herum geführet werde, und sollen darob sein, daß der Fürst der Finsterniß, der Teufel, solch Gnadenlicht, weil es ihm in die Augen scheinet und unleidlich ist, nicht auslösche; ohne dies ist ihre Regierung auch nicht allerdings nütze. Dies lehret auch das 22. Kapitel.

\section[Die dritthalb Stämme richteten im Heimkehren einen Altar]{Die dritthalb Stämme richteten im Heimkehren einen Altar auf}

\section{Summarium}

Ist eine schöne nützliche Historia, in welcher die weltliche Obrigkeit ihres fürnehmsten Amts erinnert wird, daß neben dem sie die Unterthanen schützen und ihnen treulich fürstehen soll, sie auch dem Exempel Josua nach, ihr Volk zu Gottes Wort halten, ihnen dasselbige befehlen und fleißig einbilden soll, V. 5. Denn ob wohl das Predigtamt ein unterschiedliches Amt ist von dem weltlichen Regiment, so soll doch die weltliche Obrigkeit das Wort Gottes mit Fleiß helfen fördern, und die Strafe mit Ernst gehen lassen wider alles so dem Wort entgegen und ärgerlich ist. Wo Gottes Name also von der weltlichen Obrigkeit geehrt und dem Aergerniß gewehret wird, da will Gott die Leute wiederum mit Glück und Segen ehren. Wiederum, wo weltliche Obrigkeit Gottes Ehre nicht achtet, und nimmt sich um die Kirche nicht an, da will Gott des Regiments sich auch nicht annehmen, und muß endlich alles zu Trümmern gehen, und kann weder Gewalt noch Reichthum helfen.
% Scan page 360: No non-biblical apparatus, gloss, or structural title/synopsis retained in OCR blocks.
\quellseite{361}
\section{Summarium}

Sie ist eine schöne Historia, wie die zehn Stämme besorgen, es möchten die andern jenseit des Jordan Israel sich zur Abgötterei geben, daß sie ohne Befehl einen Altar in Gilead gebauet hatten; aber da sie höreten, wie solcher Altar nicht zum Opfer gebauet wäre, sondern zum Zeugniß, daß sie auch zu GOttes Volk gehörten, ließen sie es ihnen gefallen, V. 30. Dies Exempel zeigt an, mit was Fleiß und Ernst weltliche Obrigkeit darauf sehen solle, daß allem Aergerniß gewehret werde.

\section{Votum}

Der treue, barmherzige GOtt lasse ihm unsere liebe Obrigkeit, hohen und niedrigen Standes, befohlen sein, und helfe, daß ihre Regierung auch ferner zur Erhaltung seines Worts diene, und daß allem Aergerniß gewehret werde, zu Lob und Preis seinem Namen, und den Unterthanen zur Wohlfahrt und Seligkeit, Amen.

\chapter{Das 23. Kapitel}

\section{Vorrede}

Wer nicht Gottes offenbarem Wort folget in all seinem Thun und Lassen, der gedenke nicht, daß er glücklich sein und Gottes Segen zu Leibes und der Seelen Nothdurft haben und behalten werde; denn nach dem Unglauben ist Gott keiner Sünde feinder, strafet auch keine härter und peinlicher, als wenn man wider sein Wort, darin er seinen ewigen und endlichen Willen offenbaret hat, muthwillig handelt. Dargegen daß in den schweren, betrübten, gefährlichen Zeiten Gott uns auch segne, so lasset uns nach seinem Wort und Geboten leben, und laßt uns täglich herzlich beten mit dem königlichen Propheten David, Psalm 119, 33.: Zeige mir HErr den Weg deiner Rechte, daß ich sie bewahre bis ans Ende. Unterweise mich, daß ich bewahre dein Gesetz und halte es von ganzem Herzen. Führe mich auf dem Steige deiner Gebote, denn ich habe Lust dazu. Dies erinnert auch das 23. Kapitel.

\section{Josua hält einen Landtag, das Volk zu vermahnen}
\quellseite{362}
\section{Summarium}

Hier ist eine sehr schöne Vermahnung, durch welche Josua als ein frommer und gottesfürchtiger Fürst, vor seinem Tod sein Volk vermahnet, daß sie sich nach dem Wort und Willen Gottes halten sollen, so werde Gott Glück und Heil ihnen widerfahren lassen, V. 6. Thun sie es aber nicht, sondern weichen ab von Gott und seinem Wort auf der Heiden abgöttisches Wesen, so werde Gott mit allerlei Unglück sie überschütten, bis sie endlich gar umkommen.

Solches Exempel dieses gottseligen Fürsten sollte alle Obrigkeit mit Fleiß bedenken, und ihm folgen, auf daß die Unterthanen durch ihrer Obrigkeit Exempel und ernstlichen Befehl lernen sich für allem Aergerniß hüten, und gottselig leben.

\section{Votum}

Helffe GOtt, der Herrscher aller Welt, daß Obrigkeit und Unterthanen seinem heiligen Wort folgen, auf daß beide seinen Segen haben und behalten an Land und Leuten, an Gut und Nahrung, an Leib und Seele bis zum ewigen Leben, Amen.

\chapter{Das 24. Capitel}

\section{Vorrede}

Gott pfleget die Undankbarkeit zu strafen, daß er von den undankbaren seine Gaben wiederum hinweg nimmt, wie er denn allen denen gethan hat, welchen er sein Wort gegeben, und sie es verachtet haben, und will solches noch thun. Wie davon lehret St. Paulus 2 Thess. 1, 10., da er also sagt: Dafür, daß sie die Liebe zur Wahrheit nicht haben angenommen, daß sie selig würden; darum wird ihnen Gott kräftige Irrthümer senden, daß sie glauben der Lüge: also, wenn man Gott nicht dankbar ist für den Frieden, so sendet er Krieg; wenn man nicht dankbar ist für die Gesundheit, so sendet er Krankheit, und will es dazu bei diesen zeitlichen Strafen nicht bleiben lassen, sondern er will die Undankbarkeit auch mit dem höllischen Feuer strafen. Derowegen weil wir bekennen müssen, daß bei diesen geschwinden Zeiten dennoch Gott dieser Stadt und Land viel geistliche und leibliche große Wohlthaten erzeiget hat, so laßt uns solche erkennen, und Gott mit Herz, Mund und gottseligem Leben dafür danken, auf daß er unserer Undankbarkeit wegen auch nicht vollends seine Gnade von uns wende, und denn zeitliche und ewige Strafen folgen. Darauf höret ferner an das 24. und letzte Kapitel.

\section{Letzter Landtag Josuas. Sein und Eleasars Tod}
\quellseite{363}
\section{Summarium}

Ist eine letzte Predigt, welche Josua der gottselige und fromme Fürst seinem Volk thut, daß sie bedenken wollen, was für große Wohlthaten ihnen Gott erzeiget habe, und sich hüten, daß sie nicht undankbar sein, und von solchem Gott und seinem Wort nicht abweichen, V. 2. Solcher Befehl ist trefflich nöthig, sonderlich wo neue Regenten angehen, denn da folget gemeiniglich Aenderung. Es geschieht aber solche Aenderung selten zur Besserung; wo alsdenn das zuschlägt, da man Gottesfurcht und Wort aus den Augen läßt, da muß Unglück mit Haufen folgen.
\quellseite{364}
\section{Summarium}

Ist ein fein Exempel, wie der gottseelige und fromme Fürst Josua, nachdem das Volk sich erboten, sie wollten von Gott nicht abweichen, sie auf ein neues mit einem Eide verpflichtete, auf daß sie nach seinem Tode solcher Verpflichtung gedächten, und desto weniger abweichen sollten. Solchem Exempel sollten heutiges Tages und allezeit die weltlichen Regenten folgen; aber leider es geschieht bei Wenigen. Deren aber findet man viel, die ihre Unterthanen wider Gott, wider alle Rechte und Billigkeit beleidigen, daß sie andere Papisten-Abgötterei halten, und das Evangelium nicht allein nicht annehmen, sondern auch verfolgen sollen. Und zu solchem treibet und bringet der Papst und seine Bischöfe; diese werden ihr Urtheil finden. Aber fromme gottseelige Obrigkeit soll dies Exempel wohl merken, und alle Mittel und Wege brauchen, daß Gottes Wort rein und lauter, und der Gottesdienst ungefälscht bleibe.

\section{Votum}

Also hat eure Liebe das Buch Josua gehöret, in welchem wir sehen, wie GOtt für sein Volk gestritten und die Heiden um ihrer Sünden willen gestrafet und vertilget hat; uns zum Exempel, daß wir uns bessern, und nicht auch dergleichen Strafe verdienen. Das verleihe uns unser lieber HErr JESus Christus, Amen.

\part{Das Buch der Richter}

\chapter{Das 1. Kapitel}

\section{Vorrede}

Das Buch der Richter ist ein Exempelbuch der Abwechselung, beides der Strafen und Erlösung der Kirchen Gottes. Denn wie nach Gottes Ordnung Nacht und Tag, Winter und Sommer sich abwechseln und verändern: also ist aus sonderem weisen Rath Gottes, der die Seinen wunderbarlich führet, und nichts so ernstlich und begierig sucht, als ihrer aller ewiges Heil und Seligkeit, eine stetige Abwechselung in der christlichen Kirchen, des Leides und der Freude, des Kreuzes und Trostes. Ins sonderheit aber sollen alle fromme Herzen ihnen wohl einbilden die Geschichte von Gideon und Simson, weil in diesen der HErr Christus und seine Kirche schön abgebildet sind. Denn als Esaias der Prophet im 9. Kapitel des HErrn Christi Wohlthaten beschreibt, so deutet er auf die Historia Gideons, und sagt: Denn du hast das Joch ihrer Last, und Ruthe ihrer Schulter, und den Stecken ihres Treibers zerbrochen, wie zur Zeit Midian. So wird unser HErr Christus ein Nazarener genannt, Math. 2. Welcher Name auch eigentlich zeiget auf die Geschichte von Simson. Darauf höret nun an das 1. Kapitel.

\section{Der Stamm Juda führet an Josuas Statt den Krieg}
% Scan page 365: No non-biblical apparatus, gloss, or structural title/synopsis retained in OCR blocks.
\quellseite{366}
\section{Summarium}

Hierin werden erstlich etliche Historien wiederholet, so zur Zeit Josua sich begeben haben. Darnach wird eine sonderliche Historia angezeigt, wie Gott den Tyrannen AdoniBesek habe gestraft, V. 5. Denn Gott, ob er wohl gebeut, daß weltliche Obrigkeit das Uebel mit Ernst soll strafen, so will er dennoch alle Tyrannei verboten haben. Wie es denn eine sonderliche Tyrannei an dem AdoniBesek gewesen ist, daß er siebenzig Könige an ihrem Leibe unbillig verstümmelt, und sie wie die Hunde das Brod unter dem Tisch hat lassen auflesen. Darum ward ihm vergolten, wie ers an andern verdienet hatte, V. 7. Solch Exempel soll weltliche Obrigkeit bedenken, und gleich wie sie auch nicht zu gelind wider öffentliches Aergerniß soll sein, also soll sie auch in der Strafe nicht tyrannisch sein, und den Zorn einfangen, und Maß halten. Denn es ist bald geschehen, daß man zu weit greift, wo man den Zorn läßt walten.

\section{Votum}

Der treue GOtt lindere uns die wohlverdienten Strafen, und regiere die Obrigkeit, und uns Unterthanen, daß die Obrigkeit alle Laster ernstlich strafe, und daß wir Unterthanen, gegen ihn und alle Menschen, hohen und niedern Standes, christlich leben, Amen.

\chapter{Das 2. Kapitel}

\section{Vorrede}

Auf gleiche Sünden folgen gleiche Strafen. Das lehret Gottes Wort, auch bezeugets die Erfahrung. Dennoch sind die Sünden jetzt am gemeinsten, welche Gott allzeit am greulichsten gestraft hat, als da sind Abgötterei, Hoffart, Wucher, Geiz, Unzucht, Büberei und dergleichen. Derwegen, ach ein jeder bedenke sich, wie er lebe, und ob sein Wandel Gott gefalle, ihm und andern nützlich oder schädlich sei. Der Fromme fahre fort in der Frömmigkeit nach Gottes Wort. Wer in Sünden lebet, der thue Buße, lasse von Sünden, und thue Gutes, und lebe also auch nach Gottes Gebot, daß er ihm und denen, derer er ja schonen soll, nicht Herzeleid und Strafe zurichte. Dies soll auch erinnern das 2. Kapitel.

\section{Der Kinder Israel Trägheit, Abfall, Strafe und Erledigung}

a Engel) der Priester Pinehas.

b Bochim heißt die Weinenden.
\quellseite{367}
\section{Summarium}

Ist gleich wie eine Vorrede über dies Buch, in welchem fürnehmlich wird angezeigt, wie die Kinder Israel auf der Heiden Abgötterei gerathen, und von Gott und seinem Wort sind abgewichen, V. 2. Und wie Gott solche Sünde habe damit gestrafet, daß sie ihren Feinden nicht mehr widerstehen möchten, sondern die Hand Gottes wider sie zum Unglück gewesen sei, wo sie hinaus wollten; auf daß man lerne in Gottes Furcht und Willen wandeln, sonst soll es uns auch also gehen.

Sonderlich aber ist zu merken, daß der Text spricht: Das Volk dient dem HErrn, so lang Josua lebte, und die Aeltesten nach ihm, V. 7. Denn wo fromme Regenten abgehen, die das Volk zum Wort und Gottesdienst halten, ist es leichtlich geschehen, daß das Wort sich wieder verliere. Wie wir denn heutiges Tages bei dieser jungen armen Welt auch besorgen müssen, wenn Gott diese wird hinweg nehmen, durch welche er uns sein Wort hat offenbaret.

\section{Votum}

Der barmherzige GOtt lindere die wohlverdienten Strafen, behüte uns für Abgötterei, und erhalte dieser Stadt und dem ganzen Lande sein heilig Wort bis an jüngsten Tag, und bewahre uns in wahrem Glauben und gottseligem Wandel, bis zum seligen Tod und Ende, Amen.

\chapter{Das 3. Capitel}

\section{Vorrede}

Es ist wahr, Gott erhöret die Sünder nicht, Joh. 9, 31. Aber dies soll verstanden werden von unbußfertigen Sündern, welche in offenbaren Sünden leben, und weiter zu sündigen gedenken. Denn also sagt Salomon Sprüchw. 28, 9.: Wer seine Ohren abwendet zu hören das Gesetz, deß Gebet ist ein Greuel. Aber die Sünder, die Buße thun, und in wahrem Glauben beten, die werden erhöret. Davon sagt Gott also 2 Chron. 7, 14.: Wenn sie beten, und mein Angesicht suchen, und sich von ihren bösen Wegen bekehren werden, so will ich vom
\quellseite{368}
Himmel hören, und ihnen ihre Sünde vergeben, und ihr Land heilen. So helfe der allmächtige, gütige Gott, daß wir allezeit aus bußfertigen Herzen thun unser Gebet, und mit David sprechen Psalm 5, 8.: Ich aber, HErr, will in dein Haus gehen, auf deine große Güte, und anbeten gegen deinem heiligen Tempel in deiner Furcht, das ist, mit einem bußfertigen, gläubigen Herzen. Darauf höret ferner an das 3. Kapitel.

\section{Athniel, Ehud, Samgar, und ihre Thaten}

\section{Summarium}

Dies ist eine Historia, wie Gott etliche Heiden habe lassen überbleiben, auf daß er sein Volk versuchte, ob sie auch an seinen Geboten halten, und nicht der Heiden Exempel nach Abgötterei treiben wollten, V. 1. Aber sie haben sich das Aergerniß lassen übergehen, und gesündigt, V. 6. Darum hat Gott der Syrer König wohl acht Jahr über sie herrschen lassen, V. 8. Solches soll uns reizen, daß wir Gott fürchten, und in seinem Willen und nach seinem Wort wandeln. Sintemal er die Sünde durch Krieg und fremde Herrschaft pflegt zu strafen. Wiederum ist es tröstlich, daß Gott dennoch solcher Sünder Gebet erhöret, wenn sie um Hülfe schreien, und läßt ihnen Hülfe widerfahren.
\quellseite{369}
a Die Rathstube, ubi ordinatae erant sedes.

\section{Summarium}

Ist wiederum ein Exempel, wie Gott sein Volk um der Sünde willen mit Krieg und fremder Herrschaft plaget, aber hernach, da sie um Hülfe schreien, wieder durch einen frommen Fürsten errettet habe, V. 15. Daß nun Ehud den König der Moabiter heimlich mit Eisen ersticht, V. 21., ist nicht dahin zu deuten, als sollten solche Exempel Ausführer und Mörder billigen, sondern man muß auf den Befehl, das Amt und den Beruf Ehud sehen, der war von Gott dazu berufen, daß Israel durch ihn von den Moabitern sollte errettet werden, und den König erwürgen; wie er denn zum Könige sagt, ich habe Gottes Wort an dich, V. 20., das ist, ich habe einen Befehl von Gott, der dich belanget. Wo nun solcher sonderlicher, gewisser und ausgedrückter Befehl Gottes nicht ist, da soll man sich aller Gewalt enthalten, und auch das Unrecht geduldig leiden, V. 25., bis Gott selbst durch andere Mittel von den Tyrannen erledigt und hilft.

\section{Votum}

Dem allmächtigen und barmherzigen GOtt sei Lob und Dank, der die wohlverdienten Strafen lindert; der wolle ferner unser Gebet erhören, uns in allen Nöthen beistehen, und unter dem Schutz unserer lieben Obrigkeit uns segnen und behüten, Amen.

\chapter{Das 4. Kapitel}

\section{Vorrede}

Im Psalm 50, 25. spricht der HErr: Rufe mich an in der Noth, so will ich dich erretten, so sollst du mich preisen. Hiermit lehret Gott, da jemand gesündiget hat, daß er sich seine Sünde nicht soll von Gott abschrecken lassen, sondern soll zu ihm bringen ein zerknirschtes Herz, und soll sich des HErrn Christ in wahrem Glauben trösten, und einen guten Vorsatz haben, sein Leben zu bessern, und soll also ihn anrufen, so wollte er erhören und helfen. Dies lehret auch St. Paulus 2 Timoth. 2, 19. und spricht: Es trete ab von Ungerechtigkeit, der den Namen Christi anrufet. Und Salomon in seinen Sprüchen 15, 8. sagt also: Des Gottlosen Opfer ist dem HErrn ein Greuel, aber das Gebet der Frommen ist ihm angenehm. Deswegen so laßt uns Buße thun, und Gott im Geist und in der Wahrheit anrufen. Es ist gewiß, wir werden entweder Abwendung oder ja Linderung der gegenwärtigen Strafen erlangen. Darauf hört auch an das 4. Kapitel.

\section{Debora, Barak und Jael's Heldenmuth}
\quellseite{370}
\section{Summarium}

Ist eine schöne Historia, wie Gott das Volk um ihrer Sünde willen strafet, und hernach, da sie um Hülfe schreien, erhöret, und ihnen hilft durch die Prophetin Debora, dazumal Richterin im Volk Israel, V. 4.; auf daß man sehe, wie Gottes Hülfe nicht an großer Macht und Gewalt liegt. Er kann ebenso wohl durch ein Weib erretten und helfen, als durch einen Mann, wo man sich nur seines Willens hält, und nicht Ursach zur Strafe gibt. Denn es liegt alles an dem, ob Gott den Feinden das Herz nehme oder nicht.

\section{Summarium}

Ist ein fein Exempel, wie Gott den Tyrannen eine Weile pfleget zuzusehen; wenn aber ihr Stündlein kommt, müssen sie einer leichten Ursach halben zu Boden gehen. Wie denn der tapfere Kriegsmann Sisera von einem Weib erwürget wird, unangesehen, daß zwischen ihrem Hause und dem König Jabin Friede war.
\quellseite{371}
Denn an dem war ihr, als einer Regentin, die Gottes Befehl dazu hatte, mehr gelegen, daß ihr Volk errettet würde, denn an der Privat-Verwandtnis, so zwischen ihnen war.

\section{Votum}

Der treue, gütige GOtt stehe auch heute seiner Kirche bei, steure und wehre allen deren Verfolgern und Tyrannen, und verleihe unserer lieben Obrigkeit ferner ruhige und glückliche Regierung, den armen Unterthanen zu seliger Wohlfahrt, Amen.

\chapter{Das 5. Capitel}

\section{Vorrede}

Wie wir nachlässig zum Gebet sind, also sind wir viel nachlässiger zur Danksagung. Die Undankbarkeit aber kann Gott nicht leiden, wenn dieselbige schon ein Mensch dem andern beweiset, wie wir lesen 1 Mos. 31, 12., da er die Undankbarkeit Laban gegen Jakob strafet; und zuvor 1 Mos. 16, 9., da er auch die Undankbarkeit Hagar gegen Sara strafet. Denn es stehet geschrieben in Sprüchen Salomonis 17, 13.: Wer Gutes mit Bösem vergilt, von deß Hause wird Böses nicht lassen. Viel weniger wird nun Gott leiden, daß man ihm für seine Wohlthaten undankbar ist. Derwegen, Jünglinge und Jungfrauen, Alte mit den Jungen, sollen allzeit loben den Namen des HErrn, Psalm 148, 12. Und sollen ohn Unterlaß beten, und in allen Dingen danksagen, 1 Thessal. 5, 17. Darauf höret an das 5. Capitel.

\section{Deboras und Baraks Triumphlied}

a Dies Lied will so viel sagen, daß Gott habe den Sisera geschlagen durch die geringsten Leute in Israel, daß die Geringen auch einmal hoch und groß worden sind, da die großen Geschlechter Israel stille saßen, und sie verblieben in Nöthen; daß ist das Neue, das der HErr ermahlet hat, da sind die Bauern Israel prächtig, und auch Herren worden.

b Krumme) das ist, es war kein Regiment noch Ordnung im Lande.

c Lobet) das ist, ihr Herren, Richter und gemeine Leute.

d Schöpfern) das ist, da die Schützen Sisera schrieen für Noth am Wasser Kison, da man pfleget zu schöpfen, da half Gott seinem Bauern, und ließ das Recht gehen.

e Wurzel) das ist, Josua war der erste Fürst aus dem Stamm Ephraim, der schlug Amalek, und nach ihm die andern, bis daß Sebulon auch einmal einen Josua überkommen hat wider Sisera. Und merke, sie nennet die Fürsten Regierer, die mit der Schreibfedern streiten, das ist, sie gewinnen (geistlich davon zu reden) mehr durch den Glauben an Gottes Wort, denn mit dem Schwert.

f Hürden) das ist, du bleibest daheim, ob du wohl hörtest das arme Häuflein zu Felde blasen, und hattest doch nahe zu ihm.
\quellseite{372}
\section{Summarium}

Die ist ein Danklied, daß Gott seinem Volk wider die Cananiter geholfen hat. Denn bei den Frommen sollen diese zwei Stücke auf einander gehen: daß man erstlich erkenne, wie Gott alles aus lauter Güte und Barmherzigkeit uns widerfahren lasse; das andere, daß man auch dafür dankbar sei, und solcher Wohlthat nicht vergesse.

Am Ende ist ein feiner Spruch, daß die Prophetin wünschet, daß alle Feinde Gottes sollen umkommen, wie Sisera, aber die Gott lieben, sollen aufgehen, wie die Sonne in ihrer Macht, welche niemand dämpfen kann, sondern muß sie frei gehen und leuchten lassen, V. 13. Solcher Wunsch widerfähret durchaus der christlichen Kirchen, die hat die Anfechtung und Hinderniß in der Welt sehr gewaltig, daß es scheinet, als sollte es alles zu Boden gehen. Aber wie Debora hie bittet, dringet sie doch immer fort wie die liebe Sonne, so ihr Licht nicht einfangen noch aufhalten läßt, wie Christus sagt, daß auch der Höllen Pforten sie nicht bewältigen sollen.

\section{Votum}

Lob, Ehre, Preis und Dank sei GOtt für alle seine Wohlthaten; der wolle ferner uns und die ganze Christenheit unter der vorgefallenen wunderlichen Zeit und Veränderungen segnen, behüten und erhalten, daß bis an den jüngsten Tag bleiben, die ihn erkennen und preisen, Amen.

\chapter{Das 6. Kapitel}

\section{Vorrede}

Wer Glück in seinem ehrlichen Beruf, Amt und Stande haben will, der muß dreierlei in Acht nehmen. Erstlich soll er gewiß sein, daß sein Beruf Gott gefalle; dann soll er sich in seinem Beruf alles möglichen Gehorsams gegen Gottes Wort befleißen; und zum dritten soll er Gott um glücklichen Fortgang seines Berufs täglich und herzlich anrufen. Deß ist ein Exempel der theure Fürst Gideon, von dem das 6. Kapitel.

\section{Gideon zum Richter berufen}
\quellseite{373}
\section{Summarium}

Hie fähet an die Historia des theuern und gottseeligen Fürsten Gideonis, welchen Gott dazu berufet, daß er ein Fürst sein und dem Volk Israel wider die Midianiter helfen soll. In dieser Historia sehen wir erstlich, wie Gott die Sünden durch Krieg und Unfrieden pfleget zu strafen; aber dennoch die Strafe abwendet, wenn man Hülfe bei ihm suchet. Daß also durchaus solche Historien eitel Exempel der Buße und des Glaubens sind, daß wir uns für Sünden sollen hüten, weil sie Gott mit strafen, und wiederum auf Gottes Güte sehen, und uns derselben trösten, wenn wir die Strafe empfunden und tragen; sintemal er hie so gnädig erhöret, da man ihn bittet, und sich Hülfe zu ihm versiehet. Darum ist sonderlich zu merken an Gideon, da der Engel ihn grüßet: Der HErr mit dir, daß ers beredet und spricht: Ist der HErr mit uns, warum gehet es denn so übel? V. 13.
\quellseite{374}
Denn also pflegt die Vernunft zu urtheilen, wenn es übel gehet, daß Gott gar sey von uns gewichen; so doch Gott das Unglück uns darum auferlegt, daß wir uns zu ihm halten, und bessern, und seine Hülfe suchen sollen.

\section{Summarium}

Ist ein sehr schön Exempel: ehe Gott seinem Volk durch Gideon hilft, befehlet er, daß er zuvor den Altar Baal abbrechen, und den Hain oder Forst umhauen soll. Denn wo Glück und Heil soll sein, da muß am ersten durch die Obrigkeit Abgötterei und andere Sünde abgeschafft sein.

Daß nun Gideon zwei Zeichen mit dem Fell fordert, V. 36., geschieht alles darum, daß er seines Berufs desto gewisser sei, und mit größerem Muth die Sachen angreifen möge; wie denn die Historia hernach anzeigen wird, wie ihn Gott zu einem sonderbaren hohen Werck hat brauchen wollen, daß es sehr demüthig war, daß er zuvor seines Berufs gewiß wäre.

\section{Votum}

GOtt sei uns gnädig, erhöre unser Gebet und helfe, daß aller christlichen Obrigkeit Regierung diene zur Heiligung seines Namens, daß sein Reich zu uns komme, und daß sein Wille geschehe, wie im Himmel, also auch auf Erden, von uns und allen Menschen, Amen.

\chapter{Das 7. Capitel}

\section{Vorrede}

Des Glaubens Natur und Art ist, daß er sich fröhlich erwäget, die Dinge zu glauben, die er nicht fühlet, noch mit allen seinen Kräften inwendig und auswendig ermessen kann. Welche
\quellseite{375}
Menschen nun einen solchen Glauben haben, und in dem an Christum Jesum festhalten, die macht er zu Herren, wie er selbst ist, über Sünde, Tod, Teufel, Welt und alle Dinge. Derwegen trachten wir billig nach dem wahren Glauben und seufzen: O Herr, Gott, heiliger Geist, stärk unsern Glauben immerdar, daß wir an Christo bleiben, dem treuen Heiland, und ihm aus ganzer Macht allzeit vertrauen, Amen. Dies erinnert auch das 7. Kapitel.

\section{Gideon erhält den Sieg wider die Midianiter}

\section{Summarium}

Ist eine herrliche Historia, wie die Midianiter durch Gideon erlegt sind. Denn hie stehet man,
\quellseite{376}
wie einen großen Glauben der Fürst auf Gott und seine Zusagung gehabt hat, daß er nach dem Befehl Gottes das Volk wieder heimziehen heißet, und nicht mehr bei sich behält denn dreihundert Mann, so doch der Feinde eine unzählige Menge war, V. 3.; auf daß wir lernen, wie alles an Gott und seiner Hülfe, und nicht mehr an Menschen Vernunft, Weisheit, noch Macht liege; aber solches glaubt die Vernunft nicht.

\section{Summarium}

Ist die Historia, wie die Midianiter wunderbarlich ohne allen Schwerthschlag erlegt sind. Solche Historia zeucht der Prophet Es. 9, 7. als eine Figur, auf die Erlösung, so uns durch unsern HErrn Christum widerfahren ist, da wir von dem geistlichen Midian, das ist von der Sünde, Tod und Gesetze, und des Teufels Gewalt erlöset sind, ohne unser Zuthun, allein durch Christus Sterben und Auferstehung.

Mit solcher Auslegung reimet sich auch die heimliche Deutung. Die Midianiter sind, wie gesagt, der Christen ärgste Feinde, nämlich, die Sünde, der Tod, Welt, Teufel und Fleisch. Solche Midianiter lassen sich mit dem Schwert nicht schlagen, sondern das muß die Rüstung sein, daß wir in der rechten Hand haben die Posaunen, das ist das Wort Gottes, und in der linken die Lampen des Glaubens, welche wir tragen im irdischen Gefäß, wie St. Paulus sagt. Desselbige Gefäß muß zerschlagen werden, das ist, Fleisch und Blut muß gedämpft und täglich getödtet werden, bis wir endlich gar sterben; so gehet es denn, daß unsere Feinde überwunden und geschlagen, wir aber durch Christum selig werden, Amen.

\section{Votum}

Wir danken dem HErrn Jesu Christo für seine Victori und Sieg wider seine Feinde. Der wolle unsern schwachen Glauben in uns stärken, und helfen, daß wir auch überwinden alle unsere geistlichen und leiblichen Feinde, im Leben und Sterben, Amen.

\chapter{Das 8. Kapitel}

\section{Vorrede}

Der heilige Apostel St. Jacob sagt in seiner Epistel am 1. Kapitel V. 11.: Die Sonne gehet auf mit der Hitze, und das Gras verwelket, und die Blume fällt ab, und sein schöne Gestalt verdirbet. Also hat ein jeder Mensch sein Kreuz, das von Tage zu Tage des Leibes und aller Gliedmaßen Kräfte schwächet, bis endlich folget der bittere Tod, der keinen Unterschied hält unter den Menschen, und derhalben sobald ein Junges würget, als ein Altes, dem niemand Widerstand thun kann, der die Menschen oft plötzlich überfällt und hinreißt, der auch niemand zuvor lässet wissen, wenn, wie oder wo er sterben werde. Demgegen wollen wir Buße thun, demüthig sein, Gottes Wort hören, wachen und beten, und das neben christlich leben, auf daß wir darauf auch fröhlich und selig sterben in Christo Jesu, unserm HErrn, Amen. Darauf höret weiter an das 8. Kapitel.

\section[Die übergebliebenen Midianiter werden geschlagen, und Sebah und]{Die übergebliebenen Midianiter werden geschlagen, und Sebah und Zalmunna gefangen}
\quellseite{377}
a Das ist, erfahren, inne werden.

\section{Summarium}

Ist ein fein Exempel, wie der fromme Fürst Gideon sich nicht erhebt um des großen Sieges willen, welchen ihm Gott gegeben hatte, sondern demüthiget sich, und durch solche Demuth erhält er Friede, da sonst durch Hoffart und Stolz, wie allwege geschicht, Unfriede wäre erreget worden.
\quellseite{378}
\section{Summarium}

Ist eine sonderliche Historia, wie oft etwas, daß man guter Meinung, auch aus guten Ursachen thut, übel geräth und groß Unglück verursachet, wie hie Gideon fordert die güldenen Ringe von den Kriegsleuten, und machet einen Leibrock daraus, V. 27., auf daß es solt ein Denkzeichen sein des herrlichen Sieges, welchen ihm Gott gegeben hat. Aber es gerieth solcher Anschlag übel. Denn sie mißbrauchten hernach solches Leibrocks zu Abgötterei, und hielten ihn dem gleich, welchen Gott befohlen hatte zu machen und der Hohepriester anlegen mußte. Es ist aber solcher Abgötterei halben, wie folgen wird, das ganze Volk hart gestraft worden.

Darum soll jedermann darauf sehen, daß er nichts fürnehme, noch anfange, er habe denn Gottes Wort und Befehl; wer dabei bleibt, der darf sich keines Mißbrauchs oder einiger Sünde besorgen.

\section{Votum}

Helfe Gott, daß sein Wort eine Leuchte sei unserer Füße, und ein Wegzeiger in unserm Beruf, und regiere uns, daß wir uns für allerley Abgötterei und andern Sünden und Lastern hüten, bis an unser selig Ende, Amen.

\chapter{Das 9. Capitel}

\section{Vorrede}

Niemand soll seines Nutzes wegen wider Gott und seinen Nächsten handeln, denn was dawider geschieht, obs schon eine Zeit lang einen glücklichen Fortgang hat, so wirds doch endlich gestraft. Derwegen nehme ein jeglicher wahr in seinem Berufe seines Nächsten, und schade nicht, sondern diene ihm, so sammlet er ihm und den Seinem Gottes Segen. Dieses lehrt St. Paulus 1 Corinth. 13., da er sagt: Die Liebe suchet nicht das Ihre, das ist: sie trachtet nicht auf ihr eigen Gut, Nutz, Ehr und Gemach, sondern trachtet darauf, wie sie des Nächsten Nutz, Ehre, Leib und Leben fördere. Ach, das helfe uns der treue Gott. Denn weil wir von Natur alle eigennützig sind, so ist wohl zu achten, daß wir die christliche, brüderliche Liebe nicht von uns selbst haben, sondern Gott muß sie uns geben durch den heiligen Geist. Ferner höret an den ersten Theil des 9. Capitels.

\section{Der Brudermörder Abimelech kommt um}
\quellseite{379}
a Es bestehet nicht, es ist ein Schuchal.

\section{Summarium}

Ist eine Historia, wie der Tyrann Abimelech sich des Regiments unterstehet und würget siebenzig Söhne des Fürsten Gideonis, auf daß das Regiment ihm allein bliebe, V. 5. Ob nun wohl solche Strafe Gideon verdienet habe, damit, daß er mit dem Leibrock, den er machen ließ, zur Abgötterei Ursache gab, lasse ich bleiben. Gleichwohl sagt der Text Jud. 8, 27. mit klaren Worten, daß solche That, die doch aus guter Meinung fürgenommen war, nicht allein Gideon und seinem Hause zum Aergerniß sei geraten, sondern daß auch das ganze Israel damit Abgötterei getrieben habe. Dadurch wir ja treulich gewarnet werden, daß wir bei Leibe nichts fürnehmen oder anfangen sollen aus eigener Wahl, guter Meinung und Andacht, Gott damit zu dienen und ehren, es habe einen Schein der Heiligkeit so groß und gut es kann. Geschichts ohne Gottes Wort und Befehl, so gibt es Ursache zu allerlei Irrthum und Abgötterei, wenns noch einst ein heiliger Patriarch Gideon aus sonderlicher Andacht angerichtet hätte. O wie sicher, frech und durstig hat hinein gefahren der höllische Drache zu Rom, Gott mache solches Lästernds schier ein Ende, Amen.

\section{Der andere Theil des 9. Kapitels}

\section{Vorrede}

Um die Sünde ists so ein arg Ding, daß dieselbigen von der Erde gen Himmel schreien, und für Gottes Angesicht gleich treten, und die Rache und Strafe erfordern, daß er endlich seinen Zorn darauf muß lassen ergehen, wo man sich nicht durch wahre Buße zu ihm bekehret. Die schreienden Sünden aber werden von Alten diese genennet: Todtschlag, Wucher, Unterdrückung der Wittwen und Waisen, Sodomie, und verdienten Lohn fürhalten. Derwegen wollen wir ablassen von unsern Sünden, und die wir leider gethan, wollen wir Gott abbitten im Namen und auf die Vorbitte unsers Herrn Jesu Christi, und wollen Gott und unserm Nächsten dienen im Glauben, mit wahrer Liebe, so wird es Gott eine Freude sein, und uns ein Nutzen, zeitlich und ewig. Dazu höret an den andern Theil des 9. Kapitels.
% Scan page 380: No non-biblical apparatus, gloss, or structural title/synopsis retained in OCR blocks.
\quellseite{381}
\section{Summarium}

Ist eine Historia, wie Gott eine Weile den Tyrannen und bösen Buben zu Sichem, so ihm zu seiner Tyrannei geholfen hatten, zusiehet, aber endlich sie bezahlet, wie sie verdienet haben. Und ist sonderlich hie zu merken, daß der Text mit klaren Worten sagt, daß Gott die tyrannischen Regenten mit Aufruhr strafet; und die, so zum Aufruhr helfen, wie hie Gaal thut, nur Schande und Schaden davon bringen. Solche Exempel dienen alle dazu, daß man lerne Gott fürchten und fromm sein. Sintemal Gott die Sünde nicht will ungestraft lassen, wie man hie siehet an den Sichemitern und ihrer Stadt.

\section{Summarium}

Ist eine Historia, wie, nachdem der Tyrann Abimelech die bösen Buben zu Sichem wohl gestraft hatte, er auch muß herhalten und bezahlen, was er mit dem Mord an seinen Brüdern und sonst verdienet hatte. Und ist sonderlich zu merken, daß gemeiniglich die Tyrannen schändlich und liederlich umkommen. Wie den Abimelech ein Weib zu Tode wirft und sein Knecht ihn vollends erstechen muß, V. 53. Auf daß ja die Sünde endlich nicht ungestraft bleibe, und jedermann empfähe, nach dem er verdienet hat.

\section{Votum}

Regiere uns alle der treue GOtt, daß wir der Sünden ja gram werden, und dagegen allen Tugenden hold und günstig, und daß wir leben in Heiligkeit und Gerechtigkeit, Ihm zu Lob und dem Nächsten zum Frommen, in Kraft des heil. Geistes, Amen.

\chapter{Das 10. Kapitel}

\section{Vorrede}

Gott strafet die Sünde; wie solches unsere Ohren hören und unsre Augen sehen. Viele habens auch gefühlet an ihrem Leibe, viele befindens auch an ihrer Nahrung. Ach der Blindheit und wehe der Sicherheit, daß dennoch viele von wissentlichen, muthwilligen und vorsätzlichen Sünden nicht ablassen. Derwegen wollen wir wachen und beten, Gott wolle uns in wahrer Buße erhalten, auf daß wir mit unsern Sünden nicht uns noch andern schädlich sein. Denn der Knecht, der seines Herrn Willen weiß, und hat sich nicht bereitet, auch nicht nach seinem Willen gethan, der wird viel Streiche leiden müssen, sagt unser Herr Jesus Christus Luc. 12, 47. Darauf höret nun an das 10. Kapitel.

\section[Israel wird wegen der Abgötterei von den Philistern]{Israel wird wegen der Abgötterei von den Philistern gedrängt}
\quellseite{382}
\section{Summarium}

Es ist eine Historia, wie Gott sein Volk durch die Ammoniter und Philister um ihrer Abgötterei und Sünden willen gestraft habe. Der Text zeiget hie sonderlich an, daß Gott solche Strafe ihnen widerfahren lasse, nicht daß er sie gar verderben, sondern dadurch dahin bringen will, daß sie ihre Sünde erkennen und davon ablassen sollen, V. 14. Daß aber Gott durch den Propheten sagen läßt: Er wolle nicht mehr helfen, sie sollen bei den Götzen Hülfe suchen, welchen sie gedienet haben, V. 15., geschieht darum, auf daß sie ihre Sünde mit größerem Ernst beweinten, und sich desto eher bekehrten. Also gehen gemeiniglich alle Historien dieses Buches dahin, daß man entweder Buße thue und sich bessere, oder gewarte, daß Gott zürnen und heftig durch Krieg und ander Unglück unsere Sünde strafen werde.

\section{Votum}

Der barmherzige GOtt lindere uns die Strafen, so man mit Abgötterei, mit falschem Gottesdienst und dergleichen Sünden verdienet, und helfe, daß wir in seiner Furcht leben, und ihm allzeit dienen nach seinem Wort und Willen, im geistlichen und weltlichen Leben. Amen.

\chapter{Das 11. Kapitel}

\section{Vorrede}

Wohl unserer Stadt und Land, wohl eines jeden Leib und Seele, wohl eines jeden Haus und Nahrung, da wir allezeit und sonderlich in diesen gefährlichen Zeiten fleißig, herzlich und in wahrem Glauben beten. Denn keine Noth ist so groß, und keine Gefahr so grausam, und kein Kreuz so schwer, ein gläubig Gebet kann es entweder abwenden oder lindern. Und der alte Kirchenlehrer Origenes sagt: Eines frommen Christen Gebet ist stärker, und richtet mehr aus, als die gewappnete Hand vieles Kriegsvolks. Denn dieweil das Gebet den Himmel durchbringet, sollte es dann nicht auch auf Erden mächtig sein und alle Feindschaft überwinden? Dies ist fürnehmlich ein Exempel des 11. Kapitels.

\section{Jephthahs Sieg und Gelübde}
% Scan page 383: No non-biblical apparatus, gloss, or structural title/synopsis retained in OCR blocks.
\quellseite{384}
\section{Summarium}

Ist eine Historia, wie Gott sein Volk durch den Fürsten Jephthah errettet habe, da sie sich gebessert hatten, und zu Gott um Hülfe schrieen. Es ist aber ein fein Exempel hie, daß Jephthah der Untreue vergisset, die ihm von den Seinen widerfahren war, und sie derselbigen nicht will entgelten lassen, da sie seiner Hülfe bedurften, V. 7. Denn Gott will, wo Noth vorhanden ist und wir helfen können, daß man da weder Dank noch Undank ansehen soll, sondern thun, was er befohlen hat.

Auch siehet man hie, wie Gottes Urtheil so wunderlich gehet. Jephthah war seinen Stiefbrüdern zum Mittlern nicht gut genug; aber Gott war er gut genug zum Regenten und Heiland, durch den er sein Volk Israel errettet.

a Beuget, oder demüthigest mich; Gott hat mich hoch erhebt durch diesen Sieg, daß ich mein Haupt hoch und fröhlich aufrichte. Aber du beugest mich, daß ich den Kopf muß niederschlagen mit großem Herzeleid, und solche hohe Freude zum tiefen Herzeleid kehren.

b Wie er gelobet hatte) man wil, er habe sie nicht geopfert; aber der Text stehet da klar. So siehet man auch beyde an den Richtern und Königen, daß sie nach großen Thaten haben auch große Thorheit müssen begehen, zu verhüten den leidigen Hochmuth.

\section{Summarium}

Ist eine sonderliche Historia von dem Gelübde Jephthah, welches er unbedacht that, und mit Sünden vollstreckte. Denn man ist kein Gelübde schuldig, das wider Gott und sein Wort ist. Weil nun Jephthah sein Gelübde sonderlich hat thun wollen, Gott zu Dank und Ehre, sollte er je seine Tochter nicht haben würgen lassen. Sintemal Gott den Eltern befohlen hat, daß sie ihre Kinder in Gottesfurcht aufziehen sollen.

Aber in solchen Exempeln sehen wir, daß auch hohe heilige Leute, die große wunderliche Thaten ausgerichtet, auch große Thorheit begangen haben. Es ist keiner je so heilig noch vollkommen gewesen (das Haupt allein aller Heiligen ausgenommen), der nicht oft sich versündiget, und derhalben hat beten müssen (wie der 32. Psalm V. 5. zeiget): Vergib uns unsre Schuld, und Vergebung der Sünden glauben. Wir arme, elende Leute können dem Erzvater Jephthah, der wie man aus der Historia siehet, ein hochgeleuchter Mann gewesen ist, voll Geistes und Glaubens, und doch gleichwohl schmerzlich strauchelt, das Wasser nicht reichen, noch sind wir so sicher und der Sache, wie wir uns dünken lassen, so gewiß, als könnten wir nimmermehr irren, noch fehlen. Es heißt aber: sei nicht stolz, sondern fürchte dich. Und wer sich dünken läßt, er stehe, der mag wohl zusehen, daß er nicht falle.

\section{Votum}

Helffe der treue GOtt, daß wir uns nach seinem Wort gehorsamlich verhalten, fleißig
\quellseite{385}
beten, in unserm Beruf treu sein; auch unserm Nächsten dienen, auf daß er sich auch nach seinem Wort väterlich gegen uns erzeige und verhalte, Amen.

\chapter{Das 12. Kapitel}

\section{Vorrede}

Salomon in seinen Sprüchen sagt im 16. Kapitel V. 18.: Wer zu Grunde gehen soll, der wird zuvor stolz. Hoffart und stolzer Muth kommt vor dem Fall. Und St. Paulus an die Römer im 13. V. 2. spricht: Wer sich wider die Obrigkeit setzet, der widerstrebet Gottes Ordnung: die aber widerstreben, werden über sich ein Urtheil empfangen. Derwegen menge sich niemand unter die Aufrührischen. Denn es wird kein Aufruhr ungestraft bleiben, und auch derer keiner, so sonst der weltlichen Obrigkeit guter Ordnung mutwillig und wissentlich widerstreben. Sondern jedermann sei unterthan der Obrigkeit, die Gewalt über ihn hat. Kehret und warnet abermal Salomon und St. Paulus, Spr. Salomon 24, 21. Röm. 13, 1. Darauf höret auch davon an das 12. Kapitel.

\section{Niederlage der Ephraimiter}

a Id est, qui ad Ephraim fugerunt in periculis pro salute. Non vos eripitis nos, sed nos vos fugitivos, quia sumus regia tribus. Das ist, die zu Ephraim geflohen sind in der Gefahr ihres Lebens. Nicht ihr errettet uns, sondern wir euch Flüchtlinge; denn wir sind von der königlichen Stamm.

b Schibboleth) heißt eine Kornähre, oder auch wohl eine Mundstrafe.

\section{Summarium}

Ist eine Historia, wie Gott die Ephraiter, so unbilliger Weise aus Hoffart wider den Fürsten Jephthah einen Aufruhr erregeten, hart gestraft hat, daß ihrer zwei und vierzig tausend erschlagen sind, V. 6. Also gehet es: wo man Gott nicht fürchtet, strafet er bisweilen Land und Leute mit Krieg und Aufruhr, aber demnach endlich gehet es über den Aufruhrherrn aus, daß sie müssen zu Grund gehen. Denn Gott will weder Hoffart noch Ungehorsam ungestraft lassen.

\section{Votum}

Der treue GOtt helfe, daß wir uns vor diesem und allen andern Sünden und Lastern hüten, und in Demuth und Gehorsam unserer christlichen Obrigkeit allzeit unterthan sein, derselben zur Freude, uns zu Lob und Ehren und andern zu guten Exempeln, Amen.
\quellseite{386}
\chapter{Das 13. Kapitel}

\section{Vorrede}

Viele Eltern sind in den Gedanken, es stehe ihnen frei, wie oder wozu sie ihre Kinder aufziehen. Viele lassen darüber die Hauszucht also fallen, und die Kinder untereinander aufwachsen, wie das arme dumme Vieh. Damit aber laden die Eltern gar schwerlich auf sich Gottes Zorn, Beschwerung des Gewissens, und endlich schwere Strafe und Rechnung. Dies zu verhüten, lehret die heilige Schrift in Sprüchen und Exempeln, was in dem Theil Gott und ihr Amt von ihnen fordert. Sprüche sind, an die Ephes. 6. Kapitel V. 4. sagt St. Paulus: Ihr Väter, reitzet eure Kinder nicht zu Zorn, sondern ziehet sie auf in der Zucht und Vermahnung zum Herrn. Und Salomon sagt in Sprüchen 22, 15. und 29, 17.: Züchtige deinen Sohn, denn Thorheit steckt dem Knaben im Herzen, aber die Ruthe der Zucht wird sie ferne von ihm treiben. Denn Ruthe und Strafe gibt Weisheit. Ein ausbündig Exempel aber hört an aus dem 13. Kapitel.

\section{Simsons Empfängniß und Geburt}

a Es ging wunderlich zu, wie folget, daß der Engel verschwand, und in der Flamme auffuhr.
\quellseite{387}
\section{Summarium}

Sie gehet an die Historia von dem theuren Fürsten und Helden Simson; in welcher man erstlich siehet, wie es eine große Gabe Gottes ist, wo Gott Eheleuten Kinder gibt, V. 3. Darnach ist sonderlich zu merken, wie die Eltern Simsons des Engels wieder begehren, daß sie erfahren, wie sie ihren Sohn nach dem Willen Gottes sollen auferziehen, V. 8. Denn dies will Amts halben den Eltern sonderlich gebühren, daß sie darauf Achtung haben, wie ihre Kinder in Gottes Furcht und Willen recht auferzogen werden. Und dies ist das rechte Erbe, und der höchste Schatz, welchen die Eltern ihren Kindern können lassen. Aber die Welt, wie man siehet, fraget nach solchem Amt am wenigsten. Jedermann denket, wie er seinen Kindern viel Geld und Gut lasse. Aber weil es an rechter Zucht und Furcht Gottes fehlet, sehen wir, daß weder Glück noch Heil bei den Kindern ist, und es oft viel besser wäre, es ließen die Eltern ihren Kindern weder Heller noch Pfennig, so bliebe mancher Unrath und Sünde dahinten. Darum sollen die Eltern dies Exempel des frommen Vaters Manoah fleißig merken.

\section{Votum}

Gebe GOtt, der ein Vater ist über alles, im Himmel und auf Erden, daß alle Eltern folgen, GOtt zu Ehren, ihnen zu gutem Gewissen und endlich zu ihrer, auch ihrer Kinder Seelen Heil und Seligkeit, und helfe, daß die lieben Kinder sich auch ziehen lassen und folgen, Amen.

\chapter{Das 14. Kapitel}

\section{Vorrede}

Die Kinder sollen ihre Eltern mit Gedanken, Gebärden, Worten, Werken und Gehorsam ehren, und dies darum, daß es Gott ernstlich befohlen hat. Und wenn auch gleich kein göttlich Statut oder Gebot darauf gesetzt wäre, so bringet es doch das natürliche Recht mit sich, daß wir sollen denen Gutes und Ehre erzeigen, von denen wir alle Wohltat vielfältig empfangen haben. So ist auch gewiß, was man an den Eltern thut, das wird man auch an den Kindern zu gewarten haben. Ferner soll die Kinder, daß sie ihre Eltern ehren, bewegen, die gnädige Verheißung, welche das göttliche Statut also bald mitbringt, und sagt: darum sollst du es thun, auf daß dirs wohlgehe, und lange lebest auf Erden. Ueber das soll sie verwünschen in dem Gegenteil die schreckliche Dräuung auf den Ungehorsam, als der zeitlich und ewig mit Schand und Schaden und mit höllischem Feuer gestraft wird. Dies erinnert auch das 14. Kapitel.

\section{Simsons Heldenthat, Hochzeit und Räthsel}
\quellseite{388}
\section{Summarium}

Dies ist erstlich ein fein Exempel, wie Simson als ein gehorsam Kind, ob er wohl aus Gottes Schickung ein heidnisch Weib wollt nehmen, dennoch solches ohne seiner Eltern Vorwissen nicht thun will, V. 3.

Darnach ist eine andere Historia, wie Gott den Simson, weil er ein Fürst und Hauptmann sein sollte, der Israel aus der Philister Hand erretten sollte, mit sonderlicher Leibesstärke begnadet, daß er einen Löwen mit den Händen, ohne alle Waffen, zerreißet und erwürget, V. 6.

Zum dritten siehet man, weil Gott ihn erwecket hatte, daß er die Philister, seines Volks Israel ärgste Feinde, wohl plagen sollte, gibt ihm Gott rechtmäßige, gute und auch für der Welt genügsame Ursache, daß er sich an seinen Feinden, den Philistern, rächen kann. Einmal sie ihm und seinem Weibe dräuen, und darnach ihm das Weib nehmen, und einem andern geben, V. 20.

\section{Votum}

Der treue, barmherzige GOtt stärke alle Kinder an Leib und Seele, und regiere sie durch den heiligen Geist, daß sie ihren Eltern auch gehorsam und selige Werkzeuge sein seiner Gnaden, ihm zu Ehren und ihrem Nächsten zu Nutz und Dienst, Amen.

\chapter{Das 15. Capitel}

\section{Vorrede}

Wie über den Fürsten Simson der Geist des HErrn gerathen ist, und ihn getrieben und regieret hat: also verleihe Gott, daß wir auch Tempel und Wohnungen des heiligen Geistes sein und bleiben, der uns lehre, erinnere, stärke, leite und führe. Und wie Simson ein nützliches Werkzeug Gottes und seiner Kirchen gewesen: also helfe Gott, daß wir alle und ein jedes unter uns nach seinem Beruf und Stande ihm gehorsamlich und seiner Kirchen und derselben Gliedmaßen nützlich dienen. Wie auch Simson auf Gott mit Freuden und beständig getrauet hat und also gestorben ist: also gebe uns Gott, daß
\quellseite{389}
wir uns auch fest in wahrem Glauben an ihn halten, auf daß wir fröhlich und getrost in seinem Dienst beständig bleiben, die Trübsal uns nicht lassen abtreiben, sondern daß wir hie ritterlich ringen, durch Tod und Leben zu ihm dringen. Das verleihe uns allen der treue Gott, hochgelobet in Ewigkeit, Amen. Darauf höret ferner an das 15. Kapitel.

\section[Simson thut den Philistern durch die Füchse und]{Simson thut den Philistern durch die Füchse und den Eselkinnbacken Schaden}

a Schultern und Lenden) das ist, er schlug sie, wie sie ihm vorkamen, sie waren hohen oder niedrigen Standes.

b RamathLehi) heißt ein Hinwurf des Kinnbackens.

\section{Summarium}

Ist erstlich ein fein Exempel, wie Simson die Sachen läßt wohl an sich kommen, auf daß er desto mehr Rechts wider die Philister habe, und ihnen bezahle, was sie an den Juden verübet hatten, V. 3. Darnach siehet man, wie Vater und Tochter mit einander verbrennet werden, darum, daß der Vater die Tochter dem Simson genommen und einem andern gegeben
\quellseite{390}
hatte, V. 6. Denn solche Sünde will Gott nicht ungestraft lassen.

Zum dritten sind zwei herrliche Wunderzeichen hie, daß eine, wie Simson allein, ohne alle Waffen, tausend Mann mit einem Eselsbacken totschlägt, V. 15. Das andere, daß Gott in seiner Noth ihm einen Wasserquell läßt aus einem Backenzahn entspringen, V. 19. Solche wunderliche Historien sollen uns dahin weisen, daß wir uns Gott in allen Nöthen vertrauen und Hülfe bei ihm suchen sollen. Einmal wir sehen, daß er den Seinen wunderlich helfen kann.

\section{Votum}

Der König aller Könige, unser lieber Gott, lasse ihm alle christliche Regenten befohlen sein, und stärke und erhalte sie, daß sie auch ihm dienen und seine Kirche fördern und schützen, den armen Unterthanen zu Nutz und zu Erhaltung und Ausbreitung seines heiligen Worts, Amen.

\chapter{Das 16. Kapitel}

\section{Vorrede}

Kein Mensch ist so heilig, so verständig, so fromm, wird er sicher, und weicht Gott seine Hand von ihm, so führet und stürzt ihn der Teufel in die größte Sünde. Denn der Weg, auf dem wir allhie wandeln, ist sehr schlüpfrig, und die Anfechtung mancherlei, und der Teufel unruhig, dazu wir schwach und gar zu nachlässig. Diese menschliche Schwachheit läßt uns wohl bedenken, und derwegen wer sich läßt dünken, er stehe, der soll zuschauen, daß er nicht falle. Und wie der Herr Christus vermahnet Matth. 26, 41., so mag er wohl wachen und beten, daß er nicht in Anfechtung falle. Denn fromm sein und fromm bleiben stehet ja nicht in unserer Macht, und können es ja so bald versehen, als irgend ein anderer. O wie trefflich nöthig ist die Vermahnung und Warnung. O bester Gott, daß wir die täglich betrachten und nützlich gebrauchen, Amen. Dies soll auch erinnern das 16. Kapitel.

\section{Simsons Fall, Elend und letzte Rache}
\quellseite{391}
\section{Summarium}

Ist eine sonderliche Historia, in welcher wir sehen, wie auch heilige Leute in Sünden fallen und unrecht thun. Denn unrecht ist es, daß Simson mit den Heidinnen buhlet, V. 1. Wie wir noch sehen, daß große Herren in solchem Falle wollen mehr Freiheit haben, so doch das sechste Gebot sie eben so wohl bindet als andere Leute. Darum soll diese Historia niemand dahin deuten, daß er dergleichen auch thun möge; sondern wir sollen unsere Schwachheit an solchen Exempeln recht lernen erkennen, und desto mehr auf uns selbst Achtung haben, auch Gott mit größerem Fleiß bitten, daß er uns für Sünden und Anfechtung behüten wolle; sintemal so große Leute in grobe öffentliche Sünde gefallen, und greulich darum gestraft worden sind, wie wir weiter hören werden.

\section{Summarium}

Ist eine Historia, aus welcher wir lernen sollen, wie Gott die Hurerei an Simson gestraft hat, und seine Gnade, die er bisher gehabt, ihm entzogen, seine Augen ausgestochen, V. 21., und in einem elenden jämmerlichen Leben hat unter seinen Feinden eine zeitlang bleiben lassen. Denn wie der Psalm sagt, Gott will auch der Heiligen Sünde nicht ungestraft lassen. Nun soll man aber nicht denken (wie die Philister nicht anders wußten) daß Simson seine Stärke im Haar gehabt habe; sondern weil es Gottes Befehl war, Jud. 13, 5., daß er sollte ein Nasir, das ist, ein Verlobter Gottes sein, dem kein Schermesser auf sein Haupt sollte kommen, und doch ihm aus seinem eigenen Ansagen das Haar abgeschnitten ward, hatte er in dem Gottes Willen übertreten, und mußte also um dieser und anderer Sünde willen seiner Gnade, die er bisher gehabt, eine zeitlang beraubet sein, und mit der Stärke andern Leuten gleich werden.
\quellseite{392}
\section{Summarium}

Ist eine feine Historia, wie die Philister um ihrer Abgötterei willen gestraft werden, da sie sprachen, ihr Gott hätte ihnen ihren Feind in die Hände gegeben, V. 24. Solche Gotteslästerung müßte gestraft werden. Daß nun Simson sich so herzlich begehrt zu rächen, und Gott um Gnade der Stärke bittet, so er zuvor gehabt hatte, V. 28., solches soll man nicht dahin deuten, als wäre es erlaubt, sich selbst an seinen Feinden zu rächen. Denn Simson ist ehrlich ein Regent seines Volkes gewesen, darum hat er Recht dazu, daß er sein Volk an ihren Feinden rächte. Zum andern thut er solches nicht aus eigenem Vermögen, sondern sucht die Kraft bei Gott, und Gott läßt’s ihm auch durch Wunderwerk widerfahren.

\section{Votum}

Unser lieber GOtt und HErr regiere uns, daß wir uns für Unzucht, Abgötterei und allen andern Sünden hüten, und lindere uns die wohlverdienten Sündenstrafen, damit wir seine Gnade und Barmherzigkeit loben und preisen in alle Ewigkeit, Amen.

\chapter{Das 17. Kapitel}

\section{Vorrede}

Wenn sich die Leute für Gott, für seinem Zorn, auch für christlicher Obrigkeit Strafen nicht fürchten, so ist dem größten Haufen kein Laster und Uebelthat zu klein und geringe. Daher entstehen Abgötterei, Unzucht, Diebstahl, Mord und dergleichen Sünden. Derwegen jedermann Gott vor Augen haben soll, und die Unterthanen nicht allein gehorsam und gehorsam sein sollen ihrer Obrigkeit, sondern sollen sie auch in Ehren halten, und Gott danken, daß er Obrigkeit geordnet, die denen, so Landschaden sein, steuren, und den Frommen zu gut Gottes Wort, Friede, Zucht und Gerechtigkeit schützen und erhalten. Davon unterrichtet auch das 17. Capitel.

Micha richtet Götzendienst an.

a (Geschworen) diese Frau wird solches gelobet haben zum Bilde zu geben, das wird dem Sohn zuerst nicht gefallen haben, daß ihr darum gestohlen; darnach sich ihre Klagen lassen bewegen, und wieder gegeben, und lassen machen, was sie wollte. Da es nun wohl gerieht (wie solch Ding pflegt), ließ ers ihm auch gefallen.

b (Teraphim) deutschen wir das Heilige oder Heiligthum; wie zu unserer Zeit die hölzernen Heiligen, Todtenkränze, Todtenkleider und dergleichen heiliger Tand gehalten ist für Gottesdienst. Denn dies Wort hat unterschiedliche Bedeutung.

c Das ist, er weiset ihn; wie 2 Mos. 40, 13. stehet.
\quellseite{393}
\section{Summarium}

Ist eine Historia, wie sich im Stamm Ephraim eine neue Abgötterei so liederlich hat angefangen, und ein Levit sich dazu gefunden, der Armuth halben, um des Bauchs willen, zu solcher Abgötterei sich hat brauchen lassen, V. 6. Und meldet der Text sonderlich, daß solches zu der Zeit geschehen sei, da kein König noch Regiment im Lande gewesen sei, V. 7. Denn weltliche Obrigkeit soll fürnehmlich darauf sehen und halten, daß rechte Lehre und rechter Gottesdienst erhalten, und aller Abgötterei gewehret werde. Wo nun das Regiment gefallen ist, da muß die Kirche auch mit erhalten, und sinket der Satan Mittel und Wege, wie er Irrthum und Abgötterei anrichte. Derohalben jedermann desto fleißiger für fromme Regenten und ihr langes Leben bitten solle.

\section{Votum}

Der treue GOtt erhalte unsere liebe Obrigkeit in langem Leben und beständiger Gesundheit, und verleihe, daß ihre Regierung ferner glücklich und nützlich sei, auf daß sein Lob, Güte und Wohlthaten unter uns allzeit gerühmet und gepreiset werden, bis an den jüngsten Tag, Amen.

\chapter{Das 18. Kapitel}

\section{Vorrede}

Sirach im 15. Kapitel V. 13 spricht: Der HErr hasset alle Abgötterei, und wer ihn fürchtet, der scheuet sich dafür. Das Buch der Weisheit im 14. Kapitel V. 27. lehret: Den sündlichen Götzen dienen, ist alles Bösen Anfang, Ursache und Ende. Und St. Paulus sagt 1 Corinth. 6, 9.: Die Abgöttischen werden das Reich Gottes nicht ererben. Diese göttlichen Donnerschläge sollten uns ja genug sein von Abgötterei abzustehen. Aber da ein Lämmlein für dem Wolf erschrickt, und da ein Täublein für dem Geier sich fürchtet, so laufen die Menschen, ja die ausbündige Christen sein wollen, der Abgötterei zu und helfen die bestätigen. Derwegen wende sich ein jeder vom Unrechte und hasse mit Ernst die Abgötterei, und dagegen haltet fleißig an rechter Lehre und rechtem Gottesdienst, so wird es Gott wieder mit euch halten, zu zeitlicher und ewiger Wohlfahrt. Darauf höret auch davon an das 18. Kapitel.

\section{Die Daniter gewinnen Lais, und richten Abgötterei an}

a Sie hatten weder Fürste noch Herren.
\quellseite{394}
\section{Summarium}

Ist eine Historia, wie die Abgötterei aus dem Hause Micha unter den ganzen Stamm der Kinder Dan kommen sei, V. 13., auf daß wir an solchen Exempeln lernen, daß wir bei Leib der Abgötterei und falscher Lehre nicht lang zusehen, sondern je eher je besser dazu thun und ausrotten; denn beides, die falsche Lehre und Abgötterei, hat die Art, wie St. Paulus sagt, daß es kreucht, und frißt wie der Krebs, und immer viel Leute und Anhangs zu sich zeucht. Item, es durchdringet wie der Sauerteig, darum ists Noth, daß man alsobald wehre; und gleich wie das Feuer nicht lasse Luft gewinnen, sonst ists hernach ungewehret, wenns überhand nimmt.

\section{Votum}

Wir danken Gott herzlich, der uns von der antichristlichen Abgötterei erlöset hat; der verleihe, daß wir auch alle Abgötterei allzeit ernstlich meiden, damit wir uns und gemeinem Vaterlande nicht Gottes gerechte Strafen zuziehen, Amen.
\quellseite{395}
\chapter{Das 19. Kapitel}

\section{Vorrede}

Gott hat nächst dem Unglauben keine Sünde und Laster so hart und greulich gestraft, und weniger ungestraft wollen lassen hingehen, als Hurerei und Ehebruch. Was ist die Ursache, daß die Sündflut über die Welt und das Schwefel, Feuer und Pech über Sodom kam und regnet? Keine andere, denn die Sünde wider den heiligen Ehestand, nämlich Unzucht und Ehebrecherei. Warum wurden 24,000 Israeliten in der Wüste, alle Midianitische Weiber und hernach 25,000 Benjaminiter erwürget? Um der greulichen Unzucht willen. Diese und dergleichen erschreckliche Strafen Gottes sollen Junge und Alte von aller Unreinigkeit in und außerhalb dem Ehestande abschrecken. Davon unterrichtet auch das 19. Kapitel.

\section[Die zu Gibea begehen eine schändliche und höchst]{Die zu Gibea begehen eine schändliche und höchst sträfliche That}
\quellseite{396}
\section{Summarium}

Ist eine schreckliche und greuliche Historia von den Gibeonitern, wie sie den Leviten und sein Weib nicht beherberget, sondern ihr alle Schand und Schmach angelegt haben, V. 22. Solche Exempel sollen uns zur Furcht Gottes und fleißigem Gebet vermahnen, daß uns Gott für Sünden und Schanden behüten und in Anfechtung nicht wolle sinken lassen. Denn wo Gott nicht hält und behütet, ist's sehr bald geschehen, daß man in Schand und Laster fällt; wie wir an Simson, David und andern Heiligen sehen. Darum Christus der HErr im Vater unser uns beten gelehret: Führe uns nicht in Versuchung.

\section{Votum}

Verleihe der treue, liebe GOtt, daß wir folgen; dem Teufel und unsers Fleisches bösen Lüsten, auch der Welt Anreizung Widerstand thun und in Heiligkeit und Reinigkeit allzeit leben und auch sterben, Amen.

\chapter{Das 20. Kapitel}

\section{Vorrede}

Gott läßt die Unzucht nicht ungestraft, auch an denen, die Buße thun und sich nach begangener Unzucht bekehren. Warum lassen sich denn heut so viele unter Jungen und Alten, in und außerhalb dem Ehestande, den unsaubern Geist so verblenden und verführen? Demgegen wollen wir beten, Gott der Allmächtige, der ein reines, keusches Wesen ist, und auch Reinigkeit und Keuschheit von uns haben will, wolle durch sein Wort und durch seinen heiligen Geist unsere Herzen reinigen von allen unreinen Worten, Werken und Gedanken, und daß er uns gebe, daß wir ihm mit reinem Herzen dienen, ihn anrufen und ihm danken, hie zeitlich und dort ewiglich. Denn selig sind, die da reines Herzens sind, denn sie werden Gott schauen, sagt Christus Matth. 5, 8. Das soll auch erinnern das 20. Kapitel.
\quellseite{397}
a Non petunt victoriam, sed viribus praesumunt, et justitia sua confidunt, ac animantur malitia filiorum Benjamin. D. i. sie bitten nicht um Sieg, sondern verlassen sich auf ihre Kraft und vertrauen auf ihre Gerechtigkeit, und werden muthig durch die Bosheit der Kinder Benjamin.

b Weil die Kinder Israel auch wohl diesmal nicht um Sieg den Herrn bitten, so verheißet er ihnen auch den Sieg nicht.

\section{Summarium}

Dies ist eine Historia, in welcher man siehet, wie oft der schuldige und ungerechte Haufe zu einer bösen Sache Glück hat, wider den andern Theil, der doch der Sachen Recht hat, aber sonst bußwürdig ist. Also hatten die Benjamiter gröblich Unrecht, daß sie die greuliche Sünde der Gibeoniten nicht allein nicht strafen, sondern die Uebelthäter auch vertheidigen lassen. Wiederum die andern elf Stämme thaten recht daran, daß sie die greuliche Sünde nicht wollen ungestraft lassen; und dennoch liegt Benjamin zweimal ob, und die andern elf Stämme verlieren vierzig tausend Mann in einer solchen
\quellseite{398}
guten Sache, 20, 24. Warum gehet es doch so ungleich zu? Sind doch die Gibeaen und Benjaminiten sträflich. Die elf Stämme aber haben guten Fug und Recht, und liegen dennoch unten. Woher kömmt doch solches? Antwort: Obgleich die Sache an sich selber nicht böse war, so sind doch die elf Stämme in andern Fällen sträflich gewesen, nämlich aufs erste in dem, daß sie den Sieg von Gott nicht baten, sondern sich verließen auf beides, ihre Macht und gute Sache, und daß ihr Gegenteil unrecht hatte. Darnach auch, daß sie, wie oben stehet, öffentlich Abgötterei lang her getrieben und geduldet hatten, daher ist solches Urtheil über sie gangen.

c Ich, spricht Gott, will es thun. Ziehet habt ihrs wollen thun; aber es heißet, ich, nicht ihr.

d Der Herr schlug sie, nicht Israel, denn es heißt: gloria in excelsis Deo.
\quellseite{399}
\section{Summarium}

Sie ist eine Historia, in welcher man stehet, obschon Gott bisweilen bösen Buben etwas lang zusiehet, daß er dennoch endlich mit der Strafe kommet und so viel desto härter zuschlägt. Sonderlich aber soll man nicht vergessen, daß die greuliche Unzucht, welche die Gibeoniter begangen hatten, und die Benjamiten vertheidigten, so jämmerlich und hart gestrafet wird, daß alle Städte in Benjamin verheeret, und der ganze Stamm fast gar ausgerottet ward. Denn Hurerei und Unzucht ist auch der Sünden eine, welche Gott mit Verwüstung von Land und Leuten strafen will, wo weltliche Obrigkeit nicht strafet; wie man hie und an andern Orten siehet. Darum jedermann sich mit höchstem Fleiß dafür hüten, und weltliche Obrigkeit mit harter Strafe solchem Aergerniß wehren soll.

\section{Votum}

Solches helfe der treue Gott uns und aller Obrigkeit, daß es durch seine Gnade und Kraft des heiligen Geistes geschehen möge, zur Abwendung aller Strafen, zeitlich und ewiglich, Amen.

\chapter{Das 21. Kapitel}

\section{Vorrede}

Wie eine Biene, wenn sie aus Zorn den Menschen sticht, darüber ihren Stachel verlieret, keinen Honig mehr machen kann, sondern ein steter Knecht sein, und Wasser zutragen muß: also wirket des Menschen Zorn nichts Gutes. Da vergehet alle Freude und Lust etwas Gutes zu thun, und wird der zornige Mensch ein Knecht des Teufels und der Sünden, schadet seinem Nächsten so sehr nicht als ihm selbst. Und wie ein Wasser sich selbst versiegen und endlich gar verraucht, also verlieret ein Zorniger auch die Gnade des heiligen Geistes. So ist der Zorn ihnen eine Thür und Schrotleitern zu allerlei Untugend und Schandlustern. Derwegen ein jeder Christ für ungerechtem Zorn sich fleißig hüten soll. Davon höret auch an das 21. Kapitel.

\section{Wie der Stamm Benjamin wieder erlaubet werden}
\quellseite{400}
\section{Summarium}

Dies ist eine Historia, wie die elf Stämme die übrigen vom Stamm Benjamin wieder zu Gnaden annehmen, und sie in das Ihrige einsetzen, ihnen helfen und rathen. Denn es soll mit der Strafe und dem Zorn ein Maß haben, daß man ihm nicht zu viel thue.

\part{Das Buch Ruth}

\chapter{Das 1. Kapitel}

\section{Vorrede}

Dies Buch Ruth, neben dem, daß es eine feine, liebliche und sonderlich zum Hausregiment nützliche Historia ist, dienet auch dazu, daß wir lernen, wie Gott sich zur selbigen Zeit auch der Heiden hat angenommen, als die hernach zu seinem Reich durchs Evangelium auch kommen, und des Samens Abrahams genießen sollten. Und ist sonderlich hie wohl zu merken, daß diese Ruth, eine Heidin aus der Moabiter Land, unter den Voraltern des HErrn Christi gezählet wird, wie auch Matth. 1. v. 5. zu sehen ist.

Im ersten Kapitel soll man lernen, wie eine Schnur sich gegen ihre Schwieger freundlich, demütig und dienstlich halten soll. Denn solches gefällt Gott wohl, will es auch reichlich vergelten, wie er an der frommen Ruth bewiesen hat, und wir darnach weiter hören werden.

\section{Ruth reiset mit Naemi nach Bethlehem}
\quellseite{401}
a) Naemi) heißt meine Lust.

b) Mara) heißt bitter oder betrübt.

\section{Summarium}

Bei diesem Kapitel sollen wir merken drei Stücke. Das erste ist von der Theurung im jüdischen Lande, welche Gott über Israel hat ergehen lassen von wegen ihrer mannigfältigen Sünden; denn ein jeder gelebet hat, wie es ihm gelüstet, weil sie etliche Zeit keinen Fürsten und Regenten gehabt hatten. Also straft Gott noch die Sünder der Welt, wie mit andern Plagen, also mit Theurung, wie denn Gott zum oftermal gedräuet hat, er wolle über die Sünder seinen Himmel zuschließen. Das andere Stück ist, daß Elimelech mit den Seinen kehret, daß man bei Abgöttischen und Gottlosen also wohnen soll, daß solche Beiwohnung uns gleichwohl an unserm Glauben und Religion nicht schädlich sei. Aber zu unsern Zeiten werden viele Leute gefunden, welche die Religion mit Veränderung ihrer Wohnung auch verändern. Daher etliche das Sprüchwort haben: Darum ich wohne, darnach glaube ich. Das dritte Stück ist, Naemi heißt ihre beiden Schnuren wiederum heimkehren, damit sie bei den Juden, als bei einem fremden Volke, nicht in Gefahr des Lebens und der Ehren kommen. Arpa thut's, Ruth aber bleibt bei ihrer Schwieger. Das ist ein schön Exempel, daß die Schwiegerfrauen und Schnuren mit einander treulich haus
\quellseite{402}
beln, und einander herzlich lieben, folgen und dienen sollen. Und daß Naemi ihre Schwuren segnet, und wünschet ihnen Gutes, das denn die Ruth hernach bekommen hat, lehret, daß der Eltern Segen über ihre ungehorsame Kinder nicht vergeblich ist. Wie denn auch Sirach saget Cap. 3. V. 10.: Der Segen der Eltern erbauet die Häuser der Kinder.

\section{Votum}

Der treue GOtt lindere uns die wohlverdienten Strafen, und erhalte uns bei der christlichen Religion und in herzlicher Liebe gegen Ihm, gegen den Unsern und gegen unsern Nächsten, auf daß wir einander gern und treulich dienen, Amen.

\chapter{Das 2. Capitel}

\section{Vorrede}

Wer etwa ein Fremdling ist, und wollte gern Gnade und Gunst bei fremden Leuten haben, der besiehe sich, daß er durch den Glauben einen gnädigen Gott habe, und daß er christlich lebe; so wirds Gott also schicken und machen, daß er mehr Liebe und Huld bei Fremden finden wird, denn bei seinen Freunden und Verwandten. Also wer dem Nächsten Gutes thut, ob er wohl das Guts wünschen und Danksagen der Leute dafür nicht höret, so schreibet doch das Werk an ihm selbst, da es aus dem Glauben geschieht, zu Gott. Denn der Glaube ist gleich als ein Wagen, dadurch alle unsere Werke für Gottes Angesicht geführet werden, der sie auch belohnet, wie Sirach sagt im 17. Capitel V. 18.: Gott behält die Wohltat der Menschen wie einen Siegelring, und die guten Werke wie einen Augapfel. Dies sind Exempel Ruth und Boas. Davon das 2. Kapitel.

\section{Ruth lieset auf Boas Felde Aehren auf}

a Das ist, sie ist nicht der Wehen eine, die daheim auf dem Polster sitzen und faullenzen.
\quellseite{403}
\section{Summarium}

Dies ist ein fein Exempel, erstlich an Ruth, wie sie Gott der Treue, so sie an ihrer Schwieger gethan, genießen läßt, V. 3. Darum muß folgen, wo Schnur und Schwieger, wie gemeiniglich geschieht, unter einander unfreundlich und undienstlich sind, daß Gott übel gefalle, und darum armselige, elende Leute müssen sein.

Zum andern ist ein sehr fein Exempel an Boas, der seine arme Freundin oder Schwägerin nicht verachtet, sondern ist ihr freundlich, und thut ihr alles Gutes, V. 8. Denn Gott gibt oft manchem viel Segens und Guts, allein, daß er seinen armen Freunden damit helfen, und sie erfreuen und trösten soll. Wo mans aber nicht thut, und sich armer Freunde nicht annimmt, und ihnen nicht gern hilft, solche Untugend wird Gott zu seiner Zeit wohl finden und strafen.

\section{Votum}

Verleihe GOtt, daß Freundschaft also hie gegen einander lebe, daß sie auch in ewiger Freundschaft beisammen leben; und helfe allen, tröste und fördere auch alle, um Christi JEsu willen, die armen und verlassenen Christen mit Rath und That aushelfen, Amen.

\chapter{Das 3. Kapitel}

\section{Vorrede}

Um den Ehestand ist es ein langer Kauf, darinnen erneuern und hinter sich sehen nicht gilt, noch hilft. Darum mag man wohl langsam, fürsichtig und weißlich damit umgehen, und den Ehestand in aller Furcht Gottes, und mit wohlbedachtem Muth fürnehmen. Denn er ist nicht ein Pferdekauf, sondern was einmal geschlossen wird, dabei muß man folgend beruhen. Dies sollen wohl betrachten beides Junggesellen und Jungfrauen. Denn es heißt: Ehe wiegs, denn wags. Und ein tugendsam Eheweib und frommer Mann kömmt und wird gegeben von dem HErrn, der ist Gottes Gabe und Segen, wie Salomon in seinen Sprüchen im 19. Kapitel V. 14. lehret. Und derwegen mögen junge Leute wohl beten, Augen und Ohren aufthun, sehen, hören und bedenken, was sie thun, damit sie es nicht hernach gereue. Ferner höret darauf an das 3. Kapitel.

\section{Ruth bekommt wegen künftiger Heirath guten Bescheid}
\quellseite{404}
\section{Summarium}

Ist eine feine Historia, wie die arme Ruth aus Befehl ihrer Schwieger ihren nächsten Freund Boas um die Ehe in aller Zucht und Demuth anspricht, V. 6. Und Boas fürnehmlich bedenket, daß sie sich züchtig und wohl gehalten, ein gutes Geschrei bei jedermann erhalten, gottesfürchtig und fromm gewesen sei, V. 12. Darum sollen hie beide Jungfrauen und junge Wittwen lernen, daß sie gottesfürchtig, züchtig und fromm sein, so werden sie es genießen. Desgleichen die jungen Gesellen hie sollen auch ein Exempel nehmen an Boas, daß sie, wenn sie freien wollen, mehr auf Gottesfurcht, Zucht und Frömmigkeit, denn auf Geld und Gut sehen sollen. Weil aber solches selten geschieht und wenig vermaßen in den Ehestand sich schicken, siehet man, daß desto mehr Unglücks und Unraths unter den Leuten ist.

\section{Votum}

Der Stifter des heiligen Ehstandes, unser lieber GOtt, helfe, daß junge Leute in diesen Orden in seiner Furcht, mit vorhergehendem herzlichen Gebete und mit gutem Bedacht treten, und daß sie darinnen christlich leben, und Segen an Leib, Seele und Nahrung haben und behalten, Amen.

\chapter{Das 4. Kapitel}

\section{Vorrede}

Jetzt sind die Leute gemeiniglich so geschickt, daß sie heimlich, ohne andrer Leute Beisein, auch mit Geldbuhr zusagen, das verleugnen sie hernach öffentlich. Also Geizhälse suchen allein ihrem Frommen und Nutzen nach, und lassen andern Leuten Schaden. Ferner, daß Gott Kinder bescheret, dieser Segen wird jetzt schier vor einen Fluch gehalten, welches daher kömmt, daß ihrer viel die Menge der Kinder nicht ansehen nach Gottes Wort durch den Glauben, sondern nach ihren Gedanken auf die Beschwerungen, die sich durch die Kinder und Kinderzucht zutragen. Anders lehret davon das 4. und letzte Kapitel.

\section{Die Hyrath Boas mit Ruth wird vollzogen und gesegnet}

a Wo du hin willst, denn das gebräuchliche Ploni almoni ist ein Name eines ungewissen Orts.
\quellseite{405}
b) Id est, Det Deus, ut cum illa magnificeris: ut certe factum est. Nam peperit Obed, avum Davidis. Das ist, (Gebe Gott, daß du mit ihr groß werdest; wie auch wirklich geschehen ist, denn sie hat den Obed geboren, den Großvater Davids.

\section{Summarium}

Dies ist eine Historia, wie Boas die arme Wittwe Ruth, seine Freundin, zur Ehe nimmt, V. 10. Und Gott den frommen und gottesfürchtigen Leuten einen Sohn gibt, welcher Davids Großvater gewesen ist; auf daß wir lernen, wie Gott Zucht und Frömmigkeit lieb hat und segnet, dagegen unzüchtige, furchtlose Leute weder Glück noch Heil haben, oder es zu ihrem großen Schaden und Nachtheil haben.

\section{Votum}

Gott, der HErr, verleihe, daß alle Eheleute auch in Gottes Furcht und in Zucht leben, und gebe ihnen seinen reichen Segen an Kindern, Gut und Nahrung, doch nach seinem Willen, und wie es den Eheleuten gut und selig sein soll, Amen.

\quellseite{406}

\part{Das erste Buch Samuels}

\chapter{Das 1. Kapitel}

\section{Vorrede}

Heute fangen wir an das erste Buch Samuels. In diesem Buch aber werden fürnehmlich gemeldet Exempel, wie Gott auch an gotterleuchteten und großen Leuten die Sünde gestraft, und dagegen, wie er auch an gemeinen und geringen Menschen seine Gnade bewiesen hat. Diese Exempel sollen wir zu täglicher Warnung und Trost fleißig anhören und betrachten. Denn hat Gott große, hochgeachtete Leute gestrafet; fürwahr, er wird unser auch nicht verschonen, da wir sündigen. Also hat Gott unansehnliche Leute für der Welt gesegnet; um unserer Gebrechlichkeit willen wird er unser mit seinem Segen traun auch nicht vergessen, da wir in wahrer Buße leben. Derwegen wollen wir uns alle Augenblicke für Sünden hüten, daß wir nicht straffschuldig werden, und wollen dagegen alle Tage christlich leben, daß uns Gott auch segne und behüte. O HErr JEsu Christ, das hilf uns allen, Amen. Darauf höret ferner an das 1. Kapitel.

\section{Samuel wird geboren}
\quellseite{407}

\section{Summarium}

Ist eine Historia, wie Hanna, als eine gottesfürchtige Frau, Gott um einen Sohn bittet, und Gott sie ihrer Bitte gewähret, V. 20. Auf daß wir erstlich lernen sollen, daß die Kinder eine Gabe Gottes sind, wie der 127. Psalm V. 4. auch lehret. Danach, daß Gott Kinder desto lieber geben will, wo man sie Gott zu Ehren und den Leuten zu Nutz begehrt aufzuziehen; wie denn der heiligen Hanna Herz und Gedanken gestanden, V. 11.

Solchem Fürnehmen der lieben Hanna sollen die Mütter auch folgen, ihre Kinder zu Gottesfurcht und aller Zucht halten. Denn dasselbe heißt nach unserer Sprache, dem HErrn die Kinder geben, nicht in die Klöster stoßen, wie im Papstthum die Gewohnheit war.

\section{Votum}

Helfe GOtt, der Vater aller Gnaden, daß alle Mütter solcher heilsamer Vermahnung treulich folgen, und regiere alle lieben Kinder, daß sie sich ziehen lassen und alle wohl gerathen, GOtt zu Lob und ihrem Nächsten zu Nutz, Amen.

\chapter{Das 2. Kapitel}

\section{Vorrede}

Wir sind Pflanzen Gottes, ihm zu Lob und Preis, sagt Jesaias im 61. Cap. V. 3. Wer aber Gott ehret, den will Gott wieder ehren. Ueber das alle Tage, alle Stunden, ja alle Augenblicke allerlei Gutes an Leib und Seele erzeigt uns Gott so reichlich, daß wir mit dem Patriarchen Jakob bekennen müssen: Wir sind viel zu geringe und unwürdig gegen Gottes große Barmherzigkeit, 1 Mos. 32, 10. Derwegen wir Gott für seine Wohlthaten danken sollen mit Herz und Munde, und auch mit der That, also daß ein jeder in seinem Beruf, als in der Kinderzucht und dergleichen sich christlich verhalte, das Böse lasse, das Gute thue, das Kreuz trage und das Unrecht überwinde, und alles im Glauben und Geduld. Dies erinnert auch das 2. Kapitel.

\section{Hannas Lobgesang. Auferziehung Samuels. Bosheit der Söhne Eli}

a Alte) das Feste, Gewisse, Ehrliche, wie man spricht: Gewohnheit, alt Herkommen, Landessitten und Weise; denn darauf trotzen die Leute, und sagen: Ei sicher, die alte Weise die beste, unsere Vorfahren sind auch nicht Narren gewesen. Und dennoch also wider Gottes Wort, als müßte es nicht anders noch newern.

b Das ist, sie müssen ums Brod dienen.
\quellseite{408}

\section{Summarium}

Ist ein Danklied, in welchem die heilige Frau Hanna Gott lobet, daß er allein der rechte HErr und Gott sei, der alles nach seinem göttlichen Willen könne machen, V. 3. Darum vermahnet sie, V. 6., daß man allein auf solchen Gott trauen, denselben allein fürchten und seine Hülfe allein suchen, und derselben auch mit Geduld erwarten soll. Die es aber nicht thun, sondern ihren Trost auf Geld, Gut und Gewalt setzen, wie die Welt pflegt, die müssen zu Boden gehen.

Am Ende gedenkt sie, wie Gott die Welt richten werde, V. 11., das ist, das Evangelium in aller Welt predigen lassen, und seinem Gesalbten, das ist, unserm HErrn Jesu Christo, Macht über alles geben. Mit solchem Danklied stimmet überein der Gesang der heiligen Jungfrau Maria, und ist zu beiden Theilen einerlei Meinung und schier auch einerlei Wort.

(Leibrock) das waren priesterliche Kleider, davon 2 Mos. 28, 6.

d Si Deus offenditur, et non ipse per se remiserit, non est aliquis alius vel superior, qui inter ipsum offensum et offensorem mediare possit, sicut inter homines. D. h. wenn Gott beleidiget wird, und nicht selbst durch sich selbst vergiebt, so ist kein Anderer oder Höherer, der zwischen ihm, dem Beleidigten und dem Beleidiger vermitteln könnte, wie unter Menschen.
\quellseite{409}

e (Lückest du) gleichwie Apost. 9, 5. St. Paulus wider den Stachel löcket, das ist frech und muthwillig.

f (Altar) Non auferam quidem de altari meo, sed diu non vivent. Quia ubi adoleverint, morientur.

\section{Summarium}

Ist eine Historia von dem Muthwillen und Unzucht der Söhne Eli, des Hohenpriesters, welche Gott also straft, daß, wie folgen wird, nicht allein sie wegen auf einen Tag umkommen, V. 34.; sondern Gott dräuet auch, daß ihr ganz Geschlecht soll solcher Sünde entgelten, daß wenn sie zu Mannesalter kommen, sterben sollen, und also ihrer Nachkommen keiner zum vollkommenen Alter kommen soll, V. 35. Auf daß die Eltern ja lernen ihre Kinder mit Fleiß in Gottesfurcht aufziehen, und ihnen ihren Muthwillen nicht gestatten, weil es so hart und greulich soll gestraft werden, wo sie hierin lässig sind.

\section{Votum}

Regiere uns der treue GOtt, daß wir ihm für seine Wohlthaten dankbar sein, und segne die Kinderzucht, auf daß sie bei der lieben Jugend wohl gerathe, damit sie in Verstand zunehme und in Tugenden lebe, Gott zu Ehren, ihnen selbst und andern Leuten zu Nutz, Amen.

\chapter{Das 3. Kapitel}

\section{Vorrede}

Salomon in seinen Sprüchen im 29. Kapitel V. 15. sagt: Ruthe und Strafe gibt Weisheit; aber ein Knabe ihm selbst gelassen, schändet seine Mutter. Züchtige deinen Sohn, so wird er dich ergötzen, und wird deiner Seele sanft thun. Ferner Sirach im 26. Kapitel V. 13. spricht: Ist deine Tochter nicht schamhaftig, so halt sie hart, auf daß sie nicht ihren Muthwillen treibe, wenn sie so frei ist. Wenn du merkest, daß sie frech um sich stehet, so siehe wohl darauf; wo nicht, und sie thut darüber wider dich, so laß dich es auch nicht wundern. Solche Kinderstrafe aber sollen die Eltern also mäßigen, daß die Kinder selbst merken und spüren mögen, daß ihre lieben Eltern sie nicht aus Haß und Feindschaft gegen ihre Personen, sondern aus lauter Liebe strafen, ihre Sünden anfeuren und sie gern mit ihrer Züchtigung davor verwahren wollten, daß sie der irdische oder geistliche Scharfrichter (und alsdann der höllische Henker und Peiniger) nicht strafen dürfen. Dieses erinnert auch das 3. Kapitel.

\section{Samuel wird berufen und ansehnlich gemacht}

a Theuer: es waren nicht Prediger noch Pfarrherrn genug, die Bibel ist da unter der Bank gelegen, hat niemand studiret; Samuel ist kommen, und hats wieder hervor gebracht.
\quellseite{410}
ᵇ Id est, Deus coepit apparere copioso verbo sub Samuele.

\section{Summarium}

Ist eine Historia, wie Gott sich dem Knaben Samuel offenbaret, V. 11. Und zum andernmal dem Hohenpriester Eli durch ihn dräuet, er wollt seiner Söhne Sünden strafen, und er, der Vater selbst, soll auch darum gestraft werden, daß er gewußt hat, daß seine Söhne sich so schändlich halten, und dennoch nicht einmal dazu sauer gesehen, V. 12. Solches sollen die Eltern mit Fleiß merken. Eli hatte ja seine Söhne gestraft, wie oben im andern Capitel V. 24. stehet; aber Gott will sich nicht daran genügen lassen, wenn man allein mit Worten will strafen, sondern er will, daß man mit Ernst dazu thun und mit harter Strafe der Sünden wehren soll, und dieselben in keinem Weg an den Kindern dulden. Das that Eli nicht, und sagte seinen Söhnen ein Wort, zwei, und ließ es dabei bleiben, brauchte keinen sondern Ernst wider sie, ließ sie ihm zu lieb sein. Das war unrecht, darum mußten Vater und Söhne ihre Strafe darum leiden.

\section{Votum}

HElfe GOtt, der seinen Willen will erfüllet haben, daß die Eltern in der Kinderzucht gebührlichen Glimpf und nöthigen Ernst gebrauchen, und regiere die Kinder, daß sie sich ziehen lassen und gehorsam sein und bleiben allzeit, Amen.
\quellseite{411}

\chapter{Das 4. Kapitel}

\section{Vorrede}

Der meiste Theil der Menschen gehet dahin, und bedenket nicht einmal, was Sünde sei und wie großen Schaden sie mit sich bringet. Die Jugend achtets nicht aus Unverstand und Nachlässigkeit, so wollen es die Alten nicht wissen noch erkennen, vor großer Klugheit und Sicherheit. Derwegen erzählet die heil. Schrift die schrecklichen Sündenstrafen, in viel und mancherlei Exempeln, jedermann zur Warnung. So läßt uns solche anhören, daß ein jeder vor Gottes Zorn erschrecke, Buße thue, die Sünde meide, und dagegen christlich lebe, daß Gott nicht auch also wegen der Sünde uns züchtigen dürfe, wie er andere Land und Leute derwegen grausam gestrafet hat. Darauf höret auch an das 4. Kapitel.

\section[Israels Niederlage. Entführung der Bundeslade. Todesfall Eli und]{Israels Niederlage. Entführung der Bundeslade. Todesfall Eli und seiner Söhne}

a Eben(Ezer) heißt Helfstein. f. unten Kap. 7, 12.
\quellseite{412}

\section{Summarium}

In eine Historia, in welcher man sehet, wie Gott endlich die Sünde hart und zugleich strafet. Und wird oben gemeldet, wie die Leute das Opfer des HErrn lästerten; item, wie Gottes Wort seltsam und wenig Weissagung gewesen sey. Die Sünden nun, so hernach gefolget sind, strafet Gott also, daß die Philister vier und dreißig tausend Mann erschlagen, V. 10., und die Lade des Bundes auch wegnehmen. Die zween gottlosen Söhne des Priesters Eli kommen auch um, V. 11. Der Vater Eli, da er solches höret, erschrickt also, daß er zurückfällt und den Hals bricht, V. 18. Solches sind eitel Strafen der Sünde. Sonderlich aber ist zu merken, weil die Kinder Israel nicht Buße thaten, half sie es nichts, daß sie die Lade des Bundes bei ihnen hatten. Eben wie es uns wider den Türken auch nicht hilft, daß wir Christen heißen, getauft sind, Gottes Wort und das Sacrament haben, sintemal wir unbußfertig sind, und von unserm gottlosen Wesen nicht ablassen.

\section{Votum}

GOtt erbarme dich unser, lindere die wohlverdienten Sündenstrafen; verleihe auch, daß wir in seinem Gehorsam leben, und führe uns aus dem Sündenreich dieser Welt in das Gnadenreich der Seligkeit, um JEsu Christi unsers HErrn willen, Amen.

\chapter{Das 5. Kapitel}

\section{Vorrede}

Was den Menschen unmöglich und schwer ist, das alles ist Gott möglich und leicht, und Gott kann in aller Noth erretten alle, die im wahren Glauben zu ihm treten und um Beistand ihn ansprechen. Dergleichen in schweren Sachen und großer Gefahr niemand zagen oder verzweifeln soll; sondern jedermann habe Zuflucht zu Gott, und suche durchs Gebet bei ihm Hülfe, er wird die Seinen wider ihrer Feinde sichtbare und unsichtbare List und Gewalt wohl behüten und erhalten. Diesen Trost lehret Gott selbst Jes. 43, 2. und sagt: So du durch Wasser gehest, will ich bei dir sein, daß dich die Ströme nicht sollen ersäufen, und so du ins Feuer gehest, sollst du nicht brennen, und die Flamme soll dich nicht anzünden; denn ich bin der HErr, dein Gott, der Heilige in Israel, dein Heiland. Ferner höret an das 5. Kapitel.

\section{Plage der Philister wegen der entführten Bundeslade}

a Wenn Menschen und Vernunft obliegt, so gehet Gottes Wort und alle Ehre dahin; da fället das rechte Priesteramt zurück und stirbet, und das alles aus Gottes Zorn. Aber die so gewinnen, haben darnach keine Ruhe im Gewissen. Denn wo Gottes Wort nicht recht gehet, thut es dem Gewissen alles Leid an, wie hie die Lade Gottes den Philistern, so lange, bis sie zuletzt ihre Schande erkennen müssen, daß sie Gottes Wort verachtet haben und mit Ehren wieder zurecht bringen müssen. Das ist hie bedeutet mit den äußeren Zeichen und Mäusen, welches nichts ist denn die heimliche Plage der Gewissen, die zuletzt offenbar wird durch Gottes Wort, wie St. Paulus sagt 2 Tim. 3, 9.: Ihre Thorheit wird offenbar werden jedermann.
\quellseite{413}

\section{Summarium}

Dies ist eine schöne Historia, wie es den Philistern geraten sei, daß sie die Lade des Bundes genommen haben, nämlich, daß sie Gott also mit Krankheiten und Sterben strafet, daß sie so froh werden, da sie ihrer wieder los wurden, V. 6. Solche Historia reimet sich fein mit den Ketzern, so Gottes Wort den Christen nehmen und zu eigener Ehre, um des Bauches willen, mißbrauchen. Aber endlich strafet sie Gott mit heimlichen Krankheiten, das ist, mit bösem Gewissen und Verzweiflung; das soll ihr Lohn sein.

\section{Votum}

Der barmherzige, starke GOtt steure allen Rotten und Secten, die sein Wort verkehren und seine Kirche zerstören; und behüte uns für Land-, Stadt- und Kirchen-Krieg und Empörungen, um seines Namens Ehre und der aufwachßenden Jugend Wohlfahrt willen, Amen.

\chapter{Das 6. Capitel}

\section{Vorrede}

Der wahre Gottesdienst ist, recht erkennen aus Gottes Wort das Wesen und den Willen Gottes. Das Wesen also, daß nur ein Gott ist, aber drei verschiedene Personen sind. Den Willen aber, daß es der ist, daß wir an den Sohn glauben sollen, und durch den Glauben an ihn erlangen Vergebung der Sünden, Gerechtigkeit, den heil. Geist und ewiges Leben, und diesen wahren Gott, der also erkannt wird, lieben und fürchten, auf ihn hoffen, ihm gehorsam sein im Glück und Unglück, ihn anrufen, ihm danken, seinen Namen predigen, Gottes Wort gern und fleißig hören, das Sacrament oft gebrauchen, den Eltern, der Obrigkeit, den Zuchtmeistern, Herren und Frauen Gehorsam leisten, barmherzig, keusch, züchtig, auch mäßig sein, gerecht und wohlthätig gegen jedermann, vornehmlich aber gegen Witwen und Waisen, wahrhaftig und aufrichtig, und sich unbefleckt von der Welt erhalten. Denn der wahre Gottesdienst begreift in sich alles, was in Gottes Wort, im Gesetz und Evangelio, geboten und befohlen ist. Und diesen Gottesdienst sollen wir alle in gebührlicher Ehre halten, auch christlich verrichten. Als uns denn dieß erinnert das 6. Kapitel.

\section{Die Bundeslade wird den Israeliten wieder zugestellet}
\quellseite{414}

\section{Summarium}

Ist ein fein Exempel, daß die Plage bei den Philistern nicht eher aufhöret, bis sie durch ein Schuldopfer ein öffentlich Bekenntniß über sich geben, daß sie haben Unrecht gethan und sich an dem Gott Israel versündiget, V. 4. Denn soll die Strafe aufhören, so muß man die Sünde erkennen und davon ablassen.

Das aber fünfzig tausend und siebenzig Mann zu BethSemes sterben, darum, daß sie die Lade des HErrn gesehen hatten, V. 10., solches ist nicht darum geschehen, daß sie die Lade schlecht haben angesehen; sondern die Heiden haben ein Frohlocken darüber gehabt, und gedacht, was ohnmächtigen Gott haben doch die Juden, daß er sich und sie nicht rettet? Wir haben das größte Stück von ihrem Gottesdienst und Heiligthum. So es ein rechter Gottesdienst wäre, oder ihr Gott etwas wäre und vermöchte, so hätten sie uns abgeleget. Aber da sehen wir, Gott und das Volk ist eins wie das andere. Werden also einen sondern Stolz gefasset, Gott gespottet und gehöhnet haben, wie solches der Teufel sehr wohl kann. Da läßt sich Gott sehen, mit seinem Zorn, daß ihr fünfzig tausend und siebenzig auf einmal dahin fallen und sterben, und die andern des Spottens aufhören, und wissen nicht, wo sie vor Angst bleiben sollen. Und zum Exempel, daß wir mit Gott, seinem Wort und dem rechten Gottesdienst nicht sollen leichtfertig umgehen, sondern gottesfürchtig sein und in allen Ehren haben.

\section{Votum}

GOtt, der heilige Geist, erleuchte, lehre und führe uns, daß uns der höllische Geist in Verachtung GOttes Worts und der Gottesdienste nicht verführe; und helfe, daß wir in täglicher, christlicher Buße leben, damit uns GOtt unserer Sünde wegen nicht strafen dürfe, Amen.

\chapter{Das 7. Kapitel}

\section{Vorrede}

Rechte Lehrer und gehorsame Zuhörer sind einer jeden Stadt schönster Schmuck und Zierde; und eine solche Gemeinde ist Gottes Paradies, der heiligen Engel Lusthaus und eine Schule und Zuchtspiegel der Jugend. Derwegen danken wir Gott, daß wir sein Wort rein und lauter haben, und daß dennoch dessen auch Hörer und Thäter seien. Der helfe mit Gnaden ferner; was unsere Prediger mit Worten treulich lehren,
\quellseite{415}
daß wir solches mit guten Werken fleißig verbringen, und also leben, daß Gott geprieset, die heiligen Engel erfreuet, die Jugend gebessert und wir von Gott gesegnet werden, und daß wir alle in der Himmelsburg der heiligen Engel und im Lande der Lebendigen und Auserwählten Gott loben und preisen, immer und ewiglich, Amen. Und darauf höret an das 7. Capitel.

\section[Die bußfertigen Israeliten sagen auf Samuels Gebet wider]{Die bußfertigen Israeliten sagen auf Samuels Gebet wider die Philister}

a Weinete) das ist, die Israeliten im ganzen Reich trugen Leid, und klagten dem HErrn über die Philister.

b Das ist, das noch an der Milch ist.

\section{Summarium}

Ist eine sehr schöne Historia, wie Samuel, als ein frommer, rechtschaffener Prediger, ehrlich das Volk, so unter Gottes Zorn und Strafe war, zur Buße vermahnet, daß sie ihre Sünde erkennen, die Abgötterei abthun und ihr Herz zum HErrn bekehren und ihm allein dienen sollen, V. 3. Und richtet darnach den rechten Gottesdienst wieder an und betet für das Volk.

Wo nun also rechtschaffene Buße und Besserung, reine Lehre und rechter Gottesdienst ist, daselbst will Gott gnädig sein und die Strafe abwenden, wie er denn den Kindern Israel hie hilft. Und hier sehet man sonderlich, wie viel an einem frommen, rechtschaffenen Lehrer, der Gottes Wort rein und treulich prediget, und die Leute zum selben vermahnet, gelegen sei. Darum soll man fromme, gelehrte Kirchendiener in allen Ehren und für das größte Kleinod halten und schön aufheben.

Es wird hie im Text gemeldet, die Kinder Israel schöpften Wasser und gossen es aus vor dem HErrn, V. 6. Solches deutet etliche, daß sie vor dem HErrn geweinet haben über die Sünde, damit sie solche Strafe verdienet hätten. Weil aber der Text auch des Fastens gedenkt, muß es verstanden werden vom Wasser, damit
\quellseite{416}
sie sich und ihre Kleider nach Gewohnheit des Gesetzes gereiniget haben. Das aber ist sehr fein, als die Philister ins Land fallen und jedermann furchtsam und verzagt ist, daß das Volk zu Samuel läuft und spricht: Laß nicht ab für uns zu schreien zu dem HErrn, unserm Gott, daß er uns helfe aus der Philister Hand, V. 8. Solche Zuversicht auf Gottes gnädige Hülfe bewegt Gott, daß er ein greuliches Wetter mit Donner und Blitz ließ über die Philister kommen, daß sie flohen und geschlagen worden, V. 18. Darum sollen wirs auch lernen, daß wir in Nöthen und Angst fleißig beten.

\section{Votum}

Der treue GOtt erhalte uns sein Wort und dessen treue Prediger, und regiere uns alle, daß wir recht glauben, christlich leben, das Kreuz geduldig tragen, und daß wir alle selig werden durch JEsum Christum, unsern HErrn, Amen.

\chapter{Das 8. Kapitel}

\section{Vorrede}

Wohlerathene Kinder und fromme Obrigkeit sind zwei edle, große und heilsame Gaben Gottes. Nun rumoret aber erschrecklich der Teufel im Haus- und weltlichen Regiment, und hindert gewaltig das Gute in diesen Ständen. Demgegen bedenke ein jeder gar wohl, was sein Beruf und Amt sei, treibe die treulich und fleißig und bete, daß die Kinderzucht und Regierung ferner glückselig sei, und zu Gottes Ehren und dem menschlichen Geschlecht zu Nutz gerathe und diene. O du Pflanzgarten der lieben Jugend, o du Lustgarten der lieben Obrigkeit, euch segne Gott der Vater, euer warte der treue Gärtner, unser HErr JEsus Christus, Gottes Sohn, euch begegne Gott der heilige Geist, und behüte euch vor giftigen Raupen und allem anderen Ungeziefer, und verleihe, daß ihr wachset, grünet, blühet und Früchte traget, die Gott wohl anriechen und Land und Leuten dienstlich sein, bis zum ewigen Leben, Amen. Darauf höret ferner an das 8. Kapitel.

\section{Israel begehret einen König. Dessen Recht wird angezeigt}
\quellseite{417}

\section{Summarium}

Ist erstlich ein Exempel, wie der fromme, heilige Mann Samuel ungerathene und gottlose Kinder hat, V. 3. Auf daß man sehe, wie auch heilige Leute bisweilen solch Creuz haben tragen müssen. Darum bedarfs nicht allein fleißiger Zucht, sondern auch ernstliches Beten gegen Gott, daß er fromme Kinder geben wolle.

Zum andern ist hie eine Historia, wie die Juden einen König begehren, allein daß sie in solchem dem Heiden gleich seien, V. 5. Solches zeucht Gott für eine Unehre an; denn sie sollten sich mehr ihres Gottes im Himmel, denn eines sterblichen Königs, getröstet haben, V. 7. Es zeiget auch die heil. Schrift hie an, wie die Unterthanen ihrer Obrigkeit schuldig sind zu dienen, nach ihrem besten Vermögen, mit Leib und Gut, auf daß Friede und gut Regiment möge erhalten werden. Aber jedoch soll solches Regiment keine Tyrannei sein, daß man die Leute übel halten, zu viel schätzen und dringen wollte. Denn daß der HErr spricht: Eure Töchter wird er nehmen, und eure Aecker, eure Weinberge und Oelberge, von eurer Saat und Weinberg wird er den Zehenten nehmen. Mit solchen Worten räumet er der Obrigkeit nicht ein, den Unterthanen zu nehmen, was sie gelüstet, sondern will damit anzeigen, daß die Juden solcher Last hätten können gerathen, wo sie an Gottes Schutz sich genügen und keinen König hätten haben wollen. Wo aber ein König ist, da muß jedermann mit dem Seinen zur Unterhaltung des Regiments helfen. Und doch soll des Königs oder Obrigkeit halben die Gleichmäßigkeit bleiben, daß er ein König und nicht ein Tyrann sei. Könige oder herrschafften Obrigkeit lassen den Unterthanen das Ihre, und nehmen ein Leidliches, auf daß es die Leute ertragen und sich und die Ihren nähren mögen. Die Tyrannen aber nehmen es alles und lassen nichts; das gefällt Gott übel und strafets greulich mit Aufruhr und sonst, wie die Exempel weisen.

\section{Votum}

Helffe der treue GOtt, daß sich alle Kinderlein recht ziehen lassen und als Pflänzlein GOttes wohl und nützlich gerathen; und stehe auch unserer lieben Obrigkeit ferner mit Gnaden bei, daß sie christlich und seliglich regiere, in Christo Jesu, unserm HErrn, Amen.

\chapter{Das 9. Capitel}

\section{Vorrede}

Von unsern leiblichen Eltern, als Vater und Mutter, bekommen wir Leib und Leben; von unsern Zucht- und Lehrmeistern werden wir in guten Tugenden und Künsten unterweiset; aber die Kirchendiener und Prediger sind Eltern und Lehrer der Seelen. Wie soll man sie halten? Das lehret Gott selber und sagt durch den Erzpropheten Mose, 5 Mos. 12, 19.: Hüte dich, daß du den Leviten nicht verlässest, so lange du auf Erden lebest. Davon unterrichtet auch Sirach im 7. Kap. V. 33. und spricht: Ehre den Priester und gib ihnen ihr Theil, wie dir geboten ist. Eben also redet St. Paulus, da er spricht 1 Tim. 5, 17.: Die Priester oder Aeltesten, die wohl fürstehen, die halte man zweifacher Ehren werth; nämlich einer höflichen Ehrbietung und ihres Lohnes und Unterhaltung. Wohl denen, die nach diesen Geboten thun und handeln. Denn es wird ihnen gewiß gehen nach der Verheißung des HErrn Christi, Matth. 10, 41., da er sagt: Wer einen Propheten aufnimmt in eines Propheten Namen, der wird eines Propheten Lohn empfangen, und ein Becher kaltes Wasser wird nicht unbelohnet bleiben. Dies lehret auch das 9. Kapitel.

\section{Saul, der ernannte König in Israel, ist Samuels Sach}
\quellseite{418}

\section{Summarium}

Ist eine Historia, wie Saul zum Propheten Samuel kommt; und werden hie sonderlich zwo feine Exempel angezeigt. Also erste, daß dies der feinste und beste Beruf ist, wenn man unversehens und unbedacht zu einem Amt gezogen wird, und sich nicht selbst eindringet, wie Saul hie zu einem König gesalbet wird, da er seine Eselin suchet und gar nichts vom Königreich noch andern dergleichen denket.
\quellseite{419}

\section{Summarium}

Das andere Exempel ist, daß Saul nicht zum Propheten gehen will, er habe ihm denn etwas zu geben. Denn es will Gott sonderlich, daß man die ehren und dankbar gegen ihnen sein soll, welche den Leuten im Wort und geistlichen Amt fürstehen und dienen, wie St. Paulus sehr fein vermahnet, Gal. 6, 6.: Wer unterrichtet wird mit dem Wort, der theile mit allerlei Guts dem, der ihn unterrichtet. Und 1 Cor. 9, 11.: So wir euch das Geistliche säen, ists ein groß Ding, ob wir euer Leibliches ernten? item, der HErr hatte befohlen, daß die das Evangelium verkündigen, sollen sich vom Evangelio nähren. Darum ist es eine große Sünde, die doch zu unsern Zeiten sehr gemein ist, daß man die armen Pfarrherrn auf dem Lande und sonst so kärglich und übel hält. Es wird aber sonderlich in diesem Kapitel das Predigt- oder Kirchenamt sehr hoch gerühmet, da der Text sagt, daß vorzeiten in Israel, wenn man ginge, Gott zu fragen, sprach man: Kommet, laßt uns gehen zum Seher oder Propheten, V. 9. Denn was sie von ihren Propheten und Predigern gehöret haben, haben sie anders nicht gehalten, denn als hätte es Gott selbst vom Himmel herunter geredet. Wie denn Christus will, daß man seine Prediger auch also hören soll, da er spricht: Wer euch höret, der höret mich. Euch, spricht er, das ist, die ihr mein Wort den Leuten fürtragt, und Buße und Vergebung der Sünden in meinem Namen verkündigt, und nicht dem Papst und seine Bischöfe, die das Evangelium lästern und verfolgen.

\section{Votum}

Gott regiere alle christlichen Lehrer, daß sie in ihrem Amte treu sein, und helfe, daß die Zuhörer sich auch aller Gebühr verhalten, und zeitlichen und ewigen Segen davon haben, in Christo JEsu, unserm HErrn, Amen.

\chapter{Das 10. Kapitel}

\section{Vorrede}

Es ist einem gemeinen Landrichter oder Schulheiß tröstlich, daß, wenn ihm die Sache zu hoch und über seinen Verstand und Vermögen ist, er sich derhalben, zur Erhaltung gutes Gewissens, bei seinem Amtmann, wie billig, Raths und Hülfe erholen kann. Vielmehr soll es aller hohen und niedrigen Obrigkeit tröstlich sein, weil bei den geschwinden Zeiten und unbändigen Leuten die Sachen viel zu hoch, schwer und groß sind, über allen menschlichen Witz und Vermögen, (wie alle christlichen Regenten und Obrigkeiten gern bekennen,) daß sie derhalben zu Hofe suppliciren und bei der höchsten Gewalt, bei dem HErrn, welches Reichs Amtleute sie sind, sich Raths, Hülfe und Beistand erholen mögen. Und eben das will Gott, der HErr, auch haben, und gefällt ihm von aller Obrigkeit und Menschen herzlich wohl, da man nur getrost Rath und Hülfe bei seinen göttlichen Gnaden sich erholet. So helfe Gott, der heilige Geist, daß alle Obrigkeit des sich christlich verhalte, zu ihrem und der Unterthanen Trost und Nutzen, Amen. Dies lehret auch das 10. Kapitel.

\section{Saul zum Könige gesalbet und vorgestellet}
\quellseite{420}

a Wer ist ihr Vater? das ist, lasse sie weissagen, ist's doch nicht vom Vater angeboren, sondern von GOtt, der ist der rechte Vater.

\section{Summarium}

Ist eine Historia, wie Saul zum Könige von Samuel gesalbet wird. Auf daß aber Saul solchem Beruf glaubet, sagt ihm der Prophet etliche Zeichen, auf daß er an denselbigen lerne, Gottes Wort werde ihm nicht lügen, er werde König werden, V. 2.

Zum andern ist sonderlich zu merken, ehe Gott den Saul ins Regiment setzet, gibt er ihm seinen heil. Geist und macht einen andern Mann aus ihm, V. 9. Denn zu einem Regenten, der Glück haben und den Sachen allenthalben soll recht thun, gehöret ein sonderlicher Segen Gottes. Und will sonderlich vonnöthen sein, daß er gottesfürchtig sei und den heil. Geist habe, der ihn recht leite und führe. Denn sonst sind der Sachen und Händel so viel, so sind wir Menschen so schwach und blöde, daß nicht möglich ist, daß ein Regent allzeit recht thue und alle Händel wohl ausrichte; man thut ihm entweder zu viel oder zu wenig. Darum will es vonnöthen sein, daß die Regenten Gott um seinen heil. Geist bitten und sonderlich in Gottesfurcht leben.

\section{Summarium}

Ist ein fein Exempel, ob Gott wohl dem Propheten Samuel genugsam geoffenbaret hatte, er wollte Saul zum König machen; so nimmt dennoch Samuel solche Erwählung öffentlich und nach ordentlicher Weise für, auf daß niemand sich ob solcher Wahl zu beklagen habe, und Saul öffentlich könnte beweisen.

Daß aber Saul sich verkriecht, weil man wählet, V. 23., ist eine Anzeigung, daß ihm für dem Regiment gegrauet habe, als für einem gefährlichen, mühsamen und schweren Amt, wie es denn in der Wahrheit auch ist, und ihm bald im Anfang begegnet, daß ihn etliche verachten. Aber ein weiser Regent soll sonderlich im Anfang viel
\quellseite{421}
lassen für Ohren gehen, und nichts desto weniger seines Berufs mit Fleiß warten; damit wehret man den bösen Mäulern am besten. Wiederum, wer alle Hunde, so ihn auf den Gassen anbellen, schlagen oder werfen will, der wird ihm viel zu schaffen machen und langsam heimkommen. Man lasse sie bellen und warte ein jeder seines Weges und Berufs, so hören sie von ihnen selbst auf, wie hie mit Saul geschieht.

\section{Votum}

Regiere der treue Gott ferner auch unsere liebe Obrigkeit mit seinem heiligen Geist, daß unter ihnen ihm eine Kirche gesammlet und gemeiner Friede erhalten werde, und bewahre sie mit allen Christgläubigen zum ewigen, freudenreichen Leben, Amen.

\chapter{Das 11. Kapitel}

\section{Vorrede}

St. Paulus Röm. 13, 7. lehret, daß man der Obrigkeit schuldige Ehre thun soll. Diese Ehre bestehet darinnen, daß man der Obrigkeit als Gottes Dienerin alle Ehrbietung und Dankbarkeit erzeige; darnach, daß man ihr auch zum Besten gedenke und das Ehrlichste von ihr rede; denn daß man Gott für sie anrufe und auch ihre Gebrechen zudecke, zum Besten deute und zu gut halte. Bei den Römern und Griechen hielt man daher viel darauf, daß der Obrigkeit ihre gebührliche Ehre auch in äußerlichen Geberden erzeiget würde: und wenn sich einer darinnen versehen und etwa einer Person im Regiment nicht gebührliche Reverenz mit Bücken, Neigen und Weichen erzeiget hätte, den hielt man für so gar leichtfertig, daß ihm auch keine Schande noch Laster zu viel sein würde. Gedenket, was will nun uns Christen gebühren, da wir menschliche Vernunft gebrauchen wollen und wider Gottes Gebot nicht thun sollen. Dies erinnert auch das 11. Kapitel.

\section{Sauls erster Sieg wider die Ammoniter}
\quellseite{422}

\section{Summarium}

Ist eine Historia, in welcher man siehet, wie es dem neuen König Saul im Anfang seines Regiments gegangen ist, daß er verachtet ist gewesen und der meiste Theil den Spott aus ihm getrieben. Wie denn alle Sachen, so Gott fürnimmt, zum ersten nicht unangefochten bleiben, sondern gering, schwach und verachtet anzusehen sind. Aber endlich findet man, daß Glück und Heil dabei ist. Wie denn Saul hie mit einem fürstlichen Zorn zusammen bringet und den Feinden absaget, und also auf ein neues zum König erwählet und bestätiget wird. Die aber, so ihn verachtet und übel nachgeredet hatten, müssen sich ducken, V. 15.

Merke aber mit Fleiß, daß der Text spricht: Der Geist Gottes geriet über Saul und sein Zorn ergrimmete sehr, V. 6. Denn solches ist eine Gabe und Werk des heil. Geistes, wo weltliche Obrigkeit sich mit Ernst darum annimmt, daß den Unterthanen geholfen und die bösen Buben gestraft werden. Darum jedermann solche Regenten für Augen haben, Gott für sie danken und bitten soll, daß Gott Glück und langes Leben ihnen geben und sie für Unfall behüten wolle.

\section{Votum}

Helfe Gott, der zu gebieten und zu verbieten allein Macht hat, daß der Obrigkeit und Unterthanen Herz gleich und einig sei, und daß die Obrigkeit väterlich und die Unterthanen kindlich sich gegen einander erzeigen, und gebe unserer Obrigkeit auch Glück und Sieg wider alle Feinde, Amen.

\chapter{Das 12. Kapitel}

\section{Vorrede}

Aller getreuer Prediger Amt ist, daß sie strafen, das ist, allerlei Sünde wider die erste und andere Tafel anzeigen, Gottes Zorn und Rache darüber verkündigen und zur Besserung vermahnen sollen. Denn der HErr Christus hat befohlen Luc. 24, 47, daß man in seinem Namen predigen solle Buße und Vergebung der Sünden. So kann Gottes Wort ohne Strafen nicht recht fürgetragen werden. Denn Gottes Wort zwei Hauptstücke lehret, nämlich Gesetz und Evangelium; das Gesetz soll strafen, das Evangelium soll trösten. Ueber das, so fordert's aller derjenigen Nothdurft, welche sollen wiedergeboren, gerecht und selig werden, daß sie erstlich durch Strafe des Gesetzes zur Erkenntniß der Sünden kommen, darnach, daß sie durch die Gnadenpredigt des Evangelii wieder getröstet werden. Davon lehret auch das 12. Kapitel.

\section{Wie anfänglich Samuel sein Richteramt übergeben}
\quellseite{423}

a Gott bestätiget den König und zürnet doch, daß sie ihn erwählet hatten. Das ist so viel: sie hätten übel, daß sie ihr Vertrauen von Gott auf einen Menschen und sich selbst setzten, so sie bisher so oft ohne König, durch Gott errettet waren. Dazu, weil ihnen versprochen war, Könige zu haben, läßt es Gott nicht zu, daß sie ihn wählen, sondern er selbst wählet, auf daß bestehe, daß alles, was Gott nicht ansiehet und thut, nichts gelte für Gott.

\section{Summarium}

Ist ein sehr fein Exempel von dem Propheten Samuel, in welchem man siehet, wie ein rechter Prediger oder Kirchendiener sich halten soll: nichts um Geldes und Bauches willen thun, sondern mit allem Fleiß den Leuten ihre Sünde und Untugend fürhalten, sie zur Besserung vermahnen und selbst ihnen gute Exempel geben; auf daß der Ruhm bestehen möge, daß in solchem Amt niemand ihm selbst, sondern Gott und dem Nächsten gedienet habe. Man zeuge jemand wider mich, spricht er, V. 3., ob ich jemandes Ochsen oder Esel genommen, ob ich jemand Gewalt und Unrecht gethan, ob ich jemand unterdrückt, ob ich von jemand Geschenk genommen habe, so will ich's wieder geben. Denn das sind gemeiniglich deren, so in Aemtern sind und sonderbare Verwaltung haben, böse Sünde, daß sie mit Gewalt fahren, andere drücken und sich bestechen lassen. Daher denn das Sprichwort kommt: Es sei kein Aemtlein so gering, es sei Heutens werth.

Ein Prediger aber soll sonderlich für dergleichen Aergerniß sich hüten. Denn sonst verursachet er, daß beides, das Amt und das Wort verachtet und gelästert wird, das ist eine große Sünde.

\section{Summarium}

Ist eine sonderliche Historia, wie Samuel, als ein rechtschaffener Prediger, seinem Volk ihre Sünde anzeigt. Es war aber wider Gott und seine Ehre, daß das Volk einen König wollte haben, und sich mehr eines Königs, als Gottes tröstet. Solche Sünde erkannten sie nicht, ließens ihnen auch nicht leid sein, darum zeiget ihnen Samuel solche Sünde an und sagt V. 17., Gott habe es also gestrafet, daß zu einer ungesegneten Zeit, nämlich in der Weizenernte, Regen, Hagel und Ungewitter sei kommen. Wie denn Gott noch unsere Sünde mit Ungewitter, mit großer Dürre oder Nässe, daß wir der Früchte nicht genießen können, strafet, doch ist solche Strafe weit gnädiger, denn Krieg oder langwierige theure Zeit.

Wiederum thut Samuel auch in dem, wie ein frommer Prediger, daß er sein Volk tröstet, und sie vermahnet, daß sie doch fortan wollen frömmer sein, so wolle er auch gerne mit Beten und Lehren das Beste bei ihnen thun.

\section{Votum}

GOtt, der HErr, bezahle allen christlichen Predigern und Seelsorgern ihre väterliche
\quellseite{424}
Liebe und Treue, und regiere uns, daß wir ihre Straffen annehmen, daraus uns bessern und durch das heilige Evangelium wider die Sünde uns mit Freuden trösten, im Leben und Sterben, Amen.

\chapter{Das 13. Kapitel}

\section{Vorrede}

Der königliche Prophet David sagt im 119. Psalm V. 11.: Herr, ich behalte dein Wort in meinem Herzen, auf daß ich nicht wider dich sündige. Wie nun David Gottes Wort nicht allein in Büchern und im Munde, sondern auch in seinem Herzen, als in einem Kasten behält und durch Gottes Hülfe bewahret, auf daß er nicht wider seinen lieben Gott sündige: also sollen wir auch aus Gott sein und nicht allein Gottes Wort hören, sondern auch in einem feinen, reinen, guten Herzen bewahren und darnach leben, Joh. 8, 47. Luc. 8, 15. Denn wer sich für Sünden hüten will, der kann kein besser Mittel finden, denn daß er Gottes Wort in seinem Herzen verschließe, und sehnlich daneben Gott anrufe, daß er ihn dabei gnädig erhalte, und daß er Gottes Gebot nicht aus den Augen sehe, sondern sein Leben, Wesen, Wandel und Handel als nach dem rechten unbetrüglichen Winkelmaß oder Richtscheid anstelle, daß sein Thun und Lassen allzeit mit demselbigen überein komme. Dies soll auch erinnern das 13. Kapitel.

\section{Sauls Mißtrauen im Zuge wider die Philister}
\quellseite{425}

\section{Summarium}

In eine Historia, wie König Saul sich versündiget, in dem, daß er des Propheten Samuelis nicht erwartet und für sich selbst wider den Befehl Gottes allein opfert, V. 9. Mit solcher Sünde verdienet er, daß Gott ihm und seinem ganzen Geschlecht dräuet, es soll ihr keiner zu königlichen Würden kommen, V. 13. Solches Exempel soll jedermann, insonderheit aber die Obrigkeit, wohl merken und bedenken, wie bald man Glück und Heil verscherzt habe, wenn man etwas wider Gottes Wort und Befehl fürnimmt, aus eigener Andacht. Auf daß man sich desto fleißiger an das Wort halte, und ohne und wider dasselbige seinen sondern Gottesdienst anfange.

\section{Votum}

Regiere und führe uns unser HErr JEsus Christus durch den heiligen Geist, daß wir seinem Wort folgen mit Gehorsam in diesem Leben, daß wir auch seinem Wort und Stimme folgen, mit Freuden in das ewige Leben, und allezeit bei ihm sein, Amen.

\chapter{Das 14. Capitel}

\section{Vorrede}

Geliebte Christen! gleichwie Jonathan, da er die Philister alleine mit seinem Waffenträger angreifen wollte, sich deß tröstete, daß er beschnitten, die Philister aber unbeschnitten wären, und stärket sich also mit dem Sacrament der Beschneidung in seinem Glauben, an Gottes Gnade, Verheißung und Beistand, und hat aus starkem Glauben durch Gottes Hülfe die Philister in die Flucht geschlagen und vertrieben: also sollen wir Christen auch in wahrem Glauben, in Noth und Tod, trösten unserer Taufe, welche anstatt der Beschneidung kommen ist, und sollen uns im Leben und Sterben Gott befehlen, auf seine uns darinnen gethane Verheißung, so werden wir auch siegen über Sünde, Tod, Teufel, Hölle und ewige Verdammniß, und werden die alle vertreiben und selig werden durch Jesum Christum, unsern HErrn. Dies soll man auch betrachten bei dem ersten Theil des 14. Capitels.

\section[Heldenthat Jonathans. Saul führt ihn durch unzeitigen Eifer]{Heldenthat Jonathans. Saul führt ihn durch unzeitigen Eifer in Lebensgefahr}
\quellseite{426}

\section{Summarium}

Ist eine sonderliche schöne Historia, wie Gott das Herz Jonathans erwecket und durch ihn allein die Feinde wunderbarlich dämpfet, V. 14. Denn weil Saul sich versündiget hatte an Gott, ward er solches Siegs nicht werth. Aber sein Sohn war an solcher Sünde nicht schuldig, darum gönnet ihm Gott solchen herrlichen Sieg. Und hier siehet man, wie Gott mit tyrannischen Regenten umgehet und bezahlet sie, wie sie verdienet haben. Die Philister hatten die Juden entblößet, daß sie gar keine Waffen noch Schmied im Lande hatten. Solche Tyrannei strafet Gott damit, daß die Philister ihre eigenen Waffen wider sich selbst brauchen und damit unter einander sich selbst erwürgen müssen.

\section{Votum}

Der allmächtige, barmherzige, starke GOtt stärke heute auch alle christliche Obrigkeit zu glückseliger Regierung, und steure allen Feinden seiner Christenheit, geistlichen und leiblichen, sichtbaren und unsichtbaren, und erhalte uns den leiblichen und geistlichen Seelenfrieden. Amen.

\section{Der andere Theil des 14. Kapitels}

\section{Vorrede}

Geliebte Christen! die Eide und Gelübde, so wider Gott und sein Wort, wider den christlichen Glauben, wider unser Taufgelübniß, wider alle Ehrbarkeit und Billigkeit, oder die da unmöglich sein, die soll man nicht thun noch schwören; und da solches geschehen, soll man sie keineswegs halten; denn wie die Juristen sagen: Eide sollen nicht Bande sein der Ungerechtigkeit. Und der alte Isidorus spricht: In bösem Zusagen halte nicht Glauben und verändere dein Gelübde, wenn es schädlich ist. Und da man einen Eid oder Gelübniß hält, so aus Leichtfertigkeit oder Unverstand gethan ist, und die wider Gott sein und dem Nächsten zu Schaden gedeihen möchten, so macht man die Sünde größer. Und Gott, der alle falsche Eide hasset, strafet solche Sünde an allen unbußfertigen zeitlich und ewig. Derwegen im Eide schwören oder Gelübden ein jeder christliche Bescheidenheit gebrauchen soll. Solches lehret auch der andere Theil des 14. Kapitels.
% Scan page 427: Bible verses/narrative, running heads, folios and bare cross-references only; no non-biblical apparatus retained.
\quellseite{428}

\section{Summarium}

Ist eine Historia, wie Saul als ein unbedächtiger Regent das Volk mit einem Eide verbindet, daß niemand nichts für Nachts essen soll, oder müsse sterben. Ob nun wohl Jonathan wider solchen Eid unwissend that, V. 27., schickts gleichwohl Gott also, daß ihn das Volk errettet, daß er nicht sterben mußte. Zum Zeugniß, daß solcher Eid Gott leid gewesen, und derhalben nicht geleistet hat sollen werden, V. 28. Wir sollen aber dabei lernen, wie gefährlich es sei, und so leichtlich ein Unglück möge erreget werden, wo man mit den Gelübden unweislich und nicht nach dem Wort Gottes sich hält und handelt. Denn eben darum muß die Schuld über Sauls Sohn kommen, auf daß solche unbedachte und ohne Gottes Befehl fürgenommene Geistlichkeit an Saul selbst gestraft werde, der ein Heuchler ist und gedenkt, Gott etwas damit abzugewinnen, so alles Volk den ganzen Tag faste. So doch wir, die da arbeiten sollen, müssen zu essen haben, und Gott nicht das Fasten, sondern das Vertrauen auf seine Güte und Liebe gegen den Nächsten und Besserung des Lebens fordert. Dies sind die rechten Stücke, damit man zum Sieg und Gnade Gottes kommt.

\section{Votum}

Der treue Gott behüte uns für aller Heuchelei, und regiere uns, daß wir uns im Glauben gegen ihn und in der Liebe gegen unsern Nächsten recht verhalten, und errete alle, die unschuldig in Gefahr und Noth stecken, Amen.

\chapter{Das 15. Kapitel}

\section{Vorrede}

St. Petrus spricht in seiner zweiten Epistel im 1. Cap. V. 19.: Ihr thut wohl, daß ihr auf unser festes prophetisches Wort achtet, als auf ein Licht, das da scheinet an einem dunkeln Ort. Hiermit lehret der heilige Apostel, daß die Gottesdienste nichts sind und nichts gelten und nützen ohne Gottes Wort. Dies lehret auch unser HErr Christus, und saget Matth. 15, 9.: Vergeblich dienen sie mir, dieweil sie lehren solche Lehre, die nichts denn Menschen-Gebote sind. Derwegen, wer in seinem Beruf handeln und wandeln will, daß Gott gefalle und ihm und dem Nächsten nützlich sei, der soll sich an Gottes Wort halten und nach dem allein sein Thun und seines Berufes Werk anfangen und treiben. Das lehret auch der erste Theil des 15. Capitels

\section{Saul wegen seines Ungehorsams verworfen}
\quellseite{429}

\section{Summarium}

Dies ist eine Historia, wie der König Saul zum andernmal sich öffentlich versündiget und wider den Befehl Gottes handelt, in dem, daß er der Amalekiter König nicht erwürget und das beste Vieh nicht verbannt, V. 9., wie doch Gott ihm zu thun befohlen hatte, V. 3. Und hilft ihm gar nichts, daß er fürwendet, er habe das beste Vieh behalten zum Opfer dem HErrn, V. 15. Denn der fürnehmste und größte Gottesdienst ist, wie hernach folget, gehorsam sein und thun, was Gott heißet, V. 22. Wer aber solches anstehen lässet und dem Wort ungehorsam ist, er nehme für Gottesdienst, was er wolle, so heißt es und bleibet nur ein Greuel und Sünde für Gott.

\section{Votum}

Der treue GOtt helfe und regiere uns, daß wir uns in all unserm Thun und Lassen sein Wort fürleuchten lassen, auf daß wir wider seinen Willen ja nichts handeln, und durch unser sündlich Gutdünken uns und andern Leuten nicht seinen Zorn und Strafe zuziehen, Amen.

\section{Der andere Theil des 15. Kapitels}

\section{Vorrede}

Eine fürnehme Ursache, um welcher willen Land und Leute zu Boden gehen und die Kirche Gottes Anstoß kriegt, die Regimente auch verderbt werden, ist der Ungehorsam gegen Gottes Wort. Derwegen, wenn treue Kirchendiener und Lehrer allerlei Sünde an ihren Zuhörern sehen, sollen sie, je eher je besser, mit ihnen zur Buße eilen. Und wie sie die Unterthanen ihrer Sünde wegen strafen: also sollen sie zu den Sünden der Obrigkeit auch nicht stille schweigen und sie unterschlagen, sondern sollen ihnen auch Buße predigen, damit beiderseits gelehret und gethan werde, was Gott, das Amt, das Gewissen und eines jeden Pflicht, damit man Gott verbunden ist, erfordert, auf daß man die vorherberichten Strafen abwende und durch den HErrn Christum Gott versöhne, damit er dem geistlichen und weltlichen Stande in Gnaden zu glückseligem Zustande allzeit beiwohne. Davon lehret der andere Theil des 15. Kapitels.
\quellseite{430}

a Man soll Gottes Wort nicht ändern noch bessern; es lässet sich nicht ändern.

\section{Summarium}

Hier sind zwei Exempel sonderlich zu merken. Das erste mit Saul, der thut, wie alle Heuchler pflegen, und will noch Recht haben, seine Sünde schmücken und vertheidigen, V. 20. Dagegen thut Samuel, wie einem rechten Prediger zustehet, schonet Sauls nicht, darum, daß er König war, sondern strafet seine Sünde mit Ernst und sagt: Wer Gott rechtschaffen wolle dienen, der muß seinem Wort und Befehl Gehorsam leisten, V. 22. An selbigem Dienst liegt Gott mehr, denn an allen Opfern; wiederum, wer Gottes Wort ungehorsam sei, er diene Gott, womit er wolle, so sei er um des Ungehorsams willen ein Abgöttischer und Götzendiener. Solchen Spruch sollten Mönche und Nonnen, so sie es werth wären, für sich nehmen, so würden sie befinden, daß ihr Gehorsam in Klöstern der größte Ungehorsam gegen Gott ist, würden derhalben Gelübde, Gehorsam und eines mit dem andern fallen lassen, und sich an Gottes Wort und die zehn Gebote mit rechtem Gehorsam begeben. Da könnten sie alßdenn gewiß wissen, daß sie Gott zu Gefallen und zum Gehorsam lebten, V. 33. Am Ende sehet man, wie Gott den Tyrannen Agag auch bezahlet, wie er verdienet hatte. Und lehret, daß die Feinde der christlichen Kirche endlich ein schrecklich Ende nehmen. Auch sehen wir, was Gott im andern Buch Mose im 17. Kapitel V. 14. wider die Amalekiter hat weissagen lassen; das wird alhier erfüllet. Auf daß wir lernen, wie Gottes Strafen allen Unbußfertigen gewiß folgen, unangesehen, daß sie Gott gleich eine Zeitlang aufheucht.

\section{Votum}

Der barmherzige, gütige GOtt und Vater verleihe, daß Prediger und Zuhörer, Obere und Untere, ihre Aemter treulich wahrnehmen, und daß ein jedes darin GOtt schuldigen Gehorsam leiste; und lindere die wohlverdienten Strafen um Christi willen, Amen.

\chapter{Das 16. Kapitel}

\section{Vorrede}

Man schreibt von einem Wurm, den man Spoliator, das ist, einen Räuber, nennet, daß er auf eine Schlange liegt, die im Schatten liegt, und kriecht ihr erst auf den Schwanz, und fesselt sie fein sanft, zuletzt aber zernaget er ihr das Gehirn und tödtet sie: also ist die Sünde der rechte Räuber und Mörder, sie thut den Menschen wohl sanft, weil er sie begehet nach seines argen Fleisches Lust und Wohlgefallen; aber zuletzt tödtet sie ihn. Denn wie das Quecksilber einen Tam zertreibet und läuft durch alle Glieder des Menschen, und zerreißt sie ihm, wenn es nicht wohl getödtet wird: also verderbet die Sünde den ganzen Menschen innerlich und äußerlich, wenn man die mit Gottes Wort nicht tödtet und ihr widerstrebet. Derwegen sollen wir die vergängliche Lust der Welt fliehen, 1 Pet. 2, 11. Und sollen den alten Adam kreuzigen, tödten und begraben, und sollen nicht durch die Freiheit dem Fleisch Raum geben, sondern im Geist wandeln und die Lüste des Fleisches nicht vollbringen, Röm. 6, 6., Gal. 5, 13. 16., sondern in Heiligkeit und Gerechtigkeit leben. Davon höret auch an das 16. Kapitel.

\section{Davids Beruf und erste Salbung zum Könige}
\quellseite{431}

\section{Summarium}

Ist eine Historia, wie David vom HErrn zum König erwählet und von Samuel gesalbet ist worden, V. 13. Und hie gehet die Historia an, wie die Sünde Sauls sei gestraft worden, aufs erste damit, daß GOtt ihm seinen Geist entzeucht, und lässet ihn treiben von einem bösen Geist, der ihn irre und unruhig machet, daß er hinfort nimmer zufrieden kann sein, V. 14. Auf daß wir lernen die Sünde fliehen und uns dafür hüten, sonsten sie GOtt so greulich strafet; wie wir von Saul hernach weiter hören werden. Durch diese Ursache aber kommet David gen Hof und wird bei dem König und ganzen Volk bekannt.
\quellseite{432}

\section{Votum}

Gott, der heilige Geist, regiere und stärke uns, daß wir den befleckten Rock des Fleisches hassen, und die Werke der Finsterniß ablegen und anziehen die Waffen des Lichts, damit wir nicht auch zeitlich gestraft und ewig verdammt werden, Amen.

\chapter{Das 17. Kapitel}

\section{Vorrede}

Gott ist mächtig, auch in den geringen Menschen große Dinge zu thun. Ja, Gott hat im steten Gebrauch gehabt, die allergrößten Dinge durch unansehnliche Personen auszurichten, mehr, denn durch ansehnliche. Derwegen niemand auf seine Stärke, Macht und Gewalt sich verlassen soll; auch soll niemand in seiner Schwachheit verzagen, wenn er der Dinge sich annimmt, die Gottes Ehre und der Kirchen und der Christenheit Nutz angehen. Denn Gott haffet, was die Welt hoch und groß achtet, und liebet, was für der Welt gering ist, da man nur in seinem Stande demüthig ist, sich und die Verrichtung seines Berufs Gott mit gläubigem Gebete befiehlt, und sich tröstet in wahrem Glauben, daß er durch einen jeden ausrichten könne und wolle, was zu seiner Ehre und andern zu Hülfe, Rettung und Trost dienen soll. Dies erinnert auch der erste Theil des 17. Kapitels.

\section{Davids Streit mit dem Riesen Goliath}
\quellseite{433}

\section{Summarium}

Ist eine Historia, wie David, voll Geistes und Glaubens, mit der Hülfe Gottes den Philister Goliath erwürget, V. 50. Und hie ist im ersten Theil ein Exempel an Goliath wohl zu merken, in welchem uns der Welt Art wird fürgemalet. Denn gleich wie Goliath auf seine Stärke und Größe trotzet, frech und sicher ist, gleich als bedürfe er unsers Herrn Gottes Hülfe gar nicht: also thut die Welt auch, lebt sicher ohne Gottesfurcht, denkt, wenn sie Geld und Gut genug habe, so habe sie alles und bedürfe nichts mehr. Aber wie solcher Trotz und Hochmuth ein Ende nehme, werden wir hernach hören, auf daß jedermann lerne, sich nichts anders rühmen, denn Gottes, seiner Gnade und Barmherzigkeit.

\section{Summarium}

Ist ein fein Exempel, in welchem man siehet, wie Gott seine höchste und gewisseste Hülfe erweisen anschicket, daß es für der Welt ganz und gar kein Ansehen hat. Denn weltlich davon zu reden, wäre mehr Hoffnung dabei gewesen, wenn Saul oder sonst ein großer, starker Mann sich wider den Philister zu streiten hätte brauchen lassen. Aber daß sich David unterstehet, ohne all Schwert und Harnisch mit seiner Schleuder
\quellseite{434}
und Hirtenstab, das scheinet über die Maße liederlich und kindisch. Daher schilt ihn sein ältester Bruder, V. 28. Und der König Saul hat auch geringen Glauben dazu, V. 33. Aber Gott pflegt allweg zu thun, daß er große Wunderwerke durch geringe, verachtete Leute ausrichtet. Erstlich darum, daß die Welt Ursache genug habe, sicher zu sein und Gottes, der seine Sachen schwächlich und thöricht fürnimmt, zu spotten und ihn zu verlachen. Zum andern und fürnehmlich darum, daß man sehe, ja greifen müsse, daß es Gottes Werk und Hülfe sei, und nicht menschlich Vermögen, und also an dem Gegentheil jedermann sehe, daß man mit menschlicher Weisheit, Macht und Vermögen müsse zu Schaden werden, wo Gott nicht dabei ist.

\section{Votum}

Der starke, allmächtige GOtt steure der sichern Welt Trotz und Hochmuth, und stärke seine verachtete Christen, daß sie selige Werkzeuge seiner Kirchen sein und, ihm zu Lob, und Land und Leuten zu Nutz, Freude und Trost, leben, Amen.

\section{Der andere Theil des 17. Capitels}

\section{Vorrede}

Davids Waffen wider Goliath sind ein Stab und Steine in der Hirtentasche; also sollen wir auch den höllischen Goliath, den Teufel, überwinden. Der Stab bedeutet Gottes Wort, die Steine den Glauben, die Hirtentasche das Herz. So sind nun die rechten und kräftigen Mittel, den Teufel zu schlagen und zu überwinden, daß man Gottes Wort und Verheißungen von Herzen glaube und derer wider den Tausendkünstler und alle seine Anläufe, Sturm und Anfechtungen, sich beständig, in Kraft des heil. Geistes, tröste. Derwegen wohl dem, der in Kreuz und Anfechtung wider den Satan also streitet. Denn wie der Apostel Jacobus sagt 1. Kapitel V. 12.: Selig ist der Mann, der die Anfechtung erduldet; denn nachdem er bewähret ist, wird er die Krone des Lebens empfahen, welche Gott verheißen hat denen, die ihn lieb haben. Dies soll auch erinnern der andere Theil des 17. Kapitels.
\quellseite{435}

\section{Summarium}

Ist das herrlichste und beste Stück aus dieser Historia, wie David den Goliath erwürget, nicht durch Leibes Kraft, denn Goliath war mit Stärke dem David und allem Volk weit überlegen, sondern dadurch, daß er sich seines Gottes und seiner Hülfe tröstet. Hierin war David dem Goliath überlegen, der sich allein auf seines Leibes Stärke und Größe verließ, und hatte Gott nicht bei sich. Wie denn David zu ihm sagt: Du kommest zu mir mit Schwert, Spieß und Schild; ich aber komme zu dir im Namen des HErrn Zebaoth, V. 45. Wer ein solch Vertrauen auf Gottes Hülfe und Gnade fassen kann, dem wird der Teufel nicht vermögen etwas zu thun.

Und siehet man sonderlich sein an Goliath, wie es so sehr bald um die Welt mit aller ihrer Stärke, Geld und Gut geschehen ist, wenn Gott nicht dabei ist.

Zu dem sehen wir fürnehmlich in dieser schönen Historia, wie daß David eine Figur Christi unsers HErrn und Heilandes ist, der den rechten Goliath, den Teufel, überwunden, sein Reich zerstöret, die Sünde getilget, den Tod erwürget und die Hölle zerbrochen hat, zu gut, Trost und Seligkeit allen denen, so an ihn glauben, seinen Namen bekennen und ihm von Herzen dafür danken.

Solchen großen herrlichen Sieg erlanget der HErr auch wunderbarer Weise, stellet sich noch närrischer (wie es die Vernunft ansiehet) dazu, denn David: stirbt des schmähligen Todes am Kreuz, zwischen zween Uebelthätern, als der ärgste Aufrührer und Gotteslästerer; richtet aber dadurch das aus, daß ihm Sünde, Tod und Teufel forthin ausgezogen und erwürget, in Ewigkeit zu Füßen liegen müssen und seinen Christen nicht schaden mögen.

\section{Votum}

Der treue GOtt stehe bei allen, die heute unter dem Kreuz, in Anfechtung wider den Satan streiten, und helfe ihnen und uns ritterlich überwinden, daß wir alle in Christo Jesum, der des Teufels Werk zerstöret hat, selig werden, Amen.

\chapter{Das 18. Kapitel}

\section{Vorrede}

Gott suchet heim mit dem Kreuze, beide, die Gottlosen und die Frommen. Die Gottlosen darum, daß sie ihrer Sünde wegen gestraft werden, andern zum Exempel und Warnung, damit sie sich für dergleichen Sünden hüten, sollen ihnen nicht dergleichen Sündenstrafen folgen. Die Frommen greift Gott an mit dem Kreuze, auf daß sie nicht in Hoffart gerathen und der Gaben wegen, so er ihnen verliehen, sich nicht überheben, und daß er den zukünftigen Sünden also wehre. Dieß lehret St. Paulus in seinem eigenen Exempel und sagt 2 Cor. 12, 7.: Auf daß ich mich nicht der hohen Offenbarungen überhebe, ist mir gegeben ein Pfahl ins Fleisch, nämlich, des Satans Engel, der mich mit Fäusten schlage, auf daß ich mich nicht überhebe. Und St. Petrus 1. Epistel 4. Kapitel V. 12. spricht: Wer am Fleische leidet, der höret auf von Sünden, daß er hinfort, was noch hinterstelliger Zeit im Fleisch ist, nicht der Menschen Lüsten, sondern dem Willen Gottes lebe. Davon höret auch an das 18. Kapitel.

\section{Davids Lob, Lebensgefahr und Heirath}
\quellseite{436}

\section{Summarium}

Ist eine Historia, in welcher wir sehen, wie ein greulich Ding es sei, wo man einmal von der Furcht Gottes und seinem Willen abgetreten ist und also beharret. Denn Saul, da er zum andernmal wider Gott gehandelt hatte, entzog ihm Gott erstlich seinen Geist und ließ ihn plagen mit einem bösen Geist. Derselbe treibet ihn, wie wir hören, daß er David feind wird, und ihn begehrt umzubringen, V. 29. Solche Gefahr aber verhänget Gott über David, daß er nicht stolz noch sicher würde, weil ihn Gott so mit einem herrlichen Sieg begnadet hatte; sondern er weiter Ursache hätte, Gott zu bitten um seine Hülfe und Gnade, und sich seines Worts und Zusagung zu trösten.
\quellseite{437}

\section{Votum}

Gott regiere und behüte uns alle, daß wir durch Sünde den heiligen Geist nicht verjagen, und helfe, daß wir alles Kreuz in Geduld tragen, und uns auch seiner Verheißungen, Beistandes und Hülfe allzeit trösten, Amen.

\chapter{Das 19. Kapitel}

\section{Vorrede}

Der Teufel führet die Menschen aus einer Sünde in die andere. Gott erlöset seine Gläubigen aus einem Kreuz nach dem andern. Dieses soll uns dienen zur Warnung und zu Trost. Zur Warnung, daß, wer einmal dem Teufel gefolget, daß er alsbald Buße thue, auf daß ihn der Teufel nicht in mehr Sünden und denn in ewiges Verderben führe. Zu Trost, daß Gott die Seinen in Gefahr ihres christlichen Berufs erretten und erhalten will. Diesen Trost lehret Gott durch den Propheten Esaias im 54. Kapitel V. 14., da er also saget: Du wirst ferne sein von Gewalt und Unrecht, daß du dich dafür nicht dürfest fürchten; und von Schrecken, denn es soll nicht zu dir nahen. Aller Zeug, der wider dich zubereitet wird, dem soll nicht gelingen. Das ist das Erbe der Knechte des Herrn, und ihre Gerechtigkeit von mir, spricht der Herr. Dies lehret auch das 19. Kapitel.

\section{David begibt sich, das Leben zu erhalten, ins Elend}
\quellseite{438}

a) Bloß nicht daß er nackt gewesen sei, sondern hat die königlichen Kleider abgelegt, und nur gemeine Kleider anbehalten, als ein anderer Mensch. Und fiel nieder den ganzen Tag etc. das ist, er betete mit ihnen, und wenn sie niederfielen, fiel er auch nieder mit ihnen.

\section{Summarium}

Wie ist eine Historia, in welcher man siehet, wie der Teufel nicht feiert, sondern wo er einmal hinter einen Menschen kömmt, treibet er ihn je länger je mehr in größere Sünde. Wie Saul sich hie öffentlich lässet hören, er wolle David auch erwürgen, und nimmt's auch mit der That für, V. 10. 11.

Hie siehet man auch, wie durch wunderbare Mittel Gott den unschuldigen David warnet und errettet, V. 23. Denn Gott will die nicht verlassen, so sich seiner Hülfe in einer guten Sache trösten. Darum muß Saul ehe wieder zum Propheten werden, und seine Knechte auch, ehe dem frommen David etwas von ihnen widerfahren soll.

\section{Votum}

Der treue, barmherzige GOtt lehre und regiere uns, daß uns der Teufel in Sünden nicht verführe; steure auch allen Widerwärtigen seiner Christenheit, und erhalte unter dem Kreuz alle, die seinen heiligen Namen lieben und ehren, Amen.

\chapter{Das 20. Kapitel}

\section{Vorrede}

St. Paulus 2 Tim. 3, 12. sagt: Alle die gottselig leben wollen in Christo Jesu, müssen Verfolgung leiden. Und 1 Cor. 10, 13. spricht er: Gott ist getreu, der euch nicht läßt versuchen über euer Vermögen, sondern macht, daß die Versuchung so ein Ende gewinne, daß ihr könnet ertragen. Danach St. Petrus 1 Epistel 5, 8. sagt: Euer Widersacher, der Teufel, gehet umher, wie ein brüllender Löwe, und suchet, welchen er verschlinge. Und in der Offenbarung Johannis im 12. Kapitel v. 12. spricht eine Stimme im Himmel: Wehe denen, die auf Erden wohnen, und auf dem Meer; denn der Teufel kommt zu euch hinab, und hat einen großen Zorn, und weiß, daß er wenig Zeit hat. Der Spruch St. Pauli ist ein Exempel der verfolgte David. Des Spruchs St. Petri und der Stimme im Himmel ist ein Exempel der Tyrann Saul. Von beiden lautet das 20. Kapitel.

\section{Freundschaft Jonathans und Davids}
\quellseite{439}

\section{Summarium}

Dies ist eine Historia, in welcher man siehet, wie dem frommen David so angst und bange bei den Sachen wird, daß er nirgend seines Leibes und Lebens sicher ist, V. 1. Solches Kreuz und Anfechtung kommet ihm unverdienter Sache; denn ers je um Saul nicht verdienet. Auf daß wir lernen, wie Gott seine Heiligen so mancherlei versucht, und ihnen Ursache genug gibt, daß sie beten und nicht sicher sein sollen. Aber wiederum ists tröstlich, daß ihm Gott so einen treuen Freund bescheret an Jonathan, daß er ihn warnet und für Unglück behütet, V. 4.

a Das ist, du bist nicht ein tapferer, berüchtigter, sondern eines verachteten, geringen und unartigen Weibes Sohn.
\quellseite{440}

\section{Summarium}

Ist ein schrecklich Exempel, in welchem wir sehen, wie der Teufel den König Saul aus einer Sünde in die andere wirft. Bisher hat er David nach dem Leben gestanden, jetzt schießt er nach seinem eignen Sohn, und begehrt ihn auch zu erwürgen. Darum ist es hoch vonnöthen, daß man sich mit allem Ernst in Gottes Furcht bewahre und für Sünden hüte. Denn kommet man einmal darein, so weiß man nicht, wie und wenn man wieder heraus kommet, wie dies Exempel anzeigt.

\section{Votum}

Gott regiere uns, daß wir uns den Teufel aus einer Sünde in die andere nicht führen lassen; und stehe seinen Gläubigen bei in Versuchung und Anfechtung, daß sie alles das in Geduld tragen, und ritterlich anstehen und überwinden, Amen.

\chapter{Das 21. Kapitel}

\section{Vorrede}

Die Liebe und Noth haben kein Gebot. Dieses ist wohlgeredt, man soll es aber auch recht verstehen. Vom Gehorsam gegen menschlichen Satzungen ist es ingemein wahr; aber was Gottes Gesetze anlanget, soll man das verstehen von dem Gehorsam gegen den Statuten und Ordnungen vom jüdischen Gottesdienst, und Ceremonien in der israelitischen Kirche gebräuchlich, und denn von jüdischer Polizei oder gerichtlichen Gesetzen; doch so ferne, da diese mit den natürlichen oder moralischen Gesetzen, das ist, mit den heiligen Zehen Geboten, nicht zustimmen. Denn was in diesen levitischen Kirchen Ceremonien und auch in weltlichen Gerichtordnungen der Juden mit den heiligen Zehen Geboten ingemein einstimmet, das soll auch gehalten und gethan werden. Dieser Lehre ist ein Exempel das 21. Kapitel.

\section{David hungert, fliehet und verstellet sich}

a Zeug in der Schrift heißt ein jeglicher Leib ein Zeug, wie auch Act. 9, 15. Christus von St. Paulo sagt: Er ist
\quellseite{441}

\section{Summarium}

Ist eine Historia, wie David in der Noth, da er sonst nichts hatte zu essen, die heiligen Brode isset, die doch sonst niemand denn die Priester essen durften, V. 6. Solches deutet Christus im Evangelio dahin, daß man das Gesetz so hart nicht spannen, sondern allweg auf die Noth leben und um der Noth willen etwas nachgeben soll. Auf daß wir nicht thun, wie die Phariseer, die hätten eher einen Menschen verderben lassen, ehe sie am Sabbath ihm geholfen hätten. Und die Papisten ließen eher einen Menschen sterben, ehe sie ihm am Freitag Fleisch geben. So doch Gott solches nirgends verboten hat.

\section{Votum}

Der treue Gott nehme sich unserer Noth an, als ein treuer, lieber Vater, und regiere uns, daß wir unsern Nächsten in seiner Noth auch nicht verlassen, sondern ihm allezeit mit Rath und That gern und treulich dienen. Amen.

\chapter{Das 22. Capitel}

\section{Vorrede}

Der Teufel feiret nicht, sondern wo er einmal hinter einen Menschen kömmt, so treibet er ihn je länger je mehr in andere und größere Sünde. Dergleichen ist es hoch vonnöthen, daß man sich mit allem Ernst und Fleiß in Gottes Furcht bewahre, für Sünden hüte, und Gott um Hülfe anrufe, daß der Teufel unser nicht mächtig werde, und uns aus einer Sünde in die andere, zu unserm Verderben, nicht führe. Dieses lehret der HErr Christus, da er vermahnet Matth. 26, 41.: Wachet und betet, daß ihr nicht in Anfechtung fallet. Dies lehren auch beide heiligen Apostel des HErrn Christi, Petrus und Paulus. Petrus saget 1. Epistel 5. Kapitel V. 8.: Seid nüchtern und wachet, denn euer Widersacher, der Teufel, gehet umher wie ein brüllender Löwe, und suchet, welchen er verschlinge; dem widerstehet fest im Glauben. Und St. Paulus Ephef. 6, 10. spricht: Seid stark in dem HErrn und in der Macht seiner Stärke; ziehet an den Harnisch Gottes, daß ihr bestehen könnet gegen die listigen Anläufe des Teufels. Und betet stets in allem Anliegen mit Bitten und Flehen im Geist. Dies lehret auch das 22. Kapitel.

\section[Saul läßt durch den Verräther Doeg fünf und]{Saul läßt durch den Verräther Doeg fünf und achtzig Priester tödten}
\quellseite{442}

\section{Summarium}

Ist eine schreckliche Historia, wie der böse Geist den König Saul immer in mehr und greulichere Sünde treibet, daß er ohn alle Ursache 85 Priester und eine ganze Stadt mit Weib und Kind, mit allem Vieh, erwürgen lässet, V. 18.; und solchen Mord noch mit Lügen schmücket: David stelle ihm nach Leib und Leben, V. 13. So doch David für ihn, ohn alle Ursache, weder Leibes, noch Lebens sicher war. Also siehet man, wie der Teufel ein Lügner und Mörder ist. Darum soll man sich mit allem Fleiß für
\quellseite{443}
Sünden hüten, und stets bitten, daß uns Gott nicht in Versuchung führen lasse.

\section{Votum}

Gott, der Herr Zebaoth, kenne des Teufels List und Gewalt, und regiere und stärke uns in unserer großen Schwachheit, daß wir wider des Teufels Anfechtung streiten, und uns von ihm ja nicht verführen lassen, Amen.

\chapter{Das 23. Kapitel}

\section{Vorrede}

Eines jeden Beruf, Handel und Wandel, dienet Gott und dem Nächsten, so man sie in Gottes Furcht thut, und nach Gottes Befehl und Willen; und je größere Noth darinnen fürfället, je näher ist der treue Gott mit Rath und Hülfe. Daher gehöret, daß Gott selber sagt im 91. Psalm V. 15.: Ich bin bei ihm in der Noth, ich will ihn heraus reißen und zu Ehren machen. Denn Gott hat auch mancherlei Mittel, der Gottlosen Anschläge wider die Gottseligen zu verhindern. Also gehöret auch hieher, daß der heilige Apostel Paulus sagt, 2 Cor. 4, 8. 9.: Wir haben allenthalben Trübsal, aber wir ängsten uns nicht; uns ist bange, aber wir verzagen nicht; wir leiden Verfolgung, aber wir werden nicht verlassen; wir werden unterdrückt, aber wir kommen nicht um. Davon höret auch an das 23. Kapitel.

\section[David wird von neuem verrathen, verfolgt, und wunderbarlich]{David wird von neuem verrathen, verfolgt, und wunderbarlich errettet}

\section{Summarium}

Ist eine schöne Historia, wie David, ehe er
\quellseite{444}
wider die Philister, seinem Volk zu gut, ausgezogen, zuvor den HErrn fraget, ob ers thun soll, V. 2. Denn ohne Beruf soll niemand nichts fürnehmen. Darnach ist auch ein fein Exempel, wie Gott den frommen David warnet für den untreuen und undankbaren Bürgern zu Kegila, V. 11. Auf daß wir lernen, obgleich Gott den Seinen mancherlei Kreuz und Anfechtung zuschicket, so schaffet er doch alleweg Hülfe und Rath daneben.

a Sela Mahelkoth) das heißet Scheidefels.

\section{Summarium}

Ist eine Historia, wie der fromme David von den Siphitern ausgespähet und dem König Saul verrathen ward. Aber Gott schaffts wunderbarlich, da ihn Saul umringet, und am gewissen hatte, fallen ihm die Philister ins Land, daß er muß abziehen, V. 27. Gleich wie zu unserer Zeit die Feinde des Wortes Gottes oft Arges im Sinn gehabt haben; aber Gott hats bisher allezeit also geschicket, daß sie von ihrem Fürnehmen haben ablassen müssen. Er wirds auch forthin thun, wenn wir mit Beten getrost anhalten.

\section{Votum}

Der barmherzige Gott, der treue Menschenhüter, rathe und helfe heute auch allen frommen Christen, die in Noth sind, und steure allem bösen Fürnehmen der Feinde seiner Kirchen, damit sie erhalten werden wider alle Pforten der Hölle, Amen.

\chapter{Das 24. Kapitel}

\section{Vorrede}

Rachgier ist eine Sünde, die Gott erzürnet, das Gewissen beschweret und ihre eigene Sache verderbet; dagegen, wer verhören und verzeihen kann, der gefällt Gott und den Menschen wohl, da wird Ehre aus, bleibet ungebissen und verhütet viel Schaden. Derwegen wer ein rechter Christ sein will, soll sich selbst überwinden, Geduld haben und sanften Muth, und auch dem Feinde thun das Gut. Dies lehret der HErr Christus Matth. 5, 44., da er also spricht: Ich sage euch, liebet eure Feinde, segnet, die euch fluchen, thut wohl denen, die euch hassen, bittet für die, so euch beleidigen und verfolgen, auf daß ihr Kinder seyd eures Vaters im Himmel. Dies unterrichtet auch St. Petrus 1. Epistel am 3. Kapitel V. 9., da er sagt: Vergeltet nicht Böses mit Bösem. Und St. Paulus Röm. 12, 19. vermahnet: Rächet euch selbst nicht, meine Liebsten, sondern gebet Raum dem Zorn (Gottes). Denn es stehet geschrieben, die Rache ist mein, ich will vergelten, spricht der HErr. Davon höret auch an das 24. Kapitel.

\section{David kommt bey Saul in der Höhle}
\quellseite{445}

a Sein Fuß zu decken) so pflegt es die heilige Schrift, daß sie Füße decken heißt, auf das heimliche Gemach gehen.

\section{Summarium}

Es ist eine sehr schöne Historia, wie David, da er seinen Feind, den König Saul, schon in seinen Händen hatte, dennoch sein schonet und sagt, er wolle an den Gesalbten des HErrn die Hand nicht legen, V. 7. Denn ob ihn Gott gleich zum König lassen salben, hat er ihm doch nicht befohlen, daß er Saul tödten und das Königreich ihm selbst nehmen sollte, sondern warten, bis ihm Gott in die Hände gebe. Wiederum thut David recht, daß er Saul seine Unschuld anzeiget, und ihn um sein unbillig Fürnehmen strafet, V. 11. Denn großen Herren soll man ihre Sünde eben so wenig billigen und verschweigen, als armen und geringen Leuten.

\section{Votum}

Verleihe GOtt, daß wir alle, Obrigkeit und Unterthanen, die Sündenstrafen aus GOttes Wort annehmen, und daß wir unsere Feinde auch mit christlicher Geduld überwinden, und ihnen feurige Kohlen auf ihr Haupt sammeln, Amen.
\quellseite{446}

\chapter{Das 25. Kapitel}

\section{Vorrede}

Beide die Gottlosen und Frommen haben zu Sünden Zuneigung, aber mit großem Unterschied. Die Gottlosen folgen ihren bösen Zuneigungen und Lüsten, dem Teufel zu gefallen und ihnen selbst zu Schaden. Die Frommen aber widerstreben den bösen Zuneigungen und Lüsten, Gott zu Ehren, und daß sie nicht, weder ihnen selbst, noch andern, Unglück zurichten. So lasset nun die Sünde nicht herrschen in eurem sterblichen Leibe, ihr Gehorsam zu leisten in ihren Lüsten. Auch begebet nicht der Sünden eure Glieder (vernehmet nicht allein die Augen, Ohren, Hände und Füße, sondern auch die innerlichen Sinne, das Herz und seine Gedanken) zu Waffen der Ungerechtigkeit (verstehet, Unrechtes zu thun); sondern begebet euch selbst Gott und eure Glieder (vernehmet aber, nicht allein die äußerlichen, leiblichen Gliedmaßen, sondern auch alle innerlichen Sinne und Kräfte) Gott zu Waffen der Gerechtigkeit (vernehmet, daß ihr im Glauben Gott neuen Gehorsam leistet, in aller Tugend und Ehrbarkeit): und diese Vermahnung thut der heil. Geist an alle Christen, durch den heiligen Apostel Paulus an die Römer, am 6. Kapitel V. 12. 14. Dies lehret auch der erste Theil des 25. Kapitels.

\section[Samuels Tod. Nabals Thorheit. Vorsichtigkeit Abigails. Davids Werbung]{Samuels Tod. Nabals Thorheit. Vorsichtigkeit Abigails. Davids Werbung um sie}

\section{Summarium}

Dies ist erstlich ein Exempel an Nabal, wie bisweilen grobe Leute ihnen selbst alles Unglück mit ihrer Grobheit zurichten, V. 10.

Darnach ist ein Exempel, wie David sich den Zorn übergehen läßt, und gedenkt solche Grobheit und Unfreundlichkeit an Nabal zu rächen, und sein eigener Richter zu sein, V. 13. Aber solches war Sünde und Unrecht. Denn obwohl Nabal hatte unrecht gehandelt, sollte doch David nicht sein selbst Richter sein. Also begegnet es auch den heiligen Leuten, daß sie zu jach und geschwind sind, und Unrecht zu thun fürnehmen. Aber sie lassen sichs bald gereuen, und hören auf, da die andern fortfahren, und ihnen nichts abbrechen.
\quellseite{447}

\section{Summarium}

Ist ein schön Exempel, wie Gott Mittel schicket, daß David nicht fortfähret und sich an Nabal rächt, wie er im Sinn hatte. Und hie siehet man, wie man mit Freundlichkeit und guten Worten oft eine Sache führen und Unglück verhüten kann. Darum soll jedermann sich befleißigen, daß er sich nicht bald lasse bewegen zu Zorn und Rache, daß ihn darnach gereuen möchte. Sonderlich aber ists fein, daß Abigail zu David sagt, wo er ihm selbst geholfen und sich an seinem Feind gerochen hätte, würde ihm solche Sünde ein Anstoß und Aergerniß gewesen
\quellseite{448}
sein, wenn er zum Regiment wäre kommen, V. 13. Denn Gott läßt solche Sünde nimmer ungerochen abgehen.

\section{Votum}

GOtt vergebe es allen, die ihre Hand von gottfürchtigen Herzen abziehen, und welche die verlassen, denen sie helfen sollen; und regiere uns, daß wir unserm Fleisch und Blut zu Sünden nicht folgen, sondern GOtt und seinem Wort folgen, in allem Guten, unserm Nächsten zu Nutz, Amen.

\section{Der andere Theil des 25. Kapitels}

\section{Vorrede}

Unsere lieben alten Deutschen haben recht gesagt: Wer jach, dem gebrach; wer vertrug, der hat genug; wer schweiget, der bleibet; Gott richt't, wenn niemand spricht. Denn Gott sagt selber 5 Mos. 32, 35.: Mein ist die Rache, ich will vergelten. Derwegen für Nachgier sich männiglich hüten soll. Auch sind dieses der alten Deutschen gute Sprüchwörter gewesen: Der hat genug, der sich begnügen läßt. Laß dir begnügen, so ist dir Gott und die Welt hold. Aus kleinen Brünnlein trinket man sich auch satt, sowohl als aus großen. Laß dir nicht so viel gelieben, so wird dichs nicht betrüben. Denn also sagt der Apostel Paulus 1 Tim. 6.: Es ist ein großer Gewinn, wer gottselig ist und lässet ihm genügen. Derwegen für Geiz sich auch jedermann soll hüten. Dies soll ferner erinnern der andere und letzte Theil des 25. Kapitels.

\section{Summarium}

Ist ein fein Exempel, wo man Gott läßt Richter sein und seiner Hülfe erwartet, daß Gott nicht außen bleibt, sondern die Untugend straft, wie man verdienet hat. Wie man siehet, daß Gott den David an Nabal rächet, und ihm, wie er verdienet hatte, vergilt, daß er sich zu Tode kümmeret und stirbt, V. 38. Solches Exempel gehet aber sonderlich auf die Reichen und Geizigen, so um Geizes willen andern fürhalten oder kärglich bezahlen, was sie ihnen schuldig sind, und man ihnen abverdienet hat. Darum hat dieser Nabal über die Maßen viele Brüder in allen Städten und Dörfern. Und wir sehen, wie es letztlich mit ihnen hinaus gehet, daß sie dahin fallen und sterben, und selbst ihres Guts nicht genießen noch froh werden, welches sie andern nicht gönnen; bisweilen aber wird ihnen bei alle ihrem Gut so angst und bange, daß sie sich selbst hängen oder sonst umbringen. So ein schädliches, greuliches Laster ist der Geiz.

\section{Votum}

Der treue GOtt und Vater unsers HErrn JEsu Christi segne ferner alle, die von seinem Segen andern mittheilen, und regiere uns, daß wir uns für sündlichem Zorn, Nachgier, Geiz und allen andern Sünden hüten, und aller Sünden Strafe entfliehen, Amen.

\chapter{Das 26. Kapitel}

\section{Vorrede}

Gleich wie fromme Herzen ihre Feinde mit Geduld und Wohlthaten überwinden, daß sie sich schämen müssen und ihre Uebelthat bekennen; denn das ist allein die rechte Form und Weise, an Feinden sich zu rächen: also stellen sich die Heuchler freundlich, dieweil es ihre Noth erfordert, wenn die aber ein wenig vorüber ist, da hat ihre Freundlichkeit auch ein Ende, und wird also offenbar, daß es eine gefärbte Liebe und Versöhnung sey, dahinter nichts denn Galle, Blut, Mord und Verderben steckt. Daher sagt

a Neminem invitat, nec pauperes curavit. Das ist, er hat Niemand eingeladen, und sich nicht um die Armen bekümmert.
\quellseite{449}
Salomon in Sprüchen am 29. Kapitel: Die Blutgierigen hassen den Frommen, aber die Gerechten suchen seine Seele. Derwegen niemand sich selbst an seinen Feinden soll rächen, sondern soll durch Gott und die weltliche Obrigkeit sich rächen lassen, und soll sich auch das Böse nicht lassen überwinden, sondern das Böse soll er mit Gutem überwinden, Röm. 12, 21., das ist, durch Geduld, Sanftmuth und Wohlthat gegen seinen Feinden. Dies soll auch erinnern das 26. Kapitel.

\section{David nimmt Sauls Spieß und Wasserbecher}
\quellseite{450}

\section{Summarium}

Dies ist eine Historia, wie David zum andern mal seinen Feind in seinen Händen hat, und dennoch sich enthielt, sein schonet und nicht sein selbst Richter wird, V. 10. Solches läßt Gott dem Saul zu gut geschehen, daß er den unbilligen Haß und Zorn gegen David fallen lassen, und im Werk lernen und erkennen soll, daß ihm David kein Leid zu thun begehrt. Aber die Gottlosen verachten entweder solche Warnung gar, oder haben nur eine kleine Galgen-Reue, thun hernach wie vor, wie man hie an Saul siehet.

\section{Votum}

Gott steure und wehre allen falschen Herzen, und regiere uns, daß wir treue, ungefärbte, beständige Liebe gegen unsern Freunden und Feinden tragen, daß Gott zu Ehren und ihnen zur Besserung, Nutz und Wohlfahrt geistlich und leiblich diene. Amen.

\chapter{Das 27. Capitel}

\section{Vorrede}

Der heilige Apostel Petrus 1. Epistel 3. Capitel V. 9. vermahnet und spricht: Vergeltet nicht Böses mit Bösem, sondern dagegen segnet, das ist, seid nicht rachgierig, sondern sanftmüthig, und um gemeines und eigenes Friedens willen, vertraget das Unrecht. Und setzet eine gewaltige Ursache dazu, warum solches geschehen soll, nämlich: Und wisset, daß ihr dazu berufen seid, daß ihr den Segen beerbet; das ist, weil ihr von Gott gesegnet seyd in dem gesegneten Samen Abraham, Jesu Christo, daß ihr den ewigen Segen mit seiner Herrlichkeit und Freuden habet und ererbet, ist es billig und recht, daß ihr wiederum segnet, Gutes wünschet und betet, denen, die euch fluchen. Er vermahnet auch dazu mit angehängtem herrlichen Trost, da er saget: Wer ist, der euch schaden könnte, so ihr dem Guten nachkommet? Und ob ihr auch leidet um Gerechtigkeit willen, so seyd ihr doch selig. Als wollte er sagen: Denn eure Widersacher mit ihrer Bosheit und ungerechter Gewalt euch Gottes Gnade, Hülfe und Segen, und eure Frömmigkeit nicht nehmen können. Davon höret auch das 27. Kapitel.

\section{Davids Wohnung in Ziklag}
\quellseite{451}

\section{Summarium}

Ist eine Historia, wie David sich dem Saul nicht zu weit vertraut, sondern aus dem Wege gehet, V. 1. Und nichts desto weniger seinem Volk Israel zu gut die umliegenden Heiden ausrottet und verderbet, welche Gott dem Volk Israel auszurotten geboten hatte, V. 8. Auf daß man lerne sich überwinden, und nichts Böses mit Bösem bezahle. Denn weil Saul dem David nachstellet, und ihm das Volk zu solchem half, hätte David wohl Ursache für der Welt gehabt, sich an ihnen, als seinen Feinden, zu rächen. Aber er thuts nicht, sondern wartet, bis es Gott räche, und thut ihnen dieweil alles Gutes.

\section{Votum}

Der allwissende und gütige Gott erleuchte und regiere uns, daß wir uns in dieser untreuen Welt auch wohl vorsehen, und helfe, daß wir als seine wohlgerathenen und gehorsamen Kinder seinem Exempel nachfolgen, und Freunden und Feinden gern und willig dienen, Amen.

\chapter{Das 28. Capitel}

\section{Vorrede}

Ein jedes frommes Herz und gläubiger Mensch soll täglich gedenken an König Saul, also, daß es sich dabei erinnere des großen Zorns Gottes wider die Sünde, und der erschrecklichen Strafen, die aus Gottes gerechtem Gericht allen Unbußfertigen folgen; und darnach, daß ein jedes fleißig bete, Gott wolle ihn nicht verwerfen unter die Gefäße des Zorns und unter die Landschaden des menschlichen Geschlechts, sondern daß er ihn machen wolle zum Werkzeuge der Gnaden und Barmherzigkeit, ihm, Gott dem Herrn, zu Ehren und dem Nächsten zu Heil und Wohlfahrt. Und vermögen soll aller Christen tägliches Gebetlein sein, daß sie viel und oft mit dem königlichen Propheten David seufzen sollen aus dem 119. Psalm V. 133.: Herr, laß meinen Gang gewiß sein in deinem Wort, und laß kein Unrecht über mich herrschen. Dies lehret auch das 28. Kapitel.

\section[Saul suchet Rath bei einem Zauberweibe, und wird]{Saul suchet Rath bei einem Zauberweibe, und wird vom Gespenst erschreckt}

ᵃ Das erzählet die Schrift darum, auf daß sie warne jedermann, daß er das nachfolgende Beispiel von Samuel recht verstehe, und wisse, daß Samuel todt sei, und solches der böse Geist mit der Zauberin und Saul redet und thut, in Samuels Person und Namen.

ᵇ Das Licht ist, das auf dem Brustlaz des Priesters war.
\quellseite{452}

c Götter) das ist, Richter. 2 Mos. 22, 8. Und der seidene Rock ist der priesterliche Rock.

\section{Summarium}

Ist erstlich ein Exempel, wie Saul immer in größere Sünde fällt. Er ist bisher Gott mehr denn einmal ungehorsam gewesen, hat den unschuldigen David mehr denn einmal umbringen wollen, nach seinem eigenen Sohn hat er einen Spieß geschossen. Sie sehen wir, wie er von Gott abfället und des Teufels Hülfe und Rath bei einer Wahrsagerin suchet, V. 7. Sogar ist kein Aufhören von Sünden, wo man einmal anhebet und nicht Buße thut.

Sonderlich aber ist hie zu merken das Exempel Sauls, welcher, wie hie der Text anzeiget (ohne Zweifel aus Rath und Befehl Samuelis) alle Wahrsager und Zeichendeuter aus dem Lande vertrieben hatte, V. 3. Denn weltlicher Obrigkeit gehörets ihres Amts halben, daß sie solche und andere dergleichen öffentliche Abgötterei und Aberglauben nicht leiden, sondern mit Ernst strafen und abschaffen solle.

\section{Summarium}

Es ist erstlich ein Exempel, in welchem man siehet, wie der Teufel bisweilen kann zukünftige Dinge sagen. Aber nichts desto weniger soll man solcher Weissagung nicht brauchen. Denn es ist eine Sünde wider das erste Gebot, da uns geboten wird, wir sollen uns nur eines Gottes halten, mit Glauben Hülfe und Rath suchen, und allen andern.

Daß aber hie die Schrift sagt, Samuel sei auferwecket und habe mit Saul geredet, V. 15., ist zu verstehen, nicht daß Samuel gewesen sei, sondern daß Saul solches also gemeinet und verstanden habe. Denn so mächtig ist der Teufel selbst nicht, ich geschweige denn ein Mensch, der des Teufels Hülfe braucht, daß er einen verstorbenen Menschen aus seiner Ruhe erwecken und aus dem Schooß Abraham, wie es die Schrift heißet, rücken könnte. Darum ist es gewiß, daß solches nur ein Teufelsgespenst gewesen ist, und
\quellseite{453}
nicht Samuel, der in seiner Ruhe und Schooß Abrahams blieben ist. Die Schrift aber heißt ein solch Gespenst den Samuel darum, daß Saul und die Wahrsagerin es für den Propheten Samuel gehalten haben.

\section{Votum}

Der treue Gott erbarme sich unser, behüte uns für Abgötterei, Zauberei und andern Sünden, und verleihe, daß wir uns allezeit nach seinem Wort und Willen, gegen ihm und gegen allen Menschen christlich erzeigen, bis zu unserm seligen Todesstündlein, Amen.

\chapter{Das 29. Kapitel}

\section{Vorrede}

Wie einer, den die Pestilenz angestoßen, da er ihm in 24 Stunden durch ordentliche Mittel nicht Rath schaffen läßt, sich selbst verwahrloset, daß solch Gift durch alle Adern, Geblüt und ins Herz, hinein dringet, und daß ihm darnach nicht kann geholfen werden, sondern er muß sterben: also müssen auch die, welche der geistlichen Pestilenz, der Sünden, in Zeit der Gnaden und am Tage des Heils nicht zuvor kommen, ewig daran sterben und verderben. Dies lehret St. Paulus an die Galater im 5. Kapitel V. 19. 21. und saget: Offenbar sind die Werke des Fleisches, als da sind Ehebruch, Hurerei, Unreinigkeit, Unzucht, Abgötterei, Zauberei, Feindschaft, Hader, Neid, Zorn, Zank, Zwietracht, Rotten, Haß, Mord, Saufen, Fressen und dergleichen, von welchen ich euch zuvor gesagt und sage noch zuvor, daß die solches thun, werden das Reich Gottes nicht erben. Derwegen laßt uns alle in täglicher Buße leben; und dazu soll uns auch bewegen das Gnaden- und Straf-Exempel. Davon das 29. Kapitel.

David wird von den Philistern zurückgeschickt.

\section{Summarium}

Ist eine feine Historia, weil David mit dem Könige der Philister mußte ziehen, wie Gott Mittel schaffet, daß er mit gutem Willen des Königs wieder abzeucht, V. 4. Nun ist aber kein Zweifel, wo David bei der Schlacht wäre blieben, er würde wider die Philister, mit denen er ausgezogen war, seinem Volk Israel zu gut gestritten haben, wie sie sichs denn auch besorgen, und ihn darum wieder abschaffen. Aber Gott wollte solches Fürnehmen Davids dem Saul und ganzen Israel zu gut nicht lassen kommen, sondern wollte sie um ihre Sünde, wie sie verdienet hatten, strafen. Darum verschafft er solches Mittel, daß David daheim bleibt, und die Philister das Volk Israel schlagen.
\quellseite{454}

\section{Votum}

Der barmherzige GOtt strafe uns ja nicht, wie wir leider auch verdienet haben, mit unsern Sünden; und behüte uns, daß wir uns an niemand versündigen, daß also unser Leben ihm zu Lob und unserm Nächsten zu Nutz diene, Amen.

\chapter{Das 30. Kapitel}

\section{Vorrede}

Die Gottseligen sind selten ohne Kreuz, ja ihr Kreuz wächst und nimmt täglich zu, auf daß auch ihr Glaube, Geduld, Hoffnung und Standhaftigkeit wachse und zunehme; ja zugleich Gottes Liebe gegen sie; denn welchen Menschen der HErr lieb hat, den züchtiget er, Ebr. 12, 6. Je größer aber nun die Züchtigung ist, je größer ist auch die Liebe Gottes; wie der nützliche und fröhliche Ausgang des Kreuzes aller Gläubigen beweiset. Denn die Errettung aus dem Kreuz viel herrlicher ist, denn alles vorgehende Leiden der Gottseligen, auf daß sie also auch ihrem Nächsten denn aus christlicher Liebe und brüderlicher Treue rathen, dienen und helfen können. Denn solches der Liebe Art ist, als St. Paulus lehret 1 Cor. 13, 4. 7., da er sagt: Die Liebe ist freundlich, sie suchet nicht das Ihre, sie lässet sich nicht erbittern, sie trachtet nicht nach Schaden, sie verträget alles, sie duldet alles. Dessen ist ein schön Exempel David, von dem das 30. Kapitel.

\section[David sieget wider die Amalekiter, und theilet den]{David sieget wider die Amalekiter, und theilet den Raub aus}
\quellseite{455}

\section{Summarium}

Ist eine Historia, wie dem David, weil er mit den Philistern wider Israel zeucht, die Feinde einfallen, Weib und Kind sammt aller Habe entführen und die Stadt Ziklag verbrennen. Aus solchem kömmt David zu einem andern Unrath, daß ihn sein eigen Volk will steinigen, V. 6. Aber er hält sich zu seinem Gott, und tröstet sich desselbigen, und jaget aus Gottes Befehl den Feinden nach und erobert alles wieder, V. 10. Auf daß wir lernen, wie Gott mancherlei Anfechtung den Seinen zuschickt, aber auch getreulich wiederum heraus hilft, wenn man sich seiner Gnade und Hülfe tröstet.

\section{Summarium}

Hie ist ein fein Exempel, wie David seinen Leuten zur Versöhnung räth, und Zwietracht durch billige Mittel verhütet, daß er den Raub gleich austheilt. Denn niemand soll Geld oder Gut besser und höher achten, denn Friede und Einigkeit. Wer aber solches nicht thut, und um Geld und Gut willen sich in Unfriede begibt, der kommt endlich auch ums Geld, wie die Exempel täglich für Augen sind.

\section{Votum}

Unser lieber GOtt und treuer Vater helfe heute auch, zu Trost und Rettung allen christlichen Kreuzträgern; und regiere uns, daß wir jedermann in Lieb und Leid, und nach unserm höchsten Vermögen, und auch zu Friede und Einigkeit dienen, Amen.

\chapter{Das 31. Kapitel}

\section{Vorrede}

Gott läßt wahrhaftig die Sünden nicht ungestraft, ob er gleich, als ein sanftmüthiger, barmherziger, gnädiger und langmüthiger Gott, mit der Strafe eine Zeit lang aufseucht und verschonet. Es strafet aber Gott, der Herr, nicht allein Vater und Mutter, die seine Gebote übertreten und sündigen, sondern er suchet deshalb auch ihre Kinder mit Ungnaden heim, daß sie also der Eltern Mißhandlung entgelten müssen. Dieses lehret Gott 2 Mose im 20. Kapitel V. 5., da er sagt: Ich, der Herr dein Gott, bin ein eifriger Gott, der da heimsuchet der Väter Missethat an den Kindern bis in das dritte und vierte Glied, die mich hassen. Demegegen die Eltern gottesfürchtig sein und ehrbar leben sollen, nicht allein, daß Gott sie nicht ihrer Sünden wegen strafen dürfe, sondern auch darum, daß sie ihren Kindern seinen Segen und nicht seinen Fluch und Strafe aufgeben und lassen. Dieses lehret auch das 31. und letzte Kapitel.

\section[Sauls letzter Krieg wider die Philister, darinnen er]{Sauls letzter Krieg wider die Philister, darinnen er sich selbst erstochen}
\quellseite{456}

\section{Summarium}

Ist eine Historia, in welcher man siehet, wie es endlich mit dem Saul ein Ende nimmt, daß er mit seinen drei Söhnen erstochen und das Volk Israel dazu geschlagen wird, V. 4. Darum sollen solche Exempel uns dahin halten, daß wir Gott fürchten und Tag und Nacht mit Fleiß beten, daß er uns für Sünden bewahren wolle, auf daß wir nicht endlich in dergleichen Strafe auch fallen, und so jämmerlich sammt den Unsern umkommen. Denn es ist nicht ein Geringes, daß solcher Jammer dem Saul, der so große und hohe Gaben gehabt, widerfahren, und er so greulich und tief gefallen, und die Strafe über alle seine Nachkommen gangen ist. Der allmächtige, gnädige Gott wolle mit seinem heil. Geist uns leiten und führen, daß dergleichen uns nicht auch begegne.

\part{Das zweite Buch Samuelis}

\chapter{Das 1. Kapitel}

\section{Vorrede}

Gott, der HErr, gönnet der Obrigkeit ein Stück seines Titels und Namens, und nennet sie Götter, als im Ps. 82, 1., und darum, daß sie göttlich Amt führen, und daß er sie will geehret haben. Derwegen die Unterthanen ihrer ordentlichen Obrigkeit nicht allein gebührliche, äußerliche Ehrerbietung mit Gebehrden und Worten erzeigen, sondern sie auch von Herzen in Ehren halten sollen; und wenn die Personen, welche in ordentlichem Amt und Regiment sitzen und den Stand der Obrigkeit führen, als Menschen auch ihre Gebrechen haben und nicht alle Dinge zurecht machen, sollen fromme Unterthanen wegen derselben ihrer Gebrechen Geduld und Mitleiden tragen, und sollen für die Obrigkeit beten, so oft sie zusammen kommen, insgemein und darnach auch insonderheit daheim, und was sie Gott zu Ehren und dem Regiment zum Besten thun, mit Dankbarkeit rühmen und preisen. Davon höret auch an das 1. Kapitel des andern Buchs Samuelis.

\section{Davids Traurigkeit über den Tod Sauls und Jonathans}
\quellseite{457}

\section{Summarium}

Ist erstlich eine Historia, wie David den Amalekiter erwürgen lässet, der ihm ansaget, wie er Saul erstochen hätte, V. 15. Auf daß wir lernen, wie der fromme David so hoch geachtet habe, daß Gott den Saul zum Könige gemacht hatte, und wir demselbigen Exempel nach die Obrigkeit auch lernen hoch halten und uns hüten, daß wir nicht die Hand an sie legen. Denn, so sie unrecht thut, soll man Gottes Urtheil und Strafe erwarten, und nicht wider Gottes Ordnung unser selbst Richter sein und unsern Schaden selbst rächen.

a Bogen) so heißt dies Lied, wie auch bei uns etliche Lieder Namen haben.
\quellseite{458}

\section{Summarium}

Ist ein fein Exempel, wie David den König Saul sammt seinen Söhnen beweinet, und befiehlet, daß man ihrer durch ein öffentlich Lied gedenken soll, als tapferer Helden, welche dem ganzen Lande zu gut ihr Leib und Leben wider die Feinde gewaget, und das Land gern geschützt und Frieden erhalten hätten. Denn dies sind die rechten fürstlichen Werke, welche man an großen Herren preisen und loben, und um derselbigen willen sie und die Ihren ehren und hoch halten soll. Wie es David hie auch nennet und saget, es sei der Schmuck, die schönen Kleider, die güldnen Kleinode, welche der König Saul den Töchtern Israel, das ist, den Städten und dem ganzen Lande habe angehängt, und Leib und Leben endlich darob gelassen. Solche Klage ist dem David ein rechter Ernst, denn er siehet, was für ein Ende folget, wo die Regenten dahin fallen.

\section{Votum}

Der barmherzige GOtt regiere, segne und erhalte unsere liebe Obrigkeit, daß sie christlich und wohl ferner regiere; und verleihe, daß wir sie auch herzlich ehren, ihm zum Gehorsam und uns zu gutem Gewissen, Amen.

\chapter{Das 2. Kapitel}

\section{Vorrede}

Vier Sündenstrafen und Plagen hat Gott in der Schrift namhaftig gemacht. Die erste ist die Pestilenz, die andere ist theure Zeit, die dritte ist Krieg, die vierte ist allerlei böse Wesen, als Löwen, Wölfe, Schlangen und dergleichen. Unter diesen ist Krieg und Unfriede die größte Plage und Strafe Gottes, als die da zerstöret und wegnimmt die Religion, weltlich und häuslich Regiment, daß Theurung und Pestilenz, die auch grausame und schädliche Krebse und Wesen Gottes sein, dennoch nicht zu vergleichen sind mit dem Kriege und was im Kriege vorläuft. Derwegen wer danken kann, der danke Gott täglich und herzlich, daß er unseres Vaterlandes so lange Zeit mit Krieg und andern wohlverdienten Strafen gnädiglich verschonet hat; und laßt uns bitten, daß er uns ferner den lieben Land- und Stadtfrieden erhalten, und nöthige Wohlfahrt zu Leib und Seele geben wolle, Amen. Ferner höret an das 2. Kapitel.

\section{Salbung Davids zum Könige. Jeboseths Aufruhr}
\quellseite{459}

a Helkath-Hazurim) das heißt, das Acker der Besten oder Helden.

\section{Summarium}

Dies ist erstlich ein Exempel, wie David nichts fürnimmt, er habe denn zuvor, nachdem es Gott befohlen hat, bei Gott Rath gesucht und Befehl empfangen, V. 1. Darum ordnet es Gott weiter mit ihm, daß ihn das Haus Juda, nämlich der Stamm, darin er geboren war, zum Könige annimmt und salbet, V. 14. Aber es stößet sich hie noch einmal; denn die andern elf Stämme werfen Isboseth, des Sauls Sohn, zum Könige auf, daraus denn Mord und Krieg sich erhebt, wie allweg geschiehet, wo Uneinigkeit in Regimenten fürfället.

\section{Summarium}

Ist ein fein Exempel, wie Abner, nachdem sich Uneinigkeit zwischen der zweien Königen Unterthanen erhoben hatte, zum Frieden räth und Joab ihm solches gefallen lässet und willig abwendt, V. 28., obgleich sein Bruder Asahel in dem Lärmen von Abner erstochen ward, V. 23. Denn in solchen Sachen soll man allewege mehr auf den gemeinen, denn eigenen Schaden sehen, und ums gemeinen Friedens willen des eigenen Schadens gern vergessen.

\section{Votum}

Unser lieber Herr Gott verleihe uns Friede gnädiglich zu unsern Zeiten, und regiere uns, daß wir ein geruhiglich und stilles Leben führen in aller Gottseligkeit und Ehrbarkeit, und in unserm Beruf all unser Thun und Lassen nach seinem Wort anstellen, Amen.
\quellseite{460}

\chapter{Das 3. Kapitel}

\section{Vorrede}

Gleich wie Gott den Verzug der Strafe pfleget mit Schwerheit der Strafe einzubringen, daß ist, je länger Gott verzeucht, die Gottlosen zu strafen, je greulicher er sie darnach, wenn er seine Zeit ersiehet, zu strafen pflegt: also pflegt er auch den Verzug seiner Hülfe einzubringen, und je länger er uns zu helfen verzogen hat, je reichlicher er uns, wenn er seine Zeit ersiehet, zu helfen pfleget. Derwegen sollen wir den Verzug der Hülfe Gottes nicht also verstehen, als wollte Gott uns ganz und gar nicht helfen, und sollen in keinem Wege deshalb das Gebet, Glauben und Hoffnung fahren lassen, sondern sollen vielmehr also thun, wie uns der Prophet Habakuk vermahnet Kapitel 2. V. 3.: Wenn der Herr verzeucht, so hoffe und warte auf ihn, er wird kommen und nicht außen bleiben. Denn wie David sagt Ps. 147, 11.: Der Herr hat Gefallen an denen, die ihn fürchten und auf seine Güte warten. Und St. Paulus spricht Röm. 5, 5.: Die Hoffnung läßt nicht zu Schanden werden. Dies lehret auch das 3. Kapitel.

\section{Abner von Joab umgebracht}

\section{Summarium}

Die ist eine Historia, wie es sich anspinnet, daß des Sauls Nachkommen gar zu Grunde gehen, V. 9., auf daß man lerne, es verziehe sich mit der Strafe so lange es wolle, so bleibt sie doch endlich nicht außen. Dergleichen geschieht auch mit den Frommen. Denn hie siehet man, daß sich etliche Jahr mit David verzeucht, ehe er zum Regiment kommt, daß es scheint, als würde nicht allein gar nichts daraus werden,
\quellseite{461}
sondern er müßte noch auch drüber zu Boden gehen. Aber er hält fest an Gottes Verheißung; darum muß es auch endlich ins Werk kommen und gehen, wie es Gott beschlossen hatte, daß die Bösen gestraft und die Frommen gefördert werden.

\section{Summarium}

Ist eine Historia, in welcher man abermals siehet, wie Gott sein Gericht dermaßen führet, daß endlich die Unguten gestraft werde. Es scheinet aus dem 26. Capitel V. 5. des 1. Buchs Samuels, daß Abner auch dazu geholfen habe, daß David ist verfolget worden; nach Sauls Tode aber ist Abner eine Ursache gewesen, daß Israel zertrennet und der meiste Theil David zum Könige nicht wollen annehmen. Auch verkündiget sich Abner in dem, daß er mit Rizpa, Sauls Kebsweib, Hurerei treibet. Solches und anders wird zuletzt gestraft, daß er unversehens von Joab erstochen wird, V. 27. Auf daß jedermann lerne in Gottes Furcht wandeln, so er der Strafe will überhoben sein. Denn weil Gott ein gerechter Richter ist, kann die verdiente Strafe nicht außen bleiben, es verhalte und verziehe sich gleich, so lange es wolle.
\quellseite{462}

\section{Votum}

Helfe der gerechte, gütige GOtt, daß wir das Böse fliehen und das Gute thun, auf daß er die Sünde an uns nicht strafen dürfe, und daß er das Gute aus Gnaden belohne, um Jesu Christi, seines geliebten Sohnes willen, Amen.

\chapter{Das 4. Kapitel}

\section{Vorrede}

Unser HErr JESUS Christus sagt Matth. 26, 52.: Wer das Schwert nimmt, der soll durchs Schwert umkommen. Hiermit lehret er, daß alle, die das Schwert gebrauchen, und wird ihnen nicht von Gesetzen oder von der Obrigkeit gegeben, daß sie als muthwillige Todtschläger, Mörder und Aufrührer wieder sollen getödtet werden. Derwegen hüte sich jedermann für Aufruhr, und niemand gebrauche das Schwert, andern zu schaden, ohne Beruf und ohne Willen der Obrigkeit, dergleichen ohne rechtmäßige Ursachen. Denn wer über sich haut, dem fallen die Späne in die Augen, und ist nie keiner der verdienten Strafe entronnen, der sich seiner, von Gott verordneten, Obrigkeit widersetzt hat. So ist die Geduld die beste Arznei, daß man Gott die Sache und Rache befehle, und wisse, ob man dies, was in Unrecht geschieht, gleich um die Obrigkeit nicht verdienet, so hat man es doch bei Gott wohl und oft verdienet. Davon lehret auch das 4. Kapitel.

\section{Jeboseth wird ermordet, und von David gerochen}

\section{Summarium}

Hie ist eine Historia, wie die zween Hauptleute Baena und Rechob des Sauls Sohn Isboseth, ihrem Herrn, mörderisch erstechen, V. 5. 6. und meinen, sie wollen David wohl damit zu Hofe kommen, V. 8. Aber David, als einer, dem das Schwert von Gott befohlen war, läßt solchen verrätherischen Mord nicht ungerochen, sondern heißt sie auch erstechen, V. 12. Also siebet man hie abermal, daß keine Sünde ungerochen bleibt. Sonderlich aber gebet dies Exempel wider die Aufrührer, so an ihre Obrigkeit die Hand legen. Denn solches soll nicht ungerochen bleiben.

\section{Votum}

Helfe GOtt, daß wir uns unterthänig in allem christlichen Gehorsam gegen unsere Obrigkeit erzeigen, und alle aufrührerische Anschläge und dergleichen alle andere Sünde, wider GOtt und den Nächsten, fliehen und meiden, in Kraft des heiligen Geistes, Amen.
\quellseite{463}

\chapter{Das 5. Kapitel}

\section{Vorrede}

Zu Gott fliehen unter dem Kreuz ist das Beste; und da soll ein betrübt Herz, wenn es ihm übel geht, also gedenken: Siehe, jetzt habe ich das Unglück, welches mich zwar sehr ängstet und martert; aber was ist es zu rechnen gegen die vielfältigen unzähligen Wohlthaten Gottes, welche er mir sonst mein Lebenlang bis hieher erzeiget hat, und noch erzeiget? Er hat mir ja mehr Gutes gethan, denn mir alle Teufel, die ganze Welt und böse Leute Arges anthun können. Habe ich Gutes empfangen von Gott, und sollte das Böse nicht auch annehmen? Und soll also ein jeder Kreuzträger seine Seele mit Geduld fassen und Gott, bis er zu seiner und rechten Zeit hilft, aufwarten. Dieses sind auch des Königs Davids Gedanken und Geduld gewesen in seinen mancherlei Nöthen; derwegen hat ihm Gott gnädiglich und wunderlich wieder geholfen, wie auch bezeuget das 5. Kapitel.

\section{Davids andere Salbung, glücklicher Zustand und zweyfacher Sieg}

a Diese Blinden und Lahmen sind ihre Götzen gewesen, welche sie zu Trost wider David auf die Mauren setzten, als ihre Patronen, die sie schützen sollten. Wie man sieht auch mit der Heiligen Bildern thut. Als wollten sie sagen: Du kriegest mit uns nicht, sondern mit unsern Göttern; beiß dich mit ihnen, sie werden dir wohl wehren. Nicht, daß sie Blinde oder Lahme gewesen haben, sondern der Geist Gottes in denen, so hernach die Historien beschrieben haben, heißt sie also. Ps. 115, 4. 5.

b Derer heißt ein Riß oder Klag, daher diese Stätte Baal Perazim Rißmann heißt, daß die Philister da gerissen sind.
\quellseite{464}

\section{Summarium}

Ist eine Historia, wie David vom ganzen Israel zum Könige gesalbet und angenommen wird, V. 3.; auf daß man sehe, ob gleich Gott mit seiner Verheißung lange verzieht, daß er doch endlich nicht außen bleibet. Darum soll man harren und nicht müde werden.

Daß aber David so viel Weiber nimmt, V. 13., geschieht nicht aus Fürwitz, wie es die Welt deutet, sondern der Verheißung halben, daß er wußte und glaubte, Christus würde aus dem Stamme Juda kommen. Neben dem ist auch wohl zu merken, wie er alle Dinge mit Gott, und nicht aus eigenem Gutdünken anfähet; darum ist auch Glück und Heil bei ihm, daß er seinen Feinden obliegt, V. 19. So ist es wohl zu glauben, da er in solche Hoheit kommen, daß viel armer Freunde sich zu ihm gehalten und die nächsten haben wollen sein, weil die Juden ein sonder Gesetz hatten, zu den nächsten Freunden zu heirathen, wo ohne leibliche Erben der Ehemann abgestorben war.

\section{Votum}

Gott der Allmächtige erhalte uns auch in wahrem Glauben an seine Verheißungen, und helfe uns an Leib und Seele, daß wir den Unsern und andern nothleidenden Christen wieder helfen können, in leiblicher und geistlicher Noth, Amen.

\chapter{Das 6. Kapitel}

\section{Vorrede}

Gott wendet doch endlich sein gnädig Angesicht zu den Seinen, obschon der Teufel mit seinem häßlichen schrecklichen Rachen, und der Tod mit seinem grausamen Anblick, und die Welt mit feurigen, zornigen Augen uns ansehen, anmurren und wider uns toben. Derwegen wer in wahrer Buße lebt und betet, und hat Gottes gnädige Augen, der darf sich für des Teufels, des Todes und der Welt häßlichem, erschrecklichem Gesicht nicht fürchten. Denn es heißt, wie St. Petrus sagt 1. Epistel Kapitel 3. V. 12.: Die Augen des HErrn sehen auf die Gerechten, und seine Ohren auf ihr Gebet. Das ist, die gläubigen Kreuzträger siehet Gott an mit gnädigen, lachenden Augen, und höret mit leisen, offenen Ohren ihr Klagen, Seufzen und Bitten, und errettet sie endlich mit starker Hand aus aller Noth. Das lehret Davids Exempel, von dem das 6. Kapitel.

\section{Davids Freudentanz bei Abholung der Bundeslade}
\quellseite{465}

\section{Summarium}

Ist ein fürstlich Exempel, wie David sich nicht allein des Regiments, sondern auch der Kirchen und Gottesdienstes annimmt und mit allem Fleiß und Ernst fördert, V. 2., und ist so herzlich fröhlich darüber, daß er die Lade Gottes aus dem Hause Abinadab ließe holen, daß er vor großer Freude vorher tanzet, V. 15. Solches legt ihm sein Weib Michal übel aus, V. 16. Wie noch heutiges Tages viel sind, die es für eine Unehre achten, daß große Herrn sich der Lehre und des Gottesdienstes ernstlich annehmen. Denn die Welt hält solches für eine unnöthige Fahr, der man wohl könne gerathen. Aber David sagt, er begehre noch geringer zu werden für dem HErrn, denn also werde er recht zu Ehren kommen, V. 21. Und meldet der Text sonderlich, daß die hoffärtige Michal von Gott gestraft sei, daß sie hernach ihr Lebenlang kein Kind habe gehabt, V. 23. Daher gehören auch die zwei Exempel: Usa, weil er die Lade Gottes gering hielte, und mit dem Anrühren, daß er nicht Befehl hatte, unehret, schlug ihn Gott, daß er starb. Obed-Edom, der sie ehret und herberget, ward gesegnet. Alles uns zum Exempel, daß wir uns Gottes und seines Worts nicht schämen, sondern dasselbige in Ehren halten, und uns lieb lassen sein.

\section{Votum}

Wir danken GOtt, der als ein gütiger Vater das Gute belohnet und das Böse als ein rechter Richter strafet; der lehre uns durch sein Wort und regiere uns durch den heiligen Geist, daß wir alle Sünden fliehen, und allezeit Gutes wirken, Amen.

\chapter{Das 7. Capitel}

\section{Vorrede}

Ach ein jeder nehme sich der christlichen Kirche an, so viel, als er Amts, Standes und Vermögens halben kann und soll, auf daß dieselbe mit Lehre und Gottesdienst allzeit recht und christlich bestellet sei. Denn welcher Mensch ihm läßt befohlen sein der Kirchen Wohlfahrt, und fördert die mit Beisteuer, Almosen, gutem Rath, herzlichem Gebet und andern treuen Handreichungen, den wird Gott also segnen, daß er zeitlich und ewig wird dafür zu danken haben. Denn also lautet davon die wahre Verheißung unsers HErrn Jesu Christi, Matth. 10, 42.: Wer dieser Geringsten einem nur mit einem Becher kaltes Wassers tränket in eines Jüngers Namen, wahrlich, ich sage euch, es wird ihm nicht unbelohnet bleiben. So ist dies auch gewiß, ein jeder Christ kann in seinem Beruf ja etwas der Kirchen Gottes dienen und helfen. Davon erinnert auch das 7. Kapitel.
\quellseite{466}

\section[David bekommt die Verheißung des Messias, als er]{David bekommt die Verheißung des Messias, als er Willens ist, einen Tempel zu bauen}

\section{Summarium}

Das ist ein fürstlich Exempel, wie David sich mit Ernst annimmt, daß er den Gottesdienst fördert und für unbillig achtet, daß er in einem Cedern-Hause, die Lade GOttes aber unter den Teppichen wohnen soll. V. 2. Solche schöne und feine Gedanken zeigen an ein gottseliges Herz, daß er Gottes Ehre mit allem Fleiß gesucht und gerne gefördert hat. Aber Gott läßt ihm durch Nathan sagen, er wolle seinem Samen, der von seinem Lenden kommen soll, nach ihm erwecken, der seinem Namen solle ein Haus bauen, V. 2. Das ist nicht von Salomo, sondern von Christo geredet, der ihm hie verheißen wird, daß er aus seinem Samen kommen sollte, und ein ewig Reich der Gnaden anrichten. Aus diesem Text hat David den schönen herrlichen Psalm, an der Zahl der 89., V. 5. und 30. und andere mehr gemacht.

a Das ist, so redest mit mir von solchem ewigen Reich, da niemand kann König sein, er muß Gott und Mensch sein, weil er mein Sohn und doch für und für soll König sein, welches allein Gott gebühret.
\quellseite{467}

\section{Summarium}

Es ist ein sehr schön Exempel, wie David die herrliche Wohlthat erkennet und Gott dafür so herzlich danket, daß er nicht allein ihm bisher so lange viel Gutes gethan, sondern auch aufs künftige ihm Christum verheißen. Solche unaussprechliche Gnade preiset David und bekennet, daß ers nicht verdienet, sondern Gott aus lauter Gnaden um seiner Verheißung willen gethan habe, V. 21. Und bittet, V. 25., Gott wolle solcher seiner Verheißung um seiner Güte willen nachkommen. Uns zum Exempel, daß wir auch sollen lernen, erstlich erkennen, daß alles, was uns GOtt für Gutes thut, aus lauter Gnaden thue, und wir hernach Gott dafür danken, und ferner bitten sollen, daß er bei uns sein und uns weiter erhalten, und für Sünden behüten wolle.

\section{Votum}

GOtt helfe, daß die Obrigkeit auch heute seine großen väterlichen Wohlthaten erkenne, und regiere sie, daß sie die wohl und glückselig zu seiner Ehre und den Unterthanen zu Nutz allzeit recht gebrauche, hoch gelobt in alle Ewigkeit. Amen.

\chapter{Das 8. Kapitel}

\section{Vorrede}

Die Weisheit Gottes saget in Sprüchen Salomonis im 1. Capitel V. 32.: Daß die Albernen gelüstet, tödtet sie, und der Ruchlosen Glück bringet sie um. Wer aber mir gehorchet, wird sicher bleiben und genug haben, und kein Unglück fürchten. In diesen Worten wird auch folgende Lehre begriffen: Wenn ein Mensch am sichersten ist, so ist ihm allerlei Unglück und sein Untergang am nächsten. Dagegen, wenn ein Christ am schwachgläubigsten ist, und trägt das größte Kreuz, wo er nur bei Gott mit dem Gebet anhält, so ist ihm auch am nächsten seine Erlösung und Wohlfahrt. Dieses soll schrecken die Sichern und Unbußfertigen, daß sie wachen und beten, ihre Lenden umgürtet haben und ihre Lichter brennen lassen, Luc. 12, 35. Und soll trösten die Blöden und Zaghaften, und daß sie auch wachen, im Glauben stehen, männlich und stark sein, 1 Cor. 16, 13. Davon erinnert auch das 8. Kapitel.

\section{David sieget, und bestellet die Aemter}
\quellseite{468}

\section{Summarium}

Ist eine Historia, wie Gott den frommen David, nachdem er so lange im Elend gewesen und mancherlei Unglück versucht hat, gnädiglich bedenket und viel größere, herrlichere Siege gibt wider seine Feinde, V. 1. 2. Zum Exempel aller Obrigkeit, wenn sie Glück und Heil will haben, daß sie es von Gott erwarten und sich gegen Gott dermaßen halten soll, daß er nicht Ursache habe, sie zu strafen. Darum meldet der Text zweimal: Der Herr half David, wo er hinzog, V. 7. 14. Als wollte er sagen: Wo Gott mit seiner Gnade und Hülfe bei David nicht gewesen wäre, würde er mit seiner Weisheit, Sorge und Gewalt wenig ausgerichtet haben.

\section{Votum}

Der König aller Könige, unser lieber Gott, wolle heute auch aller christlichen Obrigkeit zu Schutz und Hülfe beistehen, und wolle sie regieren, daß sie sich auch ferner also halte, auf daß sie von Gott zur ewigen Herrlichkeit erhalten und erhaben werde, Amen.

\chapter{Das 9. Kapitel}

\section{Vorrede}

Wer Gott fürchtet, der stehet nicht allein wohl bei Gott für seine Person, sondern seine Kinder, Kindeskinder und Nachkommen werden auch von Gott gesegnet, das ist, sie werden für Gott groß und bei allen Engeln und Heiligen hoch geachtet, zu großen Ehren und Würden erhoben; es gibt ihnen auch Gott zeitliche Nothdurft, also, ob sie es gleich nicht haben schlechtlich, so haben sie es doch leidlich, und ob es schon nicht fleußt, so rinnet es doch, und das Wenige segnet ihnen Gott zu ihrem und der Ihren Genügen und Auskommen. Derwegen spricht David im 112. Psalm V. 2.: Das Geschlecht der Frommen wird gesegnet sein. Und sein Sohn Salomon sagt in Sprüchen im 11. Kapitel V. 21.: Der Gerechten Same wird errettet werden. Deß ist ein Exempel das Geschlecht Jonathan, davon das 9. Kapitel.

\section{Gutthätigkeit Davids gegen Mephiboseth}
\quellseite{469}

\section{Summarium}

Ist ein schön Exempel, wie David, nachdem er das Regiment zusammen gebracht und gestillet war, die Nachkommen Sauls bedenkt um Jonathans willen, V. 1., der alle Treue an ihm gethan hatte, V. 2. Denn Gottes Segen soll man niemand zu Schaden, sondern jedermann zu Nuß brauchen, auch den Feinden, wenn sie sich erkennen und Gnade begehren.

\section{Votum}

Gott verleihe uns und den Unsern auch wahre Frömmigkeit, und segne zeitlich und ewiglich alle, die an den Seinen Treue thun, und die schützen, erhalten und fördern, Amen.

\chapter{Das 10. Kapitel}

\section{Vorrede}

Wer auf Gott trauet in billigen guten Sachen, der bauet wohl und also, daß der Teufel und die Welt nicht schaden und den Bau nicht einreißen können. Derwegen ein frommer gottfürchtiger Mensch, er sei so geringe, verachtet, elend und arm in der Welt, als möglich ist, wenn er gleich auch der elendste Bettler wäre, dennoch ist er bei Gott größer und glückseliger, und auch wider den Teufel und die Welt gewaltiger und mächtiger, denn alle Gottlosen, wie groß, hoch, ansehnlich, stark und muthig sie auch sind. Derwegen wir als christliche geistliche Kriegsleute, wenn wir streiten und kämpfen, nicht nachgeben sollen, da wir nach dem Krieg wollen gekrönet werden. Sintemal Kraft und Stärke uns mittheilen will Gott, der Herr, der unaussprechliche Schätze und Güter bereitet hat im ewigen Leben denen, die da überwinden. Davon in der Offenbarung Johannis im 2. Kapitel V. 7. also gesagt wird: Wer überwindet, dem will ich zu essen geben von dem Holz des Lebens, das im Paradies Gottes ist. Und abermals V. 11.: Wer überwindet, dem soll kein Leid geschehen von dem andern Tode. Darauf höret ferner an das 10. Kapitel.

\section{Die Ammoniter verhöhnen Davids Knechte, und werden geschlagen}
\quellseite{470}

\section{Summarium}

Ist eine Historia, wie die Gottlosen, ob sie gleich bisweilen anderer Leute halben wohl möchten Friede haben, dennoch ihnen selbst Unglück und Unruhe schaffen müssen, auf daß die Sünde ja zeitlich und wohl gestraft werde. Also hätte der Ammoniter König David wohl mögen zum freundlichen Nachbarn haben; aber er und die Seinen waren's nicht werth. Darum müssen sie nach ihrem eigenen Unglück grübeln und sich zum David nöthigen, daß sie gestraft und geschlagen werden.

\section{Votum}

Steure und wehre der starke und gerechte GOtt allen, die zu Krieg und Unruhe Lust haben, und verleihe auch unserer lieben Obrigkeit Glück und Sieg wider alle ihre Feinde, und erhalte uns den heilsamen Land- und Stadtfrieden, Amen.

\chapter{Das 11. Kapitel}

\section{Vorrede}

Der heiligen Bibel fürnehmste Kapitel sollen ihnen fromme Herzen wohl bekannt machen, daß ihnen ja an nöthiger Lehre, Trost und Warnung zu keiner Zeit mangele. Ein solch Kapitel ist das 11. des andern Buchs Samuelis, welches lehret, daß wir Menschen sehr schwach und der Teufel gar stark und listig ist, der bald einen Menschen fällen kann und zu einem großen Sünder machen. Es vermahnet auch, daß niemand in seinen Sünden verzweifeln, sondern Buße thun und bei Gott Gnade suchen solle, der werde ihm Gnade erzeigen und alle seine Sünde vergeben; wie solches der Prophet Moses von Gott rühmet, 2 Mos. 34, 6., und sagt: Herr, Herr, Gott, barmherzig und gnädig, und geduldig, und von großer Gnade und Treue, der du vergießest Missethat, Uebertretung und Sünde. Wie auch dieses lehret Sirach im 18. Kapitel V. 28. und spricht: O wie ist die Barmherzigkeit des Herrn so groß, und läßt sich gnädig finden denen, so sich zu ihm bekehren. Das gemeldete Kapitel lautet, wie folgt.

\section{Davids Ehebruch und Todtschlag}
% Scan page 471: Bible verses/narrative, running heads, folios and bare cross-references only; no non-biblical apparatus retained.
\quellseite{472}

\section{Summarium}

Die ist eine sonderliche Historia, wie David sich so jämmerlich versündiget, daß er nicht allein ein Ehebrecher wird, V. 4., sondern auch bestellet, daß der unschuldige fromme Ritter Urias, da er ihn mit Güte nicht dazu vermögen könnte, daß er sein Schanddeckel gewesen wäre, erschlagen wird, und andere mit ihm, V. 17. Dadurch denn Ursache gegeben wird den Ammonitern, Gott zu lästern. Solch Exempel soll uns erstlich eine Warnung sein, daß wir nicht sicher sein, sondern stets beten sollen, daß uns Gott für Anfechtung und Sünde behüten wolle. Zum andern soll es uns auch schrecken, daß wir dergleichen Sünde nicht begehen, weil sie Gott so greulich strafet, wie hernach wird folgen.

\section{Votum}

Der treue, barmherzige Gott steure dem boshaftigen Menschenfeinde, dem Teufel, behüte uns für Sünden, richte auf, tröste und erhalte die Gefallenen, daß sie ihn mit uns loben und preisen, auch selig werden, um Jesu Christi, seines geliebten Sohns, unsers Herrn willen, Amen.

\chapter{Das 12. Kapitel}

\section{Vorrede}

Der königliche Prophet David sagt im 84. Psalm V. 5. 11.: Wie lieblich sind deine Wohnungen, Herr Zebaoth. Wohl denen, die in deinem Hause wohnen; denn ein Tag in deinen Vorhöfen ist besser, denn sonst tausend. Ja freilich ist nichts besseres, nichts herrlicheres, nichts tröstlicheres, nichts nützlicheres, denn ein wahrhaftiger Bürger und lebendiges Gliedmaß sein der rechten und wahren Kirchen Christi, da Gott wohnet und läßt sein Wort immer fort und fort erschallen, welches der rechte lebendige Brunnen und frische Quell ist, damit der Herzen oder Seelen Durst gestillet und Seele und Leib in Anfechtung erquicket werden. Dafür sollte keiner türkischer Kaiser werden, so es in seiner Willkür stünde. Sintemal eine einige rechte Predigt und Lection Gottes Worts, darin das Herz getröstet und getroffen wird, besser ist, denn aller Welt Kleinodien, Schätze, Wollust und Freude. Das hat David gelehret und recht erfahren in seinem Elende, als davon zeuget das 12. Kapitel.

\section{Nathans Bußpredigt}

\section{Summarium}

Sie ist eine Historia, wie Gott durch den Propheten Nathan David um seine begangene Missethat strafet, V. 2. Denn gleich wie Gott dem Predigtamt befohlen hat, die Sünde zu vergeben, wo mans erkennt und Gnade begehrt,
\quellseite{473}
also hat er auch befohlen, die Sünde zu strafen, wo man’s nicht erkennet. Und hie siehet man, wie greulich Gott den Ehebruch pflegt zu strafen. Sintemal der Prophet dräuet, Gott werde solche Sünde Davids mit Aufruhr strafen, V. 11.

a Jedidia) heißt lieblich dem HErrn.

\section{Summarium}

Ist endlich ein Exempel der rechtschaffenen Buße, nämlich, daß man die Sünde bekennen und um Christus willen Gnade begehren und hoffen soll. Da will Gott Gnade widerfahren lassen, wie der Prophet zu David spricht: Du wirst nicht sterben, der HErr hat deine Sünde weggenommen, V. 14.

Zum andern ist auch hie ein Exempel, daß sonderlich in die Haustafel gehöret, nämlich, daß bisweilen die Kinder der Eltern Sünde entgelten müssen. Also mußte das Kind darum sterben, daß David die Ehe gebrochen hatte, V. 15. Auf daß die Eltern lernen ihrer Kinder halben desto gottsfürchtiger und frömmer sein und sich für Sünden hüten, welche Gott auch an den Kindern bis ins dritte und vierte Glied will strafen, mit Krankheit, Armuth und andern Unglück.

\section{Votum}

Wir danken GOtt herzlich, daß wir auch das Strafgesetz und Trostevangelium haben; der verleihe Hülfe und Gnade, daß wir beide zu christlichem Leben und beständigem seligen Trost allzeit gebrauchen, bis zum seligen Todesstündlein, Amen.
\quellseite{474}

\chapter{Das 13. Kapitel}

\section{Vorrede}

Mancher meinet, es sei alles verloren, was er gutes und christliches thut, wenns ihm in diesem Leben nicht nach seinem Willen vergolten wird; mancher meinet auch, es sei alles erlogen, was Gott der Bosheit und Sünden geträuet, wenns nicht hie zeitlich, an allen nach ihrem Verdienst, gestraft wird. Aber wer mit den Augen recht siehet, mit den Ohren recht hört und mit dem Herzen, was recht ist, versteht, der wird sehen, hören und befinden, daß, was Gottes Wort lehret und desselben Dräuungen endlich erfüllet werden, als auch dieses, daß der eifrige Gott über die so ihn hassen, die Sünde der Väter heimsucht an den Kindern bis ins dritte und vierte Glied, welches er selber träuet 2 Mos. 20, 5. Dies bezeuget auch das 13. Kapitel.

\section{Der Blutschänder Amnon wird von Absaloms Dienern ermordet}

*Ledig) das ist, sie hielt sich immer, ging nicht mehr wie die Jungfrauen, im Hause unter die Leute.
\quellseite{475}

\section{Summarium}

Hier gehet die Historia an, wie dem David, da er im Regiment Friede und Ruhe hatte, seine eigenen Kinder eine Anfechtung und Unglück über das andere machen, V. 1. Auf daß wir lernen, daß solch Kreuz mit ungerathenen Kindern bisweilen auch frommen Eltern widerfähret, da es doch an Zucht und fleißigen Vermahnungen nicht gefehlet hat. Darum bedarf es nicht allein großen Fleiß, daß man die Kinder in der Zucht und Vermahnung in dem Herrn aufziehe, sondern auch stetiges Beten, daß sie wohl gerathen.

Neben dem stehet man hie, daß der Teufel durch die schändliche That, daß Ammon seine Schwester Thamar schwächte, gedenkt Mord zu stiften zwischen Brüdern, wie hernach geschah, welches David als ein Vater ein schwer Kreuz und Herzeleid ist gewesen, wie weiter folgen wird, und wir in täglicher Erfahrung auch sehen.

\section{Summarium}

Ist das andere Unglück, daß David an seinen Kindern erlebet. Nachdem Ammon seine Schwester Thamar geschwächt hatte, erwürget ihn darum sein Bruder Absalom, welches rechte Schwester die Thamar war. Sie mögen wir lernen, wie endlich Gott Hurerei strafet. Denn obwohl Absalom Unrecht thut, so hat dennoch Ammon solche Strafe wohl verdienet.

Darum soll man sich für Hurerei und Unzucht hüten, und die Eltern nicht allein die Kinder mit Fleiß davon abziehen, und allerlei Ursache, mit Gesellschaft zechen, reden, und was mehr sein kann, verhüten, sondern auch mit Ernst beten, daß Gott für solchen gefährlichen Teufels-Stricken ihre Kinder behüten wolle, da allwege Strafe und Herzeleid darauf folget.
\quellseite{476}

\section{Votum}

Gott, der HErr, regiere alle Eltern, daß sie durch Sünde ihnen und ihren Kindern Gottes Zorn und Strafe nicht auf den Hals ziehen, und erleuchte alle liebe Kinder, daß sie mit ihrem Sündenleben ihren Eltern nicht Herzeleid machen, und sich ins Verderben führen, Amen.

\chapter{Das 14. Kapitel}

\section{Vorrede}

Die Eltern strafen mit Ruthen, die Prediger mit scharfer Gesetzespredigt, die weltliche Obrigkeit aber gebraucht das Schwert, das ist, allerlei Strafe, damit sie ihren Ernst gegen die Ungehorsamen erzeigen, an derselben Ehre, Leib und Gut. Nun sagt aber St. Paulus Röm. 13, 4.: Die Obrigkeit trägt das Schwert nicht umsonst. Das ist, die Obrigkeit soll nicht allein das Schwert in ihrer Verwahrung im Kasten liegen, oder an der Wand hängen lassen, sondern sie soll es tragen, das ist, brauchen und nicht feiern lassen, und soll es tragen nicht umsonst oder vergebens, allein zur Schau oder zur Zierde, daß es also solle wohl stehen, auch nicht ohne Unterschied, sondern die Frommen damit zu schützen und die Bösen damit zu strafen. Solcher Regenten etliche hat Gott allzeit gegeben; ihm sei Lob, Ehre und Dank dafür; und ein solcher ist auch gewesen der König David, von dem das 14. Kapitel.

\section{Absalom findet auf Fürbitte des Weibes von Thekoa Gnade}
\quellseite{477}

\section{Summarium}

Sie ist ein fein Exempel, wie David einen sondern Ernst gebraucht habe wider seinen Sohn Absalom, um der Mißthat willen, daß er seinen Bruder Amnon todt geschlagen hatte, V. 1. Und scheint, wo David von dem klugen Weib von Thekoa nicht also hinterstichlichen und ihm die Zusage verdeckter Weise wäre abgedrungen worden, V. 12., daß Absalom hätte sein Lebenlang in der Fremde und Elend müssen bleiben. Solches thut David nicht allein als ein Vater, sondern auch als ein Regent. Denn denselbigen gebührt sonderlich, daß sie keine Freundschaft ansehen, sondern das Uebel sowohl an Freunden als an Fremden ernstlich strafen sollen.

\section{Summarium}

Ist wieder ein Exempel des sondern Ernstes, so David wider seinen Sohn Absalom gebraucht hat, daß, nachdem er aus Fürbitte wiederkommen ist, dennoch David ihn nicht will für sein Angesicht kommen lassen. Wo die Obrigkeit dergleichen Ernst wider die Freunde brauchte, würden die andern desto mehr Scheu haben und sich desto steifer für aller Untugend hüten. Denn es ist sehr ärgerlich und schädlich in Regimenten, wo man mehr auf Freundschaft, denn aufs Recht siehet. Denn da begibt sichs, daß man auch Fremde muß schonen und die Untugend an ihnen nicht strafen darf. So doch Gott der Obrigkeit darum das Schwerdt in die Hände gibt, daß sie es wider alles Aergerniß brauchen und wehren soll, es sei Freund oder Feind, Bekannte oder Unbekannte. Denn da soll es alles gleich gelten, nur daß Gottes Wille geschehe und dem Aergerniß mit Ernst gewehret werde.
\quellseite{478}

\section{Votum}

Gott steure allen Aufrührern und Ungehorsamen, und regiere unsere liebe Obrigkeit, daß sie ferner alles Böse ernstlich strafe, und die Frommen treulich schütze und fördere; verleihe uns allen auch, daß wir der Obrigkeit schuldigen, christlichen Gehorsam allzeit erzeigen, Amen.

\chapter{Das 15. Kapitel}

\section{Vorrede}

Beim heiligen Propheten Micha am 7. Kapitel V. 9. sagt das nothleidende israelitische Volk Gottes: Ich will des Herrn Zorn tragen, denn ich habe wider ihn gesündiget, bis er meine Sache ausführe und mir Recht schaffe. Er wird mich ans Licht bringen, daß ich meine Lust an seiner Gnade sehe. Das ist so viel gesagt: wir haben ja wohl verschuldet und verdienet die Noth, die uns betreten hat. Denn wir sind für Gott schuldig und der Strafe würdig; aber er wird uns, die wir im Finstern, aller Trübsal und Angst sitzen, wiederum aushelfen, unser Elend wenden und enden, daß uns seine Güte und Barmherzigkeit gleich als ein neu Licht wird aufgehen, und wir aus allem Herzeleid errettet, unsern Trost und Freude an seiner gnädigen, väterlichen Hülfe sehen und haben werden. In solcher Geduld, Trost und Gehorsam soll auch heute ein jeder Mensch sein Kreuz, das er mit seinen Sünden nur wohl verdient hat, tragen, als denn auch gethan hat David, von dem das 15. Kapitel.

\section{David fliehet vor seinem aufrührischen Sohn Absalom}

a Diese vierzig Jahr rechnen wir von der ersten Salbung Davids; denn er ist zwanzig Jahr alt gesalbet, und zehn Jahr Verfolgung gelitten, und im dreißigsten Jahr König beständig und angenommen.
\quellseite{479}

\section{Summarium}

Dies ist das dritte Unglück, das David von seinen Kindern widerfähret: daß Absalom, sein Sohn, einen Aufruhr wider ihn erwecket, ihn des Landes verjaget, daß er aus Jerusalem weichen und für ihn fliehen muß, V. 14. Solches Unglück ist eine Strafe des Mordes und Ehebruchs, welchen David an Uria und seinem Weibe begangen hat, wie droben der Prophet Nathan geweissaget hatte. Darum soll jedermann, sonderlich aber die Obrigkeit, sich für solchen Sünden hüten, und dies Exempel wohl merken. Denn obgleich Gott solche Sünde dem David vergeben hatte, daß sie ihm am ewigen Leben nicht schaden und er wiederum in Gottes Gnade sollte sein, darum, daß ers ihm ließe leid sein, und glaubet der Zusagung Gottes; dennoch mußte das Aergerniß, welches er damit gegeben, öffentlich gestraft und David jedermann zum Exempel werden, daß andere für dergleichen Sünde sich hüten lernten. Insbesondere wird hie sein angezeiget, wie die Aufrührer der Obrigkeit übel nachreden und sich darnach also in die Leute flechten, als wollten sie es besser machen, da gehet denn der blinde Pöbel hinan, wie wir in der Bauern Aufruhr auch erfahren haben.
\quellseite{480}

\section{Summarium}

Ist ein sonderlich Exempel, wie David erstlich die Aufruhr willig annimmt und geduldig trägt als eine Strafe, welche er mit dem Mord und Ehebruch verdienet hatte, wie ihm denn Gott durch Nathan solches dräuen ließe. Daher sehen wir, daß Gott keines Menschen verschonet, er sei niedrigen oder hohen Standes. Wer da sündiget, der bleibt gewiß nicht ungestraft. Noch ist menschliche Natur, daß wohl Gottes schreckliche Dräuung höret, daß er das Böse strafen wolle, und dies und dergleichen Exempel, der die Bibel voll ist, lieset, ja auch täglich erfährt, daß Gott das Böse straft, so verstockt und erstarrt, daß sie sich nichts daran kehret. Das ist ja ein unsäglicher Jammer.

Zum andern ist es fein, daß David in solcher Strafe und Unglück dennoch sich aller Gnade zu Gott versiehet, und es ihm heimstellet, er mache es mit ihm, wie er wolle, so soll es ihm gefallen, V. 25. Solche Geduld und Hoffnung sollen wir im Leiden und Anfechtung auch beweisen, so wir anders rechte Christen sein.

\section{Votum}

Gott sei uns gnädig und barmherzig, die wir auch unser Kreuz mit vielen großen Sünden wohl verdienet haben, und verleihe uns allen darunter Geduld und Trost, bis zur seligen von allem Uebel Erlösung, Amen.

\chapter{Das 16. Kapitel}

\section{Vorrede}

Der Prophet Jesaias sagt im 51. Kapitel V. 7.: Fürchtet euch nicht, wenn euch die Leute schmähen; und entsetzet euch nicht, wenn sie euch verjagt machen. Denn die Motten werden sie fressen, wie ein Kleid, und Würmer werden sie fressen, wie wollenes Tuch. Und im 54. Kapitel V. 14. spricht Gott selber: Du wirst ferne sein von Gewalt und Unrecht, daß du dich dafür nicht dürfest fürchten; und von Schrecken, denn es soll nicht zu dir nahen. Siehe, wer will sich wider dich rotten und dich überfallen, so sie sich ohne mich rotten? Und Lucä am 6. V. 38. sagt der HErr Christus: Eben mit dem Maß, da ihr mit messet, wird man euch wieder messen. Auch wird gesagt im Buch der Weisheit im 11. Kapitel V. 17.: Womit jemand sündiget, damit wird er auch geplagt. Dieser Warnung und Trostsprüche Exempel ist das 16. Kapitel.

\section{Ziba Bedeutung. Simei Lästerung. Bluthunde Absalom}
\quellseite{481}

\section{Summarium}

Das ist ein sonderlich Exempel einer hohen großen Geduld, daß David hören und leiden kann, wie ihn Simei so lästerlich schmähet, und will doch nicht, daß man ihm ein Leid darum thue, V. 10. Solche hohe Geduld kommet daher, daß er erkennet, wie es eitel verdiente Strafe sei. Wie ihm nun seine Sünde von Herzen leid ist, leidet er gern und willig, was ihm der Sünden halben begegnet. Solche Geduld ist auch ein Stück, das zur rechten christlichen Buße gehöret.

\section{Summarium}

Ist eine Historia, in welcher wir sehen, wie dennoch Gott allezeit Mittel schaffet in der Anfechtung, auf daß sich die Seinen nicht beklagen dürfen, sie seien ganz verlassen. Darum schaffet Gott, daß Husai, der Arachiter, aus Anregung Davids gen Jerusalem kommet, und stellet sich, als halte ers mit Absalom, und das darum, daß er Ahitophels (der auf Absaloms Seiten war) Rathschlag zunichte machen und David warnen könnte.

Zum andern siehet man hie ein sonderlich Exempel, daß über das, daß David seines Ehebruchs und Todschlags halben des Landes verjagt wird, Gott auch denselben Ehebruch damit strafet, daß gleich wie er einem andern sein Weib heimlich genommen und sie beschlafen hatte: also auch sein Nächster (das ist Absalom) bei seinen Kebsweibern schläft öffentlich an der lichten Sonne, 2 Sam. 12, 11. Nach dem Spruch Weisheit. 11. V. 17.: Damit jemand sündiget, muß er auch gestraft werden. Auf daß jedermann sehe, wie Gott allerlei Unzucht und sonderlich dem Ehebruch feind ist, und sich desto fleißiger dafür hüte.

\section{Votum}

Der treue, barmherzige GOtt verleihe uns auch, daß wir in Geduld alle Lästerung überwinden, offenbare auch alle schädliche Rathschläge, und helfe, daß wir uns für Unzucht und andern Sünden hüten, damit wir nicht verdiente Strafe leiden dürfen, Amen.

\chapter{Das 17. Kapitel}

\section{Vorrede}

List und Gewalt ist der Feinde Gottes Rüstung, damit sie wider seine Kirche streiten. Dagegen Beten und geduldig sein, sind der Christen Waffen, damit sie wider Gottes und ihre Feinde kämpfen. Und Gott sei Lob in Ewigkeit, dieser Theil gewinnet, jener aber verlieret, obschon die Gottlosen eine Weile triumphiren, und die Kinder Gottes eine Zeitlang leiden müssen. Derwegen von beiden also sagt David im 92. Psalm V. 8.: Die Gottlosen grünen wie das Gras, und die Uebelthäter blühen alle, bis sie vertilget werden, immer und ewiglich. Und V. 13.: Der Gerechte wird grünen wie ein Palmbaum, er wird wachsen wie eine Ceder auf Libanon. Davon höret ferner an ein Exempel aus dem 17. Kapitel.

\section{Ahitophels nichtiger Rath, und verzweifeltes Ende}
\quellseite{482}

\section{Summarium}

Ist eine Historia, in welcher man siehet, mit was Ernst und großem Fleiß der Teufel dem frommen David nachstellet daß (menschlicher Weise davon zu reden) nicht möglich ist, daß er entgehen könnte. Aber Gott schaffet dennoch Mittel, daß des Teufels und böser Leute geschwinde Anschläge müssen zu nichte werden. Wie man hie an Ahitophel siehet, der einen Rath gab, der Absalom und allen Aeltesten in Israel gut däuchte; noch macht Gott durch Husai ihn zu nichten, V. 7. Darum liegt es an dem, daß jedermann auf Gott und sein Wort sehe, ihn mit Ernst anrufe und ihm alle Sache befehle; wie David hie thut, da er in dieser seiner großen Noth also zu Gott schreiet: Herr mache den Rathschlag Ahitophels zur Narrheit, 2 Sam. 15, 13., Ps. 31, 18.; so hilft denn Gott gewißlich.
\quellseite{483}

\section{Summarium}

Ist eine sonderliche Historia, wie es mit dem weisen Mann Ahitophel ein Ende nimmt, welcher ein Exempel ist aller deren, so mit menschlicher Weisheit hoch für andern begabt sind, aber doch derselbigen mißbrauchen zur Tyrannei wider das Wort Gottes und die christliche Kirche. Demselben soll es also gehen, wie Salomon sagt, Sprüche 21, 30.: Es hilft keine Weisheit, kein Verstand, keine Kunst wider den Herrn. Darum, weil sie nicht ablassen wollen, müssen sie darüber zu Boden gehen. Wie denn dem Ahitophel sein weiser Rathschlag Ursache gibt, daß er hingehet, und für Leid sich selbst erkennt, daß derselbige nicht fortgegangen war.

\section{Votum}

Gott, der Herr, mache auch heute zunichte des Teufels und aller Gottlosen böse und schädliche Anschläge, und rette und erhalte seine Christenheit und Gläubige in aller Gefahr und Noth, um Jesu Christi, seines Sohnes willen, Amen.

\chapter{Das 18. Kapitel}

\section{Vorrede}

Ungehorsame Kinder laden auf sich diese schwere Bürde und Sünden: Gottes und der Eltern Segen wird von ihnen, ihr Glück aber über sie sein, Sir. 3, 9. 10. Mit bösem Gewissen müssen sie sich tragen ihr Lebenlang; im Sterbestündlein wird ihnen ihr Licht verlöschen mitten im Finsterniß, Spr. Sal. 20, 20. Aus dem Himmel sollen sie verstoßen sein ewiglich, es sei denn, daß sie Buße thun; in der Hölle müssen sie bleiben, in Qual und Pein, in alle Ewigkeit. Das ist eins. Darnach, kein Schmerzen ist so groß, nach dem Schmerzen, denn ein Herz fühlet, wenn dasselbe Gottes Zorn, wegen begangener Sünde, brennet und drücket, als der Schmerzen ist, den die Eltern empfangen, haben und fühlen, wenn es ihren Kindern übel gehet, oder wenn sie in Noth stecken und übel umkommen, sonderlich wo man auch der Seelen Verlust muß fürchten. Beides lehret auch das 18. Kapitel.

Absalom wird an einer Eiche hangend durchstochen.
\quellseite{484}

a Das ist, Absicht.

b Vernimm, Leide thue.

\section{Summarium}

Dies ist eine Historia, wie es mit dem Aufrührer Absalom und denen, so ihm zur Aufruhr geholfen, gegangen und ein Ende genommen habe, nämlich, daß in die zwanzig tausend Mann umgekommen, V. 7., und Absalom mit dem Kopf, durch ein wunderbarlich Urtheil Gottes, an einer Eiche hangen bleibt und also dran erstochen wird, V. 9. Neben solcher Strafe zeiget der Text noch eins an, wie Absalom, der seinen Vater so übel gehalten hat, nicht werth ist gewesen, daß er einen Sohn hat sollen haben, und ist also ohne Erben ausgerottet, V. 14. Darum sollen erstlich die Kinder lernen, daß sie ihre Eltern ehren und für Augen haben. Zum andern sollen alle Unterthanen lernen sich in Gehorsam halten und für Aufruhr zu hüten, als die nimmermehr ohne sonderlichen großen Schaden Land und Leute pflegt abzugehen.
\quellseite{485}

\section{Summarium}

Ist eine Historia, wie David um seinen Sohn Absalom sehr trauert, der ihn doch erstlich hoch bekümmert, darnach des Landes verjagt und endlich ihm nach Leib und Leben gestanden hatte. Aus welchem allem man leichtlich kann abnehmen, wie er solche Sünde alle von Herzen gerne vergeben und ihn wieder zu Gnaden wollte angenommen haben. Nicht allein deshalb, daß er sein Sohn war, sondern auch darum (wie vorgemeldet), daß David solchen Aufruhr nicht anders annahm, als eine verdiente Strafe seiner Sünden. Solche Erkenntniß der Sünden macht ein freundliches Herz, das gern und willig vergibt, auch seinen ärgsten Feinden.

\section{Votum}

GOtt, der heilige Geist, regiere die liebe Jugend, daß sie ihre Eltern nicht betrüben, sondern die ehren und ihnen gehorsam sein; tröste auch und stärke alle betrübte Eltern und bekümmerte Herzen, daß sie ihr Creuz in Geduld tragen und überwinden, Amen.

\chapter{Das 19. Kapitel}

\section{Predigt}

Welcher Mensch seinem Fleisch allen Muthwillen gestattet, und in Sünden lebet, der gehört Christum nicht an. Wer aber seinen Glauben beweiset mit christlichem Leben, nach seinem Beruf, Amt und Stande, der hat Gemeinschaft an GOtt. Derwegen sollen wir auf GOttes Gnade und Barmherzigkeit nicht muthwillig sündigen, sondern sollen gottselig leben und uns aller guten Werke befleißen. Dazu soll uns auch bewegen, daß Gott ist ein eifriger Gott, der die Sünde hasset und strafet, wie er dagegen das Gute liebet und belohnet. Es soll uns auch dazu treiben, daß solches Gott zu Lob, Ehre und Preis dienet; und weil wir fürnehmlich dazu erschaffen, darum erlöst und mit Gottes Wort alle Tage unterrichtet werden, daß wir in Gerechtigkeit und Heiligkeit leben sollen. So betrachten wir allhie auch billig die Vermahnung St. Pauli an die Philipper, da er im 4. Kapitel V. 8. also saget: Was wahrhaftig ist, was ehrbar ist, was gerecht, was keusch, was lieblich, was wohl lautet, ist etwa eine Tugend, ist etwa ein Lob, dem denket nach. So wird der HErr des Friedens mit euch sein. Daher gehört auch das 19. Kapitel.

\section{David wird wieder ins Königreich eingesetzt}
\quellseite{486}

\section{Summarium}

Sie ist eine feine Historia, wie David, da ihm Gott wieder ins Königreich hilft, alle Gnade und Barmherzigkeit jedermann erzeigt, wer Gnade begehrt und nicht tyrannisch mit den Leuten verfährt, die es doch wohl verdienet hatten, als mit Simei, V. 16., Amasa und andern. Denn dies ist die rechte Krone aller fürstlichen Tugenden, daß man Gnade beweise und nicht tyrannisch fahre, wo die, so gesündiget, ihre Sünde erkennen und Gnade begehren.
\quellseite{487}

\section{Summarium}

Dies ist ein Exempel, das sonderlich gen Hof gehöret, wie Ziba seinen Herrn Mephiboseth fälschlich bei David angegeben und in Ungnade gebracht hat, auf daß er zu Hofe desto angenehmer wäre und mehr Gunst hätte. Solches beklaget sich Mephiboseth, V. 26. Aber David beharrt auf dem, daß er Ziba zugesagt hatte, V. 29. Diese Plage ist gemein zu Hofe, daß auch fromme Fürsten sich nicht allweg für Heuchlern hüten können, und bisweilen den Leuten dadurch Unrecht geschiehet, allein darum, daß die Herren zu bald glauben, wie man auch an David siehet. Aber in dem gehet Gottes Urtheil und Strafe wider Saul, daß alle seine Nachkommen solcher Sünden entgelten und weder Glück und Segen müssen haben. Darum die Eltern auch ihrer Kinder halben sich in Gottesfurcht halten und für Sünden hüten sollen, daß die Kinder derselbigen nicht entgelten müssen.

\section{Votum}

Gott, der HErr, regiere Obrigkeit und Unterthanen; die Obrigkeit, daß sie recht und christlich auch regiere; die Unterthanen, daß sie fromm und gehorsam sein, und also auch ihrer Kinder verschonen, daß dieselben ihre Sünde, mit ihrem Schaden und Herzleid, nicht büßen dürfen, Amen.
\quellseite{488}

\chapter{Das 20. Kapitel}

\section{Vorrede}

Wie keine Sünde allein ist, es folget immer eine der andern, wo man nicht in Zeiten Buße thut: also, wenn Gott anfänget, die Leute im Zorn heimzusuchen, so folget auch eine Strafe auf die andere; und wer dem vormals in dem Sündenleben fröhlich gewesen, ach, bei dem wird das Lachen sehr theuer, unter den vielen und großen Strafen. Derwegen wollen wir durch Hülfe des heiligen Geistes aufhören, Böses zu thun und von Sünden abstehen, daß Gott nicht strafen dürfe, und wollen Gutes thun und in allen Tugenden uns fleißig üben, daß Gott uns segne und an Leib und Seele behüte aus Gnaden, um Jesu Christi, seines geliebten Sohnes willen, Amen. Dies soll auch erinnern das 20. Kapitel.

\section{Seba, der Aufrührer, kriegt seinen Lohn}

a Berühmter) einer von den großen Hansen, vom hohen Adel, der großen Anhang im Volk und Ansehen oder Namen hatte, wie Catilina zu Rom.
\quellseite{489}

\section{Summarium}

Sie ist eine Historia, wie ein anderer Aufruhr wider David erreget wird, daß allein der Stamm Juda an ihm bleibt hangen, v. 2. Auf daß man lerne Gott fürchten und sich für Sünden hüten. Denn wenn Gott einmal anfähet zu strafen, hält er lange an. Es ist auch hier ein Exempel, wie es dem Aufrührer Seba gehet, nämlich, daß er den Kopf drüber lassen muß, v. 22. Denn Gott hat’s so beschlossen: Wer das Schwert nimmt, dem es nicht gehöret, soll durchs Schwert wieder umkommen.

\section{Votum}

Der allmächtige Gott steure heute auch allen Aufrührern und ungehorsamen Unterthanen, und regiere uns, daß wir ja nicht in Sünden leben, sterben und verderben; sondern helfe, daß wir in Gottesfurcht beständig bis zum seligen Simeonis-Stündlein verharren, Amen.

\chapter{Das 21. Kapitel}

\section{Vorrede}

Viele Eltern leider! erben auf ihre Kinder nicht Gottes Segen, sondern seinen Fluch und verlassen seinen Zorn und Strafe zum Gedächtniß. Denn wie die Kinder wohl über langer der Eltern Gottes-Furcht und Gehorsam gegen Gott genießen können: also müssen auch der Gottlosen Kinder ihrer Eltern Mißthat ausfeucheln, wenn es gleich noch so lange anstünde. Dieses sollen alle Väter und Mütter wohl wahrnehmen, und in allem schuldigen Gehorsam gegen Gott wandeln; und also werden sie die besten Schätze ihren Nachkommen sammeln. Dagegen sollen sie meiden alle Sünde und Mißthat. Denn sonst führen sie also Gottes Fluch und verderblichen Schaden auf dieselben. Ferner sagt David im 34. Psalm v. 20. Der Gerechte, das ist, ein gläubiger und frommer Mensch muß viel leiden; aber der HErr hilft ihm aus dem allen. Vernehmt, durch Trost, Errettung und endliche Erlösung. Beides der Lehre und Trostspruchs Exempel ist das 21. Kapitel.

\section[David wehret der Theurung; erhält Siege wider die]{David wehret der Theurung; erhält Siege wider die Philister}
\quellseite{490}

\section{Summarium}

Ist eine Historia, wie Gott nach dem Aufruhr das Land mit Hunger gestraft habe drei Jahre lang, V. 1.

* Sonderlich aber ist zu merken, wie sieben Söhne vom Geschlechte Saul gehenkt sind worden, darum, daß ihr Vater die Gibeoniter ohne alle Schuld, wider den Eid, so ihnen von dem Volk Israel geschehen, getödtet hat, V. 8. Auf daß die Eltern auch desto gottesfürchtiger und frömmer lernen sein, daß ihre Kinder ihrer Sünde nicht entgelten und desto mehr Unglück müssen leiden, wie man sehet.

\section{Summarium}

Ist eine historia, wie David weiter mit Kriegen wird angefochten und sonderlich in einer Schlacht Noth leidet, daß, wo Abisal ihm nicht geholfen hätte, er umgekommen wäre, V. 15. Auf daß man sehe, wie der fromme König so in mancherlei Gefahr und Unglück gesteckt ist, aber dennoch Gott in allen Wegen ihm gnädiglich heraus geholfen hat. Denn, ein König oder Regent sein, heißt nicht gute Tage haben, sondern seines Berufs fleißig warten, und darüber sich keiner Gefahr weigern, die sich im Handeln täglich mit Haufen finden, wie man am frommen David siehet.

\section{Votum}

Der treue, liebe GOtt regiere die Eltern, daß sie also leben, handeln und wandeln, damit
\quellseite{491}
 sie ja ihren Kindern nicht schädlich sein; und erlöse alle, die das Kreuz in wahrem Glauben und Geduld tragen, nach seinem Willen und wie es ihnen gut und selig sein soll, Amen.

\chapter{Das 22. Kapitel}

\section{Vorrede}

In der Noth sollen wir Gott um Hülfe anrufen und wenn uns geholfen ist, sollen wir auch Gott danken für erzeigte Wohlthat. Denn also befiehlet es Gott selber im 50. Psalm V. 15.: Rufe mich an, spricht er, in der Noth, so will ich dich erretten, so sollt du mich preisen. Und 1 Tim. 2 V. 1. vermahnet St. Paulus auch dazu und sagt: So ermahne ich nun, daß man für allen Dingen zuerst thue Bitte, Gebet, Fürbitte und Danksagung für alle Menschen. Für die Könige und für alle Obrigkeit, auf daß wir ein ruhiges und stilles Leben führen mögen in aller Gottseligkeit und Ehrbarkeit. Denn solches ist gut, dazu auch angenehm vor Gott, unserm Heilande. Wie solches der heilige David verstanden und gelehret, also hat er auch selber uns zum Exempel gethan. Wie bezeuget das 22. Kapitel.

\section{Lobgesang Davids für die Errettung von seinen Feinden}
\quellseite{492}

a (Sich zu lieben) das sind diejenigen, die sich mit vielen Gottesdiensten wollen um GOtt wohl verdienen, meinen herzlich und thun mit Ernst, aber ohne Gottes Wort, aus eigen erwähltem Fürnehmen, wie unsere Mönche und alle Abgöttische thun.

\section{Summarium}

Das ist ein Dank-Psalm, darin David Gott danket, daß er ihn von allen seinen Feinden erlöset und in allerlei Noth und Gefahr gnädiglich erhalten hat, wider Saul, wider Absalom, wider die Aufrührer und wider die Heiden, welche allzumal, wie ein großes Wetter, David gern zu Boden geschlagen und ersäuft hätten. Aber weil er sich an Gottes Hülfe hält, müssen sie ihn bleiben lassen, wie er spricht: Wenn mir angst ist, so rufe ich den HErrn an, und schreie zu meinem Gott, so erhöret er meine Stimme von seinem heiligen Tempel.

Solchem Exempel sollen wir auch folgen und wissen, wer sich in der Noth an etwas anders, denn an Gottes Hülfe und das Gebet hält, muß endlich zu Boden gehen. Wohl aber denen, die auf den HErrn hoffen und seiner Güte warten. Denn ob sie wohl Unglück nicht können überhoben sein, so ist doch Gott bei ihnen, der hilft in der Noth.

\section{Votum}

Wir danken GOtt auch herzlich, der uns bisher in mancher großer Noth beigestanden und erhalten hat; der wolle ferner unser treuer GOtt und lieber Vater sein und bleiben, daß wir aus aller Noth errettet und im Tode zum ewigen freudenreichen Leben erhalten werden, Amen.

\chapter{Das 23. Kapitel}

\section{Vorrede}

Das 23. Kapitel des andern Buchs Samuels ist fürnehmlich des Königs Davids Testament und letzter Wille. Denn obwohl David hernach eine Zeitlang gelebet und viel geredet, gethan und gelitten hat, dennoch hat er dies sein Testament sein und bleiben lassen. Er lehret aber darin von den fürnehmsten Artikeln christlicher Lehre, als vom wahren Erkenntniß Gottes, von beiden Naturen und großen Wohlthaten des verheißenen Messias unsers HErrn Jesu Christi, und von wahrer Anrufung Gottes. Welches so heilsame und nöthige Artikel sind, daß der HErr Christus Joh. 17, V. 3. also davon sagt: Das ist das ewige Leben, daß sie dich, der du allein wahrer Gott bist, und den du gesandt hast, Jesum Christum, erkennen. Derwegen dies Testament Davids jedermann wohl bekannt sein solle. Solches aber lautet aus dem 23. Kapitel.
\quellseite{493}

\section{Davids letzte Reden und Heldenbuch}

a Ohne Wolken) Weiß richtet des Gesetzes Reich an auf dem Berge Sinai mit Donnern, Wolken, Blitzen schrecklich. Aber dies Reich wird herrlich sein, wie es ist im Lenz, wenn es geregnet hat und die Sonne früh scheint.

b Nichts wächst) kein Königreich ist so hoch für Gott, wird auch nicht so wachsen, sondern vergehen. Allein dies Reich bestehet ewiglich.

c Belial) sind die, so dem Reich Christi feind sind, als Juden, Papst, Ketzer, Türken etc. Die wollen allein nütze und die besten sein, und sind doch die schädlichsten, darum heißen sie Belial, die Unnützen oder Schädlichen. Also sagt Jer. 23, 3. von den falschen Propheten: Sie sind mit ihrem Nützen kein Nutzen diesem Volk, das ist, sie sind die schädlichsten, eben da sie nütze sein wollen.

d Jasabeam) an diesem Ort steht es im Hebräischen also: dies sind die Namen der Helden David, Josebeth, Tachmoni der fürnehmste unter dreien. *Jofelino, daßnitz, und schlug achthundert auf einmal. Da achten wir, der Text sei durch einen Schreiber verderbet, etwa aus einem Buch unkenntlicher Schrift, und von unsern Buchstaben. Und sei also Abino für Oret, und Hachnitz für Ethanito gemacht. Denn die Hebräer wohl wissen, wie man in ihrer Handschrift kann Daleth für Resch, Nun für Zain, He für Thau und wiederum lesen. Darum haben wir nach dem Text 1 Chron. 12, 11. corrigiret. Denn der Text an diesem Ort nichts gibt. Dergleichen kann auch geschehen sein in dem Wörtlein drei, dem acht hundert, so in der Chronica dreißig Item dreihundert stehen. Doch kann das eine andere Meinung haben, wie unten 1 Chron. 12, 11.
\quellseite{494}

\section{Summarium}

Ist ersichtlich eine Weissagung Davids von seinem Königreich, daß es aufgehen, wachsen und sein werde, wie das Licht ist am Morgen, wenn die Sonne aufgehet frühe ohne Wolken, dergleichen Reich in der ganzen Welt nicht sein werde, V. 4. Solches gehet auf das Reich Christi, welcher vom Stamme David ist herkommen. Desselbigen Feinde sollen wie die Dornen ausgerottet werden. Zum andern werden hie erzählet die tapfern Helden, welche David bei ihm zu Hof gehabt hat, V. 8. Da sehet man, wie Gott erstlich viel seiner Leute gibt und einem Fürsten an die Seite setzet, wenn er einem Reich helfen und es groß will machen. Wiederum, wo es an seinen redlichen Leuten fehlet, da werden die Regimente nicht lange bestehen können.

\section{Votum}

Gott, der HErr, segne und erhalte ferner das Reich seines lieben Sohnes, unsers HErrn Jesu Christi, und gebe auch im weltlichen Regiment seiner Leute, die zu seinen Ehren und gemeinem Nutzen zum Besten herrschen und dienen, Amen.

\chapter{Das 24. Kapitel}

\section{Vorrede}

Dieweil auch nach der Wiedergeburt in allen Menschen viel übrige Sünden im Fleisch bleiben, durch welche ein Christ, wo er nicht wohl Achtung auf sich selbst gibt, in verderbliche Sicherheit kann verführet werden; als soll derhalben ein jeder desto fleißiger auf sich selbst Achtung geben, damit er durch solche anklebende sündliche Gebrechen und Zuneigungen mit diesem Laster nicht eingenommen werde, welches denn ohne Verleugnung des Gewissens, ohne Verlierung des heil. Geistes und des Glaubens, auch ohne zeitliche und ewige Strafe nicht abgehen würde. Derwegen zur Vermeidung solches Schadens soll ein jedes durch tägliche Reue und Leid, und durch ernstlichen heftigen Widerstand, neben inbrünstigem Rufen und Schreien zu Gott, dahin trachten, damit es solchen bösen Lüsten, durch Beistand des heil. Geistes möge widerstehen, auf daß solche Funken nicht zu einem schädlichen Feuer werden, das ist, es soll ein jedes, wie St. Paulus sagt Gal. 5, 24.: Sein Fleisch sammt den Lüsten und Begierden kreuzigen. Das erinnert auch das 24. und letzte Kapitel des andern Buchs Samuels.

\section[Davids Ehrgeiz in Zählung des Volks wird mit]{Davids Ehrgeiz in Zählung des Volks wird mit der Pestilenz bestrafet}
\quellseite{495}

\section{Summarium}

Das ist wiederum ein Exempel, wie auch heilige Leute, wenn sie außer der Noth und Anfechtung sind, leichtlich in Sicherheit und Sünde fallen. Derhalben ist nütze und gut, daß Gott durchs Kreuz uns in seiner Furcht erhalte. Droben haben wir oft gehört, wie David, wenn er in Anfechtung ist, vom Saul verfolget wird oder sonst mit ihm übel zugehet, betet und geduldig ist. Wie aber, nachdem er Friede hat und sich für niemand besorgen darf, wird er sicher, will wissen, wie viel er Volks in das Feld wider die Feinde vermöge, V. 2. So doch Gott verboten hat, 2 Mos. 39, 12., man sollte das Volk nicht zählen, wollte mans aber zählen, so sollte ein jeglicher Haupt sich mit einem Opfer lösen; sonst würde ihnen eine Plage widerfahren. Solches bedenkt Joab und die andern Hauptleute, und widerrathen es, V. 3. Aber David will sich nicht bereden lassen, er gedenkt, er sei Herr, andere sollen ihm folgen, er dürfe niemand folgen. Wie denn große Herren diese Anfechtung haben und wollen ihres Sinnes sein. Es geräth aber nicht allezeit wohl, wie man an David auch siehet.

\quellseite{496}

a Dieser Arauna wird der Jebusiter König gewesen sein zu Jerusalem, und hernach zu Gott bekehret, fromm und selig worden, sich des königreichs verziehen um Gottes willen.

\section{Summarium}

Hie ist eine Historia, wie David, ob er gleich seine Sünde erkennet, und ließ es ihm leid sein, dennoch eine zeitliche Strafe darum tragen muß. Und daß Gott die Sünde pflegt mit Hunger, Krieg und Pestilenz zu strafen. Darum soll man in Gottes Furcht wandeln, und sich für Sünden bewahren, auch Gott ernstlich bitten, daß er für Sünden und Bösheit wolle, so wir anders dergleichen Strafe entfliehen wollen. Unter welchem doch die allerleichteste ist, wenn Gott mit Pestilenz strafet. Denn der Hunger ist zu jämmerlich, so ist der Krieg zu greulich, als der alles mit sich reißet und verwüstet.

\section{Votum}

Der treue Gott erhalte uns in seiner Furcht, daß wir nicht in Sicherheit und andere Sünden gerathen, und strafe ja nicht unsere bisher begangene Sünde nach unserm Verdienst, und behüte uns sonderlich für Krieg und Blutvergießen, aus Barmherzigkeit, Amen.

\part{Das erste Buch von den Königen}

\chapter{Das 1. Kapitel}

\section{Vorrede}

Es ist gewiß und wahr, daß Gamaliel, der fromme Pharisäer sagt in den Geschichten der Apostel im 5. Kapitel. V. 39.: Was aus Gott ist, das kann niemand dämpfen. Denn sind die Menschen klug und stark; Gott ist weit verständiger und mächtiger, und was er will, das hat Fortgang wider aller Welt, ja auch der Höllen Pforten, Hinderniß. Derwegen sehe ein jeder in seinem Beruf auf Gott, und demselben befehle er desselben Anfang und Ausgang, und lasse sich anderer Leute Widerwärtigkeit und des Teufels Hinderniß nicht irren, lustig machen und aufhalten. Gott ist groß von Rath und mächtig von That. Jer. 32, 19. Und wem Er wohl will, dem muß alles wohl gelingen, in seinen guten Sachen, rechtmäßigem Handel und christlichem Wandel. Davon höret auch an das 1. Kapitel des ersten Buchs der Könige.

\section{Salomo zum Könige gesalbet und angeweyhet}
\quellseite{497}

\section{Summarium}

Es ist eine Historia, wie Adonia, Davids Sohn, etliche fürnehmste Amtleute an sich hänget, und ohne Davids Wissen und Willen das Königreich begehrt an sich zu bringen, V. 5. Aber es gehet, wie der Engel sagt Dan. 4, 14.: Daß der Höchste Gewalt hat über der Menschen Königreiche, und gibt sie, wem er will. Adonia wollte König sein, und ward es nicht; Salomon wollte es nicht sein und ward es. Also muß Gottes Rath für sich gehen; aber Menschen Anschläge müssen zurück gehen, sonderlich in solchen Händeln, die Land und Leute betreffen, und die Leute sich bedünken lassen, sie wollen mit eigener Klugheit viel ausrichten.
\quellseite{498}

\section{Summarium}

Sie ist eine Historia, wie Salomon ordentlicher Weise, nicht aus eigenem Fürnehmen, von David, seinem Vater, zum Könige erwählet und von dem Propheten Nathan und dem Priester Zadok gesalbet wird. Und ist hie sonderlich zu merken, wie das Volk für den neuen König betet, daß Gott bei ihm und seinem Regiment sey, V. 35. 40. 47. Denn wo weltliche Obrigkeit ihren Leuten zum Besten fürgehen soll, da ist sonderlich vonnöthen, daß Gott mit
\quellseite{499}
Segen dabei sei. Es will aber Gott um seinen Segen gebeten sein.

\section{Votum}

Wir danken GOtt für alle unsere liebe Obrigkeit hohen und niedern Standes; der wolle die ferner erhalten und durch den heiligen Geist regieren, daß alle ihre Rathschläge gelangen zu seinen Ehren und den Unterthanen zu Heil und Wohlfahrt, Amen.

\chapter{Das 2. Kapitel}

\section{Vorrede}

An wohlgezogenen Kindern erleben die Eltern Ehre und auch Freude. Und kann den Eltern keine größere Freude widerfahren, denn da sie sehen, daß ihre Kinder sind lebendige Werkzeuge Gottes, Gefäße der Ehren, der Gnaden und Barmherzigkeit Gottes, die GOtt recht erkennen, ihm zu Ehren dienen, ihn anrufen, rühmen und preisen. Darnach, aus wohlgezogenen Kindern werden auch Leute, die ihrem Nächsten wohl und fruchtbarlich dienen, oder derer Dienst dem gemeinen bürgerlichen Leben heilsam und nützlich ist. Derwegen die Eltern auf nichts mehr Mühe und Fleiß wenden sollen, denn daß sie ihre Kinder also unterweisen und ziehen, auf daß sie in diesem Leben an ihren Kindern Freude und Ehre haben, und daß nach ihrem Tode dieselben auch Gott in Heiligkeit und Gerechtigkeit dienen und andern Leuten nütze und dienstlich sein mögen. Dies lehret auch das 2. Kapitel.

\section[Davids letzter Wille und Abschied. Salomos Folge und]{Davids letzter Wille und Abschied. Salomos Folge und ernstliche Regierung}
\quellseite{500}

\section{Summarium}

Hie ist ein sehr schön Exempel, wie David erstlich seinen Sohn Salomo vor seinem Absterben mit Ernst vermahnet, daß er in der Furcht Gottes wandeln und sich in allem nach Gottes Wort halten wolle, auf daß er und seine Nachkommen Glück und Heil mögen haben, V. 3. Derhalben soll man solche Historia und Wort in aller Könige und Fürsten, ja in eines jeden Haus-Vaters Kammern malen und schreiben, daß sie mit solchem Unterricht und treulichen Vermahnen ihren Kindern fleißig vorgehen.

Zum andern ermahnet er ihn auch seines Amts, daß, weil ihm Gott das Schwert in die Hände hatte gegeben, er die Untugend wolle strafen und die bösen Buben bei dem Halse nehmen; V. 5. Denn wo weltliche Obrigkeit die Sünde nicht mit Ernst, wie sichs gebühret, strafen will, da muß Gott strafen und gehet die Strafe gemeiniglich darnach über den ganzen Haufen, das ist, über die Obrigkeit und Unterthanen zugleich.
\quellseite{501}

\section{Summarium}

Hie sind schreckliche Exempel, in welchen wir sehen, wie es mit Adonia, Joab und Simei rc. habe ein Ende genommen: nämlich, daß sie ihrem Verdienste nach endlich sind getödtet worden. Auf daß nicht allein jedermann sich für Aufruhr und Blutvergießen hüte, sondern auch die Obrigkeit lerne, wo sie Glück und Heil will haben, daß sie solche Sünde nicht ungestraft lasse, insonderheit, wo wider Gottes Ordnung Blut ist vergossen, wie Joab gethan hatte.

\section{Votum}

O Du, welcher der rechte Vater ist über alles, was da Kinder heißt im Himmel und auf Erden, helfe gnädiglich, daß auch heute die Eltern ihre Kinder fleißig unterweisen und ziehen; und regiere alle Obrigkeit, daß sie alle Laster strafen, und daß Kinder und Unterthanen in seiner heiligen Furcht und christlichem Gehorsam leben, Amen.

\chapter{Das 3. Kapitel}

\section{Vorrede}

Gleich wie die Obrigkeit selbst Gott anrufen soll, daß er zu ihrer Regierung Glück und Segen geben wolle: also sollen die Unterthanen für ihre Obrigkeit auch fleißig beten. Das lehret der heilige Prophet Jeremias die gefangenen Juden, und schreibt ihnen einen Brief ins babylonische Gefängniß, darinnen er sie, die doch von fremder heidnischer Obrigkeit gefangen gehalten waren, also anredet: Jer. 29 V. 7. So spricht der HErr Zebaoth: Suchet der Stadt Bestes, dahin ich euch habe lassen wegführen, und betet für sie zum HErrn. Denn wenn es ihr wohlgehet, so gehet es euch auch wohl. Und wer das thut, der dienet erstlich damit seiner ordentlichen Obrigkeit besser, denn mit allen Zinsen, Geschoß und Tributen; darnach bessert er das Regiment. Denn solche Gebete sind die allerbesten Anlagen und Landsteurer, welche die Unterthanen zu erlegen schuldig und die dem ganzen Regiment besserlich sind. Darauf höret ferner an das 3. Kapitel.

\section{Salomos Ehe, Gebet und Urtheil}
% Scan 502 omitted: Bible narrative only; no non-biblical apparatus or original structural heading on this scan.
\quellseite{503}

\section{Summarium}

Sie ist eine schöne Historia, welche sonderlich weltliche Obrigkeit wohl merken und bedenken soll, wie Salomo nichts anderes von Gott bittet, denn Weisheit und Verstand, auf daß er dem Regiment wohl fürstehen und jedermann recht thun möge; solches Gebet gefällt Gott wohl, V. 9. Darum wo ein Regent solches mit Ernst bittet, will es Gott reichlich geben. Wie denn Salomo geschah, daß seines gleichen für ihn nicht gewesen, noch nach ihm aufkommen ist.

Zum andern ist hie zu merken, daß Gott Salomo verheißen, wo er sich seinem Vater David nachhalte, wolle er ihm langes Leben geben, V. 14. Daraus folget, daß gottlose Regenten um ihrer Sünden willen oft vor der Zeit sterben und verderben müssen. Darum schaue jedermann, daß er gottesfürchtig sei und recht handele.

\section{Votum}

Der allmächtige Gott, in welches Hand stehen alle Regimente auf Erden, wolle unserer lieben Obrigkeit langes Leben geben und erhalten, und sie regieren, auf daß Gottes Wort ausgebreitet, die Gerechtigkeit befördert und die Bosheit verhindert und gestraft werde, Amen.

\chapter{Das 4. Kapitel}

\section{Vorrede}

Was Gott zu Josua sagt, Jos. 1, 8., das wird allen Regenten gesagt: also aber spricht Gott der Herr: Laß das Buch des Gesetzes nicht von deinem Munde kommen, sondern betrachte es Tag und Nacht, auf daß du haltest und thust allerdinge, nach dem, das darinnen geschrieben steht. Alldieweil wird dir gelingen in allem, das du thust und wirst weislich handeln können. O eine gute nöthige Lehre und tröstliche heilsame Verheißung. So helfe der allmächtige Gott, daß alle Regenten in ihrem Amt nicht ihrem eigenen verderbten Willen, auch nicht bösen Rathgebern, sondern Gottes Wort und Willen folgen, und beende mit reicher Gnade ihr Regiment, daß darinnen ihm eine ewige Kirche gesammelt und gemeiner Friede, mit Zucht und Gerechtigkeit, erhalten werde, zu seinem Lob und zu Nutz den Unterthanen, Amen. Darauf höret auch an das 4. Kapitel.

\section{Salomos Amtleute, Herrlichkeit und Weisheit}
\quellseite{504}

\section{Summarium}

Ist eine Historia, wie Salomo seinen Hof bestellet, und Gott ihm großen Reichthum gegeben, V. 7., und einen herrlichen Namen bei allen Heiden gemacht habe, V. 34. Solches alles dienet dazu, daß man lerne, daß Gott Segen, Glück und Heil den Regenten geben und widerfahren lassen wolle, die sich seines Worts halten und sich befleißigen, daß sie ihrem Amt weißlich und getreulich fürstehen.

\section{Votum}

GOtt sei Dank für alle christliche Obrigkeit und für alle Wohlthaten, so er den Unterthanen durch sie erzeiget hat; der verleihe, daß unsere Regenten alle ferner in ihrem Amt fleißig sein nach GOttes Wort, und glücklich nach GOttes Verheißungen, Amen.

\chapter{Das 5. Kapitel}

\section{Vorrede}

Den gemeinen Land- und Stadtfrieden lobet und liebet jedermann; aber gar wenig sind ihr, die ihn recht und christlich gebrauchen. Derwegen laßt uns Gott herzlich danken, daß er uns so lange Zeit Friede verleihet gnädiglich, und schaue ein jeder ja zu, daß er den Frieden in seinem Beruf gebrauche, Gott zu Ehren, seinem Nächsten zu Nutz, und ihm selber zu gutem Gewissen, auf daß man mit Mißbrauch des Friedens zu Krieg und andern Strafen nicht Ursache gebe. Ach, das merket wohl. Denn Krieg ist der größten Strafen eine, als die da die Religion, weltlich und häuslich Regiment, zerstöret und hinweg nimmt, daß Pestilenz Kinderspiel gegen Krieg zu rechnen ist. Darum auch David lieber Pestilenz wählete, denn Krieg, da er unter dreien Strafen eine wählen sollte; wollte also lieber in Gottes, denn Menschen Hände fallen, 2 Sam. 24, 14. Darauf höret ferner an das 5. Capitel.

\section{Salomos Bund mit Hiram, bei Vorbereitung zum Tempelbau}
\quellseite{505}

\section{Summarium}

Hie ist ein fein Exempel, welches sonderlich die Obrigkeit merken und ihm nachfolgen soll: daß Salomo den Frieden, so Gott in seinem Regiment ihm widerfahren läßt, dahin brauchet und wendet, daß er den rechten Gottesdienst ausrichten und fördern möge, V. 3. Denn dahin soll man nicht allein Frieden, sondern auch allen andern Segen, wie ihn Gott gibt, brauchen, daß Gottes Ehre und Wort gefördert, und die Leute dadurch gebessert werden. Da sonst die Welt den Frieden mißbrauchet, entweder zu eignem Wollust und Sicherheit, oder zu Muthwillen und unnöthigen Rumoren.

\section{Votum}

Gott, unserm HErrn, sei Lob und Dank, daß unsere liebe Obrigkeit den Frieden auch erhält; der verleihe ferner, daß wir unter ihr sein Wort und Frieden haben, und allzeit recht' und christlich gebrauchen, auf daß wir in Friede leben und sterben, Amen.

\chapter{Das 6. Kapitel}

\section{Der Tempel Salomos wird in sieben Jahren aufgebaut}

a Dreißig Ellen hoch) im andern Buch der Chronica Kap. 3, 4. spricht der Text: das Haus sei hundert und zwanzig Ellen hoch gewesen, welches ist von des ganzen Hauses Höhe geredet. Sie aber redet er vom untersten Gemach alleine, welches dreißig Ellen hoch war.

b Die Fenster, so inwendig weit und auswendig enge sind, da kann man nicht wohl noch viel hinein sehen; aber sehr hoch
% Scan 506 omitted: Bible narrative only; no non-biblical apparatus or original structural heading on this scan.
\quellseite{507}

\chapter{Das 7. Kapitel}

\section[Verfertigung der königlichen Häuser und der Gefäße des]{Verfertigung der königlichen Häuser und der Gefäße des Tempels}
\quellseite{508}

\chapter{Das 8. Kapitel}

\section{Vorrede}

Zu einem recht christlichen Leben gehören auch diese Stücke: erstlich, daß ihm niemand aus eigener Andacht Gottesdienst soll erdencken, sondern allein thun, was Gott in seinem Worte befohlen hat. Und das lehret Gott selber Ezech. 20, 18., da er spricht: Ich bin der HErr, euer Gott, nach meinen Geboten sollt ihr leben, und meine Rechte sollt ihr halten. Zum andern soll
\quellseite{509}
ein jeder Gott danken für empfangene Wohlthaten. Das lehret auch Gott, der HErr, und spricht Pf. 50, 15.: Rufe mich an in der Noth, so will ich dich erretten, so sollt du mich preisen. Und zum dritten soll man beten für seine eigene und für der Christenheit gemeine Noth. Das lehret St. Paulus 1 Tim. 2, 1. und sagt: Ich ermahne euch, daß ihr für allen Dingen thut Bitte, Gebet, Fürbitte und Danksagung für alle Menschen. Dergleichen, welcher Mensch Buße thut, der diene also Gott und seinem Nächsten; so wird’s eine rechte, Gott, dem HErrn, eine wohlgefällige, und ihm und seinem Nächsten eine nützliche Buße sein. Davon höret auch an dem ersten Theil des 8. Kapitels.

\section{Der Tempel Salomos wird eingeweihet}



\section{Summarium}

In eine Historia, wie man die Lade des HErrn sammt dem andern Kirchengeräthe in den Tempel gebracht, und da Gott geopfert und gedanket hat. Sie muß man erstlich merken, daß Salomo
\quellseite{510}
solches nicht aus eigener Andacht, sondern aus Befehl Gottes gethan habe. Sonst wäre es Gott kein Dienst noch Wohlgefallen gewesen. Zum andern muß man auch darauf sehen, was es für eine Ursache sei, daß Salomo und alles Volk so froh ist, und Gott so loben und danken; nämlich darum, daß der Gottesdienst an einem gewissen Ort verordnet war, da er allweg sollte bleiben, und die Leute gewiß wüßten, wenn sie an dem Ort beteten, opferten und dergleichen thäten, daß es Gott dienet und seinem Befehl gefolget hieß, und da Gott sich finden lassen und sie erhören wollte.

a Brand) wenn das Getreide verschienen, oder von der Sonne verbrannt ist.
\quellseite{511}

b Fremder) dies Haus soll ein Bethaus sein allen Völkern.

\section{Summarium}

Hie ist ein Gebet, daß Gott ehrlich seine Verheißung, so er David gethan, halten und vollstrecken wollte. Auf daß sein Reich so lange bestehen möge, bis daß Christus käme und das ewige Reich durch sein Evangelium anrichtete, V. 29. u. f. Zum andern, daß Gott sich in diesem Hause gnädig wolle lassen finden, und seines Volks, auch der Heiden Gebet in allen Nöthen und Anliegen erhören, V. 25. Und hie ist sonderlich zu merken, daß Salomo allweg hinzusetzt, daß sie Gott erhören wolle, wenn sie ihre Sünde erkennen und den Namen des HErrn bekennen etc. Denn wo Sünde und Unglück uns nicht drücken, daß wirs fühlen, schreien wir zu Gott um Gnade und Hülfe nicht.

\section{Votum}

Der barmherzige GOtt führe, leite und regiere uns, daß wir in seinen Geboten wandeln; und wie er uns bisher aus Gnaden viel Wohlthaten erzeiget hat, also wolle er uns ferner segnen, und mit uns alle gottselige Christen an Leib und Seele, zeitlich und ewig, Amen.

\section{Der andere Theil des 8. Capitels}

\section{Vorrede}

Gleich wie Salomo für alle seine Unterthanen bittet; also soll täglich ein jeder Christ für den andern bitten. Denn also lehret St. Jacobus Kap. 5, 16. und sagt: Betet für einander. Und derwegen leben und wohnen auch die Menschen beisammen, in Städten und Dörfern, nicht daß sie mit einander fressen und saufen, und andere Wollust treiben sollen; sondern daß einer des andern Noth sehen, ihm rathen und helfen und für ihn bitten soll. Und soll solches thun, nicht allein angelangt, sondern auch ungebeten, der anders ein Christ ist; und soll also sein als der erste Psalm sagt V. 2.: Wie ein Baum, der an die Wasserbäche gepflanzet ist, und seine Frucht bringet zu seiner Zeit. Ein guter Baum aber gibt keine Früchte, wenn sie reif sein, auch ungeschüttelt und ungerüttelt; also thut auch ein rechter Christ Gutes, auch ungebeten, zu seiner Zeit, wenn er siehet, daß sein Nächster seiner Hülfe bedarf, und bittet für ihn. Darauf höret nun an den andern Theil des 8. Capitels.
\quellseite{512}

\section{Summarium}

Ist ein feiner Beschluß des Gebets, welches Salomo für sein Reich und Volk gethan hat, da er begehret, Gott wolle solch Gebet erhören, bei ihnen sein, wie er bei ihren Vätern gewesen, sie nicht verlassen, noch die Hand abziehen, sondern ihre Herzen dahin neigen, daß sie in seinen Wegen wandeln und sich nach seinem Worte halten mögen. Auf die Weise sollte eine christliche Obrigkeit alle Tage für sich und ihre Unterthanen beten; so würde Glück und Heil beim Regiment sein.

Das aber hie stehet, daß Salomo habe das Haus also geweihet, V. 63., ist viel ein ander Ding, denn wenn unsere Papisten Kirchen weihen, mit Kreuz machen, segnen und anderm Affenspiel, welches sie ohne Gottes Befehl selbst erdichtet haben. Denn weihen heißt hie nichts anders, denn daß man da habe angefangen zu opfern, und den rechten Gottesdienst zu üben, als, Gottes Gnade und Wohlthat predigen, ihn daselbst loben, preisen, in aller Noth anrufen etc. Von solchem Weihen wissen die Papisten nichts, denn sie üben es ja nicht, und gehen derweil mit ihrem unnützen, stummen und ohne Gottes Befehl erdichtetem Gaukelwerk um.

\section{Votum}

GOtt, der heilige Geist, der ein Geist der Gnaden und Gebets ist, helfe uns allen, daß wir in diesen betrübten und gefährlichen Läuften fleißig für einander beten, und in allem Gehorsam der heiligen zehen Gebote GOttes leben, und behüte uns für aller Abgötterei und falschen Gottesdiensten, Amen.

\chapter{Das 9. Kapitel}

\section{Vorrede}

Wenn wir zu Gott, dem HErrn, unser Gebet thun, so sollen wir mit allem Fleiß einen Unterschied halten zwischen den Gütern und Gaben, die wir bitten. Denn, wenn wir Gott bitten um geistliche ewige Gaben, die zu unserm Heil und Seligkeit gehören, da stehet die Verheißung gewiß und fest, ohne alle Exception und Bedingung, daß uns solche Gott wahrhaftig will gewähren, so gewiß als er lebet und unser Vater ist. Wenn wir aber bitten um zeitliche, leibliche Gaben, da stehet zwar die Verheißung auch gewiß und fest; aber mit der Bedingung, wo dasjenige, so wir bitten, Gott zu Ehren gelanget und uns gut und selig ist. Derwegen auch die lieben Heiligen in der Schrift, wenn sie um vergängliche, leibliche Gaben bitten, allzeit den Willen Gottes zuvor behalten; wie davon lehret
\quellseite{513}
der heilige Apostel Johannes 1. Epistel 5, V. 14.: Das ist die Freudigkeit, die wir haben zu Gott, daß so wir etwas bitten nach seinem Willen, so höret er uns. Und so wir wissen, daß er uns höret, was wir bitten, so wissen wir, daß wir die Bitte haben, die wir von ihm gebeten haben. Dieß soll auch erinnern das 9. Kapitel.

\section[Der HErr erscheint Salomo zum andern mal. Feste]{Der HErr erscheint Salomo zum andern mal. Feste Städte. Schiffahrt}
\quellseite{514}

\section{Summarium}

Ist eine Historia, wie der HErr dem Salomo erschienen, V. 2., und ihm seine Bitte gewähret habe, V. 3. Daher stimmt es, daß hernach Daniel und andere Heiligen, wenn sie im Lande Babel beteten, ihr Angesicht gegen Jerusalem zum Tempel kehreten. Weil Salomo darum gebeten und Gott hernach solches zugesagt hatte, er wollte alsdann hören und helfen, wo sie an diesem Ort, oder so sie nicht im Lande wären, doch gegen diesen Ort beteten, und ihr Angesicht gegen Jerusalem wendeten. So viel aber die Verheißung das leibliche Königreich betrifft und Salomons Nachkommen, siehet man, daß es mit dem Gedinge ist zugesagt: Gott wollte sie dabei erhalten, so fern sie fromm sind und Gottes Wort halten, V. 4. Darum, da sie böse und ungehorsam worden, sind sie, wie der HErr hie dräuet, gestraft und von solcher Herrlichkeit gestoßen worden.

\section{Votum}

Der treue wahrhaftige GOtt erhöre auch unser gläubiges Gebet, daß wir nach seinem Willen thun, und segne und schütze uns und die Unsern in diesen gefährlichen und kümmerlichen Zeiten, und thue solches um JEsu Christi, unsers HErrn und Heilandes willen, Amen.

\chapter{Das 10. Kapitel}

\section{Vorrede}

St. Paulus in der Epistel an die Römer im 15. Kapitel V. 9. sagt, daß die Heiden Gott loben um der Barmherzigkeit willen. Und David im 117. Psalm V. 1. vermahnet: Lobet den HErrn alle Heiden, und preiset ihn alle Völker. Derwegen, weil wir auch geborne Heiden sind, so danken wir heute abermals Gott herzlich, daß er uns zu seinem Volke angenommen und zum Reich Christi hat kommen, und daß er uns seiner Wohlthaten läßt theilhaftig werden. Der helfe, daß wir als sein Volk alhie auf Erden in wahrem Glauben und Gehorsam, und droben im Himmel, in ewiger Seligkeit und Freude, alle leben und ihn, mit allen Auserwählten, für diese und alle andere seine Wohlthaten, loben und preisen. Davon erinnert auch der erste Theil des 10. Kapitels.

\section[Die Königin aus Reicharabien verwundert sich über Salomos]{Die Königin aus Reicharabien verwundert sich über Salomos Weisheit, Reichthum und Herrlichkeit}
\quellseite{515}

a Ebenholz) ist ein Baum in Indienland.

\section{Summarium}

Ist eine Historia, daß Salomo seiner Weisheit halben allenthalben bekannt und hoch berühmt ist gewesen, V. 1. Solches aber hat dazu gedienet, daß durch solch Gerücht etliche auch aus den Heiden gen Jerusalem kamen, und durch solche Ursache den Gottesdienst alle sehen und erkennen sollten. Denn obwohl Gott die Juden allein zum Eigenthum hatte auserwählet; hat er dennoch immerdar auch etliche aus den Heiden zu sich gezogen. Auf daß sein Volk bald merken und gewöhnen sollte, daß Gott auch der Heiden Gott sei, und sie durchs Evangelium wollte zum Reich Christi bringen.

Denn daß die Königin aus dem Reich Arabia Salomonis Weisheit rühmet, V. 6., ist nicht allein zu verstehen von weltlicher Weisheit, sondern daß er sie aus Gottes Wort, von Gott und dem rechten Gottesdienst hat unterrichtet, um welcher Ursachen willen sie am meisten gen Jerusalem kommen ist, ihn zu sehen und hören, wie man hie sehen kann, 1 Kön. 10, 1., da so geschrieben stehet: Da das Gerücht Salomo, von dem Namen des HErrn, kam für die Königin vom Reich Arabia.

\section{Votum}

Helfe und regiere uns Gott alle, daß wir seine Wahrheit in seinem heiligen Wort auch hören, und in seiner Erkenntniß täglich wachsen und zunehmen, und allzeit den rechten Gottesdienst verrichten, daß sein Name geheiliget und sein Wille geschehe, zu unserm Heil und Seligkeit, um Christi Jesu willen, Amen.

\section{Der andere Theil des 10. Kapitels}

\section{Vorrede}

Daß die Regenten Tugend und Gaben haben, die zu glückseliger Regierung von nöthen sind, das kömmt allein von Gott, dem unerschöpflichen Brunnquell aller Weisheit und Gütigkeit. Solches bezeuget Salomon aus eigener Erfahrung, und sagt an Gottesstatt, Spr. Sal. 8, 14. 15.: Mein ist beide Rath und That, ich habe Verstand und Macht, durch mich regieren die Könige, und die Rathsherren setzen das Recht, durch mich herrschen die Fürsten und alle Regenten auf Erden. Derwegen laßt uns Gott herzlich bitten, er wolle allen christlichen Regenten, als seinen und seines Reichs Amtleuten, zur Verwaltung seines Reichs und Amts, Rath und That geben, auf daß sie allzeit also regieren und herrschen, daß es zu seines heiligen Namens Ehre und zu gemeinem Nutz und Friede diene und gelange, Amen. Darauf höret ferner an den andern Theil des 10. Kapitels.
\quellseite{516}

\section{Summarium}

Sie ist eine Historia, wie Salomo ihm habe einen Königs-Stuhl gebauet, und scheinet fast, als habe er damit, gleich wie durch ein Gemälde, wollen anzeigen die Art und Eigenschaft eines vernünftigen Regenten.

Das Elfenbein, wenn mans recht zubereitet, ist es schön, weiß und glatt, bedeutet derhalben die Unschuld. Denn ein Regent soll erstlich fromm und unbefleckt sein, und für allem sich hüten, daß ärgerlich ist.

Das Gold, damit das Elfenbein überzogen ist, bedeutet die Beständigkeit, daß man sich weder Gunst noch Haß zu Unbilligkeit lasse bewegen, und fest und steif bei dem Rechten halte.

Die sechs Stufen am Stuhl bedeuten, daß ein Regent nicht zu jäh[e] sei, sondern ihm der Weile nehmen soll, auf daß niemand übereilet und er selbst durch eigen Gutdünken nicht verführet werde. Denn zu solchem Amt gehöret, daß man wohl besonnen sei.

Dazu dienet auch, daß der Stuhl hinten rund ist. Denn wo man bedächtig handelt, da vergreift man sich desto weniger.

Die zwo Lehnen, darauf man sich steuret, bedeutet das Recht an ihm selbst.

Die Löwen auf den Lehnen und Staffeln bedeuten das rechte Löwenherz, daß ein Regent muß haben, daß er ohne alle Scheu strafe, wo Strafe vonnöthen ist.

Daß nun Salomo so großen Reichthum gehabt, V. 23., ist ein Zeugniß, wo fromme und gottselige Regenten sind, daß Glück und Heil die Fülle werde da sein. Dagegen findet sich eitel Unglück, wo gottlose und ungeschickte Regenten sind.

\section{Votum}

Gott, der Vater des Lichtes, von dem alle gute Gaben von oben herab kommen, segne und mehre auch heute alle Obrigkeit, an Verstand, Tugenden und Reichthum, und gebe zu ihrer Regierung ferner alle glückselige Wohlfahrt, dem Kirchen- und Haus-Regiment zu Schutz und Erhaltung. Amen.

\chapter{Das 11. Kapitel}

\section{Vorrede}

Gleich wie die Gottesfurcht ist die fürnehmste Tugend, welche in sich begreift die erste Tafel der zehen Gebote Gottes; und wie der 111. Psalm V. 10. saget, der Weisheit Anfang, oder der Anfang und das Ende unserer Bekehrung zu Gott: also ist die Gottesfurcht der Ursprung und die rechte Probe aller Tugenden, die in der andern Tafel der zehen Gebote begriffen werden. Und weil diese Tugend sonderlich ins erste Gebot gehöret, die Gott von unsern Herzen erfordert, so haben die Alten gar fein ein Gleichniß von dem ersten Gebot geben wollen, da sie gesagt: Das erste Gebot mit dieser Tugend der Gottesfurcht ist gleich wie die Schiene an dem Kranze der Braut Christi, auf welche die schönen Blumen und Röslein aller andern Tugenden in den nachfolgenden Geboten müssen aufgebunden werden. Denn gleich wie man seinen Kranz aus Blümlein und Röslein machen kann, wo es an einer Schienen mangelt: also machen auch die Tugenden der andern Tafel keinen Christen, wo nicht kindliche Gottesfurcht vorher geht. Derwegen wir alle in Gottesfurcht uns fleißig die Tage unsers ganzen Lebens üben sollen. Dies unterrichtet auch das 11. Kapitel.

\section{Salomos Weiber, Abgötterei, Kinder und Tod}
\quellseite{517}

\section{Summarium}

Hie ist eine historia, wie Salomo erstlich wider Gottes Befehl und Warnung heidnische Weiber nimmt, V. 1., und hernach ihrem Exempel nachfolget und Abgötterei mit ihnen treibet. Auf daß jedermann lerne Gott fürchten und sich in guter Hut halte, sintemal es so leichtlich geschehen kann, daß auch große Heiligen fallen. Denn wie gehört, ist Salomo ein so hochbegnadeter und in Gottes Wort wohl geübter und erfahrener Mann gewesen, und dennoch fällt er so greuslich
\quellseite{518}
 und gefährlich, daß er seinen Weibern folget und abgöttisch mit ihnen wird.

Zum andern stehet man hie, wie Gott Abgötterei strafet mit Krieg und Zerrüttung des Regiments. Darum sonderlich vorzunehmen ist, daß weltliche Obrigkeit den rechten Gottesdienst handhabe und fördere, und aller Abgötterei mit Ernst wehre.

\section{Summarium}

Hie ist eine Historia, in welcher man klar siehet, wie Gott die Abgötterei also strafet, daß er das herrliche und schöne Königreich seines Volks, das unter David und Salomo so hoch gestiegen war, durch Aufruhr zerrütten und zertheilen läßt. Auf daß weltliche Obrigkeit sonderlich Abgötterei und falschen Gottesdienst lerne fliehen, als das schädliche Gift, das alle Regimente wüst macht, und allen Jammer erregt. Wiederum soll jedermann allen Fleiß dahin wenden, daß der rechte Gottesdienst sammt der reinen Lehre gefördert und gehandhabt werde. Denn wo das geschiehet, so gibt Gott Glück und Heil zum Regiment. Wir sollen aber solche Historia nicht allein zur Lehre, sondern auch zum Trost brauchen. Wehe thut es, daß die Bischöfe und etliche wenige Fürsten wider das Evangelium toben, die armen Pfarrherrn, so Gottes Wort wollen pflanzen, vertreiben und würgen, und an den Unterthanen alle Tyrannei üben. Aber lerne hie, sie können ihr eigen Verderbniß
\quellseite{519}
nicht baß fördern, denn auf diese Weise. Herodes wäre des Herrn Christi, und sein Sohn des heiligen Johannis des Täufers halben wohl König blieben, und in allen Würden gesessen; aber da der Vater die unschuldigen Kindlein und der Sohn den heiligen Johannem läßt würgen, eilen sie nach ihrem Verderben und müssen darum zu Boden gehen.

\section{Votum}

Der starke, allmächtige GOtt laß uns ja den Teufel und unser Fleisch und Blut von seinen Geboten nicht abführen, und regiere uns, daß wir uns für Abgötterei und allen andern Sünden hüten, und steure allen, so muthwillig die Sünden vertheidigen und Gottesfurcht in den Gläubigen verfolgen, Amen.

\chapter{Das 12. Kapitel}

\section{Vorrede}

Die Tyrannei und Abgötterei sind die größten und schädlichsten Sünden der Obrigkeit. Durch Abgötterei wird Gott gelästert und die wahre Religion unterdrückt, zu Schaden der Unterthanen Seelen. Durch Tyrannei muß die Gerechtigkeit Noth leiden, und werden die Unterthanen verderbet. Derwegen laßt uns Gott abermal danken, daß wir gnädige und gütige Obrigkeit, und auch, daß wir von Abgötterei erlöset, Gottes reines Wort und rechten Gottesdienst haben, daß wir ohne Hindernis zur Kirche gehen, mit Lehre, Lernung, Gebet und mit Gebrauch der heiligen Sacramente, Gott loben und preisen können; und daß sich ein jedes in seinem Handel, Acker und Handthierung wohl nähren, und davon durch Gottes Segen das meiste und beste Theil behalten kann. O diese Herrlichkeit und Schätze erhalte uns Gott allzeit, bis an den lieben jüngsten Tag, Amen. Darauf höret davon an das 12. Kapitel.

\section{Zertheilung des Reichs Salomo}
\quellseite{520}

\section{Summarium}

Sie ist eine sehr feine Historia, in welcher man siehet, wie die Tyrannei, zu Aufruhr Ursache gibt, wenn man die Leute zu hart dringen und sonderlich im angehenden Regiment will zu strenge sein, V. 10. Aber junge hitzige Leute, wie Rehabeam, und sein junger Rath denken, sie wollens mit Gewalt hinaus führen, wie sie es haben fürgenommen. Aber es findet sich viel anders.

Sie muß man aber nicht vergessen, daß der Text meldet, daß solcher unzeitiger und hitziger Rath, der zu Aufruhr Ursache gibt, eine Strafe sei der Abgötterei und des falschen Gottesdienstes, welchen Salomo hatte angerichtet.

\section{Summarium}

Sie ist eine sonderliche Historia, wie Jerobeam, nachdem er ins Regiment kommt, besorget, daß Königreich werde ihm nicht bleiben, wo er die zehen Stämme Israel lasse hinauf gen Jerusalem ziehen, und daselbst nach Gottes Befehl opfern. Darum erdenket er aus seinem Herzen einen neuen Gottesdienst; aber eben
\quellseite{521}
damit er sein Reich gedenkt zu erhalten, verderbet er’s, V. 28. Gleich wie heutiges Tages viel sind, die also gedenken: Wenn sie bei dem alten papistischen Wesen bleiben, so wollen sie mancher Gefahr und Unrath entlaufen. Aber damit versündigen sie sich je länger je mehr, bis sie endlich gar zu Boden gehen müssen. Denn wer ein beständig Glück will haben, der denke und halte sich zu Gott und seinem Wort, sonst soll und muß er zu Boden gehen, wenn er gleich aller Kaiser Geld und Gewalt hätte.

\section{Votum}

Der treue Gott behüte die ganze Christenheit gnädiglich für tyrannischem Regiment, und für öffentlicher und heimlicher Abgötterei; und helfe, daß Obrigkeit und Unterthanen ihm, nach seinen Geboten, gehorsamen, zu zeitlicher und ewiger Wohlfahrt, Amen.

\chapter{Das 13. Kapitel}

\section{Vorrede}

Wer da sündiget, der thut übel und unrecht; wer aber gesündiget hat, und sich Gottes Wort strafen läßt, und Buße thut, der thut wohl und recht. Derwegen verharre niemand halsstarrig in seinen Sünden, sondern, da ihm Buße von der Kanzel oder im Beichtstuhl geprediget wird, so erkenne er seine Sünde, lasse die ihm leid sein und bekehre sich zu Gott im wahren Glauben, mit beständigem Vorsatz, daß er wolle fromm sein. Denn wer in Sünden vorsätzlich verharrt, und flapp auf flapp schlägt, das ist, eines nach dem andern sündiget, dem gereicht es endlich zu zeitlichem und ewigem Leibes- und Seelen-Schaden. Denn der Tod ist der Sünden Sold, Röm. 6, 23. Das ist, die Sünde und Bosheit besoldet, belohnet und bezahlet Gott mit allerlei Strafe, mit dem zeitlichen und ewigen Tode oder Verdammniß. Dies soll auch erinnern der erste Theil des 13. Kapitels.

\section[Jerobeams Hand verdorret. Ein Prophet wird vom Löwen]{Jerobeams Hand verdorret. Ein Prophet wird vom Löwen getödtet}

\section{Summarium}

Dies ist eine Historia, wie Jerobeam seiner Abgötterei halben von einem Propheten gestraft wird, V. 2. Da siehest du, daß aller Abgöttischen und ketzer Art ist, daß sie wollen ungestraft sein und Recht haben. Derowhalben befiehlt der König auch, man soll den Propheten greifen, V. 4. Solche Gefahr müssen alle die auf sich nehmen, welchen Gottes Wort zu predigen und Abgötterei und Aergerniß zu strafen befohlen ist.

Neben dem siehet man auch, wie an einem verkehrten Menschen gar nichts hilft. Jerobeam höret den Propheten selbst, und konnte wohl gedenken, daß er von Gott zu ihm gesandt war, erstlich dadurch, da er nach ihm greift, seine Hand ihm so bald dürr wird, V. 4. Darnach auch, da der Prophet für ihn betet, daß ihm die Hand wieder zurecht gebracht wird, V. 4. Daran kehret er sich aber gar nichts, sondern bleibet
\quellseite{522}
bei seinem abgöttischen Wesen wie vor, bis er von der Erde vertilget wird. Wie man heutiges Tages dergleichen Exempel sehr viel erfähret an den Feinden des Evangelii, und täglich mehr erfahren werden. Denn Gott hat nie keines geschonet, es sei denn rechte Buße erfolget.

\section{Votum}

Helfe GOtt, der sein heilig Wort zu lehren und lernen befohlen hat, daß die Lehrer seines Worts die Sünden ernstlich und fleißig strafen, und regiere uns, daß wir uns alle aus den Straf-Predigten recht und wohl bessern, und ihn zeitlich und ewig mit einander loben und preisen, Amen.

\section{Der andere Theil des 13. Kapitels}

\section{Vorrede}

Wer von Gottes Wort abweicht, von dem weicht Gott wieder mit seinem Schutz und Segen; und wer Gottes Wort nicht folget, dem folget dagegen Unglück und Untergang. Das wird täglich gelehret, und davon gibt Gott oft Zeugniß mit wunderbaren und großen Strafen. Dennoch sind ihrer viel, die mehr anderer Leute Rath und ihrem Gutdünken folgen, denn Gott und seinem Wort und Geboten. Derwegen wollen wir unsern Rathgeber allein sein lassen das heilige Wort Gottes, Ps. 119, 24. Und wollen dem gehorsamlich folgen, auf daß wir in Gottes Segen und Schutz unter aller Gefahr geleitet und erhalten werden. Und dazu sollen alle fromme Herzen täglich beten aus dem 119. Psalm V. 4. 29.: HErr, du hast geboten fleißig zu halten deine Befehle; wende von mir den falschen Weg, und gönne mir deine Gesetze. Darauf höret auch an den andern und letzten Theil des 13. Kapitels.
\quellseite{523}

\section{Summarium}

Ist eine wunderbare Historia, wie der Prophet, welcher so eine herrliche Prophezeiung wider Jerobeam und den Altar und hernach ein herrliches Wunderwerk gethan hatte, von einem Löwen zerrissen wird, darum, daß er wider den Befehl Gottes zu Bethel gegessen hatte.

Da werden wir erstlich vermahnet, dem mit allem Ernst zu glauben und fest daran zu halten, was Gott einmal geboten und befohlen hat, und uns nicht davon lassen abwenden und reißen, wenn gleich jemand käme und gäbe für, er wäre ein Prophet, mit dem ein Engel durch des HErrn Wort geredet hätte etc. Denn was hat diesen lieben Mann Gottes anders zu diesem Fall gebracht, denn daß er gemeint hat, weil der andere so herrlich rühmet, er wäre auch ein Prophet, so müßte er Gott gehorsam sein, und ist also durch Gottes Namen betrogen worden.

Darnach hat auch Gott dem Könige Jerobeam und dem ganzen Israel durch solch Werk wollen anzeigen, wie ihnen ihre Abgötterei gerathen soll. Und wie eine große Sünde es sei, da sie ihnen Gottes Wort haben predigen lassen, und sind ihm doch nicht nachkommen, daß sie auch also sollen umkommen und zerrissen werden; denn Gott will, daß man seinem Wort folgen oder alles Unglücks gewarten soll.

\section{Votum}

Erleuchte und regiere uns, GOtt, unser lieber Vater, daß wir durch anderer Leute Schaden klug werden, und ja in all unserm Thun und Lassen unserer Beruf und Stände im Gesetze des HErrn wandeln, und allein GOttes Wort folgen, Amen.

\chapter{Das 14. Kapitel}

\section{Vorrede}

Wie der Götzendiener Hauptmann, der Teufel, Adam auch dahin brachte, daß er mehr von der Kraft eines Apfels hielt, denn von Gott und allen seinen Gaben, die er ihm hatte verliehen, und meinte, er wollte in dem Apfel alle göttliche Weisheit finden, welches Abgötterei war: also hanget uns noch heutiges Tages allen an diese Unart, daß wir von Natur nichts anderes sind, denn lauter abgöttische Menschen und Götzendiener, Röm. 6, 17. Denn was wir selber von uns und unserm Thun halten, das wissen wir wohl; wenn wir es nur wollten bedenken und uns lassen leid sein. Denn daher kömmt alle Abgötterei, daß ein jeder seinem Gutdünken mehr, denn dem klaren Wort Gottes folget, und sich auf die Creaturen mehr, denn auf seinen Schöpfer verläßt. Ach, so laßt uns alle Abgötterei fliehen, und zu Herzen nehmen diese ernste Dräuung und Warnung, als das Moses sagt in seinem fünften Buch im 8. Kapitel V. 19.: Wirst du des HErrn, deines Gottes, vergessen und andern Göttern nachfolgen, und ihnen dienen, und sie anbeten, so bezeuge ich heute über euch, daß ihr umkommen werdet. Und im 14. Kapitel des Buchs der Weisheit saget der heil. Geist V. 27.: Den schändlichen Götzen dienen ist alles Bösen Anfang, Ursache und Ende. Deß ist ein Exempel das 14. Kapitel.

\section[Weissagung wider Jerobeam. Abgötterei, Regirung; Sünde und derselben]{Weissagung wider Jerobeam. Abgötterei, Regirung; Sünde und derselben Strafe}
\quellseite{524}
ᵃ Mimesis, in futuro dicetur, his impletis. Was hat Jerobeam nun gemacht.

\section{Summarium}

Dies ist die dritte Warnung, daß Gott den König Jerobeam seiner Abgötterei halben durch den Propheten Ahia strafet. Aber es ist alle Warnung und Strafe an den Gottlosen verloren, man sage ihnen, was man wolle, so thun sie, was ihnen gefällt.

Und ist hie sonderlich die Predigt des Propheten Ahia zu merken, wie heftig er dräuet, daß Gott die Abgötterei Jerobeams nicht allein an seinem ganzen Hause, sondern auch an dem ganzen Volk Israel strafen wolle, v. 7. 10. Auf daß jedermann sich für Abgötterei hüte, und mit allem Fleiß sich zum rechten Gottesdienst halte.
\quellseite{525}

\section{Summarium}

Ist ein Exempel, wie das Königreich Juda auch Abgötterei treibet, und an andern Herrn, denn zu Jerusalem, opfert. Um solcher Sünde willen erwecket Gott den König aus Egypten, der plündert den Tempel und führet den Schatz daraus mit sich etc. Uns zur Warnung, wenn wir mit Frieden wollen leben, daß wir uns Gottes und seines Worts halten sollen. Wo man aber unrechte Lehre und Abgötterei hat, da soll Krieg und alles Unglück sich mit Haufen finden.

\section{Votum}

Wir danken Gott, dem HErrn, dem Vater aller Gnaden und Barmherzigkeit, daß wir sein heilig Wort und rechten Gottesdienst haben; der behüte uns ferner für aller Abgötterei, und regiere uns, daß wir ihn über alle Dinge fürchten, lieben und vertrauen, bis zu unserm seligen Ende, Amen.

\chapter{Das 15. Kapitel}

\section{Vorrede}

Die zeitlichen Strafen, damit der gerechte Gott Abgötterei und falsche Lehre und Gottesdienst heimsucht und straft, sind schwer; aber am allerschwersten sind die ewigen Strafen. Von denen predigt also die heilige Schrift 1 Cor. 6, 9.: Weder die Hurer, noch Abgöttischen, werden das Reich Gottes ererben. Und in der Offenbarung Johannis im 21. Capitel V. 8. wird gesagt: Der Abgöttischen und aller Lügner Theil wird sein in dem Pfuhl, der mit Feuer und Schwefel brennet. Demegegen, daß nicht zeitliche und ewige Strafen auf uns kommen, so laßt uns herzlich alle Tage beten, daß uns Gott für Abgötterei behüten und sein Wort und rechten Gottesdienst bei uns erhalten wolle. Davon höret auch an das 15. Kapitel.

\section[Regierung zweier Könige in Juda, Abiam und Asa;]{Regierung zweier Könige in Juda, Abiam und Asa; und zweier in Israel, Nadab und Baasa}
\quellseite{526}

a Miplezeth) wer der Miplezeth gewesen, ist ungewiß. Etliche sagen, er sei der Abgott Priapus gewesen.

\section{Summarium}

Dies ist eine Historia, in welcher uns zwei schöne Exempel werden angezeigt. Erstlich, daß der König Abiam nur drey Jahr regiert zu Jerusalem, V. 1.; sein Sohn Asa aber ein und vierzig, V. 10. Dieser Asa hat ein sonderlich fürstlich Werk geübet, daß er wehret dem abgöttischen Wesen und den Abgott Miplezeth ausrottet, V. 13. Denn einem Regenten (soll er anders Glück haben) stehet sonderlich zu, daß er darob sei, daß Abgötterei und falsche Lehre nicht geduldet werde. Denn wo Abgötterei und falsche Lehre ungehindert geht, daselbst schicket Gott zur Strafe Krieg und alles Unglück, wie man hie höret mit Baesa, dem Könige Israel, V. 17.
\quellseite{527}

\section{Summarium}

Ist eine Historia, in welcher uns fürgehalten wird, wie Gott mit den gottlosen Regenten, so Abgötterei und falsche Lehre handhaben, pflegt umzugehen. Jerobeams Sohn, Nadab, regiert nur zwo Jahr über Israel und wird erstochen, V. 25. Baesa, welcher an seine Statt König ward, wird mit Krieg überfallen vom Könige aus Syrien, und endlich sammt seinem ganzen Geschlecht ausgerottet, zur Strafe seiner Abgötterei, und daß er den König Nadab erwürget hat. Also ist dieß Buch durchaus ein rechter Spiegel, darin sich alle Regenten mit allem Fleiß schauen sollen, auf daß sie in Gottesfurcht wandeln und also dem Unrath entfliehen mögen, der über die gottlosen Regenten allweg kommen ist.

\section{Votum}

Der treue GOtt verleihe, daß alle Obrigkeit sein Wort liebe und fördere, und der Abgötterei und falscher Lehre steure und wehre; und regiere uns, daß wir alle Abgötterei fliehen und durch Fürbitz uns und den Unsern nicht Strafen zuziehen, Amen.

\chapter{Das 16. Kapitel}

\section{Vorrede}

Wie Gott das jüdische Königreich wegen der Abgötterei gestraft hat: also wird diese letzte Welt auch ein Schwereres noch ausstehen müssen, wegen der heimlichen und öffentlichen Abgötterei, die bisher getrieben worden und noch im Schwange gehet. Derwegen, daß Gott in gemeinem Trippel und Strafen unser gnädiglich wahrnehme, so laßt uns Buße thun in Reue und Leid über unsere Sünde, in wahrem Glauben und in Gehorsam gegen Gottes Geboten leben, und beten mit dem Könige Zedekia: Das Wort des HErrn ist gut, das du sagest; es sei nur Friede und Treue, weil ich lebe, 2 Kön. 20, 19. Wir wollen aber gleichwohl auch beten, daß Gott an unsern Nachkommen seiner Gnade und Barmherzigkeit nicht vergesse; und wolle ihm allzeit einen heiligen Samen, und ein Häuflein, das ihn erkenne, lobe und preise, und ihm diene in Heiligkeit und Gerechtigkeit, unter ihnen erhalten. Davon höret auch an das 16. Kapitel.

\section{Von vier Königen in Israel: Ela, Simri, Amri und Ahab}
\quellseite{528}

\section{Summarium}

Ist eine Historia, wie Gott das Königreich Israel um der Abgötterei willen mit Aufruhr strafet, V. 9., neben dem, daß die gottlosen Könige einer nach dem andern erstochen wird, um ihrer Sünden willen, daß sie die Abgötterei nicht wehreten, sondern mehreten, V. 10. Zum Exempel aller Obrigkeit, so sie ihr selbst glückselige Regierung und der Unterthanen Wohlfahrt
\quellseite{529}
 begehren, daß sie sich hüten für Abgötterei und Gott und sein Wort für Augen haben, und den rechten Gottesdienst (unangesehen, was andere thun, oder ihnen darüber begegnen mag) mit Ernst fördern.

\section{Votum}

Verleihe GOtt, daß sein heilig Wort unserer lieben Obrigkeit fürleuchte, auf daß ihre Regierung, wie bisher, also ferner, allzeit zur Erhaltung und Förderung seiner Kirchen und den Unterthanen zu Heil und Seligkeit in Christo JEsu, unserm HErrn, diene, Amen.

\chapter{Das 17. Kapitel}

\section{Vorrede}

Wer Gott ehret, den will Gott wieder ehren und ernähren, und wenn die Noth noch so groß wäre. Denn Gott über allen menschlichen Verstand die Seinen erhalten kann, etwa ganz und gar ohne Mittel, etwa durch geringe und kleine Mittel, wie Gottes Wort und die Erfahrung bezeuget. Derwegen sollen wir thun, als die lieben Kinderlein zu thun pflegen: dieselben verrichten mit Fleiß, was sie ihre Eltern heißen, und sorgen nicht, was sie werden zu essen und zu trinken haben, sondern sie befehlen diese Sorge den Eltern und denken, wenn die Zeit des Essens herbei kommen, so werden die Eltern wohl darauf gedacht haben; also sollen wir dem fleißig nachgehen, was uns Gott, unser himmlischer Vater, zu thun befohlen hat, sollen an seinen lieben Sohn glauben, in Gehorsam seiner heiligen Gebote leben, treulich und emsig unsern Beruf treiben, und auch Gott dafür sorgen lassen, wie er uns darinnen ernähren und erhalten werde. Ein jeder diene auch seinem Nächsten in seinem Stande und mit seinem Vermögen; so wird ihm Gott auch dienen in seinem Nothstande und Anliegen, mit Kraft und väterlicher Treue. Dies soll auch erinnern das 17. Kapitel.

\section[Elia in der Theurung von Raben gespeiset; wecket]{Elia in der Theurung von Raben gespeiset; wecket zu Zarpath einen Todten auf}
\quellseite{530}

\section{Summarium}

Hie gehet die Historia an von dem großen Propheten Elia, welchen Gott zu der selbigen Zeit in Israel hat erwecket, da die Abgötterei am größten war. Eben wie bei uns das Evangelium ist an Tag kommen, da das Pabstthum und der falsche Gottesdienst am mächtigsten war. Solches alles thut Gott darum, weil er die Abgötterei endlich strafen will und muß, daß doch die Frommen gewarnet und erhalten werden. Es ist aber hie wieder ein sonderlich Exempel, daß Gott mit einer Theurung länger denn drei Jahr das Land um der Abgötterei willen gestraft habe, V. 1. In solcher Noth gedenkt Gott an die Seinen, wie hie die Raben dem Elia Brod und Fleisch zutragen, V. 6. Und der armen Wittwe durch Eliam ihr Mehlfaß und Oelkrug also gesegnet wird, daß sie die ganze Theurung über genug hatte, V. 14. Denn Gott will den Theil mit der That und Wahrheit behalten, daß er seine Gläubigen wolle ernähren in der Theurung.

\section{Votum}

Wir danken GOtt, unserm lieben Vater, der uns speiset und ernähret; der wolle uns ferner für Theurung und andern wohl verdienten Strafen behüten, und wolle die seinen an Leib und Seele, Gut und Ehre, die seiner Kirchen Gliedmaßen in ihrer Noth sich treulich annehmen, Amen.

\chapter{Das 18. Kapitel}

\section{Vorrede}

Der erste Theil des 18. Kapitels des ersten Buchs der Könige hält uns für mancherlei Trost und Lehre. Der eine Trost ist: Wenn Gott die Gottlosen und Bösen mit Strafen im Zorn und Grimm heimsucht, so erhält er oft väterlich und wunderlich die Gläubigen und Frommen. Der andere Trost ist, davon auch Christus lehret im 11. Kapitel Johannis, des heiligen Evangelisten, da er sagt V. 9.: Wer des Tages wandelt, der stößet sich nicht, das ist, wer seines Berufs abwartet, im Glauben und mit Gebet zu Gott, der gefällt Gott wohl und dem hilft er in aller Gefahr. Und da er schon etwas Noth und Widerwärtigkeit leidet, dennoch überwindet er endlich und seine Lehre, Rath und Thaten sind dem Nächsten nütze und heilsam. Die eine Lehre aber ist: Wer von Gott gestraft wird, der lege die Schuld nicht auf andere Leute, sondern erkenne, daß er solche Strafe mit seinen Sünden selber verdienet habe, und thue wahre Buße, so wird ihm Gott auch aus Barmherzigkeit helfen. Die andere Lehre ist: Ein jeder Christ soll Gott herzlich anrufen, daß er uns mit Hunger des Leibes oder nach Brod, oder auch nicht mit Hunger nach seinem heiligen Wort zu hören, nicht strafe, wie er also gestraft hat den König Ahab, und den größten Theil des Volkes im israelitischen Königreiche, sondern daß er dem Leibe geben wolle das tägliche Brod, und der Seelen auch ihre Speise, nämlich das Wort, das aus seinem Munde gehet, und das eine Kraft ist, die da selig macht alle, die daran glauben.

\section{Der Prophet Elia schlachtet die BaalsPfaffen}
\quellseite{531}

\section{Summarium}

Dies ist erstlich ein fein Exempel, wie ein frommer Mann Obadja am Hof des gottlosen Königs Ahab ist, V. 3., durch welchen hundert Propheten erhalten und heimlich in der großen Theurung ernähret werden, V. 4., auf daß man sehe, wie Gott durch mancherlei Mittel zur Zeit der Noth die Seinen erhalten kann.

Zum andern ist sonderlich zu merken, wie die Gottlosen alles Unglück, das ihnen ihrer Sünde halben widerfähret, auf die Frommen legen. Denn Ahab gibt hie dem Propheten Elia die Schuld, daß Theurung und allerlei Unrath im Lande ist, V. 18. Eben wie zu unserer Zeit die blinden verstockten Papisten das liebe heilige Evangelium schänden und lästern, in dem, daß sie ihm Schuld geben, als sollte der Bauern Aufruhr und ander Unglück, Ketzerei, Krieg, Theurung daraus entstanden sein etc. Aber Elias spricht: Es sei solches alles der Sünden Schuld, daß Ahab und das Volk Gottes Gebot verlassen und dem Baal haben nachgewandelt.

\section{Votum}

Der treue Gott sei uns gnädig und barmherzig, und helfe, daß ein jeder Mensch seine Sünde erkenne, die Gott abbüße, davon ablasse, und nach Gottes Geboten in allem Gehorsam lebe und wandele, Amen.

\section{Der andere Theil des 18. Kapitels}

\section{Vorrede}

Die Prediger, die recht lehren und wohl leben, die springen nicht allerdings aus einem Steine oder Busche, wenn man dran schläget, wie die gottlose Welt spottet; sondern sie sind gewiß eine Gabe Gottes. Darum denn unser lieber Gott höchlich zu bitten ist, wie der Herr Christus selber lehret Luc. 10, 2., da er sagt: Die Ernte ist groß, der Arbeiter ist wenig; bittet den Herrn der Ernte, daß er Arbeiter aussende in seine Ernte. Der Herr dieser Ernte ist der Herr Christus selbst, der seinen Samen des Worts läßt ausstreuen hie auf Erden, ihm eine reiche Ernte unter dem menschlichen Geschlechte anzurichten; die Arbeiter sind alle treue Lehrer und Seelsorger, welchen, da sie den Samen göttlichen Worts unter den halsstarrigen und widerstrebenden Leuten sollen ausstreuen, es wohl so sauer wird, als irgend einem Schnitter oder Ernten-Knechte. Daß der Herr Christus solche Arbeiter aussende in seine Ernte, soll man bitten, das ist, daß er sie gebe und insinuire mit seinem heil. Geist von oben herab, und lege ihnen selbst das Wort in den Mund, damit sie Frucht und Nutzen schaffen. Helfe nun Gott, daß wir diese Erinnerung wohl in Acht nehmen und recht gebrauchen. Amen. Darauf höret ferner an den andern Theil des 18. Kapitels.
\quellseite{532}

a Hinkten die falschen Heiligen, wenn sie andächtig sein wollen, werfen sie den Kopf zu beiden Seiten, wie eine Gans gehet, daß es scheinet, wie ein hinkender zu beiden Seiten hin und her wackelt. Daher auch das Paßah kommen, 2 Mos. 12, 12. Da GOtt in Egypten durchging, und hinket, das ist, er schlug zu beiden Seiten, hie und da, wie ein trunkener Mann gehet. Sic supra: Ut quid claudicatis in duas partes? Significat impios esse vagos, et duplices animo, nihil certi habere, Eph. 4, 14. Inde etiam saltare dicuntur tales, quia motu capitum similes sunt salientibus et vere capitibus saltant. Das ist: so oben: Was hinket ihr auf beiden Seiten? zeigt an, die Gottlosen seien wankend und von doppeltem Herzen, haben nichts Gewisses, Eph. 4, 14. Daher sagt man auch von solchen, daß sie tanzen, weil sie mit der Bewegung der Häupter ähnlich sind dem Tanzenden und in der That mit den Häuptern tanzen.
\quellseite{533}

\section{Summarium}

Ist eine schöne Historia, wie Gott dem frommen Propheten Elia beistehet, und gibt ihm ein öffentlich Zeugniß seiner Lehre und Berufes, dagegen aber die Propheten Baals mit Schaden bestehen und werden um der Abgötterei willen alle erwürget. Auf daß man lerne, sich mit Ernst zum rechten Gott und seinem Wort halten, will man anders der leiblichen Strafe und dem ewigen Zorn Gottes und Verdammniß entlaufen.

Daß aber Elias die falschen Propheten selbst erwürget, thut er darum, daß es Gott mit einem öffentlichen Zeichen beweiset hatte, daß sie Verführer und Gotteslästerer wären, die dem Gesetz nach mit dem Tode mußten gestraft werden. Solches, da es der König nicht wollte thun, that es Elias, nicht aus eigenem Fürnehmen, sondern aus einem besondern prophetischen Geist und Muth, den er von Gott hatte. Im Evangelio vergleicht Christus Johannem den Täufer mit dem Elia, nicht daß er auch das Schwert geführet habe, wie Elias, sondern daß er so mit einem hohen Geist und trefflichem Ernst sein Amt geführet, die Hohenpriester und Phar isäer ohne Scheu gestrafet und die Leute zum rechten Gottesdienst gebracht habe.

\section{Votum}

Der treue GOtt regiere uns alle, daß wir von seinem Wort nicht weichen, sondern dem allzeit anhangen und gehorsamlich folgen; gebe auch allzeit christliche Lehrer und Seegen zu ihrem Amt, und steure allen falschen Propheten und Predigern, Amen.

\chapter{Das 19. Kapitel}

\section{Vorrede}

Fromme Zuhörer sind ihren Predigern verbunden, daß sie die hören sollen. Nicht allein aber, wenn die Prediger reden, was den Zuhörern wohlgefällt (wie denn St. Paulus geweissagt hat, daß am Ende der Welt der Zuhörer gar viel sein und eine Zeit sein wird, da sie die heilsame Lehre nicht leiden werden, 2 Tim. 4, 3.); sondern auch, wenn die Prediger wider die Zuhörer reden, sie dadurch ihrer Missethat zu erinnern und zu bekehren: also sollen die Zuhörer ihre Prediger hören, nicht allein, wenn sie lauter Evangelium und Gnade lehren, sondern auch, wenn sie Mosen und das Gesetz predigen. Das drücke ein jedes in sein Herz zum ewigen Gedächtniß. Denn wehe den Zuhörern, welche die Wahrheit von ihren Seelsorgern trennet und scheidet; wie die Galater die liebe Wahrheit, die ihnen zum Besten gepredigt ward, also von St. Paulo trennet, daß sie ihm todfeind worden, Gal. 4, 16. Dies soll auch erinnern das 19. Kapitel.

\section{Elias Flucht, Fasten, Trost und Nachfolger}
\quellseite{534}

\section{Summarium}

Ist eine Historia, in welcher man siehet, wie allweg die Propheten um des Worts und rechten Gottesdienstes willen sich haben leiden müssen; und hilft sie nichts, obgleich Gott mit öffentlichen Wunderzeichen vom Himmel bezeuget hat, daß ihre Lehre und Wesen recht sei. Solches macht die Heiligen oft kleinmüthig; wie Elias hie vor großer Angst und Traurigkeit begehrt zu sterben, V. 4. Aber Gott tröstet sie und hilft ihnen aus aller Noth.

Zum andern ist ein fein Exempel hie, daß Gott die sieben tausend, so Baal nicht gedienet haben, behüten will, V. 19., da er den andern, so dem Baal dieneten, dräuet, sie sollen durchs Schwert Hasael, Jehu und Elisa fallen, V. 17. Denn Gott will die Abgötterei nimmermehr ungestraft lassen.

\section{Votum}

Der barmherzige, gütige GOtt tröste, stärke und erhalte alle treue Prediger und Lehrer in ihrem schweren und gefährlichen Amt; schütze und segne auch alle fromme Zuhörer, und stürze alle Verächter und Verfolger seines heiligen Namens, Wortes und gläubiger Bekenner, Amen.

\chapter{Das 20. Kapitel}

\section{Vorrede}

Zwei nötiger Lehren sollen wir uns heute erinnern. Die erste ist: Wenn Gott die Menschen mit einer kleinen Strafe heimsucht, so folget nicht lange darauf eine große, wenn man nicht wahre Buße thut, und sich zu Gott ernstlich bekehret. Dies bestätiget die Weisheit Gottes in Sprüchen Salomonis, im 1. Kapitel, V. 24., da sie saget: Ich recke meine Hand aus, und niemand achtet darauf, und lasset fahren allen meinen Rath, und wollet meiner Strafe nicht; so will ich auch lachen in eurem Unfall und euer spotten, wenn über euch kommt, das ihr fürchtet. Wenn über euch kommt, wie ein Sturm, das ihr fürchtet und euer Unfall als ein Wetter, wenn über euch Angst und Noth kommt. Die andere Lehre ist: Wer alles haben will, der behält endlich gar nichts. Davon zeuget Sirach im 11. Kapitel V. 10. und spricht: Mein Kind, stecke dich nicht in mancherlei Händel; denn wo du dir mancherlei fürnimmst, wirst du nicht viel daran gewinnen. Wenn du gleich fast darnach
\quellseite{535}
ringet, so erlangest du es doch nicht, und wenn du gleich hie und da fliehest, so kommst du doch nicht heraus. Demgegen laßt uns dieser Lehren gebrauchen zur Buße, daß Gott unser mit seinen Strafen verschone, und zur Warnung, daß ein jeder also das Seine mehre, damit es ihm Gott auch bewahre und erhalte. Davon höret an den ersten Theil des 20. Kapitels.

\section[Zweyfacher Krieg und Sieg Ahabs wider Benhadab, den]{Zweyfacher Krieg und Sieg Ahabs wider Benhadab, den Syrer}

a Nicht rühmen) das ist: Er spreche nicht hui, ehe er über den Berg komme; denn wer genommen hat, der lege den Harnisch ab, und mag sich rühmen; wer ihn aber anlegt, hat drum noch nicht gewonnen.

\section{Summarium}

Sie ist eine Historia, wie Gott immerdar mit der Strafe innen hält, dieselbige nur von ferne weiset und wartet, ob man sich nicht bessern wolle. Also lässet er die Syrer ins Land fallen, aber doch nicht wider Israel siegen, v. 1. Zum sonderlichen Exempel den Tyrannen, wenn sie mehr, denn billig ist, wollen haben, daß sie lernen, daß Gott solche Unbilligkeit strafen wolle. Denn
\quellseite{536}
Ahab wollte den König zu Syrien zu einem Oberherrn erkennen; da aber der König zu Syrien sich an solchem nicht wollte genügen lassen, sondern alles haben, ward solche Tyrannei gestrafet.

\section{Votum}

Gott, unser lieber Vater, regiere uns, daß wir, seine Kinder, durch seine kleinen Ruthen und Heimsuchung fromm werden, auf daß uns nicht große verderbliche Straßen folgen; und helfe, daß ihm ein jeder also an dem Seinen genügen lasse, damit ihn Gott ferner an Gut und Ehre segne, aus Gnaden, Amen.

\section{Der andere Theil des 20. Kapitels}

\section{Vorrede}

Gottes-Lästerung und das Vertrauen auf Menschen-Hülfe sind zwei schändliche und schädliche Laster. Ueber die Gottes-Lästerung klaget der Prophet Jesaias im 1. Kapitel V. 4. und sagt: O weh des sündigen Volkes, des Volkes von großer Missethat, des boshaftigen Samens der schädlichen Kinder, die den Herrn verlassen, den Heiligen in Israel lästern, weichen zurück. Das Vertrauen auf Menschen-Hülfe strafet der Prophet Jeremias im 17. Kapitel V. 5. also: So spricht der Herr: Verflucht ist der Mann, der sich auf Menschen verläßt, und hält Fleisch für seinen Arm, und mit seinem Herzen vom Herrn weicht. Dagegen wollen wir Gottes Allmacht und in der Noth seinen Beistand loben und preisen mit Mund, Herz und gottseligem Wandel, und wollen in allem betrübten Zustande allein auf Gott trauen und bauen, Hülfe bei ihm suchen und auf seine Hülfe allzeit beständig warten. Denn es ist gut, sagt der königliche Prophet David Ps. 118, 8., auf den Herrn vertrauen, und nicht sich verlassen auf Menschen. Dagegen bringet Schaden, auf Fürsten sich verlassen, die Menschen sind und die ja nicht können helfen; wie abermals David spricht Ps. 146, 3. Dies lehret auch der andere Theil des 20. Kapitels.
\quellseite{537}

\section{Summarium}

Dies ist erstlich eine feine Historia, wie die Syrer geschlagen worden, nicht darum, als wäre ihr Gegentheil so fromm, sondern darum, daß die Syrer Gott gelästert und gesagt hatten, der Israeliten Gott würde ihnen nicht helfen können, wenn sie auf der Ebene eine Schlacht hätten. Auf daß man sehe, daß Gotteslästerung nicht ungestraft bleibe.

Zum andern ist eine Historia hie, in welcher man siehet, daß Ahab sich versündiget, daß er des Sieges, welchen ihm Gott gegeben hatte, nicht recht brauchet, V. 34. Er sollte den gottlosen König aus Syrien getödtet haben; so tröstet er sich seiner Macht und Hülfe, und macht einen Bund mit ihm, und vergisset also, daß er sich allein auf Gottes Hülfe sollte verlassen. Darum muß solcher Sünde sein ganzes Geschlecht hernach entgelten, V. 42. Es ist aber solches eine Fürsten- und Regenten-Sünde; die lassen sich bedünken, wo man mit Macht gefasset sei, da habe es nicht Noth. Aber Gott lehret sie anders, daß es nicht an der Macht, sondern an seiner Gnade und Segen liegt.

\section{Votum}

Der treue Gott behüte uns, daß wir ihn und seine Werke ja nicht lästern; und helfe, daß wir in Nöthen allein auf ihn sehen all unsere Zuversicht, und ihn loben und preisen, zeitlich und ewig, wegen seiner väterlichen Hülfe und treuen Errettung, Amen.

\chapter{Das 21. Kapitel}

\section{Vorrede}

Keinem armen Sünder, obgleich seine Sünden blutroth wären, und er des Tages, nicht allein sieben, sondern siebenzig mal sieben mal, gesündiget hätte, wird die Gnaden Thür verschlossen, sondern stehet ihm allezeit offen, wofern er anders den rechten Weg der Buße dazu trifft, sich demüthiget und zu Gott bekehret. Das bezeugen viele trostreiche Sprüche und gnadenreiche Exempel. Ein fürnehmer Spruch ist, daß von Gott Micha, der Prophet, im 7. Capitel V. 18. sagt: Wo ist solch ein Gott, wie du bist? Der die Sünde vergibt und erlässet die Missethat den Uebrigen seines Erbtheils, der seinen Zorn nicht ewiglich behält, denn er ist barmherzig. Er wird sich unser wieder erbarmen, unsere Missethat dämpfen und alle unsere Sünde in die Tiefe des Meers werfen. Ein schön Exempel aber, daß sich Gottes Barmherzigkeit erstrecke über alle Wolken, und erfülle Himmel und Erden, ist der bußfertige Sünder, der König Ahab. Davon das 21. Capitel.

\section{Geduld Tyranni Ahabs und Jesabels wider Naboth}
\quellseite{538}

a Wenn du) Tu nunc facis regnum super Israel id est: Bist du es doch, der es macht, daß Israel ein Königreich ist, ohne dich wäre es ein Nichts.

b Gesegnet) das ist: Er hat gelästert; denn die Schrift das grausame Laster, das Gott lästern heißt, an diesem Ort aufs züchtigste nennet, wie auch Hiob 2, 9. stehet. Gott lästern aber hatte Mose bei dem Tode verboten. So war das Königreich Israel von Gott eingesetzt, daß den König lästern auch Gott betraf.
\quellseite{539}

c Sack) das ist, betrübte, geringe Kleider, als wenn man Leid trägt, oder arbeitet.

\section{Summarium}

Es ist erstlich eine Historia, wie Ahab sammt seinem Weibe, Isabel, Naboth, den Gerechten, ohne alle Ursache, wider Gott und Recht hat lassen erwürgen, darum, daß er dem Befehl Gottes nach bei seinem Erbe bleiben und dasselbige dem Könige nicht verkaufen wollte, V. 3. Denn obgleich weltliche Obrigkeit Macht hat über der Unterthanen Leib und Gut; so soll es dennoch seine Maße haben, daß ein Unterschied bleibe zwischen einem ordentlichen Regenten und Tyrannen. Ein Tyrann nimmt, was ihm gefällt, und fraget nicht, ob ers Fug oder Recht habe. Ein Regent aber bleibt bei seinem ordentlichen Einkommen, und wo es die Noth nicht sonderlich erfordert, legt er keine neue Beschwerden auf. Solche Gleichmäßigkeit hält Ahab mit dem Naboth nicht, sondern führet tyrannisch. Ob nun wohl der fromme Naboth darum zu Boden gehet, so muß doch Ahab auch darum gestraft werden.

Zum andern ist ein sonderlich Exempel hie, daß Ahab, nachdem der Prophet Elias seiner Sünden halben ihn strafet, V. 19., Buße thut und Gott um solchen Demuth willen die Strafe bei Ahabs Leben auffhält, V. 27. Solches Exempel dienet uns dazu, daß wir im Unglück und wenn Gottes Strafe auf uns liegt, nicht sollen verzagen, sondern uns bekehren, vom Bösen ablassen und um Gnade bitten, so will Gott, wie mit Ahab, sich gnädig finden lassen, der doch ein sehr böser Mensch gewesen war.

\section{Votum}

Der barmherzige, gütige GOtt sei uns auch gnädig, vergebe uns unsere großen vielfältigen Sünden, und lindere die wohlverdienten Strafen; er wolle uns auch regieren, daß wir unserm Nächsten an seinem Leib und Gut keinen Schaden noch Leid nimmer thun, Amen.

\chapter{Das 22. Kapitel}

\section{Vorrede}

Daß die Gottlosen für den göttlichen Strafen sich nicht fürchten, ja auch dieselben nicht glauben, wenn sie ihnen angekündiget werden; diese Blindheit kommt vom Teufel, welcher die Gottlosen nicht allein so sicher macht, sondern sie auch verhindert an der Buße, daß sie Gottes Strafen nicht entrinnen mögen. Denn er weiß wohl, daß Gott durch den Propheten Esaías im 45. Capitel V. 22. saget: Wendet euch zu mir, so werdet ihr selig aller Welt Ende, denn ich bin Gott und keiner mehr. Und wie wohl sagen, daß alle, die dem HErrn widerstrecken, müssen zu Schanden werden. Verwegen wollen wir dem Teufel nicht Raum geben, wollen für Sündigkeit uns hüten, nicht in Sünden verharren und verstarren, und wollen Gottes zeitliche Strafen gebrauchen zur Bekehrung und gottseligem Wandel, daß wir nicht härter allhier und dort ewig gestraft werden wegen unserer Sünden. Darauf höret ferner an den ersten Theil des 22. Capitels.

\section{Ahabs Untergang. Josaphats und Ahabs Regierung}
\quellseite{540}

\section{Summarium}

Dies ist eine wunderbare Historia, wie Gott verhänget, daß der König Ahab, so ohne das Gottes und seines Wortes nicht viel achtet, durch falsche Propheten verführet wird und endlich umkommen, wie ihm zuvor durch einen Propheten geweissaget, v. 6. Sein Bruder, der König zu Syrien, mit dem Ahab einen Bund machte und sich seiner Macht tröstete, gibt ihm seinen rechten Lohn, darum, daß er ihn hatte los gelassen. Hie sieht man, wie die falschen Propheten mit ihren Lügen toben, und die Frommen und rechtschaffenen Propheten sich von ihnen leiden müssen. Aber es währet beides nur eine Zeitlang.

\section{Votum}

Gott steure in seinem Zorn allen falschen Lehrern, und gebe den rechtschaffenen treuen Predigern Geduld unter dem Kreuz und Verfolgung, und helfe, daß sie ihr Amt verrichten zu seinen Ehren und ihren Zuhörern zu Nutz und Seligkeit, Amen.

\section{Der andere Theil des 22. Kapitels}

\section{Vorrede}

Die ganze heil. Schrift neben allen bewährten Historien bezeuget, daß allezeit greuliche Sicherheit bei denjenigen in vollem Schwange gegangen ist, welche Gott der HErr mit seinem heiligen göttlichen Wort reichlich besucht und begabet hat; sonderlich aber ist dieses an den Zeiten des Volkes Gottes, der alten Israeliten, zu sehen, in was großer Sicherheit sie, da sie
\quellseite{541}
Gottes Wort überflüssig gehabt, gelebt haben. Daher rücket Moses im 5. Buch im 32. Kapitel V. 15. seinen Juden dieß Laster heftig auf: Da Israel (spricht er) feist und satt ward, wurde er geil. Er ist dick und stark worden, und hat den Gott fahren lassen, der ihn gemacht hat; er hat den Fels seines Heils gering geachtet. Also wenn die Sicherheit unter den Leuten am größten ist, so ist die Strafe Gottes am nächsten mit dem Garaus, das ist, mit dem letzten Untergang und Verderben; als deß ein Exempel ist der König Ahab mit seinem Untergang, und lehret also, daß ein jämmerlich Ende nehmen alle Gottlose und Abgöttische, welche Gott und sein Wort samt treuen Predigern verachten und verfolgen. Derwegen schaffet, daß ihr selig werdet mit Furcht und Zittern. Denn Gott ist, der in euch wirket, beide das Wollen und das Vollbringen. Und dies geschehe ja in seiner Furcht, und nach seinem Wohlgefallen. Davon höret auch an den andern und letzten Theil des 22. und letzten Kapitels.

a. Da das Schwert anhänget von der Achsel über her, bis auf die Hüfte.

b Huren) da die Huren waschen. Das ist, Canponas lotrices.
\quellseite{542}

\section{Summarium}

Ist erstlich eine Historia, wie der fromme Prophet Micha, darum, daß er die Wahrheit sagt, vom gottlosen König Ahab gefangen wird, V. 27. Zum andern, daß Ahab im Kriege umkommt, V. 37. Zum dritten ist eine sonderliche feine Historia, daß Josaphat die Abgötterei, so noch im Lande war überblieben, vollends ausrottet, V. 43. Denn dies stehet sonderlich christlicher Obrigkeit zu, daß sie die Abgötterei nicht dulden, sondern Gott zu Ehren und den Leuten und gemeinem Nutz zu gut, dieselbige mit Ernst abschaffen und strafen sollen.

\section{Votum}

Gott tröste und erhalte alle treuen Lehrer und Seelsorger unter ihrem Kreuze, strafe die Verfolger seiner Kirche und Christenheit; und behüte uns ferner für aller Abgötterei und andern Sünden wider Gott und wider die Menschen, Amen.

\part{Das zweite Buch von den Königen}

\chapter{Das 1. Kapitel}

\section{Vorrede}

Der königliche Prophet David vergleichet Gottes Wort einem Stecken und Stab, Ps. 23, V. 4. Nun dienet ein Stecken oder Stab einem Wandersmann zu vielem. Erstlich, so er gefallen, daß er dran aufstehe. Zum andern, daß er sich dran halte, damit er nicht auf schlüpfrigen Wegen strauchele. Zum dritten, daß er sich damit der bösen Hunde, die ihn anfallen wollen, erwehre. Zum vierten, daß er sich drauf stemme, und dran ruhe, sonderlich da er schwer zu tragen hat: Also dienet uns auch Gottes Wort, erstlich, so wir allhier, auf den schlüpfrigen Wegen gefallen, daß wir dran aufstehen von Sünden, und uns, durch den Glauben an Christum, wieder aufrichten. Zum andern, daß wir uns dran halten, damit wir nicht fallen in Sicherheit, gottlos Wesen und Verzweiflung. Zum dritten, daß wir damit Sünde, Tod, Teufel, Hölle, als den bösen Hunden, einen Streich bieten können. Zum vierten können und sollen wir Gottes Wort unter der schweren Last des Kreuzes gebrauchen und darauf ruhen und uns damit trösten. So helfe Gott, daß uns sein Wort auch ein Stecken und Stab sei, im Leben und Sterben, Amen. Darauf höret nun an das 1. Kapitel.

\section{Ahasias Krankheit. Elias Wunderwerke}

\section{2 Kön. Kap. 1. v. 1}
\quellseite{543}

\section{Summarium}

Ist erstlich eine Historia, wie der gottlose König Ahazia dem Exempel seines Vaters Ahab folget, V. 53. Weil er nun so gottlos war, daß er bei der Heiden Göttern Rath und Hülfe suchet, V. 2., und nicht bei dem Gott Israel, strafet ihn Gott, daß er in seinem Hause fällt und des Falls hernach stirbt. Auf daß wir uns lernen für falscher Lehre, unrechtem Gottesdienst und Zauberei hüten, und desto fleißiger uns zum Wort Gottes und rechtem Gottesdienst halten.

Denn weil die Abgötterei verursachet, daß Gott dadurch gelästert wird, kann sie Gott in die Länge nicht ungestraft lassen. Sonderlich aber an den Regenten, so dazu helfen, welche doch Amts halben mit allem Ernst der Abgötterei wehren sollen.

\section{Summarium}

In eine sonderliche Historia, daß der Prophet Elias heißet Feuer vom Himmel fallen über die zween Hauptmänner des Königs Ahazia, der ein jeglicher fünfzig Mann bei sich hatte. Eben wie der Apostel St. Petrus, aus göttlicher Kraft, allein mit einem Worte Ananiam und sein Weib Saphiram um der öffentlichen Lügen willen tödtet. Diese zween Hauptmänner kommen mit einem Befehl vom abgöttischen Könige, und ist gewiß, daß sie auch abgöttisch gewesen sind. Weil sie aber Eliam einen Mann Gottes heißen, V. 10., und doch den Gott, welchen Elias ehrete und von ihm predigte, nicht für den rechten Gott hielten — sie waren Baals Diener, wie ihr König — will Elias an ihnen mit solchem Wunderwerk, daß das Feuer vom Himmel fällt, beweisen, daß seine Lehre und Gottesdienst sei recht und ihrer Propheten Lehre und Gottesdienst sei unrecht und falsch. Daß also dies Wunderwerk nicht dahin gehet, daß Elias sich räche, sondern daß Gottes Name geehret und der rechte Gott erkennet werde.

\section{Votum}

Der Sohn Gottes, unser HErr JEsus Christus, regiere uns, daß wir allezeit falsche Lehre, unrechten Gottesdienst und Zauberei fliehen, und dagegen, daß wir seinem Wort folgen, ihm in Heiligkeit und Gerechtigkeit dienen, und daß wir ja mit den Gottlosen nicht umkommen und verderben, Amen.
\quellseite{544}

\chapter{Das 2. Kapitel}

\section{Vorrede}

In den Geschichten der heiligen Apostel im 14. Kapitel V. 22. sagt St. Paulus: Durch viel Trübsal müssen wir ins Reich Gottes gehen. So sagt man: Es sei ein gut Kraut, das heißt, achte es nicht, das heilet allerlei Schmerzen, wenn mans drauf leget und wohl darauf erwärmen läßt. Das rechte, achte es nicht, sollen wir aus dem himmlischen Würzgarten des göttlichen Worts nehmen, sollen es legen, in wahrem Glauben, auf unsere Noth, und den heiligen Geist bei uns wohl erwärmen lassen, alsdann wir in Christo getrost das mannigfaltige Unglück nicht achten, sondern werden uns Gottes Hülfe und des ewigen Lebens mit Freuden unter allem Kreuz trösten können. Ferner in der 1. Epistel an die Thess. 4, 8. spricht abermals St. Paulus: Wer verachtet (vernehmt die Prediger göttlichen Worts), der verachtet nicht Menschen, sondern Gott, der seinen heiligen Geist gegeben hat in euch. Ja der Herr Christus selber sagt auch Luc. 10, 16.: Wer euch (verstehet die Lehrer des Worts Gottes) verachtet, der verachtet mich; wer aber mich verachtet, der verachtet den, der mich gesandt hat. Davon erinnert auch das 2. Kapitel.

\section{Elias Himmelfahrt. Elisa's Wunderthaten}

a Von deinen Häupten) zum Häupten sein, heißt Meister und Lehrer sein; zu Füßen sein, heißt Schüler und Unterthan sein. Denn wenn der Lehrer lehret, sitzt er höher, denn die Schüler, daß er sie zu Füßen, und sie ihn zu Häupten haben. Also sagt St. Paulus Apost. 22, 3.: Er habe zu den Füßen Gamaliel das Gesetz gelernt. Und ist fast eine gemeine Weise der Schrift, also zu reden auf Hebräisch.

b Nicht wollte Elisa zwiefältigen Geist Elia haben, so es doch ein Geist ist, 1 Cor. 12, 4. in allem Heiligen, sondern einen zwiefältigen Mund desselbigen Geistes, daß er stärker und mehr predigen könnte, denn Elia, als er auch that.
\quellseite{545}

\section{Summarium}

Dies ist endlich eine Historia, wie Gott den Propheten Elia in einem Wetter von diesem Leben wegnimmt zum ewigen Leben, V. 2. 11. Zum Trost nicht allein jenen, so mit Elia recht geglaubt und den Baal nicht angebetet hatten, sondern auch uns allen, so da warten auf die selige Hoffnung und herrliche Erscheinung des großen Gottes, und unsers Heilandes Jesu Christi. Daß wir gewiß wissen und hoffen sollen, daß uns ein ewig Reich, nachdem wir aus diesem Jammerthal fahren, in jenem seligen Leben bereit sei, darin wir in ewiger Seligkeit mit unaussprechlicher und herrlicher Freude leben sollen.

Zum andern hebt sich hie an die Historia von dem Propheten Elisa, V. 12. Denn weil es mit dem Königreich Israel an dem war, daß es sollte um der Abgötterei willen gestraft werden, läßt ihnen Gott zuvor sein Wort reichlich dessen halben verkündigen. Wie aber eine große Sünde es sei, wo man Gottes Wort und desselben Verkündiger verachtet und unehret, siehet man hie an den Kindern, welche die wilden Bären zerreißen, darum, daß sie des alten Propheten spotten, V. 23., und zur Warnung, daß wir die Kirchendiener, um ihres Amts willen, lernen in Ehren halten.

\section{Votum}

Helfe der treue Gott, daß wir uns als gehorsame und dankbare Zuhörer seines Worts gegen ihn und unsere Seelsorger und Lehrer halten; und verleihe uns, daß wir alles Kreuz in Geduld überwinden, und nach diesem Leben auch das ewige Leben erlangen, um Jesu Christi, unsers Herrn willen, Amen.

\chapter{Das 3. Kapitel}

\section{Vorrede}

Gott sagt durch den Propheten Zephania, 1. Capitel V. 17.: Ich will den Leuten bange machen, darum, daß sie wider den Herrn gesündiget haben. Dieses ist erfüllet am König Joram. Dagegen spricht Sirach, 1. Kap. V. 19.: Wer den Herrn fürchtet, wenn er Trost bedarf, so wird er gesegnet. Das ist erfüllet am Könige Josaphat. Ferner im Ps. 15, 4. stehet: Ein Gläubiger achtet die Gottlosen nichts, sondern ehret die Gottesfürchtigen. Des hat sich gehalten der Prophet Elisa. Und St. Paulus Röm. 13, 2. sagt: Die der Obrigkeit widerstehen, werden über sich ein Urtheil empfahen. Das haben an sich die Moabiter erfahren. Er erzählet Sprüche. Exempel lauten aus dem 3. Kapitel.

\section{Jorams Regierung. Krieg und Sieg wider die Moabiter}
\quellseite{546}

\section{Summarium}

Es ist erstlich eine feine Historia, in welcher man siehet, was für ein Unterschied sei zwischen den Frommen und Gottlosen, zur Zeit der Noth. Die Gottlosen verzagen bald und können sich Gottes Worts nicht trösten, wie die Joram thut, V. 10., denn sie haben es verachtet und nicht gemerket. Aber die Frommen, wie Josaphat, haben Gottes Wort in Ehren, merken und behalten es, und wenn die Noth hergehet, halten sie sich daran und setzen ihr Vertrauen darauf, und finden derhalben reichen Trost darin, V. 12. Solches genießen hernach die Gottlosen auch.

Zum andern ists auch ein schön Exempel mit dem Propheten, daß er dem gottlosen König unter Augen sagt: Er wollte ihn nicht ansehen, wenn er nicht Josaphat ansähe. Warum er nicht seine Baals-Propheten um Rath fraget? V. 13. Solchen Ernst ist ein Prediger schuldig wider gottlose Obrigkeit zu brauchen, sintemal doch kein Ehren noch Vermaynen sonst an ihnen helfen will.
\quellseite{547}

a Vel quod ipsi Israel offensi, ista abominatione timerunt, ne Deus in ipsos quoque irasceretur. D. i. weil die Israeliten, selbst über solchen Greuel erzürnt, fürchteten, Gott möchte auch gegen sie selbst zornig werden.

\section{Summarium}

Es ist eine schöne Historia, wie Gott wunderbarer Weise schaffet, daß die Moabiter geschlagen werden. Solcher Sieg widerfäret dem Könige Israel und Edom allein um des frommen Josaphats willen.

Daß sie aber mit den Feinden so unbarmherzig umgehen, die Bäume umhauen, die guten Acker mit Steinen voll werfen, V. 25., geschieht nicht aus tyrannischem Fürnehmen, sondern darum, daß der Prophet Elisa solches ihnen zuvor gesagt hatte, daß es der Wille Gottes wäre.

Und hie stehet man wieder ein Exempel, wie es den Aufrührern gehet. Denn die Moabiter haben solche Strafe damit verdienet, daß sie von dem Könige Israel sind abgefallen.

\section{Votum}

Regiere uns der treue GOtt, daß wir ihm und unserer Obrigkeit schuldigen und willigen Gehorsam leisten; und helfe, daß wir uns in aller Noth seines heiligen Worts trösten und in wahrem Troste leben und sterben, Amen.

\chapter{Das 4. Kapitel}

\section{Vorrede}

Der königliche Prophet David sagt im 34. Psalm V. 20.: Der Gerechte muß viel leiden, aber der HErr hilft ihm aus dem allen. Das Kreuz ist die rechte Hof-Farbe den Christen, darin Christus seine Hofdiener kleidet. Gott salzet uns mit dem beißenden Salz des Kreuzes ein, daß nicht höllische Maden und verderbliche Sünden in uns wachsen und regieren. Gott höret uns im Ofen des Elends, daß wir nicht sinkend und faul, und demnach als untüchtig für die höllischen Hunde weggeworfen werden. Aber Gott hilft auch endlich von allem Kreuz, der da ist der einige Nothhelfer zur rechten Zeit, und der einige Meister zu helfen, Esaia 63, 1. Der alle Hülfe thut, die auf Erden geschieht, Ps. 47, 12. Weil dem also, daß uns alles Kreuz gut und nützlich ist, und daß uns Gott schuldlos will halten und endlich erretten und helfen, so sollen wir das Joch des HErrn Christi auf uns nehmen, unsere Seele mit Geduld fassen und auf Gottes Hülfe beständig warten. Davon hört auch an den ersten Theil des 4. Kapitels.

\section{Elisas fünf Wunderwercke}
\quellseite{548}

a (Ich wohne) das ist, ich habe zu Hofe nichts zu schaffen; ich wohne unter den Leuten alleine.

b Id est, More anni currentis, quo vivitur à cunctis animantibus. D. i. im Lauf dieses Jahres.

\section{Summarium}

Dies ist erstlich ein fein Exempel von einer armen Wittwe, in welchem man siehet, wie Gott, ob er gleich fromme Leute eine Zeitlang lässet Mangel leiden, dennoch endlich ihnen hilft und nicht in der Noth stecken läßt, V. 2.

Weil aber Gott der Wittwen hie solchen Segen läßt widerfahren, und heißet von dem Uebrigen die Schuldiger bezahlen, V. 7., siehet man, daß es Gottes Wille ist, daß, wer schuldig ist, dasselbe ehrbarlich und treulich bezahlen soll. Denn fremdes Gut ist durchaus verboten, daß es niemand soll inne haben. Daß aber Christus im Evangelio sagt, man soll leihen und nichts dafür nehmen, solcher Befehl gehet nicht dahin, daß man entlehnen und nichts soll wiedergeben, sondern daß die, so zu leihen haben, solches ohne ihren Nutzen und dem Dürftigen zu gut thun sollen. Denn das Entlehnte ist nicht eigen, und wo man kann, soll es wieder erlegt werden.

Zum andern ist ein schön Exempel von der reichen Frau zu Sunem, welche den Propheten Elisa beherberget und von Gott für solche Wohlthat einen Sohn erlanget. Denn Gott will sonderlich vergelten, was man denen zu gut thut, so sein Wort predigen.

\section{Votum}

Der treue, barmherzige GOtt helfe auch uns in aller Noth, und allen, die in Armuth stecken und andere Haus-Kreuze tragen; und regiere uns, daß wir auch Gerechtigkeit und Barmherzigkeit gegen unsern Nächsten üben, ihm zum Gehorsam und ihnen zu Hülfe und Trost, Amen.

\section{Der andere Theil des 4. Capitels}

\section{Vorrede}

Weil der Artikel von Auferstehung des Fleisches alle Furcht des Todes benimmt, und der heil. Geist dadurch lebendigen Trost in allen Anfechtungen in uns wirket; so bestätiget Gott diesen Artikel so reichlich mit Sprüchen und Exempeln des Alten und Neuen Testaments. Auch hat er derwegen denselben in vielen Creaturen abgebildet, als: wenn die Schnecken alt werden und fast verweset sein, werfen sie die alten Häute ab und werden wiederum aufs neue lebendig und jung. Die Seidenwürmlein werden auch wieder lebendig, ob sie schon eine Zeitlang todt gewesen sind. Alle Speise wird in das schönste rothe Blut verwandelt. Ein kranker Mensch wird gleich verzehret, bleich und erdenfarbig; aber wenn er wieder gesund wird, so wird er gleich wieder erquicket, fett und wohlgestaltet. Summa, wir sehen unterwärts, aufwärts, neben uns oder um uns, und sehen an Laub und Gras, Himmel und Erde, Menschen und Vieh, so sind sie unsere Haus-Postillen, die uns von der Auferstehung des Fleisches predigen. Gemeldeter Exempel eines erzählet auch der andere Theil des 4. Capitels.
\quellseite{549}

\section{Summarium}

Ist erstlich ein trefflich Exempel eines großen Glaubens an der Sunamitin, daß sie die Hoffnung hat, Gott werde ihren Sohn durch den Propheten vom Tod wieder zum Leben auferwecken, wie dann geschiehet.

Darnach sind zwei andere Exempel. Das erste, daß der Prophet Elisa die wilden Ranken, die vor Bitterkeit niemand essen konnte, wunderbarlich ändert, daß sie wohl schmeckten, V. 39. Das andere, daß er in der Theurung mit wenig Brod hundert Mann speiset, und noch übrig bleibt, V. 42. Solche Historien zeigen an, daß dennoch etliche fromme Leute im abgöttischen Königreich Israel sind gewesen, die an Gottes Wort gehalten haben, welchen er in der Theurung und sonst in andern Nöthen, um seines Wortes willen, geholfen hat; da die Ungläubigen haben müssen um ihrer Sünden willen Noth leiden. Uns zum Exempel, daß wir fromm sollen sein und hoffen, Gott werde es uns genießen lassen; gleich wie er die ungehorsamen bösen Buben mit der Strafe heimsuchen und ihrer auch nicht vergessen will.

\section{Votum}

Gott, der HErr, verleihe uns allen rechtschaffene Buße, erhalte uns in wahrem Glauben, speise uns an Leib und Seele, und sei unsere Hülfe in Noth und Tod, und erwecke uns auch zum ewigen Leben am jüngsten Tage, durch Jesum Christum, seinen lieben Sohn, unsern HErrn, Amen.
\quellseite{550}

\chapter{Das 5. Kapitel}

\section{Vorrede}

St. Paulus sagt Röm. 15, 9.: Daß die Heiden Gott loben um der Barmherzigkeit willen. Und Esaias spricht im 11. Kapitel V. 10.: Es wird sein die Wurzel Jesse, und der auferstehen wird zu herrschen über die Heiden, auf den werden die Heiden hoffen. Derer Sprüche ist ein Exempel Naeman, der Syrer. Ferner Salomo in Sprüchen im 21. Kapitel V. 6. sagt: Wer Schätze sammelt mit Lügen, der wird fehlen und fallen unter die, so seinen Tod suchen. Und Sirach spricht im 31. Kapitel V. 6.: Viele kommen zu Unfall um Geldes willen, und verderben drüber vor ihren Augen. Diese Sprüche sind erfüllet an Gehasi, dem Knecht Elisa. Die Exempel aber der erzählten Sprüche lauten aus dem 5. Kapitel.

\section{Naeman wird vom Aussatz gereiniget. Gehasi damit gestraft}

a Hand lehnet) das ist hebräisch geredet, wie wir Deutschen sagen: Er ist mir zur Hand, das ist, er ist um mich, thut und richtet aus, was ich ihm befehle und ich mich auf ihn verlasse.
\quellseite{551}

\section{Summarium}

Ist eine feine Historia, wie Naeman, der Syrer von seinem Aussatz rein und zum rechten Gott Israel bekehret wird, V. 14. Durch solche Historien hat Gott immerdar wollen anzeigen, daß auch die Heiden noch sollten zum Reich Gottes kommen und aus Gnaden ohne Gesetz selig werden.

Sonderlich aber ist hie zu merken, daß Gott seine großen Wunderwerke so gar einfältig und schlecht fürnimmt, daß sich die Vernunft muß darüber ärgern. Es gehe aber so schlecht und einfältig zu, als es immer wolle, wer Gottes Zusagung und Gnade gläubet und folget, der erfähret, daß sein Thun gewaltig und wunderbarlich hinaus gehet. Der Vernunft nach war der Jordan nichts besser, denn andere Wasser. Aber weil das Wort auf den Jordan weiset und auf andere Wasser nicht, wird Naeman rein, da er sich im Jordan badet. Es ist aber diese Historia auch ein Vorbild und Deutung der seligen Taufe, die im Jordan von Johannes angefangen und auf uns Heiden auch kommen ist, dadurch wir vom greulichen Aussatz der Sünden rein werden, durch den Glauben an den Tod und Verdienst Christi, wie es der heilige Basilius auch deutet.

b Wandelte nicht? das ist, hast du nirgend hin gewandelt; wie gehest denn zu, daß mein Herz wandelt, und war bei dem Wagen.

\section{Summarium}

Ist eine Historia, in welcher man siehet, wie Gott den Geiz und Lügen an des Propheten Knecht Gehasi straft, daß er und alle seine Nachkommen um solcher Sünden willen mit dem Aussatz gestraft werden, denn Gott ist allem Geiz feind. Sonderlich aber ist er dem Judas-Geiz feind, da man mit seinem heiligen Wort und Sacrament handthieret, daß man Geld, Gut, Herrschaft, Gewalt etc. dadurch erlange, wie der Erz-Judas, der Pabst, mit seinem Haufen thut, der aus Gottes Wort und den heiligen Sacramenten schlecht einen Jahrmarkt gemacht und sie um Geld verkauft hat.

\section{Votum}

Der allmächtige, barmherzige GOtt helfe heute auch an Leib und Seele allen Kranken und Kreuzträgern; und verleihe, daß wir uns für Geiz, Lügen und allen andern Sünden hüten, stürze und vertilge auch allen Judengeiz, Amen.

\chapter{Das 6. Kapitel}

\section{Vorrede}

Der heiligen Engel Schutz und das liebe tägliche Brod sind auch zwei große gnadenreiche Wohlthaten Gottes, unsers HErrn. Denn wo die heiligen Engel Hüter und Wächter sind, da können die Teufel und böse Menschen nicht schaden, weder mit List, noch mit Gewalt. Denn sie sind dienstbare Geister, die ausgesandt sind zum Dienste um derer willen, die ererben sollen die Seligkeit, Ebr. 1, 14. Und wo Gott das tägliche Brod entzeucht, da folget größere Noth, denn daß auszusprechen ist. Denn da müssen die barmherzigsten Weiber ihre Kinder selbst kochen, daß sie zu essen haben; wie zu Jerusalem geschehen, da diese Stadt verstöret war für der babylonischen Gefängniß, als davon Jeremias in Klagliedern im 4. Kapitel V. 10. sagt. Und in der letzten Zerstörung Jerusalems hat man Tauben-Mist für köstliche Würze gebraucht und theuer gekauft, und die Menschen haben ihren eigenen Mist gegessen, den Hunger zu stillen. Derwegen danken wir Gott heute herzlich, daß er uns durch seine heiligen Engel bisher beschü-
\quellseite{552}
het und für Theurung und Hunger behütet hat; der erbarme sich ferner unser aller gnädiglich, und gebe uns unsere Nothdurft und behüt uns für allem, was uns schädlich und verderblich sein mag, Amen. Dieses erinnert auch das 6. Kapitel.

\section[Vom schwimmenden Eisen. Blindheit der Syrer. Hungersnoth zu]{Vom schwimmenden Eisen. Blindheit der Syrer. Hungersnoth zu Samaria}

\section{Summarium}

Es ist eine sonderliche und schöne Historia, in welcher wir sehen, wie Gott um die, so ihn fürchten und seinem Wort anhangen, wachen, sie durch seine Engel beschützen und vor allem Uebel bewahren wolle, wider alles, was der Teufel und die Welt ihnen begehret zu thun. V. 17., wie der 34. Psalm V. 8. auch sagt: Der Engel des HErrn lagert sich um die her, so ihn fürchten, und hilft ihnen aus. Und der 91. Psalm V. 11.: Er hat seinen Engeln befohlen über dir, daß sie dich behüten auf allen deinen Wegen, daß sie dich auf den Händen tragen und du dei-
\quellseite{553}
nen Fuß nicht an einen Stein stößet etc. Durch solche Lehre und Exempel werden wir vermahnet, daß wir glauben sollen. (Denn mit den Augen werden wirs eben so wenig sehen, als des Propheten Knecht hie,) daß wir mit solchem Schutz und Hut der lieben Engel ohne Unterlaß umgeben und bewahret sind, sonst wären wir Leibes und Lebens, und alles, was wir haben, nicht einen Augenblick für dem Teufel sicher.

a Uebel ja, so gehets uns, wenn wir unserm Gott dienen, ihr heillosen Propheten. Wie viel besser hatten wirs, da wir Baal dieneten, Jer. 44, 17.

\section{Summarium}

Ist eine Historia, wie Gott den abgöttischen König mit Krieg und großer Theurung strafet, daß die Mütter ihre eigenen Kinder erwürgeten, kochten und aßen. Solches Unglück pfleget Gott zu erwecken, wenn man sein Wort verachtet und sich von falscher Lehre und unrechtem Gottesdienst nicht abwendet. Denn eben um solcher Ursache willen hatte Gott diese großen Propheten, Elia und Elisa, erwecket, daß die Abgötterei abgethan und rechter Gottesdienst wieder würde angerichtet. Aber weil sie sich an ihr Predigen und Wunderzeichen gar nichts kehren, strafet sie Gott mit solcher großen Theurung.

\section{Votum}

Gott, unser aller lieber Vater, behüte uns gnädiglich für Krieg, Theurung und andern Strafen; und beschütze durch seine heiligen Engel dieses Land und Stadt, uns und die Unsern, und verleihe, daß wir nach seinem Wort all unser Thun und Lassen allzeit anstellen, Amen.

\chapter{Das 7. Kapitel}

\section{Vorrede}

Gleich wie Gott uns tröstet und erfreuet in unserer Noth; also sollen wir auch unserm Nächsten in seinem Nothstande tröstlich sein und ihn erfreuen. Denn das ist einmal wahr: Je größer Noth, je näher Gott. Und abermal ist recht gesagt: Unser Gott ist über alle Noth. Derwegen der nothleidende Mensch gar wohl bauet, wenn er fest auf Gott trauet. So lehret auch der Herr Christus dieses Matth. 7, 12.: Alles, das ihr wollet, daß euch die Leute thun sollen, das thut ihr ihnen. Und Tobias sagt im 4. Kapitel V. 16. zu seinem lieben Sohn auch also: Was du willst, daß man dir thue, das thue einem andern auch. Welche Worte unsere Richtschnur und Regel sein sollen, darnach wir alle unsere Worte, Werke, Gedanken und Gebärden gegen unsern Nächsten anstellen sollen, zu gefallen Gott und zu Nutz andern Leuten. Dies lehret auch das 7. Kapitel.

\section{Geschwinde wohlfeile Zeit zu Samaria}
\quellseite{554}

\section{Summarium}

Dies ist eine herrliche schöne Historia, wie Gott die Syrer wunderbarlich schrecket und jaget, v. 6., und Samaria von dem jämmerlichen Hunger erlediget und nach der größten Theurung plötzlich ein unerhörter und unglaublicher Ueberfluß aller Dinge erfolget, v. 10.

Diese und dergleichen Historien dienen, daß wir an Gottes Wort fest halten und in der Noth seiner Hülfe uns gewiß trösten sollen, als deß,
\quellseite{555}
der allmächtig ist und tausenderlei Wege zu helfen hat, da wir in solcher Noth nicht einen wissen noch sehen. Auf daß wir nicht thun, wie der gottlose König und seine Ritter, die dem Propheten nicht glauben wollen. Auch ist hie zu merken das Exempel der Aussätzigen, daß sie es dafür halten, es wäre Sünde, wenn sie allein ihren eigenen Nutzen bedenken und den armen bedrängten Leuten nicht bald solche gute Botschaften verkündigen sollten.

\section{Votum}

Gott, der heilige Geist, regiere uns, daß all unser Thun und Lassen unserm Nächsten auch zur Freude und Nutz diene; und tröste und helfe allen, die in Noth und Gefahr stecken, damit sie seine Gnade und Barmherzigkeit zeitlich und ewig rühmen und preisen, Amen.

\chapter{Das 8. Kapitel}

\section{Vorrede}

Gott ist gerecht und auch gnädig, zornig und auch barmherzig. Gerecht ist er gegen die Gottlosen, die in Sünden verstarren und verharren. Gnädig erzeigt er sich gegen die Frommen, die Tugend lieben und üben. Zornig wird Gott um der Sünde willen; barmherzig ist er, wenn man die Sünde erkennt und die ihm abbittet im Namen Jesu Christi, seines lieben Sohns, des unschuldigen Lammes Gottes, das der Welt Sünde getragen und gebüßet hat, und vergibt die Sünde, wendet auch ab oder lindert die wohlverdienten Sünden-Strafen. Derwegen laßt uns in wahrer christlicher beständiger Buße leben, daß uns Gott in seinem Zorn nicht strafe, sondern aus Gnade an Leib und Seele segne und behüte, gebe und erhalte auch, was uns noth ist. Das helfe uns Gott allen, Amen. Darauf höret auch an das 8. Kapitel.

\section[Siebenjährige Theurung. Veränderung im Königreich der Syrer. Joram]{Siebenjährige Theurung. Veränderung im Königreich der Syrer. Joram und Ahasja, Könige in Juda}
\quellseite{556}

\section{Summarium}

Ist ein Exempel, daß die Sunamitin wiederum deß genießet, daß sie so gastfrei gewesen, und dem Propheten Elisa so viel Gutes gethan hatte. Denn der Prophet warnet sie, daß sie sich aus dem Lande machen und der Theurung entfliehen sollte, V. 1. Also will Gott mit reichem Segen bezahlen, was man seinem Worte zu Ehren denen Gutes thut, die es predigen und lieb haben. Darnach ist es ein Exempel, wie Gott den Tyrannen Benhadad sterben läßt, V. 10. Aber weil die Kinder Israel sich nicht bessern, kömmt Hasael, noch ein ärgerer König, an seine Statt, der ihnen mehr zu Leide that, denn Benhadad.

\section{Summarium}

Ist ein Exempel, wie der fromme König Josaphat einen ungerathenen Sohn hatte, mit Namen Joram. Darum plaget denselben Joram Gott erstlich mit Krieg, daß die Edomiter von ihm abfallen, V. 20. Darnach plagt ihn Gott mit schwerer Krankheit, der man nicht rathen konnte, daran er auch starb, da er zwei ganze Jahre mit Qual darnach lag. Davon weiter 2 Chron. 21, 8. und 14. 15. 16. Dergleichen Ahazia, sein Sohn, weil er gottlos ist, regieret auch nur ein Jahr und stirbt, V. 29. Also siehet man, daß Gott weder Glück noch Heil den Regenten gibt, so gottlos sind und Gottes Wort verachten.

\section{Votum}

Der gerechte GOtt steure auch heute allen muthwilligen Verächtern seines Wortes; und vergelte zeitlich und ewig alles, was man Gutes thut seinen Dienern, und behüte uns für Krieg, Krankheit, Theurung und anderm Elend und Herzeleid, Amen.

\chapter{Das 9. Kapitel}

\section{Vorrede}

Der heilige Prophet Jesaias sagt im 8. Kapitel V. 20.: Nach dem Gesetz und Zeugniß. Werden sie das nicht sagen, so werden sie die Morgenröthe nicht haben. Das ist wahr geworden an König Joram und Ahazia. Ferner sagt der Prophet Zacharia im 2. Kapitel V. 8.: Wer euch (vernehmt die treuen Prediger)
\quellseite{557}
antastet, der tastet Gottes Augapfel an. Dies ist erfüllet an der Königin Isebel. Auch ist's gewiß, die den HErrn Christum (die rechte Morgenröthe, den Aufgang aus der Höhe, die Sonne der Gerechtigkeit, und den ewigen Glanz der ewigen Herrlichkeit Gottes, des Vaters) und sein Wort verachten und verwerfen, die sollen hie zeitlich gestraft werden mit Hunger, Krieg, Theurung, Krankheiten und andern Plagen, sie sollen auch hernach in den Abgrund der Hölle verworfen werden, zu ewiger Marter und Pein. Dieses werden bezeugen alle Gottlose in ewiger Verdammniß. Die Geschichte gemeldeter Straßen und Untergangs lautet aus dem 9. Kapitel.

\section[Jehu zum Könige gesalbet, tödtet Joram, Ahazia und]{Jehu zum Könige gesalbet, tödtet Joram, Ahazia und Isebel}

a Rasender) Non quod attoniti aut devoti, ut Rabini delirant, sed quod impii Prophetas vocant furiosos, sicut hodie. D. i. nicht daß sie von Sinnen oder schwärmend gewesen, wie die Rabbinen träumen; sondern weil die Gottlosen die Propheten Wahnsinnige heißen, wie noch heute.

b Hohe Stufen) hic fingendum est, fuisse Sellam Magistratus elevatam in urbe, candidam quasi osseam: huc posuerunt Jehu. Sed quia pompa regia tam cito non poterat tapetis ornari, suas vestes substernebant, in pompa festinantes scilicet. D. i. hier muß man sich vorstellen, es sei ein erhöhter Amts- oder Gerichtstuhl in der Stadt gewesen, weiß wie elfenbein: auf denselben setzten sie den Jehu. Aber weil der königliche Pomp nicht so schnell konnte mit Tapeten oder Teppichen hergerichtet werden, so breiteten sie ihre Kleider unter; weil sie nämlich bei diesem Pompe eilten.
\quellseite{558}

\section{Summarium}

Ist eine Historia, wie der gottlose König in Israel, mit Namen Joram, sammt dem König Ahazia in Juda, aus Befehl Gottes um ihrer Abgötterei willen erwürget werden, V. 24. 27. Auf daß man ja lerne und es gewiß dafür halte, wer sich an Gottes Wort nicht kehren, dasselbige verachten und von der Abgötterei und falscher Lehre nicht ablassen will, daß derselbige gewiß gestraft müsse werden. Derhalben mag sonderlich die weltliche Obrigkeit, wo sie anders Glück und Heil haben will, sich an Gottes Wort begeben und demselbigen mit Herzen anhangen und allem abgöttischen Wesen mit Ernst wehren.

\section{Summarium}

Hie ist eine sonderliche Historia, wie es mit der gottlosen Isebel habe ein Ende genommen; zum Exempel allen denen, so Gottes Wort feind sind, und die, so es predigen, verfolgen; daß es Gott ihnen nicht schenken, sondern sie gewißlich strafen will, und soll sie weder ihre Gewalt noch Reichthum helfen. Isebel war eine Königin und eines Königs Tochter, aber mit ihrem gottlosen Tyrannenfürsten, wider Gott und sein Wort.

33 *
\quellseite{559}
verdienet sie, daß sie zum Fenster heraus geworfen und von den Hunden gefressen wird. Also verachtet und schändet Gott wiederum die, so ihn und sein Wort verachten und dawider toben. Geschichts nicht auf solche Weise, wie hie mit Isebel, so geschicht es doch auf andere, wie wir zu unsern Zeiten viele Exempel erlebt haben.

\section{Votum}

Der gerechte GOtt wolle noch steuren und wehren den Feinden seiner Kirche, und erhalten, die sein Wort recht und treulich lehren, und die seinen heiligen Namen bekennen, loben und preisen; und wolle uns regieren, daß wir uns für allen Sünden, wider ihn, unsere Lehrer und Obrigkeit, allzeit fleißig hüten, Amen.

\chapter{Das 10. Kapitel}

\section{Vorrede}

Höret an mit andächtigen Herzen die Worte unsers Gottes und Herrn, da er im 2. Buch Mose im 20. Kapitel V. 5. also sagt: Ich, der Herr, dein Gott, bin ein eifriger Gott, der da heimsuchet der Väter Missethat an den Kindern, bis ins dritte und vierte Glied, die mich hassen; und thue Barmherzigkeit an viel Tausenden, die mich lieb haben und meine Gebote halten. In diesen Worten dräuet Gott zu strafen alle, die seine Gebote übertreten, beides an ihrer Person und ihren Nachkommen. Darum sollen wir uns fürchten vor seinem Zorn und nicht wider seine Gebote thun; er verheißet aber auch Gnade und alles Gute allen, die seine Gebote halten, beides an ihrer eigenen Person und auch an ihren Nachkommen. Darum sollen wir ihn auch lieben und vertrauen, und gerne thun nach seinen Geboten. Davon unterrichtet auch das 10. Kapitel des 2. Buchs der Könige in den Exempeln Ahab und Jehu.

\section{Ausrottung des Geschlechts Ahabs und der Baals-Pfaffen}
\quellseite{560}

\section{Summarium}

Ist eine Historia, in welcher man die Strafe siehet, welche Gott über den abgöttischen König Ahab hat kommen lassen, daß er nicht allein, sondern auch seine Kinder und Verwandten vertilget müssen werden, V. 7. Solche Exempel sollten uns eine Warnung sein, daß wir uns für falscher Lehre und unrechtem Gottesdienst mit Ernst hüten und der gottseligen Lehre anhangen, unangesehen, daß man etwas drüber wagen muß; denn es ist besser, einen gnädigen Gott und die Welt zum Feinde haben, denn der Leute Gunst und Gottes Zorn und Ungnade tragen.

a Das ist über die Cammer.
\quellseite{561}

\section{Summarium}

Jst ein Exempel einer sonderlichen und fürstlichen Tugend, die Gott wohl gefällt und mit reichem Segen belohnet, daß Jehu alle falsche Propheten, so zum abgöttischen Dienst des Baals geholfen hatten, erwürgen lässet und die Kirche Baal zum heimlichen Gemach machet. Denn weltlicher Obrigkeit gehöret dies sonderlich zu, daß sie mit allem Ernst und ohne alle Gnade der Abgötterei wehre. Denn weil es Gottes Ehre und des Nächsten Seligkeit belanget, soll man darin nicht faul noch lässig sein. Das will Gott hernach mit glückseliger Regierung bezahlen. Wie Jehu ganzer acht und zwanzig Jahr im Regiment bleibt und sein Stamm bis ins vierte Glied. Eines solchen Jehu bedürfte heutiges Tages sehr wohl das Papstthum, sammt allen Stiftern und Klöstern, darin eben so große und greuliche Abgötterei getrieben wird, als dort von den Baalspfaffen.

\section{Votum}

Der eifrige gütige Gott wehre heute auch aller Abgötterei und allen Götzendienern, und segne an Leib und Seele, sammt allen ihren Nachkommen, alle, die ihn lieben und fürchten, und thue solches um Jesu Christi, seines lieben Sohnes, willen, Amen.

\chapter{Das 11. Kapitel}

\section{Vorrede}

Gottlose Obrigkeit ist ein allgemeiner Landschaden; deß ist ein Exempel Athalia. Dagegen gottesfürchtige Obrigkeit ist ein allgemeiner Landfrommen; das lehret das Exempel Jojada. Demwegen danken wir heute aber Gott, daß er durch unsere liebe Obrigkeit sein heilig Wort erhält, und daß durch sie die Unterthanen väterlich regiert und treulich geschützt werden. Der wolle ferner zum Regiment dieser Stadt und Landes geben Glück und Segen, Gesundheit und Leben, beides dem hochlöblichen römischen Adler und dem Rathstuhl der Stadt Obrigkeit, Amen. Ferner höret an das 11. Kapitel.

\section{Tyrannei Athaliae. Das Krönung}

a Des Sabbaths) das waren, die auf des Königes Dienst warten, eine Woche um die andere. Wenn ein Theil abging, so ging das andere an.

b Zeugniß) das war das Buch Mose, das fünfte, das dem König befohlen war.
\quellseite{562}

\section{Summarium}

Dies ist eine merkliche Historia, wie der Priester Jojada Joas, des Königs Ahazia Sohn, zum Könige machet, V. 12., und Athalia, welche das Regiment an sich gezogen und die Kinder des Königs alle erwürget hatte, umbringen heißt, zum Exempel allen Tyrannen, daß sie sich für Tyrannei hüten oder dergleichen Strafe auch warten sollen.

Darnach ist wieder ein Exempel, wie der Priester Jojada die Kirche Baals einreißet, Baals Priester erwürget und den rechten Gottesdienst, so unter dem Regiment Athalia gefallen war, wieder anrichtet, V. 18. Solche Werke gehören sonderlich den Regenten zu, so anders Glück und Heil bei ihrem Regiment sein soll.

\section{Votum}

Der allmächtige Gott steure auch heute allen Tyrannen, und regiere und erhalte unsere treue Obrigkeit zu seiner Ehre und zu gemeinem Nutz, und helfe, daß wir derselben allezeit schuldige Ehre und christlichen Gehorsam erzeigen, Amen.

\chapter{Das 12. Kapitel}

\section{Vorrede}

Vom Amt christlicher Obrigkeit sagt der königliche Prophet David im 24. Psalm V. 7.: Machet die Thore weit und die Thüren in der Welt hoch, daß der König der Ehren einziehe. Mit dem Wort Thor verstehet David die Regimente und Regenten, weil man bei den Juden in Thoren Gericht pfleget zu halten; und vermahnet alle hohen und niedrigen Standes Regenten, daß sie nicht allein darauf bedacht sein sollen, daß sie in ihrer Jurisdiction und Gebiet gute äußerliche Policei, Friede und Ruhe erhalten, sondern daß auch Christus, der Fürst des Lebens, Apost. 3, 15., und der Seligkeit, bekannt und aufgenommen werde; daß sie also für ihre Person die Wahrheit und Klarheit, die Ehre und Lehre Christi, lieb und werth halten, dieselbige bewahren und beherbergen, und daß sie auch der christlichen Kirche, Patronen, Schutzherren und Pfleger sein sollen und die ganze Religion fördern und erhalten. Denn das ist der rechte Ornat, Zierrath und Schmuck aller Regimente, wie Kaiser Theodosius gesagt hat. Ja das ist die Thür, dadurch der liebe Gott mit seinem Segen himmlischer und irdischer Wohlfahrt einzeucht bei Obrigkeit und Unterthanen. So helfe der wahre gütige Gott, daß alle Obrigkeit ihr Amt also löblich und glücklich verrichte, wie auch gethan haben Jojada und Joas, von denen das 12. Capitel.

\section{Regierung Joas, Königs in Juda}
\quellseite{563}

\section{Summarium}

Sie ist eine Historia, wie Jojada sammt dem König ordnen, daß der Tempel im Bau erhalten werde, V. 5. Da siehet man wieder, wie es weltlicher Obrigkeit Amts halben gebühren will, daß sie sich mit Ernst der Religion, und was dazu gehöret, sollen annehmen. Sonderlich aber ist zu merken, daß der Text meldet: Joas that, was recht war und dem HErrn wohl gefiel, so lange ihn der Priester Jojada lehret, V. 2. Denn wenn weltliche Obrigkeit, soll sie anders allenthalben recht fahren, muß sie die an der Hand haben, ihres Raths und Unterrichts oft brauchen, welche Gottes Wort predigen und damit umgehen. Darum stehets nimmermehr wohl, wo man die Prediger und Kirchendiener verachtet, wie doch gemeiniglich jetzund die Obrigkeit allenthalben thut.

\section{Votum}

Regiere GOtt ferner unsere und andere christliche Obrigkeiten, daß sie ihnen die Religion lassen befohlen sein, ihnen selbst zu gutem Gewissen und den Unterthanen zu Heil und Wohlfahrt, und segne und schütze sie an Leib, Seele, Land und Leuten, Amen.
\quellseite{564}

\chapter{Das 13. Kapitel}

\section{Vorrede}

Um der Abgötterei willen sind die Cananiter zu Grunde ausgerottet und die Juden verworfen und in alle Welt zerstreuet worden; also sind die mächtigsten und gewaltigsten Länder auf Erden, Syrien, Egypten, Griechenland und das ganze orientische Kaiserthum, auch um der Abgötterei willen, in der Türken Gewalt kommen. Wehe nun der unbußfertigen Welt der Abgötterei wegen, die heute im Schwange gehet und getrieben wird. Denn Gott ist der Herr, unser Gott, darum soll sich niemand zu den Götzen wenden, denn der Herr hasset alle Abgötterei und alle Abgötter müssen umkommen, 3 Mos. 19., Sir. 15., 5 Mos. 8. Derwegen wer beten kann, der bete alle Tage, christlich, herzlich und fleißig, daß Gott unser sich erbarme, und daß er uns fremder Abgötterei nicht entgelten lasse. Darauf höret auch an das 13. Kapitel.

\section{Regierung Joahas und Joas}
\quellseite{565}

\section{Summarium}

Dies ist eine Historia, wie Gott das Königreich Israel mit heftigen und großen Kriegen plaget, V. 3. Alles um der Abgötterei willen. Dennoch da sie um Hülfe bitten, schaffet ihnen Gott Hülfe durch den König Joahas, um seines Bundes willen mit Abraham, Isaak und Jakob, das ist, aus Gnade und Barmherzigkeit.

Die Historia mit dem Grab Elisa, des Propheten, V. 21., muß man deuten, wie andere seine Wunderwerke, auf daß das Volk sähe und wüßte, daß der Gott, von dem Elisa geprediget, der rechte Gott und seine Lehre die rechte Lehre wäre, und sie desto eher von der Abgötterei ließen. Dagegen die Papisten die erlogenen Wunderwerke ihrer Heiligen allein zur Bekräftigung der Abgötterei wider das Wort gebraucht haben, und wollens noch heutiges Tages mit diesen und andern Wunderwerken schmücken. Aber es reimet sich sehr übel, wie gehöret.

\section{Votum}

Der treue GOtt erleuchte alle, die in Abgötterei stecken, daß sie die abschaffen und ihm im Geist und in der Wahrheit dienen; und behüte uns für Krieg und allem Unglück, tröste auch heute alle fromme Christen in aller Noth, Amen.

\chapter{Das 14. Kapitel}

\section{Vorrede}

Laßt uns heute abermal Gott, dem Allmächtigen für die Schenkung seines Worts und Offenbarung seines Willens danken; und ein jedes kehre auch allen möglichen Fleiß an, daß wir dieser Gnade Gottes nützlich gebrauchen, nach der Vermahnung des heiligen Propheten Esaia am 55. Kapitel V. 6., da er spricht: Suchet den HErrn, weil er zu finden ist, rufet ihn an, weil er nahe ist, der Gottlose lasse von seinem Wege, und der Uebelthäter seinen Gedanken, und bekehre sich zum HErrn, so wird er sich seiner erbarmen, und zu unserm Gott, denn bei ihm ist viel Vergebung. Den rechten Weg zu glauben, und solches zur Seligkeit, den zeiget uns an das Evangelium von Jesu Christo; den Weg recht zu leben, und dies zu unserer Wohlfahrt, den zeigen uns die heiligen zehn Gebote Gottes, dazu unsere christliche Haus-Tafel und eines jeden Beruf und Stand. Diese Lehre hat Gott bestätiget mit Exempeln, beide der Belohnung und Strafe an denen, die darnach oder dawider gelebet haben. Hierher gehört auch das 14. Kapitel.

\section[Amazia und Asaria, Könige in Juda; Jerobeam der]{Amazia und Asaria, Könige in Juda; Jerobeam der andere in Israel}
\quellseite{566}

\section{Summarium}

Ist erstlich eine Historia, wie der Mord, welchen die Unterthanen an ihrem Könige Joas begangen hatten, gerochen wird. Und wird sonderlich an Amazia gelobet, daß er solche Strafe an den Thäter wendet und ihre Kinder, sammt der andern Freundschaft, die daran nicht schuldig waren, solches nicht entgelten lässet. Darnach ist eine andere Historia, wie Amazia einen Krieg wider das Königreich fürnimmt und darüber geschlagen und ein langer Ort der Mauern zu Jerusalem eingerissen und die Stadt geplündert wird, V. 13. Alles darum, daß die Sünde Amazia und Juda gestraft würde, daß sie die Götter der Edomiter gesucht und der König des Propheten Strafe und Warnung verachtet hatte, und half sie nichts, daß Israel auch abgöttisch und voll Sünde und Laster war.

\section{Summarium}

Ist eine feine Historia, wie Amazia endlich erschlagen wird, darum, daß er sich an die Abgötterei gehänget und Gottes Wort, so durch den Propheten ihm gepredigt, verachtet hatte.
\quellseite{567}
Die Historia von Jerobeam, dem Könige Israel, von dem die Schrift zeuget, daß er auch habe gethan, daß dem HErrn übel gefiel, und doch 41 Jahr regierte und viel Glück mit Kriegen wider die Syrer gehabt, V. 26. Ist darum wohl zu merken, daß man lerne, daß Gott so gnädig und gütig sei, daß er sich des Elendes, auch derer, so die Strafe wohl verdienet haben, jammern läßt, ohne Zweifel um weniger frommer Leute willen, so noch vorhanden waren, und sich an der rechten Propheten Lehre gehalten, und ihnen gefolget hatten, derselben genießet der König Jerobeam sammt dem ganzen Haufen.

\section{Votum}

GOtt, der da gnädig und barmherzig ist, erhalte uns auch allzeit einen heiligen Samen und Leute, die ihm dienen in Heiligkeit und Gerechtigkeit, und lindere uns die Strafen, die wir mit so viel und großen Sünden verdienet haben; und helfe, daß die Obrigkeit das Böse strafe und das Gute belohne, Amen.

\chapter{Das 15. Kapitel}

\section{Vorrede}

Gleich wie drei große Krebse und Staubbesen sind, als Krieg, Theurung, Pestilenz; also sind auch drei große Hauptsünden, Abgötterei, Tyrannei und Unzucht; aber Abgötterei ist die größte Sünde und wird auch am heftigsten gestraft. Man soll aber wohl lernen, was Abgötterei sei, nämlich: Abgötterei ist zum ersten, dasjenige zu Gott machen, was nicht Gott ist. Zum andern, die Ehre, so Gott allein gebühret, einem Dinge geben, das nicht Gott ist. Zum dritten, Gott an einen Ort oder Kreatur verbinden, dahin er sich nicht verbunden hat in seinem Wort. Zum vierten, vermeinte Gottesdienste dichten, erwählen und verordnen ohne Gottes Wort und Befehl, und glauben, daß man um derselben willen Gott gefalle, gerecht und selig werde. Diese greuliche und verdammte Sünde der Abgötterei laßt uns ja fleißig fliehen, auf daß wir mit den Abgöttern nicht zeitlich und ewig gestraft werden; wie solcher Strafe Exempel uns aber vorgehalten werden im 15. Kapitel.

\section{Regierung etlicher Könige in Juda und Israel}
\quellseite{568}

\section{Summarium}

Ist eine Historia, in welcher man siehet, wie Gott an beiden Königen, Israel und Juda, die Abgötterei strafet, nicht allein mit Krieg und fremdem Volk, sondern auch mit Aufruhr, daß im Königreich Israel immer einer den andern erwürget, V. 10. 14. 19. 25.

Sonderlich aber ist zu merken das Exempel mit dem Könige Asaria, welcher sonst rechtschaffen und fromm war; aber darum, daß er die Höhen nicht abthat, und das Volk nicht dazu hielt, daß sie allein zu Jerusalem opferten, da der HErr befohlen hatte zu opfern, strafet ihn Gott mit dem Aussatz, V. 5. So wird auch sein Sohn Jotham überfallen von den Syrern, V. 37. Auf daß man lerne, daß sich die welt-


\quellseite{569}
liche Obrigkeit schwerlich damit versündiget, wenn sie der Abgötterei nicht mit Ernst wehret, ob sie gleich selbst nicht Abgötterei treibet.

\section{Votum}

Wir danken GOtt, daß er uns sein Wort reichlich und richtig gegeben und von mancherlei Abgötterei gnädiglich entlediget hat; und bitten ihn herzlich, er wolle uns mit seinem Geist regieren und führen, daß wir die Zeit unserer Heimsuchung erkennen, dankbar sein und für allen Sünden uns hüten, Amen.

\chapter{Das 16. Kapitel}

\section{Vorrede}

Die Sünden und Sündenstrafen der Gottlosen sind uns in der heiligen Bibel zur Warnung vorgeschrieben: die Sünden, daß wir sie hassen und lassen sollen; die Strafen, damit wir uns nicht zeitliches und ewiges Verderben auf den Hals ziehen, zu Leibes- und Seelen-Schaden. Derwegen helfe Gott, der Brunnquell und Geber alles Guten, daß wir sein Wort also anhören, daß die Sünder sich bekehren, daß die Frommen in Gottesfurcht beständig verharren, daß die betrübten Herzen getröstet und daß wir alle Gottes Segen zeitlich und ewig haben und behalten. Und gedenket, als wenn Moses vor einem jeden jetzt stünde und ihn anspräche: So du wirst halten alle diese Gebote, die dir Gott geboten hat, daß du darnach thust, so wird aller Segen über dich kommen; wenn du aber nicht gehorchen wirst der Stimme des HErrn, deines Gottes, daß du haltest und thust alle seine Gebote und Rechte, so werden alle Flüche über dich kommen und dich treffen, 5 Mos. 28, 1. 15. Denn also können auch die Worte Mosis gebraucht werden zur Warnung. Folgendes höret an das 16. Capitel.

\section{Ahas und seine Regierung}
\quellseite{570}

a Decke) darunter sie des Sabbaths saßen oder stunden, wie jetzt Fürsten und Herren unter Teppichen oder Getäfelten sitzen.

\section{Summarium}

Ist eine Historia, wie der König Ahas Abgötterei getrieben und neue Gottesdienste aus eigenem Fürnehmen, ohne Gottes Befehl, habe angerichtet und darum vom Könige zu Syrien sei überzogen worden, V. 5. In solcher Noth suchet er nicht Hülfe bei Gott, sondern bei dem Könige zu Assyrien, V. 7. Durch denselben werden die Syrer und endlich auch das Königreich Israel gestraft, wie sie verdient hatten, V. 9. Auf daß man lerne fromm sein, Gott fürchten und nach seinem Wort sich halten, sonst würde dergleichen Unrath nicht außen bleiben.

Sonderlich aber ist das Exempel mit dem Priester Uria zu merken: der beuget dem Könige und thut in der Kirche, was ihn der König heißet; so doch das Widerspiel sollte sein, daß der König seines Regiments warten und der Priester in der Kirche nicht auf den König, sondern auf Gottes Wort sollte sehen, so ginge es zu beiden Theilen recht. Aber der Teufel siehet solche Ordnung nicht gern. Die Ordnung ist ihm viel dienstlicher, darum hilft er dazu, wie man mit des Papstes und Bischöfe Regimenten auch siehet, daß es durchaus ein umgekehrt Ding ist. Sie sollten Prediger und Lehrer sein, so wollen sie Regenten sein; Könige und Fürsten sollen Regenten sein, so wollen sie die Leute zwingen, sie sollen glauben, zum Sacrament gehen und anderes mehr thun, was ihnen gefällt. Solche Unordnung aber bringet gewissen Unrath mit sich.

\section{Votum}

Der treue GOtt lasse ihm das weltliche und geistliche Regiment in diesen letzten und gefährlichen Zeiten befohlen sein; und helfe, daß Obrigkeit und Prediger ihrem Amt recht und nützlich vorstehen, und daß Unterthanen und Zuhörer christlich und gehorsamlich folgen, Amen.

\chapter{Das 17. Kapitel}

\section{Vorrede}

Der heilige Prophet Esaiaß sagt am 30. Kapitel V. 12.: So spricht der Heilige in Israel: Weil ihr das Wort verwerfft und verlasset euch auf Frevel und Muthwillen und trotzet darauf; so soll euch solche Untugend sein, wie ein Riß an einer hohen Mauren, wenn es beginnet zu rieseln, die plötzlich unversehens einfällt und zerschmettert. Weh nun uns Deutschen; denn in Kirchen unter den Gelehrten ist groß Geschrei, Gift und Neid. Im weltlichen Regiment gehets auch seltsam, fast überall. Die Kinderzucht fällt hinweg. Jedermann beinahe wird sicher, roh und gottlos. Daher uns so viele Himmelszeichen und alle Elemente dräuen, und zeigen an, Gott habe freilich was Großes im Sinn und wolle die Ruthe und die verdienten Strafen gehen lassen. Derwegen, o HErr Gott, o Vater aller Gnaden und Barmherzigkeit, heile du uns, so werden wir heil, hilf du uns, so ist uns geholfen, Jer. 17, 14. Amen. Dies soll uns auch erinnern das 17. Kapitel.

\section[Unter Hosea werden die zehn Stämme nach Assyrien]{Unter Hosea werden die zehn Stämme nach Assyrien geführet}
\quellseite{571}

a Schmückten) sie wollten ihre Sünde vertheidigen als recht und wohlgethan, wie alle Ketzer und Abgöttische thun.

\section{Summarium}

Hie ist der fürnehmsten Historien eine, in welcher man sehet, was für ein Unglück und Jammer daraus komme, wenn man Gottes Wort verachtet, Abgötterei treibet und sich an seine Strafe nicht kehren, noch dadurch sich bessern will: nämlich, daß Gott das ganze Israel, die zehn Stämme, durch ein fremdes Volk gefangen wegführen und ihr Land verheeren läßt, V. 18. Wie wir zu unsern Zeiten leider auch sehen, und täglich, wie zu besorgen, je länger je mehr erfahren werden, daß der Türke mit uns dergleichen wird umgehen, weil des Lästerns und Verfolgens wider die heilige erkannte Wahrheit auf der Papisten Seite kein Ende ist, auf unserer Seite aber die Verachtung und der Undank gegen das liebe Wort so groß wird.

Darum soll man hie sonderlich sehen, wie der Text die Ursachen nach einander herzählet, so solche Strafe erregt haben, nämlich, daß sie fremde Götter angebetet, V. 7., die Warnung der Propheten nicht angenommen, sondern verachtet und ein jeglicher nach seinem Gutdünken ihm eine eigene Weise, Gott zu dienen, habe fürgenommen, V. 15.; wie leider zu unsern Zeiten auch geschieht.
\quellseite{572}



\section{Summarium}

Ist eine Historia, in welcher wir sehen, wie und aus was Ursachen die heidnische Abgötterei in Samaria erstlich gepflanzet und angerichtet sei, durch das fremde Volk, welches der König von Assyrien anstatt der Juden drein hat gesetzt. Solche Abgötterei strafet Gott auch an den Heiden, und schicket Löwen unter sie, welche ihrer viel umbracht(en), V. 24. Ob sie wohl nun das erkennen, daß solches eine Strafe Gottes sei, und sich darum unterrichten lassen vom rechten Gottesdienst, dennoch bleiben sie nicht dabei, sondern fallen wieder auf ihr altes abgöttisches Wesen, V. 28. 29. Denn der Teufel nimmt sich darum mit höchstem Fleiß und Ernst an, daß er die Leute vom Wort und rechten Gottesdienst abwenden möge. Darum bedarf es fleißiges Aufsehens und ernstliches Gebet, daß uns Gott dabei gnädig bis ans Ende erhalten wolle. Es ist aber eine greulichere Strafe auf solche Sünde gefolgt, denn auf die Leute zuvor gewesen war. Denn es ist auf die Abgötterei Blindheit und Verdammniß der Seelen gefolgt.

\section{Votum}

Gott erhalte uns bei seinem heiligen Wort und rechten Gottesdiensten, und lindere die wohlverdienten Strafen; helfe auch, daß wir nach seinem Gesetz wandeln, sein Zeugniß halten, ihn von ganzem Herzen suchen und kein Uebels thun, Amen.

\chapter{Das 18. Kapitel}

\section{Vorrede}

Unser lieber HErr Gott hat eine hebräische Sprache. Er lehret das hinterst zuförderst, erst das Leiden, darnach die Herrlichkeit, hat der heil. Lutherus pflegen zu sagen. Das gibt auch
\quellseite{573}
die Ordnung des Jahrs. Die Marterwoche gehet allzeit vor dem Oster-Tage her, und auf die Kreuzwoche das Pfingstfest erst folget. Wer auch in unsers Herrn Gottes Lustgärten (den er hier auf Erden hat, im Wort und Sacramenten) spazieren will, der denke nur nicht, daß er eitel Tausendschönchen darinnen finden werde, sondern er wird auch bisweilen Sauerampfer und bittere Wermuth antreffen, welche doch gleichwohl auch nicht ungesund, sondern gut ist, Leib Lust und Gift zu vertreiben, die muß man auch kosten und schmecken, sondern werden wir vom Teufel und eigenem Fleisch gar bald vergiftet und angesteckt werden. Davon lehret auch die Epistel an die Ebräer im 12. Kapitel V. 6. mit solchen Worten: Welchen der Herr lieb hat, den züchtiget; er stäupet aber einen jeglichen Sohn, den er aufnimmt. So ihr die Züchtigung erduldet, so erbeut sich euch Gott als Kindern. Das erinnert auch das 18. Kapitel.

\section[Bei der Regierung Hiskia wird Jerusalem von Sanherib]{Bei der Regierung Hiskia wird Jerusalem von Sanherib belagert}

*NehusThan) ein eherner König ist das, der die Schlange, von Gott selber in der Zeit geboten und aufgerichtet, abbricht, darum, daß sie in Mißbrauch der Abgötterei geraten war, und betrachtet darf predigen lassen, es sei NehusThan, das ist, ein ehernes Schlanglein, ein geringes Stück Erz, ein klein Kupferling, was sollte das für ein Gott sein?

\section{Summarium}

Hie gehet an die schöne Historia von dem frommen und gottseligen Könige Hiskia, welchen alle Regenten, als ein sonderlich Fürbild, mit Fleiß ansehen und seinem Exempel folgen sollen, wo sie anders Glück und Heil wollen haben. Dieser fromme König wird gerühmt aufs erste, daß er die ärgerlichen Kirchen abgethan und dem unrechten Gottesdienst mit Ernst gewehret habe. Dazu der ehernen Schlangen nicht verschonet, weil sie in einen solchen Mißbrauch kommen war, daß die Leute ihr opferten, V. 4. Zum andern, daß er Gott vertraut, an seinem Wort fest gehalten und von Gott und seinem Geboten nicht sei abgewichen, V. 5. Darum ist der HErr mit allen Gnaden und reichem Segen bei ihm gewesen, daß er von seinen Feinden wunderbarlich erlöset ist worden. Dagegen die andern, wie wir jetzt in vielen Historien
\quellseite{574}
gehöret, um ihrer Abgötterei willen in alles Unglück kommen sind.

b Rabſake heißt auf deutsch ein Erschenkne.

\section{Summarium}

Ist eine Historia, in welcher man siehet, wie der fromme König Hiskiä in großer Angst und Noth ist gewesen, daß Sanherib in Juda her-

34*
\quellseite{575}
auf zeucht und nimmt alle feste Städte ein und sendet seine Botschaft mit großer Macht für Jerusalem, dadurch er dem König Hiskia allen Trost, beide von Gott und den Menschen, ausreden will; sonderlich aber arbeitet er dahin, daß er ihm das Volk in der Stadt verzagt und abwendig mache.

In solcher Noth werden wir hören, wie der fromme König, sich an Gottes Wort gehalten und sich und seine Leute auf Gottes Hülfe vertröstet habe. Wie denn solches der gewisse und beste Trost ist, dahin wir in aller Anfechtung sehen sollen.

\section{Votum}

Helfe GOtt, daß wir auch in wahrem Glauben an ihm halten, in aller Noth ihn anrufen, und steure der Christenheit Sanherib, den Türken, und allen andern seinen und unsern Feinden, Amen.

\chapter{Das 19. Kapitel}

\section{Vorrede}

Ob die Gläubigen wohl eine Zeitlang einen harten Schweiß unter dem Kreuz und in Anfechtung leiden müssen; so haben sie doch ein Kräutlein, das heißt Heil aller Welt und heilig Wundkraut, welches ist unser lieber HErr JEsus Christus, der uns arme Heiden von allen Wunden des Teufels und Todes, die sie uns durch die Sünde gebissen, geheilet, und das rechte Heil aller Welt worden ist, und der noch hilft in geistlicher und leiblicher Noth. Denn endlich finden sich doch die liebe Angelica und die Himmelschlüssel, das ist, die Gläubigen werden zuletzt aus aller Angst gerissen, alle Thränen sollen ihnen von ihrem Angesicht abgewischt und sollen mit allen Engeln mit himmlischer Freude und ewiger Seligkeit ergötzet werden. Derhalben die Christen sich nicht sollen irren lassen, ob sie oft den Gesang traurig müssen anheben; am Ende wird sichs ausweisen, wo es hinaus will. Die Welt hebt zwar ihren Gesang hoch an, darum nicht Wunder ist, ob sie ihn endlich nicht wird erschreien können. Und wenn ihr letztlich das Halleluja wird gelegt werden, so wird sie mit Schanden müssen stille schweigen. Das unterrichtet auch der heilige Prophet Esaias am 3. Kapitel V. 10. 11. und sagt: Prediget von den Gerechten, daß sie es gut haben, denn sie werden die Frucht ihrer Werke essen. Wehe aber den Gottlosen, denn sie sind boshaftig und es wird ihnen vergolten werden, wie sie es verdienen. Davon höret auch an das 19. Kapitel.

\section[Sanheribs Trot[z] und Macht wird auf Hiskias Gebet]{Sanheribs Trot[z] und Macht wird auf Hiskias Gebet gebrochen und getrennet}
\quellseite{576}

\section{Summarium}

Dieß ist ein sehr schön und tröstlich Exempel, wie der fromme König Hiskiä in seiner größten Noth zu GOtt seine Zuflucht hat, sich vor ihm von Herzen demüthiget und um Hülfe bittet im Hause des HErrn, V. 1. Item, er sendet zu dem Propheten Jesaia, daß er auch wolle zu dem HErrn um Hülfe schreien, wider das Lästern des Königs zu Assyrien, V. 15. Solches zeiget aber an, wie ein herzliches Vertrauen der fromme König auf Gottes Güte gehabt habe. Darum hängt er mit festem Glauben am Trost und an der Zusage, die ihm der Prophet Jesaia thut.

Und dies ist der höchste und beste Gottesdienst, in Nöthen GOtt also anrufen, und sich seiner Barmherzigkeit und seines Wortes trösten, und nicht auf Menschen Gewalt, Geld noch Gunst sehen.
\quellseite{577}

\section{Summarium}

Ist eine Trost-Predigt, in welcher der Prophet Jesaia den frommen König Hiskia ermahnet, daß er auf Gottes Hülfe warten und in der Noth nicht verzagen soll. Sie ist aber sonderlich zu merken, daß Gott das Gebet derer, die auf seine Hülfe warten, hören und zu rechter Zeit helfen will. Wiederum will er den Stolz und Hochmuth der Gottlosen nicht ungestraft lassen abgehen, wie man an Sanherib und seinem Heer siehet. Sonderlich aber ist hie zu lernen, daß man in seinem Unglück nicht verzweifeln soll, sondern hoffen, Gott könne und wolle Hülfe schaffen, wo es uns gleich unmöglich dünkt. Wie denn der König Sanherib hie abgetrieben und in einer Nacht hundert und fünf und achtzig tausend Mann durch den Engel erwürget worden, und der fromme König sammt den Seinen wunderbahrlicher Weise errettet wird, V. 35. Uns zum Trost, daß wir auch lernen also auf Gottes Hülfe sehen und warten.

\section{Votum}

Der barmherzige gütige GOtt vernehme das Schreien aller, die in ihrem Elende ihn anrufen, und erhöre auch unser Gebet in unserer Noth, tröste uns und alle Betrübte durch sein Wort, und wende ab alle Gefahr, um der Vorbitte und Verdienstes unsers HErrn JEsu Christi willen, Amen.

\chapter{Das 20. Kapitel}

\section{Vorrede}

Der heilige Prophet David sagt im 41. Psalm V. 1—4.: Wohl dem, der sich des Dürftigen annimmt, den wird der HErr erretten zur bösen Zeit, der HErr wird ihn bewahren und beim Leben erhalten, und ihm lassen wohlgehen auf Erden, und nicht geben in seiner Feinde Willen. Der HErr wird ihn erquicken auf seinem Siechbette; du hilfst ihm von aller seiner Krankheit. Denn unser Glück und Unglück, Freude und Leid, Gesundheit und Krankheit, Wohlfahrt und Uebelfahrt, Leben und Tod steht nicht in des Teufels, nicht in der Welt, nicht in unserer Gewalt, sondern allein in Gottes Händen, Macht und Gewalt. Und Salomon in Sprüchen im 16. Kapitel V. 5. sagt: Ein stolzes Herz ist dem HErrn ein Greuel, und wird nicht ungestraft bleiben. Denn wie Gott die hoffärtigen Engel aus dem Himmel, und die hoffärtigen Menschen, Adam und Eva, aus dem Paradies verstoßen hat; also will er noch auf Erden die Hoffart strafen. Beide angezogene Sprüche sind erfüllet worden am Könige Hiskia, wie solches bezeuget das 20. Kapitel.

\section{Hiskias Krankheit, Lebens-Verlängerung, Ehrgeiz und Abziehen}
\quellseite{578}

\section{Summarium}

Ist endlich ein Exempel, wie Gott das Gebet der Gottseligen, so auf seine Güte hoffen, erhöret und ihnen hilft in Krankheiten und allerlei Nöthen. Sonderlich ist hie zu merken, daß es nicht Sünde ist, bitten um Gesundheit und langes Leben, und der Arznei pflegen, V. 7. So fern, daß man Gott zuvor durchs Gebet ersuche. Zum andern ist ein Exempel hie, wie es der fromme Hiskia nach seiner großen Krankheit auch verstehet, und sich versündiget mit Hoffart, darum, daß der König zu Babel seine Botschaft und Geschenke zu ihm sendet, daß er wieder gesund war worden, V. 12. Solche Botschaft und Ehre machet den König Hiskia fröhlich, daß er sich überhub und zeigte des Königs von Babel Gesandten alle seine Schätze, V. 13. Darum dräuet ihm Gott durch Jesaiam, daß alles, was er und seine Väter gesammlet haben, soll gen Babel geführet werden, und seine Nachkommen sollen des Königs von Babel Gefangene sein, V. 17. Solche Strafe nimmt Hiskia geduldig an, und erkennet, daß ers verdienet habe, V. 19.

\section{Votum}

Der treue GOtt, unser lieber Vater, erhöre auch unser Gebet, und helfe uns und allen, die seiner Hülfe bedürfen an Leib und Seele, und regiere uns, daß wir in seiner Furcht beständig leben, für Hoffart und andern Sünden uns hüten, Amen.

\chapter{Das 21. Kapitel}

\section{Vorrede}

St. Paulus sagt Röm. 2, 13.: Es ist kein Ansehen der Person vor Gott. Welche ohne Gesetz gesündiget haben, die werden auch ohne Gesetz verloren werden, und welche am Gesetze gesündiget haben, die werden durch das Gesetz verurtheilet werden; sintemal für Gott nicht, die das Gesetz hören, gerecht sind, sondern die das Gesetz thun, werden gerecht sein. Dies wiederholet er an die Colosser im 3. Capitel V. 25., da er sagt: Wer Unrecht thut, wird empfahen, was er Unrecht gethan hat, und gilt kein Ansehen der Person. Hat nun Gott die Sünde gestraft an hohen Potentaten, Königen, Fürsten und Herren, so wird er sie an gemeinen Leuten, Bürgern und Bauern freilich nicht ungestraft lassen. Derwegen wollen wir uns an andern spiegeln, wollen aus fremder Leute Schaden klug, und aus ihren Strafen fromm werden, wollen Buße thun, für Sünden uns hüten und in Gottesfurcht leben, daß wir nicht einen zornigen, sondern gnädigen Gott allzeit haben und behalten. O HErr Jesu Christe, das hilf uns allen, Amen. Darauf höret ferner an das 21. Kapitel.

\section{Historie von Manasse und Amon, den Königen Juda}
\quellseite{579}

\section{Summarium}

Ist ein schrecklich Exempel, wie der fromme König Hiskia einen ungerathenen Sohn hat, der allen Jammer und Abgötterei und Tyrannisiren wider die frommen Propheten anrichtet, V. 2. Aber hie sehet man wiederum, wie gereu-
\quellseite{580}
lich Gott solche Sünde träuet zu strafen, daß er sein Volk, gleich wie zuvor die Heiden, vertilget und in ein fremdes Land wolle lassen gefangen wegführen, V. 11. Am Ende ist eine historia, wie Amon, Manasse Sohn, König worden und von den Seinen erschlagen sei, V. 23. Aber die Aufrührer, so solchen Todschlag begangen hatten, werden darum gestraft, V. 24., nach dem Spruch: Wer das Schwert nimmt, soll durchs Schwert umkommen.

\section{Votum}

Regiere der treue liebe Gott die Jugend, daß sie in ihrer christlichen Eltern Fußstapfen treten, fromm sein, Gott ehren und ihrem Nächsten dienen; und behüte für Aufruhr, Krieg und allem andern Unglück und Sündestrafen, Amen.

\chapter{Das 22. Kapitel}

\section{Vorrede}

Wie der König Josia ihm das Gesetz-Buch Mose hat lesen lassen: also hat es im Judenthum eine sonderliche Ordnung gehabt, daß die Bücher Mose alle Sabbather für der Gemeinde sind gelesen worden, Apost. 13, 14., c. 18, 4. Dergleichen das Exempel Christi zu Nazareth, Luc. 4, 16., zeiget auch an, daß man in den Synagogen der Propheten Bücher gelesen und die Juden davon Unterredung gehalten haben. Und der Apostel St. Paulus vermahnet seine Colosser im 3. Kapitel V. 16.: Daß sie das Wort Christi reichlich unter ihnen sollen wohnen lassen in aller Weisheit, daß ein jeglicher sich selbst dadurch lehren und vermahnen möge: also wird uns auch das Alte und Neue Testament fürgelesen. Denn dieweil der Anfechtungen sehr viel sind, so ohne Unterlaß vom Teufel, der Welt und unserm eigenen Fleisch erreget werden, in Sünden und Unglauben zu stürzen, ists unmöglich, daß man sich dagegen aufhalten, wehren und trösten möge, denn allein durch das Wort Gottes. Derwegen, ach gebraucht ja fleißig das Vorlesen der heiligen Biblia zur Besserung eures Lebens, zur Stärkung eures Glaubens und zu eurem Trost in Noth und Tod, Leben und Sterben. Ferner höret nun an das 22. Kapitel.

\section{Unter Josia wird das Gesetzbuch gefunden}
\quellseite{581}

\section{Summarium}

Ist eine schöne Historia von dem frommen König Josia, wie er sich erstlich der Kirchen annimmt, V. 3., auf daß der rechte Gottesdienst gefördert werde, und wie ohngefähr das Gesetzbuch im Tempel gefunden wird, V. 8., läßt es ihm lesen und richtet alles wieder an, was von Gott geboten, aber unter der gottlosen Könige Manasse und Amons Regiment gefallen war. Daraus sehen wir, alsobald Gottes Wort von den Leuten und von der Kanzel kommt, ist’s unmöglich, daß man sich der Abgötterei und anderer Sünden, so daraus folgen, könne wehren. Sonderlich aber ist hie zu merken, daß Gott die künftige und längst verdiente Strafe gnädig darum aufhält, daß der fromme König Gottes Wort annimmt, Buße thut und sich bessert.

\section{Votum}

Wir danken GOtt für alle christliche Obrigkeit; der wolle ihm ferner auch unsere Regenten lassen befohlen sein, von ihnen und uns die verdienten Strafen abwenden, oder ja lindern, und geben, daß die heilige Biblia auch in dieser Kirche, bis an den jüngsten Tag, fleißig und fruchtbarlich vorgelesen werde, Amen.

\chapter{Das 23. Kapitel}

\section{Vorrede}

Fromme Obrigkeit, das reine und lautere Wort Gottes und christliche Kirch-Ordnung, sind große und heilsame Wohlthaten Gottes, und aller Land und Leute schönste Zierde und nützlichste Kleinod. Für diese Gnaden-Werke und Wohlthaten laßt uns heute abermal Gott dem HErrn herzlich danken und fleißig beten, daß er ihm unsere liebe Obrigkeit ferner lasse befohlen sein, auf daß wir unter ihrer christlichen Regierung Gottes Wort hören, die Gottesdienste im Glauben und mit gottseligem Leben üben mögen. Denn die Regenten sind wie die Laternen: Wie dieselbigen das Licht verwahren, daß es nicht ausgehet: also müssen nach Gott die Regenten fürnehmlich das Licht, den HErrn Christum, tragen, seine Pfleger, Nährer und Säugammen sein, soll die reine Lehre und christliche Kirchordnungen, wider alle Sturmwinde des leidigen Teufels und der argen Welt beschützt und erhalten werden. Deß sollen wir uns auch erinnern bei dem 23. Capitel.

\section[Josia reiniget den Gottesdienst. Regierung der Könige Joahas]{Josia reiniget den Gottesdienst. Regierung der Könige Joahas und Josiam}
\quellseite{582}

a Samarim) das waren sonderliche geistliche Leute, wie jetzt die allernächtigsten und strengsten Mönche sein wollen. Darum haben sie auch einen Namen Samarim, der lautet als von thätiger großer Andacht. Und das Räuchern galt bei ihnen, als ob mans der Mönche Singen und Beten in der Kirche. Denn Räucherwerk bedeutet allenthalben Gebet in der Schrift. Aber wie dieser Gebet ist, so war jener Räucherei, beide Menschenwerk ohne Gottes Geist und Wort.

b Maschith) heißt Verderbung; und man hält, es sei der Oelberg gewesen, da man die Uebelthäter hat abgethan, den wir den Galgenberg oder Rabenstein heißen.

c Grabmal) was aber solche Klammern oder Grabhügel und Grabzeichen sind, wirst du in Jeremia und Hesekiel wohl finden.

\section{Summarium}

Dies ist eine treffliche Historia, wie der fromme König Josia endlich die Predigt wieder anrichtet, daß man das Gesetzbuch öffentlich im Tempel lieset, V. 2., und hernach alle Abgötterei, Altar, Kirchen und Kirchen-Geräthe, so man zur Abgötterei brauchet, vertilget, und die abgöttischen Priester erwürget, V. 4.

Solcher Ernst gefällt Gott wohl; darum wäre es wohl zu wünschen, daß die weltliche Obrigkeit
\quellseite{583}
 dergleichen mit dem gottlosen Papstthum auch handelte, auf daß doch Gottes Zorn abgewendet und nicht der Abgötterei wiederum die Thür geöffnet würde. Weil man Klöster, Altäre, abgöttische Bilder, Bücher und dergleichen noch bleiben läßt, und nicht weg thut, wie sie werth sind. Sintemal sie zu nichts gewisser, denn zum Aergerniß dienen.

\section{Summarium}

Ist eine Historia, wie der König Josia neben dem, daß er alle Greuel und Abgötterei hat ausgefegt, auch das Osterfest wieder angerichtet habe, welches seit der Richter Zeit nie so gehalten worden. Denn weil die Leute eine äußerliche Uebung und Gottesdienst müssen haben, soll man doch in allen Dingen sonderlich auf das Wort sehen und nach demselben solche äußerliche Uebung anstellen. Solchem Exempel folgen die Söhne Josias nicht, derhalben kommt alles Unglück über sie, bis endlich das ganze Land um der Sünden willen verwüstet wird, V. 32.

\section{Votum}

Dem treuen GOtt sei Lob, Ehre, Preis und Dank, der uns von der Abgötterei auch erlöset hat; der regiere und behüte ferner Obrigkeit und Unterthanen, daß sie sich auch in der Lernung göttlichen Worts und in Gebrauch
\quellseite{584}
 christlicher Kirchordnung erzeigen, als Schülern und Dienern Christi, unsers HErrn, gebühret, Amen.

\chapter{Das 24. Kapitel}

\section{Vorrede}

Viele Sünden werden belohnet mit viel Strafen; und auf große Sünden folgen auch große Strafen; und, je mehr man sündiget an einem Ort, je näher die Strafen herzu kommen; und, je frecher und sicherer man ist im Sündenleben, je plötzlicher die Strafe herein bricht. Derwegen, ach heute ist es Zeit, daß wir alle Buße thun, und Gott um Vergebung der Sünden und Linderung der Sünde-Strafen bitten, im Namen unsers HErrn Jesu Christi. Heute ist es Zeit, daß wir von allen öffentlichen Sünden ablassen; heute ist es auch Zeit, daß man anfange, christlicher, ehrbarer, züchtiger und aufrichtiger zu leben, denn leider von vielen bisher geschehen ist. Thun wir dieses alles nicht, so folget wahrhaftig Jammer und Herzeleid mit den größten Strafen. Dieses lehret auch das 24. Kapitel.

\section[Von dreien Königen in Juda: Joakim, Jojachim und]{Von dreien Königen in Juda: Joakim, Jojachim und Zidekia}
\quellseite{585}

\section{Summarium}

Dies ist eine Historia, wie sich die Sünde und das gottlose Leben gehäufet und die Strafe je länger je näher kommen sei. Und meldet der Text fürnehmlich die Sünde Manasse, welcher Jerusalem mit unschuldigem Blut hatte voll gemacht, V. 3. Solche Sünde (spricht Gott) wolle er nicht vergeben. Darum ist gewiß zu besorgen, daß das unschuldige Blut, welches die gottlosen Papisten zu unsern Zeiten allenthalben haben vergossen, und noch, werde das Deutschland auch drücken, daß es endlich müsse zu Boden gehen.

\section{Votum}

Der treue Gott und barmherzige liebe Vater sei uns allen gnädig, lindere die wohlverdienten Strafen, die mit gottlosem Leben und unschuldigem Blutvergießen verdienet sind, und segne und behüte unsere Stadt und das ganze liebe Deutschland, Amen.

\chapter{Das 25. Kapitel}

\section{Vorrede}

Gleiche Sünden werden von dem gerechten Gott belohnet mit gleichen Strafen. Nun gehen jetzt überall nicht weniger und geringere Sünden in vollem Schwange, als im jüdischen Lande, da dies vom König NebucadNezar und hernach von Tito ist verwüstet worden. Wehe bewegen der ganzen Welt, wehe allen Land und Leuten, so man nicht heute Buße thut und in wahrer Buße auch lebet. Ach so laßt uns Gott herzlich bitten, er wolle uns unsere Sünden vergeben und die verdienten Strafen lindern und aufhalten um Jesu Christi, seines Sohnes, Verdienstes und Vorbitte willen; und ein jedes wandele und lebe nun in Seligkeit und Gerechtigkeit, in Kraft des heil. Geistes, auch wie es sein Stand und Beruf erfordert. Denn mit gläubigem Gebet löschet man Gottes Zorn, und mit Besserung des Lebens wendet man ab die Strafen oder man erlanget doch Linderung. Davon höret auch das 25. und letzte Kapitel.

\section[Von der Zerstörung der Stadt Jerusalem, und der]{Von der Zerstörung der Stadt Jerusalem, und der babylonischen Gefangenschaft}

a Jerusalem ist zuvor Jebus, item, Salem genannt gewesen, und von Sem, dem Sohn Noah, dem Könige Melchisedek, erbauet, von welches Tode bis auf diese Verwüstung 1236 Jahr sind; der Tempel, welcher im vierten Jahr Salomonis, nach dem Auszug aus Egypten, 480 Jahr erbauet, ist bis auf diese Zeit der Verwüstung gestanden 461 Jahr. Diese Verwüstung ist geschehen nach der Welt Anfang 3394 Jahr, vor Christi Geburt 566 Jahr. Die Belagerung hat anderthalb Jahr gewähret. Der Tempel ist hernachmals, nämlich nach dieser Verwüstung 123 Jahr, im 6. Jahr Darii, am dritten Tage Marti wieder vollkommen erbauet, Esra 6., vor Christi Geburt 445. In wieder sammt der Stadt Jerusalem gefangen 518 Jahr, und im andern Jahr des Kaisers Vespasiani im April, am österlichen Fest den 14. April belagert, der Tempel erobert und zum andernmal eben im Monat und auf den achten Tag des Augustmonats, wie zuvor vom Nabuchodonosor, also zum letzten von Tito, des Vespasiani Sohn, verbrannt. Die Stadt aber am achten Tage des Septembers eingenommen, geschleift und verbrennet. Dies ist geschehen nach Christi Himmelfahrt 40 Jahr, nach seiner Geburt 74 Jahr, nach dieser letzten Verwüstung 60 Jahr. Unter dem Kaiser Adriano erregten die Juden einen Aufruhr wider den Kaiser, und wollten Jerusalem und den Tempel wieder erbauen, wurden darüber fünf mal hundert und achtzig tausend Juden auf einen Tag erschlagen, Jerusalem wieder geschleift, 50 Städte und Flecken, 50 Schlösser zerstöret und verbrannt. Erobret die Verwüstung nun 1500 Jahr, und wird also ewig bleiben, wie Daniel geweissaget.
\quellseite{586}

\section{Summarium}

Ist eine schreckliche Historia, wie Gott das Königreich Juda um der Sünde willen in der Feinde Hände gegeben, verheeren und sein Volk gefangen habe lassen wegführen, V. 8. Zum Exempel uns allen, daß wir Gottes Wort folgen, in der Furcht Gottes leben, und alle Abgötterei, sammt falscher Lehre mit höchstem Ernst fliehen sollen. Denn weil Gott seines eigenen Volkes nicht verschonet hat, dem er allein, und sonst keinem andern auf Erden zur selbigen Zeit, sein Wort und rechten Gottesdienst vertrauet hatte, so dürfen wir anders nicht gedenken, er werde dergleichen Sünde auch an uns mit Krieg, Aufruhr und dergleichen Plagen, sonderlich aber mit allerlei Lügen und Irrthum strafen, wo wir uns nicht bessern.

\section{Votum}

Der barmherzige GOtt sei uns gnädig, regiere und helfe uns, daß wir uns alle bessern; und behüte uns für Krieg, Aufruhr, Irrthum und für andern zeitlichen, und dem auch ewigen, Verderben und Schaden, Amen.

\quellseite{587}

\part{Das erste Buch der Chronika}

\chapter{Das 1. Kapitel}

\section[Geschlechtsregister der Patriarchen bis auf Jakob]{Geschlechtsregister der Patriarchen bis auf Jakob}
\quellseite{588}

\chapter{Das 2. Capitel}

\section[Söhne Jakobs und Juda]{Söhne Jakobs und Juda}
\quellseite{589}

a Seher, die Propheten heißen Seher oder Schauer in der Schrift, und sehen heißt lehren oder predigen. Also ist Sobal ein Prophet oder Prediger gewesen in dem halben Manuhoth, das ist ein Name eines Ländleins, das vielleicht fein und wohl gelegen ist; darum es Manuhoth, das ist, Ruhe heißt.

\chapter{Das 3. Capitel}

\section[Register der Könige Juda]{Register der Könige Juda}

a Egla diese wird auch Davids Weib genennet, vielleicht, daß sie die erste und einige ist gewesen, da er noch der Schafe gehütet hat.
\quellseite{590}

\chapter{Das 4. Kapitel}

\section[Nachkommen Juda]{Nachkommen Juda}
\quellseite{591}

\chapter{Das 5. Kapitel}

\section[Geschlechter Simeons]{Geschlechter Simeons}

\chapter{Das 6. (sonst 5.) Kapitel}

\section[Von Ruben, Gad, Manasse]{Von Ruben, Gad, Manasse}
\quellseite{592}

\chapter{Das 7. (Sonst 6.) Capitel}

\section[Zahl und Wohnung der Kinder Levi und Aaron]{Zahl und Wohnung der Kinder Levi und Aaron}

a Erbitten siehe, wie der Glaube alles vermag, wie Christus spricht Matth. 21, 22.

b Nämlich durch Salmanasser.
% Scan page 593: omitted; Bible text/narrative only, no retained apparatus.
\quellseite{594}

\chapter{Das 8. (sonst 7.) Kapitel}

\section[Geschlechtsregister der sechs übrigen Stämme Jakobs]{Geschlechtsregister der sechs übrigen Stämme Jakobs}
% Scan page 595: omitted; Bible text/narrative only, no retained apparatus.
\quellseite{596}

\chapter{Das 9. (sonst 8.) Capitel}

\section[Das Haus Benjamin und Saul]{Das Haus Benjamin und Saul}

\chapter{Das 10. (sonst 9.) Capitel}

\section[Einwohner zu Jerusalem und Gibeon]{Einwohner zu Jerusalem und Gibeon}

a Nethinim heißt die gegebenen oder geschenkten, und waren die Leviten, die sich insonderheit zum Gottesdienst begeben hatten, wo Gottes Hütten und die Lade war unter den Priestern.
\quellseite{597}

b Durch ihren Glauben) das ist gesagt: Solche Stifte thaten sie nicht aus menschlicher Andacht und guter Meinung, sondern aus Gottes Befehl und ihrem Glauben. Denn in Gottes Sache soll keine menschliche Stiftung oder Werk gelten.
\quellseite{598}

\chapter{Das 11. (sonst 10.) Kapitel}

\section{Vorrede}

Von Strafen derer, so Gottes Wort verachten und dem ungehorsam sein, sagt Gott im Propheten Hosea am 4. Kapitel V. 6.: Mein Volk ist dahin, darum, daß es nicht lernen will. Denn du verwirfst Gottes Wort, darum will ich dich auch verwerfen. Du vergissest des Gesetzes deines Gottes, darum will ich auch deiner Kinder vergessen. Und in Sprüchen Salomo im 1. Kapitel V. 29. sagt Gott abermals: Darum, daß sie hassten die Lehre, und wollten des Herrn Furcht nicht haben, wollten meines Raths nicht, und lästerten alle meine Strafe; so sollen sie essen von den Früchten ihres Wesens, und ihres Raths satt werden. Dagegen von Belohnung derer, die Gottes Wort hoch achten und dem folgen, sagt Gott Esaia 55. Kapitel V. 2. 3. also: Höret mir doch zu, und esset das Gute, so wird eure Seele in Wollust fett werden. Neiget eure Ohren her, und kommt her zu mir; höret, so wird eure Seele leben: denn ich will mit euch einen ewigen Bund machen. Und abermals in Sprüchen Salomo 3. Kapitel V. 21. 22. spricht er: Laß sie (die Weisheit) nicht von deinen Augen weichen, daß wird deiner Seele Leben sein. Dies beides lehret auch das 11. Kapitel des ersten Buchs der Chronika.

\section[Sauls Untergang im Streit wider die Philister]{Sauls Untergang im Streit wider die Philister}
\quellseite{599}

\section{Summarium}

Dies ist eine Historie, wie es mit dem König Saul habe ein Ende genommen, darum, daß er erstlich an Gottes Wort nicht gehalten, sondern wider Gottes Befehl mit den Amalekitern gehandelt und hernach die Wahrsagerin, und nicht des HErrn Rath gefraget hatte, V. 5. Solche beide Sünden scheinen, als wären es geringe Sünden, aber die Strafe zeiget genugsam an, daß es große, schwere Sünden sind, sintemal Gott dieselbigen an allen seinen Nachkommen gestrafet hat. Darum soll jedermann allein auf Gottes Wort sehen, und demselbigen mit Fleiß folgen, und nichts dawider fürnehmen, sonderlich aber die Obrigkeit. Sonst soll es ihnen gehen, wie es Saul und seinem Geschlecht gegangen.

\section{Votum}

GOtt regiere und führe unsere Obrigkeit und uns alle, daß wir seinem Wort folgen und wider seinen Befehl nichts fürnehmen noch handeln; und lasse ihm uns und die Unsern, auch insonderheit alle christliche Regenten, in Gnaden befohlen sein, Amen.

\chapter{Das 12. (sonst 11.) Kapitel}

\section{Vorrede}

Gleich wie die Ochsen nicht von Ochsen, und die Schafe nicht von Schafen regieret werden, sondern von einer edlen Kreatur, nämlich von Menschen: also ist die menschliche Natur schwächer, denn daß sie sich selber regieren kann, das zu lassen und zu thun, was Gott geboten und verboten hat, zu seiner Ehre und zu Nutz den Nebenmenschen. Derwegen hat Gott Land und Leute zu Regenten fürgestellet, nicht gemeine Leute, sondern göttliche Menschen, nämlich die Heroas und Wunderleute, derer Regierung (denn dazu gibt Gott ihnen Glück) gottliebenden und gehorsamen Unterthanen heilsam und nützlich ist. Daher werden sie genennet Väter des Vaterlandes, Schilde der Erden und Bäume, von denen und unter denen, die Menschen Schatten, Schutz und Hülfe haben. Dieser Lehre Exempel hält uns für das 12. Kapitel.

\section[Salbung Davids zum Könige. Seine Wohnung und streitbaren]{Salbung Davids zum Könige. Seine Wohnung und streitbaren Helden}

a (Dreihundert) 2 Sam. 23, 8. stehet achthundert. Wer einen Haufen von achthundert angreift, und schlägt dreihundert todt, und die fünfhundert in die Flucht, der hat sie alle auf die Flucht geschlagen. Also da hie (unter dreißigen) steht, hat droben 2 Sam. 23, 8. der Text unter dreien. Denn welcher unter den besten dreien der Fürnehmste ist, der ist auch unter den dreißigen allen der Fürnehmste.
\quellseite{600}

b Goß das ist, er opferte es Gott zum Trankopfer.

\section{Summarium}

Ist eine Historia, wie David zum Könige erwählt wird, und ihm Gott die Gnade thut, daß er viele tapfere Helden an seinem Hofe hat, v. 3. 10. Denn wo Gott zum Regiment will Glück geben, da gibt er seine Leute, die zu Sachen geschickt und tüchtig sind. Solche Leute sollen die Regenten lieb und werth haben. Darum reuet
\quellseite{601}
es David, daß aus seinem Ausschicken die drey tapferen Helden Jasabeam, Eleasar und Samma ihr Leben so gering gewagt und in so große Gefahr gesetzt hatten, daß sie durch der Philister Lager gelaufen und David das Wasser aus dem Brunnen zu Bethlehem gebracht hatten. Und ist wahr, muthwilliges Wagen gefällt Gott nicht. Gleich wie Gott sonderlich Glück dazu gibt, wo man in der Noth auf seine Hülfe etwas waget. Aber wo man sich in Gefahr begibt, der man wohl geraten kann, da ist selten Glück bei.

\section{Votum}

Unser lieber GOtt und HErr lasse ihm unsere Obrigkeit auch befohlen sein, gebe und erhalte ihnen treue Räthe und Mitgehülffen; und erhalte bei uns sein heilig Wort und guten Frieden, und helfe, daß Obrigkeit und Unterthanen ihn allhie zeitlich und dort ewig loben und preisen, Amen.

\chapter{Das 13. (sonst 12.) Capitel}

\section{Davids Beistand an streitbaren Männern}

a Das ist, verborgen.
\quellseite{602}

b Kriegsverständig.

\chapter{Das 14. (sonst 13.) Kapitel}

\section{Vorrede}

Wiewohl der Obrigkeit Amt ist, aus Gottes Ordnung herrschen und regieren, Luc. 22, 25., so soll sie doch auch dienen; aber nicht dem Satan, nicht der Sünden, nicht Geld und Gut, nicht irdischer Wollust, sondern Gott und ihren Unterthanen soll sie dienen. Für ihre Person sollen alle die im Stande der Obrigkeit sein, Gott dienen, der sie erschaffen, durch seinen lieben Sohn erlöset, mit dem heil. Geist auf den Tag ihrer Erlösung versiegelt hat, Ephes. 4.; dem sie in der heiligen Taufe Gehorsam gelobet, und sollen ihm dienen, nach ihrer aufgerichteten Bestallung und Tabulatur der heiligen zehn Gebote Gottes. Darnach sollen sie auch ihren Unterthanen dienen in ihrem tragenden Amt, sollen dem großmächtigen Könige der Ehren ihre Thore, Land, Leute und Städte aufthun, Ps. 24, 7., sollen ihn hausen und beherbergen, Kirchen und Schulen treulich bestellen, mit höchstem Fleiß abschaffen alle Aergernisse in Lehre und Leben, nach Friede trachten, die Frommen schützen und die Bösen strafen, Ps. 34, 22. Ein solcher Regent ist gewesen Theodosius, der Kaiser, dem dieß Zeugniß gegeben wird, daß er sich in seiner Todesstunde ernstlicher und heftiger der christlichen Kirchen, als seines eigenen Leibes und Lebens angenommen hat. Ein solcher Regent ist auch gewesen der König David, von dem das 14. Kapitel.

\section[Ab- und Einholung der Lade des Bundes]{Ab- und Einholung der Lade des Bundes}
\quellseite{603}

\section{Summarium}

Ist ein fein fürstlich Exempel, wie der fromme König David sich mit Ernst darum annimmt, daß der Gottesdienst, so unter Sauls Regiment gefallen, wieder recht werde angerichtet, 2. 3. Denn die Obrigkeit soll den rechten Gottesdienst sonderlich fördern, wo sie anders wohl regieren und Gottes Segen haben will.

Daß aber der HErr den Usa tödtet, daß er die Lade anrühret (der doch von der Leviten Stamm nicht war), ist darum geschehen, daß Gott den ganzen Haufen mit solcher Strafe zur Furcht und Zucht hat gewöhnen wollen. Denn mit dem Gottesdienst soll niemand leichtfertig umgehen. Widerrum segnet Gott Obed-Edom, darum, daß er die Lade des HErrn beherberget.

\section{Votum}

Helffe GOtt, daß unsere liebe Obrigkeit die Gottesdienste auch ferner fördere, und verleihe, daß wir uns bei denen als wahre Christen erzeigen und halten; und segne reichlich, die in diesen betrübten Zeiten seiner Kirche und deren Gliedmaßen Herberge geben, tröstlich und beförderlich sein, Amen.

\chapter{Das 15. (sonst 14.) Kapitel}

\section{Vorrede}

Alle Dinge, die zu Gottes Ehren, und unserm Nächsten zu Nutz, getragen sollen, Anfang, Mittel und Ende, stehet in Gottes Händen, Gewalt und Segen. Derwegen wer regieren, predigen, handeln, obern, handthieren und arbeiten will, daß ihm gelinge und daß er Glück und Fortgang, dazu Nutz und Frommen davon habe, der fange es mit Gott an, mache auch mit Gott Feierabend, das ist, er bete, daß Gott den Anfang segne, dem Mittel beistehe und zum Ende befördere; also wird er erfahren, daß Gott ist ein Vater des Lichts, von dem alle gute Gaben von oben herab kommen, Jac. 1, 17. Und daß es umsonst ist, daß man frühe aufstehet, und darnach lange sitzet, und daß man das Brod isset mit Sorgen, daß Gott seinen Freunden solches schlafend gibt, Ps. 127, 3., das ist, daß er seinen Gläubigen und Gehorsamen über Nacht bescheret, und daß Gott mehr hab, denn er vergab; und da Tag kömmt, daß auch Rath kömmt, wie die alten Deutschen gesagt haben. Darauf höret auch an das 15. Kapitel.

\section[Davids Weiber, Kinder und Sieg wider die Philister]{Davids Weiber, Kinder und Sieg wider die Philister}
\quellseite{604}
(Baal-Prazim) Baal heißt ein Hauswirth oder Mann, Perez heißt ein Riß oder Fach, darum muß diese Stätte Baal Prazim heißen, weil die Philister da zertrennt und zerrissen sind.

\section{Summarium}

Dies ist eine Historie, wie David in allem zuvor GOtt rathfragend und sich seinem Wort nach hält, V. 10. Darum gibt ihm GOtt auch so herrlichen Sieg wider seine Feinde; auf daß man lerne, alle Sachen mit Gott anfangen, so wird Glück und Heil dabei sein, sonderlich in großen Regimenten, da menschliche Weisheit, Gewalt und Vermögen viel zu schwach ist, daß man der Sachen helfen könnte.

\section{Votum}

Der treue GOtt regiere unsere Obrigkeit, daß sie ferner Rath bei ihm suche, und helfe, daß wir auch zu unserm Beruf den Segen von ihm bitten, und daß Obrigkeit regiere, und die Unterthanen handeln, arbeiten und ordnen, daß es gut, nützlich und selig sei, in Christo Jesu, unserm HErrn, Amen.

\chapter{Das 16. (sonst 15.) Kapitel}

\section{Vorrede}

Die Regenten sollen ihrer weltlichen Hände also warten, daß sie gleichwohl daneben der Ehre Gottes nicht vergessen, und ihren Eifer gegen den Gottesdienst, den zu schützen und zu erhalten, nicht sollen fallen lassen. Sollen auch anordnen und darüber halten, daß die Kirchendiener in Verrichtung der Gottesdienste fleißig und treu sein. Darnach soll der Gottesdienst ein gemeiner Dienst für alle Menschen sein, und sollen diesen mit aller Reverenz und Ehrerbietung in Heiligkeit handeln und verrichten, Jung und Alt, Reich und Arm, geistliche und weltliche Personen, auf daß niemand durch böse Exempel geärgert, und daß Gott die Nachlässigkeit und ungebührlichen Wandel nicht strafen dürfe. Dies erinnert Davids Exempel und Usas Strafe, davon das 16. Kapitel.

\section[Die Bundeslade wird an ihren Ort gebracht]{Die Bundeslade wird an ihren Ort gebracht}
\quellseite{605}

\section{Summarium}

Hie ist eine Historie, wie David die Lade des HErrn läßt gen Jerusalem bringen, an die Stätte, die er bereitet hatte, B. 3. Und nun fort anordnet, daß die Leviten die Lade tragen sollten, B. 12. Denn weil Gott solches zuvor durch Mose geboten hatte, wollte es sich nicht gebühren, daß man anders denn solchem Befehl nachhandeln sollte. Aus diesem Kapitel scheint es, daß Usa kein Levit gewesen und wider den Befehl Gottes die Lade habe angerühret, welches allein den Leviten Amts halber gebühret und derhalben seine Strafe hat tragen müssen.

\section{Beten}

GOtt regiere unsere Obrigkeit, daß sie ferner GOttes Ehre fleißig fördere, und helfe, daß wir ihm auch dienen, damit es ihm gefalle; und er als ein gnädiger, barmherziger GOtt, die heiligen Engel und alle liebe Christen, uns auch dienen an Leib und Seele, wie es unsere geistliche und leibliche Nothdurft erfordert, Amen.

\chapter{Das 17. (sonst 16.) Capitel}

\section{Vorrede}

Wir alle sollen Gott für seine Wohlthaten dankbar sein. Zur rechten Danksagung aber geböret erstlich, daß wir erkennen, daß wir alle gute Gaben von Gott haben. Zum andern, daß wir dieselbigen nicht geringe, sondern hoch und groß achten. Zum dritten, daß wir Gott darum loben und preisen sollen, und nicht allein mit schlechten Worten, sondern auch mit freudigem Gesange. Wie gethan haben Mose, Debora, Hanna, David, Josia, Josaphat, Zacharias, Maria und andere; und solches lehret St. Paulus an die Epheser im 5. Kapitel B. 18. 19., da er sagt: Werdet voll Geistes, und redet unter einander von Psalmen und Lobgesängen und geistlichen Liedern, singet und spielet dem HErrn in euren Herzen, und saget Dank allezeit für alles Gott und dem Vater, in dem Namen unsers HErrn Jesu Christi. Zum vierten, daß wir darum Gott zu Gefallen leben, und uns nach seinen Geboten halten sollen, wie wir auch gemeiniglich geloben, wenn wir unsere Nothdurft von Gott bitten. Und zum fünften, daß wir die Wohlthaten und Güter zu dem Brauch sollen wenden,
\quellseite{606}
dazu sie uns Gott bescheret und gegeben hat. Dies sollen wir uns erinnern bei dem Dankpsalm des Königs David aus dem 17. Kapitel.

\section[Davids Lobgesang und Bestellung des Gottesdienstes]{Davids Lobgesang und Bestellung des Gottesdienstes}

a Diese drei hebräischen Wörter, Cicar, Espar, Assa, heißen nicht allein die Materie, als Brod, Fleisch, Wein, sondern auch das Maß oder Gewicht. Als so ich auf deutsch spräche: Er gab ein Loth Brod, ein Pfund Fleisch, ein Nößel Wein.
\quellseite{607}

\section{Summarium}

Dies ist ein Lied, in welchem das Volk vermahnet wird, daß sie Gottes Güte erkennen und ihm dafür von Herzen danken sollen, daß er sich ihrer so väterlich hat angenommen, seinen Bund mit ihnen gemacht und den rechten Gottesdienst ihnen geordnet hat. Darum vermahnet er, daß sie sich mit Fleiß zu solchem Gottesdienst halten und da Gott loben, seine Wohlthat rühmen und preisen sollen, B. 8. Denn sonst ist nichts, damit wir Gott können dankbar sein. Er befiehlt aber auch, daß man mit Fleiß wolle bitten, daß er sie also in seinem Wort und rechten Gottesdienst erhalten wolle, B. 34. Denn solches ist fürnehmlich vonnöthen, weil der Teufel dem Wort und rechtem Gottesdienst so feind ist.

\section{Votum}

Wir danken Gott für alle seine Wohlthaten, geistliche und leibliche, und der ganzen Christenheit erzeiget; der verleihe, daß wir derer recht gebrauchen zu seiner Ehre, auf daß er uns ferner aus Gnaden segne und behüte, und erhalte uns bei seinem Wort und rechten Gottesdiensten, Amen.

\chapter{Das 18. (sonst 17.) Kapitel}

\section{Vorrede}

Gottes Wort und gottesfürchtiger Leute Exempel und Fürbilder sollen unsers Lebens Regel, Wegweiser und Zuchtmeister sein. Von Gottes Wort lehret uns solches St. Petrus 2 Epistel 1. Kapitel V. 19. und sagt: Daß man auf das prophetische und also göttliche Wort achten soll, und im ganzen Leben und Wandel darauf sehen, als auf ein Licht, das da scheinet und uns fürleuchtet in den Finsternissen dieser Welt. Dann von Exempel der Heiligen unterrichtet St. Paulus an die Philipper 3. Kapitel V. 17. und spricht: Folget mir, lieben Brüder (verstehet auch Schwestern) und sehet auf die, die also wandeln, wie ihr uns habt zum Fürbilde. Wohl nun uns allen, so wir nach diesem Bericht leben. Denn wer Gottes Wort und christlichen Tugendexempeln folget, der hat die Verheißung und Trost, daß es Gott belohnen und daß er seine Zusagungen an seinen Gehorsamen erfüllen will. Derwegen so wollen wir christliche Tugenden aus dem göttlichen Wort und vom Könige David lernen, von welchem das 18. Kapitel.

\section[Der Tempelbau Davids wird eingestellt; und der Messias]{Der Tempelbau Davids wird eingestellt; und der Messias verheißen}
\quellseite{608}

a) (Ich) daß ich nicht meine Person, sondern meines Blutes Nachkommen, scilicet in futurum et longinquum. Der ein solcher Mensch sein wird, der in der Höhe Gott der Herr ist, Ps. 89, 7. Wer kann Gott gleich sein unter den Kindern Gottes? Er ist auch Gottes Kind, aber wohl über andere Gottes Kinder, als der selbst auch Gott ist.

\section{Summarium}

Suche dasselbe im zweiten Buch Samuelis, am 7. Capitel.

\section[Notum]{\unsicher{Notum}}

Der liebe GOtt gebe uns auch gottseelige Herzen, und helfe, daß wir den Tugenden der Heiligen GOttes folgen, ihm dankbar sein, fleißig beten und andere gute Werke üben, damit wir ihnen auch in das ewige Leben folgen, um JEsu Christi, unsers HErrn willen, Amen.

\chapter{Das 19. (sonst 18.) Capitel}

\section{Vorrede}

Selig und überselig ist das Regiment, dem solche Obrigkeit vorstehet, der Gott beistehet und derer Rathschläge er regieret, fortsetzet und zu gutem Ende führet. Denn solche sind HErr Gottes Leute, unter derer Regierung hat die Kirche ihre Herberge, und die Gerechtigkeit ihren Platz; da blühet der Unterthanen Nahrung, da sind die Stadt- und Landwächter die heiligen Engel, da hat ihr Paradies die liebe Jugend, da haben ihren Zulauf, Rettung und Platz die Armen und Verlassenen. Da ist der Friede eine Krone der Unterthanen. Derwegen laßt uns heute abermal Gott herzlich danken für unsere Obrigkeit, die er bisher gegeben, erhalten und gesegnet hat. Ach der helfe, daß bis an jüngsten Tag dies Land und diese Stadt ferner regieret werde von Obrigkeit, die Gott liebe und ehre, welche die Unterthanen ihr treulich lasse befohlen sein, und die christlich, gerecht und friedlich alle ihr Amt verrichte, Amen. Darauf höret an das 19. Capitel.
\quellseite{609}

\section[David bestellet, nach glücklich geführten Kriegen, Amtleute]{David bestellet, nach glücklich geführten Kriegen, Amtleute}

a Segnen) Glück wünschen.

\section{Summarium}

Suche dasselbe im zweiten Buch Samuelis, am 8. Kapitel.

\section{Datum}

O Gott, aller frommen Regenten Rathgeber und Schußherr, lasse ihm auch allzeit unsere liebe Obrigkeit befohlen sein, und verleihe ihr ferner glückselige Regierung, auch Sieg wider alle ihre Feinde, und den Unterthanen fromme Herzen und willige Hände, zum Gehorsam und Nothdurft, gegen der Obrigkeit, Amen.

\chapter{Das 20. (sonst 19.) Kapitel}

\section{Vorrede}

Auf daß die Sünden nicht ungestraft bleiben, so muß der Mensch selber oft Ursache zum Strafen geben. Derwegen, so wollen wir heute unsere Sünden erkennen und in wahrem Glauben herzlich bitten, Gott wolle nicht mit uns handeln nach unsern Sünden, und nicht vergelten nach unsern Missethat, Ps. 103, 10. Wolle auch nicht nach unserm Verdienst uns in seinem Zorn strafen, sondern wolle nach seiner Güte und großen Barmherzigkeit uns gnädig sein, und uns regieren, daß wir nach dem Fleische nicht leben, sondern durch den heil. Geist des Fleisches Geschäfte tödten, Röm. 8, 13., damit wir ihn mit mehr und größern Sünden ferner zu strafen nicht zwingen und bringen. Denn die Uebertretung und Ungehorsamkeit empfähet doch endlich ihren rechten Lohn, Ebr. 2, 12. Und darauf höret an das 20. Kapitel.

\section[David rächet sich an den höhnischen Ammonitern]{David rächet sich an den höhnischen Ammonitern}
\quellseite{610}

a Wie man ein Ding kehret, hinten und vorn besiehet, das man wil eigentlich erkunden.

\section{Summarium}

Suche dasselbe im zweiten Buch Samuelis, am 10. Kapitel.

\section{Votum}

GOtt sei uns gnädig, vergebe uns alle unsere Sünde, und strafe uns ja nicht, wie wir es auch mit unsern großen Sünden vielfältig verdienet haben; und helfe, daß wir allzeit üben eine gute Ritterschaft, behalten den Glauben und ein gut Gewissen, Amen.

\chapter{Das 21. (sonst 20.) Kapitel}

\section{Vorrede}

St. Paulus sagt 1 Cor. 10, 13.: Gott ist getreu, der euch nicht lässet versuchen über euer Vermögen, sondern macht, daß die Versuchung so ein Ende gewinne, daß ihrs könnet ertragen. Dieser Spruch lehret, daß der Teufel, Tod und alle Welt viel zu schwach dazu sind, einen Christen anzufechten und ein Leid zuzufügen, wenn es Gott nicht nachläßt und verhänget. Daraus folget, weil Gott über alles mächtig ist und ohne seinen Willen nichts geschehen kann; sein Wille aber allzeit gut ist: so kann uns dasjenige, so uns nach seinem Willen begegnet, nicht schädlich sein, es dünke die Vernunft so böse es gleich immer wolle, und thue Fleisch und Blut so weh es immer kann. Ferner lehret St. Paulus: Daß uns Gott nur lasse versuchen, aber nicht verderben; es soll nur ein Versuchen, ein Probieren und Bewähren sein, daß wir also geübet werden im Glauben, Liebe, Hoffnung, Geduld, Gehorsam und in dergleichen mehr göttlichen
\quellseite{611}
Stücken. Auch tröstet St. Paulus, daß Gott der Versuchung ein Ende will machen, darum, daß wir stets sein in seiner väterlichen Fürsorge. Er sorget für uns, hütet und wacht, es stehet alles in seiner Macht, und kann die Versuchung, daß Anliegen und Kreuz so stark, geschwind, mächtig und groß nicht sein, Gott ist stärker, mächtiger und größer, und weiß dem allen ein gut und selig Ende zu machen. Das alles bestätigt der König David mit seinem Exempel, von dem das 21. Kapitel.

\section[Davids Krieg und Sieg wider die Ammoniter und Philister]{Davids Krieg und Sieg wider die Ammoniter und Philister}

\section{Summarium}

Hie ist eine Historia, wie Gott dem frommen David, ob er ihn wohl mit vielen Kriegen und Widerwärtigkeiten hat betrüben lassen, dennoch immer heraus geholfen und ihm gnädigen Sieg verliehen, nach dem Spruch Ps. 34, 20.: Der Gerechte muß viel leiden, aber der Herr hilft ihm aus dem allen.

\section{Votum}

GOtt, unserm HErrn, sey Lob und Preis, der zu aller Zeit den Seinen unter dem Kreuze beigestanden ist; der helfe heute auch allen, die in Noth sind, und stehe uns allen bei in der letzten Noth, welche ist der bittere Tod, und erhalte uns zum ewigen Leben, Amen.

\chapter{Das 22. (sonst 21.) Kapitel}

\section{Vorrede}

Im andern Buch von den Maccabäern im 6. Kapitel B. 31. stehet vom heiligen Eleasar, daß seine Beständigkeit im Bekenntniß göttlicher Lehre aufgeschrieben sei, daß sie nicht allein die Jungen, sondern Jedermann zur Tugend ermahnen solle: Also ist die Geschichte von Davids Sünden und derer Strafe fürgeschrieben, daß Gottesfurcht daraus lernen soll Jedermann. Denn Gottes Wort ist geschrieben uns zur Warnung, auf welche das Ende der Welt kommen ist, sagt St. Paulus 1 Cor. 10, 11. Und 2 Tim. 3, 16. spricht er: Alle Schrift von Gott eingegeben ist nütze zur Lehre, zur Strafe, zur Besserung, zur Züchtigung in der Gerechtigkeit, daß ein Mensch Gottes sei vollkommen, zu allem guten Werk geschickt. Denn wie man sich beide gegen Gott und den Menschen halten solle, was man auch thun solle und lassen, solches alles lehret die Schrift der heiligen Biblia, wie mit Exempeln, also mit Sprüchen, durch Befehl, Neigung und warnungsweise. Derwegen höret dazu an das 22. Kapitel.

\section[Dreitägige Pestilenz durch Zählung des Volks verursacht]{Dreitägige Pestilenz durch Zählung des Volks verursacht}
\quellseite{612}

a Thörlich) David bekund die eitle Ehre, daß er wollte sehen, wie das Königreich durch ihn, als durch sein Thun, so herrlich hätte zugenommen. Da siehet du, daß fiducia operis, eitel Ehre wider Gottes Ehre und eine Abgötterei ist.

b Säcken) das ist, mit geringen, groben Kleidern, daraus man Säcke macht, wie bei uns die Kittel und grobe Leinwand.
\quellseite{613}

\section{Summarium}

Suche dasselbe im zweiten Buch Samuelis, am 24. Kapitel.

\section{Votum}

Gott sei uns allen gnädig, und weil wir auch mit unsern Sünden Strafe verdienet haben, so wolle er uns nie hindern, um Jesu Christi, unsers HErrn willen, und wolle uns regieren, daß wir das Böse lassen und das Gute thun allezeit, fleißig und beständig, Amen.

\chapter{Das 23. (sonst 22.) Kapitel}

\section{Vorrede}

Unser HErr Christus sagt Matth. 10, 41.: Wer einen Propheten aufnimmt in eines Propheten Namen, der wird eines Propheten Lohn empfangen. Und wer dieser Geringsten einen nur mit einem Becher kaltes Wassers tränket, in eines Jüngers Namen; wahrlich, ich sage euch, es wird ihm nicht unbelohnet bleiben. Mit diesen Worten lehret der HErr Christus, daß reichlich solle belohnet werden, was man gibt und mittheilet zur Erhaltung des heiligen Predigtamts. Das sollte nun alle Menschen reizen und bewegen, daß sie auf bloßen Knien Christo sollten zutragen, dadurch das Predigtamt gefördert und erhalten wird, weil der Mund und der Grund der Wahrheit verheißt, ja mit einem hohen Eide beteuert, und seine Wahrheit darüber zu Pfande setzet; und der HErr ist, der nicht verarmen kann, daß ers ihnen bezahlen wolle, und zwar mit Propheten- und gerechten Lohn, das ist, daß sie eben solche reiche Belohnung, doch aus Gnaden, bekommen sollen, als die Propheten und Gerechten selbst haben empfangen. Derwegen ein jeder nach seinem Beruf und Vermögen die Kirchen, deren christliche Lehrer und die Gottesdienste versorgen soll helfen, wie auch lehret das 23. Kapitel.

\section[Nothwendige Vorbereitung zum Tempelbau]{Nothwendige Vorbereitung zum Tempelbau}

a Friede) Salomo heißt friedsam oder friedreich.

b Ein Centner ist bei tausend Floren, doch an einem Ort geringer oder mehr \&c. In Gräcia ist 600 Kronen
\quellseite{614}

\section{Summarium}

Hier ist eine schöne Historia, wie David seinem Sohn Salomo das Regiment übergibt, und ihm sonderlich befiehlt, er solle GOtt und sein Wort für Augen haben und demselbigen mit ganzem Ernst nachleben, so werde er glücklich sein, V. 13. Insonderheit aber vermahnet er ihn, den Tempel zu bauen, an welchem David allen Vorrath verschafft hatte, V. 14.; auf daß Salomo sich den Unkosten nicht ließe zurücktreiben. Also sollten zu unsern Zeiten, beide Obrigkeit und Unterthanen, auch thun, die Kirchen vermachen bestellen und versorgen, daß die Nachkommen einen Vorrath finden. Denn es ist zu besorgen, unsere Nachkommen werden sich des Evangelii nicht sonderlich annehmen, noch dasselbige zu fördern sich etwas kosten lassen; weil wir bereits jetzt so kärglich und eigennützig damit umgehen, die wir doch wissen, was für ein Jammer unter dem Papstthum, Leibes und Seelen halber, gewesen ist.

\section{Votum}

Helfe der treue GOtt, daß wir uns nach dieser Lehre gehorsamlich verhalten, und mit Geldsteuern, Testamenten und Gebet zur Kirchen Erhaltung dienen, auf daß bis an jüngsten Tag allhie eine heilige Gemeine sei, die den HErrn Christum recht erkenne und ihn allezeit lobe und preise, Amen.

\chapter{Das 24. (sonst 23.) Kapitel}

\section[Der Leviten Zahl, Ordnung und Amt]{Der Leviten Zahl, Ordnung und Amt}

a. Dreißig) es scheinet hie die ebräische Bibel verfälschet zu sein, denn sonst allenthalben zwanzig geschrieben stehet, ut paulo infra in 24. paragrapho et in Mose.
\quellseite{615}

\section{572}

b Mutatio onerum Mosis. D. i. Aenderung der Lasten Mosis.

\chapter{Das 25. (sonst 24.) Kapitel}

\section[Obersten im Heiligthum und ihre Diener]{Obersten im Heiligthum und ihre Diener}
\quellseite{616}

\chapter{Das 26. (sonst 25.) Kapitel}

\section[Ordnung der heiligen Sänger]{Ordnung der heiligen Sänger}

*Horn zu erheben) ich achte, daß dieser Heman sei Davids Prophet gewesen im königlichen Geschäften, die das Königreich belanget haben, wie er hat streiten und regieren sollen. Denn Horn bedeutet in Regiment und Königreich.
\quellseite{617}

\chapter{Das 27. (sonst 26.) Kapitel}

\section[Bestellung der Thorhüter, Schatzmeister und Richter]{Bestellung der Thorhüter, Schatzmeister und Richter}

a Helden) Denn die Priester mußten zur Zeit des Krieges die \unsicher{Vordersten} im Heer sein mit den Trompeten.
\quellseite{618}

\chapter{Das 28. (sonst 27.) Kapitel}

\section[Davids Feldobersten, Fürsten, Amtleute und Hofräthe]{Davids Feldobersten, Fürsten, Amtleute und Hofräthe}
\quellseite{619}

\chapter{Das 29. (sonst 28.) Kapitel}

\section{Vorrede}

Von dem Nutzen der Furcht Gottes sagt Sirach im 1. Kapitel V. 10. also: Die Furcht des HErrn ist Ehre und Ruhm, Freude und eine schöne Krone. Die Furcht des HErrn macht das Herz fröhlich, und gibt Freude und Wonne ewiglich. Das ist, innwendig macht die Furcht Gottes das Herz fröhlich, macht und des Sohnes Gottes Segen theilhaftig, bringt und gibt gewissen, starken und festen Trost im Tode, und sonst in allen Aengsten und Nöthen; sie rechtfertiget das Herz aus Gnaden durch die zugerechnete Gerechtigkeit des HErrn Christi, und führet wieder an im Herzen und Leben bei den Menschen einen neuen Gehorsam, oder christlichen Wandel, gibt innerlichen Frieden und ewige Seligkeit. Von außen aber gibt die Furcht Gottes immer Glück und Wohlfahrt in allen Sachen, daß es dem Menschen wohl geht, bringet sie zu Ehren, gibt Reichthum und Vollauf, lindert und mindert manchmal viel Strafen und errettet aus leiblicher Gefährlichkeit. Derwegen wie sich jedermann der Furcht Gottes befleißen soll, also sollen die Eltern sonderlich ihre Kinder zur Furcht Gottes ermahnen und ziehen, und sollen beides ja fleißig und treulich thun. Davon höret auch an das 29. Kapitel.

\section[Davids letzter Reichstag, wegen des Tempelbaues angestellt]{Davids letzter Reichstag, wegen des Tempelbaues angestellt}
\quellseite{620}

\section{Summarium}

Sie ist ein sehr fein Exempel, an welchem alle Eltern, sonderlich aber die Fürsten und Herren, lernen sollten, wie sie ihre Kinder unterweisen und wozu sie sie fürnehmlich ziehen und halten sollen, nämlich zur Furcht Gottes, daß sie GOtt erkennen und ihm von Herzen dienen lernen. Wirst du solches thun (spricht David zu Salomo), so wirst du GOtt finden, V. 9. Das ist, GOtt wird mit Segen und Glück bei dir sein. Wirst du ihn aber verlassen, so wird dich GOtt ewig verwerfen. Aber die Welt folget solchem heiligen Exempel nicht; ja sie thut das Widerspiel, darum ist auch so viel Jammers in der Welt.

a Salomo baut aus Gottes Befehl, nach dem Wort Gottes, seinem Vater David zugesagt, 1 Chron. 18, 12. Den selbstgewählten Gottesdienst und Werk mag er nicht.

37
\quellseite{621}

\section{Psalm}

Helfe GOtt, daß die Eltern ihre Kinder in der Furcht GOttes aufziehen, und regiere die Kinder durch seinen guten Geist, daß sie ihm und den Eltern gern und beständig gehorsam sein, auf daß sie in Gottesfurcht leben, auch zeitlichen und ewigen Segen haben und behalten, Amen.

\chapter{Das 30. (sonst 29.) Capitel}

\section{Vorrede}

St. Paulus sagt von den Philippern, in der Epistel an sie geschrieben im 4. Kapitel, V. 14. 18.: Ihr habt wohl gethan, daß ihr euch meines Trübsals angenommen habt. Ihr aber von Philipp(en) wisset, daß von Anfang des Evangelii keine Gemeine mit mir getheilet hat, nach der Rechnung der Ausgabe und Einnahme, denn ihr allein. Denn gen Thessalonich sandtet ihr zu meiner Nothdurft einmal, darnach aber einmal. Ein süßer Geruch, ein angenehm Opfer Gott gefällig. Derwegen dem Predigtamt wohl zu thun und mitzutheilen vergesset nicht, solche eure Opfer werden Gott auch wohl gefallen. So hat Gott auch ein Register, darein er schreibet all dasjenige, was man thut ihm und seinem Predigtamt; und nach dem ein jeglicher Gutes oder Böses gethan wird haben, darnach will er auch einen jeglichen lohnen, wenn er über alles Rechnung halten wird am jüngsten Tage. Von dieser Versorgung des Predigtamts lehret auch das 30. und letzte Capitel.

\section[Steuer zum Bau des Tempels. Salomos und Zadoks]{Steuer zum Bau des Tempels. Salomos und Zadoks Salbung. Davids Abschied}
\quellseite{622}

\section{Summarium}

Das ist ein fein Exempel, wie David und das ganze Volk so willig und reichlich ihre Steuer geben zum Exempel, V. 6. Denn weil auf Erden kein nöthiger noch theurer Schatz ist, denn Gottes Wort und der rechte Gottesdienst, sollte solchem Exempel nach jedermann willig mit seinem Geld und Gut dazu helfen. Aber wie viel sich die Welt das Evangelium kosten lasse, und darauf wage, erfähret man leider täglich. Weil nun David sich solches freiwilligen Opfers von Herzen freuet und Gott von Herzen dafür danket, als für eine sonderliche Gnade, ist leichtlich abzunehmen, was für ein Gefallen Gott an unsern schändlichen Gaben habe, und wie er endlich den Unfleiß und Nachlässigkeit strafen werde, daß niemand mit Ernst darnach trachte, daß man die Kirchen-Aemter sammt den Schulen wohl bestellen, seine gelehrte Leute auferziehen und also Gott seinen rechten Tempel, die christliche Kirche, recht aufbauen möchte.

\section{Votum}

Regiere GOtt Obrigkeit und Zuhörer, daß sie bei dem heiligen Predigtamt thun, was zu desselben Förderung und Erhaltung fruchtbarlich dienet, und belohne solches an beiden durch Vermehrung seiner Erkenntniß, durch Stärkung ihres Glaubens und durch Segen an Leib und Seele, Land und Leuten, Gottes- und Ehre, Amen.

37*
\quellseite{623}

\part{Das zweite Buch der Chronika}

\chapter{Das 1. Kapitel}

\section{Vorrede}

Wer ihm selber wohlgefällt, der mißfällt Gott. Derwegen soll niemand in seinem Beruf, Amt und Nahrung auf seine Weisheit und Geschicklichkeit trauen und bauen, sondern auf Gott soll man allein sehen und von ihm Glück und Segen, oder guten Anfang, fröhlichen Fortgang und nützlichen Ausgang, alle Tage im Namen des HErrn Christi bitten. Dies lehret St. Jacobus im 1. Kapitel seiner Epistel, V. 5., da er also sagt: So aber jemand unter euch Weisheit mangelt (vernehmt, auf daß er in seinem Beruf recht wandele und handele, und darin Glück und Wohlfahrt bekomme und habe), der bitte von Gott, der da gibt einfältiglich jedermann, und rücket es niemand auf; so wird sie ihm gegeben werden. Er bitte aber im Glauben und zweifle nicht. Davon unterrichtet auch Salomo in seinen Sprüchen, und sagt im 2. Kapitel V. 2. also: Laß dein Ohr auf Weisheit Acht haben, und neige dein Herz mit Fleiß dazu. Denn so du mit Fleiß darnach rufest, und darum betest, alsdann wirst du die Furcht des HErrn vernehmen, und Gottes Erkenntniß finden. Eben dies lehret auch Salomo mit seinem eigenen Exempel, davon das 1. Kapitel des andern Buchs der Chronika.

\section[Salomos Weisheit und Reichthum]{Salomos Weisheit und Reichthum}
\quellseite{624}

\section{Summarium}

Suchet dasselbe im ersten Theil des 3. Capitels des ersten Buchs der Könige.

\section{Votum}

Gott, der Herr, regiere alle Obrigkeit, daß sie auch Weisheit und glückselige Regierung fleißig von ihm bitte; und erhöre unser Gebet, daß wir thun für unsere liebe Obrigkeit, daß sie Rath und That zu ihrer Regierung habe, ihm zu Ehren und uns zu geistlicher und leiblicher, zeitlicher und ewiger Wohlfahrt, Amen.

\chapter{Das 2. Capitel}

\section{Vorrede}

Ein Mensch, der sich in wahrem Glauben zu Gott bekehret hat und der Gottseeligkeit sich befleißet, so viel als möglich, der ist und wird recht genennet ein heiliger Tempel Gottes. Denn Gott wohnet in ihm, 2 Cor. 6, 16. Die steinernen Tempel oder Gotteshäuser aber werden heilig genennet, daß sie sind zugeeignet zu heiligen Werken und Gottesdiensten. Von diesem Unterschied prediget der heilige Stephanus in Geschichten der Apostel im 7. Capitel V. 47. 50., da er sagt: Salomo bauete Gott ein Haus; aber der Allerhöchste wohnet nicht in Tempeln, die mit Händen gemacht sind, wie der Prophet spricht: Der Himmel ist mein Stuhl, und die Erde meiner Füße Schemel; was wollt ihr mir denn für ein Haus bauen, spricht der Herr, oder welches ist die Stätte meiner Ruhe? Hat nicht meine Hand das alles gemacht? So helfe Gott, daß wir alle seine Tempel und Wohnungen sein und bleiben, in denen er regiere und wirke, und daß wir sein Lob preisen mit andächtigen Herzen, mit dankbaren Zungen und mit ihm wohlgefälligen Werken. Dieser unser barmherzige Gott helfe auch, daß in unsern Tempeln und Kirchen-Versammlungen allzeit sein Wort recht gelehret, die Sacramente recht ausgetheilet und die Menschen sich recht mit Gebet, Danksagung, Glauben und Gehorsam erzeigen und verhalten. Amen. Darauf höret ferner an das 2. Capitel.

\section[Nothwendige Vorbereitung zum Bau des Tempels]{Nothwendige Vorbereitung zum Bau des Tempels}

b Ebenholz soll ein Holz in Indien sein. \unsicher{Ich hieß vielleicht}, das man jetzt Sandelholz heißt.
\quellseite{625}

\section{Summarium}

Ist eine Historia, daß Salomo dem Befehl Gottes nachkommet und sich zum Bauen schicket. Es ist aber hie sonderlich zu merken, daß er dem Könige zu Tyro schreibet: Es sei nicht darum zu thun, als müßte Gott solches Hauses haben, und könnte es nicht gerathen; so doch der Himmel ihm zu enge und zu geringe sei, zu einer Wohnung, B. 6. Sondern daß er, Salomo und sein Volk, eines solchen Orts dürfen, da sie zusammen kommen, Gottes Wort hören, opfern, beten und anders thun mögen, das ihnen von Gott zu thun befohlen war. Solches war bisher an keinem beständigen gewissen Ort geschehen. Nun fortan aber bestimmet Gott den Ort, und befiehlet, daß zu Jerusalem geschehen sollte; das war eine große herrliche Wohlthat.

\section{Votum}

Wir danken Gott, der uns auch solche Wohltat erzeiget hat, und noch, was zu Seelen und Leibes Nothdurft wir bedürfen, uns reichlich mittheilet; der helfe, daß wir uns in seinen heiligen Tempeln als lebendige Gottes-Tempel verhalten, allda lernen, danken, beten und andere heilige Werke verrichten, zu seinem Lobe und uns zu Nutz, Amen.

\chapter{Das 3. Kapitel}

\section[Vom Bau des Tempels]{Vom Bau des Tempels}

a Morija) auf diesem Berge opferte Abraham seinen Sohn. 1 Mos. 22, 2.

b Man hält, Arnan sei der Jebusiter König gewesen, und bekehret zum Gott Abraham; davon in Commentaren zu reden.
\quellseite{626}

\chapter{Das 4. Capitel}

\section[Von des Tempels zugehörigen Sachen]{Von des Tempels zugehörigen Sachen}
\quellseite{627}

\chapter{Das 5. Kapitel}

\section{Vorrede}

Gleich wie ein menschlicher Leib ohne Seele kein Leben, Sinn noch Bewegung hat: also was ist ein Tempel, Kirche oder Gotteshaus ohne Gottes Wort und rechten Gottesdienst? Derwegen weil wir beides haben, Gotteshäuser und darin das s(w)ligmachende Evangelium, und heilige Ceremonien; so laßt uns auch diese Wohlthaten zu Gottes Ehre recht gebrauchen. Gottes Wort sollen wir mit Andacht hören, fleißig lernen, und fruchtbarlich gebrauchen, zu gottseligem Wandel und nöthigem Trost, im Leben und Sterben. Daneben sollen wir Gott, denn er hats verdient mit seinem und mitgetheilten Gaben; beten sollen wir mit einander, denn unsere Nothdurft erfordert es, zu des Leibes und Seelen Nutz und Gedeihen; die heiligen Sacramente sollen aus gespendet und gebraucht werden, damit wir der heiligen Dreifaltigkeit werden eingeleibt und auf daß unser Glaube werde gestärket und wir also leben mögen, wie es Söhnen und Töchtern Gottes gebühret und wohl aufstehen will. Darauf höret nun an das 5. Kapitel.

\section[Einweihung des Tempels angefangen]{Einweihung des Tempels angefangen}

\section{Summarium}

Ist eine Historia, wie man den Tempel eingeweihet hat, B. 2., das ist, den Gottesdienst darin angefangen habe. Weil aber solches alles geschehen ist nach dem Wort und ausgedrücktem Befehl Gottes, erzeiget sich Gott selbst da, daß das Haus mit einer Wolke erfüllet ward, und niemand drinnen konnte bleiben für der Herrlichkeit des Herrn etc.

Die Papisten haben auch ihre Kirchweihen, aber den Befehl Gottes haben sie nicht. Darum findet man den Teufel eher daselbst, denn unserm Herrn Gott, welchem niemand kann dienen, er folge denn lauter und bloß seinem Wort und Willen. Derhalben sind solche Papisten-Kirchen eitel Mörder-Gruben, sintemal sie nichts anderes darin thun, denn daß sie die armen Gewissen martern, und unzählig viel Seelen morden. Denn da höret man nichts von der Lehre der Gerechtigkeit vor Gott, vom Glauben, vom Gebet etc. Sondern eitel Menschen-Lehre.
\quellseite{628}
Wo aber die regieren, da müssen die Leute Noth an ihrer Seele und Seligkeit leiden. Wie denn Christus selbst klaget, daß die Phariser und Schriftgelehrten aus seinem Hause, daß ein Bethaus allen Völkern sollte sein, eine Mördergrube gemacht haben.

\section{Psalm}

Danket dem HErrn, unserm GOtt, daß wir auch unsere Tempel und rechte Gottesdienste haben; der helfe, daß wir ihn darin mit heiligen Gedanken, Worten und Werken loben und preisen, und wie wir in der Taufe und bei dem Gebrauch des heiligen Abendmahls zugesagt, für Sünden uns hüten und in guten Werken täglich fleißig uns üben, Amen.

\chapter{Das 6. Kapitel}

\section{Vorrede}

Das Beten soll eines jeden Menschen fürnehmstes und bestes Werk sein täglich. Daher sagen die Alten, daß die Hände empor gen Himmel heben (und solches, wenn man betet) sei die beste Arbeit eines Menschen. Wer aber fruchtbarlich beten will, der soll dies allein bitten, was mit Gottes Wort übereinstimmt, und soll glauben, daß Gott die Sünde vergebe und die Sündenstrafe lindern und abwenden wolle, um Christi willen, und daß er geben wolle, was zu seiner Ehre und zu des Menschen Wohlfahrt dienen soll. Davon sagt der heilige Apostel Johannes in seiner 1. Epistel im 5. Kapitel V. 14. also: Das ist die Freudigkeit, die wir haben zu Gott, daß, so wir etwas bitten nach seinem Willen, so höret er uns. Und so wir wissen, daß er uns höret, was wir bitten, so wissen wir, daß wir die Bitte haben, die wir von ihm gebeten haben. Dies betrachten wir auch bei dem 6. Kapitel.

\section[Einweihung des Tempels fortgesetzt]{Einweihung des Tempels fortgesetzt}

a Erwählet) merke hie, wie alles muß aus Gottes Befehl geschehen, auf daß ja niemand aus eigener Andacht Gottesdienst anrichte. Denn Salomo die beide die Stadt Jerusalem und die Person David rühmet, daß sie beide von Gott erwählet sind.
\quellseite{629}

b Geredet) nun laß du Salomo auch dein Wort wahr werden, das du Gott geredet hast, zu wandeln in seinem Gesetz, sonst würde Mühe werden.
\quellseite{630}

\section[- Das 2. Buch der Chronika 6. 7]{- Das 2. Buch der Chronika 6. 7}

\section{Summarium}

Suche dasselbe im ersten Buch der Könige, am 8. Kapitel.

Ist ein Gebet, daß Gott erstlich seine Verheißung, so er David gethan hat, halten und vollstrecken wolle. Siehe B. 14. 15. rc. bis zum Ende.

\section{Votum}

GOtt helfe und regiere uns, daß wir in diesen gefährlichen und betrübten Läufften fleißig beten, und erhöre unser Gebet um JEsu Christi, seines lieben Sohnes willen; und gebe im Kirchen-, weltlichen und Hausregiment, was unsere Nothdurft ist und was zu seines heiligen Namens Lob, Ehre und Preiß gelangen soll, Amen.

\chapter{Das 7. Kapitel}

\section{Vorrede}

Gott, der HErr, sagt Esaia im 65. Kapitel V. 24.: Es soll geschehen, ehe sie rufen, will ich antworten, wenn sie noch reden, will ich hören. Und der HErr Christus spricht Matth. 21, 22.: Alles was ihr bittet im Gebet, so ihr gläubet, so werdet ihrs empfahen. Und der heilige Geist sagt durch St. Paulum Röm. 10, 12.: Es ist aller zumal Ein HErr, reich über alle, die ihn anrufen. Da hören wir, daß Gott, die heilige Dreifaltigkeit, der Gläubigen Gebet erhören will, und will ihnen alles geben, was sie bitten, da es ihm zu Ehren gelangen, und dem, der da betet, und denen, für die er bittet, gut und selig sein soll. Denn Gott erhöret uns nicht allzeit nach unserm Willen, sondern nach unserer Seligkeit, das ist, darnach er erkennet, daß uns gut und selig sein soll. O gelobet und gebenedeiet sei Gott, daß er also seine Vorsichtigkeit und väterliche Treue an uns beweiset und erzeiget. Dieß erinnert auch das 7. Kapitel.

\section[Einweihung des Tempels vollzogen]{Einweihung des Tempels vollzogen}
\quellseite{631}

\section{Summarium}

Suche dasselbe im zweiten Buch der Könige, am 8. und 9. Kapitel. Ist eine Historia, wie der HErr Salomo erschienen, und ihm seine Bitte gewähret habe, B. 12. etc.

\section{Votum}

Der barmherzige GOtt erhöre auch allzeit unser christlich und gläubig Gebet, und helfe, daß wir seine Gaben, geistliche und leibliche, also gebrauchen, auf daß er uns die erhalte und mehre, wie es ihm löblich und uns gut und seliglich sein soll, Amen.

\chapter{Das 8. Kapitel}

\section[Salomos gebauete Städte, Herrschaft, Opfer und Schiffe]{Salomos gebauete Städte, Herrschaft, Opfer und Schiffe}
\quellseite{632}

\chapter{Das 9. Kapitel}

\section{Vorrede}

In Kraft des Werkes der Erlösung und des reichen Verdienstes unsers Herrn Jesu Christi, sind auch wir Heiden zum Evangelio kommen, daß auch wir sollen mit Abraham, Isaak und Jakob im Himmelreich zu Tisch sitzen, Matth. 8, 11., und der Wohlthaten Messia hie im Glauben genießen, dort aber die in Ewigkeit anschauen, 1 Cor. 13, 12. Wie denn von diesem hohen allgemeinen Beruf der Heiden zu Christo auch die Weissagungen des Alten Testaments lehren und zeugen: als daß in dem Samen Abrahams, Isaaks und Jakobs alle Völker auf Erden sollen gesegnet werden. Daß dem Helden Christo, aus dem Stamm Juda geboren, sollen die Völker anhangen. Daß alle Könige ihn anbeten, und alle Heiden ihm dienen sollten. Daß die Heiden würden in seinem Licht wandeln. Daß sie ihre Söhne und Töchter in den Armen und auf den Achseln herzu tragen und sich zu dem neuen Jerusalem versammeln werden, und dem HErrn dienen einträchtiglich. Für diese reichen, wahren, ewigen, himmlischen Güter, durch Christum auch uns, die wir Heiden gewesen, erlanget und erworben, sollen wir dankbar sein, wie St. Paulus vermahnet Col. 1, 11. und sagt: Danksaget dem Vater, der uns tüchtig gemacht hat zum Erbtheil der Heiligen im Lichte; welcher uns errettet hat von der Obrigkeit der Finsterniß, und hat uns versetzet in das Reich seines lieben Sohns, durch welchen wir haben die Erlösung durch sein Blut. Dies erinnert auch der erste Theil des 9. Kapitels.

\section[Salomo von der Königin aus Arabien besucht, hat]{Salomo von der Königin aus Arabien besucht, hat ein großes Einkommen}
\quellseite{633}

\section{Summarium}

Suche dasselbe im ersten Buch der Könige, am 10. Kapitel.

Ist eine Historia, daß Salomo seiner Weisheit halben allenthalben bekannt, bis zu B. 17.

\section[Notum]{\unsicher{Notum}}

Wir danken dem gnädigen und gütigen GOtt herzlich, daß er auch unser, als der Heiden, GOtt ist und daß wir auch Kinder des Allerhöchsten genennet werden; der helfe, daß wir wahre Gliedmaßen des Reiches seines lieben Sohnes, unsers Erlösers, des HErrn JEsu Christi, sein und bleiben zeitlich und ewiglich, Amen.

\section[Der andere Theil des 9. Capitels]{Der andere Theil des 9. Capitels}

\section{Vorrede}

Vom Amt und christlicher heilsamer Regierung der Obrigkeit lehret Gottes Wort also: Moses sagt im 5. Buch im 17. Kapitel V. 19.: Ein König soll lernen fürchten den HErrn, seinen GOtt, daß er halte seine Rechte, daß er darnach thue. Und Esaias im 49. Kapitel V. 23. spricht: Die Könige sollen deine Pfleger, und ihre Fürsten deine Säugammen sein. Und im 32. Kapitel V. 1. sagt er: Fürsten werden herrschen, daß Recht zu handhaben. Also spricht auch St. Paulus Röm. 13, 4.: Die Obrigkeit ist Gottes Dienerin, eine Rächerin zur Strafe über den, der Böses thut. Von dergleichen Eigenschaften und Amtspflichten der Obrigkeit lehret der andere Theil des 9. Capitels.
\quellseite{634}

\section{Beten}

Gott lasse ihm unsere liebe Obrigkeit befohlen sein, und helfe, daß ihre Regierung ihm, auch ihnen selbst und den Unterthanen zu Ehren, Nutz und Wohlthat gelange, in Kraft, Hülfe und Beistand des heil. Geistes, Amen.

\chapter{Das 10. Kapitel}

\section{Vorrede}

Der rechte, größte und Gott angenehmste Gottesdienst ist, die Erkenntniß des göttlichen Namens, Anrufung und Danksagung, auch Geduld und Gehorsam in Leiden und Trübsal. Wer nun seinem Gott treulich und wohl dienen will, der glaube seinem heiligen Wort und bekenne es für jedermann, sollte es auch gleich mit Gefahr Leibes und Lebens geschehen. Denn das will Gott haben, als geschrieben stehet Matth. 10, 32.: Wer mich bekennen wird vor den Menschen, den will ich bekennen vor meinem himmlischen Vater. Und St. Paulus Röm. 10. sagt V. 10.: So man von Herzen glaubet, so wird man gerecht; und so man mit dem Munde bekennet, so wird man selig. Auch soll ein jeder Christ Gott um seine Wohlthaten anrufen, und für empfangene Wohlthaten danken, als denn David lehret Ps. 50, 14. und spricht: Opfere Gott Dank, und bezahle dem Höchsten deine Gelübde. Und rufe mich an in der Noth, so will ich dich erretten, so sollst du mich preisen. Und V. 23.: Wer Dank opfert, der preiset mich. Ueber das sollen alle, die Gott recht dienen wollen, ihr Kreuz in Geduld tragen; und davon lehret die Epistel an die Ebräer am 10. Kapitel V. 36., da gesagt wird: Geduld ist euch noth, auf daß ihr den Willen Gottes thut, und die Verheißung empfahet. So helfe der treue Gott, daß wir in seinem Dienst allzeit rechtschaffen und treu erfunden werden, Amen. Ferner höret an das 10. Kapitel.

\section[Trennung des Reichs, und Abfall der zehn Stämme,]{Trennung des Reichs, und Abfall der zehn Stämme, von Rehabeam verursachet}
\quellseite{635}

\section{Summarium}

Suche dasselbe im zweiten Buch der Könige, am 12. Kapitel.

Ist eine sehr feine Historia, in welcher man siehet etc. Sammt den folgenden zweien Paragraphen.

\section{Gebet}

Gott, unser Herrscher, helfe und verleihe, daß unser Obrigkeit ferner ihre Unterthanen nach seinem heiligen Wort und Willen, väterlich und treulich regiere, und strafe ja nicht an uns der Welt Abgötterei; auch wolle er uns geben, daß wir ihm allzeit dienen in Gerechtigkeit und Heiligkeit, Amen.

\chapter{Das 11. Kapitel}

\section{Vorrede}

In 1 Johannis am 5. Kapitel V. 21. vermahnet der heil. Geist, und sagt: Kindlein, hütet euch vor den Abgöttern. Diesem ernsten Befehl Gottes sollen wir alle folgen. Denn der HErr hasset alle Abgötterei, sagt Sirach im 15. Kapitel V. 13. Ferner vermahnet der heil. Geist Ebr. 5, 14. und spricht: Lasset uns halten an dem Bekenntniß. Diesem ernsten Befehl sollen wir alle auch folgen. Denn verleugnen wir, so wird uns Gott auch verleugnen, sagt Paulus 2 Tim. 2, 12. Beides lehret auch das 11. Kapitel.

\section[Rehabeams Kriegsbereitschaft, Städte, Befestigung, Weiber und Kinder]{Rehabeams Kriegsbereitschaft, Städte, Befestigung, Weiber und Kinder}
\quellseite{636}

\section{Summarium}

Sie ist eine sonderliche feine Historia. Nachdem das Königreich Salomo durch Aufruhr zerrennet, und der König Jerobeam einen neuen Gottesdienst im Königreich Israel hatte angerichtet, lassen die Priester und Leviten, samt den andern gottesfürchtigen Leuten, so im Königreich Israel waren, Haus und Hof dahinten und ziehen gen Jerusalem, da der rechte Gottesdienst war, V. 14. Denn ein Mensch soll sich eher aller Güter, ja auch Leibes und Lebens erwägen, denn daß er sich entweder zur Abgötterei zwingen, oder vom rechten Gottesdienste sollte abdringen lassen.

\section{Votum}

Der treue Gott helfe, daß wir auch bei erkannter Wahrheit beständig bleiben, und alle Abgötterei fliehen bis an unser Ende; und verleihe, daß wir geduldig alles darüber leiden und wagen, wenn es die Noth und unser Gewissen erfordert, Amen.

\chapter{Das 12. Capitel}

\section{Vorrede}

Der HErr unser Gott ist der Armen Schutz, ein Schutz in der Noth, darum hoffen auf dich, die deinen Namen kennen, denn du verlässest nicht, HErr, die dich suchen, Ps. 9, 10. Der HErr verstößet nicht ewiglich, sondern er betrübet wohl und erbarmet sich wieder nach seiner großen Güte; denn er nicht von Herzen die Menschen plaget und betrübet, Klagl. Jer. 3, 31. Denn Gott erbarmet sich der Elenden und bedeckt uns unter dem Schatten seiner Flügel, daß wir erfahren, daß er der HErr sei, an welchem nicht zu Schanden werden, die auf ihn bauen. Wer bist du denn, daß du dich für Menschen fürchtest, die doch sterben, und für Menschen Kindern, die als Heu verzehret werden, und verlässest des HErrn, der dich gemacht hat, Esaia 51, 12. Und der uns ja behütet, wie einen Augapfel, Ps. 17. 8. Dieser Lehre und Trostes Exempel ist das 12. Kapitel.

\section[Rehabeam von Sisak überzogen]{Rehabeam von Sisak überzogen}
\quellseite{637}

\section{Summarium}

Suche dasselbe im ersten Buch der Könige, am 14. Kapitel, und hänge dies Stücklein mit daran.

Sonderlich aber ist zu merken, daß Gott die wohlverdiente Strafe gnädig abwendet, wenn man sich vor ihm demüthiget, das ist, wenn man die Sünde erkennet und lässets ihm leid sein, und hernach sich tröstet, daß Gott gnädig sei, und um solche Gnade und Hülfe bittet.

\section{Dom}

Der treue gütige GOtt leite und führe uns, daß wir alle auch rechtschaffene Buße thun, und um Linderung der verdienten Strafen täglich in wahrem Glauben und herzlichem beten; und helfe und behüte uns auch gnädiglich in diesem der Christenheit elenden und betrübten Zustande, Amen.

\chapter{Das 13. Kapitel}

\section{Vorrede}

Es verlässet Gott gar nicht, die an sein Wort sich halten und seinen Geboten gehorsam sein, sondern hilft ihnen aus in der Noth, und erhält sie wider alle ihre geistlichen und leiblichen Feinde. Davon lehret die heilige Schrift also: Durch Esaiam, den Propheten, spricht Gott im 55. Kapitel B. 2.: Höret mir doch zu, und esset das Gute; so wird eure Seele in Wollust fett werden. Neiget eure Ohren her, und kommt her zu mir; höret, so wird eure Seele leben; denn ich will mit euch einen ewigen Bund machen. Und Salomo im 3. Kapitel B. 22. seiner Sprüche sagt: Halt mein Gesetz und meinen Rath, daß wird deiner Seele Leben sein. Ferner 2 Mos. 23, 22. sagt Gott abermal: Wirst du meines Engels, den ich vor dir hersende, Stimme hören, und thun alles, was ich dir sagen werde, so will ich deiner Feinde Feind und deiner Widersacher Widersacher sein. Desgleichen lehret er im 81. Psalm B. 14. und spricht: Wollte mein Volk mir gehorsam sein und Israel auf meinen Wegen gehen, so wollte ich ihre Feinde bald dämpfen und meine Hand über ihre Widersacher wenden. Diese Lehre laßt uns ja zu unsrem Nutz gebrauchen in gegenwärtigen gefährlichen Zeiten. Davon erinnert auch das 13. Kapitel.
\quellseite{638}

\section[Abias, Königs in Juda, Krieg und Sieg wider]{Abias, Königs in Juda, Krieg und Sieg wider Jerobeam, den König in Israel}

\section[Kap. 14. v. 1]{Kap. 14. v. 1}

\section{Summarium}

Sie ist eine schöne Historia, wie Gott dem Theil, so den rechten Gottesdienst hat, Sieg gibt wider den andern Theil, der abgöttisch war, so doch allewege auf jener Seite zween Mann wi-

38*
\quellseite{639}
der einen waren, V. 17. Also strafet Gott den Jerobeam, und an dem Königreich Israel die Abgötterei. Wiederum beweiset er sich, daß ers mit denen wolle halten, die bei seinem Wort bleiben. Ob sie gleich auch nicht so rein und fromm sind, als sie wohl sollten sein. Denn auf dem andern Theil waren auch Sünder; aber Gott lässet sie seines Worts genießen.

\section{Votum}

Der barmherzige GOtt strafe nicht an uns die bisher begangene greuliche Abgötterei, und regiere uns, daß wir seinem Wort gehorchen, auf daß er uns aus Gnaden auch segne und behüte, und unsern Feinden steure und wehre, Amen.

\chapter{Das 14. Kapitel}

\section{Vorrede}

Im Buch der Weisheit im 3. Kapitel V. 9. wird gesagt: Die Gott vertrauen, die erfahren, daß er treulich hält; und die treu sind in der Liebe, läßt er ihm nicht nehmen. Denn seine Heiligen (oder Gläubigen) sind in Gnaden und Barmherzigkeit, und er hat ein Aufsehen auf seine Auserwählten. Also sagt Salomo in Sprüchen 21. Kapitel V. 22.: Ein Weiser, das ist auch ein Gläubiger, gewinnt die Stadt der Starken, und stürzet ihre Macht durch ihre Sicherheit. Daher sagt St. Johannis in der 1. Epistel im 5. Kapitel V. 4.: Unser Glaube ist der Sieg, der die Welt überwunden hat. Denn freilich ist der Glaube ein König und Herr über Sünde, Tod, Teufel, Welt und alles Unglück, so da mag genennet werden. Derwegen läßt uns Gott bitten, er wolle uns wahren Glauben geben, und helfen, daß wir in allen Nöthen Leibes und der Seele seiner Hülfe uns beständig trösten, und die auch erlangen um Jesu Christ, unsers Herrn willen, Amen. Davon höret auch an das 14. Kapitel.

\section[Asa schaffet die Abgötterei ab, und feget wider]{Asa schaffet die Abgötterei ab, und feget wider die Mohren}

\section{Summarium}

Dies ist eine sonderliche Historia, in welcher ersichtlich ein trefflich Exempel der weltlichen Obrigkeit fürgetragen wird, daß sie keine Abgötterei leiden, sondern sie mit Gewalt abthun und ausrotten soll, V. 3. 5. Solches Werk will Gott mit Frieden und Segen reichlich bezahlen, da er dagegen die Abgötterei mit Krieg und allerlei Unglück strafet. Zum andern ist ein fein Exempel hie eines
\quellseite{640}
sonderlichen Glaubens, daß der König Assa sich den großen Haufen der Feinde nicht schrecken läßt, sondern glaubt gewiß, es sei Gott eben so leicht, den Sieg zu geben wider einen großen Haufen, als wider einen kleinen, V. 11. Solcher Glaube, Vertrauen und Zuversicht auf Gottes Hülfe schlägt die Feinde in die Flucht, daß sie mit ihrer Menge wider den kleinen Haufen nichts mögen ausrichten.

\section{Votum}

GOtt, der Vater, behüte uns ferner für aller Abgötterei, und für der List der Abgötter; und JEsus Christus, der Anfänger und Vollender des Glaubens, stärke unsern Glauben, und verleihe, daß wir durch den Glauben auch alle unsere Feinde überwinden, Amen.

\chapter{Das 15. Kapitel}

\section{Vorrede}

Das Wort Gottes ist lebendig und kräftig, und ist ein Richter der Gedanken und Sinne des Herzens, und ist keine Kreatur vor ihm unsichtbar, wie in der Epistel an die Ebräer im 4. Kapitel V. 12. gesagt wird. Und St. Petrus spricht 1. Epistel 1. Kapitel: Ihr seid wiederum geboren, nicht aus vergänglichem, sondern aus unvergänglichem Samen, nämlich, aus dem lebendigen Wort Gottes, das da ewiglich bleibet. Das ist aber das Wort, welches unter euch verkündiget ist. Allhie hören wir, daß Gottes Wort die Kraft hat, wo man es mit Ernst betrachtet, höret und handelt, daß es nimmer ohne Frucht abgehet; sondern allzeit neuen Verstand, Lust, Freude und Andacht erwecket, rein Herz und Gedanken macht. Denn es sind nicht faule noch todte, sondern kräftige, lebendige Worte. So helfe der treue gütige Gott, daß wir sein heilig Wort allzeit zu unserm geistlichen und leiblichen, zu unserm zeitlichen und ewigen Nutz, zur Besserung und Trost, anhören, lernen und gebrauchen, Amen. Der angezogenen Sprüche und Lehre von Kraft des Wortes Gottes ist ein Exempel das 15. Kapitel.

\section[Assa fährt in seiner Kirchenreformation fort]{Assa fährt in seiner Kirchenreformation fort}
\quellseite{641}

\section{Summarium}

Sie ist eine schöne Predigt, in welcher der Prophet Asarja erstlich warnet, daß der Gottesdienst von Tag zu Tag abnehmen, und die Leute je länger je ärger werden, V. 3. Derhalben vermahnet er, man soll an Gott und seinem Wort fest halten, so werde Gott dem Königreich Juda gnädig sein, V. 7. Solche Predigt läßt ihm der fromme König also zu Herzen gehen, und thut alle Greuel der Abgötterei aus dem Lande und unter andern zerreißet er das Stift seiner Mutter, welches sie dem Abgott Mipßeth gestiftet hatte, V. 8. Solchem Exempel sollen alle Fürsten und Herren folgen, und nicht bei einem Vetter, dort einem Oheim hoffieren und um Menschengunst willen die Abgötterei dulden, und Gottes Wort entweder verfolgen, oder nicht mit Ernst darob halten.

\section{Votum}

Wir danken Gott herzlich, daß unsere Sünden durch sein heilig Wort und Gesetz an uns auch gestraft werden; der helfe, daß wir uns daraus bessern, und damit wir auch allem Unglück Leibes und der Seelen dieser und der zukünftigen Welt entfliehen, Amen.

\chapter{Das 16. Kapitel}

\section{Vorrede}

Der heilige Apostel Paulus sagt 1 Cor. 10, 12.: Wer sich lässet dünken, er stehe, mag wohl zusehen, daß er nicht falle. O eine nöthige Vermahnung. Denn unser Fleisch ist so verderbt, daß es bald kann zum Fall kommen, wo Gott die Hand abzieht. Auch ist der Teufel so listig und gewaltig, daß er es auf allerlei Wege versucht, damit er uns fälle. Dazu so bringets uns die Welt so nahe mit ihren Aergernissen und Bosheit, daß sie uns zu sündigen, ja zu groben Lastern leichtlich bewegen kann, wo Gott nicht uns beistehet und erhält, und wo Christus nicht für uns bittet, daß unser Glaube und Gottseligkeit nicht sinke und abnehme. Daher ist recht gesagt: Große Leute sind gefallen, darum kleine und gemeine, damit sie nicht auch fallen, sich wohl fürsehen mögen. Derwegen fromm sein und fromm bleiben in unserer Macht nicht stehet, und wir könnens ja so bald versehen und so wohl stolpern, als irgend ein anderer. Dies lehret auch das 16. Kapitel.

\section[Assa dreifacher Sündenfall]{Assa dreifacher Sündenfall}
\quellseite{642}

\section{Summarium}

Ist eine Historia, in welcher man siehet, daß der feine und fromme König Asa, der so einen festen Glauben zu Gott hatte, da er wider die Mohren krieget, aus solchem Glauben in einen Unglauben fällt und Hülfe sucht wider den König Israel, nicht bei Gott, wie zuvor, sondern bei dem König der Syrer, B. 3. Da solches der Prophet Hanani strafet, B. 7., wird er zornig über ihn, und wirft ihn in Thurm, B. 10. Aber Gott plaget ihn darum, daß er an Füßen krank wird und stirbt, B. 12. In solcher Krankheit (sagt der Text, B. 12.) suchet er den HErrn nicht, sondern die Aerzte. Die Aerzte suchen, ist nicht unrecht, denn Gott die Arznei den Menschen zu gut, sammt andern Kreaturen erschaffen hat. Aber Gott will zuvor gesucht sein, das ist, er will, daß man in Krankheit und allen Nöthen ihn anrufen und Hülfe bei ihm suchen soll, und alsdann die Creatur oder Arznei brauchen, mit der Zuversicht, Gott werde seinen Segen dazu geben.

Solche Exempel sollen uns reizen, daß wir nicht sicher sein, sondern Gott ohne Unterlaß bitten, daß er uns in seiner Furcht erhalten und mit seinem heil. Geist regieren und für Sünden gnädiglich behüten wolle. Denn wo Sünden sind, da findet sich auch die Strafe.

\section{Votum}

Helfe der treue GOtt, daß wir dieser Lehre fleißig und beständig folgen, im Glauben verharren, und in Glück und Unglück christlich uns verhalten. Er wolle ihm auch alle kranke Christen lassen befohlen sein, zu Hülfe und Trost, um JEsu Christi, seines geliebten Sohnes willen, Amen.

\chapter{Das 17. Kapitel}

\section{Vorrede}

Im 5. Buch Mose am 12. Kapitel B. 2. befiehlt Gott, und sagt der ordentlichen Obrigkeit: Verstöret alle Orte, da die Heiden ihren Göttern gedienet haben, und die Gößen ihrer Götter thut ab und vertilget ihren Namen aus demselben Ort. David, der königliche Prophet, spricht die Obrigkeit im 2 Ps. 10. 11. also an:

Und im Buch der Weisheit am 6. B. 22. lehret Gott die Obrigkeit auch also, und spricht: Wollt ihr nun, ihr Tyrannen (oder Fürsten) im Volk gern Könige und Fürsten sein, so haltet die Weisheit in Ehren, auf daß ihr ewiglich herrschet. Durch das Wort Weisheit aber soll verstanden werden Gottes Wort und das heilige Predigtamt, mit dessen treuen Dienern und Predigern. Von solchem der Obrigkeit Amt lehret auch das schöne Exempel Josaphat, von dem das 17. Kapitel.

\section[Josaphat ist bei seiner Gottesfurcht glücklich]{Josaphat ist bei seiner Gottesfurcht glücklich}
\quellseite{643}

a Vorrath b) nicht allein des Getreides, sondern auch des Zeugs, Waffen oder Rüstung.

\section{Summarium}

Dies ist eine sehr feine Historia, wie der fromme König Josaphat im Anfang seines Regiments sich vor allem besinnet, daß er die Abgötterei ausrotte, V. 7., und rechtschaffene Prediger an alle Orte seines Königreichs verordnet, V. 9. Solches ist das fürtrefflichste und ein recht fürstlich Werk, da Gott sonderlich mit Gnaden will bei sein.

Aber der Teufel läßt die Obrigkeit nicht gern zu solchen guten Werken kommen; darum treibt er den meisten Theil dahin, daß sie die frommen Prediger verfolgen und in ihren Landen und Gebieten die rechte Lehre nicht dulden, wie viele Exempel für Augen sind. Wie es aber ein Ende werde nehmen, wird die Zeit ausweisen.

\section{Votum}

Der allmächtige GOtt steure dem Zerstörer alles Guten, dem Teufel, daß er niemand das Gute hindere; und segne die Regierung unserer lieben Obrigkeit, und aller christlichen Potentaten und Herrschaft, daß dieselbe zu Abschaffung aller Abgötterei diene, und daß Gottes-Wort in aller Welt rein und lauter gepredigt werde, Amen.

\chapter{Das 18. Kapitel}

\section{Vorrede}

Die Alten sagen ja schön: Wie der zeitliche Hunger eine gewisse Anzeigung ist, daß der Mensch frisch und gesund sei: also sei es auch eine gewisse Vermuthung der Seelen Gesundheit, wenn ein Mensch Gottes Wort (welches der Seelen Speise wider den ewigen Hunger ist) lieb und werth hat, daß ihn darnach hungert und dürstet. Wie aber eines in einem lieblichen Lustgarten sich erlustiget und darauß ein schönes wohlriechendes Blümlein abbricht: also sollen wir uns auch in dem herrlichen Lustgarten des göttlichen Worts (darinnen die Blümlein des Lebens, als Herzenskraut, Angelica, Benedicten-Wurzel, Ehrenpreis, himmlischer Balsam, Herrgotts Beerlein, Teufels Abbiß, Heilwurzel, Augen-Trost, himmlisches Wundkraut und dergleichen zu finden sind) fleißig umsehen, heilsame Blumen da abbrechen und damit unsere zur Sünde geneigte und auch matte Herzen, die der leidige Satan erstänken und ersticken will, mit dem Geruch des Lebens zum Leben, 2 Cor. 2, 16., wieder erquicken, auf daß wir in wahrem Glauben und gottseligem Wandel verharren, und mit der Welt nicht aus einer Sünde in die andere fallen und damit Gott und Menschen beleidigen und beschädigen. Dazu hört auch an das 18. Kapitel.

\section[Josaphats Zug mit Ahab wider die Syrer]{Josaphats Zug mit Ahab wider die Syrer}
% Scan page 644: omitted; Bible text/narrative only, no retained apparatus.
\quellseite{645}

\section{Summarium}

Suche dasselbe im ersten Buch der Könige, am 22. Kapitel.

\section{Votum}

Der gütige, barmherzige GOtt erleuchte und regiere alle christliche Obrigkeit, daß sie sich durch falsche Lehrer nicht verführen lasse; und steure noch heute allen, die sein Wort verachten, behüte alle treue Lehrer seines Evangelii, tröste und erlöse alle Gefangenen, und segne uns alle zeitlich und ewig, Amen.

\chapter{Das 19. Kapitel}

\section{Vorrede}

Wer seiner Sünden halber öffentlich oder im geheimen, von der Kanzel oder im Beichtstuhl, gestraft wird, der soll (wes Standes und Wesens er auch ist) die Strafe annehmen, von Sünden abstehen und rechtschaffene Früchte der Buße thun. Denn ein jeder Mensch soll Gott dienen nach seinem Beruf, und soll kein Werk unterlassen, daß der christlichen Kirchen und dem weltlichen Regiment zum Besten gelangen mag. Solchen Gehorsam belohnet Gott also väterlich und reichlich, Matth. 10, 42., daß auch nicht ein Becher kaltes Wassers, der Gottes wegen gegeben wird, unvergolten bleiben soll. Ja, er wird alle Nothdurft erfüllen nach seinem Reichthum, in der Herrlichkeit, in Christo Jesu, Phil. 4, 19., denen allen, die sonderlich Gottes Wort und die rechten Gottesdienste befördern und erhalten helfen. Dies soll uns erinnern das 19. Kapitel.

Josaphat von Jehu bestraft; bestellet die Kirche und das weltliche Regiment.

\section{Summarium}

Dies ist erstlich eine Historia, wie Gott den frommen Josaphat durch Jehu, den Propheten, ernstlich strafen lässet, daß er mit dem gottlosen König Ahab sich vereiniget, und ihm wider die Syrer geholfen hat, B. 2. Denn Gott will, daß man mit denen, so öffentlich Abgötterei trei-
\quellseite{646}
ben, keine Gemeinschaft haben und ihnen gar nichts helfen soll. Denn Gott hat im Sinn, sie zu strafen ihrer Abgötterei halben, weshalben wer zu ihnen steht, sich selbst in die Gefahr gibt. Zum andern ist wiederum ein Exempel hie, dem alle Regenten folgen sollten, daß Josaphat das weltliche Regiment und die Kirchen mit tüchtigen und frommen Leuten bestellet, auf daß niemand unrecht geschehe, B. 6. Sonderlich aber, daß die Leute im Wort recht unterwiesen werden, B. 8. Solches soll sich weltliche Obrigkeit mit Ernst befleißigen, denn Gott ein sonderlich Gefallen daran hat, und gibt reichlich seinen Segen dazu.

\section{Votum}

GOtt, der heilige Geist, verleihe und wirke, daß alle Abgötter heute auch ihre Sünden erkennen; und helfe, daß unsere Obrigkeit die Kirche und Regiment ihr ferner lasse befohlen sein, auf daß sie mit allen ihren Unterthanen GOtt lobe und preise, hie und dort in alle Ewigkeit, Amen.

\chapter{Das 20. Kapitel}

\section{Vorrede}

Das 20. Kapitel des andern Buchs der Chronika ist das fürnehmste Kapitel dieses ganzen Buchs, und ist ein schön Exempel des Spruchs des königlichen Propheten Davids im 50. Psalm, V. 15., da er in Gottes Person also sagt: Rufe mich an in der Noth, so will ich dich erretten, so sollst du mich preisen. Er lehret allda, wir sollen in keiner Noth verzagen, sondern Gott um Hülfe in wahrem Glauben und beständiger Geduld bitten, es sei die Noth so groß, so schwer, so breit sie möge, er könne und wolle uns, zu seiner Zeit, daraus erretten, und also helfen, daß wir ihn genug zu rühmen und zu preisen haben werden. Und sollen aus diesem Kapitel wohl und mit Fleiß behalten werden die zween schönen Lehr- und Trostsprüche, als der erste Spruch ist, B. 12.: Wir wissen nicht, was wir thun sollen, sondern unsere Augen sehen nach dir, Gott. Der andere Spruch ist, B. 20.: Glaubet an den HErrn, euren Gott, so werdet ihr sicher sein, und glaubet seinen Propheten, so werdet ihr Glück haben.

\section[Josaphats Sieg wider die Ammoniter und Moabiter]{Josaphats Sieg wider die Ammoniter und Moabiter}
\quellseite{647}

\section{Summarium}

Ist eine schöne Historia, wie der fromme König Josaphat in der Noth sich an den rechten Trost hält und zu Gott um Hülfe schreyet, ihn seiner Zusagung und herrlichen Wohlthat erinnert und bittet, er wolle jetzt in dieser Noth auch dermaßen helfen, V. 6. Solches Gebet erhöret Gott, und schaffet wunderbarer Weise Hülfe aus aller Noth, uns zum Exempel, daß wir auch, wie Josaphat hie heißt, an den HErrn unsern Gott glauben sollen, so werden wir Schutz finden und Glück haben.
\quellseite{648}

\section{Summarium}

Ist erstlich ein sonderlich Exempel eines großen Glaubens, daß Josaphat an des Propheten Jehasiel Wort so gar nicht zweifelt, daß er sein Volk nicht anders zur Schlacht rüstet, denn als sollten sie gen Jerusalem im Tempel gehen, nicht in Streit ziehen, heißt sie singen, Gott danken und seine Güte preisen, als wären die Feinde schon geschlagen und der Sieg erobert, B. 27.

Zum andern ist hie ein Exempel, an welchem wir auch sollten Gott lernen trauen, sintemal er so leicht und unversehens, in der höchsten Noth kann Hülfe und Rath schaffen, wie er hie thut mit den Feinden Josaphat, B. 24.

Am letzten ist wieder ein Exempel, weil Josaphat sich mit dem gottlosen König Ahazia vereiniget, hat er kein Glück, B. 35. Denn Gott will weder Glück noch Heil denen widerfahren lassen, so sich an die Gottlosen, so in öffentlicher Abgötterei leben, hängen, und mit ihnen zu thun haben.

\section[Notum]{\unsicher{Notum}}

Der treue GOtt gebe und erhalte uns den Geist des Gebets, erhöre auch unser Seufzen, Rufen und Schreien in unserer Noth; und regiere uns, daß wir seiner Hülfe und Wohlthaten recht und wohl gebrauchen, ihm zu Lob, uns zu Nutz und unsern Nächsten zur Wohlfahrt, Amen.

\chapter{Das 21. Kapitel}

\section{Vorrede}

Wie der Teufel von Anfang der Welt den Menschen mit allem Fleiß nachgestellet hat, auf daß er sie fällete, ihnen zu zeitlichem und ewigem Schaden, und Gott, dem er spinnfeind ist, zu Unehren: also gehet er noch alle Augenblicke herum wie ein brüllender Löwe, und suchet, welchen er verschlinge, 1 Pet. 5, 8. Gleich wie auch Gott, so lange die Welt stehet, die Sünden an allen Unbußfertigen gestraft hat, zu Lob seiner Gerechtigkeit und zur Strafe der Gottlosen: also will er noch heute strafen alle, die da sündigen und nicht wahre Buße thun. Derwegen laßt uns wachen und beten, auf daß uns der Teufel nicht in Versuchung führe, und damit wir Gottes Zorn, Ruthen und Peitschen nicht verdienen, leiden und mit Noth und Weh fühlen dürfen. Dies lehret auch das 21. Kapitel.

\section[Jorams Brudermord, Abgötterei, Strafe und elender Tod]{Jorams Brudermord, Abgötterei, Strafe und elender Tod}
\quellseite{649}

\section{Summarium}

Dies ist eine Historia, da man erstlich siehet, daß der fromme König Josaphat einem bösen gottlosen Sohn, B. 1., welcher nicht allein allerlei Abgötterei anrichtet, sondern auch seine Brüder alle, dazu auch etliche Obersten in Israel erwürget, B. 4. Darum plaget ihn Gott auch erstlich mit Krieg, daß nicht allein die Unterthanen, Edom und Libna von ihm abfallen, B. 8., sondern auch die Araber und die Philister verwüsten und plündern das Land, B. 16. Auch wird er gestraft am Leib mit heftigen Krankheiten, B. 19. Auf daß jedermann sehe, was die zu gewarten haben, so gottlos sind, und mit Abgötterei umgehen, und wir uns desto fleißiger zu Gottes Wort und dem rechten Gottesdienst halten.

\section{Votum}

Der barmherzige GOtt, unser lieber Vater, regiere uns, daß wir uns den Teufel nicht verführen lassen; und strafe uns nicht, wie wir es alle schon wohl verdienet haben, und steure allem Feinden seiner Christenheit, geistlichen und leiblichen, Amen.

\chapter{Das 22. Kapitel}

\section{Vorrede}

Wer das 22. Kapitel des andern Buchs der Chronika ablesen höret, derselbe soll herzlich erschrecken, daß der zornige und gerechte Gott die zwei königlichen Geschlechter, als Salomonis und Ahabs, so schrecklich gestraft und vertilget hat, und soll ihm ein jedes, reich und arm, edel und unedel, Ursache daraus nehmen, daß es Gott fürchte und die Sünde meide und meide; und soll ihm diese Rechnung machen: Hat Gott nicht verschonet der Geschlechter Salomonis und Ahabs, sondern hat sie der Abgötterei und anderer Sünden wegen gar ausgerottet: so wird er auch anderer nicht verschonen, die mit dergleichen Sünden Strafe verdienet haben; und soll ein jeder also leben, lehren, regieren, handeln, erben und arbeiten, daß er nicht auch zum Verderben und Untergange seines Geschlechts diene und helfe mit seinen Sünden, Lastern und Missethaten. Darauf höret an das gemeldete 22. Kapitel.

\section[Ahazias Untergang. Athalias Mordthat]{Ahazias Untergang. Athalias Mordthat}
\quellseite{650}

\section{Summarium}

Suche dasselbe im andern Buch der Könige, am 9. Kapitel.

Ist eine Historia, wie der gottlose König in Israel, mit Namen Joram etc.

\section{Gebet}

Gott, unser lieber Vater, der es allein thun kann, wende gnädiglich von uns ab die verdienten Strafen; und regiere die Obrigkeit und die Unterthanen, daß sie Abgötterei und alle Sünden fliehen und meiden, und nach Gottes Wort regieren und leben, in Kraft des heiligen Geistes, Amen.

\chapter{Das 23. Kapitel}

\section{Vorrede}

In Sprüchen Salomonis im 24. Kapitel B. 21. sagt der heilige Geist: Fürchte den HErrn und den König, und menge dich nicht unter die Aufrührischen. Denn ihr Unfall wird plötzlich entstehen. Und im 37. Psalm B. 39. 40. spricht David: Der HErr hilft den Gerechten, der ist ihre Stärke in der Noth; und der HErr wird ihnen beistehen und wird sie erretten. Er wird sie von den Gottlosen erretten und ihnen helfen, denn sie trauen auf ihn. Und im Buch der Weisheit im 6. B. 7. stehet: Die Gewaltigen werden gewaltiglich gestraft werden. Dieser Sprüche Exempel ist das 23. Kapitel.

\section[Jojada salbet Joas zum Könige. Athalia wird getödtet;]{Jojada salbet Joas zum Könige. Athalia wird getödtet; Baal zerstöret}
\quellseite{651}

a Zeugniß sein ist dem Könige beides die Krone und das Buch gegeben, auf daß er nicht allein mächtig, sondern auch weise sein sollte, oder (wie man es reden mag) Gottes Wort und Recht wissen. So machet man jetzt Könige mit einem Schwert und Buch \&c.

\section{Summarium}

Suche dasselbe im andern Buch der Könige, am 13. Kapitel.

Ist eine merkliche Historia, wie der Priester Jojada \&c.

\section{Votum}

Der starke und gerechte GOTT steure heute auch allen Tyrannen und Abgöttern, um seines heiligen Namens Ehre willen; und erhalte uns sein Wort und Gottesdienst, ihm auch zu Ehren, uns und unsern Nachkommen zu täglicher Buße, zu täglichem Trost, und zur ewigen Seligkeit, um JEsu Christi, seines lieben Sohnes und unsers HErrn willen, Amen.

\chapter{Das 24. Kapitel}

\section{Vorrede}

Der heilige Prophet Jeremias spricht im 10. Kapitel V. 23.: Ich weiß, HErr, daß des Menschen Thun stehet nicht in seiner Gewalt, und stehet in Niemandes Macht, wie er wandele oder seinen Gang richte. Das ist so viel gesagt: Wenn ein Mensch durch seine Sünden den heiligen Geist ausgejaget und verstoßen, so wird aus dem größten Heiligen und frömmsten Christen ein großer Sünder. Und in seinen Sprüchen im 29. Kapitel V. 18. sagt Salomo: Wenn die Weissagung aus ist, so wird das Volk wild und wüst. Das ist so viel gesagt: Ohne Gottes Wort, welches lehret, wie man recht glauben und christlich leben soll, kann der Mensch nichts anders thun, denn Abgötterei und seinen Willen nach fleisches Lust in allerlei Sünden treiben. Derwegen, ach laßt uns züchtig, gerecht und gottselig leben, Tit. 3, 12., daß wir Tempel und Wohnungen des heiligen Geistes bleiben. Laßt auch Gottes Wort allein unsern Zuchtmeister und Wegweiser sein, in Thun und Lassen, nach unserm Beruf und Stande, daß wir Gott recht dienen unser Lebelang, in Heiligkeit und Gerechtigkeit, die ihm gefällig ist, Luc. 1, 75. Dies soll uns auch erinnern das 24. Kapitel.

\section[Das übrige Capitel, Abgötterei, Strafe und Tod]{Das übrige Capitel, Abgötterei, Strafe und Tod}
\quellseite{652}

\section{Summarium}

Suche dasselbe im zweiten Buch der Könige, am 12. Kapitel.

Ist eine feine Historia, wie Jojada sammt dem Könige Joas ordnen etc.

\section[Der andere Theil des 24. Kapitels]{Der andere Theil des 24. Kapitels}

Ist eine feine Historia, in welcher man siehet, wie großen Fürsten und Herren sonderlich viel daran gelegen ist, was sie für Prediger und Lehrer haben. So lange der Hohepriester Jojada lebte, that Joas recht und hielt sich wohl. Da er aber starb, wandte es sich mit Joas; denn er folgte seinem gottlosen Räthen und ließ sich von ihnen verführen, daß er abgöttisch ward und den Propheten Zachariam, des Hohenpriesters Jojada Sohn, der ihn um solches strafet, steinigen ließ, V. 21. Solches, spricht der Prophet, wird der HErr sehen und strafen, V. 22. Denn Joas ward bald hernach mit Krieg überfallen. Und neben dem, daß ihn Gott plaget mit großen Krankheiten, wird er endlich noch erstochen, V. 25. Zur Strafe, daß er den HErrn verlassen und den Propheten Zachariam unschuldig hatte erwürgen lassen. Zum Exempel allen Regenten, daß sie dergleichen nicht auch thun, oder dergleichen Strafe gewiß gewarten.
\quellseite{653}

\section{Gebet}

Gott, der heilige Geist, regiere und führe uns, daß wir in Gottes Furcht beständig bleiben, für Abgötterei uns hüten und seinem Wort allezeit gehorsam sein; und behüte uns gnädig für zeitlichem und ewigem Verderben und Untergang. Amen.

\chapter{Das 25. Kapitel}

\section{Vorrede}

Salomo sagt in seinen Sprüchen im 24. Kapitel B. 16.: Ein Gerechter fällt sieben mal. Und unser HErr Jesus Christus spricht Luc. 13, 3.: So ihr euch nicht bessert, werdet ihr alle auch also umkommen. Dervwegen, weil der, so heute stehet, morgen fallen kann, und der, so jetzt fromm ist, gar bald sündigen kann, und weil Gott die Sünden nicht will ungestraft lassen; so sollen wir vorsichtig wandeln und uns an den Stab göttlichen Wortes halten, auf daß, weil wir in dieser Welt, als auf einem glatten Eis gehen, nicht fallen; und da ja jemand strauchelt oder einen Stolper thät, der nehme die treuen Bußpredigten an, und gebrauche derer recht, der erkenne, daß er übel und unrecht gethan, der bitte, Gott wolle ihm solche Sünde um des HErrn Christi Verdienstes willen vergeben; der hüte sich für vorgethanem und allem andern Sünden, der lebe christlich und ehrbar, der bete auch, Gott, der heilige Geist, wolle ihm helfen, daß er in Gottesfurcht beständig lebe, damit Gottes Zorn nicht auch wie Wasser auf uns ausgegossen werde, und uns an Leib und Seele, Gut und Ehre zeitliche und ewige Strafe nicht folgen. Dies soll uns auch erinnern das 25. Kapitel.

\section[Regierung Amazias, des Königs in Juda]{Regierung Amazias, des Königs in Juda}
\quellseite{654}

\section{Summarium}

Dies ist eine sonderlich feine Historia, in welcher erstlich zu merken ist, wie Gott dem König Amazia warnen läßt, daß er die hundert tausend Mann, die ihm aus dem Königreich Israel zu Hülfe gezogen waren, wieder soll heim schicken, V. 10. Darum, daß Gott kein Glück geben will, wo man Hülfe bei denen sucht, so gottlos sind, unrechte Lehre und falschen Gottesdienst haben und vertheidigen. Um solcher Ursachen willen ist zu besorgen, daß wir bisher wider den Türken so wenig haben ausgerichtet.

Zum andern ist hie wieder ein schrecklich Exempel, daß uns zur Furcht Gottes halten soll. Nachdem der König Amazia so einen herrlichen Sieg durch Gottes Hülfe erlanget hatte, fällt er von Gott ab, nimmt die Götter der Edomiter (welche er geschlagen hatte), und stellet sie ihm zu Göttern, und betet an für ihnen, V. 14. Und da er von einem Propheten gestraft wird, verdroß es ihn, und dräuet dem Propheten, V. 16. Solche Sünde machts bald aus mit seinem Regiment.
\quellseite{655}

\section{Summarium}

Sie ist eine Historia, in welcher man siehet, ob wohl Amazia eine redliche Ursache hatte wider den König Israel, darum, daß sein Kriegsvolk ohne alle Ursache drei tausend erschlagen, und etliche Städte geplündert; dennoch hat er kein Glück; Ursache war diese, daß er sich an der Heiden Abgötterei gehangen hatte. Um solcher Sünde willen verlieret er die Schlacht, Jerusalem wird geplündert und zerrissen, B. 23., und er, der König Amazia, wird endlich erstochen, B. 28. Solchen Lohn sollen alle die gewarten, so Abgötterei und falscher Lehre folgen und sie handhaben.

\section{Gebet}

Helfe GOtt, daß wir uns für gottloser Leute Gesellschaft und Gemeinschaft hüten; verleihe auch, daß wir uns aus den Strafpredigten bessern und in der Frömmigkeit verharren, und lindere die wohlverdienten Strafen, um JEsu Christi, seines lieben Sohnes willen, Amen.

\chapter{Das 26. Kapitel}

\section{Vorrede}

Obschon ein Mensch eine Zeitlang gottesfürchtig und glückselig ist, und thut in seinem Beruf, was Gott gefällt und den Leuten nützlich ist; dennoch kann er fallen, sündigen, unrecht thun und in zeitlichen und ewigen Jammer kommen, wenn er aus Fürwitz oder Frevel verführet durch sein Fleisch und Blut, oder durch die Welt, oder durch den Teufel wider Gottes Wort und Gebot, in seinem Stande thut und handelt. Dieses sollen wir wohl betrachten und täglich also beten: O HErr, bestätige das gute Werk, das du in uns gewirket hast, und die Werke deiner Hände wollest du ja nicht lassen, Ps. 138, 8. Und wollest solches thun um der Ehre deines heiligen Namens willen, und daß es zu unseres Nächsten Nutz gelange, und uns zu zeitlicher und ewiger Wohlfahrt diene, um des Verdienstes JEsu Christi, deines geliebten Sohnes und unsers HErrn und Heilandes willen, Amen. Das lehret auch das 26. Kapitel.

\section[Usia räuchert und wird aussätzig]{Usia räuchert und wird aussätzig}

a (Gesichte) Gesichte sind Prophezeiungen. Er will sagen, Sacharja sei gelehrt gewesen in Mose und den Propheten, Samuel, David, Gad und dergleichen.
\quellseite{656}

\section{Summarium}

Dies ist erstlich wieder ein Exempel, daß man siehet, wie viel einem Regenten daran gelegen sei, daß er einen rechten Prediger und Lehrer habe. Denn hie stehet man, so lange Zacharias lebet, der Lehrer in den Gesichten Gottes, suchet der König Gott und gibt ihm Glück, daß ihm in allen Dingen wohl gelinget, B. 5. Da aber Zacharias stirbt, erhebet sich Usia seiner Macht und Glückes, verachtet die Priester, und will beides, König und Hohepriester sein, und selbst auf dem Altar opfern, B. 16. Um solcher Sünde willen strafet ihn Gott mit dem Aussatz sein Lebenlang, B. 19. Denn Könige und Fürsten sollen eben sowohl unter Gottes Wort sein und demselbigen gehorchen, als andere Leute; thun sie es nicht, sollen sie eben sowohl ihre Strafe darum leiden, als andere.

\section{Gebet}

GOtt, unser lieber Vater, erbarme dich unser (denn wir bisher seinem Worte auch nicht aller Dinge gefolget haben); und verleihe, daß wir nun ferner seine gehorsamen Kinder sein und bleiben, und von ihm an Leib und Seele gesegnet werden, und in unserm Beruf Glück und Wohlfahrt behalten, Amen.

\chapter{Das 27. Kapitel}

\section{Vorrede}

David sagt im 103. Psalm.: Wie sich ein Vater über Kinder erbarmet, so erbarmet sich der HErr über die, so ihn fürchten. Und B. 17.: Die Gnade des HErrn währet von Ewigkeit zu Ewigkeit, über die, so ihn fürchten. Und im andern Buch der Könige im 17. Capitel B. 39. spricht der heilige Geist: Fürchtet den HErrn, euren Gott, der wird euch erretten von allen euren Feinden. Freilich beschützt Gott alle, die ihn fürchten, wider alle Teufel und wider die ganze Welt, der er mächtig und gewaltig ist und ihnen zu gebieten und zu verbieten hat, und ein solcher HErr ist, der alle Dinge übersiehet und überhöret. Derwegen, wem er wohl ist, dem kann niemand übel; wer ihn zum Freunde hat, dem kann kein Feind schaden; wem er liebet, den kann kein Haß beschädigen. Dieser Sprüche Exempel ist das 27. Capitel.

\section[Jotham macht die Ammoniter zinsbar]{Jotham macht die Ammoniter zinsbar}
\quellseite{657}

\section{Summarium}

Ist eine Historia, daß Jotham, der König Juda, viele Siege und Glück hat, um der Ursache willen, daß er (wie der Text sagt) seine Wege richtet für Gott, seinem Herrn, das ist, er lebet in Gottesfurcht und nach dem Wort Gottes. Zum Exempel aller Obrigkeit: so sie glückselige Regierung begehrt, daß sie sich gottselig halten und nach Gottes Wort in der Furcht Gottes leben soll; wo nicht, so sollen sie alles Unglücks gewarten.

\section{Votum}

Der treue GOtt gebe der Obrigkeit und Unterthanen einerlei Herz, und regiere sie durch seinen Geist und Wort, daß sie ihn fürchten ihr Lebenlang, auf daß es ihnen und ihren Kindern wohl gehe und kein Unglück über sie komme, Amen.

\chapter{Das 28. Kapitel}

\section{Vorrede}

Gottes Wort lehrets und die allgemeine Erfahrung bezeugets, daß Gott fromme Leute segnet und erhält; und dagegen, daß Gott unbußfertige Sünder strafet und verderbet. Derwegen laßt uns alles sündliche Wesen fliehen, und aller christlichen Tugenden befleißigen, mit vorgehendem Gebet um Hülfe und Regierung des heiligen Geistes; so wird unsere Arbeit, Regieren, Lehren, Handeln und Obern nicht eine vergebene, unnütze Mühe sein, sondern wird von Gott gefördert und aus Gnaden mit Glück und Fortgang belohnet werden. Denn also lautet davon die gnädige Verheißung Gottes, Esaiæ 65, 23.: Meine Auserwählten sollen nicht umsonst arbeiten. Und St. Paulus sagt von denen, die ihren Beruf treiben in wahrem Glauben, mit Gebet und schuldigem Fleiß, 1 Cor. 15, 18.: Eure Arbeit ist nicht vergeblich in dem HErrn. Darauf höret auch an das 28. Kapitel.

\section[Von Ahas, dem Könige in Juda, und seinem]{Von Ahas, dem Könige in Juda, und seinem gottlosen Wesen}
\quellseite{658}

\section{Summarium}

Ist eine Historia, wie die Abgötterei am König Ahas greulich gestraft wird. Erstlich, daß ihm die Syrer ins Land fallen und viele Leute wegführen, V. 5. Darnach der König Israel, Pekah, fällt ihm auch ins Land und erschlägt ihm auf einen Tag hundert und zwanzig tausend Mann, und führet zwei hundert tausend mit sich gefangen weg, V. 6. An solcher Strafe sollte weltliche Obrigkeit und sonst jedermann lernen, wie greuliche Sünde es um Abgötterei, falschen Gottesdienst und unrechte Lehre sei, und derhalben mit höchstem Fleiß sich dafür hüten, und sich zu dem Haufen halten, da rechte Lehre und der Gottesdienst rein und rechtschaffen ist.

Darnach ist eine Predigt des Propheten Oded, in welcher er solche Tyrannei strafet, welche das Königreich Israel an ihren Brüdern geübt hatte, V. 9. Denn Gott will, daß die Strafe nicht gar ohne alle Maße gehen soll.

a Unsere Israel.

b Bloß hieße war, daß das Volk nicht unter Gott nach seinem Worte lebte, sondern frei, nach seinem eigenen Gutdünken im Gottesdienst, wie 2 Mos. 32, 25. Aaron das Volk entblößet.

\section{Summarium}

Dies ist eine Historia, in welcher man siehet, je mehr Gott den König Ahas strafet um seiner
\quellseite{659}
Abgötterei willen, je mehr treibet er Abgötterei. Gleich wie unsere Papisten auch thun: wenn sie Gott strafen mit den Türken und andern Unglück, richten sie Procession an, lassen Messe lesen, suchen Wallfahrten und geloben sich zu den Heiligen. Suchen also wie Ahab, des rechten Gottes Hülfe nicht, sondern setzen alle ihren Trost und Hülfe auf ihre eignen Werke und abgöttischen Ceremonien. Darum gehets auch, wie der Text hie meldet, V. 22., daß die Gottlosen in der Noth des Verzweifelns nur je länger je mehr machen; wie an Ahab zu sehen ist, welcher die Thüren am Hause des HErrn zuschloß, und machet ihm Altäre in allen Winkeln zu Jerusalem, V. 24. Also thun jetzt unsere Papisten auch: das heilige Evangelium lästern und verdammen sie als Ketzerei, wollen derhalben in ihren Landen nicht leiden. Dagegen leiden und dulden sie alle Abgötterei, Sünde und Schande. Sollte aber Gott die nicht strafen und solches rächen?

\section{Votum}

GOtt sei Lob, Ehre und Dank, der uns von Abgötterei erlöset hat, der stehe uns ferner bei; strafe an uns nicht anderer Leute Abgötterei, und helfe, daß ein jedes unter uns, nach seinem Beruf und Vermögen, sein Wort fördern helfe, Amen.

\chapter{Das 29. Kapitel}

\section{Vorrede}

Keinen Gottesdienst erfordert Gott so ernstlich, hart und heftig, als daß man die Abgötterei abschaffen und rechte Lehre und Gottesdienst fördern und erhalten soll. Derwegen alle christliche Obrigkeit die Abgötterei aus den Augen ihrer Unterthanen, die Prediger aber dieselbige aus den Herzen ihrer Zuhörer abschaffen und wegräumen sollen. Hieher gehören diese Sprüche. Von der Obrigkeit sagt das Buch der Weisheit im 6. Kapitel 22.: Wollt ihr nun, ihr Tyrannen, das ist Regenten, in Volk gerne Könige und Fürsten sein; so haltet die Weisheit in Ehren, auf daß ihr ewiglich herrschet. Von den Predigern insgemein lehret auch St. Paulus 1 Tim. 4, 16.: Hab Acht auf dich selbst und auf die Lehre, beharre in diesen Stücken. Denn, wo du solches thust, wirst du dich selbst selig machen, und die dich hören. Solche Obrigkeit und Lehrer hat der treue, barmherzige Gott dieser löblichen Gemeine auch gegeben. Dafür danken wir dir heute alle, o reicher und milder Gott ins Himmel Throne, mit höchstem Fleiß, und bitten dich herzlich, du wollest unsere liebe Obrigkeit und Prediger ferner segnen, stärken und erhalten, und ihres gleichen geben bis an den jüngsten Tag. Dazu sprech alle von Herzen: Amen. Darauf höret ferner an das 29. Kapitel.

\section[Jehiskia richtet den rechten Gottesdienst wieder an]{Jehiskia richtet den rechten Gottesdienst wieder an}
\quellseite{660}

\section{Summarium}

Ist eine sonderliche Historia, welche alle Regenten sollten mit Fleiß bedenken und sie ihnen lassen befohlen sein wie einen Spiegel, darin sie sehen sollten, was fürnehmlich ihr Amt wäre, wenn sie wollten recht und glücklich regieren: nämlich, daß sie am ersten sollten den Tempel öffnen, daß ist, Gottes Wort und den rechten Gottesdienst fördern und unter die Leute bringen, wie der fromme König Hiskiä thut, welcher die Priester und Leviten seine Söhne heißt, sie bittet und vermahnet, sie sollen nicht nachlässig sein, sondern des Gottesdienstes mit Fleiß warten, V. 11. Gegen solchem Könige sollten sich alle gottlose Regenten schämen, so entweder Gottes Wort und desselben Diener verachten, oder wohl zum Land ausjagen, verfolgen und tödten. Aber was für Glück bei ihnen sei, siehet man auch wohl, wer es anders sehen will.
\quellseite{661}

\section{Summarium}

Ist ein schön Exempel, wie der König Hiskiä selbst mit und bei dem Gottesdienst ist. Denn solches ist weltliche Obrigkeit schuldig, auf daß sie mit ihrem Exempel den Unterthanen fürgehe. Daß aber der Herr sagt vom Sündopfer für den König und das ganze Volk, ist nicht zu verstehen, als sei ihnen die Sünde um solches Opfers willen vergeben. Denn Vergebung der Sünden ist je und je gestanden auf dem einigen Opfer des Sohnes Gottes. Desselben Blut allein erwirbt Vergebung der Sünden. Die Sündopfer aber im Gesetz sind nur eine äußerliche Anzeigung gewesen, damit das Volk bekennen sollte, daß sie wider Gottes Gebot und Satzung gehandelt hätten, und also öffentlich solche Mißhandlung für der Gemeine abkehren mit gewöhnlichem Opfer, auf daß sie nicht würden von der Gemeine, als öffentliche Verbrecher, abgesondert. Also lässet der König sich und das Volk versöhnen, zum öffentlichen Bekenntniß, daß sie sich begehreten zu bessern; aber durch solches Opfern ist die Sünde für Gott nicht bezahlet worden, sondern sie haben sich mit dem Herzen zu Gott wenden und Vergebung durch den verheißenen Samen hoffen müssen. Da ist die rechte Vergebung gefolget für Gott, im Gewissen, da jene Vergebung allein für die Leute gegolten hat.

\section{Votum}

GOtt sei uns gnädig und vergebe uns alle unsere Sünde; er wolle auch ferner allzeit geben und erhalten dieser gemeinen Stadt und Vaterland christliche Obrigkeit und treue Lehrer, bis zur Wiederkunft unsers HErrn JEsu Christi, und solches zu seinem Lob und den gehorsamen Unterthanen und Zuhörern zur seligen Wohlfahrt, Amen.

\chapter{Das 30. Kapitel}

\section{Vorrede}

Man schreibt, daß Diagoras Melius, ein Epicurer, der heidnischen Opfer gespottet und gesagt: Es schmecke ihm dennoch Essen und Trinken wohl, ob er gleich den Göttern nicht opferte. Des haben die Athenienser also geeifert, daß er hat entlaufen müssen, und haben sein ganzes Vaterland, die Insel Melos, verwüstet, dieweil ein solch epicurisch Weltkind darinnen erzogen war. Sollte des nun Gott, der HErr, nicht ernstlich eifern und hart strafen, wenn man die seligmachende Lehre seines Worts verachtet? also: ein anderer, mit Namen Pherecides, ist von den Läusen aufgefressen, daß er die heidnischen Opfer verachtete. So wird freilich das höllische Feuer auffressen, die Gottes Wort abermal verachten, darinnen gepredigt wird von dem rechten Opfer für unsere Sünde, nämlich von unsers lieben HErrn JEsu Christi Person, Verdienst, Wohlthaten und Willen. Derwegen wollen wir uns alle hüten für Verachtung göttlichen Wortes, und wollen dagegen dasselbe gern hören, fleißig lernen und christlich darnach glauben und leben, in Kraft und durch Regierung Gottes, des heiligen Geistes, Amen. Folgendes höret an das 30. Kapitel.

\section[Ausrufung und Haltung des Passah]{Ausrufung und Haltung des Passah}
\quellseite{662}

\section{Summarium}

Dies ist eine schöne Historia, wie der fromme König Hiskiä Prediger ausgeschickt in das ganze Land, und vermahnet sein Volk zur Buße, V. 1. Und tröstet sie wiederum, wenn sie sich von Herzen bekehren zum HErrn, so sei Gott gnädig und werde nicht allein ihnen, sondern auch ihren Kindern, so im Gefängniß unter den Heiden waren, gnädige Hülfe erzeigen, V. 6. Solche Predigt und Vermahnung verlacht der meiste Theil, wie denn der Welt Art ist, V. 10. Dennoch waren etliche, die es annahmen und besserten sich, V. 11. Denn Gottes Wort gehet nicht ab ohne Frucht.

Darum soll dies Exempel ernstlich die Obrigkeit wohl bedenken, und wissen, daß sie für Gott schuldig sei, rechte Prediger zu bestellen, und unrechten Gottesdienst abzuthun, wie Hiskiä hie thut.

Zum andern sollen die Unterthanen sich sonderlich des befleißigen, daß sie Gottes Wort nicht verachten, sondern von Herzen und mit großem Ernst es annehmen, und sich bessern.

a Heiligen Reinigkeit) das ist, Gott siehet das Herz an; wenn das rechtschaffen ist an Gott, so fraget er nicht nach äußerlicher Reinigkeit, die nach dem Gesetze heilig ist.
\quellseite{663}

\section{Summarium}

Ist ein Exempel, welches sonderlich wohl zu merken ist. Es hatte seine sonderliche Ordnung mit dem Osterlamm, wie im andern Buch Mose stehet. Nun konnte man in der Eile nicht aller Dinge solche Bräuche und Ceremonien alle halten, wie im Gesetz geboten war.

Da greift der König Hiskia in solche Gesetzordnung, und lehret, man soll mehr darauf sehen, daß das Herz gegen Gott rechtschaffen sei und sich des äußerlichen Wesens halber nicht sonderlich bekümmern.

Denn Gott wolle gnädig sein, nicht um solcher äußerlichen Werke oder Reinigkeit willen, sondern um sein selbst willen, so fern die Herzen rein sind, das ist, so sie ihre Sünde erkennen und auf Gottes Zusagung und Barmherzigkeit trauen; zeigt dadurch an, daß es alles nicht am äußerlichen Thun und Lassen, sondern am Herzen gelegen sei.

\section{Votum}

Gott, der heilige Geist, erleuchte und regiere uns, daß wir innerlich und äußerlich mit Herz, Worten und Thaten fleißig christlich leben und die Gottesdienste üben; und helfe, daß unsere liebe Obrigkeit ferner seine Kirche fördere und schütze, und daß wir Gottes Wort hoch achten, und zu unserer Besserung gebrauchen, Amen.

\chapter{Das 31. Kapitel}

\section{Vorrede}

Gott befehlet der Obrigkeit im 5. Buch Mose im 12. Kapitel V. 3. also: Reißet um die Altäre der Heiden, und die Götzen ihrer Götter thut ab, und vertilget ihren Namen. St. Paulus Gal. 6, 6. lehret dankbare Kirch-Kinder also: Der unterrichtet wird mit dem Wort, der theile mit allerlei Gutes dem, der ihn unterrichtet. Und Sirach im 1. Kapitel V. 19. unterweiset insgemein alle Christen also: Wer den HErrn fürchtet, dem wird wohl gehen. Dieser Sprüche Exempel erzählet das 31. Kapitel.

\section[Abschaffung der Abgötterei, Versorgung der Priester und Leviten]{Abschaffung der Abgötterei, Versorgung der Priester und Leviten}

a Anhalten) auf daß sie nicht aus Mangel der Nahrung müßten die Bücher lassen, beide zu studiren und zu lehren, und ihre Nahrung suchen. Denn Kirchen-Diener sollen versorget sein und studiren, wie Nehemias und Sirach auch sagen.
\quellseite{664}
und ist noch viel überblieben; denn der HErr hat sein Volk gesegnet, darum ist dieser Haufe überblieben. *4 Mos. 5, 9.

\section{Summarium}

Die ist erstlich eine Historia, wie dem Exempel und Befehl des Königs Hiskiä nach alle abgöttischen Altäre und Kirchen im Lande sind umgerissen worden. Denn weltlicher Obrigkeit (wie oft gemeldet) stehets sonderlich zu, daß sie solches nicht dulden, sondern mit Ernst abschaffen sollen.

Darnach ist eine andere Historia, daß der König das Volk vermahnet, sie sollen ihnen die Leviten und Priester, so den Gottesdienst verwalten müssen, in lassen befohlen sein und von dem Ihren mildiglich und treulich ihnen geben, V. 4. Denn weltliche Obrigkeit soll mit allem Fleiß daran sein, daß sie nicht allein selbst gern dazu geben, daß die Kirchen-Diener mögen erhalten werden, sondern auch ihre Unterthanen dahin halten, daß sie es auch thun.

Am Ende ist ein schöner Spruch, nach welchem sich alle Menschen, sonderlich aber weltliche Obrigkeit halten soll, da er spricht: Hiskiä habe in allem, das er anfing, Gott von Herzen gesucht; darum habe er auch Glück gehabt.

\section{Votum}

Regiere und helfe uns Gott, um Jesu Christi willen, durch den heiligen Geist, daß wir alle Abgötterei fliehen, das Predigtamt nach Vermögen helfen fördern, und aus Gnaden Glück, zeitlich und ewig haben, an Leib und Seele, Amen.

\chapter{Das 32. Capitel}

\section{Vorrede}

Der königliche Prophet David sagt im 34. Psalm V. 7.: Da dieser Elende rief, hörte der HErr, und half ihm aus allen seinen Nöthen. Es ein tröstlicher Spruch für alle Nothleidende, die Gott unterm Kreuz in wahrem Glauben anrufen. Aber da möchte jemand gedenken und sagen: Ich erfahre das Gegenspiel, daß manches fleißig und herzlich betet und gleichwohl in seiner Noth stecken bleibt, als in Armuth, Verfolgung, Krankheit und dergleichen. Da sollt ihr wissen, ihr Kinder Gottes, daß wir in zeitlichen Nöthen unsern Willen in unsers lieben Gottes Willen stellen und auf ihn hoffen sollen; ob er uns gleich tödtet, und Leib und Seele ließ verschmachten. Er hat sein Bedenken, warum er also die Seinen mit der Noth ringen läßt; und langsam kömmt er, kömmt aber wohl; und, lange verzogen, heißt nicht gelogen; und im Fall, daß wir hie in diesem Leben (weil Gott weiß, daß es uns nicht zu unserer Seelen Heil und Seligkeit nutz und gut ist) nicht errettet würden; so sollen wir doch reichlich und überflüssig in jenem Leben gewißlich erhöret und unsere Bitte tausend, ja hundert tausendfältig gewähret werden. Denn Gott läßt die Zinsen zusammen über einen Haufen wachsen und wird sie uns darnach auf einmal über einen Haufen geben, im ewigen Leben. Dieses lehret auch das 32. Kapitel.
\quellseite{665}

\section[Hiskias wunderbare Beschützung wider Sanherib]{Hiskias wunderbare Beschützung wider Sanherib}

\section{Summarium}

Ist eine Historia, daß der fromme König Hiskia, der Gottes Wort und Dienst so herrlich gefördert hat, dennoch sammt den Seinen das Kreuz und sehr große Krieges-Noth tragen müssen. Da ist nun sonderlich zu merken, wie er in solcher Noth sein Volk nicht auf Menschen-Hülfe, sondern auf Gottes Hülfe weiset und tröstet, welche Hülfe allein nicht lügen und fehlen kann, V. 7.

Wiederum, was man außer derselbigen Hülfe suchet und hoffet, es sei Geld, Gut, Gewalt, Freundschaft, feste Gebäude, großer Vorrath und anders, ist alles nicht allein ungewisse, sondern auch schädlich, wo man Gott nicht dabei hat, wie alle Historien zeugen.
\quellseite{666}
*Enthielt, wie ein Hirt seine Schafe hält oder die Wölfe, und hütet, daß sie gehen hin und her zur Weide: also konnten die zu Jerusalem auch aus- und einziehen sicher etc.

\section{Summarium}

Dies ist erstlich ein Exempel, daß Gott, wenn man ihn in der Noth sucht und seiner Hülfe gewartet, das Gebet annehmen und gnädig helfen wolle, sollte es gleich durch Wunderzeichen müssen geschehen.

Darnach ist ein ander Exempel, B. 27., in dem wir sehen, wie auch große Heiligen, so durch den festen Glauben in der höchsten Noth bestanden sind, in Glück und Wohlfahrt nicht bestehen und sich durch Sicherheit und Hoffart versündigen, auf daß wir lernen in Gottes Furcht wandeln und täglich bitten, daß Gott die Hand nicht von uns abziehe, sondern uns in der Anfechtung erhalten und behüten wolle.

\section{Votum}

Helfe GOtt die hochgelobte Dreifaltigkeit, daß wir beständig bleiben in wahrem Glauben und in der Gottseligkeit; und erhöre auch unser Gebet, und helfe uns aus aller Noth, Leibes und der Seelen, Amen.

\chapter{Das 33. Kapitel}

\section{Vorrede}

Der heilige Prophet Micha sagt im 7. Kapitel B. 18. 19.: Wo ist solch ein Gott, wie du bist? Der die Sünde vergibt und erläßt die Missethat den Uebrigen seines Erbtheils. Der seinen Zorn nicht ewiglich behält, denn er ist barmherzig. Er wird sich unser wieder erbarmen, unsere Missethat dämpfen und alle unsere Sünde in die Tiefe des Meers werfen. Wenn nun euch Gott, den HErrn, als einen unbarmherzigen Stockmeister der Teufel will einbilden, der sich weder erbarmen, noch Sünde vergeben, sondern nur ewiglich Zorn halten wollte; alsdann sollt ihr ihn abfertigen mit seinem falschen Bilde, und zu ihm sprechen: Gehe hin, Teufel, und lerne das Handwerk daß, du kannst Gott nicht recht abconterfeien, noch seine rechte Gestalt treffen. Es ist ein großer Prophet, der heißt Micha, der hat mir Gott, den Allmächtigen, viel anders fürgemalet und geben ihm alle Propheten Beifall, nämlich, daß Gott Sünde vergebe und erlasse die Missethat allen denen, die sich zu ihm bekehren; er erbarme sich über uns, er dämpfe unsere Missethat, und werfe alle unsere Sünde in die Tiefe des Meers, das ist, er vertilget sie also aus Gnaden, daß er derselben nimmermehr in Ungnaden eingedenk sein will, nicht anders, als wenn er sie alle in ein Gebündlein zusammen geknüpfet und mit einander ins Meer geworfen hätte, daraus sie nimmermehr empor schwimmen, noch ans Licht
\quellseite{667}
kommen sollen. So heißt der Text, Teufel! Darum troll dich, du feindseliger Geist, wohin du gehörst; mir soll Gottes Wort, welches in Ewigkeit bleibet, viel theurer und werther sein, denn alle dein falsches Eingeben. Denn Gott lügt nicht und sein Wort und Verheißung fehlet nicht. Dies betrachtet auch bei dem 33. Kapitel.

\section[Regierung Manasses und Amons]{Regierung Manasses und Amons}
\quellseite{668}

\section{Summarium}

Dies ist eine schöne Historia, in welcher man erstlich siehet, wie der fromme König Hiskiä einen ungerathenen gottlosen Sohn hat, B. 1., der allerlei Abgötterei stiftet und auf Gottes Wort keine Achtung hat. Darum wird er von den Assyrern überfallen und gen Babel gefangen weggeführet, B. 11.

Da ist nun sonderlich ein Exempel der großen Gnade und Barmherzigkeit Gottes, daß er einen solchen gottlosen Tyrannen, nachdem er seine Sünde erkennt und um Gnade bittet, wiederum zu Gnaden annimmt und in die vorigen Würden und Gewalt setzet, B. 13. Zu Trost uns allen, daß wir von Sünden ablassen und uns bessern, und alsdann der Gnade Gottes gewarten sollen und ja nicht verzweifeln.

Sonderlich aber ist zu merken, daß Manasse, nachdem er wieder vom Gefängniß ledig war, alle Abgötterei abschaffet und dem ganzen Volk befiehlet, daß sie dem Herrn, dem Gott Israel dienen sollen, B. 14. Welches eine Anzeigung ist einer rechten Buße \&c.

Am Ende ist ein ander Exempel, von seinem Sohn Amon, der gottlos war, B. 20., aber besserte sich nicht wie sein Vater, darum treibt ers nicht lange. Denn er ward von seinen eigenen Knechten ermordet, B. 24. Aber die Aufrührer wurden gestraft, B. 25.

\section{Gebet}

Der barmherzige Gott vergebe uns auch gnädiglich alle unsere Sünde, lindere uns die wohlverdienten Strafen, helfe, daß wir allzeit christlich leben, und regiere die lieben Kinder, daß sie auch Gott und ihren Eltern gehorsam sein, Amen.

\chapter{Das 34. Capitel}

\section{Vorrede}

Wie der Prophet Esaia, also der Sohn Gottes, unser Herr Jesus Christus, nennt die Tempel oder Kirchen Gotteshäuser oder Bethäuser; darum, daß man die zu Gottes Ehre gebrauchen soll und darin sein Wort lehren und hören, seine Sacramente handeln und gebrauchen, das Gebete und andere Gottesdienste verrichten. Zu solchem Gebrauch der Kirchen- und Gotteshäuser haben allezeit alle christliche Obrigkeit und treue Seelsorger ihre Unterthanen und Zuhörer fleißig ermahnet und gehalten. Denn, anderer Exempel zu geschweigen: Im Concilio Elibertino ist beschlossen worden, wenn einer einheimisch und gesund gewesen, und in drei Wochen in keine Kirche kommen, hat man ihn unwerth geachtet, daß er zum Abendmahl des Herrn gelassen würde. Desgleichen zu Karthago, in vierter Synode, hat man also geschärfet, daß man die in Bann gehen, die unter der Predigt aus der Kirche ohne Noth gelaufen sind. Derwegen habt lieb die Stätte des Hauses Gottes, da man predigt seine Wunder und haltet euch zu seinem Altar, da man höret die Stimme des Dankens, Ps. 26, 7. Und seid fleißig in allen Diensten Gottes. Ferner höret an das 34. Capitel.

\section[Josia reiniget und bessert den Tempel]{Josia reiniget und bessert den Tempel}
\quellseite{669}

\section{Summarium}

Hie ist eine schöne Historia von dem frommen gottseligen König Josia, daß er erstlich alle Abgötterei ausrottet, V. 3., und darnach den Tempel wieder zugerichtet, daß man den Gottesdienst darin üben konnte, V. 8. Solche Exempel sollen weltliche Obrigkeit sonderlich merken, daß sie Gottes Wort und Ehre fördern und aller Abgötterei und falschem Gottesdienst mit Ernst wehren sollen; so will Gott bei ihnen sein, Glück und Heil geben.
\quellseite{670}

\section{Summarium}

Sie ist eine sonderliche Historia, wie Josias, nachdem das Gesetzbuch funden und für ihm gelesen ist worden, seine und des Volks Sünde erkennt, und lässet ihm leid sein, V. 19. Solche Reue und Demuth gefällt Gott wohl. Darum verheißt er, obwohl die Sünden müssen gestraft werden, so wolle er doch die Strafe aufhalten und Josiam in Frieden sterben lassen, V. 27. Solche Verheißung sollte uns auch zur Buße und Besserung bewegen.

Am Ende ist eine Historia, daß Josia das Gesetz-Buch öffentlich habe lesen lassen und sich und das Volk auf ein neues gegen Gott wieder verpflichtet und mit ihm verbunden.

\section{Votum}

GOtt, der Vater unsers HErrn Jesu Christi, verleihe uns auch wahre Demuth und Reue über unsere Sünde, vergebe uns auch dieselbe, und ziehe auf die wohlverdienten Strafen; und helfe, daß wir sein Wort gern hören, und darnach christlich und beständig leben, Amen.

\chapter{Das 35. Kapitel}

\section{Vorrede}

Gott will und begehret es nicht, daß wir ihm etwas umsonst thun sollen. Denn Gott ist nicht ungerecht, daß er vergesse eures Werkes und Arbeit in der Liebe, die ihr beweiset habt in seinem Namen, sagt die Epistel an die Ebräer im 6. Kapitel V. 10. Und der Prophet Maleachi 3, 16. spricht: Der HErr merket und höret, und ist für ihm ein Denkzettel geschrieben für die, so den HErrn fürchten und an seinen Namen gedenken. Und Gott träget seine Gläubige und Gehorsame auf seiner Brust, und drücket sie wie ein Siegel auf sein Herz, im Hohenlied Salomonis 8, 6. Auch lehret also der hocherleuchtete Prophet Esaias und spricht im 3. Kapitel V. 10.: Gott läßt predigen von den Gerechten, daß sie es gut haben, denn sie die Frucht ihrer Werke essen sollen. Sonderlich am jüngsten Tage, sagt St. Paulus Röm. 2, 7., wird Gott einem jeden nach seinen Werken geben, nämlich Preis und Ehre und unvergängliches Wesen denen, die mit Geduld in guten Werken trachten nach dem ewigen Leben. Derwegen wollen wir in Gottesfurcht leben, daß Gott aus Gnaden segne, Glück und Wohlfahrt beschere, schütze und behüte uns alle, und einen jeden in seinem Beruf und Stande, Handel und Wandel, Arbeit und Nahrung, Amen. Darauf höret weiter an das 35. Kapitel.

\section[Josia hält Passah, und kommt um im Streit]{Josia hält Passah, und kommt um im Streit}
% Scan page 671: omitted; Bible text/narrative only, no retained apparatus.
\quellseite{672}

\section{Summarium}

Ist eine Historia, wie der fromme Josias das Osterfest wieder anrichtet und reichlich dazu hilft, B. 1. Denn weltliche Obrigkeit soll sichs etwas kosten lassen, daß der rechte Gottesdienst, und sonderlich das Predigt-Amt, recht bestellet und angerichtet werde.

Am Ende ist eine Historia, wie der fromme König Josias sich in einen unnöthigen Krieg eindränget, da er sonst hätte mögen Ruhe haben, und wird also erschossen, B. 22. Denn es fallen auch Heilige oft. Darum bedarf es desto mehr Betens, daß uns Gott in seinem Schuz erhalten und in der Anfechtung behüten wolle.

\section{Datum}

Helffe GOtt, daß unsere liebe Obrigkeit sich ferner des heiligen Predigtamts treulich annehme, und das fördere und erhalte; und bewahre uns in seinem Gehorsam bis an unser Ende, damit wir aus Fürwitz und muthwillig zu unserm Unglück und Herzeleid nicht Ursache geben, Amen.

\chapter{Das 36. Capitel}

\section{Vorrede}

Der weise Mann Sirach sagt im 10. Capitel B. 8.: Um Gewalt, Unrecht und Geiz willen kommt ein Königreich von einem Volk aufs andere. Dieses wahren Spruchs ist ein öffentlich und klar Exempel das jüdische Königreich. Also, hat der gerechte Gott des israelitischen Volkes, das er zuvor allen Völkern, Land und Leuten, vorgezogen hat, nicht verschonet; darum, daß diese vornehmen Hauptlaster, Abgötterei, Sicherheit, Geiz, Ungerechtigkeit, Gewalt und andere Sünde mehr, wie eine Sündflut eingerissen und mit Gewalt im Schwange gegangen sein, Ezech. 16, 49.: auch so mögen wohl alle für Gottes Zorn und Strafen sich auch hüten, die dergleichen Sünden gethan haben und in solchen Lastern noch leben. Denn Gott läßt sich nicht spotten, spricht St. Paulus an die Galater im 6. Capitel B. 7. Und darum sagt Job im 9. Capitel B. 28.: Ich hab Achtung auf alle meine Werke, denn ich wußte, daß du, HErr, derer nicht schonest, die da sündigen. Dieses lehret auch das 36. und letzte Kapitel.

\section[Von Josias Nachfolgern, und der babylonischen Gefangenschaft]{Von Josias Nachfolgern, und der babylonischen Gefangenschaft}

% Source scan 673; block decisions in selection.json

\quellseite{673}

a Frühe) das ist, er sagete ihnen zeitlich zuvor, daß sie Zeit genug hätten, allemal sich zu bessern, ehe die Straffe käme.

\section{Summarium}

Ist eine Historia, wie die Kinder des frommen Josia ungerathen und abgöttisch gewesen sind, und sie Gott darum gestraft habe, bis endlich der König zu Babel das Land gar verheeret und das Volck gefangen hinweg geführet hat, V. 17. 20. Solche Straffe soll uns ja ein Exempel und Warnung sein, daß wir uns für Abgötterei hüten, Gottes Wort mit Ernst aufnehmen und demselbigen nachwandeln sollen; wo wir anders nicht gleichen Zorn und Straffe mit den Juden tragen wollen, welche siebenzig ganzer Jahre unter den Heiden in allerlei Jammer und Elend, um solcher Sünden willen mußten leiden.

\part{Das Buch Esra}

\chapter{Das 1. Kapitel}

\section{Vorrede}

Bei diesem ersten Kapitel sollen wir merken zwei Stücke. Das erste ist eine Historia, wie das jüdische Volck wieder in sein Vaterland kommt aus dem babylonischen Gefängniß. Diese lehret, daß Gottes Verheißungen gewiß sein, er selbst wahrhaftig, und daß man göttlichem Wort billig und sicher glauben möge. Auch lehret die Geschichte, daß Gott allmächtig sey und deswegen die Seinen aus Creuz und Leiden erretten könne, auch alsdann, wenn es für der Vernunft und auch menschlichen Kräften unmöglich scheinet.

Das andere Stück ist das Exempel Kores und des babylonischen Volckes, daß, wie Kores des jüdischen Volckes: also die Obrigkeit der Kirchen Gottes sich erbarmen, ihrer sich annehmen, ihr Glück und Wohlfahrt fördern und fortsetzen soll; und, wie das jüdische Volck mit Silber und Gold geholfen hat zum Bau des Tempels zu Jerusalem: also soll jedermann Hülffe thun zur Förderung göttlichen Wortes und des rechten Gottesdienstes, und derwegen von Gott reiche Belohnung erwarten, und das soll geschehen, nicht mit Unwillen oder aus Zwang, sondern freiwillig. Denn wie die Schrift sagt: Einen fröhlichen Geber hat Gott lieb, 2 Cor. 9, 7. Und Sirach im 35. Capitel V. 20. spricht: Wer Gott dienet mit Lust, der ist ihm angenehm.

\section{Kores, oder Cyrus, entläßt die gefangenen Juden}
% Source scan 674; block decisions in selection.json

\quellseite{674}

\section{Summarium}

Bisher hat eure Liebe gehört, wie die Juden sich mit Abgötterei, falscher Lehre und Verachtung des Wortes Gottes in ihren Landen versündiget haben, und darum von den Babyloniern gefangen weggeführet sind, und siebenzig Jahr im Gefängniß blieben.

Diese Zeit aber ist der Prophet Daniel in einem großen Ansehen gewesen unter den Heiden. Und sind durch seine Predigt NebucadNezar und EvilMerodach bekehret worden. Aber unter dem dritten Könige Belsazer hat die Abgötterei und andere Untugend wieder gewachsen und überhand genommen.

Darum schicket Gott über Babylonien Darium aus Meden und Cyrum aus Persien; dieselben mußten eben an den Babyloniern die Verachtung des Wortes Gottes und andere Sünde strafen, wie sie vorhin die beiden Könige von Assyrien und Babel an den Juden gestrafet hatten. Solche beide Könige, Darius und Cyrus, sind durch den Propheten Daniel auch bekehret zum rechten Gott.

Weil aber Esaias und Jeremias von Cyrus hatten geweissagt, daß er sollte das Volk in ihr Land Juda ziehen und den Tempel wieder aufbauen lassen, nimmt Cyrus diesen Befehl an und folget solcher göttlichen Weissagung. Befiehlt auch dazu seinen Landen und Leuten (welches sonderlich ein fein fürstlich Werk ist), daß sie den Juden zu solchem Werk reiche Beisteuer thun und helfen sollen, V. 6. Also hat Gott seines Volks Gefängniß unter den Heiden dahin gebracht, daß sein Wort unter ihnen gepredigt und sie auch Gottes Volk und selig worden sind. Wie nun die Juden wieder zu Haus kommen, die Stadt Jerusalem und den Tempel wieder zu bauen angefangen und wie es ihnen darüber ergangen hat, werden folgende zwei Bücher, Esra und Nehemia anzeigen.

\section{Votum}

Der treue GOtt helfe, daß Obrigkeit und Unterthanen seiner Kirchen sich treulich annehmen, belohne ihnen alle ihre Wohlthaten reichlich, und steure allen Kirchfeinden durch seinen gewaltigen Arm, um der Ehre seines heiligen Namens willen, Amen.

\chapter{Das 2. Kapitel}

\section[Zahl und Register der wiederkehrenden Juden, und ihre]{Zahl und Register der wiederkehrenden Juden, und ihre Steuer}
% Source scan 675; block decisions in selection.json

% Scan 675: no apparatus retained; only biblical narrative/running matter.
% Source scan 676; block decisions in selection.json

\quellseite{676}

\chapter{Das 3. Kapitel}

\section{Vorrede}

Wir sind zwar von den mosaischen Opfern erlöst und befreiet, dieweil der HErr Christus sich nun selbst am Stamm des Kreuzes geopfert hat; dennoch, welche durch den Glauben des Leidens und Sterbens Christi theilhaftig worden sind, die sind gleichwohl schuldig, Opfer Gott zur Danksagung zu beweisen. Derwegen sollen wir Gott opfern einen betrübten, geängsteten Geist, von wegen unserer Sünden. Wir sollen Gott unsere Vernunft opfern und wider dieselbige glauben, was zu glauben ist, die gefangen nehmen und dem Gehorsam Christi unterwerfen. Wir sollen Gott Opfer bringen mit unserm Mund und Lippen, mit Beten und Danksagung. Wir sollen Gott opfern Gerechtigkeit gegen unsern Nächsten. Wir sollen Gott opfern Barmherzigkeit und die Werke der Liebe gegen den Menschen beweisen. Wir sollen Gott aufopfern unsere Kinder, durch Gebet, durch die Taufe, durch gute Unterweisung und Zucht und daß wir sie ihm willig und gern sollen folgen lassen, wenn er sie durch den Tod zu sich fordert. Wir sollen auch Gott unsern Leib aufopfern, mit Thun und Lassen, Leiden und Sterben und sollen kreuzigen unser Fleisch, sammt den bösen Lüsten und Begierden, und den alten Adam tödten. Dies betrachten wir auch bei dem 3. Kapitel.

\section[Vom Brandopferaltar. Fest der Laubhütten und Gründung des]{Vom Brandopferaltar. Fest der Laubhütten und Gründung des Tempels}
% Source scan 677; block decisions in selection.json

\quellseite{677}

\section{Summarium}

Ist eine Historia, wie Jesua und Serubabel endlich den Altar baueten, Es. 2., auf daß, ob man mit dem andern Gebäude so bald nicht möchte fertig werden, sie dennoch opfern und den Gottesdienst üben könnten. Denn weil Gott die Opfer selbst geboten und verordnet hatte, war es ein Werk des Gehorsams. Siehet man, wie ein fein biblisch Werk das sei, dazu helfen und rathen, daß der rechte Gottesdienst, so gefallen ist, wieder angerichtet und gefördert werde.

\section{Votum}

GOtt sei Lob und Dank, der zu aller Zeit christliche Obrigkeit gibt und erhält; der segne und bewahre ferner unsere liebe Obrigkeit und alle andere christliche Regenten, die ihn ehren und seine Kirche fördern, und auf ihrer Unterthanen Heil und Wohlfahrt trachten, Amen.

\chapter{Das 4. Kapitel}

\section{Vorrede}

Daß die gläubigen Bekenner und Diener des HErrn Christi angefeindet, fälschlich angegeben, verhindert und verfolget werden, ist nicht heuer neu. Denn solches hat sich bald nach der Erschaffung und Fall des Menschen angefangen, und währet noch also und wird bleiben bis ans Ende der Welt. Aber solches sollen wir Christen nicht achten, sondern sollen unter solchem Kreuz uns trösten, daß uns Gott dennoch lieben, erhalten und fördern will, und daß er uns reichlich wird erstatten, was uns die Welt Schaden und Nachtheil zufüget. Denn Gott läßt den Seinen hinweg nehmen die zeitlichen Güter, auf daß er ihnen gebe die ewigen himmlischen Güter und Reichthümer. Gott läßt den Seinen nehmen und abziehen ein geringes Kleid, damit er sie ziere und schmücke mit dem güldenen Kleid der ewigen Herrlichkeit. Gott läßt den Seinen das zeitliche Leben nehmen, daß er ihnen gebe das ewige Leben. Das läßt auch unsern Trost sein unter allem Haß, Neid, Widerwärtigkeit und Verfolgung der Welt und ihrer Helfer. Und höret darauf an das 4. Kapitel.

\section{Hinderung des Tempelbaues, von der Juda Feinden angestellet}
% Source scan 678; block decisions in selection.json

\quellseite{678}

\section{Summarium}

Dies ist eine Historia, in welcher man siehet, daß der Teufel dem rechten Gottesdienst feind ist, denselbigen nicht leiden kann, sondern mit Lügen und wie er kann, begehrt zu dämpfen. Denn das Volk in dem Lande durch den Teufel getrieben, hindert nicht allein den Bau des Tempels, V. 4, sondern verklagt auch fälschlich die Juden als Aufrührerische, daß sie in Verdacht bei dem Könige kämen, wenn sie die Stadt befestiget hätten, daß sie alsdann abfallen und andere mit ihnen zum Abfall bereden und sich stärken würden.
% Source scan 679; block decisions in selection.json

\quellseite{679}
Solcher Nachrede und Lügen müssen die Christen noch heutiges Tages gewarten und tragen um des Evangelii und des rechten Gottesdienstes willen. Denn alsbald sie davon nicht abstehen, noch etwas begeben wollen, wie sie denn ihres Gewissens und der Ehre Gottes halben nicht können, so heißt man sie ungehorsam, aufrührerische Leute. Aber es geschieht ihnen unrecht, wie den frommen Juden hie.

\section{Votum}

Der allmächtige Gott steure und wehre allen Feinden seiner Kirchen, stärke und erhalte uns in Verfolgung und unter allem Widerstand; und helfe, daß sein Wort bei uns geprediget und die rechten Gottesdienste geübet werden, bis an den jüngsten Tag. Amen.

\chapter{Das 5. Kapitel}

\section{Vorrede}

Der königliche Prophet David sagt im 55. Psalm V. 23.: Wirf dein Anliegen auf den Herrn, der wird dich versorgen und wird den Gerechten nicht ewiglich in Unruhe lassen. Hiermit lehret der heilige David: Kreuz, Trübsal, Angst und Noth soll und muß bei den Frommen sein: darein sollen sich auch alle Gottesfürchtige freiwillig und geduldig ergeben; sintemal es Gottes Wille ist, und weil wir des Herrn Christi Malzeichen und Sterben an unserm Leibe tragen, seines Leidens viel haben, mit Thränen säen, Angst und Leid haben müssen. Aber, wie Gott demüthiget, also erhöhet er auch, daß er tröstet durch sein kräftig Wort, und Kraft und Vermögen verleihet, daß wir die Noth überwinden können; und endlich hilft er auch unter und aus allem Kreuz, und thut solches, wenn es ihm wohlgefällt und nun die größte und äußerste Noth vorhanden ist. Dieser Lehre und Trost ist ein Exempel das 5. Kapitel.

\section{Fortsetzung des Tempelbaues}
% Source scan 680; block decisions in selection.json

\quellseite{680}

\section{Summarium}

Sie ist ein fein tröstlich Exempel, nachdem die Juden von den umliegenden Heiden sind verklagt und ihnen der Bau niedergelegt wird, will Gott sein Volk in solcher Anfechtung nicht gar ohne allen Trost lassen. Schicket zween Propheten, welche sie trösten und vermahnen sollten, an dem angefangenen Werk fortzufahren, V. 8. Item, er wendet des Landpflegers und anderer Amtleute Herz, daß sie dem Könige Dario anzeigen, daß ihnen von Cyro selbst zu bauen erlaubt sei. Also siehet man, daß Anfechtung und Widerwärtigkeit in göttlichen Sachen nicht außen bleiben und doch Gott allwege Hülfe und Trost schicket, und Mittel schaffet, daß sein Wille und nicht des Teufels, noch der argen Welt Wille geschehe.

\section{Votum}

Wir danken unserm treuen lieben Gott, daß er seine Kirche und Gemeine bisher so wunderlich erhält und fördert; der wolle an uns auch ferner diese seine Gnade beweisen, und allen verfolgten Christen Trost und Hülfe mittheilen und erzeigen, Amen.

\chapter{Das 6. Kapitel}

\section{Vorrede}

Wir sollen täglich und fleißig Gott anrufen, und beten für unsere Obrigkeit. Denn die Obrigkeit solches um uns verdienet und auch bedarf. Denn wie es bei Gott stehet und wir ihn darum bitten müssen, daß er guten fruchtbaren Regen der Erde geb, soll sie anders grünen und Früchte bringen: also muß er auch seinen Segen der Obrigkeit geben, soll sie nützlich und wohl regieren. Darum sollen wir ihn auch herzlich bitten: So hat es Gott auch befohlen, als dies St. Paulus 1 Tim. 2, 1. lehret, da er sagt: So ermahne ich, daß man für allen Dingen zuerst thue Bitte, Gebet, Fürbitte und Danksagung für alle Menschen, für die Könige und für alle Obrigkeit, auf daß wir ein geruhiges und stilles Leben führen mögen, in aller Gottseligkeit und Ehrbarkeit. Denn solches ist gut, dazu auch angenehm für Gott unsern Heiland. Ueber das sollen wir für die Obrigkeit beten, daß uns unser Gewissen nicht verklage. Welches geschehen würde, da wir der Obrigkeit die Ehre nicht thäten, die wir unsern Eltern schuldig sind, darunter die Fürbitte auch wird begriffen. Dieses lehret auch das 6. Kapitel.

\section{Vollführung und Einweihung des Tempels}
% Source scan 681; block decisions in selection.json

\quellseite{681}

\section{Summarium}

Ist eine schöne historia, wie der König Darius auf ein neues befehlet, man soll nicht allein die Juden an ihrem Bau ungehindert lassen, sondern auch von dem königlichen Einkommen dazu geben, was vonnöthen ist, V. 6. Solches ist ja eine feine fürstliche Tugend, wo die Obrigkeit sich des rechten Gottesdienstes annimmt und sich seiner Unkosten dauern läßt.

Sonderlich aber ist hie zu merken, daß der gottseelige König Darius begehrt, daß man soll zu Jerusalem für sein und seiner Kinder Leben bitten, und stellet eine harte Strafe und greulichen Fluch darauf, so jemand wider solchen seinen Befehl handeln, und die Juden hindern wollte.

\section{Votum}

Wir danken GOtt, der uns auch christliche Obrigkeit gibt und erhält; der segne, regiere, stärke und erhalte sie ferner, auf daß sie ihrem Regiment möge wohl fürstehen, daß er geehrt und die armen Unterthanen geschützt werden. Amen.

\chapter{Das 7. Capitel}

\section{Vorrede}

Die Obrigkeit soll ein Schutzherr sein beider Tafeln der heiligen Zehen Gebote, und neben der weltlichen Regierung die christliche Kirche treulich befördern. Darum wird dem römischen Kaiser in der Krönung nicht allein ein Scepter, sondern auch ein Reichsapfel, darauf ein Kreuz stehet, in die Hand gegeben, zur Anzeigung, daß die Regenten die heilsame Lehre von dem gekreuzigten Christo befördern, schützen und handhaben sollen: also hat ihm der Kaiser Carolus Magnus, wenn er in seiner Herrlichkeit herein gegangen, die Bibel und ein Schwert fürtragen lassen, zur Erinnerung, daß er und alle christliche Obrigkeit beide die christliche Kirche mit dem reinen Wort Gottes versorgen, und das Schwert im weltlichen Regiment zum Schutz der Frommen und zur Strafe der Bösen führen solle. So helfe der treue Gott, daß alle Obrigkeit ihm, als ihrem Herrn, diene mit Furcht, und daß sie küssen seinen Sohn, unsern Herrn Jesum Christum, daß er nicht zürne und sie umkommen auf dem Wege. Davon höret auch an das 7. Kapitel.

\section[Ezra bekommt Erlaubniß und Beförderung, den Gottesdienst zu]{Ezra bekommt Erlaubniß und Beförderung, den Gottesdienst zu bestellen}
% Source scan 682; block decisions in selection.json

% Scan 682: no apparatus retained; only biblical narrative/running matter.
% Source scan 683; block decisions in selection.json

\quellseite{683}

\section{Summarium}

Ist ein anderer schriftlicher Befehl, den der König Artahsastha Esra, dem Schriftgelehrten, gibt, V. 11., in welchem man stehet, wie dieser heidnische König so ein feines Herz und guten Willen zu dem Gott Israel und seinem Dienst hat, daß er nicht allein aus seiner Kammer verordnet, Silber und Gold darum zu kaufen, was zum Gottesdienst gehöret, sondern auch das Gewissen hat, wo er den Gottesdienst nicht also förderte, es möchte Gottes Zorn über sein Königreich und seine Kinder kommen, V. 23.

Es ist aber sonderlich hie zu merken, daß er die Priester und Leviten befreyet von allen Zinsen und Renten, daß man sie damit unbeschweret lasse, V. 24. Auch befiehlet er dem Esra, daß er sonderlich darauf soll Achtung geben, daß man rechte Prediger möge halten, und die, so noch nicht genugsam im Gesetz erfahren, recht unterweiset und gelehret würden, V. 25. Dies sind rechte fürstliche, ja bischöfliche Werke, dafür man billig (wie Esra hie thut, V. 27.) Gott danken soll, wo eine Obrigkeit sich dermaßen in die Sache schicket.

\section{Votum}

Wir danken GOtt herzlich auch für unsere liebe Obrigkeit; der regiere und erhalte sie, daß sie ferner die göttliche Lehre und treue Lehrer schütze und fördere, und belohne reichlich ihre väterliche Treue an ihnen selbst, an ihren Erben und Kindern und an ihren Hab und Gütern, Amen.

\chapter{Das 8. Capitel}

\section{Vorrede}

Die Vollträgigen und Trunkenen beten entweder gar nicht, oder beten ohne alle Andacht, und mit bösem Gewissen. Derwegen sollen wir uns für Fressen und Saufen hüten und ein nüchtern und mäßiges Leben führen, damit wir zum Gebet desto geschickter sein. Man schreibt vom römischen Regenten Romulus, daß derselbe in einer Collation wenig Weins getrunken, und da er gefragt ward, warum er so wenig getrunken? hat er darauf geantwortet: Er hätte eine wichtige Sache zu verrichten; also vielmehr sollen wir Christen, um des Gebets willen, daran uns trefflich viel gelegen, uns mäßig im Essen und Trinken halten. Daher setzet auch die Schrift gemeiniglich zusammen: Fasten und Beten. Als Matth. 7. 21. saget Christus: Dieses Geschlecht fähret nicht aus, denn durch Beten und Fasten. Und St. Petrus spricht 1. Epistel, 5. Kapitel V. 8.: Seid nüchtern und wachet. Und der alte Tobias sagt im 12. Kapitel V. 9.: Es ist ein köstlich Ding, das Gebet mit Fasten. Dieses lehret auch das 8. Kapitel.

\section{Erste Reise gen Jerusalem wird vollbracht}
% Source scan 684; block decisions in selection.json

\quellseite{684}

\section{Summarium}

Ist eine Historia, in welcher erzählet wird, welche mit Esra aus Babel gen Jerusalem gezogen seien. Und sonderlich ist hie zu merken, was es für eine Meinung habe, daß Esra, ehe er sich auf den Weg begibt, ein Fasten ausrufet, V. 21. Wie der König Josaphat 1 Chron. 20. auch thut, da die Moabiter ihm ins Land fielen. Denn wo eine gemeine Gefahr und Noth fürfiele, da geboten die Könige und Priester ein Fasten. Nicht wie die Papisten, daß sie damit Gott wollten dienen, sondern daß jedermann zum Gebet desto geschickter wäre, und sie sich also für Gott demüthigten und mit größerem Ernst bei Gott um Hülfe ansuchten. Denn wo ein voller Bauch ist, da wird wenig ernstliches Gebets sein. Darum Christus auch mehr denn einmal Fasten und Beten zusammen setzet. Von solchem Ernst des Gebets wissen unsere Papisten nichts, und wollen dennoch aus ihrem spöttlichen Fasten einen sonderlichen Gottesdienst machen.

\section{Votum}

Helfe der treue GOtt, daß wir der heiligen Exempel fleißig folgen, und GOtt Tag
% Source scan 685; block decisions in selection.json

\quellseite{685}
und Nacht und uns selbst dienen, mit nüchternem Leben und Gebet, und allen andern Tugenden, in Kraft des heiligen Geistes, Amen.

\chapter{Das 9. Kapitel}

\section{Vorrede}

Hieran ist alles gelegen, da wir selig werden wollen, daß uns Gott die Sünde vergebe, daß wir auch solches bitten und im Glauben annehmen müssen. Denn, so wir sagen: Wir haben keine Sünde! so verführen wir uns selbst, und die Wahrheit ist nicht in uns, sagt der heilige Apostel Johannes 1. Epistel im 1. Kapitel V. 8. So wird alle Sünde den Gläubigen vergeben um des Herrn Christ willen. Daher sagt St. Petrus in den Geschichten der Apostel im 10. Kapitel V. 43.: Von Christo zeugen alle Propheten, daß durch seinen Namen alle, die an ihn glauben, Vergebung der Sünde empfahen sollen. Denn Christus ist unser Erlöser, welcher bezahlet beide, den alten Rest der Erbsünde und alle andere wirkliche Sünde. Er hat uns eine gnädige Quitanz erlanget. Denn sein rosinfarben Blut ist durch unser Schuld-Buch durchgeflossen, zur Vergebung aller unserer Sünden. Es wird auch nimmermehr keine Rechnung in die himmlische Renterei eingenommen, wo sie nicht mit der schönen Rubrik des Blutes und Todes Jesu Christi zuvor cassirt und quittirt oder ausgelöscht ist. Davon lehret auch insgemein das 9. Kapitel.

\section{Etwas Beicht- und Bußgebet}

a Nagel und Zaun) ist geredt auf Sprichwortsweise, daß alles andere, Land und Leute umkommen sind, und sie, die Israeliten, noch übrig sind blieben, als ein Nagel vom Hause und ein Zaun vom Lande.
% Source scan 686; block decisions in selection.json

\quellseite{686}

\section{Summarium}

Dieß ist eine ernstliche Beichte und feines Gebet des frommen Esra, V. 6., daß GOtt die Sünde vergeben und nicht nach dem Verdienst strafen wolle, daß die Priester und Leviten heidnische Weiber genommen hatten, V. 12. Uns zum Exempel, daß wir nach gethaner Sünde dieselbigen für Gott von Herzen bekennen sollen, und uns nichts anders trösten, denn seiner Gnade und Barmherzigkeit und auf dieselbigen unser Gebet gründen.

\section{Votum}

Der treue barmherzige GOtt sei uns allen gnädig, vergebe uns unsere Sünde und lindere uns die Strafen, so wir mit unsern Sünden wohl verdienet haben; und helfe, daß wir uns allezeit seiner Gnade und unseres Herrn Christi Verdienstes mit Freuden trösten, Amen.

\chapter{Das 10. Kapitel}

\section{Vorrede}

Es gehören zu rechter Buße und Bekehrung zu Gott drei Stücke. Das erste ist herzliche Erkenntniß der Sünden und wahre Reue und Leid von wegen der Sünden. Davon lehret Gott im 66. Kapitel V. 2. Esaia also: Ich sehe an den Elenden und der zerbrochenes Geistes ist, und der sich fürchtet für meinem Wort. Das andere Stück der christlichen Buße ist der Glaube, daß Gottes Güte und Gnade viele tausend mal größer sei, denn aller Menschen Sünde und Uebertretung, und daß er allen, die ihn darum bitten, alle Sünden vergibt, um des Herrn Christi Verdienstes und Vorbitte willen. Dieses lehret David, Ps. 147, 11.: Der Herr hat Gefallen an denen, die ihn fürchten und auf seine Güte warten. Und St. Paulus sagt Col. 1, 14.: An Christo haben wir die Erlösung durch sein Blut, nämlich die Vergebung der Sünden. Das dritte Stück der wahren Buße ist neue Gehorsam nach Gottes Wort und Gebot, daß der Mensch von Sünde ablasse und in allen guten Werken sich übe, wie es ihm in seinem Stande gebühret, damit der siebenfältige Teufel nicht wieder einschleiche und das Letzte ärger werde, denn das Erste. Dies bestätiget St. Paulus an die Römer im 6. Kapitel V. 12.: Lasset die Sünde nicht herrschen in eurem sterblichen Leibe; ihr Gehorsam zu leisten in ihren Lüsten. Auch begebet nicht der Sünden eure Glieder zu Waffen der Ungerechtigkeit, sondern zu Waffen der Gerechtigkeit. Davon lehret auch das 10. und letzte Kapitel.

\section{Die fremden Weiber werden von den Israeliten abgeschaffet}
% Source scan 687; block decisions in selection.json

\quellseite{687}

\section{Summarium}

Es ist eine Historia, daß Esra daran ist, daß alle so heidnische Weiber genommen, ob sie gleich Kinder mit ihnen gezeuget hatten, dieselbigen sammt den Kindern von sich thun sollen, V. 3. Denn weil es wider Gottes Wort und Befehl geschehen war, konnte es keine rechte Ehe sein. Hie stehet man abermal, was rechtschaffene Buße ist, nämlich, daß man ihm die Sünde von Herzen lasse leid sein und davon ablasse, und hoffe darnach, Gott werde sie gnädiglich vergeben. Ohne Zweifel wird es wohl haben gethan, sich von Weib und Kind scheiden. Aber weil Got-
% Source scan 688; block decisions in selection.json

\quellseite{688}
tes Befehl da war, und ohne daß die Buße nicht konnte rechtschaffen sein, mußten sie es thun, oder zugleich Gottes Zorn und ein böses Gewissen tragen.

\section{Votum}

Helfe der gütige barmherzige GOtt, daß wir auch wahre Buße thun, in wahrer Buße alle Tage leben unser Lebenlang, an den Herrn Christum fest glauben, und unsern Glauben beständig beweisen mit guten Werken, gegen GOtt und gegen unsern Nächsten, nach seinen Geboten und Hülfe des heiligen Geistes, Amen.

\part{Das Buch Nehemia}

\chapter{Das 1. Kapitel}

\section{Vorrede}

Ein rechter Beter soll ein Mensch sein, der nicht nur habe Mund und Zunge, sondern auch heilige Hände, übersichtige Augen, ein gut Gewissen und ein gläubig Herz, das ist: ein rechter Beter soll sein Gebet nicht allein ein Mundwerk von Worten sein lassen, sondern er soll heilige Hände aufheben, ohne Zorn, Blut und Frevel. Des Beters Augen sollen nicht unter sich sehen, wie der Zweifler, die an Gott und seiner Gnade verzagen, sondern sollen über sich sehen gen Himmel, in fröhlicher Zuversicht auf den Namen Jesu Christi und göttliche Verheißung; der Beter soll haben ein gut Gewissen und gute Sache, daß die Person mit Gott versöhnet sei. Denn ein böses Gewissen kann nicht beten. Der Beter soll auch haben ein rechtschaffen Herz, brünstig in der Anrufung und gläubig auf Gottes Zusage und Verheißung. Denn wer also betet, daß er zweifelt an der Erhörung Gottes und allein sein Gebet dahin setzt auf Ebentheuer, der begehet zwei große Sünden: erstlich, daß er sein Gebet selbst zu Sünden und Schaden macht; zum andern, daß er seinen getreuesten und wahrhaftigsten Gott für einen Lügner acht, der seinem Wort und Zusage nicht könne oder wolle Folge thun. Derwegen wer wohl und nutzbarlich beten will, der schaue wohl zu, daß er ein rechter Beter sei. Dies betrachten wir auch bei dem 1. Kapitel.

\section{Nehemias Gebet für die Kinder Israel}
% Source scan 689; block decisions in selection.json

\quellseite{689}

\section{Summarium}

Dies ist ein schön Gebet, welches Nehemias thut, weil er der Juden Sache bei dem Könige handeln wollte, daß ihm Gott zu solchem Anbringen wollte Glück geben, V. 5. Und zum Exempel, daß wir nicht allein lernen, alle Dinge mit dem Gebet anfangen, auf daß wir desto mehr Glück haben, sondern daß wir auch sehen, wie ein recht Gebet müsse gestellet sein, nämlich, daß man sich erstlich für Gott demüthige und bekenne, daß wir Sünder und aller Strafe werth sind. Zum andern, daß man auf Gottes Verheißung und Zusage das Gebet stelle, wie Nehemias die Worte erzählet, in welchen Gott Mose zugesagt hat, er wolle auch der Sünder Gebet erhören und annehmen, wenn sie sich bekehren und sein Gebot halten.

\section{Votum}

Gott, der Brunnquell alles Guten, verleihe uns auch bußfertige, demüthige Herzen, und helfe, daß wir täglich fleißig beten, für unsere und der Christenheit gemeine Nothdurft, und daß wir in wahrem Glauben seiner väterlichen Hülfe erwarten, Amen.

\chapter{Das 2. Kapitel}

\section{Vorrede}

Der königliche Prophet David spricht im 145. Psalm V. 19.: Der HErr thut, was die Gottesfürchtigen begehren. Und der weise König Salomon in seinen Sprüchen sagt im 14. Kapitel V. 27.: Die Furcht des HErrn ist eine Quelle des Lebens, daß man meide die Stricke des Todes. Und Sirach im 1. Kapitel V. 18. spricht: Wer den HErrn fürchtet, dem wird es wohl gehen, und wenn er Trost bedarf, wird er gesegnet. Und St. Petrus sagt in den Geschichten der Apostel im 10. Kapitel V. 35.: Wer den HErrn fürchtet und recht thut, der ist ihm angenehm. Dagegen spricht Salomon in seinen Sprüchen im 28. Kapitel V. 26.: Wer sich auf sein Herz verläßt, ist ein Narr. Und der Prophet Jeremias im 17. Kapitel V. 5. sagt: Verflucht ist der Mann, der sich auf Menschen verläßt und hält Fleisch für seinen Arm, und mit seinem Herzen vom HErrn weicht. Also lehret uns Gottes Wort, daß wir alle in Gottesfurcht leben und in allem Thun auf Gott allein vertrauen, und für Vermessenheit uns hüten sollen; und daß uns also werde Glück und Segen folgen. Davon lehret auch das 2. Kapitel.

\section{Nehemia erlangt Gewalt und Briefe, Jerusalem zu bauen}
% Source scan 690; block decisions in selection.json

\quellseite{690}

\section{Summarium}

Ist eine feine Historia, wie Nehemias Erlaubniß bei dem Könige Arthasasta erlanget, daß er gen Jerusalem ziehen und die Mauern und Thore bauen möge, V. 5. Auf daß wir sehen, wenn man in der Furcht Gottes und mit einem ernstlichen Gebet etwas anfähet, wie Nehemias droben gethan hat, daß Glück dabei ist. Wie dagegen Unglück sein muß, wo man ohne Gottesfurcht und aus Vermessenheit etwas anfähet, wie doch gemeiniglich jedermann thut.

\section{Votum}

Regiere uns GOtt durch seinen heiligen Geist, daß wir nichts thun aus Vermessenheit, sondern daß wir in seiner Furcht alles anfangen und vollbringen, was ein jeder in seinem Beruf schuldig ist, und lasse die Werke unserer Hände wohl gerathen, Amen.

\chapter{Das 3. Capitel}

\section{Der Bau zu Jerusalem gehet an}

a (Gewaltigen) Die Armen müssen das Kreuz tragen, die Reichen geben nichts.
% Source scan 691; block decisions in selection.json

\quellseite{691}

\chapter{Das 4. Capitel}

\section{Vorrede}

Sollen die Feinde göttlichen Wortes und Namens gefällt und erlegt werden, so müssen wir, ein jeder in seinem Beruf und Stande, mit ernstem Seufzen in unserm Vortheil des lieben Gebets liegen und Gott zum Freunde haben und behalten. Ursache: unsre Feinde alle sind nichts anders, denn eine zornige Ruthe unsers Gottes und HErrn, welche er über uns verhänget um unserer Sünde willen, welche wir anders nicht brechen werden, denn mit Büßen und Beten. Derwegen wird recht gesagt: Wider die Strafen ist eine Arznei, die wahre christliche Buße, daß sie dennoch gelindert werden, wo nicht gar abgewendet; und das Gebet der Christen ist der Rüstzeug, Wehre und Waffen wider die Feinde, und für die Freunde der Christenheit. Denn wer mit rechtem Ernst betet, der streitet auf seiner Seite für Freunde, auf der andern Seite wider die Feinde der Christenheit. So helfe Gott, daß wir heute auch also leben und beten, daß es ihm zu Ehren, und zu Nutz und unsern Feinden zum Untergange gelange, Amen. Das von höret auch an das 4. Kapitel.
% Source scan 692; block decisions in selection.json

\quellseite{692}

\section{Der Bau gehet auch unter den Hindernissen fort}

\section{Summarium}

Ist eine Historia, in welcher man siehet, daß der Teufel den Juden sonderlich feind gewesen und ihnen durch ihre Nachbarn alles Unglück zugerichtet hat, daß sie mit einer Hand arbeiten, mit der andern sich der Feinde haben wehren müssen, V. 16. 18. In solcher Noth hält sich Nehemias zu dem Gebet und tröstet sein Volk, Gott werde bei ihnen sein und der Feinde Rath zu nichte machen. Diese Historia ist ein fein Fürbild unsers Lebens hie auf Erden, da hat man immerdar wider den Teufel, die Welt und unser eigen Fleisch zu fechten, soll anders die
% Source scan 693; block decisions in selection.json

\quellseite{693}
christliche Kirche durch reine Lehre und gottseeligen Wandel recht erbauet werden.

Helfe und regiere uns alle der treue GOtt, daß wir in fürstehenden Nöthen in täglicher Buße leben, in unserm Beruf treu sein und fleißig beten; und beschütze uns auch gnädiglich wider den Türken, und alle andere unsere Feinde, Amen.

\chapter{Das 5. Kapitel}

\section{Vorrede}

Der Wucher ist wider Gott und hat den zum höchsten Zorn und manchfaltigen Strafen bisher verursacht. Der Wucher ist eine Schande und böser Fürwurf unserm heiligen Evangelio. Der Wucher ist ein Hinderniß aller Ehrbarkeit und guter Polizei-Ordnung. Der Wucher ist ein Aergerniß aller frommen Herzen. Der Wucher ist eine Zerrüttung aller göttlichen und christlichen Nahrung. Der Wucher ist ein Gift der christlichen brüderlichen Liebe. Der Wucher ist ein Stifter aller Abgötterei. Der Wucher ist ein gewisses Pulver mit Schwefel anzuzünden, das ewig brennende Feuer im Gewissen. Wohl nun zeitlich, wohl nun ewig allen, die den Wucher unterlassen. Denn denen will Gott geben allerlei leiblichen und zeitlichen Segen, an Leib und Seele, Ehre und Gut, an ihren Kindern und Nachkommen. Ja es soll ihnen auch folgen der ewige himmlische Segen, ewiges Leben und Seligkeit, als bei denen, die ihren allein seligmachenden Glauben in rechtschaffenen Früchten der Liebe gegen den Nächsten sehen lassen. Wer Ohren hat zu hören, der höre. Das soll auch betrachtet werden bei dem 5. Kapitel.

\section{Nehemia thut die Beschwerden des Volks ab}
% Source scan 694; block decisions in selection.json

\quellseite{694}

\section{Summarium}

Ist eine feine Historia, darin Nehemias gebeut, daß man in der Theurung den armen Mann mit Wucher nicht übersezen, sondern was man ihnen abgewuchert habe, wieder geben und ihnen umsonst fürstrecken und leihen soll, um der Furcht Gottes willen, V. 11. Solches wäre heutiges Tages ein sehr nöthiges Exempel, daß die weltliche Obrigkeit auch also thäte und den Reichen wehrete, daß der arme Mann, sonderlich wenn theure Zeit ist, von ihnen nicht beschweret würde. Zum andern ist ein Exempel: nachdem Nehemias zum Landpfleger verordnet ist, er solches Amt also verwaltet, daß seine Leute mit unnöthigen Unkosten nicht beschweret werden, V. 14. Zum Exempel aller Obrigkeit, daß sie dergleichen mit ihren Unterthanen auch so handeln, und mehr auf derselben Nutz, Zunehmen und Gedeihen, denn auf den eigenen Nutz sehen. Aber man erfähret leider das Widerspiel.

\section{Votum}

Gott, unser lieber Herr und Vater, regiere unsere Obrigkeit, daß sie ihres Amts treulich pflege und der Unterthanen Wohlfahrt fleißig fördere; und helfe, daß alle Wucherer und Geizhälse heut Buße thun, der zeitlichen und ewigen Strafe zu entfliehen, Amen.

\chapter{Das 6. Kapitel}

\section{Vorrede}

Afterreden, den Nächsten verleumden, belügen und verrathen ist leider gemein worden und hat als eine pestilenzische Gift und vergifteter Wurm alle Winkel durchkrochen, und ist keine Kirche so wohl bestellet, keine Festung so wohl verwahret, keine Stadt, Dorf, noch Haushaltung dazu so wohl versehen, da dieser Wurm nicht eingeschlossen wäre. Derwegen, erstlich hüte man sich männiglich für solchem Laster und rechten Malzeichen des Teufels. Denn wer seinen Nächsten heimlich verleumdet, den vertilge ich, sagt Gott, der Herr, im 101. Psalm V. 5. Und St. Paulus spricht 1 Cor. 6, 24.: Die Lästerer werden das Reich Gottes nicht ererben. Dagegen sagt David im 15. Psalm: Herr, wer wird wohnen in deiner Hütte? Wer wird bleiben auf deinem heiligen Berge? Wer die Wahrheit von Herzen redet, und wer mit seiner Zunge nicht verleumdet. Und Esaias im 33. Kapitel V. 15. spricht: Wer in Gerechtigkeit wandelt, und redet, was recht ist. Wer Unrecht hasset, samt dem Geiz, der wird in der Höhe wohnen. Fürs andere, widerspreche auch ein jedes dem Lügen und Afterreden, wenn es solches von Verleumdern und Ehrendieben müßte hören. Denn, wie Salomon sagt in seinen Sprüchen im 17. Kapitel V. 15.: Wer den Gottlosen Recht spricht, und den Gerechten verdammt, die sind beide dem Herrn ein Greuel. Dies erinnert auch das 6. Kapitel.

\section[Nehemia treibet durch seine Vorsichtigkeit die Hindernisse am]{Nehemia treibet durch seine Vorsichtigkeit die Hindernisse am Werke zurück}
% Source scan 695; block decisions in selection.json

\quellseite{695}

\section{Summarium}

Dies ist eine Historia, wie die heidnischen Amtleute der Hoffart nach sich an Nehemia versuchen, ob sie ihn mit Lügen und anderm könnten verzagt machen und abschrecken, V. 2., bestellen auch etliche falsche Propheten wider ihn, V. 14. Solche Historia zeuget ehrlich, wie mancherlei Gefahr der fromme Nehemias auf sich hat müssen nehmen, darum, daß er treulich mit Jerusalem und den Juden gemeinet hat. Aber weil er Gott fürchtet und ihm vertraut, verachtet er solche Gefahr und bleibet fest bei dem, daß er angefangen hatte. Zum Exempel aller Obrigkeit, daß sie ihres Amts mit Fleiß aus warten soll, und gar nichts darnach fragen, es lache die Welt oder sehe sauer dazu, wenn es nur Gott gefällt und mit gutem Gewissen geschiehet.

\section{Votum}

Der treue Gott verleihe auch ferner unser lieben Obrigkeit starken Muth und andere nöthige Gaben zu glückseliger Regierung, und steure und wehre allen ihren Feinden, Amen.

\chapter{Das 7. Kapitel}

\section{Wächter der Stadt; Zahl und Geschenke des Volks}
% Source scan 696; block decisions in selection.json

% Scan 696: no apparatus retained; only biblical narrative/running matter.
% Source scan 697; block decisions in selection.json

\quellseite{697}

\chapter{Das 8. Kapitel}

\section{Vorrede}

Die zwei Testamente in der heiligen Biblia, sagt St. Hieronymus, sind die zwei Mutter-Brüste der Kirche, dadurch sie ihre KirchKinder säuget und ernähret. Derwegen ist das Alte Testament in hebräischer Sprache von den Propheten geschrieben worden, dieweil es der Juden Muttersprache und jedermann bei ihnen verständlich war; das Neue Testament aber ist in griechischer Sprache gestellet, dieweil dieselbe der Zeit die allgemeinste Landsprache und fast an denen Orten allen, in welchen und an welche die Bücher daselbst geschrieben worden, bräuchlich gewesen. Auch ist das Alte Testament in die griechische, und nachmals, sammt dem Neuen, von vielen Gelehrten, und sonderlich von St. Hieronymo in die lateinische und slavonische Sprache übersetzt worden, und darauf in die deutsche vom Herrn Luther, Herrn Philippo und ihren Mitgehülfen. So danken wir heute abermal herzlich Gott, unserm lieben Vater, daß wir die heilige Biblia haben, und daß uns diese alle Tage fürgelesen wird. Derselbe helfe gnädiglich, daß wir die heilige Biblia allzeit anhören, aus herzlicher Begierde, seinem Willen und die Wahrheit zu lernen, und mit ernstlichem Fürsatz, derselben zu glauben und darnach zu leben. Darauf höret ferner an das 8. Kapitel.

\section{Nehemia und Esra richten den Gottesdienst an}
% Source scan 698; block decisions in selection.json

\quellseite{698}

\section{Summarium}

Ist eine historia, wie Nehemias verordnet, daß man predigen und dem Volk das Gesetz fürlesen soll. Denn wo das Predigtamt nicht im Schwange gehet, da ists unmöglich, daß nicht Irrthum und falscher Gottesdienst sollte folgen. Wiederum, wo Gottes Wort ist, da werden die Herzen recht unterrichtet, daß sie ihre Sünde und Gottes Barmherzigkeit erkennen lernen, und thun nicht, was sie gut dünket, sondern was sie hören, das Gott zu thun befohlen hat. Wie sie hie das Laubhüttenfest und anders anrichten, V. 16. Um solche Wohlthat soll man Gott loben und danken, wie der Priester Esra hie auch thut.

\section{Votum}

Das ewige wesentliche Wort GOttes, unser HErr JEsus Christus, erhalte uns die heilige Biblia; und verleihe, daß die in unseren Kirchen bis an den jüngsten Tag, zu Stärkung des Glaubens, zu Besserung des Lebens, und zu Trost in Noth und Tod, allen Zuhörern vorgelesen werde, Amen.

\chapter{Das 9. Capitel}

\section{Vorrede}

Eines jeden Christen Leben soll eine tägliche Buße sein. Denn wir sündigen alle leider alle Tage. So ist auch die Stunde des Todes sehr unsicher und ungewiß, wenn, wie und wo eines von dieser Welt abscheiden soll; ja niemand hat nur auf eine Stunde Siegel und Brief, dieselbige zu überleben. So sagen sehr fein die alten Lehrer: Durch die Buße wird der Mensch entbunden von seinen Sünden; er wird der christlichen Kirche wieder eingeleibet, mit Gott versöhnet, in ihm werden die geistlichen Gaben gemehret, aus einem Teufels-Kind wird ein Gottes-Kind, das Paradies wird ihm geöffnet. Ja sie sagen ferner: Die Buße erfreut die Engel,
% Source scan 699; block decisions in selection.json

\quellseite{699}
sie macht gerecht den Gottlosen, sie bringt wieder die verlornen Güter, sie schenket das ewige Leben, sie behütet für der ewigen Strafe. Ach so wollen wir alle Tage unsere Sünde erkennen, um Vergebung in wahrem Glauben bitten und denn beständig durch Gottes Hülfe in Gottesfurcht wandeln und also in täglicher Buße leben. Davon lehret auch das 9. Kapitel.

\section{Des Volcks öffentliche Buße}

\section{Summarium}

Hie ist eine Historia, wie das Volck zusammen kommen, eine Fasten gehalten, V. 1., und sich auf ein neues mit Gott verbunden und die Leviten eine schöne lange Predigt gethan haben. In welcher sie für dem Volke ehrlich Gott loben und danken für alle seine Wohlthaten, die er sonderlich den Juden je und je hat widerfahren
% Source scan 700; block decisions in selection.json

\quellseite{700}
lassen. Darnach bekennen sie, wie sie sich gegen solchem gnädigen Gott (leider) allweg ungebührlich gehalten, seines Wortes nicht geachtet und seine Propheten erwürget haben, V. 16. Darum was ihnen für Unglück und Strafe begegnet sei, das haben sie alles wohl verdienet, V. 30. Auf solche Beichte folget hernach ein Gebet, daß Gott sie wieder zu Gnaden annehmen, alles Elend wenden und ihnen gnädiglich helfen wolle.

Auf diese Weise sollen wir auch beten lernen, daß wir Gottes Güte preisen, unsere Untugend und Sünde für ihm bekennen und beichten und weiter seiner Gnade warten, und uns bessern.

a Himmels Brod.
% Source scan 701; block decisions in selection.json

\quellseite{701}

\section{Summarium}

Mit eben dem vorhergehenden Summarium.

\section{Votum}

Helffe GOtt, daß wir alle, jung und alt, reich und arm, diese nöthige und selige Betekunst wohl lernen, täglich fleißig und fruchtbarlich üben, und GOtt auch dienen in allen Ständen, in Heiligkeit und Gerechtigkeit, Amen.

\chapter{Das 10. Kapitel}

\section{Vorrede}

Der HErr Christus verbindet sich gegen uns in der Schrift, daß er die Almosen reichlich bezahlen wolle, zeitlich, geistlich und ewig. Zeitlich will Gott die Almosen belohnen. Denn also sagt Christus, Lucas im 6. Kapitel V. 38.: Gebet, so wird euch gegeben. Ein voll, gedrückt, gerüttelt und überflüssig Maß wird man in euren Schooß geben; denn eben mit dem Maß, da ihr mit messet, wird man euch wieder messen. Geistlich will Gott die Almosen belohnen. Denn also spricht abermal Christus, Matth. 5, 7.: Selig sind die Barmherzigen, denn sie werden Barmherzigkeit erlangen. Ewiglich will Gott die Almosen belohnen. Denn also saget Christus wiederum Lucä 16, 9.: Machet euch Freunde mit dem ungerechten Mammon, auf daß, wenn ihr nun darbet, sie euch aufnehmen in die ewigen Hütten. Das sind wahrlich treffliche Zusagungen. Aber sehet, wie ein wunder Ding ist es. Wenn euch ein Mensch, der doch von Art ein Lügner ist, etwas zusaget, so könnt ihr ihm glauben, ihm leihen und borgen, der doch keine Stunde seines Lebens und Gutes sicher ist. Warum thut ihr solches auch nicht gegen eurem HErrn Christo? Der wahrhaftig und ewig ist, der nicht schlechte Zusagung gethan hat, sondern ist auch darauf gestorben zu genugsamer Versicherung. Sonderlich aber soll man zur Unterhaltung des Gottesdienstes gern und treulich helfen, und solches will euch auch der HErr Christus treulich vergelten, zeitlich, geistlich und ewig. Davon höret auch an das 10. Kapitel.

\section{Versiegelung des erneuerten Bundes}
% Source scan 702; block decisions in selection.json

\quellseite{702}

\section{Summarium}

Dies ist eine Historia, wie das Volk auf ein neues mit Gott sich verbunden habe, daß sie fort nicht mehr gottlos sein, sondern nach dem Willen Gottes leben wollen, V. 33. Und ist hie sonderlich zu merken, daß sie sich neben andern auch des verpflichteten, sie wollen jährlich ein jeder zum Tempel und Unterhaltung des Gottesdienstes ein Drittteil eines Orts vom Gulden geben, V. 29. Denn Geld und Gut kann man nicht besser und seliger anlegen, denn so man damit den rechten Gottesdienst fördert.

\section{Votum}

Der Vater aller Barmherzigkeit, unser lieber GOtt, mehre und erhalte allen ihre Güter, die solche zu seiner Ehre gebrauchen, und damit seine Gottesdienste erhalten helfen, und tröste und bewahre sie zum ewigen Leben in der letzten Noth, Amen.

\chapter{Das 11. Capitel}

\section[Register der neuen Einwohner in der Stadt und]{Register der neuen Einwohner in der Stadt und auf dem Felde}
% Source scan 703; block decisions in selection.json

\quellseite{703}

\chapter{Das 12. Capitel}

\section{Vorrede}

Im 148. Psalm V. 12. sagt der heilige Prophet und König David: Jünglinge und Jungfrauen, Alte mit den Jungen, sollen loben den Namen des HErrn. Zuvor auch im 117. Psalm V. 1. vermahnet er: Lobet den HErrn alle Heiden; preiset ihn alle Völker. Denn seine Gnade und Wahrheit waltet über uns in Ewigkeit. Halleluja. St. Paulus aber lehret an die Galater im 6. Kapitel V. 6. also: Der unterrichtet wird mit dem Wort, der theile mit allerlei Gutes dem, der ihn unterrichtet. Und Sirach im 7. Kapitel V. 33. saget auch also davon: Ehre den Priester, und gib ihnen ihr Theil, wie dir geboten ist. Derwegen soll sich ein jeder Christ mit höchstem Fleiß dankbar erzeigen, gegen Gott wegen seiner Wohlthaten und dann gegen den Dienern göttlichen Wortes, wegen ihrer Mühe, Arbeit und Fleiß, so sie an dieselbe wenden, zu ihrer Bekehrung und Seligkeit, in Christo Jesu unserm HErrn. Von diesen beiden erinnert auch das 12. Kapitel.

\section[Namen der Priester und Leviten. Einweihung der Stadt.]{Namen der Priester und Leviten. Einweihung der Stadt. Ordnung der heiligen Pfleger}
% Source scan 704; block decisions in selection.json

\quellseite{704}

a) Jaddua dies ist Jaddua, der dem großen Alexander zu Jerusalem empfing. Also hat Nehemia, vielleicht Esra auch, den großen Alexander erlebt, und ist sehr alt worden, und vier oder fünf Hoherpriester überlebt.
% Source scan 705; block decisions in selection.json

\quellseite{705}

\section{Summarium}

Ist eine historia, wie das Volk Gott danket und fröhlich ist, daß die Mauren an der Stadt wieder aufgebauet und Jerusalem ein wenig gebessert ist, V. 27.

Der Text meldet hie wiederum, wie man sich sonderlich darum habe angenommen, daß die Priester und Leviten mit der Nahrung versorget würden, daß sie der Kirchen desto fleißiger warten könnten, V. 24. Denn soll der Gottesdienst recht erhalten werden, so muß man etwas darauf wagen, daß tüchtige Personen erzogen und nach Nothdurft versorget werden.

\section{Votum}

Wir danken GOtt herzlich für alle Wohlthaten, die er unserer Stadt geistlich und leiblich erzeiget hat; der helfe, daß wir ihn und seine Diener ferner ehren, und daß solches auch alle unsere Nachkommen thun mögen, bis auf den jüngsten Tag, Amen.

\chapter{Das 13. Capitel}

\section{Vorrede}

Wie allezeit böse verkehrte Leute gewesen, die in der christlichen Kirche Aergerniß angerichtet haben: also findet man solche Leute noch in allen christlichen Gemeinen und Sammlungen, und wohl solche, die feindseliger in Haß und Widerwillen, als die Heiden, ärger in Unzucht, denn die Türken, tiefer im Wucher, denn die Juden leben. So sollen nun die Christen von der Kirche Gottes nicht abfallen oder sich absondern, solchem Aergerniß wegen, sondern sie sollen ihre Herzen verwahren mit rechter Lehre und fleißigem Gebet, und also alle Aergernisse überwinden und Gliedmaßen und Bürger der Kirche und Gemeine des Herrn Christi bleiben. Wie St. Paulus vermahnet 1 Cor. 16, 13.: Wachet, stehet im Glauben, seid männlich und stark. Und unser Herr Jesus Christus lehret dieses, da er Matth. 18, 8. also sagt: So deine Hand, oder dein Fuß dich ärgert, so haue ihn ab und wirf ihn von dir. Das ist, so dich irgend ein Mensch, er sei Freund oder Feind, durch seinen sündlichen Wandel, von Gott und von der Gottseligkeit abziehen will, so sollst du ihm nicht folgen, sondern widerstreben und fliehen. Denn durch das Wort Hand verstehet der Herr Christus unsere Freunde, welcher Freundschaft uns nicht weniger vonnöthen ist, als unsere rechte Hand. Fuß aber heißt so viel als der Wandel und Thun des Menschen. Darauf höret nun an das 13. und letzte Kapitel.

\section[Nehemias Eifer in Verbesserung der Policei, des Kirchen-]{Nehemias Eifer in Verbesserung der Policei, des Kirchen- und des häuslichen Wesens}
% Source scan 706; block decisions in selection.json

\quellseite{706}

\section{Summarium}

Das sind zwei Historien. Die erste, weil im Gesetze verboten war, daß kein Ammoniter sollte in die Gemeine Gottes kommen, will Nehemias nicht gestatten, daß Tobias, der Ammoniter, einen Kasten im Tempel habe, für sein und seines Hauses Opfer, wirft solches alles aus dem Tempel, lässet sichs nicht anfechten, daß der Gottesdienst desselbigen Gelds mangeln soll, sondern reiniget den Kasten.

Die andere Historie ist, daß die Juden bald nachließen, die Zehenten und Opfer zu geben, V. 10. Daraus konnte Nehemias leichtlich abnehmen, daß solches am Gottesdienst einen Nachtheil würde bringen. Denn wer seiner Nahrung warten und nachgehen muß, der kann der Kirchen nicht warten. Darum straft er solche Sünde, und verordnet von neuem, daß den Priestern und Leviten ihr Einkommen sollte gereicht werden, V. 11. Solches wäre sehr noth, daß es in allen Dörfern und Flecken von der Obrigkeit auch also bestellet würde.
% Source scan 707; block decisions in selection.json

\quellseite{707}

\section{Summarium}

Ist ein sehr fein Exempel, daß ein Regent allenthalben ein fleißig Aufsehen soll haben, daß die Leute nicht unrecht thun, und sich versündigen. Wie Nehemias mit dem Sabbath thut und mit denen, so heidnische Weiber nahmen.

Daß er nun bittet, Gott wolle an solche und dergleichen andere gute Werke denken, V. 22., ist recht gethan. Denn eben wie man um der Furcht Gottes willen das Uebel unterlassen soll: also soll man gute Werke darum thun, daß Gott verheißen hat, dieselbigen reichlich zu vergelten. Solche Hoffnung auf Gottes Verheißung gefällt Gott wohl, und hindert das Hauptstück christlicher Lehre gar nichts. Denn Nehemias bittet nicht, daß Gott um solcher guter Werke willen seiner wolle schonen, und ihn selig machen, sondern also spricht er: Herr, mein Gott, gedenke mir deß auch und schone mein, nach deiner Barmherzigkeit. Das heißt, nicht auf eigen Werk und Verdienst, sondern auf Gottes Güte und Gnade gebaut.

\section{Votum}

Verleihe Gott, der HErr, daß unsere liebe Obrigkeit auch ferner über der wahren Religion fleißig halte; und helfe, daß wir Zuhörer das Predigtamt auch treulich fördern, und segne, zeitlich und ewig, an Leib und Seele, uns alle, Amen.

\part{Das Buch Esther}

\chapter{Das 1. Kapitel}

\section{Vorrede}

Die Historia Esther ist eine schöne liebliche Geschichte, darinnen beide, hohe und niedere Personen, Obrigkeit und Unterthanen sich spiegeln können, und daraus lernen, wie sie sich für Unglück hüten, das Kreuz, so uns Gott aufgelegt, mit Geduld ohne Murren tragen, den Glauben ohne Zweifel mit ernstigem Gebet und Hoffnung zu Gott üben, und in christlichem untadeligen Leben fortfahren sollen. Auch ist darinnen schön und ganz artig abgemalet das Wesen und der Lauf der ganzen Welt, daß man sehen mag, wie es noch heutiges Tages in allen Ständen zugehe. Welches denn nicht allein nützlich, sondern auch nöthig ist zu wissen; sonderlich weil sich die Menschen mehr nach Exempeln, denn nach bloßen Worten, Strafen, Dräuen oder Vermahnungen richten. Solche Exempel aber sollen gebraucht werden zur Stärkung des Glaubens und Besserung des Lebens. Denn wer Gott vertrauet und auf ihn bauet, wird gleichermaßen, wie andere Gläubigen, von Gott aus allen Nöthen erlöset werden. Und dagegen werden

die, welche mit den Gottlosen in gleichen Sünden leben und verharren, mit gleicher Strafe auch heimgesucht werden. Denn das ist gewisser denn gewiß, daß der gerechte Gott keine Sünde ungestraft, und daß der barmherzige Gott kein gut Werk unbelohnet läßt, nach dem Gesetze zu reden. Dies ist die Summa des Buches Esther.

\section{Vasthi wird wegen verachteter Mahlzeit von Ahasveros verstoßen}
% Source scan 708; block decisions in selection.json

\quellseite{708}

\section{Summarium}

In diesem Buch ist eine Historia, wie Gott den Juden, so im Lande unter den Heiden waren, eine sonderliche Gnade bewiesen und sie für einem großen greulichen Jammer behütet habe.

Im 1. Capitel ist eine Historia, wie es aus Gottes Gericht wunderbärlich sich zugetragen habe, daß die ungehorsame Königin Vasthi entsetzet worden ist; auf daß Esther eingesetzt, und durch sie die Juden, so noch in Babel und sonst zerstreuet waren, errettet würden. Wie wir hernach in der Historia weiter werden hören.

Sonderlich aber ist hie zu merken, daß der Ungehorsam der Weiber gegen ihre Männer eine Sünde sei. Denn Gott hat dem Weibe um der Sünde willen die Strafe aufgelegt, daß sie unter des Mannes Gewalt leben und nicht ihr selbst Herr soll sein, wie Moses meldet, 1 Mos. 3, 9.
% Source scan 709; block decisions in selection.json

\quellseite{709}

\section{Votum}

Dem frommen lieben GOtt sei Lob und Dank gesagt, der seine Kirche wunderbarlich in Gefahr und Noth allzeit erhält; der verleihe, daß alle Eheweiber in schuldiger Liebe, Demuth und Gehorsam, gegen GOtt und ihren Ehemännern, beständig leben, Amen.

\chapter{Das 2. Kapitel}

\section{Vorrede}

Gott ist ein Vater der Armen, ein Tröster und Helfer der Waisen und der Wittwen-Richter. Daher sagt David im 146. Psalm V. 9.: Der Herr behütet Fremdlinge und Waisen, und erhält die Wittwen; und kehret zurück den Weg der Gottlosen. Darnach: Gott verhänget auch über Untreue und Verrätherei, daß sie offenbar werden und ihre verdiente Strafe endlich überkommen. Daher vermahnet der Prediger Salomonis im 10. Kapitel und sagt V. 20.: Fluche dem Könige nicht in deinem Herzen, und fluche dem Reichen nicht in deiner Schlafkammer, denn die Vögel des Himmels führen die Stimme und die Fittige haben, sagen's nach. Beides lehret auch das 2. Kapitel.

\section{Esther wird zu königlichen Ehren erhoben}
% Source scan 710; block decisions in selection.json

\quellseite{710}

\section{Summarium}

Ist eine feine Historia, wie es der frommen Esther glücklich gegangen und wohl gerathen sei, V. 8. Zum Exempel allen denen, so in Gottesfurcht wandeln und fromm sind. Denn ob sie wohl im Elend und allerlei Unglück sind, so will dennoch endlich Gott sie nicht lassen, sondern mit reichem Segen bei ihnen sein.

Am Ende ist eine Historia von den zweien Aufrührern, so den König Ahasveros ermorden wollten, und darum beide an Bäume gehängt worden, V. 21. Solch ihren Rath zeiget Mardachai Esther an, durch welche es wieder der König erfähret, V. 22. Solche Treue Mardachai wird hernach in der Historia ein wunderliches Urtheil verursachen.

\section{Votum}

Gott, der Herr, der das Gute treulich belohnet und das Böse ernstlich strafet, segne und behüte heut auch alle, die ihn fürchten, und steure allen Aufrührern, beide im geistlichen und weltlichen Regiment, und thue solches um der Ehre seines heiligen Namens willen, Amen.

\chapter{Das 3. Kapitel}

\section{Vorrede}

Gott hat die hohe Obrigkeit zum Ausschuß erwählt und ihre Gewaltigen ihnen zu Mittelsherren verordnet; doch daß sie alle ihre Anschläge und Fürnehmen, Thun und Lassen, nach seinem Wort, ihm zu Ehren und nach der christlichen Liebe ihren Unterthanen und Vertrauten zu Nutz, anstellen und verrichten sollen. Aber die fürnehmste Ursache, warum Gott Könige, Fürsten und andere Obrigkeit in der Welt verordnet hat, und warum er manchen so hoch anbringt, bei großen Potentaten und Herrn, ist, daß durch ihr Aufsehen, Fleiß und Sorge, Mühe und Arbeit, Friede auf Erden erhalten werde, daß nicht durch Krieg und Aufruhr, durch Rotten und Secten, christliche Religion zerstöret, Gottes Wort und Dienst niedergelegt, Zucht und Ehrbarkeit aufgehoben und die häusliche Nahrung verderbet oder verhindert werde. Aber viele mißbrauchen ihre Gewalt und Würde, auch wider Gott und seine christliche Gemeine, oder den Frommen und Unschuldigen zum Verderben und Untergang, so viel an ihnen ist. Als davon eine Historia meldet das 3. Kapitel.

\section{Hamans Ehrgeiz und Blutdürstigkeit}
% Source scan 711; block decisions in selection.json

\quellseite{711}

\section{Summarium}

Ist eine Historia, in welcher man ein sonderlich Bild siehet eines gottlosen Tyrannen. Haman kommt bei dem Könige zu großen Gnaden und Ehren und wird gewaltig, V. 1. Solche Gewalt macht ihn stolz und hoffärtig. Darum, daß ihn Mardachai nicht wollte anbeten, wie andere thaten, ward er ihm von Herzen gram, V. 5., und gedachte, wie die Tyrannen pflegen, seine Gewalt nicht allein wider ihn, sondern auch wider alle Juden zu üben, V. 7. Und bringets auch so weit, daß er bei dem Könige ein Gebot ausbringet, daß man alle Juden in allen Landen auf einen Tag erwürgen, ihr Gut rauben und Preis machen soll. V. 12. Da siehet man, wie der Teufel allewege Gottes Volk und der Kirchen so feind gewesen und die Tyrannen wider sie gereizet hat. Aber wie hernach folgen wird, will dennoch Gott dem Teufel seine Anschläge brechen und die Seinen erretten.

\section{Votum}

Gott, der Allmächtige, steure allen Tyrannen und tyrannischen Anschlägen und Fürnehmen wider seine Christenheit; und erhalte seine Kirche und deren Gliedmaßen, zu Nutz und Freude dem menschlichen Geschlechte, um Jesu Christi, unsers HErrn willen, Amen.

\chapter{Das 4. Kapitel}

\section{Vorrede}

Der heilige Apostel Johannes sagt 1. Epistel, 3. Kapitel V. 16.: Daran haben wir erkannt die Liebe, daß Christus sein Leben für uns gelassen hat, und wir sollen auch das Leben für die Brüder lassen. Und der HErr Christus sagt Matth. 10, 39.: Wer sein Leben findet, der wirds verlieren; und wer sein Leben verlieret um meinetwillen, der wirds finden. Hiermit werden wir gelehret und vermahnet, daß ein jeder viel lieber sein Leben in Gefahr stellen solle, denn auch einen einigen Christen in Angst und Noth stecken lassen. Deß ist ein Exempel Mosis, 2 Mos. 32, 32. Welcher begehrt, aus dem Buch des Lebens ausgetilget zu werden, wegen seiner Juden, denen er doch viel Guts gethan, und Gottes Zorn um öftern von ihnen abgewandt, da sie es doch nicht erkannt, sondern ihm viel Uebels erzeigt, also, daß sie ihn auf eine Zeit gedachten zu steinigen und zu tödten. Auch ist ein Exempel der Apostel Paulus, der da sagt Röm. 9, 3.: Ich begehre ein Anathema und Fluchopfer zu werden, daß meine Brüder, die Juden, möchten selig sein. Davon lehret auch das 4. Kapitel.

\section{Trauern und Fasten Mardachai und der Juden}
% Source scan 712; block decisions in selection.json

\quellseite{712}

\section{Summarium}

Ist eine solche, schöne Historig, wie Mardachai die Königin Esther vermahnet, daß sie ihr Leib und Leben in Gefahr setzen und bei dem Könige versuchen soll, ob das tyrannische Gebot wider die Juden möchte abgeschafft werden, V. 13.

Hier siehet man einen Unterschied zwischen Frommen und Gottlosen. Da Haman im Sinn hatte, etwas wider die Juden anzufangen, warf er das Loos und brauchte Zauberei; aber Esther und Mardachai, ehe sie anfangen, die Sachen zu handeln, demüthigen sie sich für Gott mit Fasten und ernstlichem Beten, darum gibt Gott auch Glück zu ihrer Handlung.

Am Ende ist ein fein Exempel, daß die fromme Esther ihr Leib und Leben in Gefahr setzet, ihr Volk zu erretten.

\section{Votum}

Helffe uns GOtt, daß wir auch einander herzlich meinen, und einander dienen und rathen, und derwegen keine Gefahr scheuen, und daß wir all unser Thun mit Gebet anfangen und glücklich verrichten, Amen.

\chapter{Das 5. Kapitel}

\section{Vorrede}

Der Königliche Prophet David spricht im 34. Psalm V. 7.: Da dieser Elende rief, hörte der HErr, und half ihm aus seinen Nöthen. Dieses wiederholet er und sagt V. 18.: Wenn die Gerechten schreien, so höret der HErr, und errettet sie aus all ihrer Noth. Das ist, das gläubige Gebet höret und erhöret Gott, und errettet, daß man nicht verderbe, aus aller Noth, wie groß und schwer, wie dick und breit die auch sein kann und mag. Darnach im 7. Psalm sagt aber David V. 15.: Siehe, der hat Böses im Sinn, mit Unglück ist er schwanger; er wird aber einen Fehl gebären. Er hat eine Grube gegraben und ausgeführt, und ist in die Grube gefallen, die er gemacht hat. Das ist, die Sünde ergreift zuletzt den Uebelthäter. Wer einem andern einen Stein in Weg legt, der fällt selbst darüber; wer einem andern einen Strick leget, der wird darinnen gefangen; wer untreulich redet, auf den wird es fallen, Spr. Sal. 27. Denn der gerechte Gott läßt nichts Böses ungestraft, und nichts Gutes unbelohnet, es stehe an lang oder kurz. Dieser Sprüche und Lehren Exempel ist das 5. Kapitel.
% Source scan 713; block decisions in selection.json

\quellseite{713}

\section[Esther gehet zum Könige, und Haman stellt Mardachai]{Esther gehet zum Könige, und Haman stellt Mardachai nach dem Leben}

\section{Summarium}

Ist eine schöne Historia, wie Gott der Esther Glück und Heil gibt, daß sie bei dem Könige in sondern Gnaden ist, V. 2. Denn was man mit Gott und dem Gebet anfänget, da will Gott mit Glück und Segen bei sein. Wiederum ists ein wunderlich Exempel mit Haman, daß alle Gottlosen schrecken und die Frommen wiederum trösten soll. Denn da ihm sein Unglück am nächsten ist, ist er am sichersten und zornigsten und hält für gewiß, ehe er morgen das Mittagmahl esse, soll Mardachai schon erhangen sein, V. 4. Aber Gott wendet’s wunderlich, und zum Trost, wenn wir fromm bleiben und an Gottes Wort hangen, so soll alles Unglück, das unsere Feinde uns zubereiten, über ihren eigenen Kopf kommen und wir errettet werden.

\section{Votum}

Der treue gütige GOtt erbarme sich auch unser in unserer Noth, erhöre unser Gebet und steure den Feinden seiner Kirchen und Christgläubigen Anschlägen und Werken, daß sie ohne Schaden abgehen und Er allezeit gelobet und gepreiset werde, Amen.

\chapter{Das 6. Kapitel}

\section{Vorrede}

Gott pfleget aus kleinen, verachteten Leuten große und ansehnliche zu machen, und aus großen kleine. Demgegen sollen wir lernen, daß die Hand und Gewalt Gottes stark genug sei, die Hoffärtigen zu stürzen und die Verdrückten zu erhöhen. Denn es ist keine Gewalt so groß, Gott kann die leichtlich zu nichte machen; auch ist keine Noth so groß, Gott kann aus solcher die Seinen erlösen. Denn die Hand des Herrn hilft gewaltiglich und kann alles ändern, sagt der königliche Prophet David, Ps. 20, 7. und 77, 17. Und der Herr ist gütig und eine Beste zur Zeit der Noth, und kennet die, so auf ihn trauen. Wenn die Fluth überher läuft, so macht ers mit derselben ein Ende, tröstet der heilige Prophet Nahum im 1. Kapitel V. 8. Davon höret auch an das 6. Kapitel.
% Source scan 714; block decisions in selection.json

\quellseite{714}

\section{Mardachai wird zu hohen Ehren erhoben}

\section{Summarium}

Sie ist eine Historia, in welcher man fürnehmlich sehet, daß Gott alle Sachen in seinen Händen hat und wunderbärlich alle Dinge ordnen und schicken kann. Mardachai hat keine Gedancken mehr, daß er der Wohltat genießen sollte, daß er den König gewarnet hatte, sondern stund in Gefahr Leibes und Lebens, sammt allen andern Juden. Dagegen war Haman sicher, und hielts gar für gewiß, Mardachai sollte in einer Stunde an einem Baum hangen. Aber Gott wendets zu beiden Theilen wunderbärlich, daß Haman selbst den Rath gibt und Mardachai dadurch zu den höchsten Ehren kommt, V. 10. Darum soll sich niemand der gottlosen Praktiken erschrecken lassen, sondern alle Hoffnung auf Gott stellen, als der unversehens alles ändern und bessern kann.

Daneben ist ein fein Exempel von dem Könige, daß er die Zeit nicht mit Spielen noch Zechen unnütz zubringet, da er nicht schlafen kann, wie doch an Höfen gebräuchlich, sondern läßt ihm zu Bett lesen die alten Historien, und höret zu, V. 1. Solches sollten alle Fürsten und Regenten merken, und dergleichen auch thun. Aber wie man sehet, man hat zu Hofe für dem Wild, für dem Zechen und Spielen, ich geschweige anderes dergleichen Thuns, nicht dißweil, daß man mit Büchern könnte umgehen. Und solches sündliches und ärgerliches Thun liebet den Herren dermaßen, daß sie ihres Amts darüber vergessen und andern befehlen, daß sie Amtes halber selbst thun sollten, Händel verhören, Uneinigkeit stillen, die unbillig Bedrängten schützen, jedermann Recht schaffen und was mehr der Obrigkeit zustehet.

\section{Votum}

Der König aller Könige, unser lieber Gott, helfe, daß alle Obrigkeit und ihre Amtverwalter recht thun; und siehe seiner Kirchen auch bei in diesen gefährlichen Zeiten und errete aus aller Noth alle, die Ihn anrufen, und auf Ihn hoffen und vertrauen, Amen.
% Source scan 715; block decisions in selection.json

\quellseite{715}

\chapter{Das 7. Kapitel}

\section{Vorrede}

Der königliche Prophet David sagt im 34. Psalm V. 20.: Der Gerechte muß viel leiden; aber der HErr hilft ihm aus dem allen. Und der heilige Hiob spricht im 20. Kapitel V. 5. 19.: Der Ruhm der Gottlosen stehet nicht lange, und die Freude des Heuchlers währet einen Augenblick? Wenn gleich seine Höhe in den Himmel reicht, und sein Haupt an die Wolken rühret; so wird er doch zuletzt umkommen wie ein Koth; allerhand Mühe wird über ihn kommen. Denn er hat unterdrückt und verlassen den Armen. Beides von Errettung der Frommen und Untergang der Gottlosen lehret David im 37. Psalm V. 5. also: Befehl dem HErrn deine Wege, und hoffe auf ihn, er wirds wohl machen, und wird deine Gerechtigkeit herfür bringen wie das Licht, und dein Recht wie den Mittag. V. 35.: Ich habe gesehen einen Gottlosen, der war trotzig, und breitete sich aus, und grünte wie ein Lorbeerbaum. Da man vorüber ging, siehe, da war er dahin; ich fragte nach ihm, da ward er nirgend funden. Dieser Sprüche Exempel an Mardachai und Haman erzählet das 7. Kapitel.

Haman wird an einen Baum gehänget.

\section{Summarium}

Ist eine tröstliche Historia für alle gottseelige Menschen, daß Gott nicht allein die Frommen retten und ihre Unschuld an Tag bringen, und sie ehren will, sondern auch die Tyrannen strafen und alles Unglück, das sie den Frommen gedacht, über ihren eigenen Kopf bringen. Mardachai, der fromme Jude, sollte an einen Baum erhangen werden; so schaffets Gott, daß man Haman, den Tyrannen, selbst daran hanget, V. 10. Solche Aenderung kann Gott mit den Tyrannen machen. Darum soll man Gott vertrauen und sich für den bösen Praktiken der Tyrannen nicht so sehr fürchten; Gott ist allewege mächtiger und weiser, denn sie.

\section{Votum}

GOtt, der HErr Zebaoth, steure und wehre heute auch allen Feinden seiner Christenheit, und tröste und erhalte alle, die seinen heiligen Namen lieben und ehren, bringe auch ihre Unschuld an den Tag, und stärke unsern Glauben, Amen.

\chapter{Das 8. Kapitel}

\section{Vorrede}

Der heilige Geist im 37. Psalm sagt: Ergrimme dich nicht über den, dem sein Muthwille glücklich fortgehet; denn die Bösen werden ausgerottet; die aber des HErrn harren, werden das Land erben. Es ist noch um ein Kleines, so ist der Gottlose nimmer; und wenn du nach seiner Stätte sehen wirst, wird er weg sein. Aber die Elenden werden das Land erben, und Lust haben in großem Frieden. Und wie ein Palmbaum die Eigenschaft hat, daß er mit seiner Last gebeuget, viel weniger unterdrückt werden kann; sondern je mehr man ihn beschweret, je höher er
% Source scan 716; block decisions in selection.json

\quellseite{716}
sich dagegen aufbäumet: also kann kein Fall die Schriften stürzen, wie groß er auch ist. Denn die Höllenpforten können sie nicht überwältigen, Matth. 16, 18. Also, wie das Gras daher grünet, daß einem Menschen das Herz im Leibe lacht, wenn ers so lieblich stehen siehet; aber bald wird abgehauen, da liegt es und verdorret, und wird von den Pferden und Kühen aufgefressen: also wendet sich auch bald das Blatt mit den Gottlosen, daß sie heute in Freude, morgen in Leid, heute im Leben, morgen im Tode und gestorben sind. Von diesen Wunderthaten Gottes lehret auch das 8. Kapitel.

\section[Die Juden haben Erlaubniß, sich an ihren Feinden]{Die Juden haben Erlaubniß, sich an ihren Feinden zu rächen}

a. Das ist: Weil die vorigen Briefe Hamans waren mit des Königs Siegel versiegelt, hätte es die Juden nichts geholfen, wo sie nicht wären von neuem durch andere Briefe widerrufen.

\section{Summarium}

Ist eine Historia, wie sich die Sachen gar verkehren, daß der König nicht allein das Gebot
% Source scan 717; block decisions in selection.json

\quellseite{717}
widerrufet, sondern auch den Juden erlaubet, auf den bestimmten Tag sich zu rächen an ihren Feinden, die ihnen übels gedacht, und sich ihres Unglücks erfreuet hätten, V. 13. Also will Gott allewege die Seinen schützen und die Tyrannen dämpfen, so fern man nur das Vertrauen auf ihn allein setzet; aber solche Hülfe verzeucht sich eine Zeitlang. Denn die Juden haben bis in die drei Monat in der Sorge und Gefahr stehen müssen, daß die Heiden auf einen Tag sie alle erschlagen und vertilgen, das Gut, Leib und Leben, Weib und Kind und alles rauben sollten.

\section{Votum}

Der treue barmherzige Gott wolle heute auch alle Tyrannen sammt ihrem Anhang stürzen; und schütze seine arme bedrängte Christenheit, und verleihe allen Kreuzträgern unter dem Verzug der Hülfe wahren Glauben und beständige Geduld, Amen.

\chapter{Das 9. Capitel}

\section{Vorrede}

Gott sagt im 89. Psalm V. 23.: Die Feinde sollen meinen Knecht David nicht überwältigen, und die Ungerechten sollten ihn nicht dämpfen; sondern ich will seine Widersacher schlagen vor ihm her, und die ihn hassen, will ich plagen: das ist, Gott will die Gläubigen und Frommen erhalten, stärken, wider alle Feinde beschützen, ihre Widerwärtige plagen und schlagen. Auch spricht David im 50. Psalm V. 14.: Opfere Gott Dank, und bezahle dem Höchsten deine Gelübde: das ist, erkenne o Mensch, und vergiß nicht, was Gott Gutes an uns gethan, daß er uns all unsere Sünde vergibet, unser Leben vom Verderben erlöst und uns mit Gnaden und Barmherzigkeit gekrönet hat. Dem sollt du auch mit Mund, Herzen und mit deinem ganzen Leben dafür danken, ihm vertrauen und ihn für den Brunnquell aller Güter erkennen und halten. Dieser, als des Trosts und Lehrspruchs Exempel ist das 9. Capitel.

\section[Die Juden stellen nach göttlicher Rache ein Freudenfest]{Die Juden stellen nach göttlicher Rache ein Freudenfest an}
% Source scan 718; block decisions in selection.json

\quellseite{718}

\section{Summarium}

Ist eine Historia, wie Esther und Mardachai die zween Tage, an welchen ihnen Gott gnädig geholfen und sie errettet, und an ihren Feinden gerochen hatte, verordneten, daß mans feiern soll, V. 18., und solches um der Ursache willen, daß man solcher großen Wohlthat nicht vergessen, sondern Gott dafür danken, und solches sich weiter in andern Nöthen erinnern und damit trösten möge, Gott werde auch helfen, wie er diesmal geholfen habe. Solcher Dank und Zuversicht auf Gottes Hülfe ist zuvor geboten in der ersten Tafel Mose. Daß man aber die Feiertage dazu verordnet, geschieht darum, daß der gemeine Mann solcher Ceremonien darf, sich damit auf einen gewissen Tag zu erinnern der vorigen Wohlthaten Gottes. Dergleichen Feste haben die Juden hernach mehr geordnet, als sie wieder ins Land sind kommen.

\section{Votum}

Gott, der hochgelobten Dreifaltigkeit, sei Lob, Ehr, Preis und Dank, der bisher auch seine Kirche wunderlich erhält und schützet; der vertilge ferner alle der Christenheit Feinde, und helfe uns, daß wir seine Gottesdienste allzeit in wahrem Glauben fleißig verrichten, Amen.

\chapter{Das 10. Capitel}

\section{Mardachai wird gerühmet}

\quellseite{719}

\part{Das Buch Hiob}

\section{Vorrede über das Buch Hiob}

Das Buch Hiob handelt diese Frage: Ob auch den Frommen Unglück von Gott widerfahre? Sie stehet Hiob fest und hält, daß Gott auch die Frommen ohne Ursache allein zu seinem Lob peiniget. Wie Christus Joh. 9, 3. von dem, der blind geboren war, auch zeuget.

Danwider setzen sich seine Freunde und treiben ein groß und lang Geschwätz, wollen Gott Recht erhalten, daß er keinen Frommen strafe; strafe er aber, so müsse derselbe gesündiget haben. Und haben so einen weltlichen und menschlichen Gedanken von Gott und seiner Gerechtigkeit, als wäre er gleich wie Menschen sind, und sein Recht, wie der Welt Recht ist.

Wiewohl auch Hiob, als der in Todes-Nöthen kommt, aus menschlicher Schwachheit zu viel wider Gott redet, und im Leiden sündiget und doch darauf bleibet, er habe solch Leiden nicht verschuldet für andern, wie es denn auch wahr ist. Aber zuletzt urteilet Gott, daß Hiob, indem er wider Gott geredet hat im Leiden, unrecht geredet habe; doch was er wider seine Freunde gehalten hat von seiner Unschuld für dem Leiden, recht geredet habe. Also führet dieses Buches Historia endlich dahin, daß Gott allein gerecht ist, und doch wohl ein Mensch wider den andern gerecht ist auch für Gott.

Es ist aber uns zu Trost geschrieben, daß Gott seine großen Heiligen also lässet straucheln, sonderlich in der Widerwärtigkeit. Denn ehe, daß Hiob in Todes-Angst kömmt, lobet er Gott über dem Raub seiner Güter und Tod seiner Kinder. Aber da ihm der Tod unter Augen gehet, und Gott sich entzeucht, geben seine Worte Anzeige, was für Gedanken ein Mensch habe (er sei wie heilig er wolle) wider Gott, wie ihn dünket, daß Gott nicht Gott, sondern eitler Richter und zorniger Tyrann sei, der mit Gewalt fahre, und frage nach niemands gutem Leben. Dies ist das höchste Stück in diesem Buch. Das verstehen allein die, so auch erfahren und fühlen, was das sei, Gottes Zorn und Urtheil leiden, und seine Gnade verborgen sein.

Sonderlich aber ist hie zu merken, daß der Text klar anzeiget, daß Gott dem Teufel Erlaubniß gebe, auch die Frommen anzufechten, V. 6. Und daß Krieg, Feuer vom Himmel, große Winde und schädliche Wetter vom Teufel aus Verhängniß Gottes verursacht und gemacht werden. Darum soll man wider solchen Feind, der den Menschen zu schaden alle Länder durchsucht, Gott bitten, daß er uns durch seine lieben Engel wider ihn schützen und für dergleichen Jammer behüten wolle. Wie der böse Feind hie selbst zeuget, Gott habe bisher den Hiob verwahret, gleich mit einer Mauer, durch seine Engel. Begehret aber, daß Gott solchen Schutz wegnehmen und ihm erlauben soll. Solches ist nicht also zu verstehen, als habe der Teufel mit Gott geredet, und um Erlaubniß gebeten. Denn von solchem Thun kann man anders nicht reden, denn wie es unter uns Menschen gehet. Darum läßt es die Schrift bei solcher Einfalt bleiben, sonst könnten wir nichts davon verstehen. Daß der Teufel die Welt durchwandert, V. 7., ist wahr, denn er ist an allen Orten, und siehet, wo er könne Schaden thun. Daß er aber für Gott kommt, ist nicht zu verstehen, als käme er in Himmel, sondern auf das einfältigste, daß er ohne Gottes Bewilligung nichts könne ausrichten. Darum heißet Gott fragen oder zu Gott kommen, anders nichts, denn auf Gottes Willen sehen, ob ers verhängen oder zusehen wolle, wie es der Satan im Sinn hat.

\chapter[1]{1}

\section{Vorrede}

Im ersten Kapitel ist eine Historia, wie Hiob mit mancherley Unglück auf einmal umgeben, V. 15., 16., 17., 18., und dennoch Gott darin gelobet, und alles geduldig gelitten habe, V. 27. Uns zum Exempel, daß wir dergleichen, wenn es uns übel gehet, auch thun sollen.

\section{Hiobs Widerwärtigkeit und Geduld}

a Nicht daß er so reich und gewaltig sei gewesen, sondern um seiner Weisheit, Verstand und Gottseligkeit willen ist er herrlicher gehalten, denn andere.
\quellseite{720}
ᵇ Segnen) das ist: Fluchen und lästern.

\section{Gebet}

GOtt, unser treuer lieber Vater, erbarme dich unser, und beschütze uns wider den Teufel durch seine heiligen Engel, daß er nicht auch bei uns Jammer und Noth anrichte, und helfe, daß wir in all unserm Kreuz Ihn auch loben, und das in Geduld tragen und überwinden, Amen.

\chapter[2]{2}

\section{Vorrede}

Geliebte Christen, St. Petrus spricht 1. Epistel am 5., V. 8.: Euer Widersacher, der Teufel, gehet umher wie ein brüllender Löwe, und suchet, welchen er verschlinge. Derwegen wir nicht sicher sein sollen, als wäre der Teufel viele Meilen von uns; sondern sollen wachen und beten. Und St. Paulus 1 Thess. 5, 16. vermahnet: Seid allezeit fröhlich, betet ohne Unterlaß, und seid dankbar in allen Dingen; denn das ist der Wille Gottes in Christo Jesu an euch. Derwegen sollen wir Gott auch unter dem Kreuz loben. Denn was von der Hand Gottes kommt, das ist nicht böse, so wir es mit Glauben empfahen. Denn den Gläubigen muß alles zum Besten dienen, Röm. 8, 28. Ferner, der heilige Prophet Jeremias sagt im 49. Kapitel B. 12.: So spricht der Herr: Siehe, diejenigen, so es nicht verschuldet hatten, den Kelch zu trinken, müssen trinken. Derwegen, ob wir schon unschuldig sein, so sollen wir doch Gott zu ehren, was er nur will, in Geduld leiden, weil der unschuldige Sohn Gottes um unsertwillen den bittern Tod erlitten hat, daß er uns aus der Hölle und vom ewigen Tode erlösete. \unsicher{Erzählter} Sprüche und Lehren erinnert auch das 2. Kapitel.

\section{Hiob am Leibe geplaget, wird betrachtet und besucht}
\quellseite{721}

a Haut für Haut) das ist, für seine Haut läßt er fahren Kinder, Vieh, Gesinde und aller andern Haut.

b (Segne Gott) ja du thust fein, lobest und dienest Gott, und gehest drüber zu Grunde.

\section{Summarium}

Dies ist eine Historia, in welcher wir sehen, wie der Teufel anhält und wo er mit einer Anfechtung nichts kann ausrichten, sich je mehr und mehr mit höhern und schwerern Anfechtungen an dem Menschen versucht, ob er ihn zu Falle, in Ungeduld und Sünde bringen könne. Sie sollten wir von dem frommen Hiob lernen, daß wir, wenn's übel gehet, uns auch trösten, wie er sich tröstet, und des Unfalls nicht beschweren, sondern Gott auch im Leiden loben, der uns unser Lebenlang unzählig viel Gutes, beides an Leib und Seele gethan und aus viel Unglück väterlich geholfen hat.

Und hie zeiget der Text das fürnehmste Stück an, das in diesem Buch gehandelt wird, nämlich, daß Gott sagt, er habe Hiob ohne Ursache also geschlagen, B. 3. Auf daß wir lernen, daß auch die Frommen und Gottesfürchtigen unter das Kreuz müssen, ob sie es gleich nicht wie andere Leute, so ohne Gottesfurcht leben, verdienen.

\section{Votum}

Unser HErr JEsus Christus, der in die Welt kommen ist, daß er des Teufels Werk zerstöre, der wehre des Teufels List und Gewalt, und tröste und erhalte unter allem geistlichen und leiblichen Kreuz alle, die seinen heiligen Namen loben und anrufen, und helfe, daß wir allezeit in Gottesfurcht leben, Amen.

\chapter[3]{3}

\section{Vorrede}

Geliebte Christen! gleich wie der Aramant so hart ist, daß er nicht kann zerschlagen werden: also tragen die Gläubigen das Kreuz nicht allein mit geduldigem, sondern auch mit starkem und unüberwindlichem Herzen und Gemüth; gleichwie aber dennoch der starke Aramant, dem das Eisen und Feuer nichts abbrechen kann, durch frisches und warmes Bocksblut zerbrochen wird: also, wenn die Gläubigen Gottes brennenden Zorn fühlen, und das heiße höllische Feuer ihnen ein wenig ins Herz scheinet und leuchtet, so gerathen sie oft dahin, daß sie wider Gott murren und wollen lieber todt sein, denn leben. Aber also muß unser HErr Gott die großen Heiligen demüthigen, daß sie lernen sich selbst recht erkennen, wie ein elend Gemächt ein Mensch sei, wo er nicht unter dem Schutz und Schirm des Höchsten ist. Es sollen auch die betrübten, angefochtenen Herzen und verzagten Gewissen dies merken, daß sie nicht verzweifeln, wenn sie auch dergleichen Ungeduld und Murren unter dem Kreuz wider Gott fühlen; sondern sollen sich trösten, daß die größten Heiligen sowohl Fleisch und Blut gehabt, und auch, ja so leichtlich gestrauchelt und gefallen sind, als wir; Gott aber habe sich ihrer erbarmet und sie zu Gnaden angenommen; der werde uns auch mit gnädigen Augen ansehen, und uns die Sünde vergeben, um unsers HErrn JEsu Christi willen, und uns erretten und erhalten. Davon lehret auch das 3. Kapitel.
\quellseite{722}

\section{Hiobs Wehklage und Ungeduld}

a Meines Leibes) daraus ich geboren ward, das ist der Mutter Leib.

b Wüste) die mit Bauen umgehen, da zuvor nichts steht.

c Verborgenen) das ist, aus der Erden.

b Bedeckt) was soll der leben, der für Angst nicht weiß, wo aus, wohin, besser todt etc.

\section{Summarium}

Sie gehet an das rechte Leiden Hiobs. Der Teufel durch seine Anfechtung hatte das angerichtet, daß er gesucht hatte, nämlich, daß Hiob nicht mehr spricht, wie zuvor: Der HErr hats gegeben, der HErr hats genommen, der Name des HErrn sei gelobet. Item: Haben wir Gutes empfangen von Gott, warum sollten wir das Böse nicht auch annehmen? Sondern er klaget und schreiet über sein Leiden dermaßen, daß er nicht schweigt und wollte, er wäre nie geboren, B. 2. 3. 4. Also siehet man, wie auch heilige Leute, wenn das Leiden anhält, ungeduldig werden und Gottes Willen nicht können mit Gehorsam ertragen. Denn Fleisch und Blut lässet von seiner Art nicht, auch in den Heiligen. Es murret und klaget, wenn ihm übel ist, wollte des Leidens gern überhoben sein. Solcher Ungehorsam ist eine greuliche verdammliche Sünde. Aber Gott wills gerne vergeben, wenn mans nur für Sünde erkennet und um Gnade und Vergebung bittet.

\section{Votum}

Der barmherzige GOtt vergebe uns unsere Sünde, lindere uns auch die wohlverdienten Strafen, und regiere und stärke uns, daß wir unter dem Kreuz mit Murren und Ungeduld nicht sündigen, sondern allezeit in Lieb und Leid ihn loben und preisen, Amen.

\chapter[4]{4}

\section{Vorrede}

Geliebte Christen! welchen anderer Leute Unglück eine Freude ist, die sind ärger, als die un-
\quellseite{723}
vernünftigen Thiere; denn diese sind traurig, winseln und schreien, wenn sie ihre Gattung oder ihres gleichen in Noth sehen und hören; ja solche sind wahrhaftig Teufelsgenossen, denn der ist ein rechter Schadenfroher und seine Lust und Freude daran hat, wenn es den Menschen übel gehet und sie in Herzensangst, Leibesnoth und Güterverlust oder Ehrenschaden gerathen und stecken. Derwegen soll sich jedermann für solcher Unbarmherzigkeit und teuflischer Bosheit und Sünde fleißig hüten. Denn Gott will auch kurzum nicht, daß wir uns anderer Leute Unfall freuen sollen; sintemal wir alle arme Menschen sind, und kann einem jeden begegnen das, so einem andern widerfahren ist. Dazu sind wir alle Sünder und ist ein jeder werth, daß er stets von Gott heftig gestraft werde; und weil es herzlich und schmerzlich weh thut, wenn eines von denen verhöhnet und geplaget wird, die ihn billig trösten sollen; so soll auch niemand andern ihr Kreuz, Leiden und Elend größer und schwerer machen, oder sie mehr mit Worten und Thaten betrüben, wie dem lieben Hiob seine Freunde gethan haben. Als davon das 4. Kapitel also meldet.

\section{Eliphas fängt an Hiob zu bestrafen}

a Ist das) das ist, da siehet man nun, wie fromm du seist, daß dich Gott so strafet.

b Die Löwen und Löwinnen sind die Reichen und Gewaltigen auf Erden, so die Armen unterdrücken.

c Boten) oder Engeln.

\section{Summarium}

Hie gehet die Disputation an, so die Freunde Hiobs mit ihm in seinem Leiden gehalten und ihm dadurch dasselbige nur größer und schwerer gemacht haben. Des frommen Hiobs Herz stund also: er wollte den Jammer alle gerne leiden, wenn ers verdienet hätte, wie es die Welt mit ihrem gottlosen unbußfertigen Leben verdienet. Denn Gott selbst gibt ihm Zeugniß, im 2. Kapitel B. 3., daß er ihn ohne Ursache habe plagen lassen. Solche Meinung wollten ihm seine Freunde ausreden, und sprechen: Gott sei gerecht, darum thue er niemand unrecht. Weil aber Hiob so heftig gestrafet und mit so viel Jammer überschüttet werde, da müsse folgen, er habe solches wohl verdienet und geschehe ihm nicht unrecht, B. 7. Das heißt denn, das Leiden größer und schwerer machen, wenn das Gewissen dazu kommt, daß es eine verdiente Strafe der Sünden sei, und daß es Gott billig über uns habe kommen lassen.

\section{Votum}

Der heilige Geist regiere uns, daß wir betrübte Leute nicht mehr betrüben, sondern herzlich Mitleiden mit ihnen tragen, sie trösten und ihnen rathen und helfen, auf daß sie, durch seinen Beistand, ihr Kreuz in Geduld tragen und zu Ende bringen, Amen.
\quellseite{724}

\chapter[5]{5}

\section{Vorrede}

Geliebte Christen! es ist wahr, daß Gott der Herr sagt zu seinen Gläubigen Jeremiä im 30. Kapitel V. 11.: Züchtigen will ich dich mit Maße, daß du dich nicht unschuldig haltest. Aber das ist auch wahr, daß der heilige Geist durch Salomon lehret in seinen Sprüchen im 17. Kapitel V. 3., da er spricht: Wie das Feuer Silber und der Ofen Gold, also prüfet der Herr die Herzen. Also ist wahr, daß im Buch der Weisheit im 3. Kapitel V. 5. gesagt wird: Die Gerechten werden ein wenig gestraft, aber viel Gutes wird ihnen widerfahren; denn Gott versucht sie, und findet sie, daß sie sein werth sind. Dagegen wie man das Vieh, welches man bald schlachten will, nicht zu schwerer Arbeit gebraucht, und das nicht abmergelt, daß es die Haut kaum schleppen möge, sondern man mästet es: also läßt Gott die Gottlosen hie eine Zeit machen, was sie nur wollen, und in allem Bösen fortfahren; aber da sie nicht Buße thun, so belohnet ihre Sünde, nach den zeitlichen Strafen, das ewige höllische Feuer. Dieser Lehren soll uns erinnern das 5. Kapitel.

\section{Eliphas fährt in seiner Bestrafung fort}

a Heiligen das ist, zeige mir einen Heiligen, der unschuldig ist geplaget, wie du meinst. Aber die Tollen und Unwitzigen heißet er hie die losen, frechen Leute, die nach Gott nicht fragen; solche verderbet wohl der Zorn und Eifer Gottes.

b Hungrige und Durstige heißt er die Räuber und Tyrannen.

c Nicht gehet das ist, der Mensch verdienet solches mit Sünden, sonst käme es ihm nirgend her.

d Dein Bund das ist, die Steine werden dein Getreide bewahren, weil davon eine Mauer umher gemacht wird.

\section{Summarium}

Das ist gleich die Meinung des vorigen Kapitels; denn Eliphas sagt, die tägliche Erfahrung weise aus, daß die Bösen um ihrer Sünde willen gestraft werden, wiederum, daß Gott den Frommen helfe, B. 2. Darum vermahnet er Hiob, er soll ja nicht gedenken, daß er solchen Jammer nicht habe verdienet, sondern sich frei
\quellseite{725}
für einen Sünder erkennen, so werde Gott wieder mit seinem Trost kommen, und ihm helfen, V. 17. Solches ist der Vernunft nach recht geredet.

Aber Gott führet sein Werk höher, denn es die Vernunft verstehen kann, strafet nicht allein darum, daß man es verdienet habe, sondern daß er prüfen wolle, ob man auch fest am Vertrauen auf seine Hülfe und Gnade halten, und ihm den schuldigen Gehorsam auch im Kreuz und Leiden leisten und seinen Willen geduldiglich tragen wolle.

\section{Votum}

Der treue Gott verleihe, daß wir und alle Kreuzträger unter dem Kreuz seinen Willen erkennen; und helfe, daß wir in der Probe allzeit wohl bestehen, behalten den Glauben und ein gut Gewissen, Amen.

\chapter[6]{6}

\section{Vorrede}

Geliebte Christen! Gleich wie David sein Leid getragen und bitterlich geweinet hat über Saul und Jonathan, da sie durchs Schwert gefallen waren, 2 Sam. 1, 17.; und wie St. Paulus sich des kranken Epaphroditus herzlich annimmt, und seinethalben trauerte, und hoch bekümmert war, Phil. 2, 27.: also sollen wir mit unserm Nächsten in seinem Kreuz, als in Armut, Krankheit und andern Trübsalen, auch herzlich Mitleiden tragen und ihm behülflich und rathsam sein, so fern und viel wir mögen. Denn wer Barmherzigkeit seinem Nächsten weigert, der verlässet des Allmächtigen Furcht, sagt heute Hiob, B. 14. Aber es gehet, wie Sirach sagt und klaget im 37. Kapitel B. 3.: Ach wo kömmt doch das böse Ding her, daß alle Welt so voll Falschheit ist? Wenn es dem Freunde wohlgehet, so freuen sie sich mit ihm; wenn es ihm aber übel gehet, werden sie seine Feinde. Dafür hüte sich jedermann. Darauf höret an das 6. Kapitel.

\section{Hiobs Verantwortung wider Eliphas}

a Das Wild das ist, ihr habt gut trösten, euch mangelt nichts. Man isset nicht ungesalzene, wenn mans besser weiß; aber ich muß wohl selbst dies und das etc.

b Nicht verleugnet das ist, hab ichs doch nicht verdienet, daß ich so geplagt werde, wollte Gott, ich wäre doch todt.

c Verschmachten) das ist, weil meine Freunde jetzt vorüber rauschen, wie ein Wasser, und kennen mich nicht, wird sie auch einmal eine Hitze drücken, so werden sie versiegen und vertrocknen, darum, daß sie mich jetzt verlassen.

d Die Wege Thema) das ist, sie halten mit denen, die mich beraubet haben, wie oben Kapitel 1. steht, gehen denselben recht, und mir unrecht.
\quellseite{726}

\section{Summarium}

Dies ist eine Klagrede des frommen Hiobs, daß er neben allem Unglück, das Gott auf ihn geworfen hatte, auch noch dies leiden und tragen müßte, daß seine Freunde, die ihn billig trösten sollten, mit ihrem Reden nur verzagter und blöder machen, als die neben der Leibesnoth ihm auch das Gewissen beschweren und verzagt machen. Denn wo ein Berg vor zerschlagen und voll Kümmerniß ist, daselbst soll man nicht mehr martern und plagen, sondern trösten und auf Gottes Gnade und Barmherzigkeit weisen. Gleich wie man einen Wagen, der auf die Seiten hängen und fallen will, halten und nicht ganz zu Boden reißen oder stoßen soll. Denn dasselbige geschieht ohne das, wenn er hänget und sinkt, daß er leichtlich gar fällt. Solches thaten die Tröster nicht, die dem Hiob nichts anders konnten fürsprechen, denn daß er ein Sünder wäre, der solche Strafe wohl verdienet hätte, denn Gott thäte niemand unrecht etc. Aber Hiob wehret sich stark und spricht Nein dazu, solches ist nicht wahr; darum wie ihr mich tröstet, werdet ihr in eurer Noth und Jammer wieder getröstet werden.

\section{Votum}

Der barmherzige Gott erleuchte und erweiche unsere Herzen, daß wir auch mit unserm nothleidenden Nächsten Mitleiden haben, und ihm tröstlich und hülflich sein; und steure allen bösen und falschen Herzen, welche die Betrübten mehr betrüben, ängsten, quälen und plagen, Amen.

\chapter[7]{7}

\section{Vorrede}

Geliebte Christen! Es ist wohl wahr, daß viel Elend und Jammer auf Erden ist. Aber es ist alles wie nichts zu rechnen gegen der greulichen Anfechtung, wenn das Herz Gottes feuerbrennenden Zorn fühlet. Denn, wenn dem Menschen sein Herz absagt, so ist ihm immer zu Sinne, Gott habe ihn verlassen und verstoßen, und stehet noch weiß nicht, wie er aus der unaussprechlichen Angst und Noth kommen soll: Wo er sich hinwendet und kehret, stehet er nichts, denn den jämmerlichen Anblick, daß Gott zornig ist, und ihn weder sehen noch hören wolle. Und das ist ein Spiegel, wie es in der Hölle zugehen wird, in der ewigen unaussprechlichen Qual, da das Gewissen in den Gottlosen wird jämmerlich zagen, und der Zorn Gottes über ihnen bleiben; und das wird das rechte Herzeleid über alles Leid sein, und vielmehr brennen, als das höllische Feuer, wenn einem jeglichen sein eigen Herz absagen und ihn überzeugen wird, daß ihm nicht unrecht geschehe. O Herr, treuer Gott, laß uns bedenken den großen Ernst, daß wir klug werden und christlich leben. Das erinnert auch das 7. Kapitel.

\section{Hiobs Klage über der Menschen Elend}

a) Vergeblich, das ist, ich habe Ruhe und der Arbeit ein Ende gesucht, aber das ist umsonst, es bleibet noch immer Unruhe.
\quellseite{727}

\section{Summarium}

Ist eine ernstliche Klag-Rede, daß Gott so viel Jammers und Herzeleid auf die elende menschliche Natur wirft, so sie doch an ihr selbst schwach, gebrechlich, und aller Gefahr und Unglück unterworfen, dazu für dem Tod keinen Augenblick sicher ist, daß wohl Noth wäre, daß Gott schonet und hülfe. Aber Gott (spricht Hiob) legt sich noch selbst wider uns, drängt, drücket und ängstet uns. Darum wäre nichts besser, denn nur die Sünde vergeben und bald gestorben. Solche Worte alle gehen aus der angebornen Ungeduld und Ungehorsam, daß der alte Adam ihm Gottes Willen nicht gefallen lassen, sondern immer andere Mittel und Wege fürschlagen will, denn es Gott ordnet.

\section{Votum}

Gott sei uns allen gnädig, lindere uns die wohlverdienten Strafen, und erhalte uns und alle Kreuzträger in seinem Gehorsam, daß wir nicht aus Ungeduld, mit Worten oder Werken, wider ihn und unsern Nächsten sündigen, Amen.

\chapter[8]{8}

\section{Vorrede}

Geliebte Christen! Von Glück und Unglück urtheilet die Vernunft also: Daß, welchem Menschen es in seinem Leben übel gehet, der müsse einen ungnädigen Gott haben. Und wiederum: Wo alles ist, was das Herz begehret, da wohne Gott. Aber die heil. Schrift zeuget dagegen, daß welchen der Herr lieb hat, daß er denselben züchtige, Ebr. 12, 6. Und so die nur sollten einen gnädigen Gott haben, die in diesem Leben gute geruhige Tage haben, so würde daraus folgen, daß die heiligen Patriarchen, Propheten, Apostel und Märtyrer alle von Gott verstoßen wären. Daher stellet auch der Herr Christus das Urtheil von dem Blindgebornen, Joh. 9, 3., und sagt: Es hat weder dieser gesündiget, noch seine Eltern; sondern, daß die Werke Gottes offenbar würden an ihm. Dieß soll wohl betrachtet werden, daß niemand Gott in sein heimlich Gericht falle und seinen Nächsten richte und verdamme; wie mancher sich seines Nächsten Unglückes freuet, und weiß nicht, daß das seine so nahe für der Thür ist, und unterdeß grünet und blühet. Darauf höret auch davon an das 8. Kapitel.

\section{Hiob wird von Bildad der Heuchelei beschuldiget}
\quellseite{728}

\section{Summarium}

Hie ist die andere Disputation, welche Bildad mit Hiob hat; der stehet auch auf dem, Gott richte nicht unrecht; derhalben was Hiob leide, das leide er seiner Sünde halben, V. 3. Er soll aber frömmer werden, denn er bisher gewesen sei, so werde ihm Gott gnädig sein und aus dieser seiner Noth helfen, V. 5. Solche Meinung beweiset er mit der Erfahrung, daß Gott je und je die Frommen erhalten, und die Bösen verworfen und gestraft habe. Also muß die Vernunft schließen, welche nicht mehr von Gottes Willen und Werken weiß, denn im Gesetze offenbar ist; aber der arme Hiob bedarf eines andern Trostes und einer höhern Weisheit in seiner Noth und Anfechtung.

\section{Votum}

Gott, der HErr, vergebe es allen, die betrübte Kreuzträger mehr ängsten; und erleuchte sie, daß sie Buße thun und von solcher Sünde ablassen, und tröste alle Christen in ihrer Noth, und helfe uns auch gnädiglich im Leben und Sterben, Amen.

\chapter[9]{9}

\section{Vorrede}

Geliebte Christen! Wenn ein Kreuz hin ist, so denkt der Mensch, er sei nun gar hindurch. Aber in einem Nu findet sich ein anderes, daß Gott wieder vorüber gehet, ehe wirs gedenken, und verbirget sein Angesicht, läßt uns sinken, jedoch nicht gar ertrinken, und stellet sich, als kenne er uns nicht mehr. Aber das ist gewiß, wenn uns dünket zur Zeit der Noth, Gott sei am weitesten, so ist er uns am nächsten. Diese seine wunderliche Gegenwärtigkeit hat er in der Sonne uns vorgebildet. Wenn uns dünket mitten im Winter, sie sei am weitesten von uns, so ist sie etliche tausend Meilen Weges in ihrem Zirkel am Himmel niedriger, und dem Erdboden näher, denn mitten im Sommer, wenn um Johannis-Tag die größte Hitze ist. Gleich so ist der allmächtige Gott uns viel näher zur Zeit der Noth, denn wenn es gehet, wie wir selbst wollen. Ob aber wohl Fleisch und Blut solches nicht begreift, so findet sich es doch also endlich, wenn wir hindurch sind, und überwunden haben. Darauf höret nun an das 9. Kapitel.

\section{Hiobs Vertheidigung wider Bildad}

a Orion ist das helle Gestirn gegen Mittag, das die Bauern den Jacobstab heißen. Die Glucke oder die Henne sind die sieben kleine Gestirne.

b Die stolzen Junker, die sich auf ihre Macht verlassen und jedermann helfen können.
\quellseite{729}

c Scilicet, ne redeant. D. i. daß sie nicht wiederkehren.

d Kleider) das ist, meine Tugend.

\section{Summarium}

Ist eine Antwort auf Bildads Rede, in welcher Hiob noch immer darauf beharret, er sei nicht gottlos gewesen, habe auch solchen Jammer nicht wohl verdient. Daß aber Gott solches über ihn verhänget, das thue er als ein Gott, der Herr über alles ist, und mag mit jedermann machen, wie er will, dem niemand nichts einreden noch abgewinnen könne, V. 4. Solche Antwort ist etwas hart, aber der Schmerz und die große Noth und Anfechtung treiben Hiob, der von Herzen betrübt war, daß er also redete; sintemal er bekennen mußte, daß andere Leute so gottlose und böse gewesen, dennoch von Gott so hart nicht sind gestraft worden, als er, der doch gottesfürchtig und fromm gewesen war.

\section{Votum}

Gott, der heilige Geist, regiere und stärke uns alle, daß wir auch in der Frömmigkeit beständig bleiben, und was uns darüber widerfähret, daß wir es in christlicher Geduld tragen, und Gott allezeit loben und preisen, bis zu unserm seligen Ende, Amen.

\chapter[10]{10}

\section{Vorrede}

Geliebte Christen! In den Gläubigen, ob sie schon Gottes heilige und auserwählte Kinder heißen, streitet auch allwegen hienieden auf Erden in ihren Herzen der Geist wider das Fleisch, und der Glaube wider den Unglauben. Denn jetzt haben sie gewiß und glauben stark, Gott sei ihnen gnädig und barmherzig, und lache sie freundlich an; bald gerathen sie wieder in schwere Gedanken: Gott zürne mit ihnen, sei ihnen spinnefeind, habe ihnen den Rücken zugekehrt, seine Barmherzigkeit ihnen entzogen, und wolle ihr ganz und gar keine Gnade haben. Dieser Wechsel in den Christgläubigen höret nicht auf bis in die Grube. Derhalben soll niemand gedenken, daß die Christen allzeit stark im Glauben, brünstig in der Liebe und bissig im Gebet sein können und bleiben. Und fromme Christen sollen sich desto fleißiger nach dem Ziel strecken und Gott bitten, daß er ihnen Victoria und Sieg durch den Glauben verleihen und erhalten wolle, Amen. Darauf höret auch an das 10. Kapitel.

\section{Hiob klaget über seinen Jammer}
\quellseite{730}

a Um und um) nichts ist an mir, das du nicht gemacht hast, oder nicht dein sei. Noch verwirfst du mich, als hätte mich ein anderer gemacht, der dein Kind wäre, sogar nimmst du dich deines Eigenen nicht an.

b Odem) das ist, mein Leben, das der Odem anzeigt.

\section{Summarium}

Ist eine Klagrede, in welcher Hiob des fürnehmlich sich beklagt, er müsse leiden und wisse doch nicht warum. Denn er sei nicht gottlos gewesen, und wisse, daß Gott ein Wohlgefallen an den Frommen habe.

Darum klaget er, daß Gott, von welchem wir das Leben und alles haben, sich mit solcher Macht wider einen armen Menschen legt, V. 14.; und wünschet abermals, daß er nicht geboren wäre, V. 18. Also siehet man, wie Natur und Vernunft aus großer Ungeduld wider Gott murret, und des Trostes vergißet, welchen Gott allwege hinter das Leiden stellet und verberget, ob er gleich die Menschen läßt plagen, daß er doch wiederum wolle helfen, wo man sich Hülfe zu ihm versiehet, und bei ihm suchet.

\section{Votum}

Der Gott alles Trostes und der Geduld behüte uns und alle christlichen Kreuzträger für Ungeduld und allen andern Sünden; und erhalte sie und uns in wahrem Trost und allen christlichen Tugenden, bis an unser Ende, Amen.

\chapter[11]{11}

\section{Vorrede}

Geliebte Christen! Den betrübten Kreuzträgern soll man freundlich, tröstlich und herzlich zusprechen, und sie vermahnen, daß sie das liebe Kreuz, so ihnen Gott der Herr auferlegt, geduldig tragen und sich unter die gewaltige Hand Gottes, damit er unterdrückt, aber auch wieder aufhilft, demütigen sollen. Man soll sie trösten, daß solch ihr Kreuz und Leiden zu Gottes Ehre und ihrer Seelen Seligkeit durch Christum gereichen werde, und daß ihnen Gott zu seiner Zeit werde abhelfen. Und wenn sie ohne alles Kreuz wären, so möchten sie zu Weltkindern und mit der Welt verdammt werden. Auch werde ihnen solch Kreuz dazu dienen, daß sie denn ihren treuen Gott, wenn er ihnen wiederum aus allen ihren Anfechtungen herrlich helfen wird, werden von Grund ihres Herzens danken, loben und preisen, daß er sie so väterlich gezüchtiget und durch das liebe Kreuz für vielen Sünden behütet hat; und werden mit allen
\quellseite{731}
lieben Heiligen das Freuden-Liedlein aus dem 118. Psalm V. 18. singen: Der Herr züchtiget mich wohl, aber er gibt mich dem Tode nicht. Also soll man mit betrübten Christen handeln, und die nicht mehr betrüben und anschnauzen, wie Zophar dem leidenden Hiob thut, in dem 11. Kapitel.

\section{Hiob wird von Zophar der Heuchelei beschuldiget}

a Wib) das ist, frei und seines Willens.

\section{Summarium}

Ist die dritte Disputation, Zophars von Naema; der stehet auch darauf, wäre Hiob fromm gewesen, als er sich macht und hören läßt, so hätte ihn Gott nicht so gestraft. Weil ihn aber Gott also straft, so sei es ein gewisses Zeugniß, daß ers verdienet habe, und sei gottlos gewesen, V. 13. Denn Vernunft kann höher nicht kommen, denn daß sie von Gott urtheilet, wie ein Mensch von dem andern pflegt zu urtheilen. Da gehets also zu, wer unrecht thut, der muß entgelten. Wer fromm ist, der genießet. Aber Gottes Urtheil und Gericht ist viel höher. Der läßt den frommen Hiob greulicher strafen, denn alle Gottlosen. Denn er gedenket erstlich, er wolle ihm aus solchem Jammer wieder helfen, und ihn desselbigen reichlich ergetzen. Zum andern will er sehen, ob Hiob auch gerecht und schlecht bleiben, in Leiden nicht ungeduldig werden und von Gott nicht abweichen wolle. Solcher Versuchung bedarf Hiob wohl sammt allen Heiligen.

\section{Gebet}

Der allmächtige, gerechte, gütige GOtt schaffe seinen Willen auch an uns, und helfe, daß wir den erkennen, und alles, was er uns zu leiden zuschickt, mit Geduld annehmen und tragen, und im Glauben auch allezeit bis zum seligen Ende verharren, Amen.

\chapter[12]{12}

\section{Vorrede}

Geliebte Christen! Gott, der HErr, hält lieb und werth die Frommen und Gottesfürchtigen. Daher gebeut er ernstlich in seinem Wort, daß ihm dieselben niemand soll betrüben, beschädigen noch beleidigen. Denn also sagt er im 105. Psalm V. 15.: Tastet meine Gesalbten nicht an, und thut meinen Propheten kein Leid. Also spricht er auch Zach. 2, 8.: Wer euch antastet, der tastet meinen Augapfel an. Derwegen so nehme ihn nur keiner in Sinn, daß es ihm für wohl und ungestraft werde hinaus gehen, der sich an den Christen und Auserwählten Gottes vergreift, und ihnen entweder mit Worten oder Werken zu nahe und wider ist. Denn die bedrängten Christen sagen und klagen Gott ihre Noth; welch ihr Gebet und Seufzen ein unleid-
\quellseite{732}
lich Geschrei ist für den Ohren Gottes, daß ein Mensch lieber die ganze Welt zum Feinde haben sollte, denn ein einziges Christen-Gebet, das mit Ernst wider ihn wird gesprochen. Denn es sind eitel Donnerschläge, was sie in solcher Noth reden, die gewißlich ihre Feinde treffen. Dies soll auch erinnern das 12. Kapitel.

\section[Hiob fängt an sich wider Zophar und die]{Hiob fängt an sich wider Zophar und die andern Freunde zu verantworten}

a Id est, glimmend Docht.

b Das ist: Ihr saget: Weisheit sei bei den Großvätern; ich sage aber, sie sei bei GOtt, welcher gleich aller Könige, Priester, Richter, Gewalt, Kunst, Heiligkeit zu nichts machet.

\section{Summarium}

Ist eine Antwort auf Zophars Rede, in welcher Hiob sagt, er wisse es sehr wohl, daß GOtt allmächtig sei und mit uns thun könne, was er will. Denn solches sehe man an allen Creaturen, V. 6. Aber nichts desto weniger (spricht er) sehe und erfahre ichs, daß auch den Frommen übel und den Bösen wohl geht. Auf die Frage (spricht er) sollt ihr mir eine richtige tröstliche Antwort geben, so wollte ich euch für weise Leute halten, daß andere weiß ich fürwahr. In Summa, so Hiob auf andere Leute stehet, ist er frömmer, denn sie, und dünket ihn zuviel, daß es ihm sollte übler gehen, denn ihnen. Und fehlet ihm durchaus an dem, daß er nicht von Herzen beten noch sprechen kann: Dein Wille geschehe auf Erden, wie im Himmel.

\section{Votum}

Der treue, gütige GOtt helfe und verleihe, daß wir uns seinen väterlichen Willen allezeit wohlgefallen lassen, auf daß wir in christlicher Geduld alles, was er uns auflegt, willig tragen, dankbarlich rühmen und ritterlich überwinden, Amen.

\chapter[13]{13}

\section{Vorrede}

Geliebte Christen! Es ist ja ein unbeständig Ding um das menschliche Leben. Denn gleich wie der Wind im Herbst die kahlen verdorreten Blätter dahin führet; also ists auch bald um einen Menschen geschehen. Desgleichen, wenn der Böttcher ein Faß pichtet, so gehet ein Rauch auf, als wollte er die Sonne verdunkeln; aber gar bald, wenn ein kleines Windlein drein gehet, so wird er vertrieben: also auch wir, wenn uns der kalte Wind des Todes anhaucht, so ist es

44
\quellseite{733}
Alles aus. Und an dem Fluß Hypane findet man Würmlein, welche, wenn die Sonne aufgehet, lebendig werden, so bald sie aber untergehet, so sterben sie: also kurz ist auch unser Leben für Gott, wie eine Nachtwache, Ps. 90, 4., daß wir auch nicht wohl eines Tages gewiß versichert sein. Derwegen soll niemand in Tag hinein sicher leben, und meinen, er habe mit dem Tode einen Verstand gemacht, Jes. 28, 15. Sondern ein jedes soll in Gottesfurcht sich üben und sein Leben desto besser anwenden, weil es so kurz und flüchtig ist, und soll sich befleißigen, daß er einen guten Namen hinter ihm lasse, daß er Gott und den Menschen nützlich in seinem Leben allhie gewesen sei. Darauf höret ferner an das 13. Kapitel.

\section{Hiob fährt in seiner Beantwortung fort}

\section{Summarium}

Dies ist eine Widerlegung deß, daß Zophar dem Hiob hatte fürgeworfen, und ist dieß ungefähr die Meinung: Was ihr mir saget, spricht er, weiß ich fürwahr, darum will ichs nicht annehmen, daß ihr von Gott saget: Er strafe allein die, so es mit ihren Sünden verdient haben. Ihr seid unnütze, verdrießliche Aerzte, V. 4. Man soll die Leute nicht also trösten, so zuvor Gottes Strafe auf ihnen haben. Denn also macht man das Böse nur ärger und des Jammers mehr. Ja, spricht er weiter, ihr selbst, so ihr nichts anders von Gott wisset, denn das allein, daß er ein gerechter Gott sei, werdet ihr solchen Gott nicht können ertragen. Denn seine Gerechtigkeit wird euch zu schwer sein. Ich bin ein armer geplagter Mann, aber die Hoffnung habe ich, Gott werde mein Heil sein und mich nicht halten für unrecht, wie ihr mich beschuldiget.

Auf diese Rede fähret Hiob an zu beten, daß Gott mit seiner Strafe ablassen und nicht so hart gegen ihm handeln wolle, V. 20. Denn er sei ein dürres Blatt; habe er in seiner Jugend gesündiget, soll ers ihm vergeben und gnädig sein, also siehet man, daß in Anfechtung das Gebet

* Beißen) das ist, was soll ich mich viel fasten und mit wehe thun, so ich doch sterben muß und halt mich nicht. Item, meine Seele in die Hände legen, das ist, viel wagen und in Gefahr gehen.
\quellseite{734}
der größte Trost ist, an welches alle Heiligen sich gehalten haben. Und alsdann sind sie auf dem rechten Wege, da sie nicht können unrecht thun, wenn sie sich schuldig geben und Gnade begehren. Merke aber, wie auch der fromme Hiob seiner Jugend halben sich schuldig gibt, V. 26., wie David auch. Denn die Jugend ist zu schwach und das Fleisch zu stark, und kann derselben sich selbst nicht regieren. Darum wer durch solchen Strudel will, muß nicht sicher sein und fleißig beten.

\section{Votum}

O Gott regiere und stärke uns, daß wir alle die Sicherheit ja fliehen fleißig, in rechtem Glauben beten und in wahrem Troste alle Widerwärtigkeit ertragen, bis zum seligen Ende, Amen.

\chapter[14]{14}

\section{Vorrede}

Geliebte Christen! So bald wir mit Weinen, Schmerzen und Traurigkeit auf dieser Welt geboren werden, gehet eines jeglichen Menschen Angst und Noth bald an, und, je länger er lebet, je mehr Trübsal und Elend ihm unter Augen gehet; und wächst das Unglück und Kreuz mit ihm auf täglich, daß wohl recht gesagt wird: Was ist der Mensch jetzt nach dem Fall? Nichts denn Unglück hie und überall, ist doch nichts mehr das elende Leben, als daß wir stets im Tode kleben, so lang es währet. Wohl dem geschieht, der sanft im Frieden hindurch bricht. Dieses alles sollen die Christen in Geduld tragen und in Gottes Wort und bei Gott, durch Gebet, Hülfe und Trost suchen und bedenken, wenn es hoch kommt, daß gleich der Leib darüber zu Trümmern gehet und zu Asche und Pulver wird, so soll er doch herfür am jüngsten Tage und ähnlich werden dem herrlichen verklärten Leibe unsers Herrn Jesu Christ. Dieser Trost macht, daß alles Trauern verschwindet und lindert alle Schmerzen des Leibes, daß es die Christen ertragen können und überwinden. Dieses erinnert auch das 14. Kapitel.

\section{Von des menschlichen Lebens Nichtigkeit}

*Hoffnung) das ist, für dem Tode hat er keine Hoffnung in diesem Leben.

\section{Summarium}

Dies ist eine Klage, in welcher Hiob klaget über den Jammer, welchen alle Menschen auf

44*
\quellseite{735}
Erden haben, daß ihr Leben kurz und voll Unruhe ist, V. 1., und der Mensch nichts gewisser für ihm hat, denn das Sterben, und hernach im Tode liegen ohne alle Hoffnung, daß er wieder möge zu diesem Leben kommen, weil Himmel und Erde stehet.

Darnach bittet Hiob, daß Gott an jenem Tage sein in Gnaden gedenken und alle Sünden vergeben wolle, V. 13. Denn mit diesem Leben hie auf Erden sei es nichts, denn Mühe und Arbeit. Des Menschen Hoffnung, spricht er, ist verloren, das ist, weil wir hie auf Erden sind, können wir nicht hoffen, daß es mit uns besser werden soll. Es muß allerlei erlitten und erstritten und endlich gestorben sein. Gott sei uns gnädig an jenem Tage, wenn Himmel und Erde vergangen sind, Amen.

Dies Zeugniß von der Auferstehung der Todten und einem andern Leben nach diesem Leben, soll man hie fleißig merken. Denn in der ganzen Schrift des Alten Testaments kaum ein klarers ist, und ist das allererste, daß er des ewigen Lebens vermagenden gedenkt, daß er Gottes Gnade begehrt und anruft, und nicht, wie zuvor, auf seine eigene Frömmigkeit stehet, noch derselben sich tröstet. Denn er findet, daß er zu gering zu solcher Wohlthat ist, und dennoch stehet das ganze Papstthum darauf, wir können es mit eigenen Werken und Verdienst erlangen. Aber es scheinet, daß sie in der Schule nicht sind gewesen, da Hiob in ist.

\section{Votum}

Helfe Gott, der allein helfen kann, daß wir den jämmerlichen Zustand unsers Lebens erkennen, den täglich fleißig betrachten, zu täglicher wahrer Buße gebrauchen, in wahrem Glauben endlich einschlafen und am jüngsten Tage zum ewigen Leben mit Freuden auferstehen, Amen.

\chapter[15]{15}

\section{Vorrede}

Geliebte Christen! Ach schrecklich ist es, daß der Mensch ein Greuel und schnöde ist. Denn der Mensch heißt billig ein Greuel darum, daß er wider alle Gebote Gottes sündiget; schnöde und untüchtig aber darum, daß er das Gute nicht thut, so ihm Gott geboten hat. Schrecklich ist es auch, daß der Mensch wegen solcher Sünden ein Kind des Zorns, ein Gefangener des Teufels, ein Gefäß der Unehre und unter dem Fluch und Gottes Zorn ist. Und es sei einer so weise, verständig, hoch und ansehnlich, als er immer wolle, so sind alle seine herrlichen Tugenden, wie ein güldener Harnisch über einem aussätzigen Körper. Sollten doch alle Haare auf unserm Haupte und alle Blutstropfen im Leibe sich darüber entsetzen, daß die Sünde so ein greulicher Greuel ist, und so grausamen Schaden den armen Menschen zugefüget hat.

Und daraus sollen wir lernen Gott fürchten, die Sünde meiden und eine herzliche Abscheu, Furcht, Entsetzung und Erschreckung dafür haben unser Lebenlang. Ach das helfe uns allen der treue Gott, hochgelobt in Ewigkeit. Ferner höret an das 15. Kapitel.

\section{Eliphas will aus Hiob einen Heuchler machen}

a Geringe), das ist, meinst du, daß Gott die Sünder tröste und seinen Trost so geringe hinwerfe, du mußt zuvor fromm werden.
\quellseite{736}

b Brüstet sich scilicet contra Deum, sicut bos pingvis, et crassus, pugnat. D. i. nämlich gegen Gott kämpft er, wie ein fetter und wilder Stier.

\section{Summarium}

Sie fängt Eliphas zum andern mal an, wider den armen Hiob zu disputiren, und spricht: Bist du ohne Sünde, so bist du mehr und höher, denn nie kein Mensch gewesen ist. Denn keiner unter allen Heiligen Gottes ist ohne Tadel, B. 15. Du aber, spricht er, bist ein Heuchler, läßt dich wohl dünken, du seiest fromm, aber die Strafe weiset anders aus. Denn wie Gott allen Gottlosen dräuet, eben also gehests auch dir.

Also siehet man für und für in diesem Buch, daß die Vernunft von Gott nicht anders kann denken noch urtheilen, denn nach dem Gesetze, welches sagt: Gott wolle die Bösen strafen und den Frommen gnädig sein. Wer aber für Gott böse oder fromm sei, und warum Gott auch die Frommen unter das Kreuz werfe, da weiß die Vernunft nichts davon. Darum will Hiob solche Meinung ihm nicht gefallen lassen, wie folget.

\section{Votum}

GOtt der HErr, als ein lieber Vater, stärke und tröste alle betrübte Kreuzträger, und helfe, daß sie der Menschen Spott und des Teufels Anfechtung allezeit mit Geduld überwinden, und regiere uns, daß wir jedermann tröstlich und förderlich sein und bleiben, Amen.

\chapter[16]{16}

\section{Vorrede}

Geliebte Christen! Der liebe Hiob ist unter seinem schweren langwierigen Kreuz mit eitel Todesgedanken umgegangen, und hat nicht hoffen können, daß er daraus werde errettet werden können, dieweil er weder Maß noch Ziel, weder Mittel noch Wege hat mögen ersehen, wie und wodurch ihm geholfen möchte werden; sondern hat nicht anders gewußt, er müßte schon sterben. Also ist allen betrübten und trostlosen Herzen zu Sinn, daß sie gedenken: Nun ists aus. Und wenn sie keine menschliche Hülfe ersehen können, bestimmen sie ihnen selbst die Zeit, da sie nun gewiß sterben und von der Welt hinweg müssen. Dazu kommen denn auch die Gedanken, die der Satan eingibt, dagegen ihnen das Herz im Leibe erbebet und alle ihre Gebeine erzittern. In solcher Angst und Anfechtung sollen ihnen fromme Christen einbilden die Sprüche heiliger Schrift, die uns überzeugen, daß Gott, der HErr, viel reichlicher und überschwänglicher thun kann über alles, das wir bitten und verstehen, nach der Kraft, die da in uns wirket, als St. Paulus lehret Ephes. 3, 20.; und das David im 4. Psalm B. 4. sagt: Erkennet doch, daß der HErr seine Heiligen, das ist, Gläubigen, wunderlich führet. Dazu höret an das 16. Kapitel.

\section{Hiob bezeuget seine Unschuld}
\quellseite{737}

a Zeugen) das ist ihr Befehl wider mich.

b Horn) das ist, meine Gewalt, Macht und Herrschaft, und worauf ich mich verließ.

\section{Summarium}

Ist eine Antwort auf Eliphas Rede, in welcher Hiob nicht sagt, daß er allein gerecht und ohne Sünde sei. Denn seine Sünde hat er mehr denn einmal bekennet und Gott um Vergebung gebeten. Aber das sagt er, sie seien verdrießliche leidige Tröster. Denn mit solchem Trösten machen sie ihm nur des Jammers mehr, V. 1. Fähet derhalben an und machet eine neue Klage von seinem großen Leiden und Schmerzen, darinnen er ist, daß Gott ihn mit so mächtiger Hand zu Boden geschlagen habe, V. 2., und die bösen Leute ihm solches Unglück mit ihrem leidigen Disputiren noch begehren schwerer zu machen; Gott im Himmel weiß meine Unschuld. Aber es will mich gegen diese Heuchler nichts helfen, darum begehre ich zu sterben.

\section{Votum}

GOtt im Himmel, der gerecht ist, und Barmherzigkeit erfordert, steure allen, welche die Betrübten mehr ängsten und betrüben; und verleihe allen Kreuzträgern Trost und Stärke, daß sie ihr Kreuz in Geduld zu Ende bringen; und regiere uns, daß wir in Heiligkeit und Gerechtigkeit leben, Amen.

\chapter[17]{17}

\section{Vorrede}

Bei diesem Kapitel sollen wir merken das Stück, nämlich, wie sich die Frommen sollen, so in langwieriger Krankheit liegen und immer nach dem Tode seufzen und wünschen. Erstlich sollen sie ja eigentlich und gewiß das fassen, daß sie Gott in ihrer großen Angst und Nöthen durch sein Wort und heiligen Geist also herrlich stärken, tröstigen, gründen und erhalten wolle, und also mit der Anfechtung schicken, daß sie es ertragen und ausstehen werden, wie 1 Cor. 10, 13. geschrieben stehet: Gott aber ist getreu, der euch nicht läßt versuchen über euer Vermögen, sondern machet, daß die Versuchung so ein Ende gewinne, daß ihrs könnet ertragen. Und in Klagliedern Jeremiä am 3. Capitel V. 31.: Der HErr verstößet nicht ewiglich, sondern er betrübet wohl und erbarmet sich wieder, nach seiner großen Güte. Denn er nicht von Herzen die Menschen plaget und betrübet. Auch sollen sie sich trösten, wenn dermaleinst das selige Stündlein kommen wird, daß sie aus der Welt werden scheiden sollen, daß sie denn sanft werden dahin fahren und süß im HErrn Christo einschlafen, und daß das letzte Wort unsers HErrn Jesu Christi, so er mit seinem himmlischen Vater am Kreuz geredet hat, auch ihr
\quellseite{738}
letztes Wort sein wird, wo nicht auf ihrer Zunge, doch in ihrem Herzen, nämlich: Vater, in deine Hände befehle ich meinen Geist, Luc. 23, 46. Zum andern sollen sie sich herzlich freuen auf das ewige Leben und auf die himmlische Freude, so kein Auge gesehen, kein Ohr gehört und in keines Menschen Herz kommen ist, Esaia 64, 4. 1 Cor. 2, 9. Und sollen wissen, daß ihre Seele ins Paradies wird aufgenommen werden und bei dem Herrn Christo sein, und daß sie sanft und lieblich ruhen werden in ihren Gräbern, bis die letzte Posaune erschalle: Stehet auf ihr Todten; und daß alsdenn sie dünken wird, als hätten sie kaum eine Viertel-Stunde geschlafen. Zum dritten sollen sie ihnen auch gegen das gegenwärtige Elend und Trübsal dieser Welt die Freude des ewigen Lebens wohl einbilden, daß allda ein Augenblick wird besser sein, denn die hundert tausend Jahre, daß Gott alsdenn die Thränen wird abwischen von allen Angesichtern, Off. Joh. 21, 4. Daß sie die Herrlichkeit des Herrn Christi sehen werden, die er gehabt, ehe der Welt Grund gelegt ward, Joh. 17, 24. Und daß vor ihm Freude die Fülle ist und lieblich Wesen zu seiner Rechten ewiglich, Ps. 16, 11. Und daß Gott der Herr alsdenn seine Gläubige sättigen wird mit langem Leben und ihnen weisen sein Heil, Ps. 91, 16. Das ist, seinen einigen lieben Sohn, in seiner Glorie, in seiner Majestät und Herrlichkeit.

\section{Hiob erzählt sein Elend und Unschuld noch weiter}

\section{Summarium}

Ist eine Klagrede des armen Hiobs, daß er elend, mühselig und von Herzen betrübt sei, und dennoch ein gut Gewissen habe seines Lebens halben, daß er nicht gottlos sei gewesen. Solches, spricht er, können diese meine Tröster nicht verstehen, B. 10. Es thut aber den Frommen sehr wehe, wo es ihnen also begegnet, wie mir geschund. Aber nichts desto weniger werden sie ihnen solches lassen ein Exempel sein, und darum von der Gerechtigkeit nicht abweichen, ob sie gleich auch etwas leiden müssen.

Daß er aber klaget, er müsse in die Hölle fahren, B. 13. 16., hat nicht die Meinung, wie wir im Deutschen verstehen, als müsse er verdammt sein; sondern in die Hölle fahren, heißt sterben, und im Tode oder unter der Erde liegen, da wir alle hingehören. Aber doch durch Christum und Gottes Barmherzigkeit die Hoffnung haben, daß wir auch mit dem Leib wieder auferstehen und ewig leben werden.

\section{Votum}

Der große starke Gott verleihe großen Trost allen, die groß Kreuz tragen, und erhalte sie in wahrem Glauben mit Leib und Seele zur großen Freude des ewigen Lebens; und aber behüte er nach seinem gnädigen Willen für Krankheit und anderer großen Noth und Elend, Amen.

\chapter[18]{18}

\section{Vorrede}

Dies Kapitel lehret, wie Gott die Gottlosen mit allerlei Plage zu strafen pflege. Derwegen, auf daß wir uns Gottes Zorn und Strafe mit unsern Sünden nicht zuziehen, so wollen wir betrachten, wie wir uns für dem listigen Anlauf des Satans hüten sollen, auf daß er uns nicht in Sünde falle und verführe. Dazu aber gehöret endlich, daß man Gottes Wort fleißig höre, lese, betrachte und also täglich mit umgehe. Zum
\quellseite{739}
andern sollen wir anhalten mit dem Gebet, ernstlich und emsig, und sollen mit Mund und Herzen bitten: Ach Herr, führe uns nicht in Versuchung, wehre du dem leidigen Teufel, erhalte du uns, daß er uns nicht in Sünde und Schande leite und stürze. Zum dritten soll ein jedes wohl bedenken den ernsten Zorn Gottes, und wie schrecklich er den Sündern mit allerlei Plage und Strafe an Leib und Seele dräuet, 3 Mos. 26, 14., 5 Mos. 28, 15. Zum vierten sollen wir auch oft an die Exempel gedenken derer, die von Gott greulich sind gestraft worden, und jämmerlich dahin gegangen, von wegen ihrer großen und schweren Sünden. Zum fünften, wer der Sünde entgehen will, der muß auch die Oerter, Personen und Gelegenheit meiden, dadurch ihm der böse Geist nachstellet. Das sind die Mittel, dadurch uns Gott der Herr wider die List und Gewalt des schädlichen Bölands, des Teufels, bewahret. Es muß aber auch gekämpfet sein wider die Sünden, und muß ihren bösen Reizen widerstanden sein.

\section{Bildad mischet Hiob unter die Gottlosen}

a Versetzt) das ist, Gott wird mit dir nicht anders machen, denn mit allen andern, und seine Weise nicht lassen um deinetwillen.

b Fürst) das ist, die Macht und Gewalt des Todes. Also auch König des Schreckens ist die Gewalt des Schreckens, daß er muß unten liegen und nicht entrinnen kann.

c Wurzeln) Wurzel heißt er alles, was in der Erde gepflanzet ist. Ernte alles, was oben auswächset, es sei Korn, Oel, Wein etc.

\section{Summarium}

Ist eine Disputation des Bildads wider Hiob, in welcher er erzählet, wie Gott die Gottlosen pflege zu strafen mit allerlei Plage. Will also anzeigen, weil eben solche Plagen dem Hiob auf dem Nacken liegen, so sei er auch gottlos und leide solches alles um seiner Sünden willen.

\section{Votum}

Wir sagen Gott Lob und Dank, der uns diese Mittel in seinem Wort offenbaret hat; der verleihe, daß wir sie auch allezeit gebrauchen, in seiner Furcht leben, das Gute thun und das Böse lassen, und zeitlichen und ewigen Strafen entfliehen, Amen.

\chapter[19]{19}

\section{Vorrede}

Geliebte Christen! Dieweil es einem wohl gehet, so hat er viel, die sich seiner annehmen; so bald sich aber das Glück wendet, so fliehen sie davon, wie die Schwalben und etliche andere Vögel Winterszeit. Wenig sind, die in der Noth und Anfechtung bestehen. Unserm Herrn Christo selbst ist das begegnet. Denn da es an die rechte
\quellseite{740}
Noch ließen ihn seine Jünger allein in der Feinde Gewalt. St. Paulus beklagt sich auch, da er für Nero gestellet war, daß ihn seine Freunde, auch Demas, verlassen haben, 2 Tim. 4, 10. Das thut trefflich weh, sonderlich wenn wir uns unserer Freunde viel getröstet und viel auf sie gebauet haben. Aber es läßt uns Gottsolches auch zum Guten geschehen. Denn wir geben etwa den Menschen zu viel zu, so müssen wir erfahren mit unserm Schaden, wie wankelmüthig sie sind, und lernen, daß wir uns Gottes allein trösten sollen, und ihn für unsern Felsen, Schloß, Schild, Schutz und Schirm halten. Wahr ist’s, man muß mit einander haushalten und einander trauen, wie man sagt: Wenn alle Tage auch ein Mord geschähe; aber trauen soll man einem Menschen als einem Menschen. Dies soll auch erinnern das 19. Kapitel.

\section[Hiob redet von seinem Elend, und seinem Trost]{Hiob redet von seinem Elend, und seinem Trost von der Auferstehung}

a Ehre, Krone, Hoffnung ist alles geredet vom zeitlichen Leben in guter Ruhe.

b Satt werden) das ist, könnt nicht aufhören mich zu beißen und zu fressen.

c Erlöser, Retter) Vindex, Quia Christus vindicatos, contra Homicidam nostrum Diabolum. D. i. Rächer, weil Christus uns rächet an unserm Mörder, dem Teufel.

\section{Summarium}

In die Antwort auf Bildads Rede, in welcher Hiob wiederum klaget, daß seine Freunde mit solchen Reden ihm das Leiden nur schwerer machen. Gott (spricht er V. 6.) brauchet solche Macht wider mich, das muß ich leiden und dazu still schweigen. Denn was kann ich wider Gott
\quellseite{741}
ausrichten? Hebt also wieder an und klaget seine Noth. Am Ende aber tröstet er sich über die Maßen sein: wenn schon hie auf Erden muß gelitten und gestorben sein, so weiß ich (spricht er B. 25.) doch, daß mein Erlöser lebt, der wird mich an jenem Tage aus der Erde und dem Tod, zu einem andern bessern und ewigen Leben auferwecken. Auf solchen Trost wollten ihn seine Freunde nicht weisen, darum tröstet er sich selbst also. Uns zum Exempel, daß wirs ihm auch nachthun sollen. Eben auf solche Weise vermahnet Paulus Timotheum, 2 Tim. 2, 8.: Halt im Gedächtnis Jesum Christum, der auferstanden ist von den Todten, aus dem Samen David, nach meinem Evangelio, über welchem ich mich leide bis an die Bande. Denn weil es doch beschlossen ist, daß die Gottseligen dem Exempel Christi, des Sohnes Gottes, nach, hie auf Erden ihr Kreuz tragen und sich leiden müssen, ist dies der einige Trost, daß wir mit Christo wieder vom Tod auferstehen und ewig leben sollen.

\section[Notum]{\unsicher{Notum}}

Lob, Ehr, Preis und Dank sei GOtt für diesen Trost; der verleihe, daß er uns und allen unter dem Kreuz ein beständig Labsal und rechte Erquickung sei; und helfe, daß wir christlich leben und selig sterben, und mit Freuden am jüngsten Tage auferstehen, Amen.

\chapter[20]{20}

\section{Vorrede}

Geliebte Christen! Die Strafe stimmet gemeiniglich mit der Sünde überein, wie das Buch der Weisheit im 11. Kapitel V. 17. sagt: Womit jemand sündigt, damit wird er auch geplagt. Also, Antiochus Epiphanes hat viele fromme Herzen betrübet und beleidiget, darum mußte er endlich selbst für Trauern sterben, 1 Macc. 6, 13. Also, Attila, der Hunnen König, hatte viel Bluts vergossen, darum fraß ihn das strenge Gericht Gottes, daß er in seinem eigenen Blute erstickte. Also auch, es trägt sich oft zu im Ehestande, daß mancher ein böses Weib hat, manche wiederum einen bösen unartigen Mann kriegt, daß ist beiden eine Strafe der Sünden, die sie zuvor, ehe sie in Ehestand getreten sind, begangen haben. Eben also gehts den Geizigen, wenn sie den Armen ihr Gut auffressen und ihren Schweiß und Blut aussaugen, so beraubet sie Gott auch wiederum des Ihren, Jer. 30, 16. Derwegen hüte sich ein jeder für Sünden, daß ihm auf sein Sündendienst nicht Sündenlohn folge. Dies lehret auch das 20. Kapitel.

\section{Zophars Rede von der Gottlosen unbeständigem Glück}

a Wenn er Bosheit ansähet, hat er Wollust und Ruhe; aber es wird nicht währen, wird bald bitter schmecken.

b Saugen) das ist, er wird tödtliches Herzleid und Jammer leiden, und alles Guten beraubet werden.
\quellseite{742}

c Ausgeblasen) das ist, ein Feuer von Gott angezündet, nicht durch Menschen ausgeblasen.

\section{Summarium}

Ist eine Disputation des Zophars wider Hiob. Dem Gottlosen (spricht er) gehet es wohl, aber nicht lange, B. 5. Ursache, wenn sie niemand mit dem Herzen helfen und alles mit ihrem Gebet zu sich reißen, kanns Gott in die Länge nicht leiden, sondern macht sie wieder zu Bettlern, B. 19. Also, spricht er, gebets jetzt dir auch; wärst du fromm gewesen, es sollte dir wohl länger sein wohlgegangen; weil du aber gottlos bist, bleibet das Sprichwort wahr: Des Gottlosen Ruhm stehet nicht lang, und die Freude des Heuchlers währet ein Augenblick. Solche Einrede nimmt Zophar aus der Erfahrung. Denn gewiß ist es, daß Gott die Geizigen wie einen vollen Schwamm wiederum ausdrücken und sie leeren wolle, wie das Sprichwort heißt: Uebel gewonnen, übel zerronnen. Aber die Welt glaubets nicht, bis sie es erfähret.

\section{Votum}

GOtt verleihe, daß alle Geizige und andere Sünder sich bekehren; steure allen, die wider ihn und die Liebe ihres Nächsten handeln, und helfe uns, daß wir der Gerechtigkeit nachjagen und andern gerne und treulich dienen, Amen.

\chapter[21]{21}

\section{Vorrede}

Geliebte Christen! In diesem Leben gehet es einem frommen Christen so bald übel, als einem gottlosen und bösen Menschen. Derwegen ist es ein großer epicurischer Irrthum, daß man schließen will: Siehe, der lebet, wie er auf Rosen ging, derwegen muß er gewiß einen gnädigen Gott haben. Denn, wenn Gott das Böse allwege alhie zeitlich strafe, so bedürfte er des jüngsten Gerichts nicht. Dieweil es aber die Gottlosen oft alhie auf Erden am besten haben, und die Gläubigen dagegen viel Böses leiden müssen; so können vernünftige Leute daraus abnehmen und verstehen, daß ein letztes Gericht und nach diesem Leben eine Belohnung und Strafe sein müsse. Dazu grünen eine Zeitlang wohl die Gottes-Verächter in diesem Leben und stehen in voller Blüthe, aber endlich gehen sie oft auch allhier unversehens zu Grund und Boden. Die Christen aber, wie elend und voller Jammers sie sind, so haben sie doch vor sich die herrliche Erlösung und das ewige Leben. Das soll die Gottlosen schrecken, daß sie Buße thun, und das soll die Gläubigen trösten, daß sie ihr Kreuz in Geduld tragen. Davon lehret auch Hiob im 21. Capitel.

\section[Hiob widerleget die Rede Zophars von dem Glück]{Hiob widerleget die Rede Zophars von dem Glück und der Strafe der Gottlosen}

a Augenblick) das ist, sie leben bis an den Tod wohl und da ist um einen bösen Augenblick mit ihnen zu thun, so sind sie hindurch. Ich aber muß so lange Zeit Schrecken und Unglück leiden.
\quellseite{743}

b Vergelten) das ist, wer kanns urtheilen, was ihm zu vergelten sei, ohne Gott allein?

c Haufen) das ist, es ist ihm ein Grab bereit, unter andern Gräbern.

d Der Schlamm des Bachs) das ist die Wollust im Fleisch und guten Leben.

\section{Summarium}

Ist eine Antwort auf Zophars Urtheil, daß es unrecht sei. Ursache, die Gottlosen haben hie auf Erden in diesem zeitlichen Leben einen Vortheil, daß es ihnen lange wohl gehet, V. 7. Die rechte Strafe aber, die sie verdienen (spricht Hiob), ist nicht das zeitliche Unglück, sondern der Tod und die ewige Verdammniß, V. 17. Darum darf man nach dem, wie es ihnen hie auf Erden gehet (es sei böse oder gut), nicht urtheilen, wie ihr thut, wenn mirs übel gehet, so sprecht ihr, ich sei gottlos gewesen. Ich aber sehe wohl, daß das Widerspiel wahr ist, nämlich, daß die Gottlosen nicht so hart gestraft werden, als ich. Aber ihre verdiente Strafe gehöret in ein anderes Leben, davon ihr nichts wisset. Dies ist wiederum ein helles Zeugniß vom ewigen Leben und Auferstehung der Todten.

\section[Notum]{\unsicher{Notum}}

Der treue GOtt helfe, daß wir von allen Dingen nach seinem heiligen Wort urtheilen lernen! Erhalte alle Betrübte unter ihrem Kreuz zur ewigen Freude, und bestätige uns in der Frömmigkeit, zu seines heiligen Namens Lob, Ehre und Preis, Amen.

\chapter[22]{22}

\section{Vorrede}

Geliebte Christen! Wenn wir uns wegen anderer Leute Unglück freuen, aus Neid, Nachgier und anderer fleischlicher Anfechtung, so ist es unrecht und verboten. So es aber geschieht um Gottes und des Nächsten Willen, so ist es recht und zugelassen. Denn viel Heilige Gottes, wie die Schrift bezeuget, sich nicht allein gefreuet haben, wenn es den Feinden Gottes übel gegangen, sondern sie haben auch gebeten, daß er sie strafen oder gar ausrotten wolle, 2 Mos. 15, 4. 1 Sam. 25, 39. Spr. Sal. 11, 5. Das dienet aber also zur Ehre Gottes, so jedermann siehet, daß Gott zu seinen Dingen schauet, daß er seine Gewalt noch nicht von Händen gegeben, daß er seine Feinde wohl dämmen und seine Gerechtigkeit erzeigen könne. Es dienet dem Nächsten und den Einfältigen zu gut, die sich an der Glückseligkeit des Gottlosen heftig ärgern, ungeschickte Reden vermögen fahren lassen und ihrem bösen Exempel ansehen zu folgen, oder sonst ansehen zu schlottern, daß sie sich eines andern und bessern bedenken. Es dienet auch den bösen Leuten, die gestraft werden, zu gut, daß sie nur menschliche Vernunft haben, daß sie aufhören zu sündigen, da sie nicht mehr und härter sollen gestraft werden. Dies erinnert auch das 22. Kapitel.
\quellseite{744}

\section{Eliphas ziehet Hiobs Frömmigkeit abermal in Zweifel}

a Finsterniß heißt Trübsal und Unglück, wiederum Licht heißt Glück und Heil.

\section{Summarium}

Ist die dritte Disputation, in welcher Eliphas eben die vorige Meinung wiederholet, daß Gott den Frommen gnädig sein, sie behüten und ihnen helfen wolle, wiederum die Gottlosen strafen. Weil nun die Strafe auf dem Hiob öffentlich liegt, so soll ers als verdienten Lohn annehmen, V. 6., seine Sünde bekennen und sich bessern, so werde ihm geholfen werden.

Auf diese Weise muß die Vernunft von Gott urtheilen und reden, anders kann sie nicht, und ist zwar nicht ungerecht geredet; doch so fern, daß man beides bekennet, daß wir nicht recht können fromm sein, und daß Gott wolle barmherzig sein und Sünde vergeben, wenn mans begehrt und hoffet, durch Christum.

\section{Votum}

Unser treuer Heiland JEsus Christus, Gottes und Mariä Sohn, helfe, daß wir unserer Gebrechlichkeit uns täglich erinnern, in wahrem Glauben uns seiner wider unsere Sünde trösten, daß wir uns auch ferner für allen Sünden hüten, und durch ihn endlich selig werden, Amen.

\chapter[23]{23}

\section{Vorrede}

Geliebte Christen! Daß man unter dem Kreuz geduldig sei, wider Gott nicht murre und seiner gnädigen väterlichen und der nöthigen Hülfe
\quellseite{745}
erwarte, so soll ein jedes zweierlei betrachten, glauben und hoffen: Erstlich soll ein jedes herzlich bekennen, daß es sein Kreuz mit seinen Sünden wohl verdienet habe, wie Judith lehret Kapitel 8, V. 22. und saget: Wir sollen bekennen, daß es eine Strafe ist von Gott, viel geringer, denn unsere Sünden sind. Ja wir sollen bekennen, daß wir auch die größten zeitlichen, und darnach auch die ewige Strafe verdienet hätten, wenn Gott mit uns ins Gericht gehen wollte, und mit uns handeln nach unserm Verdienst. Wie wir denn solches alles bekennen müssen, wir haben auch so heilig gelebt, als wir immer vermochten. Fürs andere soll ein jedes sich herzlich trösten, daß Gott wieder helfen wolle und uns endlich von allem Uebel erlösen, wie er selber zusagt Jes. 54, 7.: Ich habe dich einen kleinen Augenblick verlassen, aber mit großer Barmherzigkeit will ich dich sammeln. Ich habe mein Angesicht im Augenblick des Zorns ein wenig von dir verborgen, aber mit ewiger Gnade will ich mich dein erbarmen. Dieses hat auch Gott der heilige Geist den königlichen Propheten David gelehret, derwegen er also im 30. Psalm V. 6. sagt: Des Herrn Zorn währet ein Augenblick, und er hat Lust zum Leben, den Abend lang wäret das Weinen, aber des Morgens die Freude. Dies soll uns erinnern das 23. Kapitel.

\section[Hiob beruft sich um seines guten Gewissens willen]{Hiob beruft sich um seines guten Gewissens willen auf Gottes Richterstuhl}

a Einig also Gal. 3, 20. Gott ist einig. Des Einigen aber ist kein Mittler.

\section{Summarium}

Hie ist eine Antwort auf Eliphas Rede, in welcher Hiob auf diesen zweien Puncten beharret. Erstlich, er sei fromm und habe solche Strafe nicht verdienet, V. 6. 11. Zum andern, daß er aber solche Strafe tragen müßte, da könnte er nicht vor sein, Gott sei allmächtig, der mache es mit den Menschen, wie er wolle.

Solche Gedanken zeiget nicht allein Ungeduld an, sondern ist auch Gott in dem zunahend, als handelte Gott mehr nach Macht, denn nach Recht und Billigkeit.

Darum fehlet es dem frommen Hiob in seiner Anfechtung an diesen zweien Stücken. Das erste, daß er die sündhafte Natur nicht eigentlich erkennt. Denn obgleich andere Leute dem äußerlichen Leben nach ärger sind, denn Hiob, so ist doch Hiob dem Urtheil Gottes damit nicht entbrochen. Denn das sündhafte Fleisch und Blut ist eben sowohl in ihm, als in andern Leuten. Das andere, daß er sich zu Gott nicht versiehet, daß er werde wieder helfen und gnädig sein.

\section{Gebet}

Der treue GOtt erleuchte unsern Verstand, und erweiche unsere Herzen, daß wir unter dem Kreuz unsere Sünden erkennen, uns demüthigen und seiner Hülfe uns trösten und freuen, bis zur seligen Erlösung; auch gläubig und fromm sein und bleiben allzeit, Amen.

\chapter[24]{24}

\section{Vorrede}

Geliebte Christen! Dieweil Gott unser trot mit der ewigen wohlverdienten Strafe, um seines lieben Sohns, unsers Herrn Jesu Christi willen, verschonen will, sollen wir uns nicht weigern, allhier ein wenig Kreuz zu leiden, sondern sollen mit dem Propheten Micha sagen, Mich. 7, 9.: Ich will des Herrn Zorn tragen, denn ich habe wider ihn gesündiget; und sollen gedenken, daß es besser sei, hie leiden, denn in jenem Leben, wie St. Augustinus betet: O Herr baue hie, schlage hie und schone dort. Wie auch Nicephorus vom Kaiser Mauritius schreibt, daß
\quellseite{746}
er solches von Gott gebeten habe, und da ihm vor seinen Augen der Tyrann Phocas sein Gemahl und Kinder umbringen ließ, und er nun auch herhalten sollte, hat er gesprochen: Herr Gott, du bist gerecht, und dein Gericht ist auch recht, Ps. 119, 137. Und derwegen sagt St. Ambrosius: Das ist eine herrliche That, daß du in Schmach Gottes Urtheil lobest, in Krankheiten seinen Willen preisest, und in Armuth seine Gerechtigkeit rühmest. Ferner höret an das 24. Kapitel.

\section[Von Gottes verborgener, doch gerechter Regierung, der Frommen]{Von Gottes verborgener, doch gerechter Regierung, der Frommen Elend und der Gottlosen Glück betreffend}

a) Die Zeiten) weil Gott die Bösen so lässet machen, wie sie wollen, so scheinet es, als wüste er nichts darum. Weil ihr denn saget, er strafe die Bösen und nicht die Frommen, so müsset ihr zugeben, daß er nicht wisse und die ihn kennen, auch nicht wissen, zu welcher Zeit er strafen werde, wie ihr euch rühmet zu wissen.

b) Das Wild) die freien, frechen Leute und Tyrannen.

c) Bauet) das ist, die der Hurerey nachgehen, bringen ihr Gut um und lassen ungebauet.

d) Auf ihr Thun) das ist, daß sie nicht ein Aufruhr wider ihn machen, dämpfet er sie immerdar, und muß also Sicherheit mit List suchen, aber es währet nicht.

\section{Summarium}

Dies ist eine Entschuldigung des armen Hiobs, in welcher er anzeigt, er sei kein Tyrann, B. 2., kein Mörder, B. 14., kein Ehebrecher, noch dergleichen Sünder gewesen, B. 15. Welche, wenn es ihnen gleich übel abgeht, haben sie es wohl verdienet, B. 20. Beruhet also immer darauf, er habe solche Strafe nicht verdienet. Das ist nun wahr, wenn man auf der Menschen Urtheil sehen will. Aber Gottes Urtheil ist höher, da sind wir allezeit bußfällig, ob uns gleich die Menschen nicht können beschuldigen.
\quellseite{747}

\section{Votum}

Der barmherzige GOtt strafe uns nicht, wie wirs verdienet haben, und helfe, daß wir das Kreuz in Geduld tragen, seine Werke allzeit rühmen und preisen, und uns, wie Hiob gethan, für allerlei Sünden fleißig hüten, Amen.

\chapter[25]{25}

\section{Vorrede}

Geliebte Christen! Ob sich schon die Christen hüten vor äußerlicher Sünde, Schande und Laster: dennoch, damit sie nicht möchten gedenken, sie wären ganz und gar ohne Sünden, und also von sich selbst gerecht, so schicket ihnen Gott, der HErr, Kreuz und Leiden zu, welches Bestrafungen oder Belohnungen der Sünden sein, Röm. 6, 23. Damit sie ihre Erbschuld und Erbsünde erkennen, darin sie empfangen und geboren und von des wegen auch von Natur Kinder des Zorns Gottes sind. Wie davon Gott, der HErr, sagt Jerem. am 30. Kapitel V. 11.: Züchtigen will ich dich mit Maße, daß du dich nicht unschuldig hältst. Derwegen sollen wir abermal unter dem Joch des Kreuzes, und unter der Last der Anfechtungen unsere Seele mit Geduld fassen, und sollen Gott recht geben, wenn er uns züchtiget und mit dem Kreuze hart drücket, und mit den Anfechtungen peinlich ängstet, und sollen mit dem lieben David bekennen, und sprechen Ps. 51, 6.: An dir, HErr, allein hab ich gesündiget, und übel vor dir gethan, auf daß du Recht behaltest in deinen Worten, und rein bleibest, wenn du gerichtet wirst. Dieses sollen uns erinnern das 25. und 26. Kapitel.

\section{Vor GOtt sind alle Menschen Sünder}

a Wer soll dir thun? Gott ist allmächtig und kann wohl steuern den Großen, wenn du nur fromm wärest. Und du meinest, er wisse es nicht, wie du jetzt neulich gesagt hast.

\section{Summarium}

Ist eben einerlei Meinung auf die Einrede, die oben Bildad that, Gott ist allein gerecht und rein, für ihm kann sich niemand seiner Reinigkeit noch Gerechtigkeit rühmen. Nein, spricht Hiob, ich bin gerecht, denn ich habe nicht so gottlos gelebet, wie andere. Daß ich aber also leiden muß, geschieht darum, daß Gott thut, was er will, das muß ich geschehen lassen, leiden, und dazu still schweigen. In Summa, es fehlet dem frommen Hiob in seiner Anfechtung an dem, daß er sich ansiehet, wie er gegen andere Leute sei, und vergißet daneben, wie er gegen Gott sei, daß er solchen Gehorsam nie habe geleistet, wie ihn Gott gefordert. Soll derhalben sich schuldig geben und nicht sagen: Gott thut mit mir, was er will. Sondern er thut mir, wie ich verdienet habe, und soll mich noch ewig verdammen, wenn er meinem Verdienst nach und nicht nach seiner Gnade mit mir wollte umgehen. Aber die arme schwache Natur kann in der Anfechtung solches nicht bedenken. Bedenket sie es aber recht, so muß sie verzweifeln, wo Gott nicht hilft.

\section{Votum}

GOtt, der heilige Geist, der in alle Wahrheit leitet, lehre und regiere uns, daß wir unser Kreuz recht ansehen und verstehen, und christlich tragen, und behüte uns und alle betrübte Kreuzträger für Verzweiflung und allen andern Sünden, wider GOtt und wider die Menschen, Amen.

\chapter[26]{26}

\section{Hiob preiset GOttes Majestät viel herrlicher, als Bildad}

a Riesen die großen Wallfische, welche bedeuten, die großen Tyrannen auf Erden.
\quellseite{748}

\chapter[27]{27}

\section{Vorrede}

Geliebte Christen! Der gerechte und starke Gott dräuet Malachiä am 3. Kapitel V. 5.: Ich will zu euch kommen, und euch strafen, und will ein schneller Zeuge sein wider die, so mich nicht fürchten. Und St. Paulus vermahnet Röm. 11, 20. und sagt: Sei nicht stolz, sondern fürchte dich. Hat Gott der natürlichen Zweige nicht verschonet, daß er vielleicht dein auch nicht verschone. Durch diese Sprüche werden wir vermahnet, daß wir uns für Sünden hüten sollen, auf daß wir mit unsern Sünden Gott nicht verursachen, daß er uns strafen müsse, und daß mit Gottes Strafe unser eigen Gewissen begangener Sünde wegen uns auch nicht nage und plage. Denn wie ein gut unverletzt Gewissen ist der edelste beste Schatz und Kleinod auf Erden und ein stetiges Paradies und Wohlleben: also ist wiederum keine greulichere, größere, schmerzlichere und ängstlichere Plage, als ein böses Gewissen, ja es ist die Hölle selber. Dieses soll uns erinnern das 27. Kapitel.

\section[Hiob hält, seine Unschuld zu retten, den Frommen]{Hiob hält, seine Unschuld zu retten, den Frommen und Heuchler gegen einander}

a Heuchler heißen in diesem Buch allenthalben einen falschen Menschen, wie sie alle sind, für Gott ohne Glauben.

b Weinen) sie werden froh werden, daß er todt ist.

\section{Summarium}

Nachdem die Freunde Hiob ausgeredet und ihn immer beschuldiget haben, er habe solche Strafe verdienet, darum soll er sich nicht für fromm halten, sondern vielmehr seine Sünde bekennen und Gott um Gnade bitten. Darauf will Hiob wider sie schließen und sagt: Sie thun ihm unrecht; er habe solche Strafe nicht verdienet, wie es die Gottlosen verdienen, welche Gott mit mancherlei Plagen pflegt zu strafen,

45
\quellseite{749}
N. 12. Weil nun solche Strafen, wie Hiob nach der Länge erzählet, über die Gottlosen sollen kommen, so soll hie billig jedermann sich für dem gottlosen Leben hüten und fromm sein. Denn obgleich das Kreuz und Unglück auch über fromme Leute kömmt, ists doch träglicher, denn ein böses Gewissen samt der Strafe haben.

\section{Votum}

Regiere und behüte uns der treue GOtt, daß wir nicht muthwillig, wissentlich und vorsätzlich sündigen, und da er uns nach seiner Gerechtigkeit begangener Sünden wegen strafte, so verleihe er uns gnädiglich Trost und Freude, unter aller Strafe bis zum seligen Ende alles Uebels, Amen.

\chapter[28]{28}

\section{Vorrede}

Geliebte Christen! Daß Gottes Gnade, Treue und Liebe gegen uns groß sei, kann man daran auch erkennen, daß er uns seine Weisheit durch sein Wort und Geist offenbaret und den theuren Schatz mittheilet. Gedenket aber, wenn ein Fürst einem seiner Unterthanen seine Rathschläge offenbarte, er würde es traun groß und hoch achten; also wenn ein König oder Herzog einen Ring von seinem Finger zöge, darin ein Türkis, Sapphir oder Smaragd versetzt wäre und den einem unter uns gebe, wie herrlich und hoch würden wir den halten? Nun stellet uns Gott in seinem Wort seinen Willen und die rechte Weisheit für, gegen welcher alles Gold, Silber, Edelgestein auf Erden nichts ist. Wer danket aber Gott für solchen großen Schatz? Wer trachtet darnach so fleißig, als wir sollten? Diese große Undankbarkeit läßt uns erkennen und läßt uns allen Fleiß anwenden, daß wir wahre Weisheit erlangen, auch mehren und behalten. Welches geschehen wird, so wir ernstlich Gottes Wort, dadurch sie gepflanzet wird, lernen und betrachten, und da wir Gott, der sie allein gibt, darum bitten werden. Denn also sagt der heilige Apostel Jacobus im 1. Kapitel seiner Epistel, V. 5.: Welchem Weisheit mangelt, der bitte sie von Gott, der da gibt einfältiglich jedermann, und rückets niemand auf, so wird sie ihm gegeben werden. Darauf höret auch an das 28. Kapitel.

\section{Lob der Weisheit an GOtt und Menschen}

a Finstern) das ist, man gräbet zuletzt so tief, daß man findet, das verborgen liegt in Finsterniß der Erde.

b Stolzen Kinder) das sind junge Löwen.
\quellseite{750}

\section{Summarium}

Ist eine feine Lehre, daß nichts so tief unter der Erde und im Meer verborgen sei, man könne ihm nachgraben und es finden, aber die rechte Weisheit könne man nirgend finden, denn allein im Wort GOttes. Dasselbige lehret, wie die rechte Weisheit nichts anders sei, denn GOtt fürchten und das Böse meiden. Solche Weisheit (will nun Hiob zu seinen Freunden sagen) habt und wisset ihr nicht; ich aber weiß und habe sie, darum dürft und könnt ihr mich nicht lehren.

\section[Notum]{\unsicher{Notum}}

GOtt, der du bist der Brunnquell aller Weisheit, gebe uns allen auch durch sein Wort und Geist die rechte Weisheit, und helfe, daß wir hier Sünder uns unser Lebenlang hüten, und in seiner Furcht gehorsamlich allzeit leben, bis zur seligen Heimfahrt aus diesem elenden Leben, Amen.

\chapter[29]{29}

\section{Vorrede}

Geliebte Christen! Der Mann Gottes, Moses, vermahnet alle Welt im 13. Kapitel B. 4. seines 5. Buchs und sagt: Ihr sollt dem HErrn, eurem GOtt, folgen, und ihn fürchten, und seine Gebote halten, und seiner Stimme gehorchen, und ihm dienen, und ihm anhangen. Dazu setzet dies St. Paulus 2 Tim. 3, 12.: Die gottseelig leben wollen in Christo Jesu, müssen Verfolgung leiden. Denn der HErr Christus lehret selbst also Matth. 10, 12.: Wer nicht sein Kreuz auf sich nimmt und folget mir nach, der ist mein nicht werth. Da hören wir, daß wir alle gegen GOtt und gegen der Welt christlich und ehrbar leben sollen, und da das liebe Kreuz uns dennoch darauf folget (wie denn alle Christträger oder gläubige Christen Kreuzträger sein müssen), daß wirs in Geduld auf uns nehmen und tragen sollen; in Betrachtung auch dieses, daß der Sohn GOttes, unser HErr Jesus Christus, den größten vollkommenen Gehorsam seinem himmlischen Vater in Thun und Leiden, nach seinem heiligen höchsten Willen geleistet, und hat dennoch das allergrößte, schwerste und längste Kreuz an Leib und Seele tragen müssen, und hat solches gethan in tieffer Demuth und unaussprechlicher Geduld. Dies soll auch erinnern das 29. Kapitel.

\section{Hiob erzählt seine vorige Glückseligkeit}

a In Butter) da ich alles übrig genug hatte, alles fett und vollauf.

b Bogen) das ist, meine Macht nahm immer zu.
\quellseite{751}

c Anlachte, freundlich, fröhlich mit ihnen war, wurden sie darum nicht kühn, mich zu verachten. Id est, familiaritas mea non peperit apud eos mei contemptum.

\section{Summarium}

Dies ist ein Klagrede, in welcher Hiob anzeigt, in was Stand und Wesen er gewesen sei, ehe er in das Unglück kommen sei. Und ist dies Kapitel ein fein Bild, weiß sich ein jeder, der da gottesfürchtig leben will, in seinem Leben befleißigen soll.

Zum andern ist auch ein Exempel, weil Hiob so fromm gewesen, und ihn dennoch Gott unter das Kreuz und Leiden geworfen hat, daß wir uns des Kreuzes auch nicht beschweren sollen, die wir doch mancherlei Ursache dazu geben. Denn, daß wir nicht allweg unser Theil nach dem, wie es äußerlich gehet, richten sollen. Sonst werden wir oft fromme Leute für böse halten und ihnen unrecht thun.

\section{Gebet}

Helfe GOtt, daß wir alle in seiner Furcht leben und wandeln, wie es eines jeden Beruf erfordert; und verleihe, daß ein jeder unter seinem Kreuz geduldig sei, und sich demüthige, bis GOtt alle aus zeitlicher Leibes- und der Seelen-Noth väterlich errettet und helfe, Amen.

\chapter[30]{30}

\section{Vorrede}

- Geliebte Christen! Es ist eine schwere Versuchung, wenn eines verneinet: Gott erhöre sein Gebet nicht. Denn er würde ihm sonst das benehmen, darum er ihn lange Zeit Tag und Nacht gebeten hat. Aber da soll man wissen, daß Gott alle, die aus wahrem Glauben zu ihm rufen, erhöre. Denn wir haben dessen vielfältige Verheißungen, und auch tröstliche Exempel. Daß er aber dergleichen thut, als ob er uns nicht erhöre, und uns gerade nicht bewilligt, daß uns an Leib und Seele beschwerlich ist; oder uns nicht gibt, daß wir von ihm begehren, das geschieht darum, daß er unsern Glauben, Hoffnung und Geduld üben will; und daß wir auf ihn warten und ihm nichts vorschreiben sollen: sondern gewiß sein, daß er zu rechter Zeit kommen und in der Weile uns nichts versäumen und verderben werde. Denn wie ein Goldschmied, der das Silber in das Feuer setzt, dabei stehet und nicht, so bald es anhebt zu schmelzen, aufhöret, am Blasebalg zu ziehen, sondern weiß, wenn es genug ist, und siehet, daß das Feuer nicht zu groß, noch zu klein sei, sondern eben recht; also ist Gott auch bei allen unsern Trübsalen, und thut allen Dingen recht zu seiner Zeit, und wenn es nur zu seinen Ehren gelanget und den Menschen nützlich ist. Dies erinnert auch das 30. Kapitel.

\section{Hiob klaget über sein Elend und Anfechtung}

a Sie die Chaldäer.

b Id est, deposuerunt, privarunt curru et aurigata, id est, dominio meo. Das heißt: abgesetzt, des Wagens und der Zügel, d. i. meiner Herrschaft beraubt.
\quellseite{752}

c) Gekleidet) das ist, mancherlei Unglück wird mir angethan gewaltiglich, daß ich mich nicht erwehren kann und gürtet mich, daß ich nicht heraus kommen kann, und muß anhaben, wie den Rock am Halse.

d) Beinhaus) das ist, im Beinhause werde ich Ruhe haben.

\section{Summarium}

Sie ist eine Klagrede, in welcher Hiob, nachdem er erzählt hat, wie wohl und glücklich es ihm vorher gegangen und wie fromm er gewesen sei, klagt, wie es jetzt alles mit einander umgekehret sei, und er solches gottfürchtigen Lebens nicht genieße. Solches, spricht Hiob, widerfähret mir darum, daß Gott nicht mehr mein Freund ist, sondern handelt gegen mir, wie ein Grausamer, wider welches Gewalt ich mich nicht kann schüßen, V. 21. Also sehen wir, wie die Vernunft pfleget zu urtheilen, wenns übel gehet, Noth und Kümmerniß für der Hand und Gottes Wort hinweg ist. Da denkt man alsobald, Gott wolle unser nicht und sei zornig über uns. So doch Gott alles leiden darum zuschicket, daß der alte Adam getödtet, der Glaube, Gehorsam und Geduld in uns probiret und gemehret, und wir also endlich selig werden.

\section{Votum}

Helfe Gott, von dem alles Vermögen kommt, Gutes zu thun, daß wir unter dem Kreuz unsere Seelen mit Geduld fassen; und verleihe, daß wir darunter allzeit als wahre Christen erfunden werden, fromm sein und auch an Leib und Seele, Amt, Nahrung, Glück und Gedeihen haben, Amen.

\chapter[31]{31}

\section{Vorrede}

Geliebte Christen! Nach dem Glauben, der die fürnehmste Tugend ist eines wahren Christen, soll ein jeder frommer Hausvater fürnehmlich sechserlei Tugenden an sich haben. Erstlich soll er sein ein keuscher Ehemann, der ihm an seinem Ehegemahl genügen läßt. Zum andern soll er sein ein Liebhaber der Wahrheit. Denn er soll nicht mit Lügen und Trügen umgehen, daß er seinen zeitlichen Nutzen schaffe, sondern er soll aufrichtig und redlich in seinem Beruf handeln. Zum dritten soll er gerecht sein gegen sein Hausgesinde, und dem erzeigen, was recht und billig ist. Zum vierten soll er gegen den Armen wohlthätig sein, und den Wittwen und Waisen mit Rath und That dienen. Zum fünften soll ein Hausvater demüthig sein in Glück und Wohlfahrt. Zum sechsten und letzten soll er gern beherbergen die Fremdlinge und zur Unschuld verjagte Leute. Das ist aller christlichen Hausväter Haustafel, darnach sie sich halten sollen. Gott, der es ihnen befohlen, zum Gehorsam: Dem heiligen Evangelio, dem sie folgen sollen; zu Ehren: Ihrem Nächsten zu Nutz und ihnen selbst zu gutem Gewissen, Nutz und Wohlfahrt, durch Gottes milden Segen. Und ein solcher Hausvater ist gewesen Hiob, von dem das 31. Kapitel.

\section{Hiobs unsträflicher Wandel}
\quellseite{753}

a Das ist, wenn mirs glücklich ginge, habe ich nicht meine Freude darinnen gehabt.

b Hand küssen: heißt sein eigen Werth preisen, welches allein Gott zugehöret.

c Das ist, mein Gesinde mußte auch nichts begehren an meine Feinde.

d Fürst) frei, unerschrocken.
\quellseite{754}

\section{Summarium}

Ist eine feine Lehre, in welcher Hiob angezeiget, daß er sein Lebenlang in Gottesfurcht und gutem Gewissen geblieben sei, B. 1. Er sei kein Hurer noch Ehebrecher gewesen; habe mit seinen Unterthanen nicht tyrannisch gefahren, B. 5. 13. Sei nicht geizig noch unbarmherzig gewesen, habe mit den Waisen nicht untreulich gehandelt, B. 16. Habe sich Geldes und Gutes nicht übernommen, B. 24. Sei keines Wanckels halben nicht hoffärtig gewesen, B. 25., habe sich ander Leute Unglück nicht gefreuet, sei gastfrei gewesen, habe niemals weder Gunst noch Ungunst angesehen in Händeln, und seinem Gesinde treulich und recht gelohnet, B. 32. 39. Solch Exempel sollen wir wohl bedenken und uns auch darnach halten. Sonderlich aber ist hie zu merken, daß er den Ehebruch ein solch Feuer heißt, welches bis ins Verderben verzehret und alles Einkommen auswurzelet, B. 12. Solches Kapitel sollten alle Hausväter in ihren Häusern haben, und täglich lesen, und ihr Leben darnach richten.

\section{Psalm}

\chapter[32]{32}

\section{Vorrede}

Bei diesem Kapitel sollen wir merken folgenden Unterricht. Die, so Gott mit Kreuz nach seinem gnädigen Willen heimsucht, soll man freilich ermahnen, warnen und strafen, wo sie in etwas, belangende Worte oder Werke, sich nicht recht erzeigen und zu viel thun; aber dabei soll man es nicht bleiben lassen, sondern man soll sie auch lehren und vermahnen, daß sie das Kreuz in Geduld tragen, weil es eine väterliche Züchtigung ist unser Sünde halben, und daß sie darunter mit gläubigem und fleißigem Gebet um Hülfe und Erlösung bei Gott, der ein Meister zu helfen ist, anhalten sollen. Denn soll man die lieben Kreuzträger trösten mit Gottes Wort, wenn sie sonderlich betrübt und kleinmüthig sind, weil das Kreuz so lange währet und ihnen oft grausame Leibes-Schmerzen anrichtet. Fürnehmlich aber soll man auch für die Kreuzträger herzlich beten, darum, daß es Gott kräftlich befohlen hat, und daß es ihnen sonderlich hoch vonnöthen ist, ihrer großen Leibes-Schmerzen und der Seele Anfechtungen halben, die oft zugleich zusammen schlagen und kommen. Derwegen soll man bitten, daß sie Gott in wahrem starken Glauben und christlicher Geduld erhalte, daß er ihnen zu seiner Zeit helfe, oder ihre Noth zu einem seligen Ende schicke. Zu solchem Gebet für alle, so unter dem Kreuz liegen, soll uns auch bewegen, daß sie ja unser Brüder oder Schwestern sein, und daß wir einen Vater haben, und ein Erbe alle zugleich im Himmel zu erwarten. Sollen uns derhalben anderer Christen Noth so fleißig, als unserer eigenen, annehmen. Neben dem soll aus dem Kapitel gemerket werden eine feine Kinderlehre, daß junge Leute mit Reden nicht behend sein, sondern den Alten weichen und zuhören sollen. Denn alte Leute in großen Sachen geübet sind, sind fürsichtig, wissen, was Nutz und Schaden bringen mag, und wegen ein Ding eher, denn daß sie es wagen. Darnach sind die Anfechtungen nicht so stark in ihnen, als in Jungen, darum sie auch bessern Rath geben können. Dieses lehret Sirach im 32. Kapitel B. 10. und sagt: Ein Jüngling mag auch wohl reden ein mal oder zwei, wenns ihm nöth ist, und wenn man ihn fragt, soll ers kurz machen, und sich halten, als der nicht viel wisse, und lieber schweigen, und soll sich nicht den Herren gleich achten, und wenn ein Alter redet, nicht drein maschen.

\section{Elihu zürnet über das Stillschweigen der Freunde Hiobs}
\quellseite{755}
* Ich ersticke sonst für großer Weisheit.

\section{Summarium}

Nachdem die drei Freunde Hiobs aufhören zu reden, darum, daß sie ihn von seiner Meinung nicht können bringen, tritt ein anderer auf, disputiret von neuem mit Hiob und strafet Hiob, daß er in seinen Leiden unrecht und wider Gott geredet habe, und läßt doch das Leben Hiobs ungetadelt, bleibet aber fürnehm(lich) darauf, Gott sei gerecht. Darum wer da achtet, es geschehe ihm unrecht, der sei bald im selbigen Unrecht, daß ers dafür hält, Gott thue ihm unrecht. Nimmt also für sich etliche Stücke, die er von Hiob gehöret und werth achtet, daß mans strafen soll.

Neben dem ist hie eine feine Kinderlehre, daß junge Leute mit Reden nicht behend sein, sondern den Alten weichen und zuhören sollen.

\section{Psalm}

\chapter[33]{33}

\section{Vorrede}

Geliebte Christen! Einen sichern ungebrochenen Menschen, der durch Kreuz und Noth nicht gedemüthiget ist, den kann der Teufel bringen und gebrauchen, wozu er will, auch zu den allerhäßlichsten und abscheulichsten Sünden und Lastern. Derwegen, daß die Menschen nicht sicher werden, und den Teufel sich verführen lassen, so legt ihnen Gott Kreuz auf, dadurch er sie in der Zucht hält und oft zu sündigen hindert, und, da er sie eine Zeitlang damit drücket und presset, so will er doch gnädig sein den Bußfertigen und endlich helfen. Derwegen soll ein jeder sein Kreuz in Geduld tragen, und soll in Gottes heiliger Gericht nicht reden, er nehme es gleich mit uns für, wie er wolle. Denn er niemand unrecht thut, als der da weiß, wer da recht oder unrecht hat. So ist er nicht allein weise, sondern die Weisheit selbst, nicht allein gerecht, sondern auch die Gerechtigkeit selbst. Dieses lehret auch das 33. Kapitel.

\section{Elihu vertheidiget Gottes Gerechtigkeit}
\quellseite{756}

a Bedenket) sicut homo post factum consuevit, paenitet, et cogitat mutare. Triumphator in Israel (inquit Samuel) non paenitet, nec mutat. Wie der Mensch nach der That überlegt, es bereut, und denkt zu verändern. Dem Held in Israel gereuet nicht (spricht Samuel), noch verändert er sich.

b Vorhaben) wie Abimelech, 1 Mos. 20, 3.

c Daß ihm's Gesicht vergehet, daß er weder siehet noch höret.

d Zwei oder drei mal) das ist, oftmals.

\section{Summarium}

Ist eine Anklage wider Hiob, daß er unrecht sei, darum, daß er über Gott geklagt habe, er thue ihm unrecht, V. 11. Denn Gott ist mehr denn ein Mensch; darum soll kein Mensch von Gott Ursache oder Rechenschaft fordern, warum er dies oder anders thut.

Darum setzet Elihu eine feine Lehre, daß Gott Kreuz und Leiden uns Menschen darum zuschickt, daß er uns vom Unglück wenden und unsre Seele vom Verderben und vor des Teufels Strick bewahren wolle, V. 16. Solches, spricht er, weiß ein Mensch von sich selbst in der Anfechtung nicht. Aber man lernets aus dem Wort Gottes, welches die Engel, das ist, die Prediger verkündigen. Dasselbige tröstet also, daß Gott wieder helfen und gnädig wolle sein, V. 23. Wer solches also weiß und für wahr hält, der fähet an zu beten und empfindet Trost und Freude, und bekennet, daß ihm Gott nicht zu viel gethan habe; sondern, daß er durch seine Sünde noch ärgers verdienet hätte.

\section{Votum}

Verleihe unser treuer lieber GOtt und Vater, daß alle, so das Kreuz tragen, Buße thun, geduldig sein und durch JEsum Christum, ihren Heiland, in wahrem Glauben der Hülfe sich trösten, unter dem Kreuz auch endlich Hülfe an Leib und Seele erlangen, Amen.

\chapter[34]{34}

\section{Vorrede}

Geliebte Christen! Die Ungeduld im Kreuz sucht das Herz oft vermessen an, und bewegt, daß es in Trübsal und Anfechtungen zu murren, zu fluchen, unleidlich zu werden und sich wider Gottes Willen aufzulehnen, und daß es darüber des Gebets und Danksagung vergisset; sonderlich wenn es siehet, wie es den Gottlosen so wohl gehet, und daß sie der betrübten und nothleidenden Gläubigen gleich spotten, als kenneten oder achteten Gott derselben nicht. Aber da ist zu wissen, daß die Ungeduld und das Murren unter dem Kreuz eine greuliche Sünde sei, und daß Gott die auch ernstlich strafe, als davon lehret die heilige Judith am 8. Kapitel B. 21. ihres Büchleins: Die andern, spricht sie, so die Trübsal nicht haben wollen annehmen, mit Gottesfurcht, sondern mit Ungeduld wider Gott gemurret und gelästert haben, sind von dem Verderber und durch die Schlangen umgebracht. Dieses wiederholet St. Paulus mit herzlicher Warnung
\quellseite{757}
und spricht 1 Cor. 10, 10.: Murret nicht, gleichwie jener etliche murreten, und wurden umgebracht durch den Verderber. Derwegen für Ungeduld unter dem Kreuz sich jedermann fleißig hüten soll. Dieses lehret auch das 34. Kapitel.

\section[Elihu schilt die Reden Hiobs, und fährt fort,]{Elihu schilt die Reden Hiobs, und fährt fort, GOttes Gerechtigkeit zu vertheidigen}

a Meinen Pfeilen) das sind Gottes Pfeile, die in mir stecken.

b Zu drängen) das ist, er lässet einen Tyrannen regieren, der das Volk mit Aufsehern und Schinden führet und quälet.
\quellseite{758}

\section{Summarium}

Ist die andere Anklage wider Hiob, daß er habe unrecht daran geredet, da er gesagt: Es gehe dem Frommen gleich wie dem Gottlosen; solches, spricht Elihu, ist eben so viel geredet, als wäre Gott auch gottlos und thäte unrecht. B. 12. Aber solches heißt Gott lästern; denn Gott ist ein Herr über alles, und könnte allwege Ursache genug wider die Menschen und alle Welt haben, daß ers in einen Haufen würfe; denn er siehet und weiß alles, warum gehet er auch also unter den Leuten um, und strafet die Gottlosen hin und wieder, wie sie es verdienen. Und du, Hiob, verdienst solche Strafe auch wohl mit deinem Lästern; und ist recht, daß Gott mit Strafe anhält, weil du von Murren und Lästern wider Gott nicht willst aufhören.

\section{Votum}

Gott, der ein Gott ist aller Geduld und Trostes, regiere alle betrübte Kreuzträger, daß sie ja unterm Kreuze nicht lästern, sondern seine heilige Gerichte lieben und loben, und in Geduld ferner Hülfe erwarten, durch Beistand und Kraft des heiligen Geistes, hochgelobet in Ewigkeit, Amen.

\chapter[35]{35}

\section{Vorrede}

Bei diesem Kapitel sollen wir merken zwei Stücke. Das erste ist, daß Gott der Gottlosen Gebet nicht erhöre, das ist, deren, die nicht ihre Sünde erkennen, lassen sie ihnen nicht leid sein, und gedenken nicht davon abzustehen, sondern vielmehr in Sünden fortzufahren. Davon sagt der Tert B. 12.: Sie werden schreien, und er wird sie nicht erhören, denn Gott wird das eitel nicht erhören. Also lehret Jesaias Kap. 59. B. 2.: Eure Untugend scheiden euch und euren Gott von einander, und eure Sünden verbergen das Angesicht von euch, daß ihr nicht gehöret werdet. Und Salomon in seinen Sprüchen im 28. Kapitel B. 9. sagt: Wer seine Ohren abwendet zu hören das Gesetz, des Gebet ist ein Greuel. Derwegen wer fruchtbarlich beten will, daß ihm und andern zu Nutz diene, der erkenne seine Sünde, und er trete ab von der Ungerechtigkeit, wie St. Paulus redet 2 Tim. 2, 19. Das ist, er bitte Gott zuvor, ehe denn er sein Gebet anfängt, um Vergebung derselben, so wird ihm widerfahren, davon Gott sagt 2 Chronika 7. Kapitel B. 14.: Wenn sie beten und mein Angesicht suchen, und sich von ihren bösen Wegen bekehren werden, so will ich vom Himmel hören. Das andere Stück ist, daß Gott endlich alles Uebel an allen Menschen strafe. Davon sagt alhie abermals Elihu B. 14. 15.: Es ist ein Gericht vor Gott, harre sein nur, ob sein Zorn so bald nicht heimsucht, und sich nicht annimmt, daß so viel Laster da sind. So bezeugets auch die Erfahrung, daß Gott den Verzug der Strafe pfleget mit der Schwerheit der Strafe einzubringen.

\section{Elihus fernere Rede von Gottes Gerechtigkeit}

a Gesänge) das ist, der Vögel Gesänge. Oder geistlich, daß man ihn lobet in Leid und Unfall, wie der 42. Psalm B. 9. auch sagt: Des Nachts singe ich ihm.
\quellseite{759}

\section{Summarium}

Dies ist die dritte Anklage wider Hiob, der Gott beschuldiget hat, daß er ihm seines gottseligen Wandels nicht genießen lasse. Solches, spricht Elihu, ist auch unrecht geredet. Ursache, wenn du gleich lange fromm bist, so hilft Gott deine Frömmigkeit nicht. Warum soll er dirs denn bezahlen. Die Leute hilft es, wenn du ihnen was Gutes thust. Wiederum, so du ihnen unrecht thust, so schadet Gott solches nicht; es schadet nur den Leuten, und Gott will solches Unrecht über sie gehen lassen und ihr Gebet nicht erhören, weil sie so eitel und gottlos sind, V. 8. Aber um solches willen soll man nicht sagen, daß Gott das Uebel nicht strafe, er hat sein Stündlein, daß er mit der Strafe kommen wird, dasselbige wird nicht außen bleiben, V. 15.

\section{Gebet}

Der allmächtige, barmherzige GOtt helfe uns, daß wir uns täglich zu ihm bekehren, und daß wir aus bußfertigen Herzen in wahrem Glauben unser Gebet thun, und lindere die wohlverdienten Strafen, um Jesu Christi, seines lieben Sohnes und unsers Herrn willen, Amen.

\chapter[36]{36}

\section{Vorrede}

Bei diesem Kapitel sollen wir lernen, daß Gott darum mit Kreuz, Leiden und Trübsal heimsucht, auf daß er dadurch den zukünftigen Sünden wehre. Denn wenn der Mensch ohne Kreuz und Widerwärtigkeit ist, und es ihm nach seinem Wunsch und Willen gehet, so wird er gemeiniglich fürwitzig, muthwillig und sicher, und gewinnet Lust und Liebe zu allerlei sündhaften Dingen. Daher kommen die Sprichwörter: Wenn dem Esel zu wohl ist, so gehet er aufs Eis tanzen. Und: Es kann ein Mensch nichts weniger erleiden, als gute Tage. Dergleichen dies: Es müssen starke Beine sein, die gute Tage ertragen können. Solches auch die Erfahrung und Exempel bezeugen. Die heilig und gottfürchtig lebte David, da ihn Saul, und darnach Absalon, sein Sohn, verfolgte. Aber da er Ruhe und Frieden hatte, fiel er und sündigte schändlich. Also that auch der König Hiskiä, da der Feind für der Stadt Jerusalem lag, war er ganz gottfürchtig und suchte bei dem Propheten Jesaia Trost und Rath. Dergleichen that er auch, da er todt krank lag; aber da sein Feind, der Sanherib, durch göttliche Macht weggetrieben und er wieder gesund worden war, geriet er in Ruhmredigkeit, welches auch seine Nachkömmlinge mußten entgelten. Daraus denn folget, daß der Mensch frömmer ist, wenn es ihm übel gehet, denn wenn es ihm wohlgehet. Derwegen so drückt uns Gott mit Kreuz, daß wir desto weniger sündigen sollen, und daß er fromme Kinder habe, die ihm dienen in Heiligkeit und Gerechtigkeit. Dies bezeugen auch St. Paulus, Röm. 6, 6. und saget: Dieser alter Mensch wird gekreuziget, auf daß der sündliche Leib aufhöre, daß wir hinfort der Sünde nicht dienen. Ja der Heyde Plinius jun. lehret in einer seiner Epistel, daß wir in Krankheiten (also unter anderm Kreuz) nicht so Lust zu sündhaften Dingen haben, als wenn wir frisch und gesund sein, und saget darauf: Wollte Gott, daß wir allezeit wären, wie wir sein, wenn wir krank liegen. Hieher gehört auch, daß St. Petrus 1. Epistel 4. Capitel V. 1. sagt: Wer am Fleisch leidet, der höret auf zu sündigen, daß er hinfort, was noch hinterstelliger Zeit im Fleisch ist, nicht der Menschen Lüsten, sondern dem Willen Gottes lebe.

\section{Ferner Beweis der Gerechtigkeit, Allmacht und Weisheit Gottes}
\quellseite{760}

a Id est, ab occidente in orientem. D. h. von West nach Ost.

\section{Summarium}

Nachdem Elihu an Hiob gestraft hat, daß er unrecht geredet und Gott gelästert habe, fähret er hie zum Beschluß eine neue Predigt an, von Gott, von seinen Werken und Gericht, B. 5., von seiner Gewalt und Allmächtigkeit, auf daß Hiob für solchem großen Gott lerne sich demüthigen, und ihn fürchten, und nicht weiter wider ihn rede.

Und ist sonderlich in diesem Kapitel zu merken, daß Elihu sagt: Gott sei gerecht, und thue niemand unrecht, B. 7. Daß er aber mancherlei Kreuz und Anfechtung den Leuten zuschicke, geschieht darum, daß sie ihre Untugend erkennen lernen, sich züchtigen lassen und von Unrecht sich bekehren. Denn wo uns Gott nicht durch mancherlei Anfechtung heimsucht, würden wir je länger, je ärger.

\section{Votum}

Wir danken GOtt herzlich, daß er durch das Kreuz unser Bestes und Nutzen suchet; der helfe allen lieben Kreuzträgern, und regiere sie, daß sie ihr Kreuz in christlicher Geduld tragen, und daß sie auch nach der Strafe frömmer werden, und beständig fromm bleiben und selig werden, Amen.

\chapter[37]{37}

\section{Vorrede}

Bei diesem Kapitel sollen wir merken folgende Lehre: Wenn der Mensch gleich kein Wort Gottes gehöret hätte, und nur das schöne herrliche Gebau des Himmels und der Erden, sammt allem, was darinnen ist und geschieht, fleißig anschauet, so kann ers nicht lassen, er muß sich darum bekümmern, wo es doch alles herkomme, es werde ja nicht ohngefähr also worden sein, es müsse dennoch irgend woher seinen Anfang haben; und schließt, daß etwas sei, nämlich Gott, der solches alles gemacht habe. Wenn der Mensch darnach weiter siehet, daß Himmel und Erden so lange Zeit, so fest und unverrücklich bestanden, Sterne, Sonne und Mond ihren Lauf so ordentlich unverhindert halten, daß Jahr seine unterschiedliche Ordnung und Abwechselung des Sommers und Winters, Herbsts und Lenzes hält, und daß eine jede Zeit ihre Eigenschaft mit sich bringet, als Regen, Schnee, Donner und Wind; und daß die Creaturen, Vögel, Fische, Thiere und Menschen, ein jedes in seinem Geschlecht und seiner Art sich vermehren und zusammen halten: Daß, zu dem auch die Gewächse, Bäume, Kräuter und andere Früchte, sich jährlich besamen; und daß eine solch schöne Ordnung allenthalben beständiglich, und solches alles dem Menschen zu gut, gehalten wird: So muß der Menschen Herz schließen, daß der, so dieses alles gemacht und erhält, und so lange Zeit so ordentlich regieret, ob er gleich unsichtbar ist, dennoch wahrhaftiger Gott und ewiger, allmächtiger, weiser, gütiger, treuer Gott sein müsse. Solche Erkenntniß Gottes findet man in dem großen Buch, welches heißt die geschaffene Welt, welches Buch so manches Blatt hat, so manche Creatur in der Welt ist, und steht an einem jeglichen Blatt geschrieben: Gott ist ewig und allmächtig, weise und gütig, gnädig und barmherzig; und sind alle Menschen darum erschaffen, und haben darum von Gott ihre Augen empfangen, daß sie in diesem schönen Buch lesen, und an solchem Werk den Schöpfer aller Creaturen erkennen, und dafür auch ihn ehren, ihm
\quellseite{761}
danken und nach seinem Willen gehorsamlich leben, und sich seiner Allmacht und Güte im Leben und Sterben freuen und trösten sollen. Derwegen, da St. Paulus zu Athen das Evangelium predigte, lehret er auch, daß man Gott in den Kreaturen suchen solle. In der Apost. 17, 27. und in der Epistel an die Römer im 1. Kapitel V. 20. schreibt er: Wie Gott diejenigen greulich gestraft, so ihn bei seinen Kreaturen und Werken nicht erkannt haben.

\section{Gottes Majestät wird aus dem Buch der Natur erkannt}

* Gold) das ist, helle Wetter, wie lauter Gold.

\section{Summarium}

Dies ist der andere Theil der Beschlußpredigt Elihu, in welcher er dem armen Hiob erzählet die herrlichen Werke Gottes, die wir alle Tage sehen, als schön Wetter, Donner, Schnee, Regen, Wind, Wolken und dergleichen, B. 2. Weil nun Hiob bekennet und sagen muß, daß solche Werke höher und größer sind, denn aller Menschen Vermögen können auswege bringen, will er ihn also vermahnet haben, daß er wider solchen hohen, großen, gewaltigen Gott nicht freventlich reden, und sein Werk, das er mit ihm fürhabe, nicht tadeln noch ausforschen soll, sondern sich für ihm demüthigen und ihm sein Werk und Willen gefallen lassen, und geduldig sein, so thut er recht.

\section{Votum}

Lob, Ehr, Preis und Dank sei GOtt für die Lehre von den Werken seiner Allmacht, Weisheit und Gütigkeit; der gebe seine Gnade, daß wir diese täglich anschauen, ihm ferner allezeit dafür danken, und dieselben wohl und christlich gebrauchen, Amen.

\chapter[38]{38}

\section{Vorrede}

Geliebte Christen! Der heilige Prophet Jeremias sagt in Klagliedern im 3. Kapitel V. 25.: Der Herr ist freundlich dem, der auf ihn harret, und der Seelen, die nach ihm fraget. Es ist ein köstlich Ding, geduldig sein, und auf die Hülfe des Herrn hoffen. Also spricht auch der Herr Christus Lucä 21, 18.: Ein Haar von eurem Haupt soll nicht umkommen. Fasset eure Seelen mit Geduld. Und St. Paulus sagt in der Epistel an die Ebräer am 10. Kapitel V. 36.: Geduld ist euch noth, auf daß ihr den Willen Gottes thut, und die Verheißung empfahet.
\quellseite{762}
Diese Sprüche lehren, daß wir unter dem Kreuz, als gottesfürchtige und gehorsame Kreuzträger, geduldig sein sollen. Damit, wenn wir auf den Zaum bringen wollten, Gott nicht verursachen, daß er uns ein schärfer Gebiß einlegen, und uns den Zaum recht machen und abrichten, oder gar verstoßen, und da wir uns nicht wollten abrichten lassen, in den Karren des höllischen Feuers anspannen müsse. Dies betrachten wir auch bei dem 38. Kapitel.

\section{Gott preiset selbst seine Allmacht und Weisheit}

a Das Siegel) das ist, ihr Stand und Wesen, daß sie gewiß sein wollen, abgebildet.

b Verborgene) das ist, ins Herz.

\section{Summarium}

Hie gehet eine andere Predigt an, welche Gott selbst dem Hiob thut, darin er alle Werke seiner göttlichen Kraft erzählet. Will dadurch den
\quellseite{763}
Hiob lernen, er habe unrecht gethan, daß er solchen großen Gott zur Rechenschaft gefordert, und wider sein Werk und Willen also gemurret habe. Uns zur Lehre, daß wir im Kreuz und Leiden geduldig sollen sein, und wider solchen Gott nicht murren, von dem wir alles haben, und ohne welchen nichts sein noch bestehen kann, beide im Himmel und auf Erden. Darum soll man solchen Gott fürchten, sein Wort und Befehl in Ehren haben und demselbigen folgen, und willig und gerne leiden, was er uns zuschickt.

\section{Votum}

Helffe GOtt, daß wir dieser nöthigen Lehre treulich folgen, und in allem Kreuz und Leiden geduldig, auch sonst zu unserer Besserung GOttes Wort gehorsam seyn, in Kraft des heiligen Geistes, um Jesu Christi, seines Sohnes und unsers Heilandes willen. Amen.

\chapter[39]{39}

\section{Vorrede}

Geliebte Christen! Im Buch der Richter im 2. Kapitel V. 1. lesen wir, daß ein Engel des HErrn, das ist, der Priester Pinehas gen Bochim kommen sei, und habe an den Leuten allda ihre Sünden durch Gottes Wort gestraft, und ihnen angekündigt, wo sie nicht in Zeiten Buße thun würden, daß sie Gott ernstlich und hart strafen wolle. Was hat aber diese Bußpredigt ausgerichtet? Davon meldet die Historia, und sagt V. 4.: Da der Engel des HErrn solche Worte geredet hatte zu allen Kindern Israel, hub das Volk seine Stimme auf und weinete, und opferten daselbst dem HErrn. Also soll man Gottes Wort hören, daß man es auch annehme, in einem feinen reinen Herzen bewahre, daraus sich bessere und dem Evangelio gehorsam sey. Dieses lehret auch David im 32. Psalm V. 8. und sagt: Ich will dich unterweisen und dir den Weg zeigen, den du wandeln sollst. Seyd nicht wie Rosse und Mauler, die nicht verständig sind. Das ist, man soll sich Gottes Wort lassen regieren, Gott fürchten und nach Gottes Wort ein gottselig Leben führen. Davon lehret auch Hiob mit seinem Exempel, in dem 39. Kapitel.

\section[An den Thieren leuchtet GOttes unerschöpfliche Weisheit hervor.]{An den Thieren leuchtet GOttes unerschöpfliche Weisheit hervor. Hiob bekennet seinen Fehler}

a) Preis) das ist, es ist nur desto trotziger und muthiger, und schnaubet, als rühmte sich, wo schreckliche Dinge, als Streit und Krieg vorhanden ist.
\quellseite{764}

b) Achtet nicht) das ist, es thut, als sei ihm nichts, das doch so schrecklich ist.

\section{Summarium}

Dies ist eine Predigt, in welcher Gott vieler Thiere gedenket, wie er sie erschaffen und einem jeden seine sonderliche Art und Eigenschaft gegeben habe, V. 1. Auf daß Hiob sich für ihm demüthige, und nicht beharre auf seinem Murren wider Gott, weil er solche Werke Gottes nicht allein nicht nachthun, sondern auch nicht verstehen kann. Solche Predigt nimmt Hiob an, bekennet seine Sünde und sagt: Er wolle es nimmer thun, V. 35.

\section{Votum}

Wir danken GOtt herzlich, daß er uns auch durch sein Wort von seinen Werken und Willen prediget; der verleihe, daß wir dies sein heiliges Wort also anhören, daß wir unsere Sünde erkennen, davon abstehen, und gegen ihm gottselig, gegen uns züchtig und gegen unsern Nächsten gerecht leben, und in solchen Tugenden selig sterben, Amen.

\chapter[40]{40}

\section{Vorrede}

Geliebte Christen! Der heilige Prophet Daniel sagt im 9. Kapitel V. 7.: Du HErr bist gerecht, wir aber müssen uns schämen. Also David im 119. Psalm V. 75. 137. spricht: HErr, du bist gerecht, und deine Gerichte sind recht. Dadurch werden wir gelehret, daß wir unter dem Kreuz gerne bekennen sollen, daß uns Gott nicht unrecht thue, sondern mit uns handele, wie wir es mit unsern Sünden wohl verdienet haben. Ferner, St. Petrus sagt 1. Epistel am 3. Kapitel V. 13. 15.: Wer ist, der euch schaden könnte, so ihr dem Guten nachkommet. Fürchtet euch für ihrem Trotzen nicht, und erschrecket nicht. Und unser HErr JEsus Christus spricht Matth. 10, 29.: Kauft man nicht zween Sperlinge um einen Pfennig? Noch fällt derselbigen keiner auf die Erde ohn euern Vater. Nun aber sind auch eure Haare auf dem Haupt alle gezählet. Darum fürchtet euch nicht; ihr seid besser, denn viel Sperlinge. Hiermit werden wir gelehret, daß uns der Teufel und die böse Welt nicht mehr schaden können, denn das ihnen Gott verhänget und zulässet. Von beiden prediget auch Gott unser HErr im 40. Kapitel Hiobs.

\section{Spiegel göttlicher Allmacht an dem Behemoth}

a. Behemoth) heißen alle großen ungeheuren Thiere. Wie Leviathan alle großen ungeheuren Fische. Aber darunter beschreibt er die Gewalt und Macht des Teufels und seines Gehülffes, des gottlosen Hausvaters in der Welt.
\quellseite{765}

\section{Summarium}

Hie ist eine harte Predigt wider Hiob. Denn hie sagt der Herr, was ihm an Hiob nicht gefallen hat. Du willst, spricht er, gerecht sein, und läßt dich dünken, ich thue dir unrecht; aber hie versuche dich, ob du die Werke kannst thun, welche ich thue, V. 2. Befindest du nun, daß du sie kannst thun, alsdann so nimm mich mit Recht für, und heiße dir Rechenschaft geben.

Daß darnach stehet von Behemoth, daß er so stark und groß sei und das ganze Erdreich ihm Speise trage, und er ganze Wasser auslaufe, V. 10.; ist zu verstehen insgemein von allen wilden, starken, ungeheuren Thieren auf Erden. Gleichwohl wird durch solchen Behemoth des Teufels und der Gottlosen Gewalt in der Welt heimlich angezeigt, daß dieselben alle ihre Gewalt nicht von ihnen selbst, sondern von Gott haben, so lange ers verhänget und duldet.

\section{Datum}

Helfe GOtt, daß wir allezeit seine Gerichte und Werke lieben und loben, derwegen auch uns unter seine gewaltige Hand demüthigen; und steure dem Teufel und der bösen Welt, daß sie uns mit ihrer List und Gewalt an unserm Leib, Seele und Leben nicht schaden, Amen.

\chapter[41]{41}

\section{Vorrede}

Geliebte Christen! Obwohl der Satan ein giftiger, listiger, mächtiger Geist ist, so kann er uns doch nichts thun ohne Gottes, des Herrn, Erlaubniß und Willen; so haben wir ja auch auf unserer Seite den großen, starken Held und Sieges-Fürsten Michael, unsern Herrn Christum, Dan. 12, 1. Der die Werke des Teufels zerstören will und kann. So bewahren uns auch die lieben Schutzengel, die heiligen und starken Engel, wenn wir in Gottesfurcht unsers ordentlichen Berufs warten, Ps. 34, 8. und 91, 11. So haben wir auch das liebe und kräftige Vater unser, und das heilige Geistes Schwert, Gottes Wort, Eph. 6, 17. Damit wir uns schützen und wehren sollen und können. Dieses sollen wir uns mit Freuden trösten wider den Teufel. Ob wir schon gegen ihm sind wie eine matte Fliege gegen einen starken Löwen, und sollen derwegen täglich in Gottesfurcht leben. Denn die Teufel fliehen für einem gottesfürchtigen Menschen, wie ein Mensch fleucht für einem grausamen Löwen, der Feuer ausspeiet, sagt der alte Irenäus. Dieß erinnert auch das 41. Kapitel.

\section{Beschreibung von des Leviathans Größe, Macht und Stärke}

a Leviathan) nennet er die großen Wasserfische im Meer, doch darunter beschreibt er der Welt Fürsten, den Teufel mit seinem Anhang.
\quellseite{766}
ᵇ Starken) das ist, die großen Fische fliehen für ihm; also auch für der Welt Gewalt fliehen die Mächtigen.

ᶜ Keine Gnade) so haben sie gesündiget, das ist, sie müssen gethan haben, und behalten, als arme Sünder.

ᵈ Weg) das ist, er schwimmet und lebet im Meer, wie er will, daß man seinen Weg von fern siehet.

\section{Summarium}

Ist der letzte Theil von der Predigt des HErrn, wider den armen Hiob, in welcher er, wie vor, dem Hiob anzeiget, er soll wider Gott nicht streiten, sondern sich leiden und geduldig sein. Den Leviathan nennt er die großen Wallfische im Meer. Aber darunter will er uns den Teufel und sein Reich fürmalen, wie gewaltig und groß es sei, daß keine menschliche Hülfe ihm kann widerstehen, aber Gott könne es und wolle es thun. Will also dem Hiob anzeigen, daß solch Leiden und Jammer der Teufel aus Verhängniß Gottes ihm habe zugerichtet. Darum soll er weder mit seiner Unschuld noch Frömmigkeit wider den Teufel fechten. Denn damit werde er nichts ausrichten. Er soll sich aber an Gott halten, und da um Hülfe und Gnade bitten, so werde der Leviathan weichen müssen.

\section{Votum}

Der HErr Zebaoth, unser HErr JEsus Christus, regiere uns, daß wir, gottesfürchtig leben; und helfe uns, daß wir mit unserm gläubigen Gebet täglich wider den Teufel streiten, wolle auch unser Gebet erhören und uns behüten für Schaden und Verderben Leibes und der Seelen, Amen.

\chapter[42]{42}

\section{Vorrede}

Geliebte Christen! Gott kann seinen Gläubigen, wenn sie gleich hinter sich kommen sind, wiederum beifür helfen, wenn sie sich nur recht nach Gottes Wort in den Handel schicken. Derwegen, wenn einer um sein Hab und Gut kömmt, durch allerlei und mancherlei Unfall, oder fällt in Verachtung oder schwere Anfechtung, so läßt es Gott nicht darum geschehen, daß er ihn gar ausschütten wolle, sondern etwa darum, daß er ihm das reichlicher wiederum gebe, daß er verloren hat. Wird aber gleich einem das Seine nicht zwiefach, so wird ihm doch oft so viel, als er vor hat, oder Gott ersetzt es in andern Wegen, daß er ihm Herz und Muth gibt, daß er ihm an wenigen genügen läßt. So übertrifft weit der geistliche Segen, den Gott den Seinen gibt, den leiblichen. So wird uns Gott auch unseres Leides erst recht ergötzen im künftigen Leben. In Summa, Gott hat niemand jemals verlassen, der ihm vertraut und sich gänzlich an ihn ergeben, und in der Noth aus wahrem Glauben zu ihm geschrieen, und ein gottesfürchtig Leben geführet hat. Davon höret auch an das 42. und letzte Capitel des Buchs Hiobs.

\section[Hiob wird auf rechte Buße von GOtt wiederum]{Hiob wird auf rechte Buße von GOtt wiederum glücklich gemacht}
\quellseite{767}

\section{Summarium}

Ist erstlich ein schön Exempel einer rechtschaffenen Buße, daß Hiob des HErrn Wort annimmt, seine Sünde erkennet und lässet ihm leid sein und bittet um Gnade, B. 2. Das aber der HErr Eliphas und die andern zwei strafet, und sagt: Hiob habe recht geredet, B. 7., gehet dahin, daß es wahr ist, Hiob habe solche Strafe nicht verdient. Denn der HErr sagt selbst: Er habe seine Hand ohne Ursache wider ihn ausgestrecket. Darum haben die Freunde Hiobs unrecht, daß sie sagen: Er müsse seiner Sünden halben leiden. Wiederum aber hat Hiob auch in dem unrecht, daß er wider Gott murret, und um solcher Strafe willen Gott für unrecht anklaget.

Zum andern ist hie ein sonderlich fein Exempel, wie Gott der Frommen Fürbitte ansehen und derselbigen andere wolle genießen lassen, B. 10.

\section{Votum}

Also hat eure Liebe die Historien vom frommen Hiob, in welcher wir fürnehmlich lernen sollen, daß GOtt auch die Frommen unter das Kreuz wirft, und der alte Adam GOttes Willen nicht kann geduldig tragen, er murret, und meinet, es geschehe ihm Unrecht. Aber tröstlich ists, daß GOtt solche Ungeduld vergeben, und aus aller Noth und Anfechtung endlich helfen, und uns alles reichlich ergänzen will. Ist derhalben eine nütze Lection für die, so in großem Unglück und Leiden stecken, daß sie hie lernen, wie sie sich für GOtt demüthigen und allein auf seine Gnade sehen, und keines Verdienstes sich rühmen sollen, so will er gnädig helfen, nicht allein hie zeitlich, sondern auch ewig. Das helfe uns unser lieber HErr Christus, Amen.
% Scan page 768: Blank family-register form excluded
% Scan page 769: No literary apparatus retained
% Scan page 770: Blank family-register form excluded
% Scan page 771: Blank family-register form excluded; Google scan watermark excluded
% Scan page 772: Blank family-register form excluded; Google scan watermark excluded
% Scan page 773: Blank family-register form excluded; Google scan watermark excluded
% Scan page 774: No literary apparatus retained
% Scan page 775: No literary apparatus retained
% Scan page 776: No literary apparatus retained
% Scan page 777: Google scan watermark excluded
% Scan page 778: Blank/non-literary scan matter excluded; Library ownership stamp matter excluded; Google scan watermark excluded
% Scan page 779: No literary apparatus retained
% Scan page 780: Library ownership stamp matter excluded; Google scan watermark excluded
% Scan page 781: Google scan watermark excluded


\end{document}
